\frontmatter

\quellseite{6}
\begin{titelseite}
{\LARGE Gebets-Schatz.}
\end{titelseite}

\quellseite{8}
\begin{titelseite}
{\LARGE Evangelisch-Lutherischer}\\[1em]
{\LARGE Gebets-Schatz.}\\[1em]
Vollständige Sammlung von Gebeten Dr. Martin Luthers und anderer rechtgläubiger, geschätzter Väter der ev.-luth. Kirche in unverändertem Abdruck.\\[1em]
Nebst einem Hausgesangbüchlein, Hundert und sechs alte rechtgläubige Lieder für den Hausbedarf enthaltend.\\[1em]
Herausgegeben von der deutschen ev.-luth. Synode von Missouri, Ohio u. a. Staaten.\\[1em]
St. Louis, Mo.\\[1em]
Zu haben bei M. C. Barthel, General-Agenten der evang.-luth. Synode.\\[1em]
1881.\\[1em]
15te Auflage.
\end{titelseite}

\quellseite{10}

\chapter{Vorerinnerungen}

Zwar gibt es bereits eine so große Menge Gebetbücher, daß eine Vermehrung derselben mit dem gegenwärtigen Manchem überflüssig zu sein scheinen mag. Allein theils sind die bereits vorhandenen durchaus elende Machwerke, theils enthalten sie neben Gutem auch Schlechtes, theils sind sie nicht vollständig. Treue Seelsorger sind daher, wenn die ihnen Anvertrauten nach einem guten Gebetbuche für allerlei Fälle fragten, bisher immer in Verlegenheit gewesen, welches Buch sie mit gutem Gewissen empfehlen sollten. Die Sammler des gegenwärtigen „Gebets-Schatzes“ hoffen durch denselben somit einem immer dringender gewordenen Bedürfnisse entsprochen zu haben. Die vorliegende Sammlung, nicht für den kirchlichen, sondern für den häuslichen Gebrauch bestimmt, ist nicht nur eine möglichst vollständige, sondern es sind auch in dieselbe nur solche Gebete aufgenommen, welche sowohl rein im Bekenntniß der Lehre des göttlichen Wortes sind, als auch die Sprache einer wahren, weder naturalistischen, noch krankhaften und schwärmerischen, Herzensandacht enthalten, so daß unser „Gebets-Schatz“ in dieser Beziehung nicht nur die neueren, sondern selbst die meisten älteren vollständigen Sammlungen übertreffen dürfte.

In Betreff des als Anhang beigegebenen „Hausgesangbüchleins“ haben wir nur so viel zu erinnern, daß darin nur solche Lieder aufgenommen worden sind, welche sich nicht schon in
\quellseite{11}
dem Kirchen-Gesangbuch der deutschen evang.-luth. Synode von Missouri, Ohio und anderen Staaten finden; daher unser „Hausgesangbüchlein“ als zweiter, für das Haus berechneter Theil des letzteren angesehen werden kann.

Die beigegebenen Register, namentlich das am Schlusse des Buches befindliche alphabetische Register, werden hoffentlich den Lesern willkommen sein und zur schnellen Auffindung eines passenden Gebetes in allerhand Fällen gute Dienste leisten.

Möge denn dieser „Gebets-Schatz“ recht vielen Heilsbegierigen ein wirklicher Schatz für ihre Seele werden und die Beter, welche sich desselben bedienen, ihn so gebrauchen, daß sie die darin enthaltenen Gebete nicht nur mit Mund und Herzen nachsprechen und Gott im Namen Jesu Christi im wahren Glauben vortragen, sondern auch daraus lernen und entzündet werden, ihr Herz betend, anrufend, lobend und dankend auch mit eigenen Worten vor dem Thron der Gnade auszuschütten. Damit dies geschehe, lassen wir Luthers, dieses gesalbten Beters, Anweisung zu einem rechten erhörlichen Gebete folgen, welche wir allen, die diesen „Gebets-Schatz“ gebrauchen wollen, zu fleißiger Durchlesung und Beherzigung hierdurch dringend empfehlen. Die Anweisung ist der Kirchenpostille Luthers entnommen.

St. Louis, Mo., im Mai 1864.

Die deutsche ev.-luth. Pastoralconferenz zu St. Louis, Mo.

\quellseite{12}

\chapter{Luthers Anweisung zu einem rechten Gebet}

Daß ein Gebet recht gut sei, und erhöret werde, ist erstlich zu merken, daß dazu fünf Dinge nöthig sind. Das erste: daß man von Gott eine Verheißung oder Zusage habe, und dieselbige zuvor bedenke, Gott derselbigen erinnere, und sich dadurch bewege, tröstlich zu bitten. Denn so Gott nicht hätte heißen bitten und Erhörung zugesagt, möchten alle Creaturen mit allem Bitten nicht ein Körnlein erlangen; daraus denn folget, daß niemand etwas von Gott erlanget, seiner oder seines Gebets Würdigkeit halben, sondern allein aus göttlicher Gütigkeit, der allem Bitten und Begierden zuvorkommt, durch sein gnädig Zusagen und Heißen uns bewege zu bitten und begehren, auf daß wir lernen, wie gar viel mehr er für uns sorget, und mehr bereit ist zu geben, denn wir zu nehmen und suchen, und also kühn werden, tröstlich zu bitten, sintemal er alles und mehr darbeut, denn wir bitten mögen.

Zum Andern ist noth, daß man je nicht zweifle an der Zusagung des wahrhaftigen und getreuen Gottes; denn eben darum hat er Erhörung zugesagt, ja zu bitten befohlen, daß man je gewissen und festen Glauben habe, es werde erhöret, als er sagt (Matth. 21. und Marc. 11.): „alles, was ihr bittet im Gebet, glaubet nur, so werdet ihr es empfangen.“ Und Luc. 11. spricht er: „ich sage euch auch: bittet, so wird euch gegeben, suchet, so werdet ihr finden, klopfet an, so wird euch aufgethan; denn wer da bittet, der nimmt, und wer da suchet, der findet, und wer da anklopfet, dem wird aufgethan. Wo bittet unter euch ein Sohn den Vater ums Brod, der ihm einen Stein dafür biete? Und so er um einen Fisch bittet, der ihm eine Schlange für den Fisch biete? oder so er um ein Ei bittet, der ihm einen Scorpion dafür biete? So denn ihr, die ihr arg seid, könnt euren Kindern gute Gaben geben, viel mehr wird der Vater im Himmel den Heiligen Geist geben denen
\quellseite{13}
die ihn bitten.“ Auf diese und dergleichen Zusagung und Befehl muß man sich tröstlich erwägen, und mit rechtem Vertrauen bitten.

Zum Dritten: so jemand also bittet, daß er zweifelt an der Erhörung Gottes, und allein dahinsiehet auf Abenteuer sein Gebet, es geschehe oder geschehe nicht, — der thut zwei böse Stücke: das erste, daß er sein Gebet selbst zu nichte macht, und umsonst arbeitet. Denn also spricht Jacobus (Kap. 1.): „wer von Gott bitten will, der bitte im Glauben, und zweifle nicht; denn wer da zweifelt, der ist gleich wie eine Woge des Meers, die vom Winde getrieben und gewehet wird; solcher Mensch gedenke nur nicht, daß er etwas von dem HErrn empfangen werde.“ Das meinet er also, daß desselbigen Menschen Herz nicht stille hält, darum kann ihm Gott nichts geben; der Glaube aber hält das Herz stille, und machet es empfänglich göttlicher Gaben.

Das andere böse Stück ist: daß er seinen allergetreuesten und wahrhaftigen Gott für einen Lügner und einen losen, ungewissen Mann achtet, als einer, der seinen Zusagen nicht möge oder wolle genugthun, und also durch seinen Zweifel Gott die Ehre und den Namen der Treue und Wahrheit beraubet, darinnen so schwerlich gesündiget wird, daß auch durch die selbige Sünde ein Christ zu einem Heiden wird, und seinen eigenen Gott verleugnet und verlieret, und, so er darin bleibet, ewiglich ohne allen Trost muß verdammt sein. Wird ihm aber etwas gegeben, daß er bittet, so wird es ihm nicht zur Seligkeit, sondern zu Schaden gegeben, zeitlich und ewiglich; auch nicht ums Gebets willen, sondern aus dem Zorn Gottes, daß er belohne die guten Werke, die da in Sünden, Unglauben und göttlicher Unehre geschehen werden.

Zum Viersten sprechen etliche: ja ich wollt’ wohl vertrauen, mein Gebet würde erhöret, wenn ichs würdig wäre, und wohl machte. Antworte ich: wenn du nicht ehe bitten willst, du wissest denn oder empfindest dich würdig und geschickt, — so mußt du nimmermehr bitten; denn wie zuvor gesagt ist: unser Gebet muß sich
\quellseite{14}
nicht gründen oder halten auf unser oder seine Würdigkeit, sondern auf die unwandelbare Wahrheit göttlicher Zusagung; und wo sich's auf sich selbst oder etwas anders gründet, — so ist es falsch und betrüget dich, ob auch für großer Andacht das Herz zerbräche, und eitel Blutstropfen weinete; denn darum bitten wir, daß wir nicht würdig sind zu bitten, und eben dadurch werden wir würdig zu bitten und erhört zu werden, daß wir glauben, wir sind unwürdig, und allein auf die Treue Gottes uns tröstlich wagen. Sei wie unwürdig du magst, so schaue darauf, und nimm mit ganzem Ernst wahr, daß tausendmal mehr daran gelegen ist, daß du Gottes Wahrheit ehrest, und nicht in deinem Zweifel seine treue Zusagung zur Lüge machest; denn deine Würdigkeit hilft dich nicht, deine Unwürdigkeit hindert dich nicht: aber das Mißtrauen verdammt dich, und die Zuversicht würdiget und behält dich.

Darum hüte dich dein Leben lang, daß du dich nimmer würdig oder geschickt achtest zu bitten oder zu empfahen, es sei denn, daß du dich findest einen freien Wagehals auf das wahrhaftige und gewisse Zusagen deines gnädigen Gottes, der dir seine Barmherzigkeit und Güte also will offenbaren, daß zugleich, wie er dir Unwürdigen, Unverdienten, aus lauter Gnaden, ungebeten, Erhörung zugesagt hat, also will er auch dich unwürdigen Beter aus lauter Gnaden, zu Ehren seiner Wahrheit und Zusagung, erhören, auf daß du nicht deiner Würdigkeit, sondern seiner Wahrheit, damit er die Zusagung erfüllet, und seiner Barmherzigkeit, damit er die Zusagung gethan hat, danksagest, auf daß bestehe der Spruch im 25. Psalm: „die Wege des HErrn sind Güte und Treue denen, die seinen Bund und Zeugniß halten.“ Güte oder Barmherzigkeit in der Zusagung, Treue oder Wahrheit in Erfüllung oder Erhörung der Zusagungen. Und im 85. Psalm spricht er: „Güte und Treue begegnen einander, Gerechtigkeit und Friede küssen sich.“ Das ist: sie kommen zusammen in einem jeglichen Werk und Gaben, die wir von Gott durch Bitten erlangen.

\quellseite{15}

Zum Fünften soll man sich in diesem Vertrauen also halten, daß man Gott nicht ein Ziel stecke, Tag oder Statt bestimme, noch Weise oder Maße setze seiner Erhörung, sondern das alles seinem Willen, Weisheit und Allmächtigkeit heimgeben, nur frisch und fröhlich warten der Erhörung: doch nicht wissen wollen, wie und wo, wie bald, wie lange und durch welche? denn seine göttliche Weisheit wird überschwänglich bessere Weise und Maß, Zeit und Statt finden, denn wir gedenken mögen, und sollten auch Wunderzeichen geschehen, wie im alten Testament. Da die Kinder Israel Gott vertraueten, er würde sie erlösen, und war doch keine mögliche Weise vor Augen, noch in allen ihren Gedanken: da thät sich das rothe Meer auf, und gab ihnen einen Weg hindurch, und ertränkte alle ihre Feinde auf einmal.

Also thät Judith, die heilige Frau, da sie höret, daß die Bürger zu Bethulia in fünf Tagen die Stadt wollten übergeben, wenn ihnen Gott indeß nicht hülfe, strafte sie dieselbigen und sprach: wer seid ihr, daß ihr Gott versuchet? Das sind nicht Fürnehmen, damit man Gnade erlange, sondern mehr Ungnade erwecket. Wollt ihr Gott eine Zeit setzen, sich euer zu erbarmen, und bestimmen einen Tag nach eurer Willkür? Darum half ihr auch Gott wunderlich, daß sie dem großen Holoferni sein Haupt abschlug, und die Feinde also vertrieben wurden.

Also spricht auch St. Paulus (Eph. 3.), daß Gottes Vermögen ist also gethan, daß er überschwänglich höher und besser thut, denn wir bitten oder verstehen. Darum soll man sich erkennen zu geringe dazu, daß wir nennen, abmaßen oder dargeben mögen Zeit, Statt, Weise, Maß und andere Umstände deß, daß wir bitten von Gott, sondern ihm ganz alles heimstellen, und unverrückt festiglich glauben, er werde uns erhören.

\quellseite{16}

\chapter[Verzeichniß der Quellen, aus welchen die in …]{Verzeichniß der Quellen, aus welchen die in diesem Gebets-Schatz enthaltenen Gebete genommen sind}

Wilh. Alard, Superintendent zu Krempe (1572—1645).

Joh. Albinus, Pastor zu Naumburg (gestorben 1679).

Bernh. Albrecht, Senior zu Augsburg (1569—1636).

Joh. Arndt, Generalsuperintendent zu Celle (1555—1621), Paradiesgärtlein.

Joh. Bugenhagen, Hauptpastor zu Wittenberg, Luthers Freund und Zeitgenosse (1485—1558).

Joh. Crüger, Pastor zu Stettin (gest. 1616), Uebung der Gottseligkeit.

Geo. Debekennus, Pastor zu Hamburg (1564—1628).

Joh. Mich. Dilherr, Oberprediger zu Nürnberg (1604—1669), Weg zur Seligkeit.

Joh. Georg Drechsler (um 1677).

Joh. Eichorn (starb 1564), Rüst- und Schatzkammer (1715).

Joh. Cundrenius, Pastor zu Nordhausen (1595—1660).

Christoph S. Facilibdes, Pastor zu Frankenberg (um 1660), Güldne Räuchwerksschale.

Joh. Finckler, Pastor zu Gleina (um 1674), Trostbüchlein.

Joh. Gerhard, Professor zu Jena (1582—1637), Uebung der Gottseligkeit.

Nic. Hags, Pastor prim. zu Bautzen (1665—1715), Das in Gott andächtige Frauenzimmer.

Joh. Habermann, Superintendent zu Zeitz (1516—1590), Christliches Gebetbuch.

Hartmann, Superintendent zu Rotenburg (1640—1684), Handbuch.

Joh. Herrmann, Pastor zu Köben (1564—1647), Buß-Leiter.

Salom. Lisecinius, Pastor zu Otterwisch (gest. 1716), Der zu seinem Jesu sich nahende Sünder.

B. C. Löscher, Superintendent zu Dresden (1672—1749), Unschuldige Nachrichten.

\quellseite{17}

Mart. Luther (1483—1546), Gebetstöcklein durch Pet. Trever.

Joh. Matthesius, Pastor zu Joachimsthal (1504—1568), Andächtige und gemeine christliche Gebete.

Joh. Olearius, Professor zu Leipzig (1611—1684).

Andr. Pancratius, Superintendent zu Hof (1531—1576), Haus- und Kirchenbuch.

Joh. Quirsfeld, Archidiaconus zu Pirna (1642—1686), Geistlicher Myrrhen-Garten.

Sigism. Scherertz, Superintendent zu Lüneburg (1584—1639), Seelen-Arznei und Spiegel der geistlichen Anfechtungen.

Georg Schimmer, Prediger zu Wittenberg (1652—1695), Biblisches Seelen-Kleinod.

Joh. Schrader (im 17. Jahrhundert).

Christ. Scriver, Oberhofprediger zu Quedlinburg (1629—1693), Gebetskleinod.

Georg Zämann, Superintendent zu Stralsund (1580—1638), Geistliches Bade Mecum.

Nürnberger Kinderlehrbüchlein.

Gebetbuch Churf. Christian des Ersten (1589)

Rigisches Gebetbuch (1719).

Himmlische Brautkammer.

Torgauischer Catechismus oder Handbuch (1676).

Mich. Cubach, Großes und vollkommenes Gebetbuch (1655).

Ein uraltes Gebethbuch (ohne Titel).

Israels Trost und Freude (1660).

Crone der Alten.

Himmlisches Freudenmahl.

Breslauer Gesangbuch (1751).

Eisenacher Gesangbuch (1760).

Freyberger Gesangbuch.

Leipziger Gesangbuch (1730).

Marburger Gesangbuch.

Eckendorffisches Gesangbuch.

Schwarzburger Agenda (1650).

\chapter[Verzeichniß der Ueberschriften sämmtlicher …]{Verzeichniß der Ueberschriften sämmtlicher Gebete}
\begin{originalinhalt}

\quellseite{18}

\inhaltsgruppe{I. Vorbereitungsgebete.}

\inhaltseintrag{1. Um wahre Andacht. (Luther.)}{1}

\inhaltseintrag{2. Um rechte Würdigkeit zu beten. (Luther.)}{1}

\inhaltseintrag{3. Danksagung, daß uns Gott zum Gebet würdig macht. (Luther.)}{1}

\inhaltseintrag{4. Trostgebet wider alle Sünde und Unwürdigkeit. (Luther.)}{2}

\inhaltseintrag{5. Gebet auf Gottes Befehl und Verheißung. (Luther.)}{2}

\inhaltseintrag{6. Versicherung der Erhörung um Christi willen. (Luther.)}{3}

\inhaltseintrag{7. Um gnädige Erhörung. (Luther.)}{3}

\inhaltseintrag{8. Vor dem Vater Unser. (Luther.)}{3}

\inhaltsgruppe{II. Dr. M. Luthers Morgen- und Abendsegen.}

\inhaltseintrag{9. Morgensegen. (Luther.)}{4}

\inhaltseintrag{10. Abendsegen. (Luther.)}{5}

\inhaltsgruppe{III. Tischgebete.}

\inhaltseintrag{11. Vor dem Essen. (Luther.)}{6}

\inhaltseintrag{12. Ein Anderes. (Luther.)}{6}

\inhaltseintrag{13. Ein Anderes. (Herrmann.)}{6}

\inhaltseintrag{14. Ein Anderes. (Breslauer Gesangbuch.)}{7}

\inhaltseintrag{15. Nach dem Essen. (Luther.)}{7}

\inhaltseintrag{16. Ein Anderes. (Torgauischer Catechismus.)}{7}

\inhaltsgruppe{IV. Gebete für jeden Tag in der Woche.}

\inhaltseintrag{17. Morgensegen am Sonntage. (Habermann.)}{8}

\inhaltseintrag{18. Ein Anderer. (Marburger Gesangbuch.)}{9}

\inhaltseintrag{19. Gebet aus der ersten Bitte des heil. Vater Unsers. (Haas.)}{10}

\inhaltseintrag{20. Abendsegen am Sonntag. (Habermann.)}{11}

\inhaltseintrag{21. Ein Anderer. (Marburger Gesangbuch.)}{12}

\inhaltseintrag{22. Morgensegen am Montag. (Habermann.)}{13}

\inhaltseintrag{23. Ein Anderer. (Marburger Gesangbuch.)}{14}

\inhaltseintrag{24. Gebet aus der zweiten Bitte des heil. Vater Unsers. (Haas.)}{14}

\inhaltseintrag{25. Abendsegen am Montage. (Habermann.)}{16}

\inhaltseintrag{26. Ein Anderer. (Marburger Gesangbuch.)}{17}

\inhaltseintrag{27. Morgensegen am Dienstag. (Habermann.)}{18}

\inhaltseintrag{28. Ein Anderer. (Marburger Gesangbuch.)}{19}

\inhaltseintrag{29. Gebet aus der dritten Bitte des heil. Vater Unsers. (Haas.)}{19}

\inhaltseintrag{30. Abendsegen am Dienstag. (Habermann.)}{21}

\inhaltseintrag{31. Ein Anderer. (Marburger Gesangbuch.)}{22}

\inhaltseintrag{(XI)}{}

\quellseite{19}

\inhaltseintrag{32. Morgensegen am Mittwoch. (Habermann.)}{22}

\inhaltseintrag{33. Ein Anderer. (Marburger Gesangbuch.)}{23}

\inhaltseintrag{34. Gebet aus der vierten Bitte des heil. Vater Unsers. (Haas.)}{24}

\inhaltseintrag{35. Abendsegen am Mittwoch. (Habermann.)}{26}

\inhaltseintrag{36. Ein Anderer. (Marburger Gesangbuch.)}{27}

\inhaltseintrag{37. Morgensegen am Donnerstag. (Habermann.)}{27}

\inhaltseintrag{38. Ein Anderer. (Marburger Gesangbuch.)}{28}

\inhaltseintrag{39. Gebet aus der fünften Bitte des heil. Vater Unsers. (Haas.)}{29}

\inhaltseintrag{40. Abendsegen am Donnerstag. (Habermann.)}{30}

\inhaltseintrag{41. Ein Anderer. (Marburger Gesangbuch.)}{31}

\inhaltseintrag{42. Morgensegen am Freitag. (Habermann.)}{32}

\inhaltseintrag{43. Ein Anderer. (Marburger Gesangbuch.)}{33}

\inhaltseintrag{44. Gebet aus der sechsten Bitte des heil. Vater Unsers. (Haas.)}{33}

\inhaltseintrag{45. Abendsegen am Freitag. (Habermann.)}{34}

\inhaltseintrag{46. Ein Anderer. (Marburger Gesangbuch.)}{35}

\inhaltseintrag{47. Morgensegen am Sonnabend. (Habermann.)}{36}

\inhaltseintrag{48. Ein Anderer. (Marburger Gesangbuch.)}{37}

\inhaltseintrag{49. Gebet aus der siebenten Bitte des heil. Vater Unsers und dem Beschluß. (Haas.)}{38}

\inhaltseintrag{50. Abendsegen am Sonnabend. (Habermann.)}{40}

\inhaltseintrag{51. Ein Anderer. (Marburger Gesangbuch.)}{41}

\inhaltsgruppe{V. Gebete für verschiedene Zeiten.}

\inhaltsgruppe{A. Tägliche Gebete.}

\inhaltseintrag{52. Eine allgemeine Beichte. (Luther.)}{42}

\inhaltseintrag{53. Gebet um ein seliges Ende. (Luther.)}{43}

\inhaltsgruppe{B. Allgemeine Sonntagsgebete.}

\inhaltseintrag{54. Vor der Predigt. (Luther.)}{44}

\inhaltseintrag{55. Danksagung für das gehörte Wort Gottes. (Luther.)}{44}

\inhaltsgruppe{C. Gebete für jeden Sonn- und Festtag des Kirchenjahrs.}

\inhaltseintrag{56. Gebet um fruchtbare Anhörung der Predigt. (Zäkmann.)}{45}

\inhaltseintrag{57. Ein Anderes. (Eubach.)}{45}

\inhaltseintrag{58. Gebet nach der Predigt. (Eubach.)}{46}

\inhaltseintrag{59. Ein Anderes. (Eubach.)}{47}

\inhaltseintrag{60. Gebet am ersten Sonntag des Advents. (Eichhorn.)}{48}

\inhaltseintrag{61. Gebet am zweiten Sonntag des Advents. (Eichhorn.)}{48}

\inhaltseintrag{62. Gebet am dritten Sonntag des Advents. (Eichhorn.)}{48}

\inhaltseintrag{63. Gebet am vierten Sonntag des Advents. (Eichhorn.)}{49}

\inhaltseintrag{64. Vorbereitungsgebet am heiligen Christtage. (Eichhorn.)}{49}

\inhaltseintrag{65. Gebet am heiligen Christtag. (Eichhorn.)}{50}

\inhaltseintrag{66. Gebet am Sonntag nach dem heiligen Christtage. (Eichhorn.)}{50}

\inhaltseintrag{67. Gebet am Sylvesterabend. (Schimmer.)}{51}

\inhaltseintrag{68. Gebet am Neujahr. (Eichhorn.)}{52}

\inhaltseintrag{69. Gebet am Sonntag nach Neujahr. (Eichhorn.)}{52}

\inhaltseintrag{70. Gebet am Epiphanias- oder Erscheinungsfeste Christi. (Eichhorn.)}{52}

\inhaltseintrag{71. Gebet am ersten Sonntag nach dem Feste der Erscheinung Christi. (Eichhorn.)}{53}

\inhaltseintrag{72. Gebet am zweiten Sonntag nach dem Feste der Erscheinung Christi. (Eichhorn.)}{53}

\quellseite{20}

\inhaltseintrag{73. Gebet am dritten Sonntag nach dem Feste der Erscheinung Christi. (Eichorn.)}{53}

\inhaltseintrag{74. Gebet am vierten Sonntag nach dem Feste der Erscheinung Christi. (Eichorn.)}{54}

\inhaltseintrag{75. Gebet am fünften Sonntag nach dem Feste der Erscheinung Christi. (Eichorn.)}{54}

\inhaltseintrag{76. Gebet am Sonntag Septuagesima. (Eichorn.)}{54}

\inhaltseintrag{77. Gebet am Sonntag Sexagesima. (Eichorn.)}{55}

\inhaltseintrag{78. Gebet am Sonntag Estomihi. (Eichorn.)}{55}

\inhaltseintrag{79. Gebet am ersten Sonntag in der Fasten. (Eichorn.)}{55}

\inhaltseintrag{80. Gebet am zweiten Sonntag in der Fasten. (Eichorn.)}{56}

\inhaltseintrag{81. Gebet am dritten Sonntag in der Fasten. (Eichorn.)}{56}

\inhaltseintrag{82. Gebet am vierten Sonntag in der Fasten. (Eichorn.)}{56}

\inhaltseintrag{83. Gebet am fünften Sonntag in der Fasten. (Eichorn.)}{56}

\inhaltseintrag{84. Gebet am Palmsonntag. (Eichorn.)}{57}

\inhaltseintrag{85. Vorbereitungsgebet auf das Osterfest. (Riegisches Gebetbuch.)}{57}

\inhaltseintrag{86. Gebet am Osterfeste. (Eichorn.)}{58}

\inhaltseintrag{87. Gebet am Ostermontag. (Eichorn.)}{58}

\inhaltseintrag{88. Gebet am Sonntag Quasimodogeniti. (Eichorn.)}{58}

\inhaltseintrag{89. Gebet am Sonntag Misericordias Domini. (Eichorn.)}{58}

\inhaltseintrag{90. Gebet am Sonntag Jubilate. (Eichorn.)}{58}

\inhaltseintrag{91. Gebet am Sonntag Cantate. (Eichorn.)}{59}

\inhaltseintrag{92. Gebet am Sonntag Rogate. (Eichorn.)}{59}

\inhaltseintrag{93. Gebet am Himmelfahrtstage. (Eichorn.)}{60}

\inhaltseintrag{94. Vorbereitungsgebet auf Pfingsten. (Riegisches Gebetbuch.)}{60}

\inhaltseintrag{95. Gebet am ersten Pfingsttag. (Eichorn.)}{61}

\inhaltseintrag{96. Gebet am zweiten Pfingsttag. (Eichorn.)}{61}

\inhaltseintrag{97. Gebet am Sonntag Trinitatis oder Fest der heil. Dreifaltigkeit. (Eichorn.)}{61}

\inhaltseintrag{98. Ein Anderes. (Eichorn.)}{61}

\inhaltseintrag{99. Gebet am ersten Sonntag nach Trinitatis. (Eichorn.)}{62}

\inhaltseintrag{100. Gebet am zweiten Sonntag nach Trinitatis. (Eichorn.)}{62}

\inhaltseintrag{101. Gebet am dritten Sonntag nach Trinitatis. (Eichorn.)}{62}

\inhaltseintrag{102. Gebet am vierten Sonntag nach Trinitatis. (Eichorn.)}{63}

\inhaltseintrag{103. Gebet am fünften Sonntag nach Trinitatis. (Eichorn.)}{63}

\inhaltseintrag{104. Gebet am sechsten Sonntag nach Trinitatis. (Eichorn.)}{63}

\inhaltseintrag{105. Gebet am siebenten Sonntag nach Trinitatis. (Eichorn.)}{64}

\inhaltseintrag{106. Gebet am achten Sonntag nach Trinitatis. (Eichorn.)}{64}

\inhaltseintrag{107. Gebet am neunten Sonntag nach Trinitatis. (Eichorn.)}{64}

\inhaltseintrag{108. Gebet am zehnten Sonntag nach Trinitatis. (Eichorn.)}{65}

\inhaltseintrag{109. Gebet am eilften Sonntag nach Trinitatis. (Eichorn.)}{65}

\inhaltseintrag{110. Gebet am zwölften Sonntag nach Trinitatis. (Eichorn.)}{65}

\inhaltseintrag{111. Gebet am dreizehnten Sonntag nach Trinitatis. (Eichorn.)}{65}

\inhaltseintrag{112. Gebet am vierzehnten Sonntag nach Trinitatis. (Eichorn.)}{66}

\inhaltseintrag{113. Gebet am fünfzehnten Sonntag nach Trinitatis. (Eichorn.)}{66}

\inhaltseintrag{114. Gebet am sechzehnten Sonntag nach Trinitatis. (Eichorn.)}{66}

\inhaltseintrag{115. Gebet am siebenzehnten Sonntag nach Trinitatis. (Eichorn.)}{67}

\inhaltseintrag{116. Gebet am achtzehnten Sonntag nach Trinitatis. (Eichorn.)}{67}

\inhaltseintrag{117. Gebet am neunzehnten Sonntag nach Trinitatis. (Eichorn.)}{67}

\inhaltseintrag{118. Gebet am zwanzigsten Sonntag nach Trinitatis. (Eichorn.)}{67}

\inhaltseintrag{119. Gebet am einundzwanzigsten Sonntag nach Trinitatis. (Eichorn.)}{68}

\inhaltseintrag{120. Gebet am zweiundzwanzigsten Sonntag nach Trinitatis. (Eichorn.)}{68}

\quellseite{21}

\inhaltseintrag{121. Gebet am dreiundzwanzigsten Sonntag nach Trinitatis. (Eichhorn.)}{69}

\inhaltseintrag{122. Gebet am vierundzwanzigsten Sonntag nach Trinitatis. (Eichhorn.)}{69}

\inhaltseintrag{123. Gebet am fünfundzwanzigsten Sonntag nach Trinitatis. (Eichhorn.)}{69}

\inhaltseintrag{124. Gebet am sechsundzwanzigsten Sonntag nach Trinitatis. (Eichhorn.)}{69}

\inhaltseintrag{125. Gebet am Geburtstage. (Eubach.)}{70}

\inhaltsgruppe{VI. Berufsgebete.}

\inhaltsgruppe{A. Nährstand.}

\inhaltseintrag{126. Gebet der Eltern für ihre Kinder. (Arndt.)}{72}

\inhaltseintrag{127. Gebet, wenn Eltern ihre Kinder zur Schule schicken. (Mathesius.)}{73}

\inhaltseintrag{128. Gebet der Eltern für ihre erwachsenen Kinder, daß ihnen Gott ein frommes Ehegemahl bescheren wolle. (Eubach.)}{73}

\inhaltseintrag{129. Gebet einer Mutter für ungerathene Kinder. (Haas.)}{74}

\inhaltseintrag{130. Trostgebet in Anfechtung bei seinem Beruf. (Luther.)}{76}

\inhaltseintrag{131. Ein Anderes. (Luther.)}{76}

\inhaltseintrag{132. Reisegebet eines Hausvaters, der mit seiner Familie seinen Wohnort verändert. (Eubach.)}{77}

\inhaltseintrag{133. Gebet eines Vormundes. (Schimmer.)}{78}

\inhaltseintrag{134. Gebet der Kinder. (Riegisches Gebetbuch.)}{79}

\inhaltseintrag{135. Gebet eines Dienstboten. (Gebetbuch Churf. Christian des Ersten.)}{80}

\inhaltseintrag{136. Gebet junger Leute, die sich in den Ehestand zu begeben Willens sind. (Luther.)}{81}

\inhaltseintrag{137. Ein Anderes. (Luther.)}{82}

\inhaltseintrag{138. Gebet junger Leute um ein frommes Ehegemahl. (Luther.)}{82}

\inhaltseintrag{139. Gebet eines jungen Mannes, der im Begriff steht, um ein Gemahl anzuhalten. (Eichhorn.)}{82}

\inhaltseintrag{140. Gebet einer Jungfrau, die zur Ehe begehret worden. (Eubach.)}{83}

\inhaltseintrag{141. Gebet im Brautstande. (Eubach.)}{84}

\inhaltseintrag{142. Gebet Verlobter am Trauungstage. (Eichhorn.)}{84}

\inhaltseintrag{143. Gebet eines Ehemannes. (Gebetbuch Churf. Christian des Ersten.)}{85}

\inhaltseintrag{144. Gebet einer Ehefrau. (Gebetbuch Churf. Christian des Ersten.)}{86}

\inhaltseintrag{145. Gebet um Ehesegen. (Schimmer.)}{87}

\inhaltseintrag{146. Gebet bei verstörtem Ehesegen. (Schimmer.)}{88}

\inhaltseintrag{147. Gebet, wenn die Frau sich vor der Geburt fürchtet. (Schimmer.)}{89}

\inhaltseintrag{148. Gebet, wenn die Stunde der Geburt naht. (Eubach.)}{90}

\inhaltseintrag{149. Gebet bei Todesgefahr in Kindesnöthen. (Eubach.)}{91}

\inhaltseintrag{150. Ein Anderes. (Schimmer.)}{92}

\inhaltseintrag{151. Danksagung nach glücklicher Geburt. (Eubach.)}{92}

\inhaltseintrag{152. Gebet eines Ehemanns nach der Entbindung seiner Frau von einem todten Kinde. (Breslauer Gesangbuch.)}{93}

\inhaltseintrag{153. Gebet einer Mutter, die ein gebrechliches Kind geboren hat. (Eubach.)}{94}

\inhaltseintrag{154. Gebet einer Mutter, nachdem ihr neugebohr. Kind ohne Taufe gestorben. (Haas.)}{94}

\inhaltseintrag{155. Gebet einer Mutter, die den ersten Kirchgang hält. (Eubach.)}{96}

\inhaltseintrag{156. Gebet eines Mannes, der eine böse Frau hat. (Eubach.)}{97}

\inhaltseintrag{157. Gebet einer Frau, die einen bösen Mann hat. (Quirsfeld.)}{98}

\inhaltseintrag{158. Gebet einer Frau, die einen wüsten und unfreundlichen Mann hat. (Eubach.)}{98}

\inhaltseintrag{159. Gebet, wenn ein Ehegemahl das andere verlassen hat. (Eubach.)}{100}

\inhaltseintrag{160. Gebet eines Witwers. (Schimmer.)}{100}

\quellseite{22}

\inhaltseintrag{\unsicher{unleserlich}tkammer.}{}

\inhaltseintrag{\unsicher{unleserlich}ft. (Heermann.)}{}

\inhaltseintrag{\unsicher{unleserlich}lard.)}{}

\inhaltseintrag{\unsicher{unleserlich}ten.)}{}

\inhaltseintrag{\unsicher{unleserlich}).}{}

\inhaltseintrag{\unsicher{unleserlich}. (Schimmer.)}{}

\inhaltseintrag{\unsicher{unleserlich}es zu verrichten vor hat. (Ein uraltes Ge-}{}

\inhaltseintrag{\unsicher{unleserlich}mung seiner Geschäfte. (Dilher.)}{113}

\inhaltseintrag{\unsicher{unleserlich}eit. (Zemmann.)}{114}

\inhaltseintrag{\unsicher{unleserlich}en in der Berufsarbeit. (Cubach.)}{114}

\inhaltseintrag{\unsicher{unleserlich}en Haushalt. (Riegisches Gebetbuch.)}{115}

\inhaltseintrag{178. Gebet Luthers über den Spruch: Alle eure Sorgen werfet auf Gott, denn er sorget für euch. (Luther.)}{116}

\inhaltseintrag{179. Gebet eines solchen, den Gott mit irdischen Gütern gesegnet hat. (Scriver.)}{117}

\inhaltseintrag{180. Gebet bei Antritt einer Reise. (Leipziger Gesangbuch.)}{118}

\inhaltseintrag{181. Gebet auf der Reise. (Quirsfeld.)}{119}

\inhaltseintrag{182. Gebet bei Ausfahrt in die See. (Cubach.)}{119}

\inhaltseintrag{183. Gebet eines Seefahrenden im Sturm. (Emdenius.)}{120}

\inhaltseintrag{184. Gebet in äußerster Wassernoth. (Nürnberger Kinderlehrbüchlein.)}{121}

\inhaltsgruppe{B. Lehrstand.}

\inhaltseintrag{185. Gebet eines Predigers. (Habermann.)}{122}

\inhaltseintrag{186. Ein Anderes. (Luther.)}{123}

\inhaltseintrag{187. Noch ein Anderes. (Luther.) 12\unsicher{unleserlich}}{}

\inhaltseintrag{188. Noch ein Anderes. (Luther.) 12\unsicher{unleserlich}}{}

\inhaltseintrag{189. Gebet eines Predigers, die heilige Schrift fruchtbarlich zu studiren. (\unsicher{unleserlich})}{}

\inhaltseintrag{190. Ein Anderes. (Luther.)}{}

\inhaltseintrag{191. Trost und Trost eines fleißigen Predigers bei seinem Studiren. (Luther.)}{}

\inhaltseintrag{192. Gebet eines Predigers vor der Predigt. (Pancratius.)}{}

\inhaltseintrag{193. Ein Anderes. (Luther.)}{}

\inhaltseintrag{194. Noch ein Anderes. (Luther.)}{}

\inhaltseintrag{195. Noch ein Anderes. (Pancratius.)}{}

\inhaltseintrag{196. Gebet eines Predigers nach der Predigt. (Luther.)}{}

\inhaltseintrag{197. Ein Anderes. (Cubach.)}{}

\inhaltseintrag{198. Gebet eines Predigers vor der Beichtshandlung. (Pancratius.)}{}

\inhaltseintrag{199. Gebet eines Predigers, wenn er etwas Schweres zu verrichten hat. (Pancratius.)}{}

\inhaltseintrag{200. Gebet eines Predigers vor dem Krankenbesuch. (Schrader.)}{}

\inhaltseintrag{201. Gebet eines Seelsorgers in Pestzeiten. (Hartmann.)}{}

\inhaltseintrag{202. Danksagung eines Predigers nach vollbrachter Amtsarbeit. (Pancratius.)}{}

\inhaltseintrag{203. Eine Andere. (Luther.)}{}

\inhaltseintrag{204. Gebet eines Lehrers. (Cubach.)}{}

\inhaltseintrag{205. Ein Anderes. (Albinus.)}{}

\inhaltseintrag{*}{}

\quellseite{23}

\inhaltseintrag{\unsicher{unleserlich} eines}{}

\inhaltseintrag{\unsicher{unleserlich}. Gebet eines Studenten der Theolog\unsicher{unleserlich}}{}

\inhaltseintrag{210. Gebet eines Studenten der Theolog\unsicher{unleserlich} nach dem Studiren. (Olearius.) ...}{}

\inhaltsgruppe{C. Wehrstand.}

\inhaltseintrag{\unsicher{unleserlich}2. Gebet einer obrigkeitlichen Person. (\unsicher{unleserlich}}{}

\inhaltseintrag{\unsicher{unleserlich}13. Gebet einer richterlichen Person. (Hee\unsicher{unleserlich}}{}

\inhaltseintrag{\unsicher{unleserlich}14. Gebet eines Geschwornen, der in einem \unsicher{unleserlich}}{}

\inhaltseintrag{\unsicher{unleserlich}15. Gebet eines Amtmanns, Juristen oder Rath\unsicher{unleserlich}}{}

\inhaltseintrag{216. Allgemeines Gebet eines Soldaten. (Luther.) ...}{}

\inhaltseintrag{217. Ein Anderes. (Eubach.) ................................}{}

\inhaltseintrag{218. Morgengebet eines Soldaten. (Riegisches Gebetbuch.) .....}{}

\inhaltseintrag{219. Ein Anderes. (Bugenhagen.) ................................}{}

\inhaltseintrag{220. Abendgebet eines Soldaten. (Riegisches Gebetbuch.) ........}{}

\inhaltseintrag{221. Ein Anderes. (Riegisches Gebetbuch.)}{146}

\inhaltseintrag{222. Gebet um herzhaften Muth. (Eubach.)}{147}

\inhaltseintrag{223. Gebet auf dem Wachposten. (Eubach.)}{148}

\inhaltseintrag{224. Gebet vor der Schlacht. (Luther.)}{148}

\inhaltseintrag{225. Stoßseufzer während der Schlacht. (Breslauer Gesangbuch.)}{149}

\inhaltseintrag{226. Gebet nach siegreicher Schlacht. (Schimmer.)}{149}

\inhaltseintrag{227. Gebet eines Verwundeten. (Breslauer Gesangbuch.)}{149}

\inhaltseintrag{228. Gebet in großen Schmerzen. (Eubach.)}{150}

\inhaltseintrag{229. Gebet vor der Operation. (Haas.)}{151}

\inhaltseintrag{\unsicher{unleserlich}30. Gebet in Krankheit. (Arndt.)}{153}

\inhaltseintrag{231. Dank- und Bußgebet nach Errettung aus mancherlei Gefahr. (Schimmer.)}{155}

\inhaltseintrag{232. Gebet christlicher Soldaten in Gemeinschaft zu sprechen. (Riegisches Gebetb.)}{156}

\inhaltseintrag{\unsicher{unleserlich}3. Gebet eines Soldaten für seine Angehörigen zu Hause. (Eubach.)}{157}

\inhaltseintrag{\unsicher{unleserlich}4. Gebet der Angehörigen zu Hause für den Soldaten im Felde. (Eubach.)}{158}

\inhaltseintrag{\unsicher{unleserlich} Gebet bei Todesfurcht. (Luther.)}{159}

\inhaltseintrag{\unsicher{unleserlich}ei Gebete in Todesnoth. (Arndt.)}{159}

\inhaltseintrag{\unsicher{unleserlich}bet um Frieden. (Luther.)}{161}

\inhaltseintrag{\unsicher{unleserlich}ie kleine Litanei}{161}

\inhaltsgruppe{VII. Fürbitten.}

\inhaltseintrag{\unsicher{unleserlich} Gebet christlicher Weiber bei einer Frau in Kindesnöthen. (Bresl. Gesangb.)}{161}

\inhaltseintrag{\unsicher{unleserlich} Gebet einer Mutter für ihren Säugling. (Scriver.)}{162}

\inhaltseintrag{241. Gebet eines, der Gevatter stehen soll. (Facilides.)}{163}

\inhaltseintrag{242. Gebet der Eltern für ihr Kind, das auf Reisen ist. (Himmlische Braut.)}{163}

\inhaltseintrag{243. Gebet für einen Sohn, der in der Fremde ist. (Haas.)}{164}

\inhaltseintrag{244. Gebet für ein krankes Kind. (Haas.)}{165}

\inhaltseintrag{245. Fürbitte kleiner Kinder für ihren auf der Reise begriffenen Vater. (Matthesius.)}{167}

\inhaltseintrag{246. Gebet frommer Kinder für ihre Eltern. (Arndt.)}{167}

\inhaltseintrag{247. Gebet der Kinder für ihren Vater, der auf Reisen ist. (Riegisches Gebetbuch.)}{168}

\inhaltseintrag{248. Gebet frommer Kinder für ihren kranken Vater. (Albrecht.)}{168}

\inhaltseintrag{249. Gebet der Kaufleute. (Eubach.)}{169}

\quellseite{24}

\inhaltseintrag{250. Fürbitte für die Confirmanden. (Schwarzburger Agende.)}{170}

\inhaltseintrag{251. Gebet eines Hausvaters für das Gesinde. (Lubach.)}{171}

\inhaltseintrag{252. Gebet für Kranke insgemein. (Habermann.)}{172}

\inhaltseintrag{253. Gebet für einen Kranken. (Scriver.)}{173}

\inhaltseintrag{254. Drei Fürbitten für einen Besessenen. (Löscher.)}{173}

\inhaltseintrag{255. Gebet für eine vom Satan besessene Person. (Luther.)}{175}

\inhaltseintrag{256. Gebet für eine kranke angefochtene Person. (Luther.)}{175}

\inhaltseintrag{257. Gebet bei einem Begräbniß. (Schimmer.)}{175}

\inhaltseintrag{258. Gebet eines Zuhörers für seinen Seelsorger. (Albrecht.)}{176}

\inhaltseintrag{259. Gebet eines Pfarrkindes. (Habermann.)}{177}

\inhaltseintrag{260. Gebet um rechtschaffene Prediger, sonderlich in Vacanz zu beten. (Habermann.)}{178}

\inhaltseintrag{261. Gebet am Bibelfeste. (Arndt.)}{179}

\inhaltseintrag{262. Gebet für Missionsstunden. (Habermann.)}{181}

\inhaltseintrag{263. Gebet für die Kirche. (Schimmer.)}{182}

\inhaltseintrag{264. Gebet für die christliche Kirche und wider ihre Feinde. (Luther.)}{183}

\inhaltseintrag{265. Gebet wider den Antichrist und Papst zu Rom mit seinem Anhang. (Luther.)}{183}

\inhaltseintrag{266. Gebet treuer Lehrer und Prediger auf Pastoral-Conferenzen. (Luther.)}{184}

\inhaltseintrag{267. Gebet um Beständigkeit und Erhaltung bei dem seligmachenden Wort am Schluß von Pastoral-Conferenzen. (Luther.)}{184}

\inhaltseintrag{268. Gebet in Gemeindeversammlungen. (Habermann.)}{185}

\inhaltseintrag{269. Gebet für die Kirche in Gemeindeversammlungen. (Israels Trost u. Freude.)}{186}

\inhaltseintrag{270. Gebet um Frieden in der Gemeinde. (Schimmer.)}{186}

\inhaltseintrag{271. Gebet um Bekehrung der Verführten. (Habermann.)}{187}

\inhaltseintrag{272. Gebet wider die Atheisterei und epikurische Sicherheit. (Schimmer.)}{188}

\inhaltseintrag{273. Gebet für den Ehestand. (Habermann.)}{190}

\inhaltseintrag{274. Gebet, wenn man pflegt neue Obrigkeit zu wählen. (Lubach.)}{191}

\inhaltseintrag{275. Gebet für die weltliche Obrigkeit. (Ein uraltes Gebetbuch.)}{191}

\inhaltseintrag{276. Gebet für das Gedeihen der Früchte des Landes. (Gebetbuch Churfürst Christian des Ersten.)}{192}

\inhaltseintrag{277. Fürbitte für die Wohlthäter. (Habermann.)}{193}

\inhaltseintrag{278. Fürbitte für die Feinde. (Gebetbuch Churfürst Christians des Ersten.)}{194}

\inhaltseintrag{279. Herzliches Gebet und Seufzen um die Zukunft Christi. (Zeä\-mann.)}{194}

\inhaltsgruppe{VIII. Buß-, Beicht- und Communion-Gebete.}

\inhaltseintrag{280. Prüfungsfragen. (Liscovius.)}{195}

\inhaltseintrag{281. Bußgebet eines nach angestellter Prüfung erwachten Sünders. (Freiberger Gesangbuch.)}{201}

\inhaltseintrag{282. Klage eines armen Sünders über seines Herzens Unbußfertigkeit. (Zeä\-mann.)}{206}

\inhaltseintrag{283. Gebet um Erkenntniß der Sünden. (Luther.)}{207}

\inhaltseintrag{284. Ein Gebet, wie sich ein armer Sünder vor Gott demüthigen soll. (Luther.)}{207}

\inhaltseintrag{285. Gebet um wahre Buße und rechtschaffene Bekehrung. (Zeä\-mann.)}{207}

\inhaltseintrag{286. Gebet, wenn man in rechte Erkenntniß der Sünde kommt. (Luther.)}{209}

\inhaltseintrag{287. Demüthige Klage eines Gefallenen, aber durch Gottes Gnade Aufgemachten. (Arndt.)}{209}

\inhaltseintrag{288. Gebet um wahren Glauben. (Luther.)}{212}

\inhaltseintrag{289. Appellation eines bußfertigen Sünders von Gottes Richterstuhl zu seinem Gnadenthron. (Luther.)}{213}

\inhaltseintrag{290. Zuflucht zur Gnade Gottes. (Luther.)}{213}

\quellseite{25}

\inhaltseintrag{291. Macht der Gnade. (Luther.)}{214}

\inhaltseintrag{292. Gebet um Gnade bei Gott. (Luther.)}{214}

\inhaltseintrag{293. Gebet um herzliche Erbarmung und Vergebung der Sünden. (Luther.)}{214}

\inhaltseintrag{294. Seufzer um Vergebung der Sünden. (Luther.)}{214}

\inhaltseintrag{295. Inbrünstiges Glaubensgebet. (Luther.)}{215}

\inhaltseintrag{296. Gebet um Vergebung der Sünden. (Eisenacher Gesangbuch.)}{215}

\inhaltseintrag{297. Gebet um Versöhnung mit dem Nächsten. (Eckendorfisches Gesangbuch.)}{217}

\inhaltseintrag{298. Gebet, wenn man zum Beichtstuhl gehen und die heil. Absolution begehren will. (Scherrenz.)}{218}

\inhaltseintrag{299. Danksagung nach der Absolution. (Lübach.)}{218}

\inhaltseintrag{300. Danksagung nach empfangener Absolution. (Luther.)}{219}

\inhaltseintrag{301. Morgengebet eines Communicanten. (Liscovius.)}{219}

\inhaltseintrag{302. Heilige Communionandachten. (Haas.)}{220}

\inhaltseintrag{303. Klage eines bußfertigen Communicanten über das Verderben seiner Natur. (Himmlisches Freudenmahl.)}{222}

\inhaltseintrag{304. Gebet um heilige Begier nach dem heil. Abendmahl. (Luther.)}{225}

\inhaltseintrag{305. Gebet um würdige Empfangung des heil. Abendmahls. (Eichhorn.)}{225}

\inhaltseintrag{306. Gebet um würdigen Genuß des heil. Abendmahls. (Himmlisches Freudenmahl.)}{226}

\inhaltseintrag{307. Gebet eines Communicanten vor Empfang des heil. Abendmahls. (Schimmer.)}{227}

\inhaltseintrag{308. Gebet in der Kirche vor der Empfangung des heil. Abendmahls. (Herrmann.)}{228}

\inhaltseintrag{309. Gebet vor dem heil. Abendmahle. (Marburger Gesangbuch.)}{232}

\inhaltseintrag{310. Gebet vor der Genießung des heil. Abendmahls. (Luther.)}{233}

\inhaltseintrag{311. Ein Anderes. (Luther.)}{233}

\inhaltseintrag{312. Klage über Unwürdigkeit. (Luther.)}{233}

\inhaltseintrag{313. Gebet, wenn man hinzugehet. (Luther.)}{233}

\inhaltseintrag{314. Kurzer Seufzer, wenn man den Leib Christi empfangen will. (Luther.)}{234}

\inhaltseintrag{315. Kurzer Seufzer, wenn man das heil. Blut Christi empfangen will. (Luther.)}{234}

\inhaltseintrag{316. Danksagung nach dem Genuß des heil. Abendmahls. (Ein uraltes Gebetbuch.)}{234}

\inhaltseintrag{317. Eine Andere. (Schimmer.)}{235}

\inhaltseintrag{318. Noch eine Andere. (Eichhorn.)}{236}

\inhaltseintrag{319. Noch eine Andere. (Luther.)}{236}

\inhaltseintrag{320. Eine Danksagung für die geistliche Vermahnung. (Luther.)}{237}

\inhaltseintrag{321. Tröstliche Versicherung vom heil. Abendmahl. (Luther.)}{237}

\inhaltseintrag{322. Eine Andere. (Luther.)}{237}

\inhaltseintrag{323. Klage über Unfruchtbarkeit in guten Werken. (Luther.)}{238}

\inhaltseintrag{324. Trost wider Unwürdigkeit und des Teufels Anfechtung. (Luther.)}{238}

\inhaltseintrag{325. Gebet um Gnade, Gottes Willen zu vollbringen. (Luther.)}{238}

\inhaltseintrag{326. Gebet, daß uns Gott vor ferneren Sünden behüte, und so wir sündigen, uns dieselben nicht zurechne. (Luther.)}{239}

\inhaltseintrag{327. Gebet nach dem Genuß des heil. Abendmahls um neuen Gehorsam. (Rieglisches Gebetbuch.)}{239}

\inhaltseintrag{328. Betrachtung der Kostbarkeit des Leibes Christi. (Arndt.)}{240}

\inhaltseintrag{329. Gebet über das Leiden Christi. (Gebetbuch Churfürst Christian des Ersten.)}{242}

\inhaltseintrag{330. Danksagung für das Leiden Christi. (Arndt.)}{243}

\inhaltseintrag{331. Gebet eines Betrübten, der sich bekümmert, als ob er das heil. Abendmahl unwürdig empfangen habe. (Debekennus.)}{245}

\quellseite{26}

\inhaltsgruppe{IX. Katechismusgebete.}

\inhaltseintrag{332. Gebet um wahre Erkenntniß Gottes. (Arndt.)}{247}

\inhaltseintrag{333. Gebet um wahre Gottesfurcht. (Arndt.)}{249}

\inhaltseintrag{334. Gebet um herzliche brünstige Liebe. (Arndt.)}{251}

\inhaltseintrag{335. Gebet um wahre beständige Hoffnung. (Arndt.)}{252}

\inhaltseintrag{336. Gebet um wahre Demuth. (Arndt.)}{253}

\inhaltseintrag{337. Gebet um wahre Geduld. (Arndt.)}{255}

\inhaltseintrag{338. Gebet um Heiligung des Namens Gottes. (Joh. Gerhard.)}{258}

\inhaltseintrag{339. Gebet um die Gnade, Gott zu loben. (Joh. Gerhard.)}{258}

\inhaltseintrag{340. Gebet um Liebe und Hochhaltung des göttlichen Worts. (Arndt.)}{259}

\inhaltseintrag{341. Gebet um Lust und Liebe zum göttlichen Worte. (Luthard.)}{260}

\inhaltseintrag{342. Gebet, wenn man in der Bibel lesen will. (Luther.)}{261}

\inhaltseintrag{343. Gebet vor dem Bibellesen.* (Schimmer.)}{261}

\inhaltseintrag{344. Gebet nach dem Bibellesen. (Schimmer.)}{262}

\inhaltseintrag{345. Gebet um die Ehre des heil. Predigtamtes. (Joh. Gerhard.)}{262}

\inhaltseintrag{346. Gebet um die Ehre der Eltern. (Joh. Gerhard.)}{263}

\inhaltseintrag{347. Gebet um christliche Freundlichkeit. (Haas.)}{264}

\inhaltseintrag{348. Gebet um Keuschheit des Herzens. (Arndt.)}{264}

\inhaltseintrag{349. Gebet um Gerechtigkeit im Handel und Wandel. (Joh. Gerhard.)}{266}

\inhaltseintrag{350. Gebet wider den Geiz. (Joh. Gerhard.)}{267}

\inhaltseintrag{351. Gebet um Gutthätigkeit. (Joh. Gerhard.)}{268}

\inhaltseintrag{352. Gebet um Wahrhaftigkeit. (Joh. Gerhard.)}{269}

\inhaltseintrag{353. Gebet um Regierung der Zunge. (Joh. Gerhard.)}{269}

\inhaltseintrag{354. Klage über das angeerbte Verderben. (Joh. Gerhard.)}{271}

\inhaltseintrag{355. Gebet um wahre Bekehrung. (Haas.)}{272}

\inhaltseintrag{356. Gebet um den rechten Glauben. (Habermann.)}{273}

\inhaltseintrag{357. Gebet um Erlösung vom Fluch des Gesetzes. (Luther.)}{274}

\inhaltseintrag{358. Danksagung für die Schöpfung. (Habermann.)}{275}

\inhaltseintrag{359. Danksagung für alle von Mutterleib an erzeigten Wohlthaten. (Luther.)}{276}

\inhaltseintrag{360. Danksagung für die Erlösung. (Habermann.)}{276}

\inhaltseintrag{361. Gebet zu Christo Jesu um seine Zukunft zum jüngsten Tage. (Schimmer.)}{278}

\inhaltseintrag{362. Gebet um baldige Zukunft zum Gericht. (Luther.)}{279}

\inhaltseintrag{363. Sehnliches Seufzen nach dem jüngsten Tage. (Luther.)}{279}

\inhaltseintrag{364. Betrachtung des jüngsten Gerichts. (Zäamann.)}{279}

\inhaltseintrag{365. Danksagung für die Heiligung. (Habermann.)}{281}

\inhaltseintrag{366. Danksagung für die Auferstehung der Todten. (Schimmer.)}{282}

\inhaltseintrag{367. Gebet nach den sieben Bitten des heil. Vater Unsers. (Schwarzburger Agenda.)}{283}

\inhaltseintrag{368. Gebet um Beständigkeit im Bekenntniß der Wahrheit in Versuchungszeit. (Schimmer.)}{284}

\inhaltseintrag{369. Ernstes Gebet um Bewahrung vor Verstockung und Abfall. (Luther.)}{285}

\inhaltseintrag{370. Gebet um ein seliges Ende. (Habermann.)}{285}

\inhaltseintrag{371. Gebet um Bewahrung vor einem bösen schnellen Tode. (Luthard.)}{287}

\inhaltseintrag{372. Gebet über die heil. Taufe. (Riegisches Gebetbuch.)}{288}

\inhaltseintrag{373. Danksagung für das Amt der Schlüssel, vor der Beichte zu beten. (Zäamann.)}{290}

\inhaltseintrag{374. Danksagung für die Einsetzung des heil. Abendmahls. (Dilherr.)}{291}

\inhaltseintrag{375. Gebet wider die Sacramentirer. (Luther.)}{294}

\inhaltsgruppe{X. Wettergebete.}

\inhaltseintrag{376. Gebet zur Zeit der Dürre. (Breslauer Gesangbuch.)}{294}

\inhaltseintrag{377. Danksagung nach erlangtem fruchtbaren Regen. (Embdenius.)}{295}

\quellseite{27}

\inhaltseintrag{378. Gebet um einen gnädigen warmen Sonnenschein zur Zeit langwieriger Nässe. (Breslauer Gesangbuch.)}{296}

\inhaltseintrag{379. Gebet für die lieben Früchte. (Luther.)}{297}

\inhaltseintrag{380. Gebet, wenn durch Gottes Verhängniß den Feldfrüchten vom Mehlthau und allerlei Ungeziefer Schaden geschieht. (Schimmer.)}{297}

\inhaltseintrag{381. Gebet bei Ungewittern. (Leipziger Gesangbuch.)}{298}

\inhaltseintrag{382. Gebet, wenn das Gewitter lange anhält. (Facilides.)}{299}

\inhaltseintrag{383. Gebet bei schrecklichem Sturm. (Riegisches Gebetbuch.)}{300}

\inhaltseintrag{384. Gebet bei Ungewittern. (Leipziger Gesangbuch.)}{300}

\inhaltseintrag{385. Danksagung, wenn das Gewitter glücklich vorüber gegangen ist. (Facilides.)}{301}

\inhaltseintrag{386. Gebet, wenn ein Komet oder andere wunderliche Zeichen am Himmel sichtbar werden. (Eubach.)}{302}

\inhaltsgruppe{XI. Kreuz- und Nothgebete.}

\inhaltseintrag{387. Gebet in großer Schwachheit des Glaubens. (Albinus.)}{302}

\inhaltseintrag{388. Klage wegen Mangels an Andacht im Gebet. (Scherertz.)}{303}

\inhaltseintrag{389. Gebet in Angst und Bangigkeit des Herzens. (Scherertz.)}{303}

\inhaltseintrag{390. Gebet in großer Sündenangst. (Albrecht.)}{304}

\inhaltseintrag{391. Gebet nach einem tiefen Sündenfall. (Riegisches Gebetbuch.)}{305}

\inhaltseintrag{392. Gebet in Anfechtung zur Verzweiflung. (Gebetbuch Churfürst Christian des Ersten.)}{307}

\inhaltseintrag{393. Zuflucht zu Christo in Sündennoth. (Luther.)}{309}

\inhaltseintrag{394. Ein anderes Gebet. (Luther.)}{309}

\inhaltseintrag{395. In großer Melancholie und Traurigkeit des Gemüths. (Scherertz.)}{309}

\inhaltseintrag{396. Gebet, wenn einer wegen der ewigen Versehung angefochten wird. (Albrecht.)}{311}

\inhaltseintrag{397. Gebet in satanischen Anfechtungen. (Scherertz.)}{312}

\inhaltseintrag{398. Kurzes Gebet wider des Satans Mordpfeile. (Schimmer.)}{313}

\inhaltseintrag{399. Danksagung nach überwindener Anfechtung. (Embdenius.)}{314}

\inhaltseintrag{400. Gebet bei Gefahr von Verführung. (Habermann.)}{315}

\inhaltseintrag{401. Gebet für unschuldig (Gefangene) Verfolgte. (Eubach.)}{316}

\inhaltseintrag{402. Gebet Vertriebener und Verfolgter. (Arndt.)}{318}

\inhaltseintrag{403. Gebet eines Christen, der unter Ungläubigen leben muß. (Dilherr.)}{320}

\inhaltseintrag{404. Gebet wider die Furcht bei unheimlichem Gepolter. (Schimmer.)}{320}

\inhaltseintrag{405. Gebet wider die Verleumder. (Haas.)}{321}

\inhaltseintrag{406. Gebet, wenn Kinder krank sind. (Luther.)}{322}

\inhaltseintrag{407. Gebet eines armen gefangenen Sünders um Vergebung der Sünden. (Eubach.)}{323}

\inhaltseintrag{408. Gebet eines Verbrechers, der zum Tode verurtheilt worden ist. (Eubach.)}{324}

\inhaltseintrag{409. Drei Gebete eines zum Tode Verurtheilten in der letzten Stunde. (Eubach.)}{325}

\inhaltseintrag{410. Gebet eines verlassenen armen Waisenkindes. (Schimmer.)}{326}

\inhaltseintrag{411. Gebet in Armuth. (Arndt.)}{327}

\inhaltseintrag{412. Trost in Armuth und Mangel. (Luther.)}{329}

\inhaltseintrag{413. Gebet um Ueberwindung in aller Noth. (Luther.)}{329}

\inhaltseintrag{414. Gebet eines Unterdrückten um der Wahrheit willen. (Luther.)}{329}

\inhaltseintrag{415. Gebet eines in einen Prozeß verwickelten Christen. (Schimmer.)}{330}

\inhaltseintrag{416. Gebet in Verlust zeitlicher Habe. (Eubach.)}{331}

\inhaltseintrag{417. Gebet in Feuersgefahr. (Olearius.)}{331}

\inhaltseintrag{418. Gebet in großer Wassernoth. (Schimmer.)}{332}

\inhaltseintrag{419. Gebet, wenn einer zum Krüppel geworden ist. (Schimmer.)}{333}

\quellseite{28}

\inhaltseintrag{420. Gebet in großer Noth und Gefahr. (Facilides.)}{334}

\inhaltseintrag{421. Ein Anderes. (Luther.)}{334}

\inhaltseintrag{422. Noch ein Anderes. (Luther.)}{335}

\inhaltseintrag{423. Noch ein Anderes. (Luther.)}{335}

\inhaltseintrag{424. Gebet in großer Widerwärtigkeit. (Luther.)}{335}

\inhaltseintrag{425. Ein Anderes. (Luther.)}{335}

\inhaltseintrag{426. Gebet in theurer Zeit. (Zeä\-mann.)}{335}

\inhaltseintrag{427. Gebet in ansteckenden Seuchen und Sterbensläuften. (Riegisches Gebetbuch.)}{337}

\inhaltseintrag{428. Gebet Eines, der Amts halber zur Zeit der Pestilenz nicht fliehen kann. (Luther.)}{338}

\inhaltseintrag{429. Wenn man nicht gebunden ist, sondern zu Pest- und Sterbenszeiten weicht und fleucht. (Luther.)}{339}

\inhaltseintrag{430. Gebet in Kriegsnoth. (Facilides.)}{339}

\inhaltseintrag{431. Ein Anderes. (Luther.)}{340}

\inhaltseintrag{432. Gebet um Frieden in Kriegszeiten. (Luther.)}{342}

\inhaltseintrag{433. Gebet bei allgemeinen Strafen und Landplagen. (Luther.)}{342}

\inhaltseintrag{434. Ein Anderes. (Luther.)}{342}

\inhaltseintrag{435. Gebet um Geduld und Ueberwindung. (Luther.)}{342}

\inhaltseintrag{436. Die Litanei, oder demüthige Fürbitte für die allgemeine Noth der ganzen Christenheit. (Pancratius.)}{343}

\inhaltsgruppe{XII. Dank- und Lobgebete.}

\inhaltseintrag{437. Gebet um christliche Dankbarkeit. (Arndt.)}{351}

\inhaltseintrag{438. Danksagung für die Liebe Gottes. (Arndt.)}{353}

\inhaltseintrag{439. Danksagung für die ewige Gnadenwahl in Christo. (Arndt.)}{354}

\inhaltseintrag{440. Danksagung für die heil. christliche Kirche. (Arndt.)}{356}

\inhaltseintrag{441. Danksagung für den Schutz der heiligen Engel. (Arndt.)}{358}

\inhaltseintrag{442. Danksagung für die gnädige Erhaltung und Behütung in so mancher Gefahr (Arndt.)}{359}

\inhaltseintrag{443. Danksagung nach der Arbeit. (Eichorn.)}{362}

\inhaltseintrag{444. Danksagung nach Vollendung eines neuen Wohnhauses. (Schimmer.)}{362}

\inhaltseintrag{445. Dankgebet nach glücklich vollbrachter Reise. (Eubach.)}{362}

\inhaltseintrag{446. Dankgebet nach Zurücklegung einer gefährlichen Seereise (besonders für Einwanderer). (Eubach.)}{363}

\inhaltseintrag{447. Danksagung für Erlösung aus Trübsal. (Eubach.)}{365}

\inhaltseintrag{448. Danksagung nach erlangtem Frieden. (Riegisches Gebetbuch.)}{366}

\inhaltseintrag{449. Ein Lob der Allmacht, Barmherzigkeit und Gerechtigkeit Gottes. (Arndt.)}{367}

\inhaltseintrag{450. Ein Lob der Herrlichkeit, Majestät und Freundlichkeit Gottes. (Arndt.)}{369}

\inhaltseintrag{451. Preis Jesu Christi von einer zu seinem Lobe entzündeten Seele. (Arndt.)}{370}

\inhaltsgruppe{XIII. Gebete für Kranke und Sterbende.}

\inhaltseintrag{452. Morgengebet eines Kranken. (Eichorn.)}{372}

\inhaltseintrag{453. Abendgebet eines Kranken. (Eichorn.)}{372}

\inhaltseintrag{454. Morgensegen eines Kranken. (Feinler.)}{372}

\inhaltseintrag{455. Abendsegen eines Kranken. (Feinler.)}{373}

\inhaltseintrag{456. Gebet in Krankheit. (Eisenacher Gesangbuch.)}{374}

\inhaltseintrag{457. Gebet eines Todkranken. (Facilides.)}{375}

\inhaltseintrag{458. Trostsprüche, den Kranken und Sterbenden vorzusprechen}{377}

\inhaltseintrag{459. Lieder und Liederverse}{395}

\quellseite{29}

\inhaltseintrag{460. Gebet eines bußfertigen Kranken. (Facilides.)}{425}

\inhaltseintrag{461. Gebet um Segen der Arznei. (Facilides.)}{425}

\inhaltseintrag{462. Kurze Stundengebete für Kranke. (Feinler.)}{426}

\inhaltseintrag{463. Gebet eines Kranken, der von Schwermuth angefochten wird. (Feinler.)}{430}

\inhaltseintrag{464. Gebet in langwieriger und schmerzhafter Krankheit. (Eisenacher Gesangb.)}{431}

\inhaltseintrag{465. Gebet um Erlösung aus Krankheitennoth. (Facilides.)}{432}

\inhaltseintrag{466. Danksagung nach Genesung. (Riegisches Gebetbuch.)}{433}

\inhaltseintrag{467. Gebet eines Kranken um Vergebung der Sünden vor Genießung des heil. Abendmahls. (Joh. Gerhard.)}{434}

\inhaltseintrag{468. Gebet eines Kranken nach Empfang des heil. Abendmahls. (Facilides.)}{435}

\inhaltseintrag{469. Gebet eines kranken Communicanten um Verlangen nach dem ewigen Leben. (Schimmer.)}{435}

\inhaltseintrag{470. Tröstlicher Zuspruch nach den drei Artikeln des christlichen Glaubens an einen sterbenden Menschen. (Eisenacher Gesangbuch.)}{436}

\inhaltseintrag{471. Seufzer in Todesnöthen. (Luther.)}{438}

\inhaltseintrag{472. Gebet um selige Auflösung durch ein sanftes Ende. (Luther.)}{439}

\inhaltseintrag{473. Gebet in Furcht und Schrecken des Todes und der Hölle. (Luther.)}{439}

\inhaltseintrag{474. Ein tröstliches Gebet in unserer letzten Stunde. (Luther.)}{439}

\inhaltseintrag{475. Gebet in Todesnöthen. (Luther.)}{440}

\inhaltseintrag{476. Gebet eines Sterbenden für die Seinen. (Dilherr.)}{440}

\inhaltseintrag{477. Gebet der Umstehenden, welche die Seele des Sterbenden Gott befehlen. (Haas.)}{440}

\inhaltseintrag{478. Ein Anderes. (Facilides.)}{441}

\inhaltseintrag{479. Eine kurze Litanei nebst Gebeten bei einem Sterbenden. (Lubach.)}{442}

\inhaltseintrag{480. Sterbensseufzer. (Scriver.)}{443}

\inhaltseintrag{481. Gebet, wenn der Kranke in dem HErrn verschieden ist. (Eisenacher Gesangb.)}{444}

\inhaltseintrag{482. Letzter Seufzer, einem Verurtheilten vorzusprechen. (Lubach.)}{444}

\end{originalinhalt}

\mainmatter

\quellseite{30}

\chapter{Vorbereitungsgebete}

\section{1. Um wahre Andacht}

Herr, weil du willst und heißest, daß ich beten und zu dir kommen soll, so will ich kommen, und zu beten genug bringen, und eben das, was mich am meisten hindert und von dir zurücktreibt, welches ist meine Sünde, die mir auf dem Halse liegt und drücket, daß du dieselbe von mir nehmen und vergeben wollest, um Christ willen. Amen.

\section{2. Um rechte Würdigkeit zu beten}

Herr, es ist deine Ehre und dein Gottesdienst, dadurch du gerühmt wirst, daß ich vor dir bettele; darum, lieber Herr, siehe nicht an, daß ich so unwürdig bin, sondern daß ich deiner Hilfe nothdürftig bin, und du der rechte einige Nothhelfer bist aller Sünder. Darum dir geschiehts zu Ehren, daß ich dich anrufe; so kann ich deiner Hülfe nicht entrathen, und du kannst und willst geben denen, die dich bitten. Amen. Amen.

\section[3. Danksagung, daß uns Gott zum Gebet würdig …]{3. Danksagung, daß uns Gott zum Gebet würdig macht}

Ich weiß gar wohl, o gnädiger Gott, daß ich ein unwürdiger Mensch bin, und würdig des Teufels, nicht Christi, noch seiner heiligen Bruder zu sein; nun aber hat Christus solches gesagt, daß ich, als für den er gestorben und auferstanden ist, sowohl als für St. Petrum, welcher auch mir gleich ein Sünder gewesen, sein Bruder sei; und will solches ernstlich von mir haben, daß ich ihm glauben soll ohne alles Zweifeln und Wanken, und nicht ansehen noch achten, daß ich unwürdig und voller Sünden bin; weil er es
\quellseite{31}
selbst nicht will ansehen noch gedenken, wie er doch billig thun könnte, und Ursach genug hätte zu rächen und zu strafen, sondern ist alles vergessen und aus seinem Herzen getilgt, ja todt, zugeschwäret und begraben; warum wollte ich es denn nicht also sein lassen, und meinem lieben HErrn nicht von Herzen dafür danken, loben und lieben, daß er so gnädig und barmherzig ist? Amen.

\section{4. Trostgebet wider alle Sünde und Unwürdigkeit}

Ich sei, wer ich wolle, so frage ich nichts danach; denn ob ich gleich ein Sünder bin, so weiß ich doch, daß darum mein HErr Christus nicht ein Sünder ist, sondern er bleibt gerecht und gnädig. Darum will ich getrost zu ihm rufen und schreien, und mich sonst an nichts kehren, denn ich habe jetzt nicht Weile zu disputiren, ob ich erwählt sei oder nicht; das aber fühle ich, daß ich Hülfe bedarf, komme derhalben und suche sie in aller Demuth. Lieber Gott, das canaanäische Weiblein war eine Heidin, konnte derhalben, ja sie mußte wohl schließen, sie wäre nicht erwählet. Tritt sie nun vor den HErrn JEsum, und läßt solchen Gedanken am Gebet sich nicht hindern, so thue ich auch also und spreche: HErr, ich komme jetzt, und muß dies und anders haben; wo will ichs sonst nehmen oder suchen, denn bei dir im Himmel, durch deinen Sohn, meinen Erlöser, Christum JEsum. Amen.

\section{5. Gebet auf Gottes Befehl und Verheißung}

Lieber HErr, du weißt, daß ich ja nicht von mir selbst und aus eigenem Vermessen, noch auf meine Würdigkeit vor dich komme; denn so ich das wollte ansehen, so dürfte ich die Augen nicht vor dir aufheben, und wüßte nicht, wie ich anfangen sollte zu beten; sondern darauf komme ich, daß du selbst geboten hast, und ernstlich forderst, daß wir dich sollen anrufen, und auch Verheißung zugesagt hast, dazu deinen eigenen Sohn gesandt, der uns gelehret, was wir beten sollen, und die Worte vorgesprochen hat; darum weiß ich, daß dir solch Gebet gefällt, und mein Vermessen, daß ich mich Gottes Kind vor dir rühmen darf, scheine wie groß es wolle, so muß ich dir gehorsam sein, der du es so haben willst, damit ich
\quellseite{32}
dich nicht Lügen strafe, und mich über andere Sünde noch schwerer gegen dich versündige, beides mit Verachtung deines Gebots und Unglauben an deine Verheißung.

\section{6. Versicherung der Erhörung um Christi willen}

O Gott, ein Schöpfer Himmels und der Erden, der du deinen Sohn, JEsum Christum, für mich in die Welt gesandt hast, daß er für mich gekreuzigt würde, stürbe, und am dritten Tage wieder auferstünde, gen Himmel führe, daß er da sollte sitzen zu deiner Rechten, und Alles in seiner Hand halten, uns seinen Geist senden, daß wir sollten warten auf seine Zukunft, zu richten beide, Lebendige und Todte, und also mit ihm erlangen das ewige Reich, unser Erbtheil, das du uns durch ihn geben willst. Dazu, o HErr Gott, hast du uns gegeben und eingesetzt die Taufe, und das Sacrament des Leibes und Blutes deines Sohnes, denn an diese seine Sacramente hat er uns Christen gebunden, und sich uns darin geoffenbaret. So wir ihn da ergreifen, so haben wir ihn gewißlich getroffen. Das Andere aber alles, das uns nicht befohlen ist, sollen wir fahren lassen; denn wir würden sonst sein weit fehlen.

\section{7. Um gnädige Erhörung}

Vater im Himmel, ich weiß, daß du mich lieb hast, darum daß ich deinen Sohn, meinen Erlöser JEsum Christum, lieb habe. In solchem Vertrauen und Zuversicht will ich dich jetzt tröstlich bitten, du wollest mich erhören, und mir geben, was ich bitte, nicht daß ich so heilig und fromm sei, sondern daß ich weiß, daß du um deines Sohnes Christi JEsu willen gern uns alles geben und schenken willst. In desselben Namen trete ich jetzt vor dich, und bitte und zweifle gar nicht, solch Gebet — ich sei, meiner Person halber, wer ich wolle — sei Ja und gewiß erhöret. Amen.

\section{8. Vor dem Vater Unser}

Ach himmlischer Vater, du lieber Gott: ich bin ein unwürdiger armer Sünder, nicht werth, daß ich meine Augen und Hände gegen dir aufhebe oder bete. Aber weil du uns allen geboten hast,
\quellseite{33}
zu beten, und dazu auch Erhörung verheißen, und über das selbst uns beide, Worte und Weise gelehrt durch deinen lieben Sohn, unsern HErrn JEsum Christ: so komm' ich auf solch dein Gebot, dir gehorsam zu sein, und verlasse mich auf deine gnädige Verheißung, und im Namen meines HErrn JEsu Christi bete ich mit allen deinen heiligen Christen auf Erden, wie er mich gelehret hat:

Vater unser! der du bist im Himmel, geheiliget werde dein Name, dein Reich komme, dein Wille geschehe, wie im Himmel, also auch auf Erden; unser tägliches Brod gib uns heute, und vergib uns unsere Schuld, als wir vergeben unsern Schuldigern, und führe uns nicht in Versuchung, sondern erlöse uns von dem Uebel, denn dein ist das Reich und die Kraft und die Herrlichkeit in Ewigkeit. Amen.

\chapter{II. Dr. M. Luthers Morgen= und Abendsegen}

\section{9. Morgensegen}

Das walte Gott Vater, Sohn und Heiliger Geist. Amen.

Ich glaube an Gott, den Vater, allmächtigen Schöpfer Himmels und der Erden.

Und an JEsum Christum, seinen einigen Sohn, unsern HErrn, der empfangen ist von dem Heiligen Geist, geboren aus Maria der Jungfrau, gelitten unter Pontio Pilato, gekreuziget, gestorben und begraben, niedergefahren zur Hölle, am dritten Tage wieder auferstanden von den Todten, aufgefahren gen Himmel, sitzend zur Rechten Gottes, des allmächtigen Vaters, von dannen er kommen wird, zu richten die Lebendigen und die Todten.

Ich glaube an den Heiligen Geist, eine heilige christliche Kirche, die Gemeinde der Heiligen, Vergebung der Sünden, Auferstehung des Fleisches und ein ewiges Leben. Amen.

Vater unser! der du bist im Himmel, geheiliget werde dein Name, dein Reich komme, dein Wille geschehe, wie im Himmel,
\quellseite{34}
also auch auf Erden; unser tägliches Brod gib uns heute, und vergib uns unsere Schuld, als wir vergeben unsern Schuldigern, und führe uns nicht in Versuchung, sondern erlöse uns von dem Uebel, denn dein ist das Reich und die Kraft und die Herrlichkeit in Ewigkeit. Amen.

Ich danke dir, mein himmlischer Vater, durch JEsum Christum, deinen lieben Sohn, daß du mich diese Nacht vor allem Schaden und Gefahr behütet hast, und bitte dich, du wollest mich diesen Tag auch behüten für Sünden und allem Uebel, daß dir all mein Thun und Leben gefalle. Denn ich befehle dir mein Leib und Seele und Alles in deine Hände, dein heiliger Engel sei mit mir, daß der böse Feind keine Macht an mir finde. Amen.

\section{10. Abendsegen}

Das walte Gott Vater, Sohn und Heiliger Geist. Amen.

Ich glaube an Gott, den Vater, allmächtigen Schöpfer Himmels und der Erden.

Und an JEsum Christum, seinen einigen Sohn, unsern HErrn, der empfangen ist von dem Heiligen Geist, geboren aus Maria der Jungfrau, gelitten unter Pontio Pilato, gekreuziget, gestorben und begraben, niedergefahren zur Hölle, am dritten Tage wieder auferstanden von den Todten, aufgefahren gen Himmel, sitzend zur Rechten Gottes, des allmächtigen Vaters, von dannen er kommen wird, zu richten die Lebendigen und die Todten.

Ich glaube an den Heiligen Geist, eine heilige christliche Kirche, die Gemeinde der Heiligen, Vergebung der Sünden, Auferstehung des Fleisches und ein ewiges Leben. Amen.

Vater unser! der du bist im Himmel, geheiliget werde dein Name, dein Reich komme, dein Wille geschehe, wie im Himmel, also auch auf Erden; unser tägliches Brod gib uns heute, und vergib uns unsere Schuld, als wir vergeben unsern Schuldigern, und führe uns nicht in Versuchung, sondern erlöse uns von dem Uebel, denn dein ist das Reich und die Kraft und die Herrlichkeit in Ewigkeit. Amen.

\quellseite{35}

Ich danke dir, mein himmlischer Vater, durch Jesum Christum, deinen lieben Sohn, daß du mich diesen Tag so gnädiglich behütet hast, und bitte dich, du wollest mir vergeben alle meine Sünde, wo ich unrecht gethan habe, und mich diese Nacht auch so gnädiglich behüten. Denn ich befehle dir mein Leib und Seele und Alles in deine Hände, dein heiliger Engel sei mit mir, daß der böse Feind keine Macht an mir finde. Amen.

\chapter{III. Tischgebete}

\section{11. Vor dem Essen}

Aller Augen warten auf dich, HErr, und du giebst ihnen ihre Speise zu seiner Zeit, du thuest deine Hand auf, und sättigest alles, was da lebet, mit Wohlgefallen.

(Darnach das Vater Unser und dies folgende Gebet:)

HErr Gott, himmlischer Vater, segne uns und diese deine Gaben, die wir von deiner milden Güte zu uns nehmen durch Jesum Christum, unsern HErrn. Amen.

\section{12. Ein Anderes}

Komm, HErr Jesu, sei unser Gast Und segne, was du bescheret hast. Amen.

\section{13. Ein Anderes}

\begin{verse}
Speis uns, o Gott, deine Kinder,\\
Tröste die betrübten Sünder,\\
Sprich den Segen zu den Gaben,\\
Die wir jetzo vor uns haben,\\
Daß sie uns zu diesem Leben\\
Stärke, Kraft und Nahrung geben,\\
Bis wir endlich zu den Frommen\\
An die Himmelstafel kommen. Amen.
\end{verse}

\quellseite{36}

\section{14. Ein Anderes}

Hütet euch, daß eure Herzen nicht beschweret werden mit Fressen und Saufen und mit Sorgen der Nahrung. Luc. 21, 34.

\begin{verse}
Geseg’n’ uns, HErr, die Gaben dein,\\
Die Speis laß unsre Nahrung sein,\\
Hilf, daß dadurch erquicket werd’\\
Der dürftig’ Leib auf dieser Erd’.\\
Denn diese zeitlich Brod allein\\
Kann uns nicht g’nug zum Leben sein,\\
Dein göttlich Wort die Seele speist,\\
Hilft uns zum Leben allermeist.
\end{verse}

\section{15. Nach dem Essen}

Danket dem HErrn, denn er ist freundlich und seine Güte währet ewiglich, der allem Fleische Speise giebt, der dem Vieh sein Futter giebt, den jungen Raben, die ihn anrufen. Er hat nicht Lust an der Stärke des Rosses, noch Gefallen an jemandes Beinen; der HErr hat Gefallen an denen, die ihn fürchten und die auf seine Güte warten.

(Darnach das Vater Unser und dies folgende Gebet:)

Wir danken dir, HErr Gott himmlischer Vater, durch JEsum Christum, unsern HErrn, für alle deine Wohlthat, der du lebest und regierest in Ewigkeit. Amen.

\section{16. Ein Anderes}

Es ist ein großer Gewinn, wer gottselig ist und lässet ihm genügen. Denn wir haben nichts in die Welt gebracht; darum offenbar ist, wir werden auch nichts hinausbringen. Wenn wir aber Nahrung und Kleider haben, so lasset uns begnügen. 1 Tim. 6, 6. 7.

\begin{verse}
Wir danken Gott für seine Gaben,\\
Die wir von ihm empfangen haben,\\
Wir bitten unsern lieben HErrn:\\
Er woll uns hinfort mehr beschern
\end{verse}

\quellseite{37}

\begin{verse}
Und speisen uns mit seinem Wort,\\
Daß wir satt werden hier und dort.\\
Ach lieber Gott, du wollst uns geben\\
Nach dieser Welt das ew’ge Leben. Amen.
\end{verse}

\begin{verse}
Ehre sei Gott dem Vater, der uns erschaffen hat.\\
Ehre sei Gott dem Sohne, der uns erlöst hat.\\
Ehre sei Gott dem Heiligen Geist, der uns geheiligt hat.\\
Ehre sei der heiligen und hochgelobten Dreieinigkeit von nun an bis in Ewigkeit. Amen.
\end{verse}

\chapter{IV. Gebete für jeden Tag in der Woche}

\section{17. Morgensegen am Sonntag}

Herr, himmlischer Vater, ewiger Gott, gebenedeiet sei deine göttliche Kraft und Allmächtigkeit, gelobet sei deine grundlose Güte und Barmherzigkeit, gepreiset sei deine ewige Weisheit und Wahrheit, daß du mich in dieser gefährlichen Nacht mit deiner Hand bedecket und unter dem Schatten deiner Flügel hast sicher ruhen und schlafen lassen, auch vor dem bösen Feind und allen seinen heimlichen Listen und Tücken bewahrt und ganz väterlich beschirmet. Darum lobe ich dich um deine Güte und um deine Wunder, die du an den Menschenkindern thust, und will dich bei der Gemeinde preisen. Dein Lob soll allweg in meinem Munde sein, meine Seele soll allezeit dich, meinen HErrn, rühmen, und was in mir ist, deinen heiligen Namen preisen, und will nimmermehr vergessen alles, was du mir Gutes gethan hast. So laß nun dir gefallen das Lobopfer aus meinem Munde, welches ich dir des Morgens früh in Einfältigkeit meines Herzens bringe. Ich rufe zu dir von ganzem Gemüth, du wollest mich heute diesen Tag auch behüten vor aller Gefahr Leibes und der Seelen und deinen lieben Engeln über mir Befehl thun, daß sie mich behüten auf allen meinen Wegen. Umsgib mich rings mit deinem Schilde und führe mich
\quellseite{38}
auf dem Steig deiner Gebote, daß ich unsträflich wandle in deinem Dienst, wie die Kinder des Tages, zu deinem Wohlgefallen. Wehre dem bösen Feind und allen Aergernißen dieser Welt, dazu steure meinem Fleisch und Blut, daß ich nicht, von ihnen überwältigt, etwa gröblich wider dich handle und dich mit meinen Sünden erzürne. Regiere du mich mit deinem heiligen Geist, daß ich nichts vornehme, thue, rede oder gedenke, denn allein, was dir gefällig, und zu Ehren deiner göttlichen Majestät gereicht. Siehe, mein Gott, ich übergebe und opfere dir mich heute ganz und gar eigen in deinem Willen, mit Leib und Seele, mit allem Vermögen und Kräften, innerlich und äußerlich. Mache du mich dir zu einem Opfer, das da lebendig, heilig und dir wohlgefällig sei, damit ich dir einen vernünftigen und angenehmen Gottesdienst leiste. Darum, du heiliger Vater, allmächtiger Gott, laß mich dein Eigenthum sein, regiere mein Herz, Seel und Gemüth, daß ich nichts denn dich wisse und verstehe. Herr, frühe wollest du meine Stimme hören, früh will ich mich zu dir schicken und darauf merken, früh will ich dich loben und des Abends nicht aufhören. Durch JEsum Christum. Amen.

\section{18. Ein Anderer}

Ewiger, allmächtiger Gott und Vater! ich sage dir von Herzen Lob und Dank, daß du mich in dieser vergangenen Nacht und alle vorige Zeit durch deine heiligen Engel für allem Uebel und Gefahr Leibes und der Seelen gnädiglich behütet hast; und bitte dich von Grund meines Herzens, vergiebe mir alle meine Sünden, damit ich dich je erzürnet habe, und zünde mein Herz an mit deinem heiligen Geist, daß ich in deinem Erkenntniß täglich wachse und zunehme. Verleihe mir auch Gnade, daß ich mich diesen angehenden Tag über für Sünden und Schanden hüte, und in deinem göttlichen Willen erfunden werde, damit ich also möge wandeln, daß ich für allem Uebel behütet, dich stets und ohne Unterlaß in meinem Herzen und Gedanken behalte, und wenn mein Sterbstündlein herzu nahet, daß ich von hinnen soll scheiden, wollest du mich in deines lieben Sohnes JEsu Christi rechtschaffenem Erkenntniß seliglich lassen einschlafen zum ewigen Leben. Amen.

\quellseite{39}

\section[19. Gebet aus der ersten Bitte des heil. …]{19. Gebet aus der ersten Bitte des heil. Vater Unsers}

Abba, heiliger Vater, deinem heiligen und herrlichen Namen zu Ehren singen die Cherubim und Seraphim droben im Himmel ohne Unterlaß mit erhabener Stimme: Heilig, heilig, heilig ist unser Gott, der Herr Zebaoth, alle Lande sind seiner Ehren voll. So gib nun Gnade, daß auch von deiner ganzen werthen Christenheit auf Erden heut und allezeit dein Name gebührend geheiligt und verehret werde. Hilf, daß ich und alle Menschen deine Allmacht recht erkennen, deine Gerechtigkeit mit Zittern scheuen, deiner Wahrheit sicherlich trauen, deine Barmherzigkeit von Herzen lieben, deine Weisheit gottseelig bewundern, deiner Güte uns immerfort freuen, und deine Ehre in allem eifrig suchen mögen.

Mache, o Gott! an diesem Sonntage sowohl meinen liebwerthen Seelsorgern, als auch andere treue Lehrer und Prediger, zu Gefäßen deiner Barmherzigkeit, und zu Werkzeugen deiner Gnade, und laß ihren Dienst zum Lobe deiner Herrlichkeit und aller ihnen anvertrauten Seelen Wohlfahrt und Seligkeit gereichen. Gib, daß ihre Lippen allezeit die Wahrheit deiner Lehre unverfälscht bewahren, und keiner teuflischen Ketzerei, oder seelenverderblichen Irrthum Raum geben. Rüste sie aus mit Kraft aus der Höhe, daß sie deiner Gemeinde, welche Christus durch sein Blut so theuer erworben hat, das Wort des Lebens rein und lauter predigen, die Unwissenden mit heiliger Lehre unterrichten, die Gottlosen mit freudigem Aufthun ihres Mundes strafen, und die Betrübten mit kräftigem Trost erquicken können. Schenke zu ihrer heiligen Arbeit dein göttliches Gedeihen, daß dein Wort niemals leer wieder komme, sondern wohl ausrichte, wozu du es sendest, und in meinem und aller Zuhörer Herzen hundertfältige Früchte bringe zum ewigen Leben. Lehre sie auch in ihrem Leben und Wandel thun nach deinem Wohlgefallen, und führe sie durch deinen guten Geist auf ebener Bahn, damit nicht von ihnen einiges Aergerniß gegeben, und dadurch ihr Amt verlästert werde.

Laß, o barmherziger Jesu! mich und andere fromme Christen allen Rath Gottes von unserer Seligkeit, den du durch diese deine Diener und Botschafter uns offenbarst, wohl lernen, und
\quellseite{40}
· selig beobachten. Erleuchte die Augen unsers Verstandes, daß wir begreifen mögen mit allen Heiligen, welches da sei die Breite und die Länge, die Tiefe und die Höhe deiner Liebe, womit du uns geliebet hast, und entzünde in uns ein recht herzliches und beständiges Vertrauen, daß wir auf dich und dein Verdienst im Leben und Tod uns festiglich verlassen, eine solche heilige und herzliche Gegenliebe, daß wir abtreten von aller Sünde und Ungerechtigkeit, weil wir nach deinem Namen genennet sind, und weder Engel noch Fürstenthum, noch Gewalt, weder Hohes noch Tiefes, weder Gegenwärtiges noch Zukünftiges uns von dir scheiden lassen.

Heilige uns sanft und sonder, o Gott Heiliger Geist, in deiner Wahrheit, und verleihe, daß wir dein theures Wort, so unsre Seelen selig machen kann, mit rechter Andacht und heiliger Ehrerbietung anhören, durch wahren Glauben in unsre Herzen einschreiben und versiegeln, deinen großen Namen zu Ehren unser Leben mit höchstem Fleiße darnach anstellen; dermaleinst aber zu deiner ewigen Herrlichkeit im Himmel gelangen, und dich in dem Chor aller heiligen Engel und Auserwählten ewig loben und preisen mögen. Amen.

\section{20. Abendsegen am Sonntag}

Ewiger Gott, barmherziger Vater, ich hebe meine Hände auf zu dir, wie ein Abendopfer und sage dir von Herzen Lob, Preis und Dank, daß du mich diesen Tag und die ganze Zeit meines Lebens vor allem Uebel und Unfall durch den Schutz deiner lieben Engel wider den bösen Feind gnädiglich beschirmet hast. Ich bitte dich, du wollest mir vergeben alle meine Sünde, wo ich Unrecht gethan habe, und mich heute diese Nacht ferner mit deinen heiligen Engeln umgeben, daß sie sich zurings um mich lagern und eine Wagenburg um mich schlagen, auf daß ich des bösen Feindes Fallstrick und arger List möge entfliehen. Laß mich, dein armes Geschöpf, deiner Güte und Barmherzigkeit befohlen sein. Schütze mich mit deinem ausgestreckten Arm, denn von Herzen begehre ich dein des Nachts, dazu mit meinem Geist in mir wache ich alle Zeit zu dir. Ich harre deiner Güte und meine Seele hofft auf dich, lebendiger Gott, denn du bist meine Zuflucht und treuer Heiland.

\quellseite{41}

Siehe, HErr, wir schlafen oder wachen, so sind wir dein, wir leben oder sterben, so bist du ja unser Gott, der uns geschaffen hat. Darum schreie ich zu dir, laß deine Gnade nicht ferne von mir sein, beschirme mich mit deinem Schilde. Erhalte mich, daß ich geruhiglich liege und sanft einschlafe und gesund wieder erwache. Decke mich in deiner Hütten zur bösen Zeit, verbirg mich heimlich in deinem Gezelt und erhöhe mich auf einen Felsen, so werde ich mir nicht grauen lassen. Und ob ich schon wanderte im finstern Thal, werde ich kein Unglück fürchten, denn du bist bei mir, dein Stecken und Stab trösten mich. Verleihe mir nur deine Gnade, daß mein Leib schlafe, die Seele aber allzeit zu dir wache, und ich dich stets in meinem Herzen habe und du nimmermehr aus meinem Gemüthe kommest, auf daß mich die Nacht der Sünden nicht überfalle. Bewahre mich vor bösen schändlichen Träumen, vor unruhigem Wachen und unnützen Sorgen, vor argen schweren Gedanken und vor aller Quaal. Siehe, mein Gott, in deinen gnädigen Schutz befehle ich dir meinen Leib und meine Seele, (meine lieben Geschwister) und alle, die mit mir durch Blutfreundschaft oder sonst verwandt sind. Laß dirs gefallen, HErr, daß du uns errettest, wende deine Barmherzigkeit nicht von uns, laß deine Güte und Treue allweg uns behüten. Darum bedecke uns in dieser Nacht mit deiner Güte, umfasse uns mit deiner Gnade auf allen Seiten, unterlege uns mit deiner Barmherzigkeit, daß wir an Seel und Leib behütet werden. Amen.

\section{21. Ein Anderer}

HErr, allmächtiger Gott und himmlischer Vater, ich danke dir von Grund meines Herzens für alle deine Güte und Wohlthat, die du mir heut diesen Tag so väterlich erzeiget und mich vor allem Uebel frisch und gesund behütet hast. Und bitte dich ferner, als dein liebes Kind, du wollest mich bei deinem allein seligmachenden Wort bis an mein letztes Seufzen gnädiglich erhalten, und mein Herz mit deinem heiligen Geist erleuchten, daß ich möge erkennen, was gut oder böse ist. Und wollest alle meine Sünde, so ich heute wissentlich gethan habe, auch die verborgenen Fehler, durch deine Barmherzigkeit vergessen und mir die zukünftige Nacht eine christliche
\quellseite{42}
Ruhe verleihen, daß ich frisch und gesund zu deinem Lobe wieder aufstehe. Hilf auch, lieber Vater, daß ich durch deine Gnade ein neu, dir wohlgefällig Leben anfange, zu meiner Seelen Seligkeit, in Jesu Christo, deinem lieben Sohn, unserm einigen Helfer. Amen.

\section{22. Morgensegen am Montag}

O du unsterblicher, gütiger Gott! weil du deinem Volk im Gesetz geboten, dir alle Morgen ein Brandopfer zu thun, dich damit zu loben und dir zu danken für deinen gnädigen Schutz, so bringe ich dir jetzt auch das Lobopfer, das ist die Frucht meiner Lippen, und preise deinen heiligen Namen, daß du mich durch deine große Gnade und Barmherzigkeit heute diese Nacht vor allem Uebel und Schaden an Leib und Seele behütet und wider meinen Widersacher, den bösen Feind, ganz väterlich erhalten hast. Denn wo du nicht mein Schild und Beistand wärest gewesen, so hätte mich unzähliges Unglück verderbet, daß ich nicht gesund hätte mögen aufstehen. Darum danke ich dir je billig für deine gnädige Beschützung. Weiter schreie ich zu dir aus Grund meines Herzens, und mein Gebet kommt früh vor dich; früh suche ich dich und bitte, du wollest mich mit Allem, was mir zuständig ist, heut ferner behüten vor der List und Gewalt des Teufels, vor Sünden, Schulden und allem Uebel. Komm du mir zuvor in dieser Frühstunde mit deiner Gnade, sintemal ich ohne dich nichts vermag, und hilf, daß ich an diesem Tage alle meine Werke in deinem Namen christlich anfange und seliglich vollführe, zu Ehren deiner göttlichen Majestät und zu Nutz meines Nächsten. Bewahre meine Seele, Verstand, Vernunft, Sinne und Gedanken, all mein Thun und Lassen, daß der leidige Teufel mir nicht könne schaden. Beschütze mich vor der Seuche, die im Mittag verderbet; behüte mich vor allen meinen Feinden, sichtbaren und unsichtbaren, daß sie mich mit ihren Listen und Tücken, mit Bezauberung und Vergiftungen, mit all ihrer Gewalt und Bosheit heimlich und öffentlich nicht können verletzen noch beschädigen. Herr Gott Vater und Herr meines Lebens, behüte mich auch vor unzüchtigem Gesichte und wende von mir alle böse Lust. Laß mich nicht in Schlemmerei und Unkeuschheit geraten
\quellseite{43}
und behüte mich vor unverschämtem Herzen. Hilf, daß ich durch deine Gnade das Auge, so mich ärgern will, ausreiße und von mir werfe durch Ablegung böser und unreiner Begierden des Herzens. Nimm von mir hinweg Alles, was dir mißfällt und mir schädlich ist; gib mir, was dir wohlgefällt und mir nützlich ist, auf daß ich dir in wahrem Glauben diene. Siehe mich an, du Heiland aller Welt, mit den Augen deiner Barmherzigkeit und erleuchte mein Herz und Augen, daß ich wandele in dem Glanz deiner Gnaden, so über mir aufgehet, auf daß ich doch, das ewige Licht, nimmermehr verliere. Amen.

\section{23. Ein Anderer}

Ach du frommer und getreuer Gott und himmlischer Vater! ich lobe, ehre und preise dich von Grund meines Herzens, daß du mich diese vergangene Nacht hast sicher ruhen und schlafen lassen und durch deine väterliche Liebe mich frisch und gesund wieder erwecket hast. Ich bitte dich von Herzen, du wollest mich sammt meinen lieben Eltern, Brüdern und Schwestern und allen frommen Christen, heut diesen Tag und alle Zeit vor allem Uebel und Gefahr Leibes und der Seelen auch gnädiglich behüten, damit ich alle Tage in deinem Willen möge erfunden werden. Denn ich befehle mich, meinen Leib und Seele, Herz, Sinn, Muth und Gedanken, all mein Dichten und Trachten, mein Gehen und Stehen, mein Sitzen und Liegen, meinen Eingang und Ausgang, mein Leben und Sterben, und alles, was ich bin und vermag, in deinen göttlichen Schutz und Schirm. Dein heiliger Engel sei und bleibe bei mir, daß mir kein Unglück an Seele und Leib widerfahren möge. Solches verleihe mir um Jesu Christi, deines lieben Sohnes willen. Amen.

\section[24. Gebet aus der zweiten Bitte des heil. …]{24. Gebet aus der zweiten Bitte des heil. Vater Unsers}

Abba, allmächtiger Vater und Herr Himmels und der Erden! Ich preise dich von ganzem Herzen, daß du mich und meine gläubige Mitchristen (da wir von Natur Kinder der Sünden und Sclaven des Teufels waren) errettet hast von der Obrigkeit der Finsterniß, und hast uns versetzet in das Reich deines lieben Sohnes Jesu
\quellseite{44}
Christi, an welchem wir haben die Erlösung durch sein Blut, die Vergebung der Sünden, die Gerechtigkeit des Glaubens, den Frieden des Gewissens, und die gewisse Hoffnung des ewigen Lebens. Ach allgewaltiger Gott, erhalte und befestige doch mit deiner starken Hand und allmächtigem Arm dieses Reich, so du unter uns hast, daß die Pforten der Höllen dasselbe niemals überwältigen mögen. Wohne allezeit, bis an das Ende der Welt, bei und unter uns mit deinem seligmachenden Worte, welches unsers Herzens Freude und Trost ist, und laß der heiligen Sacramente, als der theuersten Schätze unserer Seelen, uns niemals beraubet werden. Beschütze und bewahre deine christgläubige Gemeine vor aller falschen Lehre und ärgerlichem Sündenwesen, als dem rechten Greuel der Verwüstung; tritt unter unsere Füße den leidigen Satan, steuere den mächtigen Feinden unsers Glaubens, wehre den wüthenden Verfolgern deiner Kirche, und gib nicht zu, daß dein Volk, so du dir zum Eigenthum erwählet, und nach deinem Namen genennet hast, von denselben untergedrückt werde.

Vermehre aber auch und erweitere durch deine überschwängliche Gnade dieses dein Gnadenreich, und bringe zu demselben die armen blinden Heiden, Juden und Türken, und andere Kinder des Unglaubens, so noch unter des Teufels Reich leben und schweben. Siehe, o liebreicher Heiland, mit erbarmendem Herzen und mitleidigen Augen ihren Jammer und Elend an, thue hinweg die Decke, so für ihren verblendeten Herzen hänget, und bekehre sie, wo sie anders zu bekehren sind, von denen Wegen, welche sie führen ins Verderben und Verdammniß. Sie sind ja nicht zu dem ewigen Tode erschaffen, sondern nachdem sie darein durch des Teufels Neid gerathen sind, hast du, mein Jesu, sie durch dein blutiges Verdienst, so wohl als uns, wiederum erlöset. Ach so laß derowegen deine Mühe und Arbeit, deine Marter und Pein, dein Kreuz und Tod an ihnen nicht umsonst und vergeblich sein, sondern hilf, daß sie zur Erkenntniß der Wahrheit kommen, und ewig selig werden mögen.

Mich aber und andere Christen, welche du in der heiligen Taufe zu deinen Reichs-Genossen angenommen hast, regiere, Herr Gott Heiliger Geist, daß wir verleugnen alles ungöttliche Wesen, und
\quellseite{45}
die weltlichen Lüste, und meiden alle Wercke des Fleisches, als da sind Ehebruch, Hurerey, Unreinigkeit, Unzucht, Abgötterey, Zauberey, Feindschaft, Hader, Neid, Zorn, Zanck, Zwietracht, Rotten, Haß, Mord, Saufen, Fressen und dergleichen, von welchen du uns durch deinen außerwählten Rüstzeug Paulum zuvor gesaget hast, daß die solches thun, sollen das Reich Gottes nicht erben. Stärke uns im rechten Glauben, und in der wahren Gottseligkeit, daß wir in dem Reich Christi auch christlich leben, und ihm dienen in Heiligkeit und Gerechtigkeit, wie es ihm gefällig ist. Endlich nimm uns auf in das ewige Freuden- und Ehren-Reich des Himmels, und laß unsere Seele einst bey ihrer Ausfahrt aus dem Leibe die Gnadenstimme unsers JEsu hören: Kommet her, ihr Gesegneten meines Vaters, ererbet das Reich, das euch bereitet ist von Anbeginn der Welt. Amen.

\section{25. Abendsegen am Montag}

O du starker, lebendiger Gott, Vater unsers HErrn JEsu Christi! ich danke dir, daß du mich heute diesen Tag durch deinen göttlichen Schutz vor allem Schaden und Gefahr gnädiglich behütet hast, daß ich nicht bin kommen in Feuersnoth, daß mich Wasserfluthen nicht ersäuft haben, daß mich wilde Thiere nicht zerrissen haben, daß ich nicht bin gefallen in die Schärfe des Schwerdts, daß mich die Feinde nicht ermordet und böse Leute nicht verwundet haben, daß ich von Dieben und Räubern nicht erschlagen noch beschädigt bin worden, dazu auch, daß ich mich nicht habe zu Tode gefallen und sonst unvorsichtiglich gestoßen oder in andern Wege bin verletzt worden, in Summa, daß ich nicht kommen bin in Fährlichkeit Leibes und Lebens. Das Alles habe ich allein deiner Barmherzigkeit zu danken, der du mich auf allen meinen Wegen und Stegen behütest. Ich bitte nun ferner, du wollest mir vergeben alle meine Sünden, die ich wider dich gethan habe, und mich sammt Allen, die mir verwandt sind und zugehören, heute diese Nacht und die ganze Zeit unseres Lebens auch gnädiglich bewahren vor aller Angst und Beschwerniß, vor des Teufels List und Geschwindigkeit, damit er uns Tag und Nacht gedenckt zu bestricken. Behüte uns auch vor der schädlichen Pestilenz, die im Finstern
\quellseite{46}
schleichet; bewahre uns vor dem Strick des Feindes; daß wir nicht kommen in das Netz des Jägers, so u nachstehet. Beschütze uns auch vor schwerem Sich Schrecken des Satans, vor allen Uebeln Leibes und der Denn du bist unsere feste Burg, unsre Wehr und Waffen; steht all unsre Hoffnung und Zuversicht. Darum, du getreue laß deine Augen über uns offen sein und bewahre uns i Nacht vor aller Gewalt und Anfechtung des bösen Feind unser Wächter und Hüter, umgib uns mit deinem Schutz, a uns der Teufel nicht könne beschädigen. Denn bei dir allein unser Heil, zu dir allein hebe ich meine Augen auf, daher Hilfe kommt. Meine Hilfe kommt vom Herrn, der Himmel Erde gemacht hat. Siehe, wie die Augen der Knechte auf Hände ihrer Herren sehen, wie die Augen der Mägde auf die Hä ihrer Frauen; also sehen unsre Augen auf den Herrn, un Gott, bis er uns gnädig werde. Sei uns gnädig, Herr, sei gnädig! denn wir sind arm und elend. Erhebe über mich das deiner Antlitzes und erleuchte meine Augen, daß ich nicht im entschlafe, der du lebest und regierest in Ewigkeit. Amen.

\section{26. Ein Anderer}

Barmherziger Gott, gnädiger Vater! ich danke dir von ganzem Herzen, daß du mich heut diesen Tag bis auf gegenwärtige Stunde vor allem Unglück Leibes und der Seelen gnädiglich behütet hast, und bitte dich um JEsu Christi, deines lieben Sohnes, willen, du wollest mir aus Gnaden vergeben alle meine Sünde, so ich heut diesen Tag abermal in Gedanken, Worten und Werken wider dich gethan habe, und derselben in Ewigkeit nicht gedenken. Und weil ich mich nun will zur Ruhe niederlegen, wollest du mich heut diese Nacht vor allem Uebel Leibes und der Seelen auch väterlich behüten, daß ich vor des Teufels List und Gewalt, vor unnützen, schändlichen Träumen behütet, sicher ruhen und und zu deinem Lobe frisch und gesund wieder erwachen mög in deine gütliche Hand befahr.

\quellseite{47}

\section{27. Morgensegen am Dienstag}

\unsicher{unleserlich}et sei Gott, der Schöpfer Himmels und der Erde, gelob\unsicher{unleserlich} Herr, der allein Wunder thut, und gelobet sei sein herrs\unsicher{unleserlich} Name ewiglich. Der den Tag und die Nacht durch seine \unsicher{unleserlich} Weisheit geschaffen und unterschieden hat, dazu verordnet, \unsicher{unleserlich} lange die Erde steht, sie nicht sollen aufhören, auf daß der \unsicher{unleserlich} des Nachts möge seine Ruhe haben und am Tage wieder seine Arbeit gehen. Herr, wie sind deine Werke so groß und \unsicher{unleserlich} Du hast sie alle weißlich geordnet, und die Erde ist voll deiner \unsicher{unleserlich} Für solche Wohlthaten sollen wir dir danken, ehe die \unsicher{unleserlich} aufwacht, und vor dich treten, wenn das Licht hervorbricht. \unsicher{unleserlich} preise ich dich jetzt auch, daß du mich in dieser Nacht hast \unsicher{unleserlich} und schlafen lassen und wiederum fröhlich und gesund \unsicher{unleserlich} dazu vor aller des Feindes Gewalt und Bosheit väter\unsicher{unleserlich}schirmet. Ich bitte dich, du wollest meine Seele in deine \unsicher{unleserlich} einschließen, meinen Leib nach deinem Wohlgefallen gesund erhalten und mich vor allem Uebel und Schaden behüten. Sei du mein gewaltiger Schutz, meine große Stärke, mein Schirm wider die Hitze, meine Hütte wider den heißen Mittag, meine Hut wider \unsicher{unleserlich} Straucheln, meine Hilfe wider den Fall, auf daß mich kein \unsicher{unleserlich} treffe. O gütiger Gott, weil die Stunde auch da ist, aufzustehen vom Schlaf der Sünden und aller Ungerechtigkeit, sintemal unser Heil jetzt gar nahe ist, denn die Nacht ist vergangen, aber der Tag herbeigekommen: so hilf, daß wir ablegen die Werke der Finsterniß und anlegen die Waffen des Lichts, auf daß wir ehrbarlich wandeln als am Tage, nicht in Fressen und Saufen, nicht in Hader und Neid, sondern daß wir anziehen Jesum Christum in wahrem Glauben und christlichem Wandel. Derhalben wecke mich alle Morgen, wecke mir das Ohr, daß ich mit gläubigem Herzen \unsicher{unleserlich} höre, dasselbe in meinem Gedächtniß behalte, \unsicher{unleserlich} treten und flehen der Noth nicht verlasse. \unsicher{unleserlich} wenn ich \unsicher{unleserlich} zu dir rufe, so erhöre du \unsicher{unleserlich} Seufzen
\quellseite{48}
Neige deine Ohren zu meinem Schreien. Erfülle uns früh mit deiner Gnade, so wollen wir rühmen und fröhlich sein in dir unser Lebelang, durch denselben deinen lieben Sohn, JEsum Christum, unsern HErrn. Amen.

\section{28. Ein Anderer}

HErr JEsu Christe, du einiger Heiland der Welt! zu dir erhebe ich mein Herz, Muth und Sinn, und dancke dir abermal, daß du mich durch deine unerschöpffliche Barmherzigkeit und Liebe wider des bösen Feindes List und Gewalt frisch und gesund behütet hast. HErr JEsu Christe, du mein Gut und mein Erbtheil, mein Heil stehet in deinen Händen, ich weiß auch von keinem Helfer weder im Himmel, noch auf Erden, denn von dir allein. Darum bitte ich dich um deiner unaussprechlichen Marter, Angst und allerschmählichsten, bittersten Todes willen, den du, liebster HErr JEsu Christe, für mich armen Sünder aus großer Liebe hast gelitten, du wollest mir gnädig und barmherzig sein, und mich heut diesen Tag und die ganze Zeit meines Lebens segnen, fristen, behüten, und bewahren vor Sünden und vor allem Uebel von dieser mühseligen elenden Zeit an, bis du mich zur ewigen Freude und Seligkeit gnädiglich abforderst, um deines allerheiligsten Namens willen. Amen.

\section[29. Gebet aus der dritten Bitte des heil. …]{29. Gebet aus der dritten Bitte des heil. Vater Unsers}

Abba, lieber Vater! Ich erkenne in kindlicher Demuth mich schuldig, alles, was in und an mir ist, deinem göttlichen Wohlgefallen gehorsamlich zu unterwerfen, und deinen heiligen Willen in meinem Stande und ganzem Leben nach dem Exempel der auserwählten Himmelsbürger mit willigem Herzen zu erfüllen. Aber ach, du weißt auch, daß in mir und in meinem Fleisch nichts gutes wohnet. Wollen habe ich wohl, aber vollbringen das Gute (nach deinem Willen) finde ich nicht, denn das Gute, das ich will, das thue ich nicht, und das Böse, das ich nicht will, thue ich. HErr! deine allsehenden Augen erkennen, wie die alte Schlange, die da heißt der Teufel und Satanas, der die ganze Welt verführet durch seine listigen Anläufe; wie die Welt, die ganz im Argen liegt, mit ihrer Fleischeslust, Augenlust und hoffärtigem Leben,
\quellseite{49}
mich immerdar reizen, nach ihrem bösen Willen zu leben, und deinem allzeit guten Willen zu widerstreben. Demnach rechne mir, mein Gott, meinen Ungehorsam nicht zur Verdammniß zu, sondern sei mir gnädig nach deiner großen Güte, und tilge die Sünde meiner Widerspenstigkeit nach deiner unendlichen Barmherzigkeit.

Jesu, mein Jesu, der du deinem himmlischen Vater gehorsam warest bis zum Tode am Kreuze, ersetze du, was mir fehlet, aus dem Reichthum deines vollkommenen Gehorsams, und stehe mir in Kraft deines Verdienstes wider den Teufel und die Welt gewaltig bei, daß sie mich nicht mehr hindern mögen, den Willen des Allerhöchsten zu vollbringen. Heiliger Geist, gib mir ein neues Herz und einen neuen Geist, und mache hinfort einen ganz andern Menschen aus mir, der in Gottes Gebot wandele, seine Rechte halte und darnach thue.

Es ist, Herr, du getreuer Gott, dein väterlicher Wille, daß ich auch in Trübsalen, in Aengsten und Nöthen mich als deinen Diener beweisen soll. Mir hat mein Jesus gesagt: So ich sein treuer Liebhaber sein wolle, so müsse ich mich selbst verleugnen, täglich mein Kreuz auf mich nehmen, und ihm nachfolgen. Siehe, mein Geist ist willig, um deiner Ehre willen alles zu leiden, was dein Rath über mich beschlossen hat, aber das Fleisch ist schwach. Jedoch will ich nichts begehren wider deinen Willen, sondern opfre dir mich hiermit ganz und gar auf, und untergebe mich dir in Lieb und Leid, in Ehre und Schande, im Leben und Tod. Ich lege zu deinen Füßen mein Leib und Seele, mein Herz und Muth, meine Sinnen und Verstand, meine Kräfte und Vermögen, mein Hab und Gut, und alles was ich habe. Du bist der Herr, mache es mit mir, wie es dir gefällt, du wirst mir doch nicht mehr auflegen, als ich ertragen kann, und endlich alles zu meinem seligen Besten richten. Gib aber, nach dem Reichthum deiner Herrlichkeit, mir und allen deinen Kindern Kraft, von Tag zu Tag stärker zu werden an dem inwendigen Menschen, und so viel möglich ist, nach deinem Willen hier auf Erden zu leben, bis du uns in dein himmlisches Jerusalem, da die rechte Gerechtigkeit wohnet und dein Wille vollkommenlich beobachtet wird, einholen und versetzen wirst. Amen!

\quellseite{50}

\section{30. Abendsegen am Dienstag}

HErr, gütiger Gott, heiliger Vater! Des Tages schrey ich zu dir mit meiner Stimme; wenn mir Angst ist, rufe ich zu dir; und des Abends gedenke ich an deine Güte und Treue, die du mir erzeiget hast; und sonderlich preise ich dich jetzund, daß du mich aus lauter Gnade und Barmherzigkeit, ohne all mein Verdienst und Würdigkeit, heute diesen verschienenen Tag vor unzähligem Schaden und Gefahr, die mich hätten treffen können, so ganz väterlich behütet hast, daß mich der Satan nicht mit einem bösen, schnellen Tod umgebracht hat; daß ich von Ungewitter, Hagel und Blitz nicht bin versehrt worden; daß mich große Sturmwinde nicht beschädiget haben, und daß mich der böse Feind im Essen und Trinken mit Gift und andern seinen tausendkünftigen Listen nicht verderbet hat; daß ich auch keinen Arm oder Bein entzwei gefallen habe. Dafür sage ich dir jetzt und allezeit Lob und Dank, und bitte dich um des bittern Leidens JEsu Christi willen, du wollest mir vergeben alle Mißhandlungen, so ich heute wider dich gethan habe, und mich die zukünftige Nacht fürbaß gnädiglich bewahren vor meinem Widersacher, dem leidigen Teufel, vor Schrecken und Entsetzen des Nachts, daß mich kein Ungethüm noch Phantasie bethöre noch beschädige. Behüte mich mit Allem, was ich habe, vor Wassers- und Feuersnöthen, vor allem Uebel Leibes und der Seele. Laß mich geruhiglich ohne alle Sorge und Bekümmerniß einschlafen, auf daß ich auch in Finsterniß das Licht deines göttlichen Glanzes über mich scheinend möge sehen mit den Augen meines Herzens. Denn du bist ein helles und wahrhaftes Licht, welches da erleuchtet alle Dunkelheit, die uns umfangen hat. Du, HErr Gott, bist bei mir, du bist mein Fels, meine Burg, mein Erretter, mein Hort, auf den ich traue, mein Schild und Horn meines Heils und mein Schutz. HErr mein Gott, zu dir breite ich des Nachts meine Hände aus; komm zu mir, wie ein Spätregen, der das Land befeuchtet. Bleibe bei uns, HErr; denn der Tag hat sich geneiget, und ist niemand, der uns in dieser Finsterniß schütze, denn du allein, unser Gott. Eile mir zu helfen und schütze mich in dieser Nacht, daß mein Geist in Sünden nicht einschlafe und meinem Leibe kein
\quellseite{51}
Uebel widerfahre. Wecke mich zu rechter und bequemer Zeit und laß mich hören Freude und Wonne; denn ich habe Lust zu deinem Wort und Zeugnissen, die sind meines Herzens Trost. Gib, daß vor meine Ohren nicht komme eine traurige Botschaft und wende die Angst meiner Seele, denn du allein kannst mein Leben fristen und mich mit allerlei Segen erfüllen in Christo Jesu, unsrem HErrn. Amen.

\section{31. Ein Anderer}

O ewiger, barmherziger und reicher Gott vom Himmel, der du aus großer Gnade und väterlicher Fürsorge mich abermal diesen Tag frisch und gesund hast lassen vollenden, dafür danke ich dir billig von ganzem Herzen. Und weil ich leider diesen Tag nicht zu deinem Lob und Ehre und zu Nutz meines Nächsten zugebracht habe, sintemal meine verderbte Natur allezeit eher zum Bösen, als zum Guten geneiget ist, so hilf du getreuer Gott, daß ich, der ich aus sündlichem Samen gezeuget bin, meine Gebrechlichkeit erkenne, und deiner göttlichen Gnade theilhaftig werde, und lehre mich, daß ich oft und vielmal an mein Ende gedenke, daß ich mich zu demselben in wahrer Buße gefaßt mache, und wenn dasselbige herzu nahet, daß ich möge getrost und selig aus diesem Elend abscheiden, und mit allen gläubigen Christen in das himmlische Paradies versetzet werden. Indeß wollest du mich, so lang ich noch hie zu leben habe, in deinen allmächtigen Schutz nehmen, und vor allem Schaden und Gefahr an Seel und Leib gnädiglich behüten und bewahren um Jesu Christi, deines lieben Sohnes, willen. Amen.

\section{32. Morgensegen am Mittwoch}

Allmächtiger, barmherziger Gott! Nachdem alle deine Creaturen dich ohne Unterlaß loben und preisen sollen; welches auch thun die Vöglein unter dem Himmel, die früh gegen dem Tage mit ihren Zünglein und Stimmlein aufs lieblichste dich, ewiger Gott, als ihren HErrn und Schöpfer rühmen, so danke ich dir jetzt auch von Herzen, daß du mich die vergangene Nacht und alle vorige Zeit meines Lebens unter deinem Schutz und Schirm bis auf diese gegenwärtige Stunde erhalten hast und mich aus dem Schlafe und Finsterniß dieser Nacht erweckt, fröhlich und gesund hast lassen aufstehen.

\quellseite{52}

Ich bitte dich durch die heilwärtige Auferstehung Jesu Christi von den Todten, du wollest mich fürbaß auch behüten vor allem Unglück und Uebel, sammt allen, die mir mit Blutfreundschaft und sonsten verwandt sind. O Herr, hilf deinem Volk und segne dein Erbe und weide sie und erhöhe sie ewiglich! Du wollest mich auch in dieser Frühstunde mit deiner Barmherzigkeit erfüllen, daß ich heute mit Freuden diesen ganzen Tag in deinen Geboten ohne alle Todsünde zubringe. Erzeige mir deine Gnade wie eine Thauwolke des Morgens, und wie ein fruchtbarer Thau, der früh Morgens sich ausbreitet und das Land befeuchtet, also wollest du deine Güte über mich ausstrecken, mein träges Gemüth erfrischen, daß ich wacker und fröhlich deinen Willen thue. Regiere mich mit deinem heiligen Geist, daß ich dir mit reinem Herzen möge dienen in rechtschaffener Gerechtigkeit und Heiligkeit, die dir gefällig ist. Behüte mich, daß ich heute nicht wider dich sündige und etwa mein Gewissen beflecke mit fleischlichen Lüsten, welche wider die Seele streiten. Bewahre auch meine Zunge vor Bösem und meine Lippen, daß sie nicht falsch reden oder trügen. Schandbare Worte und Narrentheidung, die den Christen nicht geziemen, laß ferne von mir sein. Hilf, daß ich mit meinen Lippen niemand ärgere noch afterrede, verurtheile oder verdamme, auch niemand lästere oder schmähe. O daß ich könnte ein Schloß an meinen Mund legen und ein festes Siegel auf mein Maul drücken, daß ich dadurch nicht zu Fall käme und meine Zunge mich nicht verderbete! Gib Gnade, daß ich mich selbst in meinen Gebrechen erkenne und strafe, damit ich nicht in dein gestreng Urtheil und Gericht falle. Solches verleihe mir, ewiger Gott, durch Christum, deinen lieben Sohn. Amen.

\section{33. Ein Anderer}

O allmächtiger, barmherziger Gott, gnädiger Vater im Himmel, nachdem du abermal durch den Schutz deiner himmlischen Frohngeister, der heiligen Engel, mich diese Nacht also behütet hast, daß ich diesen Tag frisch und gesund habe erlebet, so sage ich dir dafür von Herzen Lob, Preis und ewigen Dank, und führe dich nochmals, lieber himmlischer Vater, in dessen Willen ich all mein Thun und Lassen, Anfang und Ende setze, du wollest dich meinerauch
\quellseite{53}
forthin gnädiglich erbarmen, und all’ mein Dichten und Trachten, Herz, Sinn, Muth und Gedanken, alle meine Worte und Werke, mit deinem Heiligen Geist regieren, daß ich verstehen möge, was gut oder böse ist, und daß ich in dieser bösen und verkehrten Welt heut also möge wandeln und leben, daß ich vor allem Dingen los von meinen Sünden, nach dem ewigen Vaterland, welches mir Christus, mein Heiland, erkauft und erworben hat, ein herzliches Verlangen trage, und dasselbe mit meinem bösen, sündlichen Leben nicht versch herze. Dazu wollest du mir gnädiglich helfen mit deiner göttlichen Liebe und Kraft des Heiligen Geistes, um JEsu Christi, deines lieben Sohnes willen. Amen.

\section[34. Gebet aus der vierten Bitte des heil. …]{34. Gebet aus der vierten Bitte des heil. Vater Unsers}

Abba, allergütigster Vater! du weißt alle Dinge, du siehest auch, daß ich deiner väterlichen Güte zu meiner Unterhaltung in diesem zeitlichen Leben alle Tage, ja alle Stunden und Augenblicke höchst bedürftig bin; so laß demnach dieselbe diesen und alle Morgen mit der Sonne neu über mir aufgehen. Dir ist am besten bekannt, wie viel mir in gegenwärtigem Jammertal von nöthen ist; darum wie die Augen der Knechte auf die Hände ihrer Herren, und die Augen der Mägde auf die Hände ihrer Frauen sehen, also sehen auch meine Augen auf dich, mein Gott! daß du aus deinen milden Segens-Händen mir alles zufallen lässest, was zu meines elenden Leibes Nahrung und Nothdurft gehöret. Du hast gesagt: Du wollest die, so auf dich trauen und auf deine Güte warten, nicht verlassen noch versäumen, solches wirst du auch an mir in Gnaden erfüllen, und mich keine Noth und Mangel bei deinem Dienst leiden lassen.

Aber zweierlei bitte ich von dir, das wollest du mir nicht weigern, ehe denn ich sterbe: Armuth und Reichthum gib mir nicht, laß mich nur mein bescheiden Theil Speise dahin nehmen; ich möchte sonst, wo ich zu satt würde, dich verleugnen und sagen, wer ist der HErr? oder wo ich zu arm würde, möchte ich stehlen, und mich an dem Namen meines Gottes vergreifen. Ich habe nichts mit in diese Welt gebracht, werde auch nichts mit hinaus nehmen; so hilf bewegen, daß ich an deinen Gaben und Segen
\quellseite{54}
mir allezeit wohl genügen lasse, und solche wieder anwende, deine Ehre zu befördern, mich mit Nahrung und Kleidern zu versorgen, und meinem armen Nächsten zu helfen und zu dienen, damit ich einst, wenn ich Rechenschaft hiervon thun soll, vor deinem Gerichte wohl bestehen möge.

O du Liebhaber des Lebens! friste mir und allen den lieben Meinigen das Leben in dieser Welt, so lange es dir gefällt; und erhalte uns bei guter Gesundheit, auf daß wir dir in unserm Stande allezeit ungehindert dienen, und mit vereinigtem Muth und Munde dich stets loben und ehren können. Du Gott des Friedens! verleihe immerdar den edlen Frieden zu unsern Zeiten, daß ein jeder unter seinem Weinstock und Feigenbaum sein Lebenlang sicher wohnen möge. Mache mit deinen heiligen Engeln eine feurige Mauer um unser ganzes Land her, daß kein Feind zu dessen Gränzen sich nahen dürfe. Ach Herr der Herrlichkeit, überschütte unsern Landesherrn und alle christliche Obrigkeit hier und an andern Orten mit deinem Segen aus der Höhe, und erfülle alle ihre guten und wohlgefälligen Anschläge. Gib ihnen langes und vergnügtes Leben, glückselige Regierung, getreue Diener und Räthe, fromme und gehorsame Unterthanen, mächtigen Sieg wider ihre Feinde, ja alles, was ihr Herze Gutes wünschet, und laß uns unter ihrem Schutz und Schirm ein stilles und geruhiges Leben führen in aller Gottseligkeit und Ehrbarkeit.

O Vater aller Gnaden! beschere allen Hausvätern und Hausmüttern glücklichen Fortgang zu ihrer Nahrung. Segne die Frucht ihres Leibes, daß sie sehen, wie ihre Kinder zunehmen an Weisheit, Alter und Gnade bei Gott und den Menschen; segne die Frucht ihres Landes, und thue auf die Fenster des Himmels, daß Früh-Regen und Spät-Regen herab komme zu rechter Zeit, und behüte die Ernte treulich und jährlich; segne die Frucht ihres Viehes, daß sie ihren Acker wohl bauen, ihrer sauren Mühe genießen, und ihr Brod mit gutem Vergnügen essen können. Segne die Werke ihrer Hände, daß sie nicht umsonst arbeiten, weil sie sind der Saame der Gesegneten des HErrn. Gesegnet laß mich und alle deine Christen sein mit allerlei geistlichem und leiblichem Segen, in irdischen und himmlischen Gütern, durch JEsum Christum unsern HErrn. Amen.

\quellseite{55}

\section{35. Abendsegen am Mittwoch}

O du heilige Dreifaltigkeit in Einem göttlichen Wesen, du bist mein Leben, Heil und ewiger Trost! Dir sage ich mit Mund und Herzen Lob und Dank, daß du mich diesen Tag über gnädiglich behütet hast. Ich bitte deine göttliche Güte, du wollest alle meine Missethat bedecken, allermeist aber, was ich heute wider dich und deine heiligen Gebote mit meiner Zunge, mit unnützen, vergeblichen Worten, mit Afterreden und sonst gesündigt habe, und wollest mich diese zukünftige Nacht auch bewahren vor allem Schaden und Gefahr, denn zu dir allein habe ich all mein Vertrauen. Gott, wie dein Name ist, so ist auch dein Ruhm bis an der Welt Ende. Deine Rechte ist voll Gerechtigkeit. Darum befehle ich meinen Leib und Seele in deine Hände. So gebenedeye mich nun die göttliche Majestät und beschirme mich die heilige Dreifaltigkeit und erhalte mich die ewige Einigkeit. Es beschütze mich die unermäßliche Barmherzigkeit, es vertheidige mich die unaussprechliche Mildigkeit, es erfreue mich die unendliche Süßigkeit, mich bedecke die höchste Wahrheit Gottes, mich stärke die tiefste Erkenntnuß Christi, mich bewahre die grundlose Gütigkeit des HErrn. Die Gnade des Vaters regiere mich, die Weisheit des Sohnes erquicke mich, die Kraft des heiligen Geistes erleuchte mich. Mein Schöpfer stehe mir bei, mein Erlöser helfe mir, mein Tröster belohne mir. Der HErr segne mich und behüte mich, der HErr erleuchte sein Angesicht über mich und sei mir gnädig, der HErr erhebe sein Angesicht auf mich und gebe mir Frieden. Dieser Schutz und Segen der einigen Gottheit sei heute und allezeit zwischen mir und allen meinen Feinden, sichtbaren und unsichtbaren, daß sie sich zu mir nicht können machen, noch mich beschädigen. Gleichwie die Wolkensäule in der Wüste sich machte zwischen das Heer der Egypter und das Heer Israel, daß diese und jene nicht konnten zusammen kommen, damit den Kindern Israel kein Leid widerführe: also wollest du zwischen mir und allen meinen Feinden eine Feuermauer und Unterschied sein, daß sie mich nicht berühren. Erhalte mich auch an meinem letzten Stündlein. Wenn meine Augen nimmer sehen, meine Ohren nimmer hören und meine Zunge nimmer redet,
\quellseite{56}
wenn meine Hände nimmer greifen und die Füße nimmer gehen mögen, so stehe bei mir, du hochgelobte Dreifaltigkeit, daß der böse Feind keine Macht an mir finde. Amen.

\section{36. Ein Anderer}

Ich danke dir, o großmächtigster Gott und gütigster Vater, daß du mich heute abermal diesen Tag über an Leib und Seele für allem Schaden und Uebel gnädiglich behütet hast. Ich bitte dich durch deine göttliche Güte, du wollest alles, was ich heute mit Herz und Mund und sonst wider dich und deine heiligen Gebote gesündiget habe, zudecken, und mich, dein Geschöpf, welches du mit dem Blut deines lieben Sohnes, Jesu Christi, theuer erkauft hast, nicht verlassen; und mich diese Nacht unter deinen allmächtigen Gnadenflügeln wider den bösen Feind, welcher um mich hergehet, wie ein brüllender Löwe, mich zu verschlingen, beschützen, daß ich unter dem Schatten deiner Güte und Barmherzigkeit sicher ruhen und schlafen möge, auf daß sich der böse Feind nicht zu mir dürfe nahen, und mir keinen Schaden thun könne. Amen.

\section{37. Morgensegen am Donnerstag}

Herr Jesu Christe! der du bist das ewige, wahre Licht, welches da vertreibet die Finsterniß des Nachts und Schatten des Todes, deinen Namen will ich rühmen, dir will ich lobsingen und danken, daß du mich in dieser Nacht so gnädiglich behütet und aus der Finsterniß an das Licht gebracht hast. Mich hast du bewahrt vor dem Grauen des Nachts, vor des Teufels Schrecken und Gespenstern, vor der schädlichen Pestilenz, die im Finstern schleicht, vor mancherlei Seuchen und Krankheiten, die mir hätten können widerfahren. Auch hast du meine Seele mit deinem Schild umringt und bewacht, wie ein Hirte seine Heerde bewacht. Dazu alles, was ich habe, ist durch deine große Barmherzigkeit unversehrt behütet worden. Für solchen gnädigen Schutz und alle deine Wohlthaten sei dir Lob und Preis gesagt; von deiner Macht will ich reden und des Morgens deine Güte rühmen, denn du bist meine höchste Zuversicht, meine feste Burg, meine starke Hülfe, mein treuer Gott, auf den ich traue, du erfreuest mein Herz
\quellseite{57}
und machst mir fröhlich das Angesicht. Ich bitte dich durch deine heilige Geburt und Menschwerdung, du wollest an diesem Tage deine Barmherzigkeit über mich lassen aufgehen und hervorbrechen wie die schöne Morgenröthe und zu mir kommen wie der Frühregen. Erleuchte meine blinde Natur und verdunkelt Herz mit deinem Glanz, auf daß in meinem Herzen aufgehe der rechte Morgenstern und das wahrhaftige Licht, welches erleuchtet die Menschen zum ewigen Leben. Behüte mich auch heute vor allem Uebel. Sei mir gnädig, Herr, denn auf dich harre ich, meine Seele wartet auf dich von einer Morgenwache bis zur andern. Sei du mein Arm früh, dazu mein Heil zur Zeit der Trübsal. Beschirme mich an Leib und Seele, daß mir kein Uebel begegne und keine Plage zu mir nahe. Treibe fern von mir alle bösen Geister, stehe du bei mir wider die Boshaftigen, tritt zu mir wider die Uebelthäter und schütze mich, daß die Hand der Widermärtigen, mich nicht berühre. O Herr, unser Gott, fördere das Werk unserer Hände bei uns, ja das Werk unserer Hände wollest du fördern und unsere Arme stärken, auch unsere Finger lehren halten deine Gebote, daß wir heute nicht wider dich sündigen. Solches verleihe uns um deiner Barmherzigkeit willen, welche für und für währet in Ewigkeit. Amen.

\section{38. Ein Anderer}

Gott Vater, Gott Sohn, Gott Heiliger Geist, du hochgelobte Dreifaltigkeit! dir ergebe ich mich mit Leib und Seele von nun an bis in Ewigkeit, und sage dir herzlich Lob und Dank, daß du dem bösen Feinde nicht hast zugelassen, daß er mich diese Nacht hat können beschädigen, sondern durch den Schutz deiner lieben Engel bin ich frisch und gesund behütet worden. Was soll ich dir vergelten, womit soll ich dich dafür loben? Ich will dir ein geängsteter und zerschlagenes Herz geben voller blutrothen Sünden, mit Reue und Leid, daß wollest du gnädiglich annehmen, mit dem edlen Blut deines lieben Sohnes, meines Erlösers, schneeweiß waschen und in seine heiligen, unschuldigen Wunden verbergen und also mir Vergebung aller meiner Sünden gnädiglich widerfahren lassen. Und hilf, daß ich heute und alle Tage in christlicher Bereitschaft sei, denn ich nicht wissen kann, wann du kommen, wie und
\quellseite{58}
wo du mich von hinnen abfordern wirst, daß ich zur ewigen Freude seliglich geführet werde. Solches verleihe mir, gnädiger Gott und Vater, um deines lieben Sohnes, JEsu Christi willen. Amen.

\section[39. Gebet aus der fünften Bitte des heil. …]{39. Gebet aus der fünften Bitte des heil. Vater Unsers}

Abba, barmherziger Vater, ich bekenne mit reuigem Herzen, daß ich deiner Wohlthaten nicht würdig bin, keine Gnade, sondern Zorn und Strafe verdienet habe, weil ich lauter Uebels vor dir gethan und mich zeitlebens an deinen heiligen Geboten mit unzähligen schweren Sünden verschuldet. Meiner Missethaten, womit ich dich beleidiget, sind mehr denn Haare auf meinem Haupt, mehr denn Sandes am Meer. Meine Schuld, womit ich dir verhaftet bin, ist groß bis in den Himmel und will als eine unerträgliche Last mich in die Hölle drücken. Wolltest du mit mir nach deiner strengen Gerechtigkeit handeln, so müßtest du mich von deinem Angesicht zu den abscheulichen Teufeln an den Ort der ewigen Qual verstoßen. Aber, ach HErr, habe Geduld mit mir, und deines Rechtes Strenge mit Gnad’ vermenge. Ach Gott, Gott, sei dem armen Sünder gnädig, der vor dir fußfällig lieget und seine Zuflucht zu deiner väterlichen Liebe nimmt! Erinnere dich des theuern Eides, den du geschworen hast: So wahr ich lebe, ich habe keinen Gefallen an dem Tode des Gottlosen, sondern daß er sich bekehre und lebe. Gedenk’ an deines Sohnes bittern Tod, sieh’ an seine heiligen Wunden roth, die sind ja für die ganze Welt die Zahlung und das Lösegeld; deß tröste ich mich allezeit und hoffe auf Barmherzigkeit.

Ach JEsu, mein Erlöser und Seelenbräutigam, nimm dich doch meiner armen sündigen Seele herzlich an und laß sie nicht in Sünden sterben und verderben. Durchstreiche das schwarze Sündenregister und die scharfe Handschrift des Gesetzes, so wider mich ist, mit deinem rosinfarbenen Blut, mit dem Blut, wodurch du für mich bezahlet hast, auf daß meiner Sündenschulden vor Gottes Gerichte und Gerichte nicht mehr gedacht werde. Um deiner vollgültigen Genugthuung willen, womit du den gerechten Gott an meiner Statt vergnüget hast, wende alle zeitlichen und ewigen Strafen meiner Missethaten von mir gnädiglich ab und schenke mir deine theuer erworbene Gerechtigkeit und die ewige Seligkeit.

\quellseite{59}

Aber, erneuere hinfort mein Herz, o Gott heiliger Geist, daß ich die schweren Sündenschulden nicht auf’s neue aufwachsen lasse oder mit Unbußfertigkeit den gerechten Zorn Gottes über meine Seele häufe, sondern vielmehr mein Lebenlang des Herrn Wort vor Augen habe und in seine Sünde willige wider sein Gebot. Wo ich aber aus menschlicher Schwachheit durch Verführung des Satans und meines Fleisches und Blutes sündige, so ermuntere mich unverzüglich zu wahrer Buße und hilf, daß ich die Vergebung meiner Uebertretung nicht nur demüthigst suchen, sondern auch gewiß finden möge. Reiß aus meinem Gemüthe alle Feindseligkeit und Unversöhnlichkeit und entzünde in demselben eine rechte brünstige Liebe gegen meinen Nächsten, daß ich meinen Brüdern und Schwestern von Herzen ihre Fehler vergebe, segne, die mir fluchen, wohlthue denen, die mich hassen, für die bitte, die mich beleidigen und verfolgen, und so als ein rechtes gehorsames Kind meines Vaters im Himmel mich bezeige.

Erleuchte und bekehre aber auch alle anderen Sünder und Sünderinnen, alle ruchlosen, verstockten und unbußfertigen Seelen, daß sie ihre großen Sünden rechtschaffen erkennen, den feuerbrennenden Zorn Gottes fürchten, vor der unendlichen Höllenpein erschrecken, ihre Missethaten sich herzlich leid sein lassen, das blutige Verdienst Jesu in wahrem Glauben ergreifen, ihr gottloses Leben ändern und zur ewigen Seligkeit erhalten werden mögen, um Christi Blutes und Wunden willen. Amen.

\section{40. Abendsegen am Donnerstag}

Gelobet sei Gott, der Vater, durch Jesum Christum im Heiligen Geist, ein einiger, ewiger Gott, der durch seine mannigfaltige Güte mich armen Sünder und elenden Menschen heut diesen Tag gnädiglich bewahret hat vor allen feurigen Pfeilen des Satans, die des Tages fliegen, vor der Seuche, die im Mittag verderbet, dazu vor dem jähen schnellen Tod und vor allem Schaden väterlich behütet. Herr, deine Güte reicht, so hoch der Himmel ist und deine Wahrheit so weit die Wolken gehen. Du bist gnädig und barmherzig, alle deine Werke sind löblich. Ich bitte dich, mildreicher Gott, du wollest mir aus Gnaden verzeihen alles, was
\quellseite{60}
ich heute wider dich gethan habe, es sei mit Werken, Worten und Gedanken. Wollest auch deine Barmherzigkeit zu mir richten, und mich die zukünftige Nacht lassen einschlafen und ruhen, daß ich dich, der du bist die ewige Ruhe, nun und nimmermehr verlasse, sondern in dir bleibe durch den Glauben und unter deinem Schirm sicher wohne, auf daß sich der böse Feind nicht dürfe zu mir nahen und mir keinen Schaden könne zufügen. Herr, du bist mein Licht und mein Heil, vor wem sollt ich mich fürchten? Du bist meines Lebens Kraft, vor wem sollte mir grauen? Auf dich verlässet sich mein Herz und mir ist geholfen. Du bist mein Trost und gewaltiger Schutz, deine rechte Hand stärket mich, deine Rechte tröstet mich und unter dem Schirm deiner Arme habe ich Zuflucht. Siehe, mein Gott, des Tages rufe ich, so antwortest du mir, und des Nachts schweige ich auch nicht und du erhörest mich. Wenn ich mich zu Bette lege, so denke ich an dich, wenn ich erwache, so rede ich von dir, denn du bist mein Helfer und unter dem Schatten deiner Flügel rühme ich, meine Seele hanget dir an, deine Rechte erhält mich. Wenn ich im Finstern sitze, so ist doch der Herr mein Licht und mein Heil. O gütiger Gott, verleihe mir Gnade, daß, wenn mein Sterb-Stündlein herzu nahet und ich in’s Todtbette zur ewigen Ruhe mich soll niederlegen, ich durch deine Hülffe im rechten festen Glauben getrost und unverzagt möge seliglich zum ewigen Leben einschlafen. Indes enthalte mich dir, daß ich alle Zeit wache, nüchtern und mäßig lebe und in christlicher Bereitschaft erfunden werde, sintemal ich nicht wissen kann, zu welcher Stunde du, unser Gott, kommen wirst und mich von hinnen abfordern, auf daß ich würdig werde, zu stehen vor des Menschen Sohn und nicht zu Schanden werde in seinem Gericht, der mit dir lebet und regieret in Ewigkeit. Amen.

\section{41. Ein Anderer}

Barmherziger, sanftmüthiger Gott, ewiger Vater, wie trägst du doch so eine herzliche Liebe und väterliche Sorge für mich armen Sünder, indem du mich alle Tage und Stunden von meiner Jugend an bis auf gegenwärtige Zeit für allerlei List des Teufels und der gottlosen bösen Welt Gefahr und Schaden
\quellseite{61}
behütet hast; und bitte dich demüthiglich, du wollest alles dessen, was ich heute wider dich gethan habe, nach deiner väterlichen Liebe, so du gegen mich hast, nun und in alle Ewigkeit nicht mehr gedenken, sondern aus Gnaden erlassen, und mir um deines lieben Sohnes, Jesu Christi, willen, welcher für alle meine Sünde ist Bürge worden, gnädig sein, und diese Nacht mich und alle die Meinen vor einem bösen schnellen Tod, vor Feuers- und Wassernoth, Pestilenz und vor allem Unglück sicher behüten. Darum befehle ich mich, mein Leib und Seel und alles, was ich habe, in deinen väterlichen Schutz; dein heiliger Engel sei mit mir, daß ich kein Unglück fürchte. Amen.

\section{42. Morgensegen am Freitag}

Gebenedeiet sei Gott, mein Schöpfer, gebenedeiet sei Gott, mein Heiland, gebenedeiet sei Gott, mein höchster Trost, der mir Leib, Gesundheit, Leben und Segen, der mein Schutz und Hülfe ist, so mich an Leib und Seele nach seiner großen und hochberühmten Barmherzigkeit in dieser vergangenen Nacht vor mancherley Beschädigung des Satans behütet und gesund an diesen Tag hat kommen lassen. Ich bitte dich, himmlischer Vater, durch den blutigen Schweiß Jesu Christi, deines lieben Sohnes, du wollest mich heute diesen Tag auch in deinen göttlichen Schutz nehmen, mich beschirmen und handhaben innerlich und äußerlich, daß mir kein Arges widerfahre. Denn in deine Hände thu’ ich befehlen heut’ und alle Tage meine arme Seele, meinen elenden Leib, mein dürftiges Leben, meine Sinne, Vernunft, Verstand und Anschläge, alle meine Gedanken, Worte und Werke, mein Thun und Lassen, meinen Eingang und Ausgang, mein Gehen und Stehen, mein Sitzen und Liegen, meinen Willen und Rathschlagen, mein Dichten und Trachten, meinen Glauben und Bekenntniß, und was ich auswendig und inwendig bin und vermag, das Ende meines Lebens, den Tag und die Stunde meines Todes, mein Sterben und Auferstehen. O Herr Gott, schaffe du es mit mir, wie du willst und weißest, was zu deiner Ehre und zu meiner Seligkeit das Beste ist. Erhalte mich in deiner Furcht und wahrer Erkenntniß. Behüte mich vor den Werken der Ungerechtigkeit, und so ich etwa wider dich aus
\quellseite{62}
Gebrechlichkeit würde sündigen, so bitte ich dich, du wollest deine Barmherzigkeit nicht von mir nehmen, deine Gnade nicht von mir wenden, deine Hülfe mir nicht entziehen; denn es ist sonst kein anderer Gott noch Helfer, vor dir ist kein anderer gemacht, so wird auch nach dir kein anderer sein. Du, du bist der Erste und der Letzte und außer dir ist kein Gott! Darum rufe ich zu dir allein, laß deine Güte über mich walten, laß mich früh hören deine Gnade, denn ich hoffe auf dich. Leite meinen Fuß auf rechter Straße, daß ich nicht wandle im Rath der Gottlosen, noch trete auf den Weg der Sünder, noch sitze, da die Spötter sitzen; sondern daß ich all meines Herzens Lust und Liebe habe zu deinem Wort und Geboten und in denselben mich übe Tag und Nacht. Durch unsern Herrn JEsum Christum, der mit dir und dem Heiligen Geiste lebet und regieret in Ewigkeit. Amen.

\section{43. Ein Anderer}

In deinem Namen, du gekreuzigter HErr JEsu Christe, bin ich armer Sünder jetzt aufgestanden, der du für mich am Stamm des heiligen Kreuzes, als das rechte geduldige Schlachtslämmlein, den allerschmerzlichsten Tod erlitten, und mich mit deinem rosinfarbenen Blut von allen meinen Sünden, Tod, Teufel und Hölle erlöst hast: regiere mein Herz durch deinen Heiligen Geist, erfrische es mit dem himmlischen Thau deiner Gnade, bewahre mich mit deiner göttlichen Liebe heute diesen Tag und verbirg mich mit Leib und Seele in deine heiligen Wunden; wasche mich von allen meinen Sünden rein ab, und erhalte mich in allen guten Werken, und führe mich aus dem Jammertal dieser Welt in die ewige Freude und Herrlichkeit, du getreuester Heiland JEsu Christe, mein einiger Trost, Hoffnung und Leben. Amen.

\section[44. Gebet aus der sechsten Bitte des heil. …]{44. Gebet aus der sechsten Bitte des heil. Vater Unsers. (Haas.)}

\unsicher{[Redaktioneller Hinweis: Die Überschrift zu Gebet 44 ist aus dem Originalinhaltsverzeichnis übernommen. Der Gebetsanfang fehlt im Scan; der erhaltene Text folgt auf der nächsten Scan-Seite.]}

\quellseite{63}

liche Sünden und Laster verführen und zuletzt in die endliche Verzweiflung und das ewige Verderben stürzen möge. Ach, schütze doch mit deiner starken Hand das Volk deiner Rechten und bewahre mich, dein Kind, und alle frommen Herzen, daß wir nicht in solche Anfechtung und Versuchung des höllischen Bösewichts fallen. Willt du, aber nach deinem gerechten Gerichte dem Satan über uns einige Gewalt verstatten, so laß uns nur nicht versucht werden über Vermögen, sondern mache, daß die Versuchung so ein Ende gewinne, daß wir’s ertragen können.

JEsu, mein JEsu! du bist erschienen, des Teufels Werke zu zerstören, und hast uns zu gute den Satan überwunden und zu deinen Füßen geleget; darum sei unser Helfer und beschirme uns unter dem Schatten deiner Flügel, bitte für uns droben im Himmel, daß unser Glaube nicht aufhöre oder unterliege, sondern durch dein Blut den Sieg erhalte.

Heiliger Geist! wenn das böse Stündlein kömmt, so ziehe uns an den Harnisch Gottes, daß wir bestehen können wider die listigen Anläufe des Teufels; rüste uns mit dem Schild des Glaubens, daß wir auslöschen können alle feurigen Pfeile des Bösewichts; lehre uns führen das Schwert des Geistes, welches ist das Wort Gottes, daß wir alles wohl ausrichten und das Feld behalten mögen.

Stehe aber indessen, o dreieiniger Gott! in Gnaden bei allen geängsteten Seelen, welche in dem Anfechtungsstiebe des Satans allzu sehr sieden und schwitzen. Neige deine Ohren zu ihrem Schreien und Flehen und schweige nicht zu ihren heißen Thränen. Sei mit deiner Kraft in ihrer Schwachheit mächtig und laß das glimmende Docht ihres Glaubens nicht auslöschen. Tröste sie einstens zu rechter Zeit wieder mit deiner gnädigen Hülfe, und dein freudiger Geist erhalte sie. Amen.

\section{45. Abendsegen am Freitag. (Habermann.)}

\unsicher{[Redaktioneller Hinweis: Die Überschrift zu Gebet 45 ist aus dem Originalinhaltsverzeichnis übernommen. Der Gebetsanfang fehlt im Scan; der erhaltene Text folgt auf der nächsten Scan-Seite.]}

\quellseite{64}

zu mir; wenn ich bete, merket er auf meine Stimme. Der HErr ist meine Zuversicht und Stärke, meine Hülfe in den großen Nöthen, die mich troffen haben. Darum sage ich dir, ewiger Gott, Lob und Preis, daß du mich heute diesen Tag väterlich behütet hast vor allem Unfall und Schaden, so mir hätte können begegnen. Mein Herz ist fröhlich und meine Seele preiset dich um alle deine Güte und Barmherzigkeit. Meine Zunge soll ihr Gespräch von dir haben und immer sagen: „Hochgelobet sei Gott, gesegnet sei sein heiliger Name.“ Ich bitte dich, du wollest mir aus Gnaden nachlassen alles, was ich heute wider dich gesündigt habe, und mich diese zukünftige Nacht sammt allem, was mir zustehet, auch schützen. Sei du mein Schild und mein Schatten über meiner rechten Hand. O HErr, behüte mich vor allem Uebel, behüte meine Seele, sei mir gnädig, denn auf dich allein traue ich. Ich hoffe auf den HErrn und rufe zu Gott, dem Allerhöchsten, zu Gott, der meines Jammers ein Ende macht. Siehe, der mich behütet, schläft nicht. Siehe, der Hüter Israel schläft noch schlummert nicht, er wird meinen Gang erhalten auf rechter Bahn, daß ich nicht strauchle und meine Tritte nicht wanken; er wird meinen Fuß nicht gleiten lassen, denn sein Wort ist ein Licht auf meinen Wegen. Darum wenn ich mich lege, so werde ich mich nicht fürchten, sondern süße schlafen, und werde mich nicht fürchten vor plötzlichem Schrecken, noch vor dem Sturm der Gottlosen, wenn er kommt; denn du behütest meinen Fuß, daß er nicht gefangen werde, und errettest mich vor den Stricken des Todes. O HErr Gott, erhebe über mich das Licht deines Antlitzes, auf daß ich ruhe und schlafe ganz mit Frieden und sicher wohne unter deinem Schirm, denn du allein hilfst mir. Auf deinen Namen will ich mich jetzt zur Ruhe niederlegen und meine Augenlider lassen schlummern. Du, HErr Gott, wirst mich fröhlich wieder erwecken zu Lob und Ehre deiner göttlichen Majestät. Durch JEsum Christum, deinen lieben Sohn, unsern HErrn, der mit dir lebet und regieret in Einigkeit des Heiligen Geistes. Amen.

\section{46. Ein Anderer}

O HErr JEsu Christe, du geduldiges Schlachtlämmlein und heiliges Versöhnopfer für alle meine Sünden, nicht allein aber für
\quellseite{65}
meine, sondern auch für der ganzen Welt Sünde! dir danke ich abermal von Grund meines Herzens, daß du mich an Leib- und Seele durch deinen Schutz diesen Tag so väterlich behütet hast; und bitte dich, du wollest mir alle meine Sünde, so ich heute diesen Tag aus Schwachheit meiner verderbten Natur und Anreizung des bösen Geistes gethan habe, welche mein Herz und Gewissen sehr beschweren und drücken, gnädiglich verzeihen und vergeben; und weil ich mich nun zur Ruhe und Schlaf will niederlegen, wollest du deine Gnadenflügel über mich ausbreiten, und helfen, daß ich darunter in Friede und Ruhe diese Nacht mit dem Leibe schlafe, mit der Seele aber allezeit zu dir wache und deiner herrlichen Zukunft zum Gericht wahrnehme, und mit herzlichem Seufzen auf dich warte, bis ich dermaleinst gar zu dir seliglich von hinnen fahre.

Dazu hilf mir, o treuer Gott, mit deinem lieben Sohn und Heiligem Geist hochgelobet in alle Ewigkeit. Amen.

\section{47. Morgen segen am. Sonnabend}

O du wahrer, unsterblicher Gott, Vater unsers HErrn JEsu Christi! zu dir erhebe ich mein Gemüth mit schuldiger Dankbarkeit, deine Gerechtigkeit will ich nicht verbergen in meinem Herzen, von deiner Wahrheit und von deinem Heil will ich reden, ich will nicht verhehlen deine große Güte und Treue vor der großen Gemeinde, und alles, was du mir Gutes gethan hast, will ich nicht verschweigen. Denn es ist ein köstlich Ding, dem HErrn danken und lobsingen deinem Namen, du Höchster; des Morgens deine Gnade und des Abends deine Wahrheit verkündigen. Darum preiset dich meine Seele, daß du mich in dieser Nacht durch deine überschwängliche Barmherzigkeit beschützt hast. Gesegnet bist du HErr Zebaoth, der du dich gnädig erzeigest allen, die nach dir fragen und dein Heil lieben! Gesegnet ist dein Name in allen Landen, der unser Schutz und Hülfe ist. Gesegnet sind alle deine Werke, die du an den Menschenkindern thust. Ich bitte dich, du wollest mich heute diesen Tag auch behüten, daß mir der böse Feind keinen Schaden zufüge und die Hand der Gottlosen mich nicht berühre. HErr Gott, mein Heiland, früh wache ich zu dir, früh rufe ich zu dir, hilf, daß ich die Werke meines Berufs und was mir befohlen
\quellseite{66}
ist, fleißig und treulich ausrichte, zu deinem Lob und meines Nächsten Besserung, damit ich das Licht dieses Tages und deine Creaturen nicht mißbrauche zur Sünde oder zur Eitelkeit, dich nicht beleidige mit meinem Thun und Lassen und den Bund meiner heiligen Taufe nicht übertrete. Verleihe mir auch Gnade, daß ich mich hüte vor den sechs Stücken, die du hassest und vor dem siebenten, darum du einen Greuel hast, als da sind hohe Augen, falsche Zunge, Hände, die unschuldig Blut vergießen, ein Herz, das mit bösen Tücken umgeht, Füße, die behende sind, Schaden zu thun, falscher Zeuge, der frech Lügen redet, und der Hader zwischen Brüdern anrichtet. Vor solchen und dergleichen Lastern bewahre mich, mein Gott, daß ich nimmermehr darein gerathe oder bewillige, sondern lehre mich thun nach deinem Willen, denn du bist mein Gott und Herr, dein guter Geist führe mich auf ebener Bahn, daß ich dir diene in unsträflichem Wandel, dazu all mein Thun und Leben dir gefalle in Christo Jesu, deinem lieben Sohn, unserm Herrn, der mit dir lebet und regieret in Ewigkeit des Heiligen Geistes von Ewigkeit zu Ewigkeit. Amen.

\section{48. Ein Anderer}

Ach du getreuer Vater im Himmel! daß ich dich genugsam sollte loben und dir könnte danken für alle Wohlthaten, so du mir die Zeit meines Lebens bis auf diese Stunde väterlich erzeiget hast. Das ist in meinem Vermögen und Kräften nicht. Denn ich bin Fleisch und Blut, welches nichts denn Böses thun kann. Du aber lässest mir täglich über die Maaßen viel Gutes widerfahren; und sonderlich, wo du in dieser Nacht nicht wärest mein Schild und Beistand gewesen, so hätte mich des Teufels Gewalt vielfältig beschädigt, daß ich nicht gesund hätte wieder aufstehen mögen. Aber durch deinen gnädigen Schutz bin ich unversehrt behütet worden. Ich bitte dich ganz inniglich; du wollest mir an diesem Tag deine Gnade widerfahren lassen, und mich, dein durch Christi Blut erworbenes Gut, fortan bis in’s ewige Leben gnädiglich behüten. Amen. Amen. Herr Jesu, nimm meine Seele in deine Hände, und laß sie dir befohlen sein. Amen.

\quellseite{67}

\section[49. Gebet aus der siebenten Bitte des heil. …]{49. Gebet aus der siebenten Bitte des heil. Vater Unsers und dem Beschluß}

Abba, getreuer Vater! Du bist der Erlöser und Nothhelfer derer, die auf dich trauen, und hast gesagt: Rufe mich an in der Zeit der Noth, so will ich dich erretten. Mein Herz hält dir vor dein Wort, da du befohlen hast: Kehre dich zu mir, denn ich erlöse dich. Ach, so erhöre demnach die Stimme meines Flehens und verachte mein Gebet nicht, welches ich anjetzo um die gnädige Erlösung aus allen Nöthen zu dir abschicke. Wende doch von mir und den Meinigen gnädiglich ab alles Böse, so uns an der Seelen Seligkeit schaden kann. Vor aller Sünde und Ungerechtigkeit, vor des Teufels List und Macht, vor Verlockung und Unbußfertigkeit, vor Unglauben und Verzweiflung, vor einem bösen schnellen Tod und der ewigen Höllengluth behüte uns, lieber Herr und Gott! Beschütze unsern Leib und Leben vor der Pestilenz, die im Finstern schleicht, und andern ansteckenden Krankheiten, vor Krieg und Blutvergießen, vor Theuerung und Hungersnoth und allen Plagen, welche uns an unserer zeitlichen Wohlfahrt hinderlich sein mögen. Rette unser Hab und Gut von Feuers- und Wassersnoth, Hagel und Ungewitter, Raub und Diebstahl und anderm verderblichen Unglück, so desselben uns berauben könnte. Bewahre endlich auch unsere Ehre und guten Namen vor Lügen und Verleumdung, vor öffentlichem Spott und Schande, vor Lästerung und Verachtung und aller Schmach, so denselben kränken mag.

Gefällt dir aber, mich oder die Meinigen mit Kreuz und Elend heimzusuchen und zu züchtigen, so gib, daß ich nicht murre wider deine Züchtigung, sondern die väterliche Ruthe demüthigst küsse; nicht kleinmüthig werde unter den Trübsalen, sondern dieselben freudig ertrage; nicht verzage über den Verzug deiner Hülfe, sondern darauf in stiller Geduld warte. Stärke durch deine Kraft aus der Höhe meines Fleisches Blödigkeit, erquicke mit deinen süßen Tröstungen meine bekümmerte Seele, nimm zu rechter Zeit die aufgelegte Last wieder von meinen Schultern hinweg und laß mir nach dem trüben Unglückswetter endlich die Sonne deiner Gnade helle scheinen.

\quellseite{68}

Mein Gott! dieser heutige letzte Tag in der Woche erinnert mich, daß mein Leben ein Ziel habe und ich davon müsse. Wenn nun der Tag und die Stunde meines Todes kommt, so verleihe, daß ich durch ein vernünftiges und seliges Ende allen Jammer dieses Lebens beschließe, im Glauben auf das blutige Verdienst Jesu sanft einschlafe und aus diesem Thränenthal zu deiner ewigen Freude und himmlischen Ruhe eingehen möge. O Jesu, du blutiges Gotteslamm, erscheine mir in dem Bilde bei meiner letzten Noth, wie du am Kreuz so milde dich hast geblut’t zu Tod’. Stärcke mich mit deinem Freudengeist, heile mich mit deinen Wunden, wasche mich mit deinem Todesschweiß in solcher letzten Stunde und nimm mich so, wenn dirs gefällt, in wahrem Glauben von der Welt zu deinen Auserwählten. Sei, o werthester Heiland, mein Fürsprecher am Tage des strengen Gerichts Gottes und hilf mir auch zur selben Zeit von wegen deiner Wunden, daß ich im Buch der Seligkeit werd eingezeichnet funden; daran ich denn auch zweifle nicht, denn du hast ja den Feind gericht’t und meine Schuld bezahlt. Verlaß mich nicht, Gott Heiliger Geist, bis an meinen letzten Lebensblick, sondern wohne mir beständig bei mit deinem kräftigen Trost wider alle Anfechtungen des bösen Geistes, sei meine Stärke in der größten Schwachheit und gib, daß meine letzte Speise in dieser Welt sei das heilige Abendmahl, mein letzter Gedanke mir vorstelle den gekreuzigten und sterbenden Jesum, mein letztes Wort aus meinem Munde rufe: Vater, in deine Hände befehle ich meinen Geist.

Ich bitte dich aber auch, o allerhöchster und hülfreicher Gott, für alle meine nothleidenden Brüder und Schwestern, sie mögen sein wer oder wo sie wollen, und befehle sie dir zu deiner Güte und Erbarmung. Ach! richte auf und befestige durch deine Gnade alle Schwachen im Glauben. Erleuchte, bekehre und bringe wieder zurecht alle Irrigen und Verführten. Erfreue, beruhige und tröste alle geängsteten Herzen und betrübten Gewissen. Sei der Kranken Arzt, Pfleger und Helfer. Ernähre, versorge und erhalte die Armen, Trostlosen und Verlassenen. Sei ein Vater, Richter und Berather der Wittwen und Waisen. Mache los und befreie die in barbarischer Dienstbarkeit oder sonsten unschuldig Gefangenen. Sei ein Retter, Schutz und Zuflucht der um deiner Ehre und Lehre
\quellseite{69}
wällen Bedrängten und Verfolgten. Erlöse und entbinde glücklich die in schweren Kindesnöthen Arbeitenden. Erhalte in wahrer Buße und Glauben, Liebe und Hoffnung die mit der letzten Todesangst Ringenden, verkürze ihnen ihre Quaal und führe durch deine Engel ihre Seelen in den seligen Schooß Abrahams zur ewigen Erquickung. Ach HErr, höre und erhöre mich und alle Menschen, so oft wir in unsern Nöthen und Aengsten zu dir schreien, und hilf uns, wo, wie und wenn wir deiner Hülfe bedürfen und begehren.

Ja, HErr, du willst solches thun, weil du unser Vater bist, und uns nicht allein zu deinen Kindern erschaffen, erlöset und geheiliget, sondern auch uns väterlich zu lieben, zu schützen und zu helfen versprochen hast. Du kannst solches thun, weil du alles vermagst im Himmel und auf Erden und nach deiner Allmacht überschwänglich zu thun weißt über alles, was wir bitten und verstehen. Du wirst es auch gewißlich thun, weil du dafür von uns sollst herzlich geliebet und herrlich gelobet und gepriesen werden, hier in der Zeit und dort in Ewigkeit, Amen. Mein HErr und mein Gott, sprich du auch dazu Amen, Amen.

\section{50. Abendsegen am Sonnabend}

Lob sei dir, du höchster und unsterblicher Gott, Lob sei deiner milden Güte und Barmherzigkeit, Lob sei deiner ewigen Weisheit und Wahrheit, der du mich diesen Tag vor allem Schaden und Uebel behütet hast. Ich bitte dich, du wollest deine Güte, so du an mir angefangen hast, gnädiglich vollenden, und mich heut diese Nacht auch lassen ruhen unter deinem höchsten Schirm, und mich mit deinen Fittigen bedecken. Laß meine Zuversicht sein unter dem Schatten deiner Arme, daß ich kein Unglück fürchte. Behüte mich, mein Gott, wie einen Augapfel, beschirme mich unter dem Schatten deiner Flügel. HErr, du bist mein Gut und mein Erbtheil, mein Heil steht in deinen Händen. Hilf mir durch deine Güte, daß nicht Furcht und Zittern über mich komme, und mich Grauen des Nachts nicht überfalle. Sei mir gnädig, denn auf dich trauet meine Seele und unter dem Schatten deiner Flügel habe ich Zuflucht. Ich suche den HErrn in meiner Noth, meine Hand ist des Nachts ausgereckt und läßt nicht ab. Denn meine Seele hat sonst
\quellseite{70}
keinen Trost, so weiß ich auch von keinem Helfer im Himmel und auf Erden, denn von dir allein. Zu Mitternacht, so ich erwache, gedenke ich an deinen Namen, daß er so lieblich ist, und an deine Güte und Treue, die du mir erzeiget hast, und danke dir für die Rechte deiner Gerechtigkeit. Wenn ich betrübt bin, so denke ich an Gott, wenn mein Herz in Aengsten ist, so rede ich von meinem Heiland. Denn er führet meine Seele aus dem Verderben und errettet mich von den Banden des Todes. Herr Gott, mein Heiland, ich schreie Tag und Nacht vor dir, daß du mir vergeben wollest alle meine Mißhandlungen, damit ich diese Woche und heut wider dich gethan habe. O Herr, errette meine Seele um deiner Barmherzigkeit willen. Du bist gnädig und gerecht und unser Gott ist barmherzig. Der Herr behütet die Einfältigen. Wenn ich niederliege, so hilft er mir auf. Darum will ich fröhlich sein und dich preisen und rühmen auf meinem Lager. Denn die Zeit meines Lebens wird mir aufgehen wie der Mittag, und das Finstere wird mir ein lichter Morgen sein und werde mich des trösten, daß du, Gott, meine Hoffnung und Ruhe bist im Leben und Tod. Ich werde mich legen und niemand wird mich aufschrecken. So befehle ich nun meinen Leib und Seele in deine Hände, du treuer Gott, du hast mich erlöset durch Jesum Christum, unsern Herrn. Amen.

\section{51. Ein Anderer}

Barmherziger, gnädiger Gott und Vater! ich sage dir abermal von Herzen Lob und Dank, daß du dich meiner die Zeit meines Lebens so väterlich angenommen, und mich diesen Tag über und sonderlich die ganze Woche bis auf diese Stunde vor allem Unglück Leibes und der Seele gnädiglich behütet hast. Ich bitte dich ferner, du wollest alle meine Sünden, so ich heute und die ganze Woche wider dich und meinen Nächsten wissentlich und unwissentlich gethan, durch deine Gnade, welche alle bußfertigen Sünder bei dir haben, zudecken, und derselben nimmermehr gedenken. Und hilf gnädiglich, daß ich aus der alten Woche in ein neu christlich Leben trete, das dir (und allen Auserwählten im Himmel) gefallen möge. Und laß mich diese angehende Nacht in deine gnädigen Arme befohlen
\quellseite{71}
sein, daß ich sicher ruhen und schlafen, frisch und gesund zu deinem ferneren Lob wieder aufstehen möge. Und wenn mein Stündlein vorhanden ist, nimm mich zu dir, HErr JEsu Christ, denn ich bin dein und du bist mein, wie herzlich gern wollt ich bald bei dir sein. Amen.

\chapter{V. Gebete für verschiedene Zeiten}

\section{A. Tägliche Gebete}

\section{52. Eine allgemeine Beichte}

Mein lieber Vater, ich bekenne allewege, du siehest es auch, daß ich je meinethalben, wie ich gehe oder stehe, inwendig und auswendig, mit Haut und Haar, mit Leib und Seel in das ewige höllische Feuer hinein gehöre; daß auch in Summa, weißt du, mein Vater, meinethalben nichts Gutes in mir ist, nicht ein Haar auf dem Haupt droben, es gehöret doch alles mit einander hinein in den Abgrund der Höllen, zu dem leidigen Teufel. Was soll ich viel Worte davon machen? Aber, mein lieber Vater, ich bitte wiederum hiergegen allewege, ich sei meinethalben was ich wolle, so bitte ich dich dennoch und will es von dir auch haben allewege, daß du dein Aufsehen und dein Aufmerken auf mich nicht wollest heben, und wollest deine Augen auf mich nicht kehren und wenden. O es ist sonst mit mir verloren und verdorben, und wenn hunderttausend Welten auf mir wären; sondern da bitte ich dich, daß du wollest dein Aufsehen und dein Aufmerken haben und wollest deine Augen kehren, wenden und richten in das Angesicht deines lieben Sohnes JEsu Christi, deines Gesalbten, meines Mittlers, Hohenpriesters und Fürsprechers, meines Heilandes, Erlösers und Seligmachers, und wollest mir um seinetwillen, und nicht um meinetwillen (bitte ich dich, mein Vater) gnädig und barmherzig sein, und wollest mir um deines lieben Sohnes JEsu Christi willen verleihen ein seliges Ende und eine fröhliche Auferstehung, hier helfen dem Leibe, und dort in
\quellseite{72}
jener Welt der armen Seele! Und um seines rosinfärben Blutes wegen, daß er da mildiglich am Galgen des Kreuzes zu Verzeihung und Vergebung meiner Sünde vergossen hat — dein Sohn JEsus Christus — bitte ich dich jetzo und, mein Vater, daß du dasselbige Blut JEsu Christi, deines lieben Sohnes, an mir armen Creatur meiner mannigfaltigen Sünden halben, die da nicht auszureden noch auszuforschen sind, nach deiner Gerechtigkeit nicht wollest anders machen und umkehren, sondern wollest es, nach deiner grundlosen Barmherzigkeit, den Nutz und die Frucht lassen an mir schaffen und ausrichten, dazu es denn in Ewigkeit von dir verordnet und von deinem lieben Sohne JEsu Christo an dem Galgen des Kreuzes auch vergossen ist, als nämlich, daß du mir es je wollest gereichen und kommen lassen zur Verzeihung und Vergebung meiner Sünden, auf daß, welche Stunde, welchen Augenblick bei Nacht oder bei Tag du kommest und klopfest an und willst wiederum meinen Geist, welchen du mir erstlich hast eingeblasen, hinwegfordern, so bitte ich dich allwege, mein Vater, daß du dir denselben meinen Geist, das ist meine Seele, wollest je lassen in deine Hände befohlen sein um deines lieben Sohnes JEsu Christi Bluts, Leidens und Sterbens willen. Amen.

\section{53. Gebet um ein seliges Ende}

Ach lieber Gott und HErr, ich lebe, aber ich weiß nicht, wie lange? Ich muß sterben und weiß nicht, wann? Du, mein himmlischer Vater, aber weißt es allein. Wohl an! soll dieser Tag (diese Nacht) der letzte Tag (die letzte Nacht) meines Lebens sein, HErr, dein Wille geschehe, der ist allein der beste, nach demselben bin ich bereit, in wahrem Glauben an Christum, meinen Erlöser, zu leben und zu sterben. Allein, lieber Gott, gewähre mir nur diese Bitte, daß ich nicht möge plötzlich in meinen Sünden sterben und verderben. Gib mir rechtschaffene Erkenntniß, Reue und Leid über meine begangene Sünde und stelle sie mir noch in diesem Leben unter Augen, damit sie mir nicht am jüngsten Tage unter Augen gestellet und ich dadurch vor Engel und Menschen zu Schanden werden möge; sondern verleihe mir so viel Zeit und Raum zur Buße, daß ich meine Uebertretung von Herzen erkennen, bekennen und berselben
\quellseite{73}
Vergebung und Trost aus deinem seligmachenden Wort erlangen möge. Ach, barmherziger Vater! verlaß mich nicht und nimm deinen Heiligen Geist nicht von mir. Mein Herz und dessen auf Jesum gerichtete Zuversicht ist dir, o Herzenserkündiger! wohl bekannt; in derselben erhalte mich zum ewigen Leben. Laß mich sterben, wann du willst; nur, so es möglich ist, verleihe mir ein vernünftiges, sanftes und seliges Ende. Amen.

\begin{verse}
Herr Jesu Christ, du höchstes Gut,\\
Ich bitte durch dein theures Blut,\\
Mach’ nur mein letztes Ende gut!
\end{verse}

Ach Gott, sei mir armen Sünder gnädig und barmherzig. Amen.

\section{B. Allgemeine Sonntagsgebete}

\section{54. Vor der Predigt}

Lieber Gott, du sprichst durch deinen lieben Sohn die selig, so dein Wort hören. Wie viel billiger wäre es, daß wir dich, o ewiger barmherziger Vater, ohne Unterlaß mit fröhlichem Herzen selig preißeten, dir dankten und dich lobten, daß du dich so freundlich, ja väterlich gegen uns arme Würmlein erzeigest, und mit uns von der größten und höchsten Sache, nämlich vom ewigen Leben und Seligkeit, redest. Gleichwohl unterlässest du es nicht, uns freundlich zu locken durch deinen Sohn, dein Wort zu hören, da er spricht: Selig sind, die Gottes Wort hören und behalten; als könntest du unseres Gehöres nicht entbehren, und wir, die wir Erd und Asche sind, nicht viel tausendmal mehr deines seligen Wortes bedürften. O wie unaussprechlich groß und wundersam ist deine Güte und Geduld! Wiederum Ach und Wehe über die Undankbarkeit und Starrblindheit derer, die dein Wort nicht allein nicht hören wollen, sondern es auch muthwillig verachten, verfolgen und lästern.

\section{55. Danksagung für das gehörte Wort Gottes}

O Gott, Vater aller Barmherzigkeit, wir danken dir herzlich und allezeit, daß du uns, nach dem überschwänglichen Reichthum
\quellseite{74}
beiner Gnaden, hast gebracht zu dem Schatz deines Wortes, darinnen wir haben Erkenntniß deines lieben Sohnes, das ist, ein sicher Pfand unsers Lebens und Seligkeit, die im Himmel zukünftig ist, und bereitet allen, die in reinem Glauben und brünstiger Liebe beständig bis ans Ende beharren; wie wir denn hoffen und bitten, du werdest uns, barmherziger Vater, erhalten und vollkommen machen mit allen Auserwählten, in einem Sinne, zu gleichem Bild deines lieben Sohnes Jesu Christi, unsers HErrn. Amen.

\section[C. Gebete für jeden Sonn- und Festtag des …]{C. Gebete für jeden Sonn- und Festtag des Kirchenjahrs}

\section{56. Gebet um fruchtbare Anhörung der Predigt}

Gütiger Gott! dieweil es dir wohlgefallen, durch die Predigt des heil. Evangelii selig zu machen die, so daran glauben, so bitte ich dich, du wollest meinen Verstand erleuchten, und mein Herz eröffnen, daß ich dein Wort nicht allein mit leiblichen Ohren höre und verstehe, sondern auch ins Herz fasse und behalte, und darnach thue. Laß mich nicht einen vergeblichen Hörer, sondern einen Thäter des Worts sein. O HErr, ich habe lieb die Stätte deines Hauses und den Ort, da deine Ehre wohnet. Laß deine Lehre aus dem Mund des Predigers auf mich triefen wie den Regen, und deine Rede fließen wie den Thau, daß ich sie mit Sanftmuth annehme und in meinem Herzen bewahre, auch gute Früchte des Glaubens bringe. Wehre dem bösen Geiste, daß er den edlen Samen nicht wiederum aus meinem Herzen reiße, noch im Kreuz und Verfolgung mich abfällig mache, oder durch die Dornen der weltlichen Sorgen, Reichthum und Wollust mein Herz tödtlich verwunde, sondern gib, daß ich durch dein lebendig Wort, als den unvergänglichen Samen, wiedergeboren und zum ewigen Leben erhalten werde. Amen.

\section{57. Ein Anderes}

Ach barmherziger Gott und Vater! du siehest, wie ich wegen meiner verderbten Natur keine Lust zu deinem Worte habe, und wie ich mich leicht vom Teufel, meinem Fleisch und Blut, von den
\quellseite{75}
Weltkindern, auch falschen Lehrern und Predigern, ja andern gar liederlichen Ursachen davon gar so leicht lasse abhalten, unfleissig, faul und träge bin, dein Wort zu hören und zu bewahren. Darum bitte ich nun, o ewiger Gott, verzeihe mir diese meine angeborne Faulheit und Trägheit, und gib mir doch ein williges und munteres Herz, dein Wort zu hören und zu betrachten. Erwecke in mir eine solche herzliche Begierde, daß ich begierig sei nach der vernünftigen lautern Milch göttlichen Worts, als die jetzt gebornen Kindlein. Hilf, daß ich meine beste Lust und Ergötzlichkeit an deinem Wort habe. Laß in dieser Welt mir nichts Lieberes als dein Wort sein, daß ich dein Wort liebe über Gold und über viel feines Gold und dasselbige allezeit für meinen besten Schatz achte. Dieweil ich auch leider die Zeit erlebt habe, von welcher dein lieber Sohn selber geweissaget, daß falsche Christi aufstehen werden und Wunder thun, daß verführet würden in den Irrthum, wo es möglich wäre, auch die Auserwählten: so bitte ich dich, du wollest mich ja für Irrthum und falscher Lehre gnädiglich behüten und bewahren. Erhalte mich bei deiner Wahrheit, denn dein Wort ist die Wahrheit, daß ich demselben als der himmlischen Wahrheit anhangen, auch darinnen bis an mein letztes Ende standhaftig verharren möge. Und dieses wollest du thun um deines allerheiligsten, hochgelobten Namens Ehre willen. Amen.

\section{58. Gebet nach der Predigt}

„Wir loben und preisen dich, HErr Gott, Vater des Himmels und der Erden, daß du uns an diesem heutigen Tage aus lauter Güte und Barmherzigkeit reichlich versorget hast, nicht allein dem Leibe nach mit dem täglichen Brot, sondern auch an unsrer Seele mit der himmlischen Speise deines heiligen Wortes. Verleihe, o treuer Gott, daß es in unsern Herzen verbleibe und großen Nutzen schaffe zu deiner Ehre und unsrer Seligkeit. Wir bitten auch deine väterliche Gnade, du wollest uns ja dein rein Wort und Kirchendienst nicht entziehen; wollest uns auch unsre Schuld und Sünden, mit welchen wir diesen heutigen Tag möchten verunheiliget haben, nicht entgelten lassen, sondern dieselbigen uns nach deiner großen Güte und durch das bittere Leiden und Sterben deines
\quellseite{76}
lieben Sohnes, Jesu Christi, vergeben. Wollest endlich heute und die ganze übrige Zeit unsers elenden Lebens in diesem Jammertal uns dir in deinen göttlichen Schutz und väterlichen Segen befohlen sein lassen und deine heiligen Engel um uns her lagern, damit der böse Feind keine Macht an uns finde. Weil aber nunmehr der Abend dieser Welt herbei nahet, und sich der Tag geneigt hat, so verleihe uns die Gnade deines Geistes, damit wir stets wachen, und im lebendigen Glauben, in brünstiger Liebe, fester Hoffnung und heiligem Wandel erwarten der seligen Erscheinung der Herrlichkeit des großen Gottes, unsers Heilandes, Jesu Christi, wenn er kommen wird, zu richten die Lebendigen und die Todten, auf daß wir ihn alsdann würdig und mit Freuden empfahen, ins ewige Leben mit ihm eingehen, und den ewigen Feier- und Freudentag mit allen Auserwählten erreichen, und ohne Ende halten mögen. Erbarme dich, o Herr! der ganzen Christenheit, und entlade deine arme verfolgte Gemeine aller Bedrängniß, Spottes und Tyrannei, tröste alle geängstete, betrübte Gewissen durch den wahren Tröster, den Heiligen Geist, und sende uns allen deinen Frieden durch Jesum Christum, unsern Fürsprecher und einigen Mittler. Amen.

\section{59. Ein Anderes}

Ich danke dir, mein himmlischer Vater, für diese große Gnade, daß du mir armen Sünder zu Trost dein heiliges Wort so reichlich hast verkündigen und fürtragen lassen. Drücke und versiegle du, o Gott, dasselbe in unsere Herzen, daß wir darnach thun, leben und sterben mögen. Gib, daß wir allezeit in deiner Erkenntniß wachsen und zunehmen. Laß es nicht allein bei uns, sondern auch bei andern Völkern bekannt und ausgebreitet werden, damit es unser und aller Menschen Herzen rühren und befeuchten möge, gleichwie der Regen und Schnee die Erde befeuchtet und fruchtbar macht, auf daß wir dich, ewigen Gott, lernen erkennen, anrufen, loben, ehren und danken, daß wir alles, was wir thun und lassen, nach deinen heiligen Geboten anstellen, damit wir allezeit wie gehorsame Kinder erfunden, endlich des ewigen Lebens und Erbtheils, das uns durch Jesum Christum erworben ist, theilhaftig werden mögen. Amen.

\quellseite{77}

\section{60. Gebet am ersten Sonntag des Advents}

Herr Gott himmlischer Vater! Wir danken dir ja billig, loben und preisen dich in Ewigkeit, daß du uns arme Menschen, so wir der Sünden halber unter des Todes und Satans Tyrannei haben sein müssen, deinen Sohn zu einem König gesetzt hast, der ein gerechter König und Heiland sei, das ist, uns von unsern Sünden durch seinen Tod retten und helfen soll. Wir bitten dich auch, du wollest durch deinen Heiligen Geist uns erleuchten, regieren und führen, daß wir an diesen rechten König und Heiland uns allein halten und an seiner armen Gestalt und verachteten Wort uns nicht, wie die Welt, ärgern, und also durch ein festes Vertrauen auf Ihn ewig selig werden. Amen.

\section{61. Gebet am zweiten Sonntag des Advents}

Herr Gott, himmlischer Vater! der du durch deinen Sohn uns hast offenbart, wie Himmel und Erde vergehen, unsere Leiber wieder auferstehen und wir alle vor Gericht sollen gestellet werden: wir bitten dich, du wollest durch deinen Heiligen Geist uns in deinem Wort und rechtem Glauben erhalten, vor Sünden gnädig behüten und in aller Anfechtung bewahren, daß wir unsere Herzen mit Fressen und Saufen und Sorgen der Nahrung nicht beschweren, sondern allezeit wacker sein und beten und in gewisser Zuversicht deiner Gnade die Zukunft deines Sohnes mit Freuden erwarten und durch ihn ewig selig werden. Amen.

\section{62. Gebet am dritten Sonntag des Advents}

O Herr Gott, himmlischer Vater! der du deinen Sohn, unsern Herrn Jesum Christum, Mensch werden und darum in diese Welt hast kommen lassen, daß er des Teufels Tyrannei wehren, uns armen Menschen wider Sünde und Tod helfen und ewig selig machen solle: wir bitten dich, du wollest mit deinem Heiligen Geist unsere Herzen also führen und leiten, daß wir auf anders nichts, denn auf dein Wort sehen und also allem Aergerniß, da unsere Vernunft uns auf weiset, entfliehen und unter dem Häuflein mögen gefunden werden, welches sich an deinem Sohn Christo Jesu nicht geärgert, sondern durch ihn selig werden. Amen.

\quellseite{78}

\section{63. Gebet am vierten Sonntag des Advents}

Herr Gott, himmlischer Vater! wir danken dir ja billig, daß du die selige gnadenreiche Taufe durch den heiligen Täufer Johannem hast erstlich angerichtet und uns auch dazu kommen lassen, in welcher du uns um deines lieben Sohnes Christi Jesu willen zusagst Vergebung der Sünden, den Heiligen Geist und das ewige Leben: wir bitten dich, erhalte uns in solchem Vertrauen auf deine Gnade und Barmherzigkeit, daß wir an solcher Zusagung nimmermehr zweifeln, sondern derselben in allerlei Anfechtung uns trösten, und gib durch deinen Heiligen Geist, daß wir vor Sünden uns hüten und in solcher Unschuld, dazu wir durch die Taufe kommen sind, bleiben mögen oder, wo wir ja fallen und aus menschlicher Schwachheit nicht können bestehen, daß wir in Sünden nicht liegen bleiben, sondern durch rechtschaffene Buße uns bekehren und deiner Zusage uns wieder trösten und also durch deine Gnade ewig selig werden. Amen.

\section{64. Vorbereitungsgebet am heiligen Christtage}

O Herr Jehova, du heiliger Gott, wie herzlich freuen wir uns, daß die heilige Zeit wieder herzukommt, daß wir in dein Haus gehen sollen, zu schauen die schönen Gottesdienste, und deinen Tempel zu besuchen, und also in öffentlicher Versammlung hören von dem kündlich-großen Geheimniß, daß Gott geoffenbaret ist im Fleisch, nämlich dein Sohn, so von dem Samen David nach dem Fleisch geboren ist. Wir bitten dich, du wollest uns Leben und Segen verleihen, daß wir diese Tage über, wie auch allezeit uns halten zu deinem Altar, da man höret die Stimme des Dankens, und da man prediget alle deine Wunder, und daß wir in großer Menge und herzlicher Andacht zu deinem Hause wallen mit Frohlocken und Danken, unter dem Haufen, die da feiern. Gib unsern Predigern und Seelsorgern Mund und Weisheit, daß sie uns predigen, was du sie heißest, und mit fröhlichem Aufthun ihres Mundes uns also unterrichten, daß ihre Lehre daher fließe, wie eine lebendige Quelle, auf daß sie sich selbst und uns selig machen. Herr, thue auf unsere Herzen, daß wir mit Freuden und herzlicher

Gebetbuch. 4
\quellseite{79}
Andacht hören und fassen die tröstlichen und holdseligen Predigten von der seligmachenden Geburt und Menschwerdung deines lieben Sohnes, unsers Heilandes Jesu Christi. Ach hilf, du frommer Vater, daß wir uns freuen im Herrn, und unsere Seele fröhlich sei in unserm Gott und Heiland, und mit Freuden singen in den Hütten der Gerechten. Schaffe in uns reine Herzen, daß wir uns mit rechtschaffener Buße zu dieser bevorstehenden heiligen Zeit gebührlich und würdig bereiten und schicken. Behüte uns diese Feiertage und allezeit für Feuersbrunst, Aufruhr, Mord und Todschlag, steure allen Werken des Fleisches, als da sind Ehebruch, Hurerei, Unreinigkeit, Unzucht, Abgötterei, Zauberei, Feindschaft, Hader, Mord, Saufen und Fressen. Wehre dem Teufel, der als ein Lügner und Mörder uns gern in guter Andacht hindern, und zu diesen Sünden treiben wollte, und hilf, daß wir diese Feiertage mit guter Gesundheit, Ruhe und Frieden vollbringen, zu unserm Trost und seliger Erbauung, damit wir alle Wohlthaten unsers lieben neugebornen Christ-Kindleins zu unsrer ewigen Seligkeit genießen mögen. Amen.

\section{65. Gebet am heiligen Christtage}

Herr Gott, himmlischer Vater! wir danken deiner großen Barmherzigkeit, daß du deinen eingebornen Sohn hast lassen Mensch werden, und dadurch uns von Sünden und dem ewigen Tod gnädiglich hast helfen lassen, und bitten dich, erleuchte unsere Herzen durch deinen Heiligen Geist, daß wir für solche deine Gnade dir dankbar seien, und derselben uns in allerlei Noth und Anfechtung trösten, und also durch deinen Sohn, unsern Herrn Christum, ewig selig werden. Amen.

\section[66. Gebet am Sonntag nach dem heiligen …]{66. Gebet am Sonntag nach dem heiligen Christtage}

Herr Gott, himmlischer Vater! der du durch deinen heiligen Simeon hast geweissaget, wie Christus, dein Sohn, zum Fall und Auferstehung Vieler in Israel gesetzt sei. Wir bitten dich, erleuchte unsere Herzen durch deinen Heiligen Geist, daß wir deinen Sohn, Christum Jesum, recht erkennen, und uns in allerlei Noth und Gefahr an ihn halten, und an ihm auferstehen, und nicht mit
\quellseite{80}
der sichern blinden Welt uns an ihm stoßen, auch in allerlei Leiden und Anfechtung, so deines Worts und Bekenntniß halber uns begegnet, getrost und muthig sein, sintemal gewiß ist, daß die, so dein Wort verfolgen, zu deinem Reich nicht gehören, die aber daran glauben und dabei bleiben, werden ewig selig sein. Amen.

\section{67. Gebet am Sylvesterabend}

Herr Gott, du bist unsere Zuflucht für und für, und ein Gott von Ewigkeit zu Ewigkeit, ehe denn die Berge worden, und die Erde und die Welt geschaffen worden. Du hast gleich in der Schöpfung gemachet Lichter an der Veste des Himmels, die da scheiden Tag und Nacht, und geben Zeichen, Zeiten, Tag und Jahre. Wir erkennen solche deine Güte sonderlich anjetzo beim Ausgang des alten und Eingang des neuen Jahres, da wir rühmen müssen nicht allein deine Weisheit, welche durch diese Lichter uns gegeben dieses Jahr und darinnen Zeichen und Zeiten, Tag und Nacht: sondern preisen auch deine große Fürsichtigkeit, daß du uns das ganze Jahr so wunderlich versorget, gespeiset, getränket, bekleidet, genähret, gewehret, und zu allem unsern guten Vornehmen deine Gnade und Segen bescheret. Wir bitten demüthiglich, laß auch weiter über uns walten deine Gnade und Wahrheit, daß mit dem alten Jahr unsre Missethaten zugestellet, und mit dem neuen Jahr die Gnadenthür uns auf ein neues aufgeschlossen werde; damit wir mit demselben auch bekommen neue Kräfte des Leibes und der Seelen, und dabei in einem neuen Leben wandeln mögen, daß es uns und unsern Kindern wohl gehe ewiglich. O getreuer Gott, gleichwie du weiland den frommen Noah aus dem Kasten geführet hast, seiner Betrübniß mit dem alten Jahr ein Ende gemacht, und mit dem neuen ihm auch eine neue Freude aufgeben lassen: also führe auch uns aus dem Kasten der Trübsal, und verleihe Gnade, daß mit dem alten Jahr das alte Ungewitter vorüber gehe, und mit dem neuen auch ein neuer fröhlicher Sonnenschein aufgehe. Und wie du vor Zeiten deinen reichen Segen in Egypten nicht nur ein Jahr, sondern sieben Jahr nacheinander gegeben: also laß deinen Segen auch immerdar ob uns schweben, damit wir denselben alle Jahr, im alten und neuen, rühmen und preisen müssen. Dessen
\quellseite{81}
mir nicht allein christlich wollen gedenken, sondern auch dafür dir allezeit danken durch Christum Jesum, der da ist das A und O, der Anfang und das Ende. Amen.

\section{68. Gebet am Neujahr}

Herr Gott, himmlischer Vater! wir danken dir für deine große unaussprechliche Gnade, daß du der Armen dich angenommen, und deinen Sohn uns zu gut unter das Gesetz hast gethan, auf daß er, durch seinen vollkommenen Gehorsam, deinen Zorn gestillet, und unsern Ungehorsam geheilet. Wir bitten dich, du wollest durch deinen Heiligen Geist unsre Herzen also erleuchten, daß wir solches Gehorsames uns wider unsre Sünde und böses Gewissen trösten, und, durch Hülfe deines Heiligen Geistes, auch anfangen, gehorsame Kinder zu sein, und endlich durch Christum Jesum ewig selig werden. Amen.

\section{69. Gebet am Sonntag nach Neujahr}

Herr Gott, himmlischer Vater! der du dich mit deinem Heiligen Geist in allen Gnaden bei der Taufe deines lieben Sohnes, der unser aller Sünde auf sich geladen, hast offenbaret, und mit deiner Stimme zu ihm gewiesen, daß wir durch ihn Vergebung der Sünden und alle Gnade haben sollen: Wir bitten dich, erhalte uns in solcher Zuversicht, und weil wir, nach dem Befehl und Exempel deines Sohnes auch getauft sind, so stärke unsern Glauben durch deinen Heiligen Geist, und hilf uns zum ewigen Leben. Amen.

\section[70. Gebet am Epiphanias- oder …]{70. Gebet am Epiphanias- oder Erscheinungsfeste Christi}

Herr Gott, himmlischer Vater! der du dein liebes Wort, den rechten Stern, der das Kindlein Jesus zeiget, uns hast lassen leuchten: wir bitten dich, gib deinen Heiligen Geist in unsre Herzen, daß wir ja solches Licht annehmen, und seliglich brauchen, daß wir, wie die Weisen, solchem Stern nachziehen, keine Mühe noch Gefahr uns lassen schrecken, sondern mit ganzem Herzen uns an deinen Sohn, Christum Jesum, als den einigen Heiland, halten,
\quellseite{82}
und unser Zeitliches dazu verwenden, daß deine arme Christenheit dadurch gebessert, und deinem Sohn, Christo Jesu, gedienet werde. Amen.

\section[71. Gebet am ersten Sonntag nach dem Feste …]{71. Gebet am ersten Sonntag nach dem Feste der Erscheinung Christi}

O Herr Gott, himmlischer Vater! ich danke dir für deine Gnade, daß du zur heiligen Taufe und zum Erkenntniß deines göttlichen Wortes und Willens mich hast kommen lassen. Ich bitte dich, gib deinen Heiligen Geist in mein Herz, daß ich mich deines Wortes befleißige, es nicht versäume, noch verachte, sondern wohl merke und daß in mir die rechte Furcht Gottes zunehme und wachse, und ich endlich seliglich in deinem Wort sterbe und selig werde. Amen.

\section[72. Gebet am zweiten Sonntag nach dem Feste …]{72. Gebet am zweiten Sonntag nach dem Feste der Erscheinung Christi}

Herr Gott, himmlischer Vater! ich danke dir, daß du mir deine Gnade gegeben, und mich in den heiligen Ehestand kommen lassen, und dadurch für Sünden behütet hast. Und bitte dich, gib deinen Heiligen Geist, der mich also darinnen leite und führe, daß ich mit meinem Ehegemahl einig sei, keine Ursache zum Zorn und Widerwillen gebe, sondern ein freundliches, gotteliges Leben führen möge, auf daß, obgleich Anfechtung nicht außen bleiben wird, wir dennoch immerdar deine gnädige Hülfe und Segen spüren, unsre Kinderlein nach deinem Willen erziehen, und in aller Gottseligkeit unter einander leben, und endlich durch deinen Sohn, Jesum Christum, ewig selig werden. Amen.

\section[73. Gebet am dritten Sonntag nach dem Feste …]{73. Gebet am dritten Sonntag nach dem Feste der Erscheinung Christi}

Herr Gott himmlischer Vater! der du aus solcher Liebe und Barmherzigkeit uns deinen Sohn gegeben, und für uns Mensch werden, und am Kreuz hast sterben lassen: Gib uns deinen Heiligen Geist in unsre Herzen, daß wir all unser Vertrauen allein auf ihn setzen, und durch ihn Vergebung unsrer Sünden, und ewiges
\quellseite{83}
Leben ja so fest glauben, als der Hauptmann hie glaubet, er werde seinem kranken Knecht mit einem Wort helfen. Wir haben ja sein Wort: Wer an ihn glaube, soll den Tod nicht sehen ewiglich. Gib nun, lieber Vater im Himmel, daß wir es mit unzweifelhaftem Herzen glauben, und an unserm letzten Ende fest dabei bleiben. Amen.

\section[74. Gebet am vierten Sonntag nach dem Feste …]{74. Gebet am vierten Sonntag nach dem Feste der Erscheinung Christi}

O HErr Gott, himmlischer Vater! der du aus väterlichem guten Rath uns, deine Kinder, hie auf Erden unter das Kreuz stellest, allerlei Ungewitter läßt über uns kommen, der Sünde damit zu wehren, und uns im Glauben, Hoffnung und Gebet zu üben. Wir bitten dich, du wollest dich unser erbarmen, unser Gebet in aller Anfechtung und Noth erhören, und gnädige Hülfe schaffen, auf daß wir deine Gnade und väterlichen Beistand erkennen, und dich mit allen Heiligen in Ewigkeit loben und preisen, der du mit dem Sohn und Heiligen Geist allein der ewige, allmächtige Gott bist. Amen.

\section[75. Gebet am fünften Sonntag nach dem Feste …]{75. Gebet am fünften Sonntag nach dem Feste der Erscheinung Christi}

O HErr Gott, himmlischer Vater! wir danken dir, daß du den guten Samen, dein heiliges Wort, in unsre Herzen hast ausgestreuet, und bitten dich, du wollest durch deinen Heiligen Geist solchen Samen in uns lassen lebendig werden, und Frucht bringen, und vor dem bösen Feinde uns behüten, daß er uns nicht Unkraut mit einsäe, behüte uns vor Sicherheit, erhalte uns in deiner Furcht, und hilf in aller Anfechtung, bis wir endlich selig werden. Amen.

\section{76. Gebet am Sonntag Septuagesima}

O HErr Gott, himmlischer Vater! der du durch dein liebes Wort uns in deinen Weingarten berufen hast: Wir bitten dich, gib deinen Heiligen Geist in unsre Herzen, daß wir in solchem Weingarten treulich arbeiten, vor Sünden und allem Aergerniß uns hüten, uns deines Worts und Willens befleißen, und alle
\quellseite{84}
Hoffnung allein auf deine Gnade, so du uns durch deinen Sohn, Christum Jesum, reichlich erzeiget hast, setzen, und also ewig selig werden. Amen.

\section{77. Gebet am Sonntag Septuagesima}

Herr Gott, himmlischer Vater! wir danken dir, daß du dein liebes Wort durch deinen Sohn, Christum Jesum, unter uns hast ausgesäet, und bitten dich, du wollest durch deinen Heiligen Geist unsere Herzen dermaßen zubereiten, daß wir solches Wort mit Fleiß hören, und in einem feinen guten Herzen behalten und in Geduld Frucht bringen, daß wir der Sünde nicht nachhängen, sie durch deine Hülfe dämpfen, und in allerlei Verfolgung deiner Gnade und ewigen Hülfe uns gewiß trösten und versehen. Amen.

\section{78. Gebet am Sonntag Estomihi}

Herr Gott, himmlischer Vater! der du dem blinden Menschen durch deinen Sohn, Christum Jesum, gnädiglich die Augen geöffnet, und das Licht hast sehen lassen. Wir armen Sünder bitten dich, du wollest durch dein liebes Wort unsere Herzen erleuchten, daß wir dich durch Christum, deinen Sohn, der für uns ist am Kreuz gestorben, und für unsre Sünden bezahlt hat, recht lernen erkennen, und in allen Nöthen und Anfechtungen allein auf deine gnädige Hülfe und Barmherzigkeit sehen, dieselbe durch ein gläubiges Gebet suchen und also Trost und Rettung wider den Teufel, Sünde und Tod finden und selig werden. Amen.

\section{79. Gebet am ersten Sonntag in der Fasten}

Herr Gott, himmlischer Vater! wir bitten dich, weil der arge Feind für und für uns nachschleichet, und wie ein brüllender Löwe umhergehet, und suchet, ob er uns könne fressen, du wollest um deines Sohnes, Christi, willen, mit dem Heiligen Geist zu uns sehen, und durch dein Wort unsre Herzen stärken, auf daß der Feind uns nicht obliege, und wir in deiner Gnade für und für bleiben und endlich selig werden. Amen.

\quellseite{85}

\section{80. Gebet am zweiten Sonntag in der Fasten}

Herr Gott, himmlischer Vater! wir bitten dich, du wollest durch deinen Heiligen Geist den Glauben und Hoffnung auf deine Gnade und Barmherzigkeit in unsern Herzen fest und gewiß machen, auf daß, ob wir gleich unsers Gewissens, Sünde und Unwürdigkeit halben Ursache haben uns zu fürchten, dennoch in deiner Gnade, wie das Cananäische Weiblein, fest hangen bleiben, und durch Christum Hülfe und Rettung in allerlei Anfechtung und Nöthen bei dir finden. Amen.

\section{81. Gebet am dritten Sonntag in der Fasten}

O Herr Gott, himmlischer Vater! der du deinen Sohn, unsern Herrn Christum Jesum, darum hast lassen Mensch werden, daß er des Teufels Tyrannei wehren, und uns Menschen wider den argen Feind schützen soll. Wir danken dir erstlich für solche gnädige Hülfe, und bitten darnach, du wollest in allen Anfechtungen gnädig zu uns stehen, vor Sicherheit uns behüten, und in deinem Wort und deiner Furcht durch deinen Heiligen Geist gnädiglich erhalten, auf daß wir bis an’s Ende für den bösen Feind befriedigt bleiben und durch deinen Sohn ewig selig werden. Amen.

\section{82. Gebet am vierten Sonntag in der Fasten}

Herr Gott, himmlischer Vater! der du durch deinen Sohn in der Wüste fünftausend Mann, Weib und Kind nicht mit gerechnet, nur mit fünf Broden und zween Fischen reichlich gespeiset hast: Wir bitten dich, sei mit deinem Segen auch gnädig bei uns, behüte uns für dem Geiz, und Sorgen der Nahrung, daß wir dein Reich und Gerechtigkeit am ersten suchen, und deine milde väterliche Güte an allem, was wir für Leib und Seele bedürfen, gnädiglich spüren, durch Jesum Christum unsern Herrn. Amen.

\section{83. Gebet am fünften Sonntag in der Fasten}

O Herr Gott himmlischer Vater! dein lieber Sohn und unser Erlöser, Christus Jesus, sagt: „Wer mein Wort hält, der wird den Tod nicht sehen ewiglich:“ Wir bitten dich, gib deinen
\quellseite{86}
Heiligen Geist in unsere Herzen, weil du zu deines Sohnes Wort uns kommen, und dasselbe hast hören lassen, daß wir es auch behalten, von ganzem Herzen uns darauf verlassen, und mitten im Tod desselben uns trösten, und den Tod in Ewigkeit nicht sehen, um deines lieben Sohnes, und unsers Erlösers willen. Amen.

\section{84. Gebet am Palmsonntag}

O Herr Gott, himmlischer Vater! der du aus väterlicher Gnade deines eingebornen Sohnes nicht verschonet, sondern in den Tod, und an das Kreuz hinan gegeben hast: Wir bitten dich, gib deinen Heiligen Geist in unsere Herzen, daß wir solcher Gnade uns herzlich trösten, vor Sünden ferner uns hüten, und was du uns zu leiden zuschicken, geduldig tragen mögen, auf daß wir durch Ihn mit dir ewiglich leben. Amen.

\section{85. Vorbereitungsgebet auf das Osterfest}

Herr Gott, dich loben wir, Herr Gott, dir danken wir, daß du den heiligen Ostertag in diesem Jahre uns hast in Frieden und guter Gesundheit erleben lassen; insonderheit preisen wir dich, daß du uns den Sieg wider alle unsere Feinde durch die Auferstehung Jesu Christi, deines lieben Sohnes, gegeben hast. Herr Jesu, wir danken dir für deine triumphirende Auferstehung vom Tod, damit du uns den Himmel eröffnet, den Tod zerstört, und uns den Trost erworben hast, daß du seist die Auferstehung und das Leben. Gott, Heiliger Geist, wir danken dir für deine Osterbotschaft, und bitten dich, du wollest die Osterfreude in unsere Herzen einpflanzen und dieselbe in uns mehren und erhalten und gleichwie durch deine Kraft die betrübten Osterherzen der Jünger sind erfreut worden, also tröste uns mit Christi Auferstehung, daß wir einen gnädigen Gott haben, daß wir ewig leben sollen, daß die Hölle und der Satan an uns keine Macht hat, auch die Welt uns nicht schaden könne. O du heilige Dreifaltigkeit, Gott Vater, Sohn und Heiliger Geist, erhalte deine Kirche wider alle geistliche und leibliche Feinde, hilf, daß wir nach diesem, in deinem Freudenleben, alle bei dir fröhliche Ostern halten und dich mit allen Engeln loben, rühmen und preisen mögen, hochgelobet in Ewigkeit. Amen.

\quellseite{87}

\section{86. Gebet am Osterfeste}

HErr Gott, himmlischer Vater! der du deinen Sohn um unsrer Sünde willen hingegeben, und um unsrer Gerechtigkeit willen auferwecket hast: Wir bitten dich, du wollest deinen Heiligen Geist uns schenken, durch ihn uns regieren und führen, im wahren Glauben erhalten, und für allen Sünden, in einem neuen reinen Leben behüten, und durch deinen Sohn, Christum JEsum, uns nach diesem Leben auferwecken zu dem ewigen Leben. Amen.

\section{87. Gebet am Ostermontag}

HErr Gott, himmlischer Vater! der du deinen Sohn, Christum JEsum, den zweien Jüngern zu Emmaus hast offenbaret: Wir bitten dich, du wollest durch dein Wort und Heiligen Geist unsre Herzen auch erleuchten, daß wir im Glauben stark und fest werden, an deinem Wort treulich halten, gerne davon reden, und fleißig damit umgehen, auf daß, ob wir gleich dem Exempel Christi nach allerlei hie auf Erden müssen leiden, dennoch durch dein Wort den rechten Trost haben und behalten, daß wir durch deinen lieben Sohn nach diesem Leben zum ewigen auferweckt werden. Amen.

\section{88. Gebet am Sonntag Quasimodogeniti}

HErr Gott, himmlischer Vater! wir danken dir für deine große unaussprechliche Gnade, daß du durch deinen Sohn, unsern HErrn Christum, das heilige Evangelium und die heiligen Sacramente verordnet hast, uns zum Trost, daß wir da Vergebung der Sünden sollen finden: Wir bitten dich, gib deinen Heiligen Geist in unsre Herzen, auf daß wir deinem Worte von Herzen glauben, und durch die heiligen Sacramente von Tag zu Tag den Glauben stärken, bis wir endlich selig werden durch JEsum Christum, deinen Sohn, unsern HErrn. Amen.

\section{89. Gebet am Sonntag Misericordias Domini}

HErr Gott, himmlischer Vater! der du uns elende Menschen väterlich bedacht, und deinen Sohn zum Hiter über uns gesetzt hast, daß er nicht allein mit deinem Wort uns weiden, sondern auch
\quellseite{88}
durch seinen Schutz wider Sünde, Tod und Teufel uns erretten und erhalten soll: Wir bitten dich, gib durch deinen Heiligen Geist, daß, gleichwie dieser Hirte uns kennet, und unsrer Noth sich annimmt, wir wiederum in allerlei Anliegen ihn auch erkennen, uns an ihn halten, Hülfe und Trost bei ihm suchen und gewarten, seiner Stimme von Herzen folgen, und also durch ihn ewig selig werden. Amen.

\section{90. Gebet am Sonntag Jubilate}

Herr Gott, himmlischer Vater! der du aus väterlichem Wohlmeinen uns, deine Kinder, hie auf Erden unter der Ruthen willst haben, auf daß wir deinem eingebornen Sohn, beides im Leiden und hernach in der Herrlichkeit, sollen gleich sein: Wir bitten dich, tröste uns in der Anfechtung und im Kreuz mit deinem Heiligen Geist, daß wir nicht verzagen, sondern nach der Zusagung deines Sohns fest an diesem Trost halten, es soll das Leiden nur ein kleines sein, und darnach die ewige Freude folgen, auf daß wir also durch diese Hoffnung alles Unglück in Geduld überwinden, und durch Jesum Christum selig werden. Amen.

\section{91. Gebet am Sonntag Cantate}

O Herr Gott, himmlischer Vater! der du durch deinen Sohn uns den Heiligen Geist verheißen hast, daß er die Welt um die Sünde, Gerechtigkeit und das Gericht strafen soll: Wir bitten dich, erleuchte unsre Herzen, daß wir unsre Sünde erkennen, und durch den Glauben an Christum, zur ewigen Gerechtigkeit kommen, und in allerlei Noth und Anfechtung diesen rechten Trost fassen, daß Christus der Herr über Teufel, Tod und Alles sei, und uns mit Gnaden von allem Unglück helfen, und ewig wolle selig machen. Amen.

\section{92. Gebet am Sonntag Rogate}

O Herr Gott, himmlischer Vater! der du durch deinen Sohn uns verheißen hast, was wir in seinem Namen bitten, das wollest du uns geben: Wir bitten dich, erhalte uns in deinem Worte, gib deinen Heiligen Geist, der uns regiere und führe nach deinem Willen
\quellseite{89}
, behüte uns vor des Teufels Reich, vor falscher Lehre, unrechtem Gottesdienst, behüte auch unsern Leib und Leben vor allem Unglück, gib deinen Segen und Frieden, auf daß wir allenthalben deine gnädige Hülfe erfahren, und dich hie und dort in Ewigkeit, als unsern gnädigen Vater durch Christum preisen und loben. Amen.

\section{93. Gebet am Himmelfahrtstage}

Jesu Christe, du Sohn des allmächtigen Gottes! der du fortan nicht mehr auf Erden arm und elend, sondern zur Rechten deines Vaters, ein Herr über alles bist: Wir bitten dich, du wollest deinen Heiligen Geist uns senden, fromme Kirchendiener uns geben, ob deinem Wort halten; dem Satan und allen Tyrannen wehren, und dein Reich auf Erden gewaltig erhalten, bis alle deine Feinde zu deinen Füßen liegen, und wir durch dich über Sünde, Tod und alles siegen. Amen.

\section{94. Vorbereitungsgebet auf Pfingsten}

Herr Gott Heiliger Geist, der du am heiligen Pfingsttag die Herzen der lieben Apostel mit neuen Gaben erfüllet hast, erleuchte auch unsere Herzen und gib, daß wir andächtige Pfingsten halten. Komm du edler Gast unsers Herzens, komm du einiger Schatz unserer Seelen, wohne in uns als in deinem Tempel, reinige unsere Herzen von Sünden, laß uns alle sein die Gefallten des Herrn und hilf, daß wir durch dich wider alle unsere geistliche Feinde streiten und den Sieg erlangen. Vertritt uns alle mit unaussprechlichem Seufzen und Bewegniß zum Gebet und Danksagung. Tröste alle betrübten Pfingstherzen und laß die Gebeine fröhlich werden, die du zerschlagen hast. Erhalte uns durch deine Kraft wider die Anklage unseres Herzens und Gewissens, daß an uns nichts Verdammliches sei. Dein Friede, welcher höher ist denn alle menschliche Vernunft, bewahre unsere Herzen und Sinne zur ewigen Freude. Gib und verleihe uns, o Gott Heiliger Geist, herzliche Pfingstfreude, und laß uns nach diesem Leben deine Wunder und Thaten ewig verkündigen, hochgelobet sammt dem Vater und Sohn in Ewigkeit. Amen.

\quellseite{90}

\section{95. Gebet am ersten Pfingsttag}

Herr Jesu Christe, du Sohn des allmächtigen Gottes! wir bitten dich, du wollest durch dein liebes Wort deinen Heiligen Geist in unsre Herzen geben, daß derselbe uns regiere und führe nach deinem Willen, und in allerlei Anfechtung und Unglück uns tröste, und in deiner Wahrheit, wider allen Irrthum uns leite, auf daß wir im Glauben fest bestehen, in der Liebe und allen guten Werken zunehmen, und durch eine gewisse Hoffnung deiner Gnade, so du uns mit deinem Sterben erworben hast, ewig selig werden, der du regierest mit dem Vater und Heiligen Geist. Amen.

\section{96. Gebet am zweiten Pfingsttag}

Herr Gott, himmlischer Vater! der du aus väterlicher Liebe gegen uns arme Sünder deinen Sohn uns geschenket hast, daß wir an ihn glauben, und durch den Glauben sollen selig werden: Wir bitten dich, gib deinen Heiligen Geist in unsere Herzen, daß wir in solchem Glauben bis an’s Ende beharren, und selig werden durch Christum unsern HErrn. Amen.

\section[97. Gebet am Sonntag Trinitatis oder Fest der …]{97. Gebet am Sonntag Trinitatis oder Fest der heil. Dreifaltigkeit}

O allmächtiger, ewiger Gott, Gott Vater, Gott Sohn, Gott Heiliger Geist! der du deine Gnade uns Menschen in dem erzeiget, und deine Barmherzigkeit reichlich hast sehen lassen, daß erstlich du, ewiger Vater, uns zu Menschen geschaffen, und du, ewiger Sohn Gottes, für uns Mensch worden, und für unsre Sünde am Kreuz bezahlet, und du, ewiger Geist, der du vom Vater und Sohn in Ewigkeit ausgehet, durch das heilige Evangelium uns gläubig gemacht und geheiliget hast: Wir bitten dich, o einiger, ewiger, allmächtiger Gott, Gott Vater! Gott Sohn! Gott Heiliger Geist! erhalte uns in deiner Gnade, bis an das Ende, und mache uns ewig selig. Amen.

\section{98. Ein Anderes}

Herr Gott, himmlischer Vater! wir armen Sünder bekennen, daß nichts Gutes an uns ist, und wir unserthalben müssen in Sünden
\quellseite{91}
und ewigen Tode verderben, sintemal was vom Fleisch geboren wird, ist Fleisch, und kann dein Reich nicht sehen. Wir bitten dich aber, sei gnädig und barmherzig, und sende deinen Heiligen Geist um deines Sohnes, Christi Jesu, willen in unsere Herzen, und mache neue Menschen aus uns, daß wir Vergebung der Sünden durch Christum, wie uns in der Taufe zugesagt, gewiß glauben, und in der Liebe gegen unsern Nächsten, und allen andern christlichen Tugenden von Tag zu Tag zunehmen, bis wir endlich selig werden. Amen.

\section{99. Gebet am ersten Sonntag nach Trinitatis}

HErr Gott, himmlischer Vater! wir bitten dich, du wollest unsre Herzen durch deinen Heiligen Geist also führen und regieren, daß wir dein liebes Wort nicht wie der reiche Mann ohne Frucht hören, und mit dem zeitlichen Gut nicht also umgehen, daß wir des Ewigen dabei vergessen, sondern armen Leuten gern und mildiglich nach unserm Vermögen helfen, und mit Hoffart und Ueberfluß uns nicht versündigen, auch wo wir mit dem Kreuz und Unglück bedrängt werden, nicht verzagen, sondern alle Hoffnung auf deine ewige Hülfe und Gnade haben, und in Geduld alles überwinden. Amen.

\section{100. Gebet am zweiten Sonntag nach Trinitatis}

HErr Gott, himmlischer Vater! wir danken dir, daß du uns zu deiner Mahlzeit, dem ewigen Leben, durch dein liebes Wort hast lassen laden, und bitten dich, du wollest durch deinen Heiligen Geist unsere Herzen erwecken, daß wir dein Wort nicht vergeblich hören und das Zeitliche uns nicht daran hindern lassen. Amen.

\section{101. Gebet am dritten Sonntag nach Trinitatis}

HErr Gott, himmlischer Vater! wir sind, leider! wie Schafe in der Irre und haben uns durch den Satan und unser sündhaftiges Fleisch vom rechten Wege lassen abführen. Wir bitten dich aber, sei gnädig und vergib all unsere Sünde, um deines Sohnes Christi Jesu willen, und erwecke unsere Herzen durch deinen Heiligen Geist, daß wir an deinem Wort fest halten und im
\quellseite{92}
Schaffsal deiner christlichen Kirche in rechter Buße und wahrem Glauben bis an unser Ende verharren und ewig selig werden. Amen.

\section{102. Gebet am vierten Sonntag nach Trinitatis}

HErr Gott, himmlischer Vater! der du barmherzig bist und uns durch deinen Sohn JEsum Christum alle unsere Sünden zu vergeben versprochen hast, wir bitten dich herzlich, du wollest uns nicht richten nach unserm bösen Leben, unrechten Wandel und vielen aus menschlicher Schwachheit begangenen Sünden, sondern dieselben alle, als unser himmlischer Vater, dem unsere Schwachheit wohl bekannt, nach deinem göttlichen Vaterherzen verzeihen und vergeben, die Krone der ewigen Seligkeit uns zulegen, auch unsere harten, steinernen Herzen lindern und mindern, damit wir unsern Nächsten gern mögen vergeben, ihnen mit Barmherzigkeit erscheinen, niemand mit frühzeitigem, unbegründetem, schnellem Urtheil unschuldig, wie oft geschieht, richten und verdammen, sondern jedermänniglich von ganzem Herzen gern verzeihen und vergeben, über uns, unsere Gewissen, Herzen und Seelen selbst erstlich rechte, einige, wahre, gute Richter sein, auf daß wir also in deiner Furcht christlich leben und seliglich sterben mögen. Amen.

\section{103. Gebet am fünften Sonntag nach Trinitatis}

O JEsu Christe, du Sohn des lebendigen Gottes! der du uns dein liebes Wort gegeben und mit allerlei leiblichem Segen begnadet hast: wir erkennen, daß wir zu solchem Allen unwürdig und wohl Aergers verdienet hätten, und bitten dich, du wollest uns, wie Petro, unsere Sünden vergeben und Glück und Heil zu unserm Beruf verleihen, auf daß wir, durch dich zeitlich und ewig erhalten, beschützet, dich in Ewigkeit loben und preisen. Amen.

\section{104. Gebet am sechsten Sonntag nach Trinitatis}

O HErr Gott, himmlischer Vater! wir bekennen, daß wir leider! arme, elende Sünder sind, an denen gar nichts Gutes ist; denn unser Herz, Fleisch und Blut ist dermaßen vergiftet, daß wir ohne die bösen Lüste nimmer sein können. Darum, lieber Vater,
\quellseite{93}
vergib uns solche Sünde und reinige unsere Herzen durch deinen Heiligen Geist, daß wir zu deinem Wort Lust und Liebe haben und uns darnach halten und durch Christum in deiner Gnade ewig bleiben mögen. Amen.

\section{105. Gebet am siebenten Sonntag nach Trinitatis}

HErr Gott, himmlischer Vater! der du durch deinen Sohn in der Wüste vier tausend Mann, Weib und Kind nicht mitgerechnet, nur mit sieben Broden und ein wenig Fischlein reichlich gespeiset hast: wir bitten dich, sei mit deinem Segen auch gnädig bei uns, behüte uns vor dem Geiz und Sorgen der Nahrung, daß wir dein Reich und deine Gerechtigkeit am ersten suchen und deine milde, väterliche Güte in allem, was wir zu Leib und Seele bedürfen, gnädiglich spüren durch Christum, unsern HErrn. Amen.

\section{106. Gebet am achten Sonntag nach Trinitatis}

HErr Gott, himmlischer Vater! wir danken dir von ganzem Herzen, daß du uns hast lassen zur Erkenntniß deines lieben Wortes kommen. Wir bitten dich aber auch, du wollest uns dabei gnädig erhalten und seliglich darin sterben lassen, wollest auch alle frommen Prediger, die dein Wort treulich vortragen, vor allem Aergerniß und anderm Unglück gnädig behüten und ihnen ihr Leben fristen, die andern aber, so untreulich mit deinem Wort umgehen und sich als fromme Schäflein sehen lassen, sind aber reißende Wölfe, die wollest du ihrem Verdienst nach strafen und ihnen wehren und deine arme Christenheit gnädig vor ihnen behüten durch Jesum Christum, deinen Sohn, unsern HErrn. Amen.

\section{107. Gebet am neunten Sonntag nach Trinitatis}

O HErr Gott, himmlischer Vater! der du uns reichen Segen und das tägliche Brod überflüssig hast gegeben: wir bitten dich, behüte uns vor Geiz und erwecke unsre Herzen, daß wir armen Leuten gern geben und mittheilen, auf daß wir als treue Haushalter über deine Güter erfunden werden und nicht müssen in Ewigkeit darben, wann wir von solcher Haushaltung abgefordert und vor dein Gerichte kommen. Amen.

\quellseite{94}

\section{108. Gebet am zehnten Sonntag nach Trinitatis}

Allmächtiger, ewiger Gott! der du dein Wort von deinem Sohn, Christo Jesu, durch deinen Heiligen Geist uns geoffenbaret und kund gethan hast: wir bitten dich, erwecke unsere Herzen, daß wir es mit Ernst annehmen und nicht in den Wind schlagen oder unfleissig hören, wie dein Volk, die ungläubigen Juden, auf daß wir in deiner Furcht und rechtem Vertrauen auf deine Barmherzigkeit leben und täglich zunehmen und endlich durch deinen Sohn, Christum Jesum, ewig selig werden. Amen.

\section{109. Gebet am eilften Sonntag nach Trinitatis}

O Herr Gott, himmlischer Vater! ich bitte dich, du wollest durch deinen Heiligen Geist mich also leiten und führen, daß ich erstlich meiner Sünde nicht vergesse, noch sicher werde, sondern in steter Buße stehe und darnach mich allein desselben tröste, daß du um deines Sohnes Christi Jesu willen mir gnädig sein, alle Sünden vergeben und mich selig wollest machen. Amen.

\section{110. Gebet am zwölften Sonntag nach Trinitatis}

Allmächtiger, ewiger Gott, du gnädiger Schöpfer aller Dinge! ich danke dir, daß du mir gesunden Leib gegeben und meine Zunge, Ohren und andere Gliedmaßen gnädiglich vor dem bösen Feind behütet hast; und bitte dich, gib mir deine Gnade, daß ich meine Ohren und Zunge recht brauchen, mit den Ohren dein Wort fleißig hören und wohl merken und mit meiner Zunge deine Gnade preisen und rühmen möge, auf daß niemand durch meine Zunge geärgert, sondern jedermann dadurch gebessert werde. Amen.

\section[111. Gebet am dreizehnten Sonntag nach …]{111. Gebet am dreizehnten Sonntag nach Trinitatis}

Herr Gott, himmlischer Vater! wir danken dir von Grund unseres Herzens, daß du uns diese selige Zeit hast erleben lassen, daß wir die selige Lehre, dein Evangelium, hören, deinen Willen und väterliches Herz dadurch erkennen und deinen Sohn Jesum Christum sehen können. Wir bitten dich durch deine grundlose

Gebetsschatz.

\quellseite{95}

\section{06. Gebete für verschiedene Zeiten}

Barmherzigkeit, du wollest solches seliges Licht deines Wortes uns gnädig erhalten und durch deinen Heiligen Geist unsre Herzen also leiten und führen, daß wir nimmermehr davon abweichen, sondern fest daran halten, eher alles verlassen und endlich dadurch selig werden. Amen.

\section[112. Gebet am vierzehnten Sonntag nach …]{112. Gebet am vierzehnten Sonntag nach Trinitatis}

HErr Gott, himmlischer Vater! der du hast durch deinen Sohn Christum Jesum uns Gläubigen allen von dem greulichen Ausfluß der Sünde durch dein Wort und heilige Taufe gnädiglich geholfen, und beweisest noch täglich gnädige Hülfe in allem, was wir bedürfen: wir bitten dich, erwecke unsere Herzen durch deinen Heiligen Geist, daß wir solcher Wohlthat nimmer vergessen, sondern in deiner Furcht und herzlichem Vertrauen auf deine Barmherzigkeit, für und für leben, und mit fröhlichem Herzen dir immerdar danken und dich loben. Amen.

\section[113. Gebet am fünfzehnten Sonntag nach …]{113. Gebet am fünfzehnten Sonntag nach Trinitatis}

HErr Gott, himmlischer Vater! wir danken dir für alle deine Wohlthat, daß du uns Leib und Leben gegeben, und bisher gnädiglich erhalten hast, und bitten dich, nimm deinen Segen nicht von uns, sondern behüte uns vor dem Geiz, auf daß wir dir allein dienen, dich lieben, und dir anhangen, und uns nicht mit Abgötterei und dem schädlichen Mammonsdienst versündigen, sondern alle Hoffnung, Trost und Zuversicht allein auf deine Güte und Gnade setzen, durch Jesum Christum, unsern HErrn. Amen.

\section[114. Gebet am sechzehnten Sonntag nach …]{114. Gebet am sechzehnten Sonntag nach Trinitatis}

HErr Gott, himmlischer Vater! der du deinen Sohn hast lassen Mensch werden, daß er erstlich mit seinem Tode für unsre Sünde bezahlen, und darnach uns vom ewigen Tode helfen soll: Wir bitten dich, erhalte uns in solcher Hoffnung, auf daß wir gar nicht daran zweifeln, wie unser lieber HErr Christus der Wittwen Sohn durch sein Wort auferweckt, daß er uns desgleichen am jüngsten Tage auch auferwecken, und ewig werde selig machen. Amen.

\quellseite{96}

\section[115. Gebet am siebenzehnten Sonntag nach …]{115. Gebet am siebenzehnten Sonntag nach Trinitatis}

Herr Gott, himmlischer Vater! wir bitten dich, du wollest uns durch deinen Heiligen Geist also regieren und führen, daß wir in deiner Furcht uns halten, und nicht hoffärtig werden, sondern mit ganzem Herzen dein Wort hören und annehmen, und den Sabbath recht heiligen, auf daß durch dein Wort wir auch geheiliget werden, erstlich also, daß wir auf deinen Sohn, Jesum Christum, all unser Vertrauen und Hoffnung setzen, daß er allein sei unsere Gerechtigkeit und Erlösung, und darnach unser Leben nach deinem Wort auch bessern, und vor allem Aergerniß uns hüten, bis wir durch deine Gnade in Christo ewig selig werden. Amen.

\section[116. Gebet am achtzehnten Sonntag nach …]{116. Gebet am achtzehnten Sonntag nach Trinitatis}

O Herr Gott, himmlischer Vater! wir sind ja arme elende Sünder, deinen Willen wissen wir, aber wir sind gar zu schwach, wir können ihn nicht erfüllen, Fleisch und Blut ist uns zu stark, so läßt der leidige Feind, der Teufel, uns auch nicht zufrieden. Darum bitten wir dich, du wollest deinen Heiligen Geist in unsere Herzen ausgießen, daß wir mit festem Glauben an deinen Sohn, Jesum Christum, uns halten, seines Todes und Sterbens uns trösten, und volle Vergebung aller Sünden durch ihn glauben, und also hie auf Erden in deinem Willen und Gehorsam heilig leben, und auf deine Gnade durch Jesum Christum seliglich sterben mögen. Amen.

\section[117. Gebet am neunzehnten Sonntag nach …]{117. Gebet am neunzehnten Sonntag nach Trinitatis}

Allmächtiger, ewiger Gott! der du durch deinen Sohn, Jesum Christum, den gichtbrüchigen Menschen an Leib und Seele hast gnädiglich helfen lassen: Ich bitte deine grundlose Barmherzigkeit, sei mir auch gnädig, und vergib mir alle meine Sünde; stärke meinen Glauben durch dein Wort und Heiligen Geist, und leite mich also, daß ich nicht durch Sünde zu Krankheit und andern Unglück Ursach gebe, sondern in deiner Furcht mich halte, und also vom zeitlichen und ewigen Zorn und Strafe erlediget werde. Amen.

\section[118. Gebet am zwanzigsten Sonntag nach …]{118. Gebet am zwanzigsten Sonntag nach Trinitatis}

Herr, Gott, himmlischer Vater! ich danke dir für deine große Barmherzigkeit und Gütigkeit, daß du mich, in deinem lieben
\quellseite{97}
Wort, auf die fröhliche Hochzeit kommen lassen, und durch deinen Sohn mir alle meine Sünden vergeben hast, weil aber der täglichen Anfechtungen, Aergernisse und Gefahr so viel ist, und wir für uns selbst sehr schwach, und zur Sünde sehr geneigt sind, bitte ich dich, du wollest durch deinen Heiligen Geist, mich für allem Unfall gnädig behüten, und aber, wo ich falle, und also mein Kleid, welches mir dein lieber Sohn, JEsus Christus, angezogen hat, beflecke, gnädig wieder aufhelffen, und mich zur rechten Buße und wahrem Glauben auf deine Barmherzigkeit leiten, und in Ewigkeit nicht sinken lassen. Amen.

\section[119. Gebet am einundzwanzigsten Sonntag nach …]{119. Gebet am einundzwanzigsten Sonntag nach Trinitatis}

Allmächtiger, ewiger Gott! der du durch deinen Sohn Vergebung unserer Sünden, Gerechtigkeit und ewiges Leben uns hast verheißen, wir bitten dich, du wollest durch deinen Heiligen Geist unsere Herzen also führen und erwecken, daß wir solche Hülfe durchs tägliche Gebet, und sonderlich in aller Anfechtung bei ihm suchen, und durch einen festen Glauben auf seine Zusagung und Wort gewiß finden und erlangen, und endlich durch ihn ewig selig werden. Amen.

\section[120. Gebet am zweiundzwanzigsten Sonntag nach …]{120. Gebet am zweiundzwanzigsten Sonntag nach Trinitatis}

Allmächtiger, ewiger Gott! wir erkennen, daß wir Sünder sind, und tief in deinem Register stehen, wir danken dir aber von Herzen, daß du solche Schuld von uns genommen, und auf deinen lieben Sohn, Christum JEsum, geleget, und ihn für uns hast zahlen lassen, und bitten dich auch, du wollest in solchem Glauben uns gnädig erhalten, und durch den Heiligen Geist hie auf Erden also leiten und führen, daß wir deinem Willen nachleben, und unsern Nächsten alle Liebe, Dienst und Hülfe gern beweisen. Wiederum allen Zorn, Ungeduld und Nachgierigkeit fallen lassen, auf daß wir deinen Zorn wider uns nicht erregen, sondern immerdar einen gnädigen Vater durch Christum JEsum an dir haben, der mit dir und dem Heiligen Geist lebet und regieret von Ewigkeit zu Ewigkeit. Amen.

\quellseite{98}

\section[121. Gebet am dreiundzwanzigsten Sonntag nach …]{121. Gebet am dreiundzwanzigsten Sonntag nach Trinitatis}

O Herr Gott, himmlischer Vater! wir danken dir, daß du bisher gnädigen Frieden geben, uns vor Krieg und fremder Herrschaft väterlich behütet hast. Wir bitten dich, gib deine Gnade ferner, daß wir in deiner Furcht und nach deinem Willen leben, und durch Sünde nicht Ursach zu Krieg und andern dergleichen Strafen geben. Regiere und führe unsre Obrigkeit also, daß sie deinen Gehorsam nicht hindere, sondern fördere, und wir desto mehr Segens und Glücks unter ihrem Regiment haben mögen. Amen.

\section[122. Gebet am vierundzwanzigsten Sonntag nach …]{122. Gebet am vierundzwanzigsten Sonntag nach Trinitatis}

O allmächtiger Gott! der du durch deinen Sohn Vergebung der Sünden und Rettung wider den ewigen Tod uns hast zugesagt, wir bitten dich, stärke uns durch deinen Heiligen Geist, daß wir in solchem Vertrauen auf deine Gnade durch Christum täglich zunehmen, und die Hoffnung fest halten, daß wir nicht sterben, sondern sanft einschlafen, und durch deinen Sohn, Christum Jesum, am jüngsten Tage zum ewigen Leben wieder sollen auferwecket werden. Amen.

\section[123. Gebet am fünfundzwanzigsten Sonntag nach …]{123. Gebet am fünfundzwanzigsten Sonntag nach Trinitatis}

Herr Gott, himmlischer Vater! wir danken dir von Herzen, daß du uns aus dem schädlichen Irrthum des Papstthums durch dein Wort geführet, und zum rechten Licht deiner Gnaden gebracht hast, und bitten dich, du wollest uns in solchem Licht gnädig erhalten, vor allem Irrthum und Ketzerei bewahren, und ja davor behüten, daß wir nicht, wie die Juden, dir undankbar werden, dein Wort verachten und verfolgen, sondern daß wir es mit ganzem Herzen annehmen, unser Leben darnach bessern, und alles Vertrauen allein auf die Gnade und das Verdienst deines lieben Sohnes, Jesu Christi, setzen, der mit dir und dem Heiligen Geist lebet und regieret von Ewigkeit zu Ewigkeit. Amen.

\section[124. Gebet am sechsundzwanzigsten Sonntag …]{124. Gebet am sechsundzwanzigsten Sonntag nach Trinitatis}

O Herr Gott, himmlischer Vater! es befiehlt uns ja dein lieber Sohn, unser Herr Jesus Christus, wir sollen den Hungrigen
\quellseite{99}
speisen, den Durstigen tränken, den Nackenden kleiden, den Kranken besuchen, und also Liebe und Treue an unserm Nächsten beweisen. Du weißt aber ja auch, lieber Vater! wie kalt, träg und faul, leider! unsere Herzen von Natur zu solchem Allen sind, und wie heftig auch der leidige Teufel die Liebe in uns hindert: bitten dich demnach von Herzen, du wollest uns verzeihen und vergeben alles, damit wir bisher diesem guten und heilsamen Befehle zuwider gehandelt haben, und uns durch deinen Heiligen Geist regieren, auf daß wir in brüderlicher Liebe täglich zunehmen, und darinnen bis ans Ende verharren, auch endlich mit allen Heiligen eingehen in das Reich, das ist die ewige Freude und Seligkeit. Amen.

\section{125. Gebet am Geburtstage}

Mein HErr und mein Gott, in dir lebe, webe und bin ich, todt und lebendig. Ich danke dir, daß ich wunderbarlich gemacht bin, wunderbarlich sind deine Werke, und das erkennet meine Seele wohl. Du warest über mir im Mutterleibe; es war dir mein Gebein nicht verborgen, da ich im Verborgenen gemacht ward, da ich gebildet ward unten in der Erde; deine Augen sahen mich, da ich noch unbereitet war, und waren alle Tage auf dein Buch geschrieben, die noch werden sollten, und derselben noch keiner da war. Ich erinnere mich heute des Tages meiner Geburt, welcher ein Anfang ist nicht allein des zeitlichen, sondern auch des geistlichen Lebens, da ich bald zur heiligen Taufe getragen, mit Gott versöhnet, durch’s Wasserbad im Wort gereiniget, von Sünden abgewaschen worden und den Heiligen Geist empfangen habe; dafür sage ich dir herzlich Lob und Dank. Ich preise dich, Gott, mein Schöpfer, daß du mich heute als einen Menschen, wiewohl mit Schmerzen und Weinen, aus meiner Mutter Leibe gezogen hast; daß du mich auch an einen Ort gebracht hast, da dein Name rechtschaffen erkannt und dein heiliges Wort rein gelehrt wird, dadurch ich mit deiner Hülfe zum wahren Glauben kommen bin. Ich danke dir auch, o HErr JEsu Christe, für alle Wohlthaten, damit du mich an Leib und Seele häufig überschüttet, von Mutterleibe an bis auf diese gegenwärtige Stunde, in dem, daß du nicht allein meinen Leib
\quellseite{100}
ernähret und erhalten hast, und das zeitliche Leben auf mancherley Weise durch die heiligen Engel, welche allezeit das Angesicht Gottes im Himmel sehen, beschirmet, mich auch täglich mit Speise, Trank und Kleidern reichlich versorget hast, sondern daß du auch hierüber geistliche und unvergängliche Gaben durch dein Wort mir mitgetheilet, und in der heiligen Taufe, als dem Bad der Wiedergeburt bezeuget, daß ich eine neue Creatur und Erstling Jesu Christi worden, mir auch Vergebung der Sünden, den Heiligen Geist und Erbschaft des ewigen Lebens geschenket hast und da ich oft aus Schwachheit des Fleisches, oder muthwillig von Gott gewichen, seiner Gebote vergessen und mich mit viel schrecklichen Sünden mehr denn siebenzig mal sieben mal im Himmel wider Gott, und auf Erden wider den Nächsten vergriffen, daß du nicht mit mir bist in das Gericht gegangen, und mich mit der plötzlichen wohlverdienten Strafe überfallen. O Gott Heiliger Geist, wie schändlich habe ich mein Leben zugebracht, wie in manchem Jahr, Ort und Stelle habe ich schrecklich gestrauchelt und große himmelschreiende Sünde begangen! Aber du, o Gott Heiliger Geist, hast mich verlorenes Schaf so vielmal gesucht, wieder auf den rechten Weg gebracht, geheiliget, getröstet, regieret, und in meinem Beruf erhalten, dem Teufel gesteuret und unzählige leibliche und geistliche Gefahr abgewendet, darin ich längst hätte verderben müssen, wenn du mir nicht beigestanden hättest. O heilige Dreifaltigkeit, geboren werden hat seine Zeit, hilf aller Leibesfrucht und gib ihr einen fröhlichen Geburtstag, Vergebung der Sünden, und das ewige Leben. Laß mich heute meinen Geburtstag, wie auch alle künftigen Tage meines Lebens ohne verdammliche Sünde begehen und mich meiner geringen Ankunft und Elendes erinnern. Erhalte mich bis an das gesteckte Ziel bei Gesundheit und Leben. Verleihe mir Geduld in allem Kreuz und Leiden. Tröste mich in allem Unfall und Nöthen. Vergib mir täglich meine Sünden, und verleihe mir Beständigkeit zu meiner Seelen Seligkeit, daß ich im rechten seligmachenden Glauben und gutem Gewissen, dir bis an’s Ende diene durch Jesum Christum. Amen.

\quellseite{101}

\chapter{VI. Berufsgebete}

\section{A. Nähstand}

\section{126. Gebet der Eltern für ihre Kinder}

Ach getreuer, lieber Gott und Vater, Schöpfer und Erhalter aller Creaturen: Ich danke dir von Herzen für die Leibesfrüchte, so du mir durch deinen Segen gegeben hast, und bitte dich herzlich, weil du gesagt hast: du wollest deinem Heiligen Geist geben allen, die dich darum bitten: Begnadige auch meine armen Kinder mit deinem Heiligen Geist, der in ihnen die wahre Furcht Gottes anzünde, welche ist der Weisheit Anfang, und die rechte Klugheit, wer darnach thut, deß Lob bleibet ewiglich. Befelige sie auch mit deiner wahren Erkenntniß, behüte sie für aller Abgötterei und falscher Lehre, laß sie in dem wahren seligmachenden Glauben und in aller Gottseiligkeit aufwachsen, und darin bis ans Ende verharren. Gieb ihnen ein gläubiges, gehorsames, demüthiges Herz, auch die rechte Weisheit und Verstand, daß sie wachsen und zunehmen an Alter und Gnade bei Gott und den Menschen. Ich pflanze in ihr Herz die Liebe deines göttlichen Wortes, daß sie seien andächtig im Gebet und Gottesdienst, ehrerbietig gegen die Diener des Worts, und gegen jedermann aufrichtig in Handlung, schamhaftig in Geberden, züchtig in Sitten, wahrhaftig in Worten, treu in Werken, fleißig in Geschäften, glücklich in Verrichtungen ihres Berufs und Amts, verständig in Sachen, richtig in allen Dingen, sanftmüthig und freundlich gegen alle Menschen. Behüte sie für allen Aergernissen der argen Welt, daß sie nicht verführet werden durch böse Gesellschaft, daß sie nicht in Schlemmerei und Unzucht gerathen, daß sie ihnen ihr Leben nicht selber verkürzen, auch andere nicht beleidigen, sei ihr Schutz in aller Gefahr, daß sie nicht plötzlich umkommen, laß mich ja nicht Unehre und Schande, sondern Freude und Ehre an ihnen erleben, daß durch sie auch dein Reich vermehret, und die Zahl der Gläubigen groß werde, daß sie auch im Himmel um deinen Tisch
\quellseite{102}
hersitzen mögen als die himmlischen Oelzweige, und dich mit allen Auserwählten ehren, loben und preisen mögen, durch JEsum Christum, unsern HErrn. Amen.

(Hierauf kann der 127. und 128. Psalm gebetet werden.)

\section[127. Gebet, wenn Eltern ihre Kinder zur …]{127. Gebet, wenn Eltern ihre Kinder zur Schule schicken}

Getreuer Gott und himmlischer Vater, der du uns Eltern mit Ernst in deinem Wort befohlen, unsere Kinder dir wohlgefällig zu erziehen, ich bitte dich von Herzensgrund, laß mich und männiglich gebührlich erkennen, was an Auferziehung der jungen Kinder für ein Großes gelegen, und verleihe, daß ich den Meinen in solchem gebührlich fürstehe, und weil ich zu desselben Beförderung dies mein Kind, in deinem Namen, du Allerhöchster, zur Schulen zu schicken, entschlossen, so bitte ich, erfülle es mit deines Heiligen Geistes Gaben, regiere es zu deinem Erkenntniß und Furchten, zu fleißigem Lernen und Gehorsam, daß es also zu deinen Ehren, zu seinem und des Nächsten Frommen täglich durch deine Gnade zunehmen möge, um deines Kindleins JEsu Christi, meines HErrn willen. Amen.

\section[128. Gebet der Eltern für ihre erwachsenen …]{128. Gebet der Eltern für ihre erwachsenen Kinder, daß ihnen Gott ein frommes Ehegemahl bescheren wolle}

Allmächtiger, gütiger, barmherziger Gott! der du den heiligen und züchtigen Ehestand selbst gestiftet, befohlen und gesegnet hast und willst, daß der Mensch nicht soll sein Lebenlang allein bleiben, und hast darum einem jeden seinen eigenen Gehilfen geschaffen und verordnet, mit dem er sich in diesem Leben in freundlicher Liebe und Treue nähren und in Zucht und Ehren sich mehren soll. Weil du auch uns Eltern im heiligen Ehestande Kinder, Söhne und Töchter, gegeben und gnädiglich bescheeret hast, die nun täglich daher wachsen, zu ihren mannbaren Jahren kommen, auch heute oder morgen nach deinem göttlichen Willen und Befehl in den Ehestand treten werden, und aber in deinem göttlichen Wort den Eltern auferleget und befohlen ist, daß sie fürnehmlich dafür sorgen, und ihren Söhnen Weiber nehmen, und ihren Töchtern Männer geben sollen: so bitten wir dich, gütiger und gnädiger Gott, erleuchte und regiere unsere Herzen und Gedanken mit deinem heiligen Geiste,
\quellseite{103}
daß wir unsere lieben Kinder mit gottseeligen und frommen Ehegenossen wohl versorgen mögen, und gib uns dazu gute Gelegenheit, und zeige uns zu rechter Zeit die Personen, so du hierzu in deinem Rath ersehen hast. Und weil ja ein frommer Ehegenoß allein von dir herkömmt, so befehlen wir unsere Kinder, du frommer und getreuer Gott! dir in deine göttliche Vorsorge, und bitten dich, du wollest jedem sein Theil zu rechter Zeit selbst zuschicken. Unterdeß behüte sie, lieber Gott, für allem Uebel, für Aergerniß, für böser Gesellschaft, daß sie ja nicht verführet werden, sondern in deiner Furcht sich keusch und züchtig halten, und die rechte Zeit zu ihrem Ehestande in Ehren erwarten, alles um deines lieben Sohnes, unseres HErrn JEsu Christi willen. Amen.

\section{129. Gebet einer Mutter für ungerathene Kinder}

Allmächtiger Gott, barmherziger Vater! wende dich zu dem Gebet und Flehen deiner armen Magd und vergieb mir um deines liebsten Sohnes JEsu Christi willen alle meine Sünde. Zürne nicht, daß ich mich unterwinde mein Herz vor dir auszuschütten, und in meinem großen Kummer und Traurigkeit mit dir zu reden. Denn es ist außer dir kein Gott, der mir helfen kann: ich weiß sonst keinen, zu dem ich meine Zuflucht nehmen könnte. Auf dich aber bin ich geworfen von meiner Mutter Leibe, du bist mein Gott von meiner Mutter Leibe an. HErr! Du erkennst bereits die Ursach meiner Schmerzen: vor dir ist alle meine Begierde und mein Seufzen ist dir nicht verborgen, du siehest ja und schauest den Jammer, welchen mein mütterliches Herz empfindet über den Ungehorsam und die greuliche Bosheit meines Sohnes (meines Kindes.) Ich habe denselben (dasselbe) von deiner Hand als eine theure Gabe empfangen, und dir hinwieder in der heiligen Taufe aufgeopfert. Ich habe ihn (es) durch herzliches Gebet dir täglich befohlen und, so viel an mir gewesen, nicht unterlassen ihn (es) zu allem Guten zu ermahnen, daß er (es) seinen Weg unsträflich gehen und sich halten soll nach deinem Wort. Aber wie schändlich hat er (es) seinen Weg verderbet, und durch den Teufel und böse Gesellschaft sich verführen lassen. Er (es) ist aus einem Gefäße der Ehren, so dir geheiliget war, ein Gefäß des Zornes worden, und wird ewig verderben müssen
\quellseite{104}
, wo du nicht seiner dich erbarmen, ihn (es) im Wasser deiner Gnaden waschen, mit dem Blut deines liebsten Sohnes JESU Christi reinigen, und die arme Seele retten wirst. Nun Herr! so stehe ihm (es) denn, um deines allerliebsten und gehorsamsten Kindes JESU Christi willen, gnädig an; gedenke nicht der Sünden seiner Jugend und aller seiner Uebertretung. Nimm dich seiner armen Seele herzlich an, daß sie nicht verderbe, sondern als ein Brand aus dem Feuer gerissen und für der ewigen Verdammniß bewahret werde.

JEsu, liebster Heiland! du hast keinen Gefallen am Tode des Gottlosen, sondern willst, daß er sich bekehre von seinem bösen Wesen und lebe. Darum bist du auch deinem Vater im Himmel gehorsam worden bis zum Tode am Kreuz, und hast für aller Menschenkinder Ungehorsam gebüßet und bezahlet, und mit deinem Blut am Kreuz versöhnet alles, es sei auf Erden, oder im Himmel. Ach so laß doch dein Blut und Tod an diesem gottlosen Menschen nicht verloren sein, sondern regiere ihn (es) durch deinen Geist, daß er (es) seine Uebertretungen erkenne, die zeitlichen und ewigen Strafen der ungehorsamen Kinder und unbußfertigen Sünder wohl bedenken und sich rechtschaffen zu dir bekehren möge.

Herr! du ließest ja vormals dein Herz brechen und weintest selbst über deinen Freund, den Lazarus, als er gestorben war und schon im Grabe lag, wie willig wärest du, denselben zu diesem Leben wieder aufzuwecken? Ach! so wirst du demnach die Thränen und Seufzer deiner armen Magd dir lassen zu Herzen gehen, indem ich mit Rahel meinen Sohn (mein Kind) beweine, weil er (es) in Sünden lebendig todt ist. JEsu, du kannst meine Thränen stillen, meine Traurigkeit und Herzeleid in Freude verkehren; du hast den Schlüssel des Todes und der Höllen in deiner Hand; so du willst, kannst du diesen Sünder, der in Sünden und Uebertretung todt ist, wohl wieder lebendig machen. Nun so rufe denn kräftig in sein Herz: wache auf, der du schläfest, und stehe auf von den Todten, daß dich JEsus Christus erleuchte! Hilf, daß er (es) solches höre, das Laster- und Sündengrab verlasse, ein neues Leben anfange, und durch dein Verdienst zum ewigen Leben im Himmel erhalten werde.

Herr Gott Heiliger Geist! der du aller Menschen Herzen in
\quellseite{105}
deinen Händen hast, und sie regierest, wie du willst. Du hast wohl eher einen verlornen Sohn dahin gebracht, daß er seine lange getriebene Bosheit erkennet und herzlich betrauert hat, zu seinem Vater wieder umgekehret, und ein anderer Mensch geworden ist. Dich rufe ich von Grund meines Herzens an, neige doch meines ungerathenen Sohnes (Kindes) Herz und Gemüth zum Guten, ändere seinen Sinn, richte seine Füße wieder auf den Weg der wahren Gottseligkeit, und mache einen solchen Menschen aus ihm, der in deinen Geboten hinführo wandele, deine Rechte halte und darnach thue. Gieb, daß er (es) die böse Gesellschaft der verkehrten Menschen verlasse, ehestens mit jenem verlornen Sohn wiederkomme und sage: ich habe gesündigt in den Himmel und vor dir, und bin hinfort nicht werth, daß ich dein Kind heiße; damit ich Ursache habe mit allen, die in meinem Hause sind, mich zu freuen und zu sagen: Laßt uns fröhlich sein, denn dieser mein Sohn (dieses mein Kind) war todt, und er (es) ist lebendig worden, er (es) war verloren, und ist wiedergefunden worden. So wird auch Freude sein im Himmel vor den Engeln Gottes über diesen Sünder; so will ich dich für diese Gnade aus allen Kräften preisen und rühmen in Zeit und Ewigkeit.

HErr, du dreieiniger Gott! erbarme dich über mich; erhöre mein Gebet, tröste und stärke mich in meinem großen Kreuz. Erfreue mich wieder mit deiner Hilfe, um deines Namens Ehre willen. Amen.

\section{130. Trostgebet in Anfechtung bei seinem Beruf}

Lieber HErr, ich habe ja dein Wort und bin in dem Stande, der dir gefället, das weiß ich. Nun siehest du, wie es allenthalben mangelt, daß ich keine Hülfe weiß, ohne bei dir; darum hilf du, weil du gesagt und befohlen hast, daß wir sollen bitten, suchen und klopfen, so sollen wirs gewißlich empfahen, finden und haben, was wir begehren. Amen.

\section{131. Ein Anderes}

Lieber Gott, was ich jetzt thue, das will ich im Namen Jesu thun und in dem Gehorsam, darein ich von Gott gesetzt bin, und will's mit Freuden thun; ob mir etwas darüber widerfä[h]ret und der Teufel mir zusetzet, was schadet mir's? Dennoch bin ich in dem
\quellseite{106}
Stande, da Gottes Wort mich lehret und tröstet, was ich thue oder leide, daß sei wohlgethan, und Gott Wohlgefallen daran haben und mit Gnaden bei mir sein. Amen.

\section[132. Reisegebet eines Hausvaters, der mit …]{132. Reisegebet eines Hausvaters, der mit seiner Familie seinen Wohnort verändert}

Getreuer, himmlischer Vater! der du vor Zeiten dem Abraham aus seinem Vaterlande und von seiner Freundschaft in ein anderes, ihm unbekanntes Land zu reisen befohlen, und ihn daselbst zu einem großen Volke gemacht hast, auch den Jakob aus dem gelobten Lande nach Mesopotamien, dann nach Egypten und seine Nachkommen, die Kinder Israel, mit starker Hand und ausgestrecktem Arm durchs rothe Meer und durch die grausame Wüste wieder herausgeführet, und ihnen auf der ganzen Reise selbst des Tages in einer Wolken- und des Nachts in einer Feuersäule vorgegangen; desgleichen die Weisen aus Morgenland bis gen Bethlehem vermittelst des Wundersternes sicher begleitet hast! Du hast mir auch nach deinem heiligen, wunderbaren Rath und Willen eine schwere, ferne, gefährliche Reise auferlegt, und heißest mich ziehen in ein solches Land, das ich nicht kenne. Weil es aber dein göttlicher Beruf und Ordnung ist, so will ich dir, nach dem Exempel Abrahams, durch den Glauben gehorsam sein, und die Wallfahrt sammt meinem lieben Eheweibe und Kindern in deinem Namen anfangen, führen und vollenden. O HErr! behüte unsern Aus- und Eingang, sei du unser Schatten über unserer rechten Hand, daß uns des Tages die Sonne nicht steche, noch der Mond des Nachts. Laß deine heiligen Engel sich um uns her lagern, und eine Wagenburg um uns schlagen, daß uns zu Land und Wasser kein Uebel begegne, und keine Plage sich zu uns nahe. Befiehl ihnen, den starken Helden, daß sie uns behüten auf allen unsern Wegen, daß sie uns auf den Händen tragen, und wir unsere Füße nicht an einen Stein stoßen, oder in die Hände böser Leute gerathen. Beschere uns hingegen fromme Christen, die sich unser treulich annehmen. Ach herzlieber Vater, wir sind ja schwach und blöde, und wissen nicht, was wir thun und wie wir diese schwere Reise angreifen sollen. Aber unsre Augen sehen allein nach dir, HErr, du bist unsre Stärke und Zuversicht,
\quellseite{107}
unser Fels, unsre Burg, unser Erretter, unser Gott, unser Hort, auf den wir trauen, unser Schild und Horn unsers Heils und unser Schutz, der du alle unsre Gänge, ja alle unsre Härlein auf dem Haupte gezählet, daß deren keines ohne deinen Willen umkommen, noch auf die Erde fallen kann. Darum, ob wir schon wandern im finstern Thal, und viele hundert Tausend sich umher wider uns legen, fürchten wir doch kein Unglück, denn du, HErr, bist bei uns; du bist der Schild für uns, und der uns zu Ehren setzt, und unser Haupt aufrichtet. Ach HErr! der du aller Menschen Herzen in deiner Hand und Gewalt hast, neige die Herzen der Leute, zu denen wir kommen sollen, zu uns, daß sie die Werke der Barmherzigkeit an uns erzeigen. Erfreue uns nun wieder, nachdem du uns so lange geplaget, nachdem wir so lange Unglück gelitten. Sei und bleibe, liebster Vater, bei uns, deinen Kindern, auf dem ganzen Wege, und erhalte uns durch die rechte Hand deiner Gerechtigkeit; sei du unser treuer Geleitsmann und Wegweiser, daß wir mit Friede und guter Gesundheit kommen an den Ort, so du uns verordnet hast, auch daselbst, wenn es dein heiliger Wille, in gutem Alter unser Leben beschließen, durch JEsum Christum. Amen.

\section{133. Gebet eines Vormundes}

Allmächtiger, ewiger Gott, der du bist ein Vater der Waisen und Richter der Wittwen, auch ein schneller Zeuge sein willst wider die, so Gewalt und Unrecht thun den Wittwen und Waisen; und daher es dir gefallen läßt, wenn sich welche finden lassen, die den Waisen Recht schaffen und der Wittwen Sachen helfen, auch durchgehends solcher Leute Nothdurft sich annehmen, und, was dieser Geringsten einem geschieht, annehmen willst, als wäre es dir selber geschehen. Siehe, ich habe mich überreden lassen, dergleichen Sorgfalt auf mich zu nehmen und zu sorgen für Kinder, welche von Vater und Mutter verlassen sind. Dieweil du nun von solchen Haushaltern nicht mehr erforderst, denn daß sie mögen treu erfunden werden: so erfülle auch selbst mein Herz gegen diese meine Pflegkinder mit Liebe und Treue. Für allen Dingen laß mich wohl zusehen, daß sie aufgezogen werden in der Zucht und Vermahnung zum HErrn, damit sie zunehmen in der Gottesfurcht und
\quellseite{108}
wachsen in allen Stücken an dem, der das Haupt ist, Christus Jesus, vornehmlich da auch die Gottseeligkeit zu allen Dingen nütze ist und die Verheißung hat dieses und des zukünftigen Lebens, auch den Kindern nichts Besseres kann eingeschärft werden, als die Erkenntniß Gottes, in welcher das ewige Leben beruht und ist. Neige mein Herz zu deinen Zeugnissen und nicht zum Geiz, daß ich nicht Gewalt treibe mit ihrem Erbe, noch zu mir reiße ihre Aecker und nehme ihre Häuser, sondern mit ihrem Gut und Vermögen so umgehe, wie ich es getreu gegen dir, allwissender Gott, und den Menschen zu verantworten vermag, wenn ich werde müssen für diese meine Haushaltung Rechnung thun. Laß mich vielmehr ihr Erbgut wider alle Gewalt nach Möglichkeit beschützen und sie erretten von der Hand aller, die sie hassen, damit ihnen niemand Unrecht thue, noch sie jemandes Beute sein dürfen. Und weil du nicht verachtest des Waisen Gebet, so erhöre sie von deiner heiligen Höhe und segne sie mit Segen oben vom Himmel herab. Sonderlich aber erhalte sie in deiner Furcht, daß sie all ihr Leben lang dich vor Augen und im Herzen haben und in keine Sünde willigen, noch wissentlich thun wider deine Gebote, damit ich dermaleins diese meine Mündlein mit in Himmel bringen und sagen kann: Siehe, hier bin ich und die Kinder, die mir der HErr gegeben hat, und sie nebst mir eingehen in das ewige Reich deiner Herrlichkeit. Amen.

\section{134. Gebet der Kinder}

Allmächtiger Gott, himmlischer Vater! dieweil du mir im vierten Gebot so ernstlich befohlen, meinen Vater und meine Mutter zu ehren, auch verheißen hast, solchen Dienst mit langem Leben und Wohlergehen aus Gnaden zu belohnen: so bitte ich dich von Herzensgrund, du wollest mir geben ein gehorsames Herz, daß meinen lieben Eltern und allen denen, die an der Eltern Statt mir vorgesetzt sind, in allen Dingen mit gebührlicher Unterthänigkeit gehorche, sie mit der That, mit Worten und Geduld ehre, und sie ja nicht betrübe, verachte und verspotte, so lange ich lebe, auf daß der Segen über mich komme. Gib, daß ich ihre Worte, Gebote und Gesetze wohl bewahre, und in meinem Herzen fest behalte, ja auf mein Herz allezeit binde, und an meinen Hals hänge, auch ihre
\quellseite{109}
Strafe annehme, und ihre Lehre ja nicht vergesse, damit ich desto länger lebe. Regiere mich auch mit deinem guten Geist, daß ich der Zucht meines Vaters gehorche, und das Gebot meiner Mutter nicht verlasse, damit ich ihnen eine Freude und Trost sein möge. Gib, daß ich mir in allem meinen Thun und Lassen das Exempel meines HErrn und Heilandes fleißig einbilde, der allen gottseligen Kindern zum Vorbilde seiner lieben Mutter, der Jungfrau Maria, und seinem Pflegvater, Joseph, unterthan gewesen: Verleihe mir auch ein dankbar und demüthiges Herz, daß ich meinen lieben Eltern gleiches vergelte, und wieder diene, sonderlich aber ihrer im Alter pflege, und ihnen zu gute halte, wenn sie kindisch und wunderlich werden, und ja nicht vergesse, daß ich von ihnen geboren bin. Vergib mir auch gnädiglich, was ich aus kindlicher Schwachheit oder Unverstand wider sie gethan. Behüte mich für böser Gesellschaft, und leichtfertigen Leuten, daß, wenn mich die bösen Buben locken, ich ihnen nicht folge, noch den Weg mit ihnen wandele, sondern meinen Fuß vor ihrem Pfad wahre, damit ich nicht in Schand und Laster gerathe. O HErr! du hast mich bisher gezogen, wie ein Mann seinen Sohn zeucht, gib, daß ich meine Ohren neige, und das Wort der Weisen höre, und also in meiner Jugend wohl lerne, damit ich vernünftig und vorsichtig werde. Hilf, daß ich zunehme an Weisheit, Alter und Gnade bei dir und den Menschen, um JEsu Christi willen. Amen.

\section{135. Gebet eines Dienstboten}

Barmherziger Gott, der du mich durch deines lieben Sohnes JEsu Christi theures Leiden und Sterben von der ewigen Dienstbarkeit, von der Gewalt der Sünden, von der Obrigkeit der Finsterniß, von der grausamen Tyrannei des Teufels erlöset und befreit, zum HErrn über Tod und Hölle gemacht hast: ich bitte dich, du wollest mir Gnade geben, daß ich an meinem Stand der leiblichen Dienstbarkeit, darein du mich hier auf Erden nach deinem Willen und Wohlgefallen gesetzt hast, keinen Verdruß gewinne, und nicht etwa wider deine Ordnung in Ungeduld murre, auch andern Leuten ihren höheren Stand nicht mißgönne, sondern daß ich solchen deinen Willen thue von ganzem Herzen mit gutem Willen und nicht anders
\quellseite{110}
gedenke, als diente ich dir, Gott im Himmel, und nicht den Menschen auf Erden. So hilf nun, lieber Gott, daß ich dir, dem höchsten HErrn im Himmel und auf Erden in wahrer Erkenntniß und rechter Furcht diene, dich über Alles liebe, alle meine Hoffnung und Seligkeit in dich setze und in deinen Geboten untadelich wandele, darnach auch, daß ich meinem leiblichen Herrn und Frauen, nicht allein den gütigen und gelinden, sondern auch den wunderlichen und ungeschlachten, in Geduld gehorsam sei in allen Dingen, die nicht wider dich sind, mit aller Furcht in Einfältigkeit meines Herzens, als Christo, meinem HErren, nicht mit Dienst allein vor Augen, als den Menschen zu gefallen, sondern von Herzensgrund um deines Befehles und Gebotes willen. Gib Gnade, daß ich treu erfunden werde in Allem, das mir befohlen ist und angethan wird, nichts verwahrlöse und zu Schanden mache, auch keinen Unrath und Schaden durch Unachtsamkeit geschehen lasse, dazu daß ich mich fremdes Guts, so mir vertraut und unter die Hände gegeben, nicht lasse gelüsten, oder dasselbige veruntreue. Bewahre mir meine Gesundheit, stärke meine Gliedmaßen und alle Kräfte meines Leibes, dazu begabe mich mit Vernunft und Weisheit, daß ich die Arbeit meines Herren und Frauen könne versorgen, ihre Nahrung durch deine göttliche Hülfe bessern, auf daß sie durch meinen Fleiß gesegnet und in die Menge ausgebreitet werden, und Alles, was ich thue und ausrichte, daß es gereiche zu deinen göttlichen Ehren und mir zu einer seligen Uebung meines Glaubens um JEsu Christi, deines lieben Sohnes willen, welcher mit dir lebet und regieret in Einigkeit des Heiligen Geistes, von Ewigkeit zu Ewigkeit. Amen.

\section[136. Gebet junger Leute, die sich in den …]{136. Gebet junger Leute, die sich in den Ehestand zu begeben Willens sind}

Siehe, lieber Gott, ich höre, daß der Ehestand dein geschaffen Werk ist und dir wohlgefällt; derohalben will ich mich auf dein Wort darein begeben, es gehe mir darin, wie du willst, so soll es mir alles gefallen und behagen. Ach lieber Gott und Vater unsers HErrn JEsu Christi, beschere und gib mir armen Kinde einen frommen Mann (ein frommes Weib), mit dem ich göttlich durch die Gnade des Heiligen Geistes im Ehestande leben mag. Amen.

Gebetsschatz.

\quellseite{111}

\section{137. Ein Anderes}

HErr Gott, du hast mich zu einem Manne geschaffen, du siehest, daß ich nicht keusch leben kann; ich rufe dich an und bitte, du wollest mein Vorhaben regieren und Glück dazu geben. Gib du guten Rath und hilf mir, erwähle du mir Eine, mit der ich ehrlich leben möge und dir dienen, und durch den Glauben und das Gebet Unglück und Beschwerung, so sich im Ehestande findet, mag überwinden. Amen.

\section{138. Gebet junger Leute um ein frommes Ehenmahl}

Du allmächtiger, ewiger Gott! Schöpfer, Erhalter und Mehrer des ganzen menschlichen Geschlechtes, der du den heiligen Ehestand im Paradiese eingesetzt, auch denselben durch deinen geliebten Sohn, Jesum Christum, unsern Heiland, zu Cana in Galiläa mit seinem ersten Wunderzeichen geehret hast, auch mich dazu erschaffen, und mir das Alter gegeben hast: ich bitte dich herzlich, du wollest mir einen vernünftigen, christlichen und gottesfürchtigen Ehegenossen bescheren, welchen ich herzlich lieb habe, und bei dem ich in wahrer Gottesfurcht und christlichem Wandel friedsam und einträchtig auf Erden leben, und mich ehrlich nähren möge, dir zu Lob, andern zum guten Exempel, uns auch und jedermann zu Nutz und Wohlfahrt an Leib und Seele. Du bist ein Herzenskündiger, weißt aller Menschen Gemüth und Eigenschaft; von dir kommt ein vernünftiges Gemahl, denn es ist ein Geschenk und Gabe des Allerhöchsten: derohalben schreie ich zu dir, solchen zu bekommen, und bitte herzlich, du wollest, lieber HErr Gott, mein und meiner lieben Eltern Herz recht zu erwählen und zu erhalten, mit deinem Heiligen Geist erleuchten, und dir die Sache in deine väterliche Sorge lassen befohlen sein um Jesu Christi willen. Amen.

\section[139. Gebet eines jungen Mannes, der im …]{139. Gebet eines jungen Mannes, der im Begriff steht, um ein Gemahl anzuhalten}

Gnädiger, barmherziger Gott! himmlischer Vater! ich habe nun, nach lang gehaltenem Gebet, mir ein Ehegemahl ausersehen, mit dem ich mich (meines Bedünkens) wohl nähren und eine christliche
\quellseite{112}
Ehe besitzen wollte. Sollte es nun diese Person N. N. sein, die du mir zugeordnet hast, und soll solche Ehe gereichen zu deinen göttlichen Ehren und andern Leuten zum guten Exempel, mir und ihr zur Wohlfahrt an Leib und Seel’, so wollest du Gnade dazu geben, daß es möge durch ordentliche christliche Mittel vollzogen werden. Wo es aber sollte dir, o Gott! zu Unehren, andern zum Aergerniß, und mir selbst zur Hinderniß meiner Seelen Seligkeit gereichen, so ändere, Herre Gott! mein Herz und Gemüth, daß ich wider deinen Willen nichts begehren, noch unterstehen möge. Dein Wille, Herr Gott! geschehe, und laß du dir diese Sache befohlen sein, und verleihe deinen Heiligen Geist dazu. Amen.

\section[140. Gebet einer Jungfrau, die zur Ehe …]{140. Gebet einer Jungfrau, die zur Ehe begehret worden}

Ewiger, allmächtiger und gütiger Gott, mein himmlischer Vater, weil du mich hast zum Ehestand erschaffen, und fället mir jetzt dazu eine Gelegenheit vor, daß dieser N. meiner in Ehren begehret. Ist es nun mit diesem Gesellen dein göttlicher Wille, und sollte er es sein, den du mir zum lieben Ehemann versehen und erwählet hast, und soll es auch dir zu Ehren und mir zur Wohlfahrt gereichen, so bitte ich dich herzlich, du wollest dazu Gnade und Segen geben, daß diese Ehrensache einen glücklichen Fortgang gewinne. Und weil ich merkte, daß dieses ehrlichen Gesellen Herz und Sinn in Liebe und Ehren zu mir stehet, so wende gleicher Gestalt auch mein Herz und Gemüth zu ihm, daß ich ihn wieder in Zucht und Ehren lieb habe, und auf meiner lieben Eltern Rath und Gutachten mich ganz willig darein ergeben möge. Wenn es aber nicht sein sollte, und ich mit diesem Gesellen nicht wohl versorget sein möchte, Herr, so hindere dieses Werk, und wende seine Gedanken von mir ab, und beschere ihm und mir sonst jedem sein Theil nach deinem gnädigen Willen und Wohlgefallen, dahin ich Alles stelle, und befehle dir, himmlischer Vater, und meinen lieben Eltern diese Sache ganz und gar; schicke es, wie du es auf beiden Theilen für’s Beste erkennst. Was du mir nur durch meine herzlieben Eltern hierin rathen und sagen lassen wirst, dem will ich gern folgen, und nicht zweifeln, du werdest mich auf mein demüthiges Gebet wohl versorgen um deines allerliebsten Sohnes JEsu Christi willen. Amen.

\quellseite{113}

\section{141. Gebet im Brautstande}

O du gnädiger, gütiger und barmherziger Gott, Vater unsers Herrn Jesu Christi, Schöpfer Himmels und der Erden, der du den heiligen Ehestand aus sonderbarem Rathe deiner göttlichen Weisheit hast eingesetzt und verordnet, daß Mann und Weib in ehelicher Gemeinschaft beisammen wohnen sollen, vornehmlich zur Erschaffung der Menschen und zur Erfüllung des Himmelreichs. Wir sagen dir von Herzen Lob und Dank, daß du uns dazu bereits so weit gebracht, daß wir nun als eheliche Verlobte zur Vollziehung unsers Ehestandes mehr nicht als der ordentlichen, öffentlichen Zusammengebung oder Kopulirung durch des Priesters Hand zu gewarten haben. Darum bitten wir dich demüthig, erhalte uns unterdeß beiderseits bei guter Gesundheit, daß wir unsern christlichen Ehrentag und öffentlichen Kirchgang mit Freuden erleben und christlich verrichten können. Wollest uns auch dazwischen Gnade und Segen geben, daß wir als zukünftige Eheleute bedenken und zu Gemüthe führen, welcher Gestalt wir unsern ehelichen Stand christlich anfangen, und selig darin leben, und dir und unsern Nächsten mit gutem Gewissen darin dienen und dich allezeit ehren und preisen mögen durch unsern Herrn Jesum Christum, der mit dir lebet und regieret, gleicher Gott, hochgelobet in Ewigkeit. Amen.

\section{142. Gebet Verlobter am Trauungstage}

Allmächtiger Gott, himmlischer Vater! der du den heiligen Ehestand selbst eingesetzt, verordnet und gesegnet hast, wir sagen dir Lob und Dank, daß du uns auch dazu berufen, und bis auf diese Zeit gnädiglich für Schaden behütet hast. Weil wir aber wissen, daß der Teufel dieser deiner Ordnung mächtig feind ist, so bitten wir heut zu Tage, du wollest uns deinen heiligen Engel zuordnen, der uns für des Teufels oder Asmodi Mord und Lügen, in diesem unsern Stande, schütze und bewahre. Wollest uns beiden auch durch deinen heiligen Geist dahin weisen und regieren, auf daß wir für allen Dingen deine göttlichen Gebote lieben, dir mit treuem Herzen dienen, und daß wir beide Eheleute, ein Theil dem andern alle eheliche Treue, Liebe, Ehre und Gutes erweisen, lange Zeit in
\quellseite{114}
Gesundheit, Friede, Ruhe und Einigkeit vertraulich mit einander leben, nach deinem göttlichen Willen Kinderlein zeugen, und alles unser Thun und Lassen allein zu deinem göttlichen Lob, Ehre und Preis richten, auf daß wir nach diesem Leben, durch deine göttliche Gnade, das ewige Leben mit einander besitzen mögen. Amen.

\section{143. Gebet eines Ehemanns}

Allmächtiger, gütiger Gott! der du den heiligen Ehestand selbst eingesetzt und durch deines lieben Sohnes Jesu Christi erste Wunderzeichen geehrt und gezieret hast als einen Stand, der dir angenehm ist, in welchem auch viel heilige Erzväter und Propheten gottseelig gelebt, und dir wohl gefallen haben. Weil denn du mich auch in der heiligen Ehe berathen, zur Haushaltung verordnet und ein sonderliches Wohlgefallen an den dreien Stücken hast, nämlich: wenn Brüder eins sind, und die Nachbarn sich lieb haben, und Mann und Weib sich mit einander wohl begehen, so bitte ich dich von Herzensgrund, verleihe mir, daß ich in christlicher Liebe und Einigkeit mit Vernunft bei meinem Weibe als dem schwächsten Werkzeug wohne, derselben ihre Ehre, als auch Miterbin der Gnade des Lebens gebe, sie sammt Kindern und Gesinde ziehe zu deiner Erkenntniß und göttlichen Ehre in aller Zucht und Ehrbarkeit. Dazu so gib Gnad, daß sie mir in allem Guten und zu aller Gottseligkeit folgen und sich ziehen lassen. Wehre dem Eheteufel, daß er nicht Zwietracht und Zank zwischen uns einmenge, und, wo wir etwa, aus Schwachheit übereilt, uneins würden, so hilf, daß wir uns bald wieder mit einander versöhnen. Gib mir Gnade, daß ich mich keines andern Ehegemahls und Weibes gelüsten lasse, oder dieselbe mit einem bösen Auge ansehe, ihrer zu begehren. Behüte mich, mein Weib, Kinder und Gesinde vor Krankheit nach deinem göttlichen Willen. Du wollest auch mir, deinem Knechte, verleihen, daß ich meines Berufes fleißig warte, im Schweiße meines Angesichts mein Brod esse, und mich es nicht verdrießen lasse, ob es mir sauer werden muß, denn du hast es also geschaffen. Verleihe auch Glück und Heil zu meiner Nahrung, daß dieselbe durch deinen Segen ohn’ anderer Leute Schaden und Nachtheil gemehret werde. Beschwere
\quellseite{115}
mir frommes Gesinde und treue Arbeiter. Behüte mir Haus und Hof und alles, was du mir gegeben hast. Hilf uns auch das Kreuz in unserm Stande geduldig tragen, und nach diesem Leben versammle uns in dein Reich zu allen gottfeligen Eheleuten. Amen.

\section{144. Gebet einer Ehefrau}

Gütiger Gott, getreuer Vater! weil du mich aus Gnaden in den heiligen Ehestand zur Haushaltung berufen hast, darinnen ich dir auch, meinem Gott, dienen und gefallen kann, so gib mir, deiner Dienerin, Gnade, daß ich deine göttliche Furcht stets vor meinen Augen habe und dich, meinen Schöpfer und Erlöser, über alle Dinge liebe, und dir vertraue; nach dir aber meinen Mann fürchte, ehre und liebe, und mich keines andern gelüsten lasse. Hilf, daß nach deinem Gebot mein Wille meinem Mann unterworfen sei, demselben zu gehorchen in aller Willigkeit, und daß der verborgene Mensch meines Herzens unverrückt mit sanftem, stillen Geiste und mit allerlei Tugenden geschmückt sei, wie vor Zeiten die heiligen Weiber gewesen, die ihre Hoffnung auf Gott setzten und ihren Männern gehorsam waren. Gib mir, deiner Magd, einen keuschen, züchtigen Wandel in der Furcht und Demuth, daß ich in aller Gottseligkeit mit freundlichen und sittigen Worten den Zorn und Unmuth meines Hauswirthes ablehnen und versöhnen, und ihm mit Glimpf begegnen könne, dazu meine Kinder und Gesinde mit Sanftmuth zu Lob und Ehre deines heiligen Namens aufziehe, und daß sie mir folgen mit willigem Herzen zu allem Guten. Hilf auch, daß ich meinem Manne eine getreue Gehülfin sei, und die Nahrung, so du uns aus Gnaden bescherest, fleißig zusammenhalte, und nichts verwahrlos[e], sondern arbeiten und schaffen helfe, was mir zuständig ist in meinem Beruf, auf daß ich habe zu geben den Dürftigen, und meine Hand ausbreiten könne zu den Armen. Behüte uns vor untreuen Arbeitern und bösem Gesinde, die unsere Nahrung schmälern und vergeuden. Verleihe mir auch Gnade, daß ich das Kreuz im Ehestand geduldig trage und nicht sobald schüchtern oder abgeschreckt werde, wo sich ein Leiden erhöhe, indem durch Anfechtung und Widerwärtigkeit unser Glaube bewähret wird. O Herr Gott! in deine allmächtige Gewalt befehle ich dir
\quellseite{116}
mich, meinen lieben Ehemann, alle meine Kinder und Gesinde, du wollest uns behüten vor Sünden, Schaden und allem Leid, durch Jesum Christum, unsern Herrn. Amen.

\section{145. Gebet um Ehejeegen}

Herr, du bist der allmächtige Gott und die Erde ist voll deiner Güte, all dein Thun ist Wahrheit und deine Wege sind recht. Du hast selber geredet: Seid fruchtbar und mehret euch und erfüllet die Erde, welches auch nicht soll aufhören, so lange die Erde steht, dannenhero ist es zuweilen deine Strafe, unfruchtbare Leiber und versiegende Brüste haben und die heiligen Weiber haben es vor Zeiten, ehe Gott im Fleisch geoffenbart ist, für Elend und Schmach gehalten, wenn sie ohne Kinder dahingegangen. Darum haben sie, wenn ihr Herz um deßwillen sich so übel gehabt, zu dir, dem Herrn, gebetet, und ihr Herz vor dir ausgeschüttet und du hast ihnen ihre Bitte gegeben; denn du, Herr, kannst schaffen, was du willst, und kein Ding ist bei dir unmöglich. Du bist es, der die Unfruchtbare im Hause wohnen macht, daß sie eine fröhliche Kindermutter wird, ja du kannst auch denen, so alt und betagt sind, ein Lachen zurichten. Denn Kinder sind, o Herr, deine Gabe und Leibesfrucht ist dein Geschenk. Deswegen kehre auch ich mein Angesicht zu dir, mein Herr, zu dir hebe ich meine Augen auf, der du überschwänglich thun kannst, mehr als wir bitten oder verstehen. Ich, die ich im Geschrei bin, daß ich unfruchtbar sei, erkenne dich für einen Gott, der mir meine Leibesfrucht wohl geben kann, und erbitte es von dir nach deinem Willen, mein Gott, erhöre mich. Herr Zebaoth, du wollest an mich gedenken und deiner Magd nicht vergessen. Wunderlich sind ja deine Werke, das erkennet meine Seele wohl. Deine Augen sehen, was noch unbereitet ist, so befehle ich auch hierinnen dir meine Wege und hoffe auf dich, du wirst es wohl machen. Du bist der Herr, thue, was dir wohlgefällt. Wirst du nun deine Magd ansehen, so will ich dir, mein Gott, lobsingen. Soll ich aber kein Kind haben bis an den Tag meines Todes, so geschehe, Herr, dein Wille. Du weißt alles, was heimlich und verborgen ist. Es ist doch besser ohne Kinder sterben, als gottlose Kinder haben, fürnehmlich in diesen letzten Tagen, da greuliche
\quellseite{117}
Zeiten gekommen, in welchen viel leben, die den Eltern ungehorsam sind, da die Bosheit den Verkauf verkehrt und falsche Lehre die Seele betrügt. Auch bei unschuldigen Herzen, da auch die Mütter durchs Schwert an manchen Orten beraubet werden, also daß sie oft anfangen und sagen: Selig sind die Unfruchtbaren und die Leiber, die nicht geboren haben und die Brüste, die nicht gesäuget haben. Für allen Dingen aber hilf mir, daß Jesus Christus eine Gestalt in mir gewinne, und daß ich fruchtbar sei zu allen guten Werken, bis wir in der Auferstehung gleich sein werden den Engeln Gottes im Himmel. Höre, Gott, das Gebet deiner Magd und ihr Flehen um des Herrn willen. Deinem Namen sei ewig Ehre und Lob. Amen.

\section{146. Gebet bei verstürtem Ehelegen}

Herr, allmächtiger Gott, Schöpfer Himmels und der Erden! Du hast deinen göttlichen Segen über den heiligen Ehestand gesprochen, und hast gewollt, daß von einem Blute aller Menschen Geschlechter auf dem ganzen Erdboden wohnen. Ja, du bist allein derjenige, so die Unfruchtbare im Hause wohnend machet, daß sie eine fröhliche Kindermutter wird. Denn Kinder sind eine Gabe des Herrn, und Leibesfrucht ist ein Geschenk. Solches deines Segens, der noch unter meinem Herzen verborgen lieget, bin ich mit Ehren durch deine Gnade auch theilhaftig geworden. Dieweil du aber dem weiblichen Geschlechte um der Sünde willen auferlegst, daß sie sollen mit Schmerzen Kinder gebähren, und es mit einem schwängern Weibe gar gefährlich stehet, so bitte ich deine milde Güte, im Namen deines einigen und liebsten Sohnes, Jesu Christi, welcher auch mir und meiner Frucht zu Gute, aus Ueberschattung des Heiligen Geistes, von der reinen und keuschen Jungfrau Maria empfangen und geboren worden, und dadurch unsre sündhafte Empfängniß und Geburt geheiliget hat, du wollest mich deinen Willen lassen erkennen, mir meine Schmerzen in Geduld helfen tragen, meine Leibesfrucht stärken, dieselbe segnen, und mit deinem Heiligen Geiste begaben. Gib mir deine lieben Engel zu Geleitsleuten. Bewahre mich auf Wegen und Stegen vor allem Unfall. Laß mich nicht fallen oder unvorsichtig stoßen. Wende meine Augen
\quellseite{118}
ab, daß ich mich nicht versehe an irgend einer Ungestalt. Laß mich keine Mißgeburt zur Welt gebähren, oder meine Frucht vor der Zeit abgehen. Entbinde mich zur rechten Zeit in Gnaden. Erfreue mich und meinen Ehemann mit einem fröhlichen Anblick, und laß dem armen Kindlein die heilige Taufe widerfahren. Denn ich ergebe mich dir ganz und gar mit Leib und Seele, mit meiner Frucht und allen Gliedern. Du hast mich erschaffen und erlöst, dein bin ich, auf dich hoffe ich, mein Heil steht in deinen Händen, dir will ich leben und sterben. O HErr! erhöre mich, und laß mich in meiner Hoffnung nicht zu Schanden werden. Amen.

\section[147. Gebet, wenn die Frau sich vor der Geburt …]{147. Gebet, wenn die Frau sich vor der Geburt fürchtet}

HErr JEsu Christe! der du lebendig bist von Ewigkeit zu Ewigkeit, und hast die Schlüssel des Lebens und des Todes, siehe! meine Stunde ist kommen, die Traurigkeit bricht herein, ach! wie bange ist mir, daß mir das Herz im Leibe davon wehe thut. Ich fühle billig den Schmerz, der dem weiblichen Geschlecht verkündigt ist, ehe noch ein Weib geboren hatte, denn durch dasselbe ist die Uebertretung in die Welt gekommen; dir sei es geklagt, du treuer Heiland! der du durch deine heilige Geburt und Menschwerdung die sündliche Geburt der Menschen geheiligt, ja der du die ganze werth(e) Christenheit durch das Exempel eines gebährenden Weibleins aufgerichtet hast, daß sie in ihrer Traurigkeit nicht verzagen, sondern der Freude, so bald darauf folgen wird, erwarten sollen, ach! stärke mich, HErr! in dieser Stunde, wende dein Angesicht nicht von mir, laß mich dich wiedersehen, und sich mein Herz freuen, daß du so gerne hilfest. Es ist allein dein Werk, du bist allein der mächtige HErr, der uns aus Mutterleibe ziehen kann. Darum fliehe ich zu dir, und bitte dich um deines Blutes und Todes willen, stehe mir bei in diesen Nöthen, wo Menschen-Hülfe nichts nütze ist. Sende mir die Kraft des heiligen Geistes, daß er in mir erwecke ein herzliches Vertrauen auf deine Gnade, der meinen Glauben stärke und christliche Geduld verleihe. Allmächtiger und gnädiger Gott! der du befohlen hast, wir sollen dich anrufen in der Zeit der Noth, siehe an meinen Jammer und Elend, laß mich nicht trost- und hülflos, du siehest ja, denn du schauest das Elend und Jammer,
\quellseite{119}
es stehet in deinen Händen, darum erweise deine wunderliche Güte und Allmacht an mir. Der du mich mit Liebesfrucht gesegnet, und mich mit Hoffnung erfreuet hast, laß mich darinnen nicht zu Schanden werden, beschere mir einen fröhlichen Anblick, und nach dem Heulen und Weinen überschütte mich mit Freuden. Wäre es aber ja dein heiliger Wille, daß mein letztes Stündlein vorhanden sei, und ich über dieser Geburt meinen Geist aufgeben sollte, ei so verleihe Beständigkeit, schone des armen Würmleins, das du erlöst hast, laß es lebendig zur Welt kommen und durch die heilige Taufe in das Buch des Lebens eingezeichnet werden. Tröste und erquicke auch diejenigen, so durch meinen Abschied betrübt werden. Sei du, HErr, ihr Vater, versorge sie und pflege ihrer, als ein treuer Vormund, und nimm meine Seele auf in deine Gnadenhände, der du gesagt hast: ein Weib wird selig durch Kinderzeugen, so sie bleibt im Glauben und in der Liebe. Ich glaube, lieber HErr, stärke meinen Glauben, versiegele in mir deine Gnade und Verheißung, daß du mich, als ein Kind Gottes, auferziehen wirst in die ewige Freude und Herrlichkeit. Ach, HErr, höre, ach HErr, sei gnädig! ach HErr, merke auf! und thue es, siehe! ich bin des HErrn Magd, mir geschehe nach deinem Willen. Amen.

\section{148. Gebet, wenn die Stunde der Geburt naht}

Ach du lieber Gott! nun habe ich keine Stunde mehr vor mir, sondern muß alle Augenblicke aufwarten, wenn du mich zu arbeiten in der Kindesnoth erfordern wirst; dir ergeb ich mich ganz und gar mit Leib und Seele in deinen allmächtigen Schutz und Schirm, schaffe es mit mir nach deinem väterlichen Willen und Wohlgefallen, wie du weißt, daß es mir zu Leib und Seele nütz und gut ist. O HErr Jesu Christe! du hast ja nach der Einsetzung des heil. Abendmahls in der tröstlichen Valetpredigt, die du deinen Jüngern gethan, des gebärenden Weibes Kindesnoth und Traurigkeit gedacht, auch wiederum der Freude auf die Geburt, und hast damit die Christenheit nach der Traurigkeit trösten wollen; ich bitte dich durch deinen Todeskampf und blutigen Schweiß, durch dein Kreuz und Tod, stehe mir bei durch Hülfe und Beistand, Trost und Stärke deines Heiligen Geistes, daß ich auch nach ausgestau-
\quellseite{120}
[91]

deiner Traurigkeit und Schmerzen mit einem fröhlichen Anblick meiner Geburt und gnädiger Entbindung meiner Schmerzen möge erfreuet werden; wo aber mein Stündlein, von dieser Welt abzuscheiden, vorhanden wäre, und ich in der Geburt meinen Geist aufgeben sollte, so bitte ich dich von Herzen, du wollest dein bitter Leiden und Sterben an mir armen Sünder nicht lassen verloren sein. Laß dir meinen Mann (und unerzogene kleine Kindlein), weil du selbst dich zu einem Vater der Waisen angegeben hast, befohlen sein, sintemal ich nur, weil ich lebe, zu einer Pflegerin darüber gesetzt bin. Darum so übergebe und befehle ich dir alles in deine väterliche Fürsorge, wie du es machest, so ist es alles sehr gut. Lebe ich, so will ich dir dienen, so lange ich hier bin; sterbe ich, so nimmst du mich in meinem Beruf in das ewige Leben. Dir, ewigen Vater, Sohn und Heiligen Geist, sei allezeit Lob, Ehre und Preis, heut und in Ewigkeit. Amen.

\section{149. Gebet bei Todesgefahr in Kindnöthen}

Allmächtiger, barmherziger Gott und Vater im Himmel! ich bitte dich durch den Todeskampf und ängstigen Schweiß und Schmerzen meines einigen Heilandes und Seligmachers, deines lieben Sohnes, JEsu Christi, du wollest mir schwachem und mattem Weibe Stärke, Kraft, Geduld und Trost in dieser sehr großen Noth verleihen, und mir einen fröhlichen Anblick bescheren, denn du allein thust alle Hülfe, die auf Erden geschieht. Soll ich aber sammt meinem Kindlein bleiben, ach! so kaufe doch diese meine Leibesfrucht mit dem theuren Blute deines lieben Sohnes, und beschere mir ein seliges Ende. Laß mich fröhlich leben und selig sterben. Nimm meine Seele in deine göttlichen Hände, denn Christus ist mein Leben; ich begehre auch nach deinem heiligen Willen aufgelöst zu werden, und bei Christo JEsu, meinem Bruder und Immanuel, zu sein, aus dessen Hand mich niemand reißen wird. O HErr JEsu Christe, du Sohn des lebendigen Gottes, hilf mir, und tröste mich in dieser großen Noth, und gib mir und meiner Leibesfrucht, so deiner Hände Werk und dein Segen ist, Stärke und Kraft, daß wir nicht verderben! O Gott, Heiliger Geist, du höchster Tröster in aller Noth, tröste mich mit deinem göttlichen
\quellseite{121}
Trost, laß mich nicht verzagen, sondern stehe mir bei, du starker Gott, um deiner Wahrheit und deines Namens Ehre willen. Amen, du heilige Dreifaltigkeit! Amen.

\section{150. Ein Anderes}

Ach HErr, wie lange willst du mein so gar vergessen? wie lange verbirgst du dein Antlitz vor mir? Hast du denn, liebster Vater, vergessen gnädig zu sein, und dich verwandelt in einen Grausamen, um deinen Grimm an mir zu zeigen mit der Stärke deiner Hand? Ach daß Gott erbarm! sollte ich nicht mit Zion winseln: Der HErr hat mich verlassen, der HErr hat mein vergessen! Schau doch, du frommer Gott! und siehe, ob irgend ein Schmerz sei, wie mein Schmerz, der mich troffen hat. Denn mir ist bang, daß ich kaum Odem holen kann, weil das Kind an die Geburt kömmt, und dennoch keine Kraft zu gebären bei mir ist. Ach Jesu, du Sohn des lebendigen Gottes! erbarme dich mein und hilf mir, es gerathe mir zum Leben oder zum Tode. Und weil ich weiß, daß das Weib selig wird durch Kinderzeugen, wenn sie bleibet im Glauben, und die Stunde meines Abschieds ohnedem scheint vorhanden zu sein, so geschehe zwar, HErr! dein Wille auf Erden, wie im Himmel; nur erbarme dich meiner armen Seele, die du nicht mit vergänglichem Golde oder Silber erlöset hast, sondern mit deinem theuren Blute, und nimm meinen Geist auf. Ists nicht möglich, daß mein Schmerzenskind zur Welt soll geboren werden, sondern sein Grab in meinem Leibe haben, oder doch außer demselben ohne Leben sein soll, so befehle ich dir hiermit seine und meine Seele in deine Hände, dir leb ich, dir sterbe ich, dein bin ich todt und lebendig. Amen.

\section{151. Danksagung nach glücklicher Geburt}

O gütiger und barmherziger Gott! der du auch an uns sonderlich und gnädiglich bewiesen hast, daß du allein der einige Gott bist, der aus Nichts Etwas und aus dem Tode lebendig und wohlgeschickt macht, der du uns eine lebendige, rechtschaffene Frucht nach deinem Bildniß ohne alle unser Verdienst gnädiglich
\quellseite{122}
gegeben hast, darum wir dir nicht weniger pflichtig sind, dankbar zu sein, denn vorzeiten verpflichtet sind gewesen Elkana und Hanna, Zacharias und Elisabeth, Joachim und Anna, und andere dergleichen, denen du auch wunderbärlich und gnädiglich zu seliger und fröhlicher Geburt verholfen hast. Nun aber sagen wir dir demüthiglich Lob und Dank solcher unaussprechlichen Gnade und Wohlthat halben mit höchster Bitte, du wollest Gnade verleihen, daß dieß unser Kindlein sammt andern, die du uns nach deinen Willen hinfüro würdest bescheren, in deiner göttlichen Ehre, Furcht, Tugend, Frömmigkeit und Dienstbarkeit, in ziemlicher Gesundheit aufwachse, und bis an sein Ende zunehme, daß du daran ein Wohlgefallen und wir eine selige Freude ohne alle Schaude und Laster daran haben mögen und dich, unsern HErrn und Schöpfer, hier und dort zu ewigen Zeiten loben, ehren und preisen. Amen.

\section[152. Gebet eines Ehemanns nach der Entbindung …]{152. Gebet eines Ehemanns nach der Entbindung seiner Frau mit einem todten Kinde}

Ach du mein lieber frommer Gott! weil es denn also deiner Weisheit gefallen, daß du unsre Leibesfrucht in der Geburt von diesem Leben abgefordert und zu dir genommen hast, wohlan, so müssen wir es uns auch also gefallen lassen. Und wollen wir nicht zweifeln, du habest es auf unser herzliches Gebet um deines lieben Sohnes willen zu allen Gnaden an- und aufgenommen, und es sei nun gewißlich ein liebes Kindlein der ewigen Seligkeit, da es auch freilich am besten versorget und verwahret ist vor allem Uebel, dessen es viel hätte ausstehen müssen in dieser argen Welt. Ich danke dir aber daneben von Herzen, daß du in diesen gegenwärtigen schweren Kindesnöthen dennoch mein liebes Weib beim Leben erhalten und gleich mitten aus dem Tode gerissen und mir (auch meinen unerzogenen Kindern) zum besten gleich wieder aufs neue geschenket hast. Ach hilf, daß sie wieder zu ihren Kräften, die sie fast in so großer Noth und schwerer Arbeit zugesetzt und verloren hat, kommen möge. So will ich dich loben und preisen und deine Barmherzigkeit rühmen, hier zeitlich und dort ewiglich, durch Ehristum, unsern HErrn. Amen.

\quellseite{123}

\section[153. Gebet einer Mutter, die ein …]{153. Gebet einer Mutter, die ein gebrechliches Kind geboren hat}

O gnädiger, barmherziger Gott, ich sage dir von Herzen Dank, daß du mich aus Leibs- und Lebensgefahr durch deine allmächtige Hand erlöset, und meiner weiblichen Bürde mich mit Gnaden entbunden. Weil du mir aber eine gebrechliche und unförmliche Creatur bescheeret hast, und zweifelsohne damit mir meine geistliche Gebrechen, als die Unform der verderbten Natur und Abscheulichkeit meiner großen Sünde zu verstehen gibst; so hilf auch, du getreuer Gott, daß ich solche meine Gebrechen erkenne, und mich hinfüro für Sünden hüte. Du wollest auch mein gebrechlich Kindlein durch das heilige Bad der Taufe heilen und gesund machen, auch seine leibliche Gebrechen mit Gottesfurcht, auch andern schönen Gaben und Tugenden ersetzen, und es zum Himmelreich tüchtig und geschickt machen; mir auch Gnade verleihen, daß ich selbiges seiner Gebrechlichkeit halben nicht hasse, sondern als dein Geschöpf liebe, nähre und zu allem Guten, wie einer christlichen Mutter gebühret, auferziehe; um deines lieben Sohnes Jesu Christi willen. Amen.

\section[154. Gebet einer Mutter, nachdem ihr …]{154. Gebet einer Mutter, nachdem ihr neugeborenes Kind ohne Taufe gestorben}

Allmächtiger Gott, barmherziger Vater! du siehest die große Traurigkeit deiner Magd und weißt, daß ich in diesem meinem betrübten Zustande nicht so wohl über meines Kindes Tod, als vielmehr und am allermeisten um dessen Seligkeit bekümmert bin. Ich erkenne und verehre deine Macht, die du über meine Leibesfrucht als dein Geschöpf hast; daß, wie du ihr nach deinem Wohlgefallen Leben und Odem gegeben, dir auch freigestanden habe, beides ihr wieder zu nehmen. Dieses aber betrübet mich herzlich und dringet wie ein Schwert durch meine Seele, daß dieses arme Kind verschieden ist, ehe es durch die heilige Taufe dem Bund deiner Gnaden hat einverleibt werden können. Ach gnädiger Gott! nimm doch diese Traurigkeit aus meinem Herzen und laß den Trost aus deinem
\quellseite{124}
Worte mich erquicken: du habest mein Gebet und Flehen, welches ich für diese meine Leibesfrucht gethan, gnädig erhöret, es sei mein Kind, kraft der Verheißung, die du dem Abraham und einem jeden gläubigen Christen gegeben hast, in deinen Gnadenbund demnach eingeschlossen worden, weil du gesprochen: ich bin dein Gott und deines Samens nach dir; ja du könnest im Nothfalle auch ohne die Taufe die Menschen selig machen.

Jesu, liebster Heiland, der du ein herzlicher Kinder-Freund bist und sprichst: Lasset die Kindlein zu mir kommen und wehret ihnen nicht, denn solcher ist das Reich Gottes. Wie herzlich gern hätte ich mein liebes Kind in der heiligen Taufe zu dir kommen lassen und in deinen Gnadenschooß legen wollen, wenn solches nicht durch den so frühzeitigen Tod mir wäre verwehret worden. Ich weiß aber aus deinem göttlichen Munde, es sei nicht der Wille des Vaters im Himmel, daß jemand von diesen Kleinen verloren werde; und glaube festiglich, es habe Gott dieses mein Kind nicht zum ewigen Verderben erschaffen, du, o eingeborner Sohn Gottes, habest es mit deinem theuren Blut von Sünden abgewaschen und mit deinem Kreuzestod von der ewigen Todesgewalt erlöset, wollest dich also seiner Seele herzlich annehmen, damit sie nicht verloren gehe. Hast du dem bekehrten Schächer am Kreuz, der niemals des Sacraments der heiligen Taufe theilhaftig worden ist, das Paradies versprochen, warum hättest du nicht auch dieses Kind meines Leibes ohne Taufe in dein ewiges Reich versetzen können. Aber mein Heiland, stärke doch in mir den schwachen Glauben und gib, daß ich mich hiermit allezeit trösten und aufrichten kann.

Herr Gott Heiliger Geist, du Gott der Geduld und des Trostes! vertreibe mit deinem Gnaden- und Trostlicht die finstern Wolken des Betrübnisses aus meiner Seele. Laß mich wohl betrachten die herrlichen Wohlthaten, die Gott durch den Tod an meinem Kinde gethan, daß er dasselbe so geschwind von der Last der Sünden und aller Mühseligkeit erlöset, aus dem Tode ins Leben, aus dem Elend in die Glückseligkeit, aus der Angst in die Ruhe geführt, die Seele zu der Gesellschaft der heiligen Engel aufgenommen und alsobald im Anfang ihres Laufes gekrönet hat. Hilf, daß ich in Betrachtung dessen bald meines Leibes vergesse, und nicht
\quellseite{125}
trauere, wie die Heiden, die keine Hoffnung haben, sondern mich mit einem gottseligen David gewiß in meinem Herzen versichert halte, daß, ob wohl mein Kind nicht wieder in diese Welt zu mir kommt, so werde ich doch einstens zu ihm kommen und in himmlischer Ehre und Freude es wieder sehen; dafür aber dich sammt Gott dem Vater und meinem Heiland Christo rühmen und preisen in alle Ewigkeit. Amen.

\section[155. Gebet einer Mutter, die den ersten …]{155. Gebet einer Mutter, die den ersten Kirchgang hält}

Gott Vater, Sohn und Heiliger Geist, einiger ewiger Gott, dir sei Lob und Dank gesagt, daß du mir nicht allein in meinem Kreissen zu der Stunde, da ich habe gebären sollen, mit deiner Kraft und Stärke beigestanden und eine gesunde Leibesfrucht gnädiglich bescheeret hast, sondern auch dieselbe zur heiligen Taufe kommen und deiner christlichen Kirche als ein wahres Glied hast einverleiben lassen; auch noch dazu mich, nebst meinem lieben Kinde in meinem Wochenbett ganz väterlich erhalten, versorget, erquicket und gestärket hast. Nun komme und erscheine ich vor deinem göttlichen Angesicht und opfre dir mein liebes Kind, das du mir gegeben hast — wie die liebe Hanna that — mit herzlicher Bitte, du wollest sein lieber Gott und Vater sein und bleiben. Und wie du es in der heiligen Taufe zu deinem Kinde hast auf- und angenommen, also wollest du es nun auch einen Erben deiner himmlischen Güter und der ewigen Seligkeit in Christo JESU, deinem eingebornen Sohne, sein und bleiben lassen; wollest es auch die Zeit seines Lebens vor allem Uebel Leibes und der Seele bewahren, und mir deine göttliche Hülfe und Segen verleihen, daß ich dir mein liebes Kind in allen christlichen Tugenden zu deiner Ehre, zu meiner und meines lieben Eheherrn Freude, zu des Nächsten Nutz und zu seiner zeitlichen und ewigen Wohlfahrt, als eine treue Mutter, fleißig auferziehen möge. Das wollest du thun um deines lieben Sohnes JESU Christi willen, in welchem dir alle getauften Christenkindlein herzlich gefallen und angenehm sind; dem sei mit dir und dem Heiligen Geiste Lob, Ehr und Preis von nun an bis in Ewigkeit. Amen.

\quellseite{126}

\section{156. Gebet eines Mannes, der eine böse Frau hat}

Herr Gott, himmlischer Vater! ich weiß und bin gewiß, daß dieser Stand, darein du mich gesetzet, dir gefällig sei, und daß dir alles mein Thun und Leiden, in diesem meinem Stand und Beruf, ein angenehmes Opfer sei, von wegen deines lieben Sohnes, meines Erlösers und Seligmachers Jesu Christi, den ich ehre, liebe und bekenne. Nun siehest du, in was für Jammer und Elend ich stecke, und daß ich in solcher Gefahr stehe, daraus keine menschliche Hülfe mich erretten kann. Ich armer und betrübter Mann habe leider eine unglückliche Ehe mit meinem Eheweibe, wegen ihrer Halsstarrigkeit, Ungehorsam, Bosheit und Unfleiß in der Haushaltung, und fället mir über die Maßen schwer, in solchem Wesen mit meinem Ehegenossen in die Länge zu leben und zu hausen, denn ich habe doch keine fröhliche, ruhige und friedliche Stunde bei ihr, und muß besorgen, es werde dadurch meine ganze Haushaltung zu Grunde und zu Boden gehen. Derowegen klage ich dir solche Noth mit betrübtem Herzen, und rufe dich flehentlich an, du wollest hierinnen ein nothwendig, väterlich und gnädig Einsehen haben, solches bei Zeiten wenden, und meinem Eheweibe Gnade und Geist verleihen, daß sie mit rechter Vernunft und Verstand mir christlich beiwohne, sich freundlich, geduldig und sanftmüthig jederzeit gegen mir erzeige, mich ehrlich liebe und werth halte, und nicht gar zu widerspenstig, strenge und ernsthaft gegen mich sei, damit dem argen bösen Feinde nicht Ursach und Anlaß gegeben werde, noch mehr Uneinigkeit und Unfrieden zwischen uns anzurichten, dadurch wir in unserm Gebet und also in allem Segen und gedeihlicher Wohlfahrt zu Leib und Seele verhindert werden. Ach Herr! lehre uns thun nach deinem Wohlgefallen, und dein guter Geist führe uns auf rechter Bahn, daß wir in rechter, unverfälschter, ehelicher Liebe einherwandeln, und sonderlich dem löblichen Hausfrieden nachjagen, eines dem andern etwas zu gute halte und nachgebe, damit wir in gutem Frieden und löblichem Wohlstande, als christlichen Eheleuten gebühret, unser Leben vollstrecken mögen, bis wir unsern Ehestand allhier aufgeben, und dort mit dir in deinem himmlischen Reich ewiglich leben, auch deine herrlichen Wundertaten rühmen und preisen mögen. Amen.

Gebetbuch.

\quellseite{127}

\section{157. Gebet einer Frau, die einen bösen Mann hat}

Ach barmherziger Gott und Vater, du allwissender HErr, du kennest mein Elend besser, als ich dirs sagen und klagen kann, denn mein Ehegatte, den du mir einmal zugestellet, ist leider nicht also, wie es billig sein sollte, und habe ich mit ihm, wie dir wohl bekannt, mein schweres und fast tägliches Hauskreuz. Daß ich also nicht weiß, was ich in solchem elenden und betrübten Zustande soll anfangen, dieweil der leidige Satan mir meinen Ehestand solcher Gestalt mit meinem unartigen Ehegatten sehr sauer machet. Von einander können wir nicht kommen, weil wir einmal von dir zusammengefüget sind, daß uns nichts, als nur allein der Tod, scheiden kann. Nun, ich wende mich zu dir, HErr mein Gott, ohne dessen Rath und Willen und ohne dessen Zulassung nichts geschehen kann. Ich muß bedenken, daß du mich mit der bösen Unart meines Ehegatten, als mit einem sonderlichen Kreuz willst züchtigen, so ich anderweit um dich wohl habe verschuldet. Aber gib mir doch Geduld dazu und hilf mir solche Last ertragen mit deinen liebreichen Vaterhänden. Laß uns auch nicht immer in solcher Widerwärtigkeit mit einander leben, sondern regiere das Herz meines Ehegatten, daß er anders werde, und sich also bezeige, wie einem frommen, christlichen Ehegatten zusteht und gebühret. Steure dem leidigen Eheteufel und andern bösen Mäulern, dadurch mein Ehegatte leider oft verführet wird; und pflanze rechtschaffene Gottesfurcht in sein Herz, auch ehrliche, aufrichtige Liebe gegen mich, damit wir hinfüro in Friede und Einigkeit bei einander wohnen, und unser Hauswesen und Nahrung also mit einander fortsetzen mögen, daß Gott und Menschen an uns ein Wohlgefallen haben. Segne auch uns von oben herab mit deinem himmlischen Segen, daß wir als gehorsame Kinder nach deinen Geboten wandeln und als die Gesegneten des HErrn sein und bleiben mögen, hier zeitlich und dort einmal auch ewiglich. Amen.

\section[158. Gebet einer Frau, die einen bösen und …]{158. Gebet einer Frau, die einen bösen und aufrührerlichen Mann hat}

HErr Gott, himmlischer Vater! Ich weiß und bin gewiß, daß dieser Stand, darein du mich gesetzet, dir gefällig sei, und daß
\quellseite{128}
dir alles mein Thun und Leiden, in diesem meinem Stand und Beruf ein angenehmes Opfer sei, von wegen deines lieben Sohnes, meines Erlösers und Seligmachers Jesu Christi, den ich ehre, liebe und bekenne. Nun siehest du, in was für Jammer und Elend ich stecke und daß ich in solcher Gefahr stehe, daraus keine menschliche Hülfe mich retten kann. Ich armes und betrübtes Weib habe leider eine unglückliche Ehe mit meinem Ehemann wegen seiner Halsstarrigkeit, Ungehorsam, Bosheit und Unfleiß in der Haushaltung, wegen seines Schwelgens, Spielens, Wüthens und Hinlässigkeit in seinem Beruf und fällt mir über die Maßen schwer, in solchem Wesen mit meinem Ehegenossen in die Länge zu leben und zu hausen, denn ich habe doch keine fröhliche, ruhige und friedliche Stunde bei ihm, und muß besorgen, es werde dadurch meine ganze Haushaltung zu Grunde und zu Boden gehen. Derowegen klage ich dir solche Noth mit betrübten Herzen, und rufe dich flehentlich an, du wollest hierinnen ein nothwendig, väterlich und gnädig Einsehen haben, solches bei Zeiten wenden, und meinem Ehemann Gnade und Geist verleihen, daß er mit rechter Vernunft und Verstand mir christlich beiwohne, sich freundlich, geduldig und sanftmüthig jederzeit gegen mich erzeige, mich ehrlich liebe und werth halte, und nicht gar zu widerspenstig und streng gegen mich sei, damit dem argen bösen Feinde nicht Ursache und Anlaß gegeben werde, noch mehr Uneinigkeit und Unfrieden zwischen uns anzurichten, dadurch wir in unserm Gebet und also in allem Segen und gedeihlicher Wohlfahrt zu Leib und Seele verhindert werden. Ach HErr! Lehre uns thun nach deinem Wohlgefallen, und dein guter Geist führe uns auf rechter Bahn, daß wir in rechter, unverfälschter, ehelicher Liebe einherwandeln, und sonderlich dem löblichen Hausfrieden nachjagen, eines dem andern was zu gute halte und nachgebe, damit wir in gutem Frieden und löblichem Wohlstande, als christlichen Eheleuten gebührt, unser Leben vollstrecken mögen, bis wir unsern Ehestand allhier aufgeben, und dort mit dir in deinem himmlischen Reiche ewiglich leben, auch deine herrlichen Wunderthaten rühmen und preisen mögen. Amen.

\quellseite{129}

\section[159. Gebet, wenn ein Ehegemahl das andere …]{159. Gebet, wenn ein Ehegemahl das andere verlassen hat}

Allmächtiger Gott! ich klage dir meinen großen Schmerz, welchen ich mit Thränen nicht genugsam beweinen kann, wie mein Ehegemahl so treulos an mir worden und mich, ohne einige erhebliche Ursache, so schändlich verlassen; niemand weiß mir zu helfen oder zu rathen, als du, mein getreuer Gott. So gib mir nun deine Gnade, daß ich dies schwere Hauskreuz als deine väterliche Züchtigung erkenne und mich deinem väterlichen Willen in kindlicher Demuth unterwerfe. Du Herzenkündiger und Herzenzreger, stehest sein Herz, ißt noch zur Besserung zu lenken, ach! so regiere und leite es, verknüpfe es wieder mit rechter, ehelicher Liebe und Treue mit meinem Herzen, daß wir noch, als rechten Eheleuten gebühret, dir zu Ehren bei einander wohnen mögen. Ist aber sein Herz nicht zu gewinnen, so nimm dich meiner an und versorge mich, wie du armen Wittwen (Wittwern) in deinem Wort versprochen hast, um deines hochwerthen Namens Ehre willen. Amen.

\section{160. Gebet eines Wittwers}

HErr, allmächtiger Gott, du Vater der Barmherzigkeit und Gott alles Trostes, dessen Hand verwundet und heilet. Ich klage dir mit betrübtem Herzen mein unaussprechliches Leid und Jammer, indem du mir meine Augenlust, meine getreue Gehilfin, die Zierde meines Hauses, darauf sich mein Herz in allen meine zeitliche Wohlfahrt betreffenden Dingen verlassen konnte, durch den zeitlichen Tod von der Seite gerissen. HErr, siehe an mein Weinen, höre mein Seufzen und Wehklagen, siehe an meinen Jammer und Elend, und vergib mir alle meine Sünde. Wende dich zu mir, und sei mir gnädig, denn ich bin einsam und elend, die Angst meines Herzens ist groß, führe mich aus meinen Nöthen. Heile mich, HErr, erquicke meine Seele, laß mich nicht in meinem Kummer verschmachten, laß mich nicht, und thue nicht von mir die Hand ab, Gott mein Heil. Ich bitte dich demüthiglich, laß dein heiliges Wort meinen Trost sein, damit ich nicht vergehe in meinem Elende. Erhalte ferner, was du mir und den Meinen bisher gegeben, (erbarme dich der armen und unverzogenen Waisen), beschütze mein Haus
\quellseite{130}
und Hof, Hab, Ehr und Gut, gib mir getreu Gesinde, wende ab alle gefährliche Krankheit, Diebe, Räuber, Wasser- und Feuersnoth, sammt allem Unfall Leibes und der Seelen. Auf dich, HErr, traue ich, laß mich nimmermehr zu Schanden werden, sei du mein Trost, mein Schutz, meine Stärke, meine Hilfe, meine Freude, daß ich dich für deine Treue lobe und preise in Ewigkeit. Amen.

\section{161. Gebet einer Wittwe}

Ach du getreuer, hülfreicher, freundlicher und gnädiger Gott, der du dich selbst einen Vater der Waisen und Richter der Wittwen genannt hast; einen Vater der Waisen, dieweil sie keinen Vater auf Erden haben, der sie so herzlich liebet als seine leiblichen Kinder, der so emsig und fleißig für sie sorget, als ein leiblicher Vater, der ihnen ihre Nothdurft schaffet als ihr Vater, der sie erziehet, sie tröstet, sich über sie in Krankheit und Elend erbarmet, wie sich ein Vater über seine Kinder erbarmet. Dieses freundliche und holdselige Vateramt hast du dich erboten auf dich zu nehmen, darum wirst du es auch an mir, an meinen und allen armen Waisen ganz väterlich und treulich verrichten, meine und alle armen Waisen väterlich lieben, versorgen, ernähren, erziehen, sie trösten und dich über sie als ein herzlicher Vater erbarmen. Einen Richter der Wittwen hast du dich genannt, weil dieselben viel von der ungerechten Welt leiden müssen, und keinen Schutz noch Trost haben, und als Elende, Trostlose, Verlassene alle Wetter müssen lassen über sich ergehen. Darum bitte ich dich demüthiglich, weil du gesagt hast: Der HErr, der dich gemacht hat, ist dein Mann, der aller Welt Gott genannt wird; du wollest auch mich, als eine Trostlose und Verlassene, in deinen allmächtigen Schutz nehmen, und nicht zugeben, daß mir Gewalt und Leid von der bösen Welt geschehe; sondern mich und mein betrübtes Haus verwahren und umschränken mit deiner heiligen Verheißung: Du sollst die Wittwen nicht beleidigen. Sie werden zu mir schreien und ich werde ihr Gebet erhören und dein Weib zur Wittwe und deine Kinder zu Waisen machen. Laß dein Wort eine starke Mauer um mich und meine armen Waisen her sein. Erwecke fromme Herzen,
\quellseite{131}
die da bedenken, daß ein reiner und unbefleckter Gottesdienst sei, Witwen und Waisen in ihrem Trübsal besuchen, das ist, mit Rath, Hülfe und Trost besuchen. Behüte mich, lieber Vater, vor den Lügenmäulern und falschen Zungen. Gib mir deine Gnade, daß ich in dieser meiner Einsamkeit meine Hoffnung auf dich allein setze, an dir alle meinen Trost habe, dich mit wahrem Glauben und starker Zuversicht ergreife. Du wollest auch mein Mehlsäßlein segnen, und nicht lassen leer werden, wie der Wittwe zu Sarepta, die den Propheten Eliam aufnahm; und mein Oelkrüglein nicht lassen zerrinnen und vertrocknen, wie der armen Wittwe geschah, welche nach des Propheten Elisä Segen aus ihrem Oelkrüglein alle ihre Gefäße füllete, und allein durch Gottes Wort und Segen mit ihren Kindern ernähret und erhalten wurde; auf daß ich erkenne, daß du mein Vater, Erhalter, Trost und Schutz bist, durch JEsum Christum, deinen lieben Sohn, unsern HErrn. Amen.

\section{162. Gebet einer hochbetrübten Wittwe}

Barmherziger Gott! du hast mich nach deinem heiligen Willen sehr betrübet, indem du mich zur Wittwen gemacht, daß ich nun von der Welt verlassen bin, und gleichsam ein zweischneidiges Schwert durch meine Seele dringet. Ach nun ist meine Sonne bei hohem Tage untergegangen, daß beide, mein Ruhm und Herzensfreude ein Ende hat; ja, die Krone meines Hauptes ist abgefallen, mein Neigen ist in Wehe und Klage verkehret. Dieweil du dich aber, lieber HErr, einen Richter der Wittwen nennest, und die Wittwe, wenn sie klaget, nicht verachtest, auch ihre Thränen, so die Backen herabfließen, über sich zu dir schreien, so habe ich meine Zuflucht allein zu dir, HErr! und bitte dich herzlich, du wollest deine Hülfe zu mir neigen, und meine Noth dir zu Herzen gehen lassen. Wende dich zu mir, und sei mir gnädig, denn ich bin einsam und elend, die Angst meines Herzens ist groß, führe mich aus meinen Nöthen. Erzeige dich gegen mich wie der rechte Helfer, und gegen meine armen Waiselein, wie der rechte Vater. Denn auf Menschenhülfe ist sich ganz nicht zu verlassen, insonderheit müssen Wittwen und Waisen anderer Leute Fußtrit scheuen, und einen Weg über sich machen lassen. Darum ergebe ich mich dir
\quellseite{132}
ganz und gar, laß mich deinem väterlichen Schutz befohlen sein, und laß meine Funken nicht ganz und gar verlöschen. Siehe, hier bin ich, lieber himmlischer Vater! und die Kinder, die du mir gegeben hast. Hilf du, HErr, schütze du, versorge du, segne du mich, ich kann mir selbst weder rathen noch helfen. Das weiß ich aber fürwahr, wer Gott dient, der wird nach der Anfechtung getröstet, und aus der Trübsal erlöset, und nach der Züchtigung findet er Gnade; denn nach dem Ungewitter lässest du die Sonne wieder scheinen, und nach dem Heulen und Weinen überschüttest du uns mit Freuden. Beschere mir ein Stück Brots, und weise mir getreue Leute, die sich mit Rath und That meiner annehmen, gleichwie du in der großen Theuerung die Wittwe zu Sarepta versorget und ihr den Propheten Eliam zugeschicket hast. Insbesondere aber regiere mich mit deinem Heiligen Geist, daß ich in aller Noth meine Hoffnung auf dich, lebendigen Gott, setze, und mit dem Gebet Tag und Nacht anhalte. Desgleichen, daß ich mich eines gottseligen, keuschen, züchtigen, eingezogenen Lebens und Wandels befleißige, auf daß ich nicht mit Geberden, Worten und Werken jemand Aergerniß gebe, und Sünde auf mich lade, und also lebendig todt erfunden werde, sondern mich also verhalte, daß niemand von mir übel rede. Behüte mich vor den Lügenmäulern und falschen Zungen. Sei mir gnädig, HErr! sei mir gnädig, denn ich bin einsam und elend. Die Angst meines Herzens ist groß. Erbarme dich meiner um JEsu Christi, deines lieben Sohnes willen. Amen.

\section{163. Gebet einer Hebamme}

HErr, allmächtiger Gott, du wunderbarer Schöpfer der Menschen! ich danke dir, daß du mich auch zu einem Werkzeuge deiner göttlichen Kraft und Macht verordnet und mir den Beruf gegeben hast, daß ich den mit Leibesfrucht gesegneten Frauen in ihrer Entbindung beistehen und über der Geburt ihnen förderlich und dienlich sein solle. Daneben aber bitte ich dich auch, du wollest deine Gnade dazu verleihen, daß meine Berufsarbeit jederzeit glücklich und wohl möge verrichtet, und mit Freuden möge vollendet werden. Lieber Gott! laß mich ja nichts versehen, das in solchen
\quellseite{133}
Fällen von Nöthen ist, damit Niemand versäumet, und Mutter oder Kind am Leben gefährdet werde. Der du selber die Kinder im Mutterleibe bildest, und bis an die Zeit ihrer Geburt darinnen erhältst, ach, lege auch deine Hand selbst mit an, wenn sie jetzt sollen geboren und an das Tageslicht gebracht werden, daß unsere Arbeit desto besser wirken möge. Laß mich Alles mit gutem Bedacht und Verstand thun, damit ich weder mit Worten, noch mit Werken möge beschwerlich werden. Laß mich nicht, ehe die Zeit kommt, die Kreisserinnen veranlassen zu arbeiten, daß sie unnöthig ermüdet werden, und dann zu rechter Zeit keine Kraft mehr haben. Auch laß mich nicht mit harten, sondern sanften, gelinden und freundlichen Worten gegen Jedermann verfahren, daß man sich nicht über mich zu beschweren habe. Behüte mich vor starkem Trinken, daß ich nicht, damit überfüllet, unverantwortliche Fehler begehe. Heimliche Zauberei, Aberglauben und Mißbrauch deines göttlichen Namens laß ferne von mir sein. Wenn es auch je zuweilen hart hergehet, laß mich nicht ungeduldig darüber werden, oder etwa gar die Leute verlassen, um zu andern zu laufen, wo mehr Verdienst zu hoffen ist; sondern Armen und Reichen mit gleicher Willfährigkeit dienen, weil es nicht um meinen Gewinn, sondern um Erhaltung der Mutter und des Kindes zu thun ist. Und wenn dann das Kind zur Welt geboren ist, so laß mich nichts an ihm versäumen, was zu seiner geistlichen Wiedergeburt gehöret, und mich ja der Nothlaufe nicht muthwillig unterfangen, damit deine Ordnung nicht verrücket werde. Allerliebster Herr und Gott! laß mich in allem meinen Thun christlich, gottesfürchtig, züchtig, voll heiligen Glaubens, im Gebet und rechtschaffener Andacht erfunden werden, damit das Werk, das ich vorhabe, desto glücklicher und gesegneter sein möge. Amen, Amen.

\section{164. Gebet junger Leute}

O Vater aller Gnaden und Barmherzigkeit! der du ein heiliger Gott bist, du willst, daß auch wir heilig, ohne Lustseuche und Begierde unsers sündlichen Fleisches leben sollen. Ich klage dir meine Schwachheit und angeborne sündliche Gebrechen. Ich habe zwar das Wollen, aber zu vollbringen das Gute finde ich
\quellseite{134}
nicht. Da stürmen als drei mächtige Feinde auf mich, Fleischeslust, Augenlust, und hoffärtiges Leben, und wollen mich von deiner Gnade abziehen. Darum bitte ich, mein Vater, schaffe und bewahre in mir ein rein Herz. Gieb, daß ich alle Oerter und Stellen, so Gelegenheit zur Unreinigkeit geben, von Herzen meide, und mich fürchte, aus den Gliedern meines Leibes, als Christi Jesu, meines Heilandes Gliedern, Hurenlieder zu machen. Laß mich vielmehr in feuriger Andacht und sehnlichem Himmelsverlangen allezeit brünstig sein, daß ich als ein Tempel Gottes des Heiligen Geistes lebe, und zu lieben nicht aufhöre Jesum Christum, meine Liebe, die da gekreuziget ist, und sich selbst für mich in den Tod gegeben hat. Laß mich lieben die himmlische Weisheit, damit ich ja nimmermehr auf verführerische Umwege gerathe. O du gerechter und heiliger Gott! du prüfest beide Herzen und Nieren, gib, daß ich in dieser Welt dir also anhange, daß ich auch in Ewigkeit von dir ungeschieden bleibe. Amen.

\section{165. Gebet eines Jünglings}

O heiliger und keuscher Gott! der du ein Wohlgefallen hast am heiligen Wandel und gottseeligen Wesen, und also auch an keuschen und züchtigen Herzen, und willst, daß junge Gesellen und Jungfrauen heilig und keusch sein sollen beide am Leibe und auch am Geiste, daher auch dein Knecht Paulus schreibet: keusch die Lüste der Jugend, jage aber nach der Gerechtigkeit, dem Glauben, der Liebe, dem Frieden, mit allen, die den HErrn anrufen mit reinem Herzen, und ermahne die jungen Leute, als Brüder und Schwestern, mit aller Keuschheit. Diesen nachzukommen, so bitte ich dich inniglich, du wollest in mir schaffen ein reines Herz, daß ich keusch und züchtig lebe, in Gedanken, Worten und Werken, und dir diene mit reiner Seele und unbeflecktem Leibe, wie vorzeiten der züchtige Jüngling Joseph in Egypten dir mit reiner Seele und unbeflecktem Leibe gedient, da er sich weigerte, bei seines Herrn Weibe zu schlafen.

Behüte mich vor unzüchtigem Gesichte, und wende von mir ab alle böse Lust. Laß mich nicht in Schlemmen und Unkeuschheit gerathen, und behüte mich vor unverschämtem Herzen. Behüte mich
\quellseite{135}
vor aller Leichtfertigkeit in Worten und Gebärden, vor faulem Geschwätz, schandbaren Worten, Narrentheidung und Scherz, welche uns Christen nicht ziemen; vor Müßiggang und Faulheit, als Stricken und Netzen des Teufels; vor Schwelgerei und Trunkenheit, daraus ein unordentlich Leben folget; vor Hurerei und aller unreinen Vermischung, so gemeiniglich aus Müßiggang, als einem Rissen des Teufels, und aus Trunkenheit, als einer Mutter alles Bösen, entspringet. Lösche aus in mir die Flammen fleischlicher Begierden, und tödte in meinen Gliedern Hurerei, Unreinigkeit, schändliche Brunft und böse Lust, auf daß ich mich nicht ergebe der Wollust des Fleisches, Schande und Laster zu treiben mit unkeuschen Leuten, und nicht aus Christi Gliedern Hurenglieder mache. Bewahre mich auch für böser Gesellschaft, und für solchen Oertern, welche Ursache und Anreizung zu Sünden und Schauden geben; und wenn mich die bösen Buben locken, so gib, daß ich ihnen nicht folge.

Weil auch ein fromm und vernünftig Ehegemahl von dir, frommer Gott, herkommt, so soll es billig bei und von dir erbeten werden; wenn’s dann deinem heiligen Willen gefällig, daß ich auch dermaleinst in den Ehestand treten soll, so beschere mir doch ein gottseeliges, frommes und vernünftiges Ehegemahl, mit welchem ich die Zeit meines Lebens in Gottseeligkeit und Ehrbarkeit, in Friede, Liebe und Einigkeit zubringen möge, und gib, daß ich denselben Stand in deinem Namen und in deiner Furcht, auch mit Rath und Bewilligung meiner Eltern und Freunde anfange, nicht aus Vorwitz und um Unzucht willen, sondern um deiner Ordnung willen, und aus Begierde der Frucht, dadurch dein heiliger Name ewiglich gepreiset und gelobet werde. Gefällt dir’s aber, daß ich im ledigen Stande, ohne Ehe bleiben soll, wie dein Diener Paulus, und andere mehr ohne Ehe geblieben sind, so hilf mir, daß ich auch Pauli Worten folge, da er schreibet: Wer ledig ist, der sorget, was dem Herrn angehöret, wie er dem Herrn gefalle. In diesem allen wollest du, himmlischer Vater, schaffen, was dir gefällig und mir nützlich ist, der du aller Menschen Herzen in deiner Hand hast, und neigest sie, wohin du willst, hochgelobet in Ewigkeit. Amen.

\quellseite{136}

\section{166. Gebet um christlichen Hausstand}

Du allmächtiger, ewiger und barmherziger Gott, dieweil du mit Worten und Werken genugsam bewiesen hast, daß du als ein getreuer Vater für uns sorgest, und hast uns gnädiglich zu Kindern angenommen und einem jeden seiner Beruf gegeben, darinnen er dir und dem Nächsten dienen soll, so bitten wir dich von Herzen, lieber Vater, gib Gnade, daß wir unsers Berufs fleißig wahrnehmen, und im Gehorsam als deine treuen Kinder allzeit erfunden werden. Fördere die Werke unsers Haushaltens also, daß wir ja unsere Herzen nicht an die Güter dieser Welt hangen, noch jemand dadurch Ärgerniß geben, sondern laß uns alle zeitlichen Güter und Gaben, so wir durch deinen Segen empfangen haben, in stillem Wesen mit täglicher Danksagung genießen und gebrauchen; wollest auch von uns abwenden allen Müßiggang, übrige Sorge der Nahrung, und alles, was dir mißfällt, und bei uns und in uns fördern alles, was dir wohlgefällt, auf daß wir in all unserm Thun deinem Befehle nachgehen, und alle Sorge und Anliegen aus rechtem Glauben auf dich werfen. Denn du weißt alles, was uns mangelt und noth ist; das wollest du uns gnädiglich verleihen. Amen.

\section{167. Gebet eines Hausvaters, der ein Farmer ist}

Herr, allmächtiger Gott! der du einem jeglichen Menschen gibst, wie dich gut deucht, du hast Adam, den ersten Menschen, vor dem Fall zum Erben, Herrn und Hausvater der Welt gemacht, und in den Garten des Paradieses gesetzt, daß er den bauen und bewahren soll, auch nach dem Fall auferleget, im Schweiß seines Angesichtes und seiner Hände Arbeit sein Brod zu essen; durch deinen Segen habe ich auch eine Haushaltung und Nahrung, Acker und Feldbau, und nähre mich mit meiner Hände Arbeit. Lieber Vater! der du uns gesegnet hast mit allerlei geistlichem Segen und himmlischen Gütern durch Christum, gereicht es zu deiner Ehre und meiner Seelen Seligkeit, so segne auch meine zeitlichen Güter, daß meine Nahrung zunehme und sich ausbreite, gebauet dem Segen, daß er bei mir sei in allem, das ich vornehme. Gib mir von dem Thau des Himmels und von der Fettigkeit der Erde und
\quellseite{137}
Laß mich die Ströme und Wasserbäche, so mit Honig und Butter fließen, sehen. Der Heiden Götter können nicht Güter geben, du aber füllest die Erde mit deinen Gütern, und segnest uns allenthalben, daß wir des Landes Gut genießen; dein Segen macht reich ohne Mühe, und gibst es den Deinen im Schlafe. Thue uns deinen Schatz auf, und gib wahrhaftige Güter, die da bleiben und immerdar gedeihen. Laß uns ordentlich haushalten, vorsichtig und weislich handeln, und unsere Sachen so verrichten, daß wir niemand Schaden thun, und eine Last seien. Laß uns nicht Wein und Oel lieben, in Wollust und Müßiggang leben, unsere Nahrung gering achten, deine Güter verprassen, und übel umbringen, daß uns nicht Noth und Armuth übereile. Lieber Vater! der du der rechte Hausvater im Himmel und auf Erden bist, laß mich dein Wort und Befehl an alle Hausväter zu Herzen nehmen und meine Sachen darnach richten, daß ich meinen Kindern und Gesinde und allen, die mir angehören, deine Gebote schärfe und zu Haus und Felde, wenn ich mich niederlege und aufstehe, mit ihnen davon rede, daß sie, HErr, deine Wege halten und thun, was recht und gut ist, auf daß ihnen alles das Gute widerfahre, das du uns geredet hast, daß ich den Herrn ehre, von meinem Einkommen den Armen und Dürftigen Gutes thue, auf daß der Segen in meinem Hause bleibe und ein bleibend Wesen habe. Laß mich des gegenwärtigen Gutes zum Leben und nicht zur Sünde brauchen, daß ich von deinem Segen in deiner Furcht esse und trinke, meiner sauern Arbeit genieße und nicht ein anderer verzehre. O reicher Gott im Thron! Armuth und Reichthum gib mir nicht, laß mich aber mein bescheiden Theil Speise dahinnehmen, ich möchte sonst, wo ich zu satt würde, verleugnen und sagen: Wer ist der HErr? ich habe genug, wie kann mir's fehlen? Oder wo ich zu arm würde, möchte ich stehlen, und mich an dem Namen meines Gottes vergreifen, zürnen, Gott und meiner Obrigkeit fluchen. Lieber HErr und Meister! der du alles mit Lindigkeit und vielem Verschonen regierest, laß mich meinem Hause wohl fürstehen, Kinder und Gesinde christlich regieren, Arbeitern und Tagelöhnern nicht unrecht thun, damit nicht der abgebrochene Lohn der Arbeits- und Ackersleute und derer, die unser Land eingeerntet haben, wider mich gen Himmel schreie, und mir
\quellseite{138}
Sünde sei, sondern des Tages gebe, daß die Sonne nicht darüber untergehe. Laß mich nicht Wucher und Ueberfaß von meinem Nächsten nehmen und ihn beleidigen, daß mein Segen nicht verflucht werde, und mir alles unter den Händen zerrinne, daß ich mir genügen lasse an dem, das da ist, und gefallen, ich habe wenig oder viel; auch thue mir deine Gnade, daß ich mich in allem meinem Thun und Leben als ein christlicher Hausvater erzeige, und in deinen Gütern Treue beweise, damit ich in der Rechnung wohl bestehe und als ein frommer und treuer Knecht über viel gesetzt werde, zu meines HErrn Freude eingehe und selig werde um Christi willen, so will ich dir zeitlich danken, daß du mich mit deinen Gütern gesättigt hast, und dort mit Abraham, Isaac und Jakob, deinen Dienern, im Himmelreich dich ewig preisen. Amen.

\section{168. Gebet eines Kaufmanns}

O HErr, der du rechtmäßige Händel und Gewerbe dir wohlgefallen läßt, und mich zum Kaufmannstande berufen hast, ich bitte dich, du wollest mir verleihen, daß ich mich mit Fleiß in deiner Furcht halte, und ja nicht zu weit greife, und meinen Bruder verwortheile, sondern mit dem zeitlichen Gut recht handle im Kaufen und Verkaufen. Behüte mich, daß ich meinem Nächsten nicht vorsätzlicher Weise Unrecht oder Schaden thue, noch mein Gut mit Wucher oder Ueberfaß mehre, und noch dazu trüge, sondern, was ich will, daß mir geschehe, ich hingegen auch andern Leuten thue. Bewahre mich vor der leidigen Geldseuche, daß ich nicht das Zeitliche mehr als dich, der du bist mein höchstes Gut, liebe, und durch unordentliche Mittel mich zu bereichern begehre, auf daß ich nicht falle in Versuchung und Stricke und viel thörichter und schädlicher Lüste, welche versenken die Menschen ins Verderben und Verdammniß. Gib, daß ich mich begnügen lasse, und vielmehr das gegenwärtige Gut brauche, denn nach andern gedenke. Regiere mich, daß ich mir zuvörderst Schätze im Himmel sammle und am ersten nach deinem Reiche und seiner Gerechtigkeit trachte, auf daß mir nichts mangele an irgend einem Gut in Christo Jesu. Amen.

\quellseite{139}

\section{169. Gebet einer alten verlassenen Person}

O Herr Jesu Christe! mein einiger Heiland! der du der Elenden-Seufzen nicht verschmähest, und der betrübten Herzen Verlangen nicht verachtest. Dir klage ich aus hochbekümmertem Gemüthe, daß ich in meinem hohen Alter und Unvermöglichkeit von Jedermann verlassen, so gar einsam und elend bin, daß ich auch meinen Verwandten und Bekannten fremd bin worden. Denn meine Freunde haben mich vergessen, und meine Nächsten haben sich entzogen. Dich bitte ich demüthig, du wollest solch mein Elend in Gnaden ansehen, und ja nicht, wie andere Freunde, in dieser Einsamkeit mich verlassen, sondern nach deiner wahrhaften Verheißungsfest bei mir halten. Herr, wenn ich nur dich habe, so frage ich nichts nach Himmel und Erde. Und wenn mich schon alle Welt verläßt, so bist du dennoch meines Herzens einiger Trost und gewisser, treuer Freund im Leben und Sterben.

\begin{quote}
Auf dich hoffe ich, lieber Herr, In Schanden laß uns nimmermehr. Amen.
\end{quote}

\section{170. Gebet einer alten Person}

Herr Gott, unsere Zuflucht für und für, unser Leben und langes Alter, in dessen Händen unsere Zeit stehet, siehe, es sind die Jahre mir kommen, davon ich sagen möchte, sie gefallen mir nicht. Sintemal meine Kräfte in diesem meinem hohen Alter merklich abgenommen, und allerhand Beschwernisse und Schwachheiten sich sehr gehäufet haben, dich bitte ich demüthiglich, nachdem du mich vom Mutterleibe so gnädiglich und väterlich geführet, und bis hieher von Jugend an meine Hoffnung gewesen bist, du wollest mich ja nun im Alter nicht verlassen, da ich grau und schwach geworden bin, sondern ferner bis an meines Lebens seliges Ende mich heben, tragen und erretten. Insbesondere aber bitte ich dich, o gnädiger Vater, du wollest mit deinem heiligen und guten Geist mich regieren und führen, daß ich die noch übrige Zeit meines Lebens mit christlichem Wandel, und Beten und Seufzen, dir ganz ergeben und zum seligen Ende je mehr und mehr mich christlich bereiten
\quellseite{140}
möge, auf daß, wenn heut oder morgen meines Lebens Ziel kommet, ich mit dem alten Simeon meine Friedefahrt halten, und aus dieser Welt zum ewigen Leben abscheiden möge. Amen.

\section{171. Gebet bei Wahl eines Berufs}

O weiser und frommer Gott! der du unter deinen Creaturen alles in Zahl, Maaß und Gewicht geschaffen, sonderliche Stände verordnet hast, und willst, daß der Mensch in einem gewissen Berufe lebe, und auf seinen Wegen gehe, du hast auch zugesagt, demselben Glück und Segen zu verleihen. Nachdem ich denn meine Jahre erreicht, und gesonnen bin, in einen gewissen Stand zu treten und etwas Redliches darinnen zu ergreifen, damit ich einmal mich und die Meinen ehrlich nähren und mit gutem Gewissen und deinem Segen fortbringen könne; und aber in Niemandes Macht und Vermögen steht, wie er wandeln soll, sondern du mußt ihn berathen: so bitte ich dich kindlich, gib mir, daß ich erstlich meines Herzens Gedanken zu dir richte, ehe ich hierinnen etwas vornehme und anfasse. Laß mich deinen Willen durch das Gebet erforschen, und mit dir Rath schlagen, und Rath nehmen, daß ich darnach erwähle, was recht ist und dir wohlgefällt. Laß mich auch anderer guter und verständiger Leute Rath hören, und nicht verachten, und ja nicht etwas über mein Vermögen vornehmen. Sonderlich aber regiere mein Gemüth, daß ich in diesem Vorhaben mehr sehe auf deine göttliche Ehre und meines Nächsten Nutzen, als auf eigene Ehre, eingebildeten Reichthum und andere Eitelkeiten, und erkenne, was gut oder böse, löblich oder schändlich sei. Wann ich nun in meinem Vorhaben schlüssig sein und einen gewissen Lebensberuf erwählt haben werde: so gib dazu dein göttliches Gedeihen, und regiere mich mit deinem Geiste, daß all mein Thun und Lassen einen christlichen Anfang, glücklichen Fortgang und seligen Ausgang gewinnen möge, um Jesu Christi willen. Amen.

\section{172. Gebet, wenn einer ein Geschäft anfängt}

Ich danke dem Herrn von ganzem Herzen und erzähle alle deine Wunder, ich freue mich und bin fröhlich in dir und lobe deinen Namen, du Allerhöchster, daß du nicht allein von Kindesbeinen
\quellseite{141}
an mit mir gewesen und mich behütet auf meinen Wegen, sondern mir auch allezeit Brot zu essen und Kleider anzuziehen gegeben hast, und noch überdies mich bis hieher gebracht, daß ich nun hier selbsten mein Gewerb treiben und der zeitlichen Nahrung warten soll. Ach, lieber Gott! ich bin zu gering aller Treue und aller Barmherzigkeit, die du mir thust, bevorab, wenn ich gedenke an die Sünden meiner Jugend, und wie ich deine Gebote zurückgeworfen und nicht, wie ich gesollt, dich allezeit vor Augen gehabt. O, wie groß ist deine Güte über mich gewesen, daß du mir bisher geholfen und mich nicht aus dem Lande vertrieben, daß ich um meiner Sünden willen unstät und flüchtig sein müssen auf Erden, sondern mir vielmehr Ort und Stelle angewiesen, da ich im Lande bleiben und mich ehrlich nähren kann. Nun ich fange, mein lieber Gott, diesen meinen Beruf und Nahrung in deinem Namen an und will auch auf dein Wort das Netz auswerfen. Segne du mich nur mit Segen oben vom Himmel herab und laß dir hinfort die Werke meiner Hände befohlen sein. Sende mir Hilfe von deinem Heiligtum und stärke mich aus Zion. Fördere von nun an das Werk meiner Hände, ja das Werk meiner Hände wollest du fördern, damit ich etwas Redliches schaffe und zu geben habe den Dürftigen. Und weil zum Laufen nicht hilft schnell sein, zur Nahrung nicht hilft geschickt sein, zum Reichthum nicht hilft klug sein, sondern es liegt alles an der Zeit und Glück, so zeige mir selbst bequeme Zeit und Gelegenheit und verleihe mir Glück und Fortgang, daß ich meine Sachen recht ausrichten und alles, was ich thue, wohl gerathen möge. Laß deine Furcht allezeit in meinem Herzen und für meinen Augen sein, daß ich niemand Gewalt und Unrecht thue, noch einige Missethat in aller meiner Arbeit gefunden werde. O, du Vater des Lichts, von welchem alle gute und vollkommene Gaben von oben herab kommen, gib mir auch zu diesem meinem Beruf Gesundheit, Segen und Leben, damit ich desto eher meinen Zweck erlangen und, was du mir befohlen, verrichten kann. Und ob es mir schon ein wenig sauer werden sollte in meiner Nahrung, so laß mich es doch nicht verdrießen, sondern bedenken, daß es dein Wille ist, daß ich und alle Menschen im Schweiß unseres Angesichts unser Brod essen sollen. Hilf, daß ich nicht irdisch gesinnt sei und
\quellseite{142}
mein Herz mit Sorgen der Nahrung beschwere, sondern am ersten trachte nach deinem Reiche und deiner Gerechtigkeit, ja nach dem, was droben ist, da Christus sitzet zur Rechten Gottes; so weiß ich, daß, was das Zeitliche betrifft, es nach deiner milden Güte mir werde zufallen und ich keinen Mangel haben soll an irgend einem Guten. Nun, HErr, hilf, o HErr, laß wohl gelingen. Amen.

\section[173. Gebet, wenn einer etwas Wichtiges zu …]{173. Gebet, wenn einer etwas Wichtiges zu verrichten vor hat}

HErr, der du große Dinge thust, die nicht zu forschen, und Wunder, die nicht zu zählen sind, dessen Kraft auch in den Schwachen mächtig ist: ich rufe dich an, du Allerhöchster, laß doch mein Thun gelingen, und hilf mir gnädiglich das Werk vollbringen, so ich mit ganzem Vertrauen auf dich habe vorgenommen. HErr Gott, stärke mich in dieser Stunde, und verlaß nicht deinen Knecht, der alle seine Hoffnung auf dich setzet und ohne dich nichts thun kann. Rüste mich aus mit Kraft aus der Höhe, und gib mir die Weisheit, die stets um deinen Thron ist, daß sie bei mir sei, und mit mir arbeite. Es liegt alles an deinem göttlichen Segen; darum gedenke jetzt meiner, mein Gott, im Besten, und segne mein Thun und Lassen, daß ich alles klüglich anfange und glücklich ausrichte. Sprich Ja zu meinen Thaten, hilf selbst das Beste rathen, den Anfang, Mittel und Ende, ach HErr, zum Besten wende. Denn dir befehle ich meinen Ausgang und Eingang, dir will ich auch allen Ruhm und Fortgang meiner Verrichtung dankbarlich zuschreiben. Amen. In Jesu Namen. Amen, Amen!

\section[174. Gebet um rechte Verrichtung seiner …]{174. Gebet um rechte Verrichtung seiner Geschäfte}

Barmherziger, getreuer Gott! du hast befohlen, daß ein Jeder in einem gewissen Stande sich befinden, und desselben sich treulich annehmen soll. Nun siehst und weißt du, allwissender HErr, was mir obliegt, und wie ich ohne deine Gnade, Kraft und Segen sogar nichts vermag, und in wie viel Wege ich irren kann. Deshalb bitte ich dich von Herzen, daß du mir Verstand und Stärke verleihen, und einen tapfern und unerschrockenen Muth geben wollest, damit ich vernünftig handle, alle aufgestoßene Wi-

Gebetschatz. 8
\quellseite{143}
derwärtigkeit getrost überwinde, und durch deine väterliche Regierung einen gewünschten seligen Zweck erlange. Mein Gott, ich hoffe auf dich, laß mich nicht zu Schanden werden, daß sich meine Feinde nicht freuen über mich, denn keiner wird zu Schanden, der deiner harret. HErr, zeige mir deine Wege, und lehre mich deine Steige, denn du bist der Gott, der mir hilft, täglich harre ich dein. Bewahre meine Seele, errette mich, und hilf mir in allen meinen Anliegen. Laß mich nicht zu Schanden werden, denn ich traue auf dich. Du HErr, mein Gott, sei nicht fern von mir, eile mir beizustehen, HErr, meine Hilfe. Amen, Amen, Amen.

\section{175. Gebet vor der Arbeit}

Ewiger, barmherziger Gott! ich will nun wieder frisch an meine Arbeit gehen, und die Werke meines ordentlichen Berufs, in den du mich, deinen Diener, gesetzt hast, getrost angreifen, und nach deiner Gabe ausrichten, meinem Nächsten dadurch zu dienen, und mein Brot zu erwerben. Darum bitte ich dich herzlich, du wollest mir zu solcher vorgenommenen Arbeit und Werk, deiner gnadenreichen Verheißung nach, Sinn, Witz, Verstand, Glück, Heil und Segen, auch Gesundheit verleihen, damit ich es nicht allein wohl anfangen, sondern auch durch deine väterliche Hilfe in deinem Namen nützlich vollführen, und mein Amt jederzeit in deiner Furcht mit reinem Gewissen treulich abwarten möge. Dir befehle ich mein Werk, laß meine Anschläge fortgehen und laß mir gelingen, was ich mir in deinem Namen vornehme, daß das Licht auf meinen Wegen immer scheine, dir zu Lob und Preis, mir und den Meinigen, auch meinem Nächsten zur Wohlfahrt und Frommen durch deinen lieben Sohn, Jesum Christum, unsern HErrn. Amen. O HErr Jesu, in deinem Namen und auf dein Wort will ich mein Netz auswerfen. O HErr, hilf, o HErr, laß wohlgelingen. Amen.

\section{176. Gebet um Segen in der Berufsarbeit}

HErr, der du einem Jeglichen Arbeit auferlegst nach seinem Maß und willst, daß wir im Schweiß unseres Angesichts, bis wir wieder zur Erde werden, unser Brot essen, uns unsrer Hände Arbeit nähren, und es dabei gut haben sollen: siehe, auf dein Wort
\quellseite{144}
will ich das Netz meines Berufs auswerfen, laß doch meine Mühe und Arbeit nicht verloren sein, sondern hilf, daß ich mit meinen Händen etwas Redliches schaffe, auf daß ich habe zu geben den Dürftigen. Laß allezeit Glück und Segen sein, daran alles liegt, daß ich meine Sache ausrichte, und Niemand Unrecht thue, damit ich den HErrn ehre von meinem Gut und die Meinigen, sonderlich meine Hausgenossen versorge. Segne uns je mehr und mehr, segne unsern Korb und unser Uebriges, denn dein Segen ohne Mühe der fleißigen Hand macht reich. HErr, erfreue mein Herz, mache mein Angesicht fröhlich, gib mir Gesundheit, Leben und Segen. Ich erbitte aber dies Alles nach deinem Willen, HErr, erhöre mein Gebet. Dir befehle ich meine Wege, auf dich hoffe ich, du wirst es wohl machen. Gelobt sei der HErr täglich. Amen.

\section{177. Gebet für den Haushalt}

Getreuer Gott, himmlischer Vater, ein Liebhaber der Ordnung und Gott des Friedens! der du die Menschenkinder auf Erden in gewisse Stände gesasset, und dieselben also in einander gerichtet hast, daß immer einer dem andern die Hände bieten, und dem menschlichen Geschlecht in Leibes- und Seelenmühen behülflich sein solle. Ich danke dir dafür von Herzen, daß du mir zur Seelenspeise und ewigen Wohlfahrt täglich lässest predigen dein heilig Wort, und durch den Schutz der weltlichen Obrigkeit darneben mich also bewahrest, daß ich dabei ein stilles und geruhiges Leben führen kann in aller Zucht und Ehrbarkeit. Und bitte dich, laß deine Gnade auch ferner walten über meinem Haus, daß es durch deinen Segen je länger, je mehr erbauet werde; es hilft nichts, daß ich baue, wo du, HErr, nicht selber bauest: es hilft nichts, daß ich arbeite, pflanze und begieße, wo du, HErr, das Gedeihen nicht selber giebst; ja es hilft nichts, daß ich frühe aufstehe, hernach lange sitze, und mein Brot mit Sorgen esse, wo du, HErr, mir es nicht, wie andern, deinen Freunden, schlafend bescherst.

Ach, HErr Jesu Christ! es ist meinem Hause durch deine seligmachende Erkenntniß bisher Heil genug widerfahren, laß ihm nun weiter an deinem Segen auch nichts mangeln, alles was ich thue, das will ich thun in deinem Namen, und danken Gott und
\quellseite{145}
dem Vater durch dich. Dein Name ist mir ein festes Schloß, mit dem Gerechten laufe ich dahin, so werde ich erhalten. O, edler Tröster, Heiliger Geist! laß mein Haus, wie bisher, also noch weiter deine Werkstatt sein, daß ich und meine Hausgenossen, von dir geführet auf ebener Bahn, lerne thun nach Gottes Wohlgefallen. Wir für uns selber sind nicht tüchtig zum Guten, und haben aus uns selber weder das Wollen noch Vollbringen, daß wir aber tüchtig werden, das kommt von Gott durch deine Gnade. Darauf verlassen wir uns, und bitten demüthiglich, tröste uns in aller Mühe und Arbeit auf Erden, welche der HErr um unserer Sünden willen verflucht hat. Vereinige unsere Herzen in christlicher Liebe, daß wir uns dieses Leben nicht durch Uneinigkeit noch mühseliger machen, sondern im Frieden beisammen wohnen und des göttlichen Segens reichlich genießen mögen.

Insbesondere regiere, o Gott! das Hausgesinde, daß sie dich fürchten, und gedenken, daß sie Gott, und nicht den Menschen dienen, daß auch ein jeglicher, was er Gutes thun wird, von dir empfangen soll, er sei ein Knecht oder Freier; gib ihnen Verstand und Geschicklichkeit, daß sie alles, was dem Hause Schaden bringen kann, klüglich abwenden, und was ihm nützlich sein mag, heilsamlich befördern mögen. Regiere ihre Herzen, daß sie ihren Beruf mit uns treulich ausrichten, und thun, was du ihnen befohlen hast, so werden auch unsere Herzen stets zu dir gerichtet sein, und wie die Augen der Knechte auf die Hände ihrer Herren sehen, wie die Augen der Mägde auf die Hände ihrer Frauen, also sollen auch unsere Augen sehen auf den HErrn unsern Gott, bis er uns gnädig werde. Sei uns gnädig, HErr, sei uns gnädig, und segne unser Hauswesen mit allen Hausgenossen und aller Zugehör, wie du weißt, daß es erfordert unser zeitliches und ewiges Wohlergehen, nach deiner väterlichen Güte, durch JEsum Christum in der Gnade und Kraft des Heiligen Geistes. Amen.

\section[178. Gebet Luthers über den Spruch: Alle eure …]{178. Gebet Luthers über den Spruch: Alle eure Sorgen werfet auf Gott, denn er sorget für euch}

Himmlischer Vater, du bist ja mein HErr und Gott, der mich geschaffen, da ich nichts war, dazu mich erlöset hast durch deinen
\quellseite{146}
Sohn. Nun hast du mir dieß Amt und Werk befohlen und auferlegt, da geht es nicht, wie ich will, und ist so viel, daß mich drücken und ängsten will, daß ich bei mir selbst weder Rath noch Hülfe finde; darum lasse dir auch solches befohlen sein, gib du Rath und Hülfe, und sei selbst alles in diesen Sachen. Amen.

\section[179. Gebet eines solchen, den Gott mit …]{179. Gebet eines solchen, den Gott mit irdischen Gütern gesegnet hat}

O mein gütigster und barmherzigster Gott, dein heiligster Name sei gelobet, daß du mich, unwürdigen Menschen und Sünder, mit allem Ueberfluß zeitlicher Güter angesehen und bescheret hast. Was habe ich mehr um dich verdient, als andere neben mir, die in Mangel und Dürftigkeit ihr elendes Leben kümmerlich führen und zubringen müssen? Weil ich aber in deinem heiligen Worte, nicht ohne Bewegung meines Herzens, finde und lese, daß ich bei meinem Reichthum in großer Gefahr, meine Seele zu verlieren, lebe, indem dein liebster Sohn selber spricht: daß ein Reicher schwerlich ins Himmelreich kommen werde, und daß es leichter sei, daß ein Kameel durch ein Nadelöhr gehe, denn daß ein Reicher ins Reich Gottes komme; so erschrecke ich hierüber von Herzen, und möchte wünschen, der Allärmste auf Erden zu sein. Es tröstet mich aber, daß eben mein JEsus bald darauf hinzu thut: bei den Menschen sei es unmöglich, aber bei Gott sind alle Dinge möglich. Darum so bitte ich dich, mein Gott, von ganzem Herzen, in tiefster Demuth, um JEsu Christi theuerstes Verdienstes und Blutes willen, du wollest mir Gnade und Barmherzigkeit erzeigen, daß ich bei meinem Reichthum, den du mir zu meiner Verdammniß nicht wirst gegeben haben, nicht ewig verloren werde. In deinen heiligen Willen stelle ich es, mein Gott, ob du mich bei dem, was ich durch Erbschaften, Heirathen, Handlungen und so fort, vermittelst deines göttlichen Segens, empfangen habe und besitze, ferner lassen und erhalten, oder alles, wenn es mir an meiner Seligkeit hinderlich und schädlich sein sollte, hinwegnehmen, und mich lieber blutarm, doch dabei selig machen wollest. Denn was hülfe es mir, so ich die ganze Welt gewönne, und nähme Schaden an meiner Seele, und was könnte ich geben, meine Seele wieder zu erlösen? Es ist
\quellseite{147}
gewiß lauter nichts, ob mich gleich alle Welt glückselig preiset, und ich mich selbst sehr vergnügt in meinem Zustand befinde! Wenn mich nicht deine allgewaltige Hand und große Güte aus den Netzen des Satans errettet, und vor der Hölle bewahret, so bin ich der Allerunglückseligste, und wäre besser, daß ich nie einen Groschen in der Welt gehabt hätte, ja, daß ich nie geboren wäre. Darum bitte ich nochmals, durch deine Barmherzigkeit und durch meinen Herrn JEsum Christum, der mich so theuer erkauft hat, laß mich durch deinen Heiligen Geist in meinem Herzen meiner Seelen ewige Seligkeit vergewissert sein.

Würdige mich deines heilsamen Kreuzes, daß sich ein geistlich-armes, zerbrochenes und demüthiges Herz bei mir finde. Sei du mein allerhöchstes, einiges Gut, daß ich dich über alles herzlich liebe, und aus Liebe zu dir auch meinen armen Nächsten Liebe erweise, auch der Heiligen Nothdurft annehme, und zu den Werken der Barmherzigkeit willig sei. Fället mir denn aber Reichthum zu, daß ich mein Herz nicht daran hänge, sondern an dir, mein Gott, an deiner Furcht, Anbetung, Lob und Preis, mehr Freude und Trost als in allen zeitlichen Gütern suche und finde, erhöre und gewähre mich dieses, du Vater der Barmherzigkeit, und Gott, reich über alles, um JEsu Christi, meines Seligmachers willen. Amen.

\section{180. Gebet bei Antritt einer Reise}

Mein Gott! dir lebe ich, dir sterbe ich, dein bin ich todt und lebendig. In deinem Namen habe ich auch meine Reise angefangen. Deiner väterlichen Fürsorge vertraue ich mich, auf deine Allmacht verlasse ich mich. Und dir übergebe ich nun alles mein Gehen und Stehen, meine Heimfahrt und Wiederkunft. Segne du meinen Ausgang und Eingang, und hilf Alles vollenden zu deiner Ehre und meinem Besten. Sei mit mir und zeuch vor mir her auf dem Wege, den ich reisen werde. Bestell auch deinen Engeln über mir, daß sie mich auf demselben begleiten und allenthalben sich um mich her lagern, wo ich bleiben werde. Beschere mir verständige Wegweiser, die mich anleiten, gute Führer, die mich fort bringen, treue Reisegefährten, die mit mir ziehen, gutherzige
\quellseite{148}
Leute in der Herberge, die mich aufnehmen, und bringe mich also zur rechten Zeit an den Ort, wohin ich verlange. Erhalte meine Gesundheit auch in der fremden Luft, segne mir allenthalben Speise und Trank, so gut ich es antreffe. Bewahre meinen Vorrath, den ich bei mir habe. Laß mich überall klüglich handeln unter den Leuten, zu denen ich komme, und vorsichtig sein bei denen, die ich nicht kenne. Hilf aber, daß ich überall Gnade finde bei den Menschen. Beschere mir nothdürftigen Aufenthalt auch außer meiner Heimath. Laß mich alle meine Geschäfte wohl ausrichten. Gib, daß ich allenthalben etwas lerne aus alle dem, was ich bei andern Leuten sehe und höre; und hilf mir endlich, wenn es Zeit ist, mit Freuden wieder nach Hause zu den Meinigen. Amen.

\section{181. Gebet auf der Reise}

O HErr JEsu Christe, du allergütigster HErr und Heiland! der du um unserwillen ein Gast und Fremdling geworden bist auf Erden, ich komme zu dir mit demüthigem Herzen, und bitte dich, du wollest dir mein ganzes Leben und Wandel in dieser Welt lassen befohlen sein. Behüte mich auch auf dieser meiner Reise durch deiner lieben Engel Schutz auf allen meinen Wegen gnädiglich vor allem Uebel und Unfall, vor Unglück und Gefahr an Seele und Leib, damit ich alles glücklich ausrichten, und endlich auch frisch und wiederum gesund zu den lieben Meinigen gelangen möge. Dir, o HErr, allein befehle ich meinen Leib und meine Seele, mein Hab und Gut auf meiner ganzen Reise. Richte du meinen Gang und Wandel, wie es mir gut und selig ist. Nimm dich auch unterdessen der lieben Meinigen zu Hause treulich an, damit wir einander glücklich und gesund wieder mögen antreffen. Und wenn ich endlich auch meine Lebenswallfahrt hier in dieser Welt vollendet habe, so nimm mich zu dir durch einen sanften und seligen Tod in deine Herrlichkeit, als mein himmlisches Vaterland, auf und an. Amen.

\section{182. Gebet bei Ausfahrt in die See}

Siehe, ewiger Gott, ich fahre jetzt auf die Höhe, und ergebe mich der wilden, rohen See, da schwebe ich jetzt daher zwischen Himmel und Erde auf dem Wasser, auf diesem geringen und
\quellseite{149}
schwachen Gebäude und Schifflein, das sich selbst nicht, viel weniger mich schützen noch bergen kann, worauf soll ich armer Mensch mich denn verlassen? Ach, allein auf dich und deine väterliche Allmacht und Fürsorge, o Gott und HErr! denn weder Mensch, noch je eines deiner Geschöpfe können mich erretten. Der Wind bläset, wo er will, das Meer hebet empor seine Wellen, unter mir sehe ich keinen Grund noch Boden, oben aber den blößen Himmel, Menschen stehen mir zur Seite in gleicher Gefahr und Sorgen; darum Gott, hebe ich meine Augen auf allein zu dir, von dem mir Hilfe kommt, meine Hülfe kommt vom HErrn, der Himmel und Erde gemacht hat! Dir ergebe ich mich zu eigen mit Leib und Seel, Gut und Blut! Auf dich sehe ich allein, dir traue ich, meine Hoffnung stehet zu dir, Menschenhülfe ist keine! Trage du dies Schifflein auf deinen Flügeln, erhalte mich darin, und sei mit mir im Schiffe drinnen, denn sonst, o HErr, vergehen wir; wenn ich aber nur dich darin habe, so frage ich nichts nach Himmel und Erde! Ja, wenn ich jetzt vom Lande abfahre, und die Erde hinter mir lasse, so hilf, daß ich auch alle Zuversicht irdischer und menschlicher Hülfe verlasse, ihr absage, und mich dir allein ergebe. Siehe, HErr, ich ergebe mich ganz und gar in deinen Willen, mache du es mit mir, wie du willst und weißt, daß es zu meiner Seelen Seligkeit das Beste ist. Ich ergebe mich in deine Hände, denn du hast mich in deine Hände gezeichnet, und Niemand wird mich aus deinen Händen reißen, weil ich ein Schäflein Christi und seiner Weide bin, darauf ich gerne will leben und sterben durch Beistand des Heiligen Geistes. Amen.

\section{183. Gebet eines Seefahrenden im Sturm}

Allmächtiger Gott, du Gott der Gerechten! der du das Meer und Trockne gemacht, und Alles in deiner Hand hast, und erregst einen Sturmwind, der die Wellen erhebet, daß man mit dem Schiffe jetzt gen Himmel und dann in den Abgrund führet, daß man davon taumelt und wanket wie ein Trunkener, und keinen Rath mehr weiß, wie der Heilige Geist redet. In solcher großen Noth und Gefahr sind wir jetzt auch in dieser unserer Schiffahrt. Weil du nun gesagt: „Rufe mich an in der Noth, so will ich dich erretten,“ so halten wir dir in unsrer großen Noth, daraus uns keine Kreatur
\quellseite{150}
erretten kann, als du allein, auch mit gläubigem Herzen vor dein Wort, und bitten dich demüthiglich, du wollest nach deiner großen Barmherzigkeit deine Allmacht sehen lassen, Wind und Meer bedrohen, daß es ganz stille werde, vor Schiffbruch und Schaden uns gnädiglich behüten, mit deiner Gnadenhand ob uns halten und walten, aus dieser Gefahr uns gnädiglich erlösen: so wollen wir deinen Namen rühmen unser Leben lang, dieser Wohlthaten nimmer vergessen, unser Leben auch bessern, und hinfüro nach deinem heiligen Wort anstellen. Ach HErr, sei uns gnädig, und vergib uns unsre Sünde um JEsu Christi willen. Ach hilf uns, aber wir verderben! Ach HErr, erlöse uns um deiner Güte willen! Sollte es aber dein Wille sein, unser Leben hier zu enden, so verleihe uns christliche Geduld, erhalte uns im wahren Glauben, und gib uns das Ende des Glaubens, das ist der Seelen Seligkeit. HErr, sei bei uns! Verlasse uns nicht! Du kannst aus dem Tode erretten; so kannst du auch aus dieser Gefahr erlösen! Erhöre uns, erhöre uns, erhöre uns, HErr Gott Vater, um JEsu Christi willen! HErr JEsu, um deiner heiligen fünf Wunden willen! HErr Gott, Heiliger Geist, um deiner Güte und Gnade willen! Hilf uns, HErr Gott Vater! Hilf uns HErr Gott Sohn! Hilf uns, HErr Gott Heiliger Geist um dein selbst, um deines Namens Ehre und um deiner grundlosen Barmherzigkeit willen. Amen.

\section{184. Gebet in äußerster Wassernoth}

HErr! du hast mir versprochen: So du durch’s Wasser gehst, will ich bei dir sein, daß dich die Ströme nicht sollen ersäufen. Du kannst dieses als ein allmächtiger HErr wohl thun, mir in dieser äußersten Gefahr wohl helfen, und mich lebendig zum Lande bringen, denn du bist der HErr, mein Gott, der Heilige in Israel, mein Heiland: und ob ich ja hier versinken sollte nach deinem Willen, so weiß ich doch, daß du mich, als deinen Diener Mosen aus dem Schilf, könntest hervorziehen, und wo nicht hier zeitlich, jedoch ewiglich mit dir zum himmlischen Leben einführen. Hilf, HErr, hilf, denn wir vergehen. Ich bin in tiefem Wassern und die Fluth will mich ersäufen, wenn du aber hilfst, so ist mir geholffen, darum befehle ich dir meine Seele, Leib, Leben, Blut und
\quellseite{151}
Gott in deine Hände, dir lebe ich, dir sterbe ich, dein bin ich todt und lebendig, denn du bist mein und ich bin dein, und wo du bist, da will ich sein, wo ich bleibe, wirst du mich nicht verlassen, denn du bist Immanuel, Gott mit uns. Amen.

\section{B. Lehrstand}

\section{185. Gebet eines Predigers}

O getreuer Gott, barmherziger Vater! der du mich armen unwürdigen Diener zum heiligen Predigtamt berufen, und zu einem Menschenfischer gesetzt hast, daß ich viele Seelen zum Himmelreich fahen soll, und mich durch ordentliche Mittel zu diesem meinem Pfarrvolk, das Evangelium zu verkündigen, ausgesondert hast: Ich bitte dich, du wollest mich untüchtigen Menschen mit deiner Gnade tüchtig machen, zu führen das Amt des Neuen Testamentes, auf daß ich ein treuer Diener und Haushalter über deine göttliche Geheimnisse sei und meinem Amte wohl vorstehe, nicht gezwungen, sondern willig, nicht um schändlichen Gewinnes willen, sondern von Herzensgrund, aus Liebe deines heiligen Namens, und daß ich meine anbefohlenen Schäflein mit heilsamer Lehre treulich weide, der Schwachen warte, die Kranken heile, das Verwundete verbinde, das Verirrte wieder hole, und das Verlorene suche, und die, so von einem Fehl übereilt sind, wieder zurecht bringe mit sanftmüthigem Geist. Sende mir deinen heiligen Geist, der mir Mund und Weisheit gebe, zu reden, und dein Wort unterschied lich zu theilen. Laß mich nicht ein leer Instrument und Werkzeug deiner Gnade sein, sondern wirke du durch mich kräftiglich und gib das Gedeihen reichlich. O lieber Gott! du hast mich zum Wächter gesetzt über dein Volk: hilf, daß ich den Gottlosen warne, und die Uebertreter deine Wege lehre, daß sich die Sünder zu dir bekehren, und ihrer Viele zu deinem Reiche gewonnen werden. Verleihe auch, daß ich deinem Worte gleichförmig und gemäß lebe, und ein Fürbild werde der Heerde, auf daß ich nicht andern predige und selbst verwerflich werde, auf daß ich nicht mit unchristlichem Wandel Ursache gebe den Feinden, deinen Namen zu lästern, oder sonst
\quellseite{152}
Aergerniß anrichte. Gib mir und allen meinen Zuhörern deinen Heiligen Geist, daß wir dich, daß du allein wahrer Gott bist, und den du gesandt hast, JEsum Christum, erkennen, im rechten wahren Glauben, Furcht und Liebe wachsen und in fester Hoffnung mit bußfertigem Herzen beständig bis an unser Ende verharren, und also mit einander die unverwelkliche Krone der Ehren empfangen, wenn erscheinen wird der Erzhirt, JEsus Christus, unser Heiland, welcher mit dir lebet und regieret in Einigkeit des Heiligen Geistes, in Ewigkeit. Amen.

\section{186. Ein Anderes}

O lieber Gott! gnädiger Vater, du hast mich zu einem Prediger deines heilsamen Worts berufen, und weißt, daß ich deine Geheimnisse, ohne deine göttliche Hülfe und Beistand, würdiglich zu handeln, nicht vermag. Darum bitte ich dich herzlich, du wollest mir geben deinen Heiligen Geist, der mich leite, führe und regiere, daß ich in solchem meinem Beruf nichts anders bedenke, rede, vornehme oder thue, denn das du gereichet dir, o HErr! zu Lob und Ehren, und meinen befohlenen Christen zu Nutz und zu Trost.

Behüte mich, gnädiger Gott! für dem Vater der Lügen, der ein abgesagter Feind ist des heiligen Predigtamts, daß ich ja nichts anders lehre, denn dein Wort, welches allein die Wahrheit ist, und das Leben mit sich bringet. Daß ich auch ein solch Leben führe, dadurch meine lieben Christen nicht geärgert; sondern vielmehr gebessert werden, und nicht eine Ursache sei, daß dein heiliges gnadenreiches Wort von den Ungläubigen gelästert werde.

Insunderheit aber bitte ich dich, mein HErr und mein Gott, du wollest deine gnädige väterliche und tröstliche Verheißung, die du durch deinen Diener Jesaias hast verkündigen lassen, nämlich, daß dein heiliges Wort, so aus deinem Munde gehet, nicht wieder leer kommen soll; sondern thun, was dir gefällt, und soll ihm gelingen, dazu du es sendest, heut zu Tage an mir, deinem Knecht, desgleichen an denen, so es hören werden, gnädiglich erfüllen, und wahr lassen werden, also, daß meine Predigt gereichen möge dir, o HErr Gott, zu Lob und Danksagung, mir aber und allen Zuhörern zur Besserung unsers Lebens, zu Trost unserm schwachen Gewissen, zur Stärfung
\quellseite{153}
unsers Glaubens, und endlich zu unser aller Seelen Heil und Seligkeit, durch Jesum Christum, deinen geliebten Sohn, unsern HErrn. Amen.

\section{187. Noch ein Anderes}

HErr Gott, du hast mich in der Kirche zu einem Bischof und Pfarrherrn gesetzt, du siehest, wie ich so ungeschickt bin, solch groß und schwer Amt recht auszurichten, und wo es ohne deinen Rath gewesen wäre, so hätte ich es schon vorlängst alles miteinander verderbet. Drum rufe ich dich an. Ich will zwar gerne meinen Mund und mein Herz dazu leihen und neigen, ich will das Volk lehren, ich will selbst auch immer lernen und mit deinem Worte umgehen, und demselben fleißig nachdenken; brauche du mein als deines Werkzeugs; lieber HErr, verlasse du mich nur nicht, denn wo ich werde allein sein, so werde ichs leichtlich alles miteinander verderben. Amen.

\section{188. Noch ein Anderes}

Lieber Gott, ich habe angefangen zu predigen und das Volk zu lehren. Es will aber nicht fortgehen, es stößt sich hie und da; aber das schadet nicht. Weil du mir befohlen hast, dein Wort zu predigen, will ich davon nicht ablassen; mißrath es, so mißrath es dir; geräth es, so geräth es mir und dir. Hier ist dein Befehl, darauf gehe ich hin und werfe mein Netz aus und lasse dich sorgen, wie es gerathen werde. Amen.

\section[189. Gebet eines Predigers, die heilige …]{189. Gebet eines Predigers, die heilige Schrift fruchtbarlich zu studiren}

Ach lieber Gott, gefällt es dir, durch mich etwas auszurichten zu deinen Ehren, und nicht zu meinem oder einiges Menschen Ruhm, so verleihe mir aus lauter Gnade und Barmherzigkeit den rechten Verstand deines Wortes. Amen.

\section{190. Ein Anderes}

Gib, lieber HErr Gott, deine Gnade, daß ich dein Wort recht verstehe, und viel mehr, daß ich es auch thun möge. Siehe
\quellseite{154}
hoch, allerliebster Herr Jesu Christe, sollte dies mein Studiren nicht zu deinen Ehren allein gereichen, so lasse mich lieber keinen Buchstaben verstehen, und gib mir nur, so viel mir armen Sünder nutz ist zu deinen Ehren. Amen.

\section[191. Trost und Trost eines fleißigen …]{191. Trost und Trost eines fleißigen Predigers bei seinem Studiren}

Lieber Herr und Gott, ich schlafe in deinem Namen, und weiß, daß auch mein Schlaf dir wohlgefällt; wenn ich aber wache und meine gewöhnliche Arbeit thue in meinem Beruf, mit Schreiben, Lesen, Meditiren oder Betrachten und mit Beten, zweifele ich daran nicht, solche Arbeit ist dir auch angenehm, und wenn ich wüßte, daß es dir mißfällig wäre, wollte ich mich dessen viel lieber enthalten. Ich bin aber deß gewiß, daß ich dir wohlgefalle mit all meinem Thun, nicht um meinetwillen, der ich solches thue, sondern um deinetwillen, der du dich mein erbarmest, mir die Sünde vergiebest, mich liebest, führest und mit dem Heiligen Geist regierest. Amen.

\section{192. Gebet eines Predigers vor der Predigt}

Ewiger Gott, himmlischer Vater, ich unwürdiger, elender Diener deines Evangelii bekenne, daß ich von mir selbst zu diesem hohen und heiligen Predigtamt viel zu gering und ungeschickt bin, und daß ich ohne Hülfe und Beistand deines Sohnes Geistes keine bleibende Frucht bringen möge, darum rufe ich zu dir von Grund meines Herzens, im Namen und auf's Wort deines ewigen Sohnes und Redners, unsers einigen Hohenpriesters und Bischofs, du wollest deinen Gnadengeist heute durch mich reden lassen, und mir meine Lippen aufthun und meinen Mund füllen, daß meine Zunge deine Ehre und dieser Kirchen Bestes suche, und ich dein werthes Wort richten, unterscheiden und den Meinigen vortragen und weislich und bescheidenlich lehren, und meine Schäflein Lehre, Trost und Vermahnung draus lernen und behalten können, damit dein Lob in aller Welt gepreiset, des Satans Mord und Lügen durch den Mund der Säuglinge und Albernen zerstört, Glaube, Liebe, Trost, Leben, Geduld, Zucht, Unterthänigkeit und Gehorsam durch deines Geistes
\quellseite{155}
Amt in meiner befohlenen Schäflein Herzen gepflanzt, und ich nicht andern predige und selbst verwerflich werde, der du durch deinen lieben Sohn allein gute Gaben und Kraft zu deinem Evangelio gibst, hochgelobt in alle Ewigkeit.

Vater unser! der du bist im Himmel, geheiliget werde dein Name, dein Reich komme, dein Wille geschehe, wie im Himmel, also auch auf Erden; unser tägliches Brod gib uns heute, und vergib uns unsere Schuld, als wir vergeben unsern Schuldigern, und führe uns nicht in Versuchung, sondern erlöse uns von dem Uebel, denn dein ist das Reich und die Kraft und die Herrlichkeit in Ewigkeit. Amen.

\section{193. Ein Anderes}

O ewiger Gott, himmlischer Vater, ich elender Diener deines heiligen Evangelii bekenne, daß ich zu diesem hohen Amt von mir selbst viel zu gering und ungeschickt bin und ohne Hülfe und Beistand deines Heiligen Geistes nichts Gutes stiften noch ausrichten kann. Darum rufe ich zu dir von Grund meines Herzens und bitte, du wollest deinen Gnadengeist jetzt durch mich armes schwaches Werkzeug reden lassen und mir meine Lippen aufthun und meinen Mund füllen, daß meine Zunge deine Ehre und dieser Kirche Bestes suche, durch JEsum Christum, deinen lieben Sohn, unsern HErrn. Amen.

\section{194. Noch ein Anderes}

Ach mein HErr JEsu Christe! Gott und unser Bruder, Gottes und des Menschen Immanuel und Heiland, der du gesagt hast, ohne mich könnet ihr nichts thun: Ich trete nun in deinem Namen und Befehl auf die Kanzel, und bitte dich, verleihe mir deinen Heiligen Geist zu dieser meiner Predigt, welchen du von deinem Vater ohne alle Maaßen empfangen hast. Thue meine Lippen auf, daß mein Mund deinen Ruhm verkündige, stärke mein Gedächtniß, Verstand, Aussprechen und schwache Leibeskräfte, daß dein heiliges Wort rein, deutlich und mit Nutz vorgetragen, des Satans Mord und Lügen zerstöret, der seligmachende Glaube in meiner Zuhörer Herzen gepflanzt, und in uns allen die ewige Seligkeit befördert werde. Amen.

\quellseite{156}

\section{195. Noch ein Anderes}

Veni Sancte Spiritus! veni dator munerum! dona mihi os et sapientiam, ad annunciandum laudem tuam; effice, ne quod verbum mihi imprudenter excidat, quo possit aut nominis tui gloria, aut conscientia mea lædi; sed, ut cum fructu doceam. Da successum, accende corda auditorum, ut attente et magna cum reverentia verbum tuum audiant, atque inde emendentur. Amen.

Cor, mentem, linguam, Tu rege, Christe! meam.

\section{196. Gebet eines Predigers nach der Predigt}

O du wahrer Sohn Gottes und ewiger Redner, ich armer Diener habe auf deinen Befehl dein Wort ausgesäet und meinen lieben Pfarrkindern ausgetheilt, was du mir gebrochen und gegeben hast: nun bitte ich dich, mein lieber HErr JEsu, du wollest deinem Donner und Wort Kraft geben, und bestätigen in meinen Schäflein, was du durch mich in ihnen gewürkt hast, und durch deinen Geist Gedeihen zu meinem Pflanzen und Begießen sprechen, auf daß wir mit einander viele und bleibende Frucht bringen und dein Wort in ihnen wachse und wuchere. Vertheile, HErr, dein Wort in ihren Herzen, damit dein Name geheiliget werde, und ich mit meinem Pflündlein dir viel Seelen gewinne, die ich mit mir vor dein Angesicht bringe, wenn du die unverwelkliche Krone treuen Lehrern und fleißigen Zuhörern austheilen wirst; denn wir bekennen, daß ohne deinen Geist und Segen alle unsre Mühe und Arbeit vergebens ist; drum schließ du ihre Herzen auf und verwahr in ihnen dein Wort, daß es keine Angst, Noth und alle Teufel nicht wieder herausreißen können, der du bei deinem Predigtamt kräftig und geschäftig bist, und sammelst dir durch das mündliche Wort eine ewige Kirche aus allen Nationen, die dich hier und in alle Ewigkeit, mit deinem Vater und Geiste, preisen und loben. Amen, lieber HErr JEsu! Amen!

\section{197. Ein Anderes}

Gnädiger Gott, gnädiger Vater, ich danke dir von ganzem Herzen, daß du meinen Mund gefüllt und mich aus Gnaden rücktig
\quellseite{157}
gemacht, dein Wort zu verkündigen, und bitte dich um deines lieben Sohnes willen, du wollest dein Wort durch deinen Geist in meiner lieben Pfarrkinder Herzen bekleiben, wachsen und wuchern lassen, und dein Gedeyen dazu sprechen, daß es in ihnen viel Frucht bringe, und ihr Herz umzäunen, daß es der böse Geist nicht durch Sorge dieser Welt in ihnen dämpfe oder gar aus ihren Herzen reiße, damit Lehrer und Zuhörer Glauben und gut Gewissen bewahren und bei deiner Wahrheit bis an’s Ende beharren, und wir sämtlich vor deinem Angesicht in Ehren erscheinen und die unverweltliche Krone empfangen, die du aus Gnaden um deines Sohnes Blutes willen allen deinen gläubigen Dienern und Jüngern bereitet hast, der du dir durch dein Evangelium eine ewige Kirche aus menschlichem Geschlecht und dieser Gemeinde sammelst, und kräftig bist mit deines Sohnes Geiste beim gepredigten und geglaubten Wort, hochgelobt von allen seligen Predigern in alle Ewigkeit! Amen.

\section[198. Gebet eines Predigers vor der …]{198. Gebet eines Predigers vor der Beichthandlung}

O allmächtiger, ewiger Gott, ich soll jetzt verrichten das große, hohe Werk, darob sich nicht allein alle Vernunft, sondern ohne Zweifel auch die lieben Engel verwundern, nämlich an deiner Statt Sünde vergeben. Nun weißt du, wie ganz unwürdig und ungeschickt ich mich zu solchem erkenne, dieweil ich selbst ein armer Sünder bin, und es auch so schwerlich mißbraucht wird; du aber so treulich befohlen hast: wir sollen das Heiligtum nicht den Hunden geben und so edle Perlen nicht vor die Säue werfen. Und ob du wohl mir deinen Heiligen Geist zur Verrichtung meines Amtes verliehen hast, so muß ich doch leider bekennen, daß ich denselben mit meinen vielfältigen Sünden oft betrübt habe. Weil du aber deren Blut, die unterdessen durch mich verwahrlos’t worden sind, von meinen Händen fordern willst, ist es mir desto treulicher und herzlicher leid, daß ich dich jemals beleidigt habe. Und nachdem du nicht willst, daß wir unser Amt unterlassen sollen, bitte ich dich von Grund meines Herzens, um Jesu Christi, deines allerliebsten Sohnes willen: du wollest mir alle meine Missethat, Fehler und Gebrechen verzeihen, vergeben und vergeben, und deinen Heiligen Geist
\quellseite{158}
nicht von mir nehmen, sondern vielmehr mich mit denselben begnaden, erleuchten, heiligen, leiten und führen, damit ich, da ich den Leuten nicht ins Herz sehen kann, die erschrockenen Herzen trösten und stärken, die unbußfertigen aber schrecken möge, auf daß sie durch deines Heiligen Geistes Kraft und Wirkung zu beiden Theilen deiner Gnade und Barmherzigkeit herzlich froh werden, derselben desto weniger mißbrauchen, und sich mit Herzen, Mund und der That je länger je mehr und lieber zu dir bekehren und sich dankbar erzeigen, zu Lob und Ehr deines heiligsten Namens, ihrer eignen Wohlfahrt und Seelen Seligkeit und mir zu gutem Gewissen, andern aber zur Besserung und gutem Exempel, durch deinen lieben Sohn, unsern HErrn JEsum Christum. Amen.

\section[199. Gebet eines Predigers, wenn er etwas …]{199. Gebet eines Predigers, wenn er etwas Schweres zu verrichten hat}

Barmherziger Gott, du siehest, welch ein schweres Werk mir zu verrichten bevorsteht, und weißt auch, daß ich viel zu schwach und zu gering bin, solches ohne Schimpf und Schaden hinaus zu führen. Derwegen bitte ich dich demüthiglich, daß du mir darinnen gnädiglich beistehen wollest. Ach HErr, verlaß den nicht, der alle seine Hoffnung und sein Vertrauen auf dich setzt, sondern stärke mich mit himmlischer Kraft und rüste mich aus mit heiliger Weisheit, daß ich nicht kleinmüthig werde, oder unvernünftig handele, sondern das Werk, das mir obliegt, ersprießlich vollziehe. Denn dir befehle ich meinen Ausgang und Eingang, dir will ich auch allen Ruhm und Fortgang meiner Verrichtungen dankbarlich zuschreiben. Amen, in JEsu Namen, Amen, Amen.

\section{200. Gebet eines Predigers vor dem Krankenbesuch}

Allmächtiger, ewiger Gott, ein Vater unsers HErrn JEsu Christi, und Brunnen alles wahren Trostes, siehe, ich armer Diener gehe jetzt hin, wie mein Beruf erfordert, eine kranke, geplagte und vielleicht sterbende Person zu trösten. Weil denn solches nicht bei menschlicher, sondern göttlicher Hülfe stehet, so bitte ich, verleihe mir deinen Heiligen Geist, welcher deine Worte in meinen Mund legen, und meine Zunge und Lippen also regieren möge, daß meine

Gebetsbüch.

\quellseite{159}

Rede fließe und schmecke wie Honig, auch den Kranken erquicke, und im Glauben, Hoffnung und Geduld stärke, oder sanft zu sterben hurtig und fröhlich mache, um deines eingebornen Sohnes, unsers Herrn und Heilandes Jesu Christi willen. Amen.

\section{201. Gebet eines Seelsorgers in Pestzeiten}

Herr, allmächtiger Gott, es ist ja böse und gefährliche Zeit, und nimmt die auffallende Seuche unserer Sünden halben täglich überhand; weil aber der verflucht ist, der dein Werk lässig treibt (Jer 48, 10.), und ich Amts halben als ein Hirt billig bei meinen Schafen stehe, so bitte ich dich, du wollest mir armen Sünder gnädig sein und mich behüten, wie einen Augapfel, mich beschirmen unter dem Schatten deiner Flügel. Herr, du hast meinen Odem und alle meine Wege in deiner Hand, und sendest die Engel als dienstbare Geister aus zum Dienst um derer willen, die ererben sollen die Seligkeit, wie du Daniel mitten unter den Löwen, Jonas im Bauche des Wallfisches, Aaron, David und Jesaias in der grausamen Pestilenz mitten unter den Todten und Lebendigen erhalten, ach Herr! stehe mir auch bei in diesen meinen Nöthen, befehl deinen Engeln über mir, daß sie mich behüten auf allen meinen Wegen, daß mir kein Uebels begegne und keine Plage zu meiner Hütte sich nahe. Gib mir Herz und Muth und den Geist des Raths und der Stärke, daß ich nicht erschrecken müsse vor dem Grauen des Nachts, vor der Pestilenz, die im Finstern schleicht, vor der Seuche, die im Mittag verderbet, und hilf, daß ich als ein treuer Diener die Waffen meines Amts, nämlich das Gebet, redlich gebrauche, und also für die meinigen streite, und sonst unerschrocken thue, was mir Amts halben befohlen, und mein Leben wie eine Beute davon bringe. Herr, du bist meine Stärke und Kraft, meine Zuflucht in der Noth, du bist bei mir, daß du mir helfest und mich errettest, ob tausend fallen zu meiner Seiten und zehntausend zu meiner Rechten, so wird es doch mich nicht treffen. Auf dieses dein mir versprochenes Geleit traue ich von Herzen; denn deine Wahrheit ist Schirm und Schild. Wenn es aber ja anders von dir beschlossen, und es sollte die Reihe jehund an mich auch kommen, und ich meine Hütte ablegen, wie es denn ohne das ein Ende mit mir haben muß, und mein Leben ein
\quellseite{160}
Ziel hat und ich davon muß, denn meine Tage sind einer Hand breit bei dir und mein Leben ist wie nichts vor dir; so geschehe, Herr, dein Wille. Hilf nur, daß ich im Glauben lauter und unanstößig bleibe, das Ende desselbigen, der Seelen Seligkeit davon bringe, und lasse dir die meinigen befohlen sein, du Richter der Witwen und Vater der Waisen, um Jesu Christi willen. Amen.

\section[202. Danksagung eines Predigers nach …]{202. Danksagung eines Predigers nach vollbrachter Amtsarbeit}

Lieber Herr Jesu Christe, der du uns gerecht und selig gemacht und mir Kraft und Macht verliehen hast, die heilige Schrift auszulegen, und auch meinen Zuhörern, dieselbe anzuhören, erhalte und stärke mich und meine Zuhörer in dieser Lehre und gib Gnade, daß wir je länger je mehr wachsen und zunehmen in der Erkenntniß deiner Gnade und ungefärbtem Glauben, und behüte uns vor Sekten und falscher Lehre, auf daß wir mögen untadelig und unsträflich auf den seligen Tag unserer Erlösung erfunden werden. Dir sei mit dem Vater und Heiligen Geist Lob und Dank gesagt in Ewigkeit. Amen.

\section{203. Eine Andere}

Barmherziger, gnädiger Gott, der du der Anfang und das Ende bist, durch welchen alle Dinge geschaffen sind und noch erhalten werden, dir sage ich Dank, daß du mich dieses Werk, das ich in deinem Namen angefangen habe, durch deine Gnade hast zu Ende bringen lassen. Gib du mir auch diese Gnade dazu, daß solches Werk möge nützlich und ersprießlich sein. Ist etwas mangelhaft daran, so ist solches von mir; denn was gut daran ist, kommt von dir allein. Vergib du mir gnädiglich meine Mängel und ersetze dieselbigen mit deiner großen Vollkommenheit, bis du mich in dem Himmel auch zu einer Vollkommenheit wirst gelangen lassen. Amen, um Christi Jesu willen, Amen, Amen.

\section{204. Gebet eines Lehrers}

Barmherziger, gnädiger Gott und Vater! ich danke dir von Herzen, daß du mich nach deinem gnädigen Willen und Wohlgefallen
\quellseite{161}
in dieses Amt gesetzt hast, daß ich deine Täummer weiden, und die mit deines lieben Sohnes theurem Blute erkauften Kinder zu allem Guten anführen soll. Ich bitte dich demüthiglich, regiere mein Herz mit deinem Heiligen Geiste, auf daß ich früh und spät mit schuldigen Fleiß und Ernst all mein Thun und Lassen, Wissen, Wollen und Können also einrichte, damit alle meine Untergebenen täglich mögen zunehmen an Weisheit, Alter und Gnade bei Gott und den Menschen, und in wahrer Gottseligkeit, Frömmigkeit und Geschicklichkeit mögen wachsen als die schönen, an den Wasserbächen deines heiligen Wortes gepflanzten Bäume zu deines Namens Lob und Preis und ihrer ewigen Seligkeit. O mein Gott und Vater, schenke allen diesen deinen Kindern deinen Heiligen Geist, gib ihnen ein gehorsames Herz, und dabei beständige Gesundheit, Weisheit, Verstand und gut Gedächtniß, daß Gott und Menschen an ihnen ein herzliches Wohlgefallen haben, und sie fleißig studiren, löblich zunehmen, sich christlich und wohl verhalten, und mit der Zeit in nützlichen Verrichtungen und Zuständen wohl gebraucht werden mögen. Ach laß ja keinen von ihnen verloren werden, laß mich an keinem einzigen eine Unehre oder Schande, sondern an allen und jeden ferner Ehre und Freude, nach deinem väterlichen Willen und Wohlgefallen, zu deiner Ehre und ihrer ewigen Seligkeit, erleben, damit ich sie dir am jüngsten Tage fröhlich vorstellen und sagen möge: Hier bin ich und die Kinder, die du mir gegeben hast! Das verleihe mir Gott der Vater, durch Jesum Christum, in Kraft des Heiligen Geistes. Amen.

\section{205. Ein Anderes}

O du unsichtbarer und allweiser Gott! es ist mir das Amt befohlen, eine gewisse Lehre zu geben, und klärlich zu unterrichten die Unerfahrenen, die sich zu mir machen, und in die Schule kommen, zu lernen, was ihnen fehlt. Ja, du selbst hast mich ausgesondert und dargestellt, daß ich bauen und pflanzen, und die Jugend deine Furcht lehren soll. Nun bin ich selbst ein schwacher Mensch und gering im Verstande, dein aber ist beides Weisheit und Stärke. Du giebst den Weisen ihre Weisheit und den Verständigen ihren Verstand. Darum flehe ich vor deinem Angesicht von ganzem
\quellseite{162}
Herzen: Ach HErr, hilf mir, daß ich meine Arbeit mit Weisheit, Vernunft und Geschicklichkeit thun möge, deine Herde weide, und wohl zusehe, doch nicht um schändlichen Gewinnes willen, sondern von Herzensgrunde, und für allen Dingen weise, wie sich junge Leute halten sollen nach deinem Wort, wenn sie ihren Weg unsträflich gehen wollen, und lehre sie dein Wort halten von Jugend auf, daß sie dich fürchten, und die Sünde meiden. Sintemal die Furcht des HErrn der Weisheit Anfang ist, auf daß du vornehmlich, o HErr unser Herrscher, aus dem Munde der Unmündigen gepriesen werdest durch JEsum Christum. Und hilf, daß ich mit der Gabe, die ich empfangen habe, andern diene als ein guter Haushalter deiner mancherlei Gaben. Verleihe mir aber auch, daß die Jugend, wenn ich es ihnen mannichfaltig fürschreibe, rathe und lehre, den Hals unter der Weisheit Joch ergebe und sich ziehen lasse. Verleihe ihr, o Vater der Herrlichkeit, den Geist der Weisheit und der Offenbarung, zu deiner selbst Erkenntniß. Gieb ihnen Kunst und Verstand in allerlei Schrift und Weisheit, daß sie zu Vieler Verwunderung lernen und studiren. Erleuchte ihre Seelen durch die Weisheit, daß sie gedenken, wie das Dichten und Trachten ihres Herzens böse sei von Jugend auf, damit sie sich nicht dünken weise sein, sondern dich fürchten und vom Bösen weichen, daß sie sich gerne strafen lassen, und dadurch zur Klugheit kommen, daß sie dir allezeit danken, und dich bitten, daß du sie regierest. Laß sie in allem ihren Vornehmen deinem Worte folgen, auf daß sie hier und dort unsre Hoffnung und Krone des Ruhmes, ja unsre Ehre und Freude seien, die Eltern ergötzen, die Freunde erfreuen, und sie selbst endlich durch die Weisheit selig werden mögen. Wenn aber insgemein derjenige viel leiden muß, der viel lehrt, und ich etwa auch Undank für meine Wohlthat haben sollte: so verleihe, daß ich dennoch mit stillem Wesen arbeite, und nicht verdrossen werde, Gutes zu thun, sondern mein Amt redlich ausrichte, daß ich geduldig sei gegen Jedermann, und nicht sehe auf das Sichtbare, sondern auf das Unsichtbare und gedenke, mein Lohn werde im Himmel groß sein. Amen. Lob und Ehr, Weisheit und Dank, Preis und Stärke sei unserm Gott von Ewigkeit zu Ewigkeit. Amen.

\quellseite{163}

\section{206. Gebet eines Schülers}

Herr Gott, himmlischer Vater! sehe du mir bei mit deinem Heiligen Geist, und segne meine Arbeit, Studiren und Vorhaben, damit ich täglich zunehmen möge an Gottseligkeit, Frömmigkeit und Geschicklichkeit, ja an Weisheit, Alter und Gnade bei Gott und den Menschen. Nicht uns, Herr, nicht uns, sondern deinem Namen gieb die Ehre um Jesu Christi willen. Amen.

\section{207. Gebet eines College-Schülers}

Herr! gütiger Gott, barmherziger Vater! Es gefällt deinem göttlichen, allweisen und heiligen Willen, daß du dir ein Lob und eine Macht deiner Ehre aus dem Munde der Jungen und noch Unerwachsenen zubereitest, und sie durch gottselige Obrigkeit und Leute in Kirchen und Schulen väterlich erhältst und ernährst; unter denselben befinde ich mich auch durch deine Gnade, liebster Gott! und schreie demnach aus gläubigem Herzen zu dir. Herr, gieb mir deinen Heiligen Geist, und erfülle mein Herz mit deinem Gnadenglanz, daß ich dir mit Herz und Mund eifrig und andächtig diene, und alle Gutthaten, die ich empfange, zu deines heiligen Namens Ehre, deiner christlichen Kirche und Gemeinde Erbauung, und meiner Seelen zeitlichem und ewigem Heil anwende. Erleuchte meinen Verstand, barmherziger Vater! und stärke alle meine Kräfte des Leibes und des Gemüths durch dein Licht und Kraft, zu herrlicher Ausbreitung deines Lobes, und zur Erbauung deiner christlichen Gemeine, daß ich der Gottseligkeit, und gebührender Besuchung der Kirchen und Schulen, und anderer christlichen Tugenden mit Fleiß und Eifer mich befleißige, durch steten Fleiß und tägliche Uebung nach allem Vermögen die Weisheit, dich und deine Werke zu erkennen, und deinem heiligen Namen zu lobsingen, erforsche und suche, alle Tage und Stunden in heilsamer Wissenschaft und nützlichem Gebrauch deines heiligen Wortes und der edlen Musst wachse, und die Schwachheit meiner Jugend mit Eifer und Fleiß überwinde. Laß dir, liebreicher Vater! meine Jugend, Studiren und Thun in Gnaden befohlen sein, und komme mir in allen Stücken zu Hülfe, mit deiner Gnade und Heiligen Geist, daß ich meine jungen Jahre durch rechtschaffene Gerechtigkeit und Heiligkeit
\quellseite{164}
, dir aufopfere, in der Weisheit, und aller heiligen Sachen Wissenschaften bis zu erwünschtem Zweck fortfahre, und durch christlicher und gutthätiger Leute milde Hände ferner gewünschte Erleichterung meiner Armuth erlange. Ach! behüte meine Seele, liebster Gott, vor allerhand Lastern und Sünden, dadurch deines heiligen Namens Ehre verhindert, und meine Wohlfahrt beschädiget wird, und erhalte mich in deiner Furcht und wahrem Erkenntniß, ja in gottfeliger und rühmlicher Uebung aller Gebühr und Schuldigkeit, damit ich dich bis an mein Ende zu einem gnädigen Gott und liebreichen Vater habe, und nachmals von meinem allerliebsten Erlöser JEsu Christo zu der ewigen und unaussprechlichen Freude erhoben werde. Das alles laß mir aus Gnaden widerfahren, liebster Gott! um deiner Ehre und meiner Wohlfahrt willen. Amen.

\section[208. Gebet eines Studenten, der durch Armuth …]{208. Gebet eines Studenten, der durch Armuth aufgefordert wird, seine Studien aufzugeben}

Herr JEsu Christe, dem der himmlische Vater alles in seine Hände gegeben hat, ohne welchen wir auch nichts gedenken oder thun können, du weißt, mein Heiland, daß ich mein Studiren in deinem Namen angefangen habe, nicht um weltlicher Wollust, Ehre, Reichthums oder anderer nichtigen Dinge willen, sondern allein dir zu Ehren, damit dein Name unter uns geheiliget, und des Nächsten Nutzen gefördert werde. Weil dir aber auch daneben die menschliche Schwachheit wohl bewußt, sonderlich wie meine große Armuth mich an meinem guten Vorsatz hindern und abhalten will: so bitte ich dich, du wollest mit deiner vollkommenen und reichen Gnade bei mir sein, meinen Verstand, Kraft und Vermögen, das du mir gegeben, in solchem vorgenommenen Werk ferner richten, daß ich es seliglich und glücklich vollenden möge. Gib mir Beständigkeit, Mittel und Rath, daß es die Elenden sehen, und sich freuen, daß du deinem Knechte wohl thuest, und daher auch in ihrem Vornehmen gestärkt werden. Ja, sei du meine Stärke und Trost, der Anfang und das Ende, daß ich dich nach gnädiger geschehener Hülfe lobe und preise in alle Ewigkeit. Amen.

\quellseite{165}

\section{209. Gebet eines Studenten der Theologie}

Mein Gott, ich habe bisher nach Weisheit, Wissenschaft, Künsten und Sprachen getrachtet, gib du mir ferner Gnade, daß ich meinem Vorhaben mit Ernst und Eifer nachsetze und meine Zeit wohl anwende. Laß mich wohl bedenken, welch eine schwere Verantwortung ich haben werde, wenn ich die beste Blüthe meines Alters mit Müßiggang, Saufen, Fressen, Tanzen, Spielen und andern sündlichen Thorheiten hinbringe, der Eltern Schweiß und Blut verschwende und die Hoffnung meines Vaterlandes zu nichte mache. Laß mich auch daneben bedenken, daß es nicht genug sei, wenn ich meine natürliche Gabe wohl anlege und mich in allerlei Wissenschaften übe und fertig mache, sondern, wenn ich alles weiß, was mich die Welt lehren kann, daß das Vornehmste dennoch übrig, nämlich die himmlische Weisheit, die seligmachende Erkenntniß JEsu Christi, die Kunst des Glaubens, das Gesetz der Liebe, die Uebung der Gottseeligkeit, die Wissenschaft, selig zu sterben. Mit den andern Dingen kann ich mich bei der Welt beliebt machen, außer diesen aber kann ich dir, meinem Gott, nicht gefallen. Was hilfts mir aber, wenn ich der Welt gefalle und vor dir ein Greuel bin? Darum bitte ich dich demüthiglich im Namen JEsu Christi, du wollest mit deinem himmlischen Gnadenlicht mein Herz erleuchten, meinen Fleiß und Arbeit heiligen und segnen, und mich, nach deinem Wohlgefallen, zu deinem heiligen Dienst bereiten. Gib, daß ich mich befleißige, ein gut Gewissen allenthalben, beide gegen dich, o Gott, und den Menschen zu haben, daß ich in deiner göttlichen Furcht wandele, mich meines Taufbundes öfters erinnere, in Demuth, Sanftmuth, Keuschheit und Mäßigkeit lebe und mich unbefleckt behalte von dem heutigen gottlosen Wesen der Welt. Ich bin einer von den Nasiräern deiner Kirche, ein Verlobter Gottes zu deinem heiligen Dienst ausgesondert. Ich bin ein junger Baum, von welchem so viel einfältige Seelen dermaleinst viel Früchte der Lehre, des Trostes und Erbauung hoffen. Derowegen, mein Gott, gib mir ein erneuertes, erleuchtetes, gottseeliges Herz, daß ich mich eines heiligen Wandels befleißige. O Gott, mein Vater und Herr aller Güte, gib mir die Weisheit, die stets um
\quellseite{166}
deinen Thron ist, denn ich bin ein schwacher Mensch und kurzes Lebens, und zu gering am Verstand des Rechts und Gesetzes, und wenn gleich einer unter den Menschen vollkommen wäre, so gilt er doch nichts, wenn er ohne die Weisheit ist, die von dir kommt. Sende sie herab von deinem heiligen Himmel und aus dem Thron deiner Herrlichkeit, sende sie, daß sie bei mir sei, und mit mir arbeite, daß ich erkenne, was dir wohlgefalle. Ist es auch dermaleinst, mein Vater, dein göttlicher Wille, daß du mich in einen gewissen Stand setzen, oder auch, worauf ich denn, durch den innerlichen Trieb deines Heiligen Geistes, meine Gedanken und Studia gerichtet, in deiner Kirche dir zu dienen, schicken willst, so mache mich auch vollends dazu geschickt. Ich weiß, o allmächtiger Gott, daß du in deinem ewigen allweisen Rath allbereit die Zeit, den Ort und die Gelegenheit versehen hast, darum so laß ferne von mir seyn, daß ich dir solche bestimme und vorschreibe, oder auch durch solche Mittel, die die heutige Welt beliebet, und die von dir selbst also genannten Läufer gebrauchen, dir aber und deinem heiligen Wort nicht gemäß, sondern zuwider sind, dazu gelangen möge und mich nicht der Welt bewegen gleichstelle und mich nach ihrem Sinn füge. Sende mich, wo du hin willst, nur gib, daß deines hochheiligen Namens Ehre und meine und aller, die mich hören, Seligkeit befördert werde. Laß mich alsdann sein voll Kraft und Geistes, voll Rechts und Stärke, voll lebendigen Erkenntnisses, voll Glaubens, voll Liebe, voll gottseeligen Eifers, voll himmlischer Weisheit, voll Gottseligkeit, Demuth, Sanftmuth, Keuschheit, Mäßigkeit und anderer Tugenden, daß ich werde und sei eine neue Kreatur in Christo, ein wahrhaftig erleuchteter, bekehrter und heiliger Mann, ein Fürbild der Gläubigen im Wort, im Wandel, in der Liebe, in der Geduld, daß ich das Geheimniß des Glaubens in reinem Gewissen habe, und also die Lehre Gottes unsers Heilandes in allen Stücken zieren könne. Amen.

\section[210. Gebet eines Studenten der Theologie vor …]{210. Gebet eines Studenten der Theologie vor dem Studiren}

Stehe mir bei, mein Gott, du Vater des Lichts, von welchem alle guten und vollkommenen Gaben zu uns herabkommen, und regiere mich mit deinem Heiligen Geiste, damit ich dasjenige, was
\quellseite{167}
wahr, was richtig, was deinem göttlichen Wort gemäß, was der christlichen Kirchen erbaulich, mir selbst nützlich und selig ist, möge lesen, bedenken und vollbringen, zu deines allerheiligsten Namens Lob und Preis, um JEsu Christi willen. Amen.

\section{211. Nach dem Studiren}

Ehre sei Gott in der Höhe für seine Gnade und Beistand des Heiligen Geistes, welche er mir auch zu diesem Mal so reichlich verliehen hat. Fahre weiter fort, o Gott, lieber Vater, weil ohne des Heiligen Geistes Licht nichts Gutes in dem Menschen geschieht, und bestätige das Gute, das du in mir angefangen hast. Erleuchte meinen Verstand, stärke mein Gedächtniß, regiere meinen Willen und Begierden, damit ich dir möge gehorsam, und dich preisen, und alles, was ich lese, recht verstehen, was ich verstehe, wohl behalten, und was ich behalte, mir zu Nutz wohl anwenden. Mache mich zum heilsamen Gefäß deiner Gnade, zu deines Namens Lob und Preis und deiner Kirchen Aufnehmen, um JEsu Christi willen. Amen.

\section{C. Wehrstand}

\section{212. Gebet einer obrigkeitlichen Person}

O Gott, mein Vater, und HErr aller Dinge, der du alle Dinge durch dein Wort gemacht, und den Weisen bereitet hast, daß er herrschen soll über die Kreatur, so von dir gemacht ist, daß er die Welt regieren sollte mit Heiligkeit und Gerechtigkeit, und mit rechtem Herzen richten, du hast unter andern auch mich erwählet zum Regenten über dein Volk, zum Richter über deine Söhne und Töchter: ich bin aber ein schwacher Mensch, und kurzes Lebens, und zu gering im Verstande des Rechts und Gesetzes, und weiß weder meinen Eingang noch Ausgang. Ja, wenn gleich einer unter den Menschenkindern vollkommen wäre, gilt er doch nichts, so er ohne die Weisheit ist, so von dir kömmt. Darum flehe ich vor deinem Angesichte von ganzem Herzen, und rufe dich an mit meinem Gebet, denn du allein gibst Weisheit, und aus deinem Munde kömmt Erkenntniß
\quellseite{168}
und Verstand, du bist groß von Rath und mächtig von That. So laß nun über mir ruhen den Geist des HErrn, den Geist der Weisheit und des Verstandes, den Geist des Raths und der Stärke, den Geist der Erkenntniß und der Furcht des HErrn. Verleihe deinem Knechte ein gehorsames Herz, daß er dein Volk richten möge, und verstehen, was gut und böse ist. Gib mir Weisheit und Erkenntniß, daß ich vor diesem Volk ein- und ausgehe, und demselben fürstehe, wie sich’s gebühret. Sende die Weisheit herab von deinem heiligen Himmel, und aus dem Throne deiner Herrlichkeit, sende sie, daß sie bei mir sei und mit mir arbeite, daß ich erkenne, was dir gefalle; daß ich verstehe Gerechtigkeit und Recht, Frömmigkeit und allen guten Weg, und dich mit Ernst fürchte. Beschirme mich für Hoffart, daß ich sie weder in meinen Worten, noch in meinem Herzen lasse herrschen, dieweil sie der Anfang ist alles Verderbens, denn du willst, daß auch ein König sein Herz nicht erheben soll über seine Brüder. Neige mein Herz zu deinen Zeugnissen, und nicht zum Geiz, damit ich nicht Geschenke nehme, und das Recht beuge, und dadurch den Weg der Elenden hindere, sondern daß ich den Armen gerne höre, und ihm freundlich und sanft antworte; daß ich Recht schaffe den Armen und Waisen, und helfe dem Elenden und Dürftigen zum Recht; daß ich keine Person im Gericht ansehe, sondern den Kleinen höre, wie den Großen, und für Niemandes Person mich scheue, dieweil das Gerichtamt dein ist, und wir solches nicht den Menschen, sondern dir halten, der du mit uns im Gerichte bist. Hilf, daß die Lügner bei mir nicht gedeihen, welche der Wahrheit widerstehen, und Tücke erfinden, zu verderben die Elenden mit falschen Worten, und dadurch der Gottlosen Sache gut machen, daß ihr Recht erhalten wird. Laß mich also thun in der Furcht des HErrn treulich und von rechtem Herzen, so wird dir mein Werk angenehm sein und werde dein Volk recht richten. HErr Gott, verleihe, daß wir und alle Obrigkeit, so von dir verordnet ist, Land und Leute mit Rath und Verstand der Schrift, und im Frieden regieren, und stets an den Namen des HErrn gedenken. Hilf auch, daß dein Volk unterthan sei aller menschlichen Ordnung um deinetwillen, und mit rechtschaffenen Herzen aufrichtig wandele, und für uns bete, auf daß sie also ein geruhiges und stilles Leben
\quellseite{169}
führen mögen in aller Gottseligkeit und Ehrbarkeit. Segne uns, o unser Gott, und mehre unsere Grenzen; laß deine Hand mit uns sein, schaffe es mit dem Uebel, daß es uns nicht kümmere. Herr, laß kommen, was ich bitte. Ich hoffe auf dich, du Allerhöchster, darum werde ich durch deine Güte fest bleiben. Dir, als dem König aller Könige, und Herrn aller Herren, der allein Unsterblichkeit hat, sei Ehre und ewiges Reich. Amen.

\section{213. Gebet einer richterlichen Person}

Herr, allmächtiger Gott! der du unter den Göttern Richter bist, fromme und getreue Richter gibst, und auch hinwegnimmst, durch deine göttliche Vorsehung bin ich von Mittelspersonen dazu erwählet und berufen, daß ich das Recht aussprechen, und alle Uneinigkeit und Irrungen, die zwischen Bruder und Brüder, Freund und Freund, und also zwischen einem Jeden und seinem Nächsten vorfallen, entscheiden und richtig machen soll: Herr, ich erkenne mein Unvermögen und meinen Unverstand, mache du mich, du Gott des Gerichtes, tüchtig zu diesem Amt, verleihe mir Verstand und gib mir einen Muth, daß ich den Kleinen, wie den Großen, und den Armen, wie den Reichen höre und richtig antworte, daß ich den unbewußten Sachen wohl nachforsche, und sie recht erfahre, ob es so sei, wie ich berichtet werde, oder ob es nicht so sei, daß ich es wisse, daß ich keine Person im Gerichte ehre, ansehe, achte und vorziehe, daß ich mich für keiner Person scheue, nach Recht zu strafen, oder jemand durch die Finger sehe. Du gerechter Richter! der du Herzen und Nieren prüfest, und einem jeden ins Herz siehest, einem jeglichen däucht sein Weg recht sein, aber du Herr allein machst die Herzen gewiß, leite mich in deiner Furcht, die weislich thun macht in allen Sachen. Laß mich Gerechtigkeit und Wahrheit lieben, dem Geize feind sein, daß mich die Geschenke nicht blenden, das Recht zu beugen, gerechte Sachen unrecht zu machen, und jemand mit Unrecht zu verdammen, daß ich nicht, wie die Richter zu Bersaba und Babylon, mit Lügen umgehe und verkehrte Urtheile spreche. Gib mir nicht zu, daß ich jemand mit unbilligen Auflagen und übrigen Unkosten beschwere, und mit dem ungerechten Richter die Sache aufziehe, sondern den rette, dem Gewalt geschieht, von dem,
\quellseite{170}
der ihm Gewalt thut, dem Unschuldigen seiner Unschuld Zeugniß gebe und für unrechter Sache mich hüte, und also mein Amt in Gottes Furcht führe und verrichte, und nicht mit Grimm strafe. Du Richter der Welt! laß nicht Zorn und Neidischard Richter sein, daß ich nicht etwa den alten Haß büße, und mich auf jemand erhitze, und mein Müt(h)lein fühle, sondern ein recht Gericht richte, dem Rechten nachjage, und ob demselben halte, niemand zu unrechter Sache diene, noch das Recht des Gerechten von ihm wende, daß ich nicht nach dem Ansehen, sondern nach Maaß und Zahl der Missethat strafe, und also alle Sachen nach Gottes Wort richte, und das vollziehe. Laß mich bedenken, daß mich meine Sünden, da ich unrecht Urtheil gebe, den Schuldigen losspreche, und den Unschuldigen verdamme, endlich treffen würden, und daß du zu seiner Zeit fragen werdest, was ich gehandelt, und forschen, was ich geordnet, und daß die, welche ihre Gewalt mißbrauchen, gewaltig werden gestraft werden. Du hast gesagt, sei ferne von falschen Sachen, laß mir dein Gebot vor Augen, und mein Herz ferne von falschen Sachen sein, daß ich nicht das Unrecht verschlinge und alles, was richtig ist, verkehre. Du Richter der Lebendigen und der Todten! laß mich unter den frommen Richtern der heiligen Schrift sein und gezählet werden, und mit ihnen eine herrliche Belohnung in dieser und der zukünftigen Welt erlangen, um dein selbst willen. Amen.

\section[214. Gebet eines Geschwornen, der in einem …]{214. Gebet eines Geschwornen, der in einem Criminal-Falle richten soll}

HErr, du eifriger starker Gott, der du willst, daß man die Bösen von sich thun, und, nachdem ihre Verbrechen, sie an Leib und Leben strafen soll, doch auch hinwiederum des Menschen Blut hoch und theuer achtest, daß du es nicht unschuldig vergießen lassen willst. Ich bitte dich, weil es die Ordnung meines Berufs also erfordert, daß ich bei einem Urtheile, das Leib und Leben antrifft, sitzen, und mich auf meinen Eid und Seligkeit darüber verbinden muß, daß solches Urtheil vor Gott und den Menschen recht und billig sei, du wollest mir den Geist der Weisheit jetzt mittheilen, damit ich mit gutem Verstande die Sachen oder Verbrechen der versündigten Person genugsam erkennen, auch das verfaßte Urtheil,
\quellseite{171}
ob es auf solches abschwebende Verbrechen nicht zu scharf sei, wohl bedenken möge. Ach Herr, mein Gott! laß mich nicht bei einem ungerechten Urtheile sitzen und dem Unschuldigen das Leben absprechen helfen. Wann ich es aber in meinem Gewissen recht befinde, so gib mir, Herr, auch einen freudigen Geist, zu reden, was mein Gewissen mich unterweiset, und in das verfaßte Urtheil beherzt und getrost zu willigen, dieweil es besser ist, ein faules Glied am Leibe abzuschneiden, und einen bösen Menschen aus der Gemeine durch Recht beiseit zu thun, als durch Gelindigkeit der Bösen zu schonen, und deinen Zorn zu reizen, daß er über die ganze Gemeine zur Rache und Strafe entzündet werde. O Herr! gib Weisheit und Verstand, Muth und Herz in allen Stücken meines Berufs. Amen, Amen.

\section[215. Gebet eines Amtmanns, Juristen oder …]{215. Gebet eines Amtmanns, Juristen oder Rathsherrn}

Allmächtiger, ewiger Gott, himmlischer Vater, du hast also geordnet und befohlen, daß ich Jura studiren, lernen und sprechen soll, was recht ist; darum gib du Gnade und Segen dazu, daß ich allein die Wahrheit suche und finde und thue nur, was dir gefällig ist, dir zu Ehren, und Land und Leuten zu Nutz, um Christus willen, deines lieben Sohnes, meines Herrn und Heilandes. Amen.

\section{216. Allgemeines Gebet eines Soldaten}

Himmlischer Vater, hier bin ich nach deinem göttlichen Willen in diesem äußerlichen Werk und Dienst meines Oberherrn, wie ich schuldig bin, dir zuvor und demselben Oberherrn um deinetwillen, und danke deiner Gnade und Barmherzigkeit, daß du mich in solch Werk gestellt hast, da ich gewiß bin, daß es nicht Sünde ist, sondern recht und deinem Willen ein gefälliger Gehorsam ist. Weil ich aber weiß, und durch dein gnadenreiches Wort gelernt habe, daß keines unserer guten Werke uns helfen mag, und Niemand als ein Krieger, sondern allein als ein Christ muß selig werden, so will ich mich gar nicht auf solch einen Gehorsam und Werk verlassen, sondern dasselbige deinem Willen frei zu Dienst thun, und glaube von Herzen, daß mich allein das unschuldige Blut deines lieben Sohnes, meines Herrn Jesu Christi, selig mache, welches er für
\quellseite{172}
mich, deinem gnädigen Willen gehorsamlich, vergessen hat. Da bleibe ich auf, da lebe und sterbe ich drauf, da streit’ und thu ich Alles auf. Erhalte, lieber Herr Gott Vater, und stärke mir solchen Glauben durch deinen Geist. Amen.

(Willst du darauf den Glauben und ein Vaterunser sprechen, magst du thun und lassen damit genug sein, und befehl damit Leib und Seele in seine Hände, und zeuch dann vom Leder und schlage drein in Gottes Namen.)

\section{217. Ein Anderes}

Herr Jesu Christe, du himmlischer Siegesfürst und Ueberwinder der Sünden, des Todes und des Teufels! Ich bitte dich, du wollest zu meiner Kriegsverrichtung deinen gnädigen Segen verleihen und mir wegen meiner unzähligen Sünden, für welche alle du reichlich genug gethan, denselben nicht versagen. Gib mir auch ein gottselig Herz, daß ich mich für Fluchen, Schwören, Schwelgen, unziemendem Raub, Plünderung und andern Sünden hüte, meinen Vorgesetzten treu und gehorsam verbleibe, und für meinem Feind nicht erschrecke, sondern in deinem Namen frisch und freudig anfalle, mich in aller Gefahr mit festem Glauben an dich halte, und von dir mich weder Wunden noch Tod lasse wankelmüthig machen. Gefällt es dir, daß ich soll gefangen oder verwundet werden, oder halb todt auf der Wahlstatt liegen bleiben, so gieb mir in meinen Schmerzen Geduld, und laß mich alles für deinen väterlichen Willen erkennen und aufnehmen. Wolltest du aber mich entweder durch Feuer und Schwert, oder was es sein möchte, ganz von hinnen hinweg nehmen: so gieb mir um deines schmerzlichen und schmählichen Todes willen ein seliges Ende, und laß mich nimmermehr von dir geschieden werden. Amen.

\section{218. Morgengebet eines Soldaten}

Herr, allmächtiger Gott, barmherziger lieber Vater, wir danken dir von Grund unsrer Herzen, daß du die Zeit unsres Lebens, wie auch diese vergangene Nacht über, so gnädiglich uns erhalten und behütet hast; und bitten dich in dem Namen Jesu Christi, deines Sohnes, du wollest um desselben theuren Leidens und Sterbens willen uns Allen unsre Sünden verzeihen, und uns
\quellseite{173}
durch deinen Heiligen Geist je länger, je mehr erleuchten und regieren, daß wir von Herzen Buße thun, in deiner Furcht wandeln, dir allein vertrauen, dich lieben, und für aller Schande und Laster uns hüten. Bewahre uns auch heute diesen Tag vor allem Unfall und segne unsern Ausgang und Eingang, Thun und Lassen, daß es gereiche zu deines heiligen Namens Ehre, deiner Kirche und des Vaterlandes Rettung und zu unserer aller zeitlichen und ewigen Wohlfahrt. Insbesondere regiere mit deinem Geiste und beschütze wie deinen Augapfel unsern Herrn General, sei du, o HErr, dessen Schirm und Schutz, und vollbringe seine Anschläge; mache zu nichte alle bösen Anschläge der Feinde und ihre große Gewalt. Verhüte und zerstöre alle Untreue. Regiere, bewahre und segne unsre übrigen hohen und niedern Offiziere und das ganze Kriegsheer, HErr Zebaoth, du mächtiger Kriegsfürst, stehe uns bei. Verleihe uns allen einen unerschrockenen, getrosten Muth, ziehe aus mit unserm Heere, beschere uns erwünschten Sieg, und laß über uns Alle walten deine väterliche Gnade um deines allerliebsten Sohnes, Jesu Christi, unsers einigen Heilandes willen. Amen.

\section{219. Ein Anderes}

Allmächtiger, barmherziger Gott und Vater! ich sage dir Lob, Ehr und Dank, daß du mich armen Sünder also väterlich beschirmst und erhältst, obgleich ich so mannigfaltig wider deinen göttlichen Willen, Gnade und Güte gesündigt. Und bitte dich, HErr, mein Gott, du wollest mir gnädig und barmherzig sein, und mir verzeihen und vergeben alle Missethat. Nimm weg mein steinern Herz, zeuch ab von mir den alten Menschen, erneuere mein Herz und Geist, daß ich dich, meinen höchsten Schatz, erkenne, vollkommiglich lieb habe, deine heiligen Gebote emsiglich betrachte, und in ihnen mit höchstem Fleiß und heiligem Gehorsam unsträflich wandle. HErr, allmächtiger, treuer Gott, ich bitte dich durch das bittere Leiden und Sterben deines lieben Sohnes Jesu Christi, behüte mich heut und allewege vor Verunehrung deines heiligen Namens und Uebertretung deines göttlichen Willens, nimm mich mir und meiner bösen vergifteten Natur, und gieb mich dir, daß ich mein selbst in rechter Gelassenheit ganz und gar ledig stehe, mich
\quellseite{174}
an dich ergebe durch wahren lebendigen Glauben, starke Hoffnung und brünstige Liebe. O HErr JEsu Christe, mein Herzog und Hauptmann, ich bitte dich durch die Liebe deines himmlischen Vaters, gehe mir vor mit deiner Gnade, führe mich in der rechten Straße deines heilsamen Wortes durch diese irrige und schalkhaftige Welt, daß ich nicht von jemand verführt, von deinem Worte abfalle, dein vergesse und unchristlich lebe, sondern behüte mich, mein lieber HErr JEsu, und alle Menschen vor Sünde und Unglauben, vor bösen Gedanken, Worten und Werken und ärgerlichem Wandel, auf daß dein himmlischer Vater in und durch uns gelobt, geehrt und gepriesen werde mit dir und dem Heiligen Geiste, der uns aus herzlicher, inbrünstiger Andacht beten und sprechen mache im Geist und in der Wahrheit: Vater unser, der du bist im Himmel, geheiliget werde dein Name, dein Reich komme, dein Wille geschehe, wie im Himmel, also auch auf Erden; unser täglich Brod gib uns heute, und vergib uns unsere Schuld, als wir vergeben unsern Schuldigern; und führe uns nicht in Versuchung, sondern erlöse uns von dem Uebel, denn dein ist das Reich, und die Kraft, und die Herrlichkeit, in Ewigkeit. Amen.

\section{220. Abendgebet eines Soldaten}

Allmächtiger, gnädiger Gott und Vater, wir sagen dir von Herzen Lob und Dank, daß du auch diesen vergangenen Tag so väterlich uns erhalten, geführet und behütet hast. Und bitten dich ferner im Namen deines allerliebsten Sohnes, Jesu Christi, du wollest um seinetwillen gnädig sein unsrer Missethat, die da sehr groß ist, und uns auch diese Nacht für aller List und Gewalt des Satans, für Sünden und Schanden, desgleichen für feindlichem Einfall und allem Schaden behüten. Bekehr uns auch durch die Kraft deines Heiligen Geistes und gib, daß wir im wahren Glauben und gottseligen Leben, wie Christen gebühret, dir dienen und ein unverletztes gutes Gewissen behalten mögen. HErr, unser Gott, der du nicht schläffst, noch schlummerst, laß dir in deinen allmächtigen Schutz gnädiglich befohlen sein unsern Herrn General, regiere denselben allezeit durch deinen Geist, behüte ihn durch deine heiligen Engel für allem Uebel und segne alle seine guten Anschläge.

\quellseite{175}

Bewahre auch unsere übrigen hohen und niedrigen Offiziere und das ganze Kriegsheer. Zerstöre alle List und Gewalt unsrer Feinde, und behüte uns auch vor aller Untreue und bösen Anschlägen. Sei du eine feurige Mauer um uns her. Denn wo du uns, HErr, nicht bewahrest, so wachen die Wächter umsonst, wo du uns nicht hilfst, so ist mit menschlicher Macht nichts gethan. Darum, o himmlischer Vater, in deine Hände befehlen wir unsere Leiber, Seele und Alles, was du uns gegeben hast, im Namen deines Sohnes, unseres einigen Heilandes, Jesu Christi, welcher uns also hat gelehret zu beten: Vater unser! der du bist im Himmel, geheiliget werde dein Name, dein Reich komme, dein Wille geschehe, wie im Himmel, also auch auf Erden; unser tägliches Brod gib uns heute, und vergib uns unsere Schuld, als wir vergeben unsern Schuldigern, und führe uns nicht in Versuchung, sondern erlöse uns von dem Uebel, denn dein ist das Reich und die Kraft und die Herrlichkeit in Ewigkeit. Amen.

\section{221. Ein Anderes}

Allmächtiger, barmherziger Gott und Vater, ich sage dir herzlich Dank, daß du mich armen, sündigen Menschen heut diesen Tag ganz väterlich behütet und mein Leib und Seele, Leben und Gesundheit bewahret hast. Du hast mich mit deinen Fittigen bedecket, daß mir kein Uebels begegnet, und deinen Engeln befohlen über mir, daß sie mich wider des Satans Trutz vertheidigt haben. Deß danke ich dir und lobe dich, Gott, daß du mein Schutz und Hilfe bist, und bitte dich um deiner grundlosen Barmherzigkeit willen, vergib mir meine Sünden, mit welchen ich dich in Worten oder Werken, Gedanken oder Geberden mag erzürnt haben. Laß meine Sünde im Untergang, deine Gnade aber im Aufgang sein von wegen der überreichen Bezahlung deines geliebten Sohnes, meines einigen Erlösers, Jesu Christi, in welchem du uns alle Sünden vergeben willst. Ich bitte dich aber ferner, liebe mich bis an das Ende, und erzeige mir Güte und Treue, die mich behüten. Bewahre mich diese Nacht für allem Schaden und Gefahr, und laß mich ganz mit Frieden sicher schlafen und ruhen. Oder so ich ja meines Standes wegen eilend auf sein sollte, so verleihe mir wackere
\quellseite{176}
Augen, ein getrostes Herz und fertige Hände und Füße, auf daß ich meinen Dienst unverdrossen und glücklich ausrichten möge. So ich auch gar diese Nacht sollte übereilet werden, und den Weg aller Welt gehen, und unversehens sterben müßte, so ergötze doch meine Seele durch deinen Trost, und laß mir ohne deinen Willen nichts Gefährliches widerfahren, sondern mache es, daß ich’s ertragen und selig abfahren kann. Wenn es aber deinem Willen wohlgefällig und meiner Seele nützlich ist, so errette mein Leben vom Verderben und erwecke mich zu rechter und bequemer Zeit, daß ich die liebe Sonne frisch, gesund und fröhlich anschauen, mit Freuden an meinen Beruf wieder gehen, und dir und meinem Nächsten treulich dienen möge. Dafür will ich dir zeitlich und ewiglich Lob und Dank sagen, und diesen meinen Abendsegen schließen und sprechen: HErr, segne und behüte mich; HErr, erleuchte dein Angesicht über mich und sei mir gnädig; HErr, erhebe dein Angesicht auf mich und gib mir Friede. Mich segne Vater und der Sohn, mich segne Gott der Heilige Geist von nun an bis in Ewigkeit. Amen.

\section{222. Gebet um herzhaften Muth}

Starker großer Gott! der du gibst Herz und Muth nach deinem Wohlgefallen denen, so auf dich trauen, ich bitte dich, weil ich jetzt in einem sehr gefährlichen Stande bin, da ich wohl Kraft, Stärke und Muth bedarf, du wollest von mir reißen allen natürlichen Kleinmuth und Blödigkeit, dazu ich von Natur als ein sündlicher Mensch geneigt bin. Strafe mich nicht mit einem erschrockenen Herzen und verzagten Geist, darüber ich mich sehr betrüben würde. Doch so dir nach deinem göttlichen Rath gefiele, mein Herz blöde zu machen, damit ich mich nicht auf meine Stärke, sondern auf deine Hilfe verlassen lernte, so laß mich nicht gar verzinken, gib mir mitten im Kleinmuth einen Blick deiner Stärke, daß ich mich dessen tröste, du seiest mit mir und wollest mich gewaltig in aller Gefahr beschirmen und von mir alles Herzeleid abwenden. Laß mich auch nicht in Tollkühnheit und Frechheit gerathen, daraus Schaden und Verderben folget, sondern laß mein Herz dich fürchten, daß ich dich vor Augen habe, und gewiß glaube, du leitest mich bei deiner rechten Hand, daß ich nicht falle und umkomme.

\quellseite{177}

Ich weiß, daß des Menschen Thun nichts ist, und er von sich selber nicht getrost sein kann ohne deinen väterlichen Beistand, darum rufe ich dich an um diese hohe Gabe, die ich in meinem Amt und Stand täglich bedarf, damit ich ungescheuet, was recht ist, zu Lob deines Namens verrichte. Herr, du bist mein Fels, meine Burg, darauf ich mich jetzt und an meinem Ende herzlich verlasse. Amen.

\section{223. Gebet auf dem Wachtposten}

Lieber Herr Jesu Christe, ich stehe hier an der Spitze des Lagers und soll die Wache halten wider unsere Feinde. Aber wo du, Herr, uns nicht behütest, so wachet der Wächter umsonst. Ich bitte dich, vergib mir alle meine Sünden, und decke uns mit deiner Gnade, wie mit einem Schilde. Laß deine heiligen Engel um uns sein, die uns schützen und bewahren, daß wir nicht unversehens von unsern Feinden überfallen werden. Laß mich das Grauen der Nacht nicht erschrecken, behüte mich vor den Phantasien und Blendwerk des Teufels. Öffne meine Augen und Ohren, daß ich die Ankunft der Feinde von ferne merke, und mich und das ganze Heer wohl in Acht nehme; bewahre mich in meinem Dienst vor Schlaf und Faulheit. Laß mich mit meinem Herzen beständig zu dir rufen, weil ich an diesem stillen Orte einsam stehe, und du neige dich zu mir mit deiner mächtigen Gegenwart. Sei du bei mir, wenn ich allein bin, und stärke mich an Leib und Seele, daß ich in Frost und Hitze und allem Unwetter meine Kraft behalte, und gesund wieder zu den hellen Haufen meiner Brüder komme; so will ich deinen Namen loben und mich deines Schutzes rühmen, weil ich lebe. Amen.

\section{224. Gebet vor der Schlacht}

Herr, in deiner Gewalt steht alle Kraft und Sieg! Herr, hilf du mir! Herr, die Victorie, der Triumph und Sieg steht in deiner Gewalt; so du sie mir geben willst, will ich dir darum danken; wo du aber unsere Sünden mit einem solchen Schaden und Jammer strafen willst, Herr, so bin ich da, und wills geduldig leiden. Amen.

\quellseite{178}

\section{225. Stoßenser während der Schlacht}

\begin{verse}
Herr, meinen Geist befehl ich dir,\\
Mein Gott, mein Gott, weich nicht von mir,\\
Nimm mich in deine Hände!\\
O wahrer Gott,\\
Aus aller Noth\\
Hilf mir am letzten Ende!\\
Herr Jesu, dir leb’ ich,\\
Herr Jesu, dir sterb ich,\\
Herr Jesu, dein bin ich\\
Todt und lebendig:\\
Mache mich selig. Amen.
\end{verse}

\section{226. Gebet nach siegreicher Schlacht}

Jauchzet dem Herrn, alle Welt, und kommet vor sein Angesicht mit Frohlocken, schauet die Werke des Herrn, was er für ein Zerstören auf Erden angerichtet: er hat dem Kriege gesteuert, der Feinde Spieße zerschlagen, ihre Bogen zerbrochen, und ihre Wagen mit Feuer verbrannt. O barmherziger Vater, wie können wir dir genugsam vergelten die große Wohlthat, daß du uns, die wir nicht würdig sind, unsere Augen zu dir im Himmel aufzuheben, so gnädig angesehen, und nicht mit uns nach unsern Sünden gehandelt, und uns nicht vergolten nach unserer Missethat, sondern unser Gebet in Gnaden erhöret und gethan hast, was wir von dir begehret und gebeten haben. O Herr Jehovah, sei noch weiter mit uns, o Gott Jakob’s, sei unser Schutz. Erhalte und bewahre fernerhin deine liebe Kirche und Christenheit auf Erden. Behüte sie wie einen Augapfel im Auge und beschirme sie unter dem Schatten deiner Flügel, und gib uns Friede zu unsern Zeiten, denn es ist doch kein anderer nicht, der für uns könnte streiten, denn du, unser Gott, alleine. Dir gebühret Ehre und Preis, Macht und Herrlichkeit von nun an bis in Ewigkeit. Amen.

\section{227. Gebet eines Verwundeten}

Ach, du getreuer Gott! Hier liege ich mit Schwachheit umfangen; niemand aber kann mir helfen, als du allein. Ich habe
\quellseite{179}
wohl mehr, als dieses väterliche Mitlein mit meinen ungeheuren Sünden verdient; aber du, HErr, züchtigest mich in Gnaden, wie ein Vater sein Kind. Ist es nun dein väterlicher Wille, ist es dir gefällig und mir seliglich, so kuriere und heile mich, denn du bist mein bester Arzt, du kannst helfen; Tod, Leben, Schwachheit und Gesundheit stehet in deinen Händen, du weißt allein, was mir am Besten ist: wie dir’s gefällt, so mache es; willst du mich länger am Leben haben, daß ich dir und meinem Nächsten in diesem Stande ferner nützlich und dienstlich sein kann, so kannst du mit Einem Worte mich bald gesund machen. Willst du mich aber durch diese Verwundung hinwegfordern, HErr, ich bin bereit und willig, auch in die selig zu sterben. Amen.

\section{228. Gebet in großen Schmerzen}

Ach lieber Vater im Himmel! Wie große Schmerzen, Angst und Noth habe ich an meinem Leibe, die ich leider und noch mehr mit meinen Sünden wohl verdienet habe. Lieber Vater, schaue an, wie meine Krankheit frißt, tobt, wüthet, peiniget und hitzt, und sogar anhält, und laß dichs erbarmen; denn ich bin ja dein liebes Kind und ein Bruder deines geliebten Sohnes JEsu Christi. Thut’s doch einen irdischen Vater weh und geht ihm durchs Herz, wenn er sein Kind stehet Noth und Schmerzen leiden, das wollest du vielmehr an mir thun, und aus solcher Erbarmung mir zu Hülfe kommen in dieser großen Leibesnoth. Lind’re meine Schmerzen, und wenn dir’s gefällt, so nimm sie gar weg. Ach HErr, wie? willst du mein so gar vergessen? Eile mir zu helfen, ich verzage sonst in dieser Angst. Hast du mich doch heißen, dich um Hilfe anzurufen in der Noth, und versprochen, daß du daraus erretten willst, und gesagt, weil ich deiner begehre, so wollest du mir aushelfen. O du gütiger Gott! so komme mir zu Hilfe, und wo ich ja mein Leben mit solchem Schmerz und großer Noth beschließen soll, so wende doch dein väterliches, freundliches und gnädiges Angesicht nicht von mir, sondern hilf mir durch den Heiligen Geist in solcher meiner Noth, daß ich das unaussprechliche große Leiden und die große Qual meines HErrn Christi vor Augen habe, und ihm nach mein Leiden mit Geduld trage, damit ich nicht solcher Leibesnoth halber
\quellseite{180}
mich verkündige, sondern beständig bleibe im Glauben und Gebet bis an mein Ende, wie die lieben Märtyrer, um Jesu Christi, deines lieben Sohnes, willen. Amen.

\section{229. Gebet vor der Operation}

Mein Herr und mein Gott! Zu wem soll ich fliehen in dieser großen Noth, die mich nach deinem heiligen Willen betroffen hat, als zu dir, der du befohlen hast: „Rufe mich an in der Noth, so will ich dich erretten und du sollst mich preisen!“ Wo soll ich Hülfe suchen in meinem Elend, als bei dir, der du alle Hülfe thust, so auf Erden geschiehet? Du hast mich zerrissen, du kannst mich auch wieder heilen. Deine Hand hat mich geschlagen, die wird mich auch verbinden. Du hast mich auf das Schmerzensbette niedergelegt und in Gefahr des Todes gerathen lassen, du kannst mich auch wieder aufrichten, daß ich vor dir leben möge.

Ich bekenne vor deinem heil. Angesicht, daß ich mit meinen unzähligen, schweren Sünden nicht nur diese deine Vater-Ruthe, womit du mich bisher gezüchtiget, sondern auch die Strafe der ewigen Verdammniß, die Marter und Pein der Höllen verdient habe. Allein es ist mir von Grund des Herzens leid, daß ich deine hohe göttliche Majestät erzürnet und beleidiget; meine Seele ängstet sich über die Menge und Größe meiner Mißhandlungen, meine Augen fließen mit Thränen, daß ich leider! dein Gesetz nicht gehalten habe. Doch ich weiß, daß du, o barmherziger Vater, nicht Gefallen hast am Tode des Gottlosen, sondern willst, daß der Sünder sich bekehre von seinem bösen Wesen und lebe. Darum kehre ich mich auch jetzt zu dir mit geängstetem Geiste, und bitte zuvörderst um gnädige Vergebung aller meiner Sünden. Ach wende dich doch, Herr, zu mir, und sei mir gnädig, denn ich bin einsam und elend, die Angst meines Herzens ist groß, führe mich aus meinen Nöthen. Zu dir flieh’ ich, verstoß’ mich nicht, wie ich’s wohl hab’ verdienet. Ach Gott, zürn’ nicht, geh’ nicht ins G’richt, dein Sohn hat mich versöhnet.

Ich glaube festiglich, daß, ob jemand sündiget, haben wir einen Fürsprecher vor deinem Thron, Jesum Christum, der gerecht ist, welcher ist die Versöhnung für unsere Sünden. Auf diesen
\quellseite{181}
meinen JEsus, auf sein blutiges Leiden und schmerzliche Wunden, auf seine vollkommene Genugthuung und kräftigste Fürbitte verlaß ich mich in dieser Stunde ganz und gar, und zweifle nicht, ich werde vor dir Gnade finden und Barmherzigkeit erlangen. So bin ich auch im Glauben gewiß, daß ich dein Kind sei und bleibe, weil du den Geist deines Sohnes in mein Herz gesandt hast, der da unablässig rufet: „Abba, lieber Vater!“ und mich bei dir vertritt mit unaussprechlichem Seufzen. Demnach, weil du mein Gott und Vater bist, dein Kind wirst du verlassen nicht, du väterliches Herze! Ich bin ein armer Erdenkloß, auf Erden hab’ ich keinen Trost. Erbarme dich über mich armseligen und elenden Menschen, wie sich ein Vater über seine Kinder erbarmt, und laß dein Herz über meiner Noth brechen.

Du, HErr, bist der Nothhelfer aller derer, die auf dich trauen, und hast deine Hülfe versprochen Allen, die dich darum ersuchen und aufschreien: nun so sei jetzt auch mein Helfer. Ach HErr, sei nicht ferne von mir, denn Angst ist nahe, eile mir beizustehen, HErr meine Hülfe! Laß dir’s gefallen, daß du mich errettest. Mein Gott, verzeuch nicht! Du gibst den Müden Kraft und Stärke genug den Unvermögenden. Die auf dich harren, kriegen neue Kraft. Ach siehe doch, HErr, wie müde ich vom Seufzen und Schmerzen bin, wie die Krankheit mich entkräftet hat! Meine Hoffnung und Vertrauen steht allein zu dir, darum laß deine Kraft in meiner Schwachheit mächtig erscheinen. Sprich meiner Seelen tröstlich zu: „Fürchte dich nicht, ich bin mit dir, weiche nicht, ich bin dein Gott. Ich stärke dich, ich helfe dir auch durch die rechte Hand meiner Gerechtigkeit.“

O getreuer Arzt Israels, gib Gnade und Segen vom Himmel zu den ordentlichen Mitteln, welche ich in deinem Namen zur Wiedererlangung meiner Gesundheit gebrauchen will. Verleihe dem Arzt zu dem schweren und schmerzhaften Werke Weisheit und Verstand, und führe und regiere mit deiner Gnadenhand seine Hand, damit alles glücklich und wohl abgehe. Mich aber begnade, o du Gott der Geduld und des Trostes, mit einem rechten Heldenmuth und christlicher Geduld, daß ich die so heftigen Schmerzen getrost überwinden und ausstehen möge.

\quellseite{182}

O Herr, hilf, o Herr, laß alles wohl gelingen um der großen Schmerzen und bittern Todesangst meines Erlösers, Jesu Christ, willen! Laß mich in deiner väterlichen Liebe, in den blutigen Wunden deines Sohnes und unter den Gnadenflügeln des Heiligen Geistes Trost, Ruhe und Erquickung finden, bis die Angst und der Schmerz vorübergehe. Sollte ich aber bei der Heftigkeit der Schmerzen aus menschlicher Schwachheit mit einem Fehler übereilet werden, und wider dich, o großer Gott, mit meinen Lippen mich versündigen, ach, so rechne solches mir nicht zur Verdammniß zu, und wende darum dein Gnadenantlitz und hilfreiche Hand nicht von mir ab, sondern nimm auch wegen dieser Sünden meines Jesu Blut zur Versöhnung an. Und so ich nach deinem Rath in bevorstehender großer Noth mein Leben enden, und den Geist aufgeben soll, so binde meine Seele ins Bündlein der Lebendigen und Auserwählten, daß sie nicht umkomme und verloren werde. Erhalte mich in wahrem Glauben an meinen Erlöser und Seligmacher Jesum Christum, daß ich an nichts, als an sein Kreuz und schmerzlichen Tod in meiner letzten Noth gedenke, und durch seine offnen Wunden zu dir in das himmlische Paradies eingehe. Laß deinen Heiligen Geist das Zeugniß geben meinem Geist, daß ich ein Kind Gottes und der Himmel mein Erbtheil sei. Mein letztes Wort müsse sein, mit welchem dein liebster Sohn am Kreuz seinen Mund geschlossen: „Vater, in deine Hände befehle ich meinen Geist!“ Und wenn ich nicht mehr reden kann, so nimm, o Herr, den letzten Seufzer an durch Jesum Christum. Amen.

\section{230. Gebet in Krankheit}

O du barmherziger, gnädiger Gott und Vater! Ich klage und bekenne dir, daß ich bis daher mein Leben, so du mir verliehen hast, übel habe zugebracht und nicht zu deiner Ehre, sondern zur Vollbringung meines Fleisches Lüsten, ich habe mehr mir selber und der Welt gelebt, denn dir. Ach Herr, wie herzlich leid ist’s mir, daß ich die kurze Zeit meines Lebens so übel habe angelegt; darum erkenne ich nun, daß du mich um meiner Sünden willen heimsuchest, und diese Krankheit mir zuschickest, auf daß du mich zur Buße rufest. Du richtest mich billig, Herr, und züchti-

\quellseite{183}

\subsection{Berufsgebete}

gest mich, auf daß ich nicht mit der gottlosen Welt verdammet werde. Du erinnerst mich durch diese Krankheit meiner Sterblichkeit, und willst mich lehren, daß mein Leben ein Ziel hat, und ich davon muß. Ach, wie gar nichts sind doch alle Menschen, die so sicher leben; sie gehen dahin und sind wie ein Schatten; sie sammeln und sorgen und wissen nicht, wer es kriegen wird. Und nun, HErr, weß soll ich mich trösten? Errette mich von meinen Sünden. Ach HErr, ich weiß, daß mein Leben in deiner Hand stehet, du hast alle meine Tage auf dein Buch geschrieben, die noch kommen sollen, als deren noch keiner da war. Es war dir mein Gebein nicht verborgen, da ich gebildet ward im Mutterleibe. Weil es nun dein Wille ist, daß ich krank sein soll, so laß mich diesen deinen Willen gerne thun. Ach HErr, deine Pfeile stecken in mir, und deine Hand drücket mich, es ist nichts Gesundes an meinem Leibe für deinem Dräuen, und ist kein Friede in meinem Gebeine für meiner Sünde, denn meine Sünden gehen über mein Haupt, und wie eine schwere Last sind sie mir zu schwer geworden. Ach mein Gott, bewahre meine Seele für allen Anfechtungen, laß diese meine leibliche Krankheit sein eine Arznei meiner Seele, daß meine Seele dadurch gesund werde. Bewahre ja mein Herz und Sinn, und durch den Frieden meines HErrn Jesu Christi bedecke mich mit deinem Schirm und Schild wider alle feurigen Pfeile des Teufels. Du weißest ja, HErr, daß ich deiner Hände Werk bin, daß du mich aus Leimen gemacht hast; wir sind Thon, du unser Töpfer. Ach wie bin ich zerbrochen! Darum schone mein, o lieber HErr; du wirst ja nicht verwerfen und im Zorn zerbrechen das Werk deiner Hände, du hast mich ja aus Mutterleibe gezogen, und bist meine Zuversicht gewesen, da ich noch an meiner Mutter Brüsten war; du bist ein Ursprung und HErr meines Lebens und hast Macht, es wieder zu nehmen, wenn du willst. Ist es nun dein Wille, so thue die Barmherzigkeit an mir, wie an dem Könige Hißkia, welchem du fünfzehn Jahre zu seinem Leben zusetztest. Siehe! um Trost ist mir sehr bange, nimm dich meiner Seelen herzlich an, daß sie nicht verderbe, und wirf alle Sünden hinter dich zurücke; ist mir diese Krankheit nicht zum Tode, so richte mich wieder auf und heile mich, HErr, denn ich bin sehr schwach. O du
\quellseite{184}
himmlischer Arzt, HErr JEsu Christe, nimm mich in deine Kur, lindere meine Schmerzen, und heile mich durch deine Wunden, laß dein bitter Leiden und Sterben meine Arznei sein. HErr, laß mich deines Kleides Saum anrühren, das ist dein Wort und Sacrament, darein du dich verkleidet, und heile mich, daß ich durch deine Kraft, die von dir ausgehet, gesund werde. Ist es aber dein Wille, daß ich durch diese Krankheit hingehen soll den Weg aller Welt, wohlplan, so bin ich nicht besser denn meine Väter; laß nur, HErr, deinen Diener in Frieden fahren, wie du gesagt hast, denn meine Augen haben dich, meinen Heiland, gesehen, welcher bereitet ist für alle Völker, ein Licht, zu erleuchten die Heiden und zum Preis deines Volks Israel. Amen.

\section[231. Dank- und Bußgebet nach Errettung aus …]{231. Dank- und Bußgebet nach Errettung aus mancherlei Gefahr}

Du Vater der Barmherzigkeit und Gott alles Trostes, der du mein Gott bist von Mutterleibe an, und mich die Zeit meines Lebens unter dem Schatten deiner Flügel wunderbarlich beschirmet und beschützet hast, dich lobe ich jetzt in der Stille, und preise deinen Namen ewiglich. Nimmermehr kann ich dir vergelten alle deine Wohltat, die du an mir gethan hast, daß du mich aus mancher Gefahr errettet, und täglich viel Gutes an Leib und Seele erwiesen hast. Denn während sonst im Kriege, zu welchem du mich verordnet hast, das Schwert bald diesen, bald jenen getroffen, und die Schützen von den Mauern und im Felde geschossen, und viel tapfere Soldaten getödtet haben, also daß ihrer viel vor mir und neben mir zu beiden Seiten gefallen sind, auch oftmals zwischen mir selbst und dem Tode nur ein Schritt gewesen ist, so hast doch mich deine Hand bedecket, daß ich leben geblieben, und mein Leben, wie eine Beute, davon gebracht habe. Du allein bist es, der mich ausgerüstet mit der Kraft, und meine Hände streiten gelehret, auch Stärke genug gegeben hast, daß ich öfters mit dir über die Mauern der Gefahr springen, und manchem Unglück habe entrinnen können, also daß ich noch diese Stunde lebe, und von dem wunderbahrlichen Sieg in den Hütten der Gerechten singen kann, daß es die Elenden hören und sich freuen. Und darum danke ich dir, o Gott, billig,
\quellseite{185}
daß du mich aus so vieler und großer Noth erlöset hast, und lobe den Namen deiner Herrlichkeit, daß du bisher gewesen bist meines Lebens Kraft und mein Heil. Weil ich aber nicht würdiglich gewandelt habe meinem Beruf, noch dich allezeit für Augen und im Herzen gehabt, sondern manche Sünde begangen, davon ich nun ein Brandmal im Gewissen habe: so stehe ich dich kindlich an, daß du mir wollest gnädig sein nach deiner Güte und meine Sünden tilgen nach deiner großen Barmherzigkeit. Wasche mich durch das Blut Jesu Christi, welches uns rein macht von allen unsern Sünden, wohl von meiner Missethat, und reinige mich von meiner Sünde; denn ich erkenne meine Missethat und meine Sünde ist immer vor mir. Und wo ich zu unbarmherzig gewesen, so laß doch um Christi willen nicht wieder ein unbarmherzig Gericht über mich ergehen, sondern errette mich von diesen meinen Blutschulden, Gott, der du mein Gott und Heiland bist, daß meine Zunge deine Barmherzigkeit rühme hier zeitlich und dort ewiglich. Amen.

\section[232. Gebet christlicher Soldaten in …]{232. Gebet christlicher Soldaten in Gemeinschafft zu sprechen}

HErr Himmels und der Erden, du großer und starker Gott, wir danken dir von Herzen durch Jesum Christum, daß du uns bisher so gnädiglich behütet und erhalten hast. Wir bitten dich, getreuer himmlischer Vater, du wollest uns armen Sündern unsre schweren und vielfältigen Sünden um Jesu Christi willen gnädig vergeben. Bekehre und regiere unsre Herzen durch deinen Geist, daß wir dem bösen Exempel des großen Haufens nicht nachfolgen, nicht heidnisch und gottlos, sondern christlich und gottseelig leben, dich lieben und fürchten, und allenthalben ein gut Gewissen zu behalten uns befleißigen. Und verleihe uns allen auch Weißheit, Stärke, Gesundheit, Standhaftigkeit und einen Heldenmuth, ritterlich wider die Feinde zu kämpfen für das Vaterland. Und da menschliche Macht, Hülff, Rüstung und Anschläge ohne deinen Beistand durchaus nichts vermögen, so streite du für uns, HErr, du starker Gott, und segne unsern Ausgang und Eingang von nun an bis in Ewigkeit. Beug für uns her und treib unsre Feinde zurück, gib ihnen verzagte Herzen und matte Hände, damit unser armes Vaterland errettet, deine liebe Kirche erhalten, die Betrübten getröstet
\quellseite{186}
, und die, so niedergeschlagen und unterdrückt sind, wieder aufgerichtet werden mögen. Wir zwar, o HErr, sind ja freilich unwürdig an uns selbst solcher deiner Hilfe, denn wir haben alle schwerlich gesündigt; aber, ach HErr, siehe an das heilige theure Verdienst Jesu Christi, deines Sohnes, siehe an den Jammer und das Elend deiner lieben Kirche, und erhöre das Seufzen so vieler armen, betrübten, verderbten, nothleidenden, frommen Herzen, die ohne Unterlaß mit heißen Thränen im kindlichen Vertrauen zu dir um Hülfe und Rettung schreien. Um derentwillen, HErr, stehe uns bei, und verleihe uns Glück und Sieg, damit deine Ehre kund werde in aller Welt, und alle Menschen erkennen, daß du der HErr bist, der rechte einige Helfer. In der Noth, der den Kriegen wehret, Spieß und Bogen zerbricht, alle Anschläge wendet, wie er will, und der endlich den Seinigen zu Hilfe kommt, auch wenn man sich's am wenigsten versiehet. Laß dir auch, o HErr, in deinen Schutz befohlen sein alle christlichen Regenten, desgleichen unsre Kriegsobersten, ja das ganze Heer. Geleite und beschütze uns alle durch deine heiligen Engel, und welchen du ihre Zeit bestimmt hast in diesem unsern Kriegsheer zu sterben, denen wollest du beständigen Glauben und wahre Buße gnädiglich verleihen; sie auch an ihrem letzten Ende trösten durch deinen Heiligen Geist und ihre Seelen zu dir nehmen in die ewige Freude und Seligkeit um Jesu Christi, deines lieben Sohnes willen, welcher uns in seinem Namen befohlen zu beten: Vater unser! der du bist im Himmel, geheiliget werde dein Name, dein Reich komme, dein Wille geschehe, wie im Himmel, also auch auf Erden; unser tägliches Brod gib uns heute; und vergib uns unsere Schuld, als wir vergeben unsern Schuldigern, und führe uns nicht in Versuchung, sondern erlöse uns von dem Uebel, denn dein ist das Reich und die Kraft und die Herrlichkeit in Ewigkeit. Amen.

\section[233. Gebet eines Soldaten für seine …]{233. Gebet eines Soldaten für seine Angehörigen zu Hause}

Vielgütiger HErr und einziger Beschirmer aller Frommen! Ich danke dir von ganzem Herzen, daß du mich bisher in allen Gefahren meines Berufes gnädig behütet und erhalten hast, und bitte dich, wollest mich auch fernerhin in deinen väterlichen Schutz
\quellseite{187}
nehmen. Auch sehe ich an deine unendliche Güte und Barmherzigkeit, du wollest mein Haus und alles, was mir zuständig ist, ringsumher bewahren, damit kein Unglück zu meiner Hütte sich nahe, und mein armes Vermögen, noch was ich zurückgelassen habe, nicht beschädigt, noch mir genommen werde. Insbesondere aber befehle ich dir in deine Hände, was mir am liebsten ist, nämlich meine liebe Ehegattin und Kinder, die mir mein Herz beugen und mein größter Reichthum sind. Sei du selbst ihre Hilfe und ihr Schutz, daß sie nicht fallen, und wende ab von ihnen alle Leibes- und Seelen-Gefahr. Bewahre sie heute und allezeit, daß sie deine Majestät fürchten, und nicht einwilligen in vorsätzliche Sünden, noch thun wider deine Gebote. Befiehl deinem Engel, daß er sie behüte auf allen ihren Wegen, und tröste sie selbst in aller ihrer Mühe und Arbeit. Gib, daß sie anhalten am Gebet, und ausüben alle möglichen Werke der Gottesfurcht. Und so du willst, so bringe mich nach glücklich vollendeter Kriegsgefahrt gesund wieder heim, und gib, daß ich auch die Meinigen gesund finde, so wollen wir dann mit einmüthigem Munde und Herzen dich loben und preisen für alles das Gute, das du uns an Seele und Leib erzeiget hast. Amen.

\section[234. Gebet der Angehörigen zu Hause für den …]{234. Gebet der Angehörigen zu Hause für den Soldaten im Felde}

O allerliebster Gott und Vater, wir sind schuldig, nicht nur für uns, sondern auch für die Unsrigen, ja für alle Menschen zu bitten. So bitte ich nun von Grund meines Herzens für meinen lieben Eheherrn (Vater, Sohn, Bruder, Freund), welcher in deinem Namen vermöge seines Berufs in den Krieg gezogen ist. Weil ich nun diesen Tag und Zeit nicht kann bei ihm sein, so gern ich auch wollte: ach Herr, der du im Himmel und auf Erden bist, so sei und bleibe du bei ihm heute und allezeit, schütze und beschirme ihn an Leib und Seele, erhalte ihn bei guter Gesundheit, bewahre ihn vor allem Uebel, versorge ihn als ein Vater, führe ihn als ein treuer Geleitsmann, schütze ihn in aller Gefahr, segne ihn in allen seinen Thun und Geschäfte. Und dieweil ja der Satan durchwandelt alle Lande, und sonderlich deiner
\quellseite{188}
frommen Kinder Gange Netze stellet, daß er sie in Unglück bringen möge: ach HErr, du starker Gott! so bewahre meinen lieben Mann (Vater, Sohn, Bruder, Freund) vor allen listigen Anläufen, Wegen und Stegen des leidigen Satans und seines verfluchten Anhangs. O HErr der Heerschaaren! befehle deinen heiligen Engeln, daß sie sich um ihn lagern, und ihn auf allen seinen Wegen behüten. Regiere und führe ihn durch deinen heiligen Geist heute diesen Tag und allezeit, bringe ihn auch wieder glücklich zu uns, daß wir ihn unversehret und unbeschädiget wieder sehen mögen. So wollen wir, treuer Gott, dich sämmtlich mit Freuden loben unser Leben lang, und dich dort preisen ewiglich. Amen.

\section{235. Gebet bei Todesfurcht}

Lieber HErr JEsu Christe, ob ich gleich das Gesetz nicht erfülle und obschon noch Sünde vorhanden ist, und ich mich vor dem Tode und der Hölle fürchte, so weiß ich doch dies aus dem Evangelio, daß du mir alle deine Werke geschenkt und gegeben hast: deß bin ich gewiß, du lügst nicht, deine Zusage wirst du wahrhaftig halten, und zum Zeichen habe ich die Taufe darauf empfangen. Darauf verlaß ich mich. Weil du denn, lieber Gott, mein bist, so will ich gerne sterben, denn also gefällt es dir, lieber Vater, und der Tod kann mir nicht schaden: er ist verschlungen in den Sieg, und dir, lieber HErr Gott, sei Dank, daß du uns den Sieg gegeben hast durch unsern HErrn JEsum Christum. Amen.

\section{236. Drei Gebete in Todesnoth}

\subsection{Das erste zu Gott dem Vater}

Allmächtiger Gott! ein Vater der Barmherzigkeit und Gott alles Trostes! erbarme dich meiner, wie sich ein Vater über seine Kinder erbarmt. Gedenke, daß dein lieber Sohn, JEsus Christus, um meinetwillen ist Mensch geworden. Gedenke, daß du die Welt also geliebet hast, daß du deinen eingebornen Sohn hast gegeben, auf daß Alle, die an ihn glauben, nicht verloren werden, sondern das ewige Leben haben. In diesem Glauben kommt meine Seele zu dir und bringt mit sich deinen lieben Sohn, JEsum Christum,
\quellseite{189}
welchen du mir geschenkt hast; um desselben willen wollest du an mich denken und mich zu Gnaden annehmen. Dir befehle ich meine Seele. Laß mich, dein liebes Kind, ewiglich in deiner Gnade bleiben, und thue Barmherzigkeit an mir, weil ich noch lebe und wenn ich sterbe, um deines lieben Sohnes, unseres HErrn JEsu willen. Amen.

\subsection{Das andere in Todesnoth zu Gott dem Sohn}

HErr JEsus Christus, mein einziger Trost und meine einzige Hilfe! Der du meine Krankheit hast getragen und bist um meiner Sünden willen verwundet! Du Lamm Gottes, der du trägst die Sünden der Welt! Sei du meine Gerechtigkeit, meine Heiligung und meine Erlösung. Heile meine Schmerzen durch deine Wunden. Laß dein unschuldiges Blut an mir nicht verloren sein; nimm deinen Heiligen Geist nicht von mir. Verlaß mich nicht in meiner letzten Noth, sondern hilf mir bald hindurch zum ewigen Leben. Gedenke daran, daß ich dein armes Schäflein war und du mein getreuer Hirt. Nimm mich auf deine Arme und trage mich in das freudenreiche Paradies. Laß mich auch dein freundliches Angesicht und deine Herrlichkeit sehen. Amen.

\subsection{Das dritte zu Gott dem Heiligen Geist}

O Gott Heiliger Geist, du einiger Trost aller Betrübten! stärke und tröste mein kleinmüthiges und trauriges Herz. Bewahre meinen Glauben, den du in mir wirkest, und laß das glimmende Töchtlein nicht auslöschen. Behüte mich vor des bösen Feindes Anfechtungen. Laß mein Herz vor dem Tode nicht erschrecken, noch verzagen; wehre aller Ungeduld; mache mich willig, Gott meinem HErrn gehorsam zu sein. Hilf mir mein Leben selig und fröhlich beschließen. Bewahre mein Herz und meinen Sinn in dem Frieden meines HErrn JEsu Christi. Und wenn ich nicht mehr reden kann, so bewahre und versiegele das Gedächtniß des heiligen Namens JEsu Christi in meinem Herzen, und laß mich mit demselben selig einschlafen. Amen.

\quellseite{190}

\section{237. Gebet um Frieden}

\begin{verse}
Verleih uns Frieden gnädiglich,\\
Herr Gott! zu unsern Zeiten,\\
Es ist doch ja kein andrer nicht,\\
Der für uns könnte streiten,\\
Denn du, unser Gott, alleine.
\end{verse}

Gib unserm Land und aller Obrigkeit Fried’ und gut Regiment, daß wir unter ihnen ein christlich, ehrbar, geruhig Leben führen mögen, in aller Gottseligkeit und Wahrheit. Amen.

\section{238. Die kleine Litanei}

Herr Gott, Vater im Himmel, du Schöpfer aller Dinge, erbarme dich über mich und erhöre mich!

Herr Gott Sohn, du Heiland aller Welt, erbarme dich über mich und bitte für mich!

Herr Gott Heiliger Geist, du Tröster in aller Noth, erbarme dich über mich und seufze in mir!

O du heilige, heilige, heilige, hochgelobte Dreieinigkeit, erbarme dich über mich und sei mir armen Sünder gnädig, vergib mir alle meine Sünde, erhalte mir dein reines seligmachendes Wort und schenke mir deinen zeitlichen, geistlichen und ewigen Frieden. Amen!

\chapter{VII. Fürbitten}

\section[239. Gebet christlicher Weiber bei einer Frau …]{239. Gebet christlicher Weiber bei einer Frau in Kindes-Nöthen}

O allmächtiger, starker Gott, wir bitten dich in dieser so großen Noth durch Jesum Christum, deinen lieben Sohn, hilf dieser Kindes-Mutter, die hier liegt in ihrem weiblichen Beruf, nach deinem göttlichen Willen, und erbarme dich deines Geschöpfs und schicke es mit Mutter und Kinde, wie es ihnen an Leib und Seele mag nützlich und dienlich sein; der du keinen Menschen
\quellseite{191}
jemals verlassen hast, welcher zu dir geschrieen, und hast allen Kindes-Müttern endlich geholfen, die in ihren Nöthen ihre Zuflucht und Vertrauen auf dich gesetzt. HErr, du gnädiger Helfer in aller Noth! hie ist Noth, hie ist Angst, hie ist Zeit zu retten und zu helfen: wenn du, HErr, nicht hilfst, so ist es mit aller Welt Hülfe verloren. Denn das Kind ist bis an die Geburt kommen, und ist Niemand, der helfen kann, denn du, HErr unser Gott, allein. Du weißt am besten die rechte Stunde und Augenblick. Wir harren deiner Hülfe und hoffen auf dein Wort von einer Stunde bis zur andern. Und ob es währt bis in die Nacht und wieder an den Morgen, soll doch unser Herz an deiner Macht verzweifeln nicht, noch sorgen. Deine Hand zu helfen hat kein Ziel, wie groß auch sei der Schmerzen, das glauben wir von Herzen. Daran wir keinen Zweifel tragen, dein Wort kann uns nicht betrügen. Denn du weißt wohl, wenn's am besten ist, und brauchst an uns kein arge List, das wollen wir dir vertrauen. Ob sich's anließ, als wolltest du nicht, so lassen wir uns doch nicht erschrecken. Denn wo du bist am besten mit, da willst du es nicht entdecken. Dein Wort lassen wir uns gewisser sein und ob das Herz sprach lauter Nein, so lassen wir uns doch nicht grauen. Ach Gott, stärke uns nur den Glauben und hilf uns um deines Namens Ehre willen. Amen. Lasset uns mit einander beten: Vater unser! der du bist im Himmel, geheiliget werde dein Name, dein Reich komme, dein Wille geschehe, wie im Himmel, also auch auf Erden; unser tägliches Brod gib uns heute, und vergib uns unsere Schuld, als wir vergeben unsern Schuldigern, und führe uns nicht in Versuchung, sondern erlöse uns von dem Uebel, denn dein ist das Reich und die Kraft und die Herrlichkeit in Ewigkeit. Amen.

\section{240. Gebet einer Mutter für ihren Säugling}

Ach mein HErr JEsu, ich danke dir, daß du mich gewürdigt hast, eines Kindes Mutter zu sein, und gibst mir die Gnade, daß ich es mit meinen Brüsten stillen kann! Ich gebe meinem Kinde die Muttermilch ein, flöße du ihm, mein Erlöser, deine Liebe und heilige Furcht ein. Ach laß es ja allezeit ein lebendiges Glied an deinem Leibe bleiben! Dämpfe du selbst durch deinen Heiligen Geist
\quellseite{192}
die Erbsünde und alle Bosheit in ihm. Gib ihm deinen Heiligen Geist, regiere sein Herz und Sinn, bewahre es vor der sündlichen, bösen Welt Aergerniß und Verführung, heilige und segne es an Leib und Seele, laß es täglich zunehmen, mit dem leiblichen Wachsthum, an Weisheit, Alter und Gnade bei dir, o Gott im Himmel, und den Menschen auf Erden. Amen.

\section{241. Gebet eines, der Gevatter stehen soll}

O dreieiniger Gott, ich übergebe dir hiermit diesen meinen Pathen, und bitte, du wollest ihm seine sündliche Unreinigkeit, darinnen er empfangen und geboren, nicht zur Verdammniß zurechnen, sondern solche mit Christ, deines Sohnes, Heiligkeit und Gerechtigkeit zudecken, ihn in die Zahl der Christen aufnehmen und annehmen, in deiner Gnade beständig erhalten, und nach diesem Leben die rechte Seligkeit geben und schenken. Amen.

\section[242. Gebet der Eltern für ihr Kind, das auf …]{242. Gebet der Eltern für ihr Kind, das auf Reisen ist}

O allmächtiger, barmherziger, gütiger Gott und Vater, der du wunderbarlich den jungen Tobias einen langen und fernen Weg durch deinen Engel Raphael begleitet, und vor aller Gefahr bewahret, auch ihm Gnade verliehen, daß er mit allem Glück, Segen und Benedeyung wiederum gesund anheim zu seinen Eltern (die seiner Zukunft mit großem Verlangen erwartet) ist kommen, durch den Namen deines geliebten Sohnes, unseres HErrn Jesu Christi, der da ist der Weg, die Wahrheit und das Leben, ich bitte dich, allmächtiger, getreuer Vater, du wollest meinem lieben Kinde auf seiner Reise Gnade, Glück, Heil, Segen und Benedeyung verleihen, dasselbe vor aller Gefahr, Anstoß und Unglück, beide an Seel und Leib, auf allen seinen Wegen und Stegen, zu Wasser und Land, bewahren, und ihm deinen heiligen Engel als einen glückseligen Geleitsmann senden, der es hin und wieder beschütze und beschirme vor alle dem, so ihm an Seel und Leib, Ehr und Gut mög schädlich sein, und es also mit gesundem Leibe glücklich wiederum zu Hause kommen möge. Dazu helf ihm der göttliche Segen der heiligen Dreifaltigkeit, Gottes des Vaters, des Sohnes und des Heiligen Geistes. Amen.

\quellseite{193}

\section{243. Gebet für einen Sohn, der in der Fremde ist}

Barmherziger Gott und Vater über alles, was Kinder heißt im Himmel und auf Erden! Erhöre gnädig das Gebet und Flehen, das deine arme Magd jetzt vor dir thut. Herr! Du weißt, wie mit viel Thränen ich meinen Sohn von mir gelassen und in die Fremde geschickt, und mit welchem Kummer oftmals mein Herz seinetwegen erfüllet ist. Ich wünsche nichts mehr, als daß er dich fürchte, und Redliches und Nützliches lerne, und im gesunden Leibe eine unverletzte Seele und reines Gewissen behalte. Und darum rufe ich dich (von welchem alle guten und vollkommenen Gaben herab kommen) in tiefster Demuth und von Grund meines Herzens an: Lieber Vater! regiere ihn doch mit deinem Heiligen Geiste, daß er an allen Orten und zu aller Zeit dich vor Augen und dein Wort in seinem Herzen habe, damit er nicht gottlos werde, und in Sünde willige wider deine Gebote. Laß ihn bei dir in Gnaden und Barmherzigkeit sein, halte deine Augen über ihn offen, leite und schütze ihn durch deine Engel, die starken Helden, auf allen seinen Wegen, daß ihm kein Uebel begegne, und kein verderbliches Unglück sich zu ihm nahe.

Herr Jesu, du Heiland derer, die dir vertrauen! der du von deinem himmlischen Vater auf die Erde gesendet worden, den Menschenkindern zum Besten. Laß meinen Sohn deiner Fürsorge, Aufsicht, Liebe und Huld befohlen sein. Gib, daß er deinen Fußstapfen in allem, wie du ihm ein Vorbild gelassen hast, nachfolge, und von Tag zu Tag zunehme, nicht allein an Alter, sondern auch an Gottseeligkeit und Tugend, an Weisheit und Geschicklichkeit, (im Handel und Wandel etc.), an Gnade bei Gott und den Menschen. Segne und benedeie sein Thun und Vornehmen, und laß wohlgerathen, was er in deinem Namen anfängt. Bringe ihn auch endlich gesund und mit Freuden wieder heim, daß er einst nach deinem Willen der Stab und Trost meines Alters sein, vornehmlich aber ein Erbe des Himmels und der ewigen Seligkeit durch dein Verdienst bleiben möge.

Gott heiliger Geist, du Regierer und Führer der Jugend! lehre diesen Jüngling seinen Weg unsträflich gehen, und leite ihn auf ebener Bahn nach deinen heiligen Geboten zu leben. Gib, daß
\quellseite{194}
er mit aufrichtigem Herzen vor dir wandle, und gewiß glaube, daß die Gottseeligkeit zu allen Dingen nütze sei, und er sich ohne Unterlaß übe, zu haben ein unverletztes Gewissen, beide gegen Gott und Menschen. Neige zu ihm die Herzen vornehmer, redlicher und frommer Leute, daß sie sich seiner annehmen, wo er ihrer bedarf, und ihm alles Gutes erweisen. Bewahre ihn aber auch vor des Satans Trug und List, und der Welt ärgerlichem Wesen, vor falscher Lehre und verführerischer Gesellschaft, vor Sünden und Schanden, Räubern und Mördern, vor Schaden und Gefahr, und einem bösen und schnellen Tod. Mache aus ihm einen solchen Menschen, der deiner Ehre dienen, dem Nächsten nützlich sein und dich für diese und alle deine Gnadenwohlthaten immer und ewig mit Herz und Mund preisen möge. Amen.

\section{244. Gebet für ein krankes Kind}

Allmächtiger Gott, lieber Vater! du hast nicht allein zu dem Abraham gesprochen: ich bin dein Gott und deines Samens nach dir; sondern auch uns durch deinen Apostel versichert, daß uns und unsern Kindern diese Verheißung geschehen sei. Darum zürne nicht, daß ich mich unterwinde, jetzt mit dir zu reden, und mein Gebet und Seufzen für mein krankes Kind in deinen Schooß zu schütten. Herr! es ist dein Geschöpf und deiner Hände Werk; du hast es im Mutterleibe bereitet, mit Leib und Seele begabet, und Leben und Wohlthat an ihm gethan. Du hast es an das Licht dieser Welt gebracht und in der heiligen Taufe zu deinem Kinde angenommen, auch bis hieher in deiner Gnade erhalten. Nun, so laß dir demnach auch jetzt dasselbe in seiner größten Schwachheit befohlen sein. Ach! rechne ihm nicht zu die Sünde, in welcher es empfangen und geboren ist, und womit es in seinem Leben deine heiligsten Gebote übertreten hat. Erbarme dich sein nach deiner großen Güte, lindere ihm seine Schmerzen und Krankheit, und so du willst, so mache es bald wieder gesund. Denn, Herr! in deiner Hand und Macht stehet, demselben die völlige Gesundheit wieder zu erstatten, und das wirst du auch thun, wenn du findest, daß es zu deiner Ehre, meiner Wohlfahrt und seiner eigenen Seligkeit nützlich ist. Willst du aber nach deinem
\quellseite{195}
Heiligen Rath anders mit ihm verfahren, daß es entweder länger auf dem Siechbette liege, oder frühzeitig sterben soll, so geschehe nicht mein, sondern dein Wille, der allzeit der beste ist. HErr JESU, du theures Gottes-Lamm, das der Welt Sünde trägt! Siehe an dieses arme Lamm, welches du mit dem allerkostbarsten Lösegeld zu deinem Eigenthum erkauft hast. Wasche es und reinige es mit deinem unschuldigen Blut von allen seinen Sünden, stärke es durch deinen Heiligen Geist in wahrem Glauben an dich und dein Verdienst, in christlicher Geduld unter dem Kreuz, das du ihm aufgelegt, und in der guten Hoffnung deiner gnädigen Hülfe. Ist es dir gefällig, daß ich dieses Kind in deiner Furcht auferziehen und sehen soll, wie es durch deinen Segen ferner von Jahr zu Jahr in allen christlichen Tugenden zunehme, so will ich deinen Namen dafür preisen immer und ewig. HErr! sprich nur ein Wort, so wird es gesund werden und leben! Hast du aber beschlossen, es durch einen seligen Tod zu dir zu nehmen (wie du denn mehr Recht und Macht über dasselbe hast, als ich), so bin ich auch damit zufrieden, und weiß, du wirst seine Seele zu den Geistern der vollkommenen Gerechten in den Himmel durch deine heiligen Engel einführen; den kranken Leib aber an jenem Tage in der Auferstehung der Gerechten, mit seliger Unsterblichkeit bekleiden und deinem verklärten Leibe ähnlich machen.

Gott, Heiliger Geist, du Geist der Kindschaft und Gott alles Trostes! erfülle meine traurige Seele und meines kranken Kindes Herz mit deinem kräftigen Trost. Hilf durch deine himmlische Kraft seiner Schwachheit auf, und vertritt es bei Gott mit unaussprechlichem Seufzen. Verleihe auch mir die Gnade, daß ich mich dem Willen Gottes gehorsam unterwerfe, unter deine gewaltige Hand demüthige, mit stillem Geist des Ausgangs dieser Krankheit erwarte, und nichts mehr sage, als dieses: Du bist der HErr, thue, was dir wohlgefällt. Ich habe bisher in der Gottseligkeit und guten Art dieses Kindes gesehen, daß es dir zugehöre; dir soll es demnach auch ergeben und zugeeignet sein im Leben und im Sterben. Wird es durch den Tod aus meinem Schooß und Armen genommen, so versehe du es in den Schooß deiner Gnaden, und wirst es im Himmel sättigen von den Brüsten deines Trostes,
\quellseite{196}
daß es genieße Freude die Fülle und liebliches Wesen zu Gottes Rechten ewiglich.

Amen! Herr, du dreieiniger Gott! Sprich selbst zu meinem Gebet und Flehen: Amen. Ja, ja, es soll also geschehen. Amen!

\section[245. Fürbitte kleiner Kinder für ihren auf …]{245. Fürbitte kleiner Kinder für ihren auf der Reise begriffenen Vater}

Ewiger Sohn Gottes, o treuer Heiland aller, die dich im Glauben anrufen, wir deine Kindlein, die auf dein Blut getauft und mit deinem Geist zu königlichen Priesterlein geweihet und deine Brüder und Miterben aus Gnaden erwählet sind, schreien mit unsern unschuldigen Zünglein zu dir, bitten herzlich, du wollest unsern lieben Vater, der seines Berufs und Nahrung halber über Land gereiset, auf seinen Wegen gnädig behüten und ihn seine Sachen wohl und schleunig ausrichten lassen, und ihn gesund und mit Freuden wieder zu uns bringen, der du deinen Diener Jacob auch mit deinen Engelein auf seiner Reise behütet und begleitet hast, du ewiger und treuer Geleitsmann und Gefährte aller deren, die dich fürchten, und auf deine Güte warten. Amen.

\section{246. Gebet frommer Kinder für ihre Eltern}

Ach gnädiger, barmherziger Gott, lieber Vater, der du bist der rechte Vater über alles, das Kinder heißet im Himmel und auf Erden: Ich danke dir herzlich, daß du mir meine liebe Eltern, Vater und Mutter gegeben, und bis daher in guter Gesundheit und Wohlstand erhalten hast, dir sei Lob, Ehr und Dank für diese deine große Wohlthat, und bitte dich, du wollest mir meinen Ungehorsam, damit ich mich gegen meine liebe Eltern oft versündigt habe, aus Gnaden vergeben, und die Strafe von mir abwenden, die du im vierten Gebot dräuest. Gib mir aber ein gehorsam und dankbar Herz gegen sie, daß ich sie ehre, fürchte, liebe, mit meinem Gehorsam und deiner göttlichen Furcht erfreue, daß ich sie für Gottes Ordnung erkenne, und ihre väterliche, wohlmeinende Strafen ge-duldig annehme. Lehre mich auch bedenken, wie sauer ich meiner Mutter worden bin, und mit was großer Mühe und Arbeit sie mich erzogen, laß mich dieselbe wieder ehren mit Gehorsam, Liebe, Demuth
\quellseite{197}
, Furcht in Worten und Werken, auf daß ich den Segen und nicht den Fluch ererbe, sondern ein langes Leben. Laß das Exempel des Gehorsams meines Herrn Jesu Christi immer für meinen Augen stehen, welcher seinem Vater gehorsam gewesen bis zum Tode am Kreuz. Gib mir den Gehorsam Isaaks, die Furcht Jakobs, die Zucht Josephs, die Gottesfurcht des jungen Tobia, und gib meinen lieben Eltern den Glauben Abrahams, den Segen Isaaks, den Schutz Jakobs, die Glückseligkeit Josephs, und die Barmherzigkeit des alten Tobia. Laß sie in einem feinen geruhigen Leben, in Fried und Einigkeit alt werden, lindere ihnen ihr Kreuz, und hilfs ihnen tragen, erhöre ihr Gebet, und segne ihre Nahrung, behüte sie für allem Uebel Leibes und der Seelen, und wenn ihre Zeit vorhanden ist, so laß sie sanft und stille einschlafen, und nimm sie zu dir ins ewige Vaterland durch Jesum Christum. Amen.

\section[247. Gebet der Kinder für ihren Vater, der …]{247. Gebet der Kinder für ihren Vater, der auf Reisen ist}

Allmächtiger, barmherziger, gnädiger Gott und Vater, der du uns nicht allein erschaffen und gemacht hast, sondern auch noch täglich erhältst und alles gibst, was wir bedürfen an Leib und Seel, ich bitte dich von Herzen, du wollest dich meines Vaters, welcher jetzund, anderweit seine Geschäfte zu verrichten, verreiset ist, gnädig annehmen, ihn erhalten und bewahren für aller Gefahr und Schaden, durch deine heiligen Engel begleiten und ihn nach verrichteter Sache glücklich wiederum zu Haus bringen. Ach HErr, erhöre mich! Erhöre, o HErr! mein Gebet und Seufzen und sei selber sein Wegweiser und Geleitsmann und gib Glück und Segen zu seinem Vorhaben. Erhalte ihn bei gutem Wohlstand und Gesundheit, auf daß wir uns bald wiederum sehen und in dir miteinander ergötzen mögen um deines Sohnes willen. Amen.

\section[248. Gebet frommer Kinder für ihren kranken …]{248. Gebet frommer Kinder für ihren kranken Vater}

Allmächtiger, ewiger Gott, Vater unsers HErrn Jesu Christi! du sprichst in deinem heiligen, göttlichen Worte: Rufe mich an in der Zeit der Noth, so will ich dich erretten, und du sollst mich preisen. Dieweil denn uns arme Kinder jetzt große Noth betroffen hat, daß du unsern lieben Vater (Mutter) mit großer
\quellseite{198}
Leibesschwachheit, um unserer Sünde willen, hast heimgesucht; so kommen wir mit demüthigen Herzen vor deine göttliche Majestät und rufen dich an im Namen deines lieben Sohnes, Jesu Christi, unsers lieben HErrn, Seligmachers und Erlösers, strafe uns nicht in deinem Zorn, und züchtige uns nicht in deinem Grimm, sei uns gnädig, HErr, sei uns gnädig, und lindere das schwere Kreuz, so du unserm lieben Vater (Mutter) zugeschicket hast. Tröste ihn (sie) durch deinen Heiligen Geist, verleihe ihm (ihr) wiederum Gesundheit des Leibes, Stärke und Kraft, und entzeuch ihn (sie) uns noch unerzogenen Kindern ja nicht so bald; sondern lasse ihn (sie) nach deinem göttlichen gnädigen Willen noch länger bei uns auf dieser Erde leben, damit er (sie) unserer lieben Mutter (Vater) zu Trost und Hilfe, uns in deiner Furcht und wahrer Erkenntniß auferziehen möge, und wir mit der Zeit deiner christlichen Kirche und der Gemeinde dienen können, dich, unsern Gott loben, ehren, rühmen und preisen. Gib, daß unser lieber Vater und unsere liebe Mutter, welchem (welcher) du wieder gnädiglich zu seiner (ihrer) Gesundheit helfen wollest, Trost und Freude an uns haben, und wir ihrer wiederum in ihrem Alter können pflegen und warten, und ihnen alles Liebes und Gutes erzeigen, wie sie uns in ihrer Jugend bewiesen haben. Ach HErr, sei uns gnädig, rechne uns nicht zu die Sünde und Missethat unserer Jugend, damit wir dich erzürnet und solche Strafen wohl verdienet haben; sondern stehe uns an mit den Augen deiner Barmherzigkeit, in deinem lieben Sohne Jesu Christo, und richte es mit der Schwachheit unsers lieben Vaters (Mutter) dahin, daß es gereiche zu deines Namens Lob, Ehre und Preis, und zu unsers lieben Vaters (Mutter) Leibes und der Seele Heil und Wohlfahrt. Erhöre uns, deine lieben Kinder gnädiglich, um Jesu Christi, deines lieben Sohnes, unsers lieben HErrn und Seligmachers willen. Amen.

\section{249. Gebet der Taufpathen}

HErr Jesu Christe, du einiger Heiland und Seligmacher des ganzen menschlichen Geschlechts! Du hast gesagt: wer da glaubet und getauft wird, der wird selig werden. Und abermal: Es sei denn, daß jemand geboren werde aus Wasser und Geist, so
\quellseite{199}
kann er nicht in das Reich Gottes kommen. Derselben danke ich dir von Herzen, daß du beide meine liebe Gevattern mit der Leibesfrucht so väterlich und gnädiglich gesegnet, und mir die Gnade verliehen hast, daß sie ihre christliche Zuneigung in diesem hochheiligen Werke auf mich gewandt, und mich für andern ihnen zu ihrem Gevatter belieben und gefallen lassen: also, daß ich unwürdiger, anstatt des Unwürdigen, meines Pathen, dem bösen Feind und allem seinem Anhang widersagen, und dich, den wahren Gott, mit Herz, Mund und wahrhaftigem Glauben bekennen, und dir meinen lieben Pathen mit Leib und Seel, zeitlich und ewig, in deiner allergnädigsten Schutz und Schirm allerdings befehlen und überantworten soll und will, inmaßen denn du befohlen hast: Lasset die Kindlein zu mir kommen, denn ihr ist das Reich Gottes; also übergebe ich und befehle ich dir, o du Gott Abrahams, Isaaks und Jakobs, du Gott unserer Väter, dieses unwürdige Kindlein, meinen lieben Pathen, daß du ihn mit dem Lichte deiner väterlichen Barmherzigkeit erleuchten, segnen, behüten, und dein Angesicht von ihm nicht abwenden wollest, sondern in die Anzahl aller deiner Christgläubigen aufnehmen, und ihn allhier zeitlich und dort ewig (als den du mit deinem rosinfarbenen Blut so theuer erkauft hast) in deiner göttlichen Gnade erhalten, und ihm nach diesem Leben die immerwährende Freude und Seligkeit geben wollest. Amen.

\section{250. Fürbitte für die Confirmanden}

O allmächtiger, wahrhaftiger, lebendiger Gott, ewiger und barmherziger Vater unsers Heilandes Jesu Christi, der du dich uns aus großer Gütigkeit geoffenbaret hast, du sprichst von deinem lieben Sohne Jesu Christo: Das ist mein lieber Sohn, den sollt ihr hören — und er, dein lieber Sohn, hat selbst gesagt: Lasset die Kindlein zu mir kommen und wehret ihnen nicht, denn solcher ist das Reich Gottes. Weil denn diese gegenwärtigen Kinder zu deinem lieben Sohn Jesu Christo in der heiligen Taufe gebracht und in seinem Tod getauft sind, mit seinem Blute gewaschen und durch die heilige Taufe mit Christo begraben worden: so bitten wir dich, barmherziger Gott und Vater, du wollest aus milder Gütigkeit in diesen deinen Kindern den Heiligen Geist, den du über sie in der
\quellseite{200}
heiligen Taufe reichlich ausgegossen hast, wiederum erneuern, daß ihre Herzen mit dem Schein des heiligen Evangeliums erleuchtet werden. Vermehre in ihnen, o gnädiger Gott, die reine Erkenntniß deines lieben Sohnes Jesu Christi und den rechten Glauben, daß sie an dich, den wahrhaftigen Gott, und an Jesum Christum, deinen Sohn, den du gesandt hast, fest glauben und an dir mit festem Vertrauen beständig bleiben. Erlöse sie aus der Gewalt der Finsterniß und setze sie in das Reich deines lieben Sohnes, darin wir haben die Erlösung und Vergebung der Sünde. Gib ihnen den Frieden Christi und die Freude des Heiligen Geistes in ihr Gemüth und die Liebe Gottes und des Nächsten. Erfülle sie mit den Gaben des Heiligen Geistes, der sie lehre und leite in alle Wahrheit, und tüchtig mache, dir in deiner heiligen Gemeine zu dienen, die Geschäfte des Fleisches zu tödten, den Teufel in aller seiner List und Anfechtung zu überwinden, dir in Heiligkeit und Gerechtigkeit, die dir gefällt ist, zu dienen, daß sie die Zukunft deines lieben Sohnes, sammt allen Gläubigen mit fröhlichem Herzen und wachendem Gebete in Nüchternheit, Gerechtigkeit und Gottseiligkeit mögen erwarten und hoffen, zu deinem Ehren und Besserung deiner armen Kirche, der du lebst und regierest mit Jesu Christo, deinem eingebornen Sohne, in Kraft und Einigkeit des Heiligen Geistes, von Ewigkeit zu Ewigkeit. Amen.

\section{251. Gebet eines Hausvaters für das Gesinde}

O treuer Gott, du himmlischer Hausvater! Du hast in deinem heiligen Wort allen Knechten und Mägden, auch allem Hausgesinde ihr gebührend Amt, wie sie sich gegen ihre ordentliche Herrschaft christlich verhalten sollen, ausdrücklich fürgebildet. Wir finden aber zu diesen unsern letzten Zeiten bei dem größern Theil eben das Widerspiel; darum bitten wir dich, lieber Vater! gib deinem Heiligen Geist, daß sich das Hausgesinde in seinem Beruf, wie sichs gebühret, halte, erstlich zu allem christlichen Gehorsam in allen ehrlichen Dingen sich ihrer Herrschaft ergebe, folgendes ihre auferlegte Dienst und Arbeit im Schweiß ihres Angesichts, nicht aus Furcht, Zwang und Drang, sondern aus christlicher Liebe, mit gutem Willen, der Herrschaft zu Nutze und ihnen
\quellseite{201}
selbst zu zeitlicher und ewiger Wohlfahrt fleißig ausrichten, sich auch letzlich für aller Leichtfertigkeit, Untreue und Falsch ernstlich hüten, auf daß sie hier zeitlich ein gut Gewissen behalten, und also im Todtenbette, sonderlich aber am jüngsten Tage, wenn Christus, dein lieber Sohn, selbst die Rechnung ihrer Haushaltung von ihnen erfordern wird, bestehen und endlich die fröhliche Stimme anhören mögen: Komm her, du getreuer Knecht, du bist im Geringen getreu gewesen, ich will dich über viel setzen, gehe ein in die Freude deines HErrn, durch denselben deinen lieben Sohn, unsern einigen HErrn und Heiland. Amen.

\section{252. Gebet für Krancke insgemein}

O du starker, gewaltiger Gott, der du bist eine Kraft, welche in den Schwachen mächtig ist, eine Stärke aller, die auf dich hoffen, eine treue Hülfe in aller Noth zur unangenehmen Zeit. Ich bitte dich für alle kranken Menschen, die mit Leibesschwachheit beladen, oder in Todesnöthen sein, du wollest ihnen geben, daß sie deinen heiligen Willen erkennen, mit Geduld denselben annehmen, leiden und tragen. Wollest sie auch durch deinen heiligen Geist stärken und trösten, sie erquicken auf ihrem Siechbette, und ihnen helfen von allen ihren Krankheiten. Erbarme dich ihrer nach deiner großen Güte, verzeihe ihnen, um des bittern Leidens Jesu Christi, deines lieben Sohnes willen, alle ihre Sünde und Missethat, damit sie dich erzürnt und etwa solche ihre Krankheiten verursachet haben, und wenn es deinem heiligen göttlichen Willen und ihrer Seelen Seligkeit nicht zuwider ist, so wollest du ihnen wiederum aufhelfen, sie an Leib und Seel gesund machen und stärken; denn so du willst, kannst du sie wohl reinigen; du bist ja unser lieber Arzt, unsre beste Zuversicht, und wenn aller Menschen Hilfe aus ist, so thust du am liebsten helfen, derohalben so schaue vom Himmel und siehe an das Elend und Jammer aller Kranken, laß es dir zu Herzen gehen, und schaffe Rath und ordentliche Mittel durch deine göttliche Weisheit, thue Hülfe mit deinem ausgestreckten Arm, daß sie mögen genesen und gesund werden. Welchen es aber nützlicher ist, aus diesem Elend nach deinem göttlichen Willen jetzund abzuscheiden, denn allhie in diesem Jammertahl länger zu wallen, die wollest du
\quellseite{202}
im rechten festen Glauben und starker Hoffnung mit einem seligen Ende zu dir in die ewige Freude versammeln. Stehe ihnen bei mit deinem Heiligen Geiste an ihrem letzten Stündlein, gib ihnen Gnade, daß sie sich zum Sterben wohl und willig bereiten, auch mit Empfahung des allerheiligsten Abendmahls Christi, so viel möglich versorget werden, und bei ihm sein in der Schaar aller Auserwählten. Verkläre ihnen ihres Leibes Pein und Qual, die sie jetzt leiden; hilf ihnen überwinden alle Schmerzen des Todes, daß sie hindurch bringen zum ewigen Leben. Verkehre ihre Traurigkeit in Freude, die nicht von ihnen mag genommen werden; nimm ihre Seele in deine Hände, welche du geschaffen und erlöst hast, durch Jesum Christum, unsern Herrn. Amen.

\section{253. Gebet für einen Kranken}

Herr, siehe, den du lieb hast und der dich lieb hat, der ist krank; ach! erbarme dich über ihn, hilf ihm, so ist ihm geholfen an Leib und an der Seelen. Amen.

\section{254. Drei Fürbitten für einen Besessenen}

\subsection{Die Erste}

Herr Jesu Christe, der du bist niedergegestiegen zur Hölle und den Satan gebunden hast, wir bitten dich, fahre doch in den Abgrund der Höllen, die der Satan in des Besessenen Herzen angerichtet hat, zerstöre sie, binde da den bösen Feind, daß er nichts schaffen könne, treib ihn von ihm, zerstöre sein Reich und Werk in ihm und sprich in seine Seele deinen Segen, deinen Trost und deinen Frieden, deine Lust und Freude. Amen.

\subsection{Die Zweite}

Herr Jesu Christe, der du einem armen stummen und tauben Menschen mit deinem allmächtigen Wort gnädiglich geholfen hast und daneben unser aller Elend herzlich beseufzet: Wir klagen dir, der Satan hat diesen Jüngling auch beleidiget am Leibe und zu allem Guten schrecklich verderbet und niemand kann ihm helfen, denn du Herr Jesu Christ alleine. Wir bitten dich, erbarme dich seines Elendes, rühre an sein Herz mit dem Finger des Heiligen
\quellseite{203}
Geistes, rede dein kräftiges Wort in seine Seele, sprich Hephata zu ihm, daß alsbald zerspringen die eisernen Banden der Sünden, damit er umgeben ist. Eröffne seine Ohren, dein heiliges Wort seliglich zu hören. Eröffne seinen Mund zu reden, was dir gefällig und ihm seliglich ist, und daß er allezeit deinen heiligen Namen mit dankbarem Herzen, mit heiligem Munde und fröhlicher Stimme loben und preisen und sagen möge das schöne Wort: „Er hat Alles wohl gemacht im Himmel und auf Erden.“ Das hilf ihm, HErr JEsu Christe, gelobet in Ewigkeit. Amen.

\subsection{Die Dritte}

O allmächtige, ewige, heilige Dreifaltigkeit, Vater, Sohn und Heiliger Geist, du einiges ewiges göttliches Wesen, wir danken dir von Herzen für alle Wohlthat und für den Reichthum deiner Gnade an uns gewendet, und bitten dich durch deine ewige Gottheit, Heiligkeit und Herrlichkeit, erbarme dich unser und erhalte uns bei deinem Wort und Wahrheit, bei der rechten Erkenntniß deines Wesens und Willens in festem Vertrauen, in heiliger Liebe deines Namens und Gehorsam deines Willens, daß wir dich mit den Engeln loben und ehren und dich fürchten. O HErr! deine große Liebe und Barmherzigkeit tilge des Besessenen Sünde! Deine unaussprechliche Güte komme ihm zu Hülfe in all seinem Elende, deine allmächtige Gewalt sei seine Stärke wider den Satan und alles Unglück, deine Treue helfe ihm aus aller Noth, deine Huld, Gunst und Güte stehe ihm bei und tröste ihn.

O du ewige Kraft Gottes, Vater, Sohn und Heiliger Geist, stärke ihn. O du ewige Gottheit Gottes erbarme dich seiner und sei ihm gnädig. O du allerheiligste Majestät Gottes hilf ihm in allen Nöthen, erhalte ihn zu deines Namens Ehre und seiner Seelen Seligkeit durch JEsum Christum. Amen.

HErr Gott Vater, der du deinen Sohn gesandt hast, des Satans Werke zu zerstören und uns in dein Reich zu bringen: wir bitten dich, mache doch in dem Besessenen den seligen Anfang deines göttlichen Reiches und zerstöre in ihm des Satans Werk. Amen.

Ewiger, wahrhaftiger Sohn Gottes, JEsu Christe, unser HErr und König, ein Ueberwinder des Satans, zu dir schreien
\quellseite{204}
wir und bitten demütiglich, treib aus von ihm den bösen Geist mit deines Fingers Kraft, verleihe uns Stärke, Wehr und Waffen, dem Satan zu widerstehen. Hilf uns selber kämpfen, streiten, siegen und überwinden. Amen.

O heiliger Geist erfülle uns mit Kraft vom Himmel, dem Satan zu widerstehen und ihm abzusiegen. Amen.

\section{255. Gebet für eine vom Satan besessene Person}

Herr Gott, himmlischer Vater, der du uns durch deinen lieben Sohn befohlen hast zu beten und das Predigtamt in der heiligen christlichen Kirche geordnet und eingesetzt hast, daß wir die Brüder, so etwa durch einen Fehl übereilet werden, mit sanftmüthigem Geist unterweisen und wieder zurecht bringen sollen; und Christus, dein lieber Sohn, sagt selber, er sei nicht kommen, denn nur allein um der Sünde willen. Darum bitten wir dich für diesen deinen Diener, du wollest ihm seine Sünde vergeben, und ihn in den Artikel der Vergebung der Sünden wieder einschließen und in den Schooß der heiligen Kirche wieder annehmen, um deines lieben Sohnes willen, unsers Herrn Christi. Amen.

\section{256. Gebet für eine kranke aufgeschreckte Person}

Herr Gott, himmlischer Vater, der du uns für die Kranken hast heißen beten: wir bitten dich durch Jesum Christum, deinen lieben Sohn, daß du diese deine Dienerin von ihrer Krankheit und von des Teufels Banden erlösen wollest. Schone doch, lieber Herr, ihrer Seele, die du sammt ihrem Leib durch deines lieben Sohnes Jesu Christi Blutvergießen erworben und errettet hast von der Sünden, des Todes und des Teufels Gewalt. Amen.

\section{257. Gebet bei einem Begräbniß}

Allmächtiger Gott, lieber Vater, wie gar nichts sind alle Menschen, die in der Welt leben. Du hast uns nun ein Ziel unsers Lebens gesetzet, welches wir nicht übergehen noch überschreiten können. Wir haben anjetzo aber einen Mitbruder (eine Mitschwester) zum Grabe begleitet. Ach Herr! lehre uns doch, daß wir sterben müssen. Regiere uns durch deinen Heiligen Geist, daß
\quellseite{205}
wir uns allezeit bereit halten, wenn du kommen und auch uns abfordern wirst. Laß uns nicht erschrecken vor dem Tode, sondern mit freudigem Herzen unserer Erledigung erwarten und laß uns nach diesem elenden Leben bei dir ewig leben. Verwandele alsdann unser Leiden in Freude und wische ab unsere Thränen. Hole heim, HErr, deine Knechte, wenn dirs gefällt; gib nur ein seliges Stündlein. Dem nun Verbleichenen verleihe eine sanfte Ruhe im Schooß der Erden und an dem Erscheinungstage deines geliebten Sohnes JEsu sammt allen Auserwählten eine fröhliche Auferstehung. Die Hinterbliebenen tröste kräftiglich mit deinem Heiligen Geist; sei ihr Gott, ihr Vater und Beschützer zu Lob und Ehr deines heiligen Namens, um deiner großen Barmherzigkeit willen. Amen.

\section{258. Gebet eines Zuhörers für seinen Seelsorger}

Allmächtiger, gütiger Gott, himmlischer Vater! ich bitte dich für meinen lieben Seelsorger, daß du ihm gebest mit allerlei Freudigkeit zu reden dein Wort wider falsche Lehre und Irrthum und das Geheimniß deines Evangelii zu verkündigen. Erhalte ihn in reiner Lehre und heilsamem Leben. Gib ihm gesunden Leib, Kraft und Stärke, daß er uns lange Zeit mit Nutz und Frucht vorstehen könne. Laß ihn alle Sünde und Laster strafen, wie sich’s gebühret, damit ich und viele Andere zu deinem Reich gewonnen werden. Eröffne mir mein Herz und Ohren, daß ich mit Lust und Liebe, mit Andacht meines Gemüthes und herzlichem Aufmerken meinem Seelsorger zuhöre, und nach seiner Lehre und Predigt in rechtem Glauben gottseilig wandle und Frucht bringe, damit dein Wort nicht leer abgehe, sondern thue, was dir gefalle, und ihm gelinge, dazu du es gesandt hast. Hilf, daß ich alle väterliche Strafe und Vermahnung von meinem Prediger gutwillig aufnehme, und mich daraus bessere, von meinen Sünden abstehe, und in deinem Gehorsam lebe. Erhalte uns allesamt in wahrem Glauben und christlichem Leben, daß wir darinnen täglich zunehmen, und bis ans Ende beharren, und alsdann durch deinen lieben Sohn JEsum Christum ewig selig werden. Amen.

\quellseite{206}

\section{259. Gebet eines Pfarrkindes}

Ewiger, gütiger Gott, himmlischer Vater! ich bitte dich für meinen Seelsorger, daß du ihm gebest dein Wort mit freudigem Aufthun seines Mundes, unerschrocken wider allen Irrthum, falsche Lehre und Mißbräuche zu reden, auf daß er uns das Geheimniß deines Evangelii offenbare und verkündige, allen falschen Wahn aus unsern Herzen reiße. Erhalte ihn bei reiner, heilsamer Lehre und in christlichem Wandel, daß er uns vorgehe zum ewigen Leben. Behüte seinen Leib und Gesundheit, auf daß er uns lange Zeit mit Frucht und Nutz könne vorstehen und dein göttliches Wort ohne alle Scheu, Furcht und Entsetzung, ohne Heuchelei, nicht aus Gunst, Haß, Neid oder Eigennutz möge predigen, sondern die reine Wahrheit lauter sagen und die Laster strafen, wie sichs gebühret, damit ich und andere viel zu deinem Reich gewonnen werden. Eröffne mir mein Herz und Ohren, daß ich mit Lust und Liebe, mit Andacht meines Gemüthes und herzlicher Aufmerckung deinem heilwärtigen Worte zuhöre, und nach demselben in rechtem Glauben gottseelig wandle, und Frucht bringe zu deinem göttlichen Ehren, auf daß dein Wort nicht leer abgehe, sondern thue, das dir gefalle, und ihm gelinge, dazu du es gesandt hast. Deshalb nimm von mir allen Ueberdruß und des Herzens Trägheit, und senke in mein Gemüth einen rechten Hunger und ernstliches Verlangen nach dem überschwänglichen Reichthum deiner Gnade; so uns in der Predigt vorgetragen wird. Gib mir Gnade, daß ich meinen Seelsorger erkenne und halte für deinen Diener und Haushalter deiner göttlichen Geheimnisse, auf daß ich dein Wort aus seinem Munde annehme, nicht anders als von dir selbst, und nicht verachte das heilige Amt, welches du den Menschen befohlen hast; auch daß ich dein Wort nicht geringer achte von wegen etlicher Gebrechen dessen, der es predigt und verkündiget. Hilf, daß ich alle väterliche Strafen und Ermahnungen von meinem Prediger gütlich und willig aufnehme, dieselbigen mit Wohlmeinung verstehe, mich daraus bessere, von meinen Sünden abstehe und nicht die Züchtigung hasse noch verwerfe, oder diejenigen, so mich strafen, anfeinde und lästere. Erhalte uns allesamt in wahrem Glauben und christlichem Leben,
\quellseite{207}
daß wir darinnen wachsen, täglich zunehmen und beständig bis an unser Ende verharren, und durch deinen lieben Sohn Jesum Christum ewig selig werden. Amen.

\section[260. Gebet um rechtschaffene Prediger, …]{260. Gebet um rechtschaffene Prediger, sonderlich in Vacanz zu beten}

O allmächtiger Gott! der du durch deinen lieben Sohn Jesum Christum uns befohlen, dich zu bitten um Arbeiter in deine Ernte, und uns Erhörung gnädiglich zugesaget hast. Ich bitte dich um desselben deines lieben Sohnes willen, du wollest deiner lieben Christenheit rechtschaffene Diener und Prediger deines heiligen Wortes mit großer Anzahl und Schaaren geben. Sende uns Weise und Schriftgelehrte, begnade sie mit deinem Heiligen Geist, welche du uns zuschickest. Erfülle sie mit wahrem Verstand deines Worts, daß sie uns den Weg zum ewigen Leben recht weisen, die Schrift unverfälscht auslegen, dazu uns für unrechtem Glauben und Menschentand, sammt allem Gift der Seelen warnen, auf daß wir in reiner Lehre und Glauben lauter erfunden werden. Denn wie möchten wir sonst glauben, von dem wir nichts gehöret hätten? wie sollen wir aber hören ohne Prediger? Wie sollen sie aber predigen, wenn sie nicht gesandt werden? Darum, o herzallerliebster Vater! setze in deine Gemeinde gottesgelehrte Bischöfe, christliche Pfarrherrn, reine Prediger, treue Hirten, fromme Lehrer und Seelsorger, welche durch deinen Heiligen Geist die Gemeinde, so dein lieber Sohn mit seinem eignen Blut erworben hat, treulich weiden, das heilige Evangelium, dein theuer werthes Wort zu Ehren deines göttlichen Namens mit vielem Nutzen und Frucht verkündigen und ausbreiten, die Sündhaftigen zur Buße bekehren, die Schwachen im Glauben stärken, die blöden Gewissen trösten wider alle Anfechtungen des Satans und der argen Welt, und die da mächtig sind, zu ermahnen durch die heilige Lehre und zu strafen die Widersprecher, auf daß sie den Widersachern können gewaltig Widerstand thun, sie eintreiben und überwinden, des Teufels Reich zerstören und verwüsten. Welche auch der Heerde und Gemeinde Gottes dermaßen vorgehen mit gutem Exempel und gottseeligem Wandel, daß alle Aergerniß verhütet, und nicht jemand durch ein
\quellseite{208}
gottloses Leben vom Glauben abgehalten werde. O barmherziger Gott, schicke uns nicht ins Land einen Hunger und Mangel deines Worts, daß wir der reinen Prediger nicht beraubet werden. Laß uns nicht entgelten unsere Undankbarkeit. Behüte uns, lieber Vater, für den bösen Arbeitern, für den untreuen Miethlingen, und allen, so von ihnen selbst kommen gelaufen, und nicht von dir gesendet werden, durch JEsum Christum deinen lieben Sohn, unsern HErrn. Amen.

\section{261. Gebet am Bibelfeste}

Ach du getreuer, gnädiger und barmherziger Gott! Wir danken dir für die große Gnade, Liebe und Treue, die du uns armen elenden Menschen erzeiget hast durch die Offenbarung deines heiligen göttlichen Wortes und die Einsetzung deiner hochwürdigen Sacramente. Der Geist des HErrn hat durch die Propheten geredet, und seine Rede ist durch ihren Mund geschehen. Dadurch hast du dich uns zu erkennen gegeben, in welcher deiner Erkenntniß steht das ewige Leben; dadurch hast du uns berufen zur Gemeinschaft deiner heiligen Kirche, dein heiliges Evangelium ist auch eine Kraft Gottes, selig zu machen alle, die daran glauben; dadurch erweckest du in uns den Glauben, giebst uns den Heiligen Geist, lässest uns Vergebung der Sünden verkündigen, dadurch empfangen wir wahrhaftigen lebendigen Trost in allem Kreuz, in Todesnoth; dadurch stillest du unser trauriges Gemüthe, dadurch erfreuest du unser trauriges Herz, dadurch speisest du unsere Seelen, als mit dem rechten Himmelsbrod, dadurch machest du uns zu neuen Creaturen, dadurch kommst du zu uns und machest Wohnung bei uns, dadurch linderst du unsere Todesnoth, daß wir den Tod nicht schmecken ewiglich, dadurch erleuchtest du unsere Seelen, dadurch regierest du und leitest uns in unserm Beruf und in dem Lauf unsers ganzen Lebens, als mit einer Leuchte auf unserm Wegen und als mit einem Lichte auf-unsern Fußsteigen, daß unsere Tritte nicht gleiten. Ach HErr! wenn dein Wort nicht wäre unser Trost gewesen, so wären wir vergangen in unserm Elende, denn dein Wort ist unsers Herzens Freude und Trost. Durch Gottes Wort werden wir geheiliget und an Leib und Seele gesegnet; darum, o HErr,
\quellseite{209}
wer dein Wort verwirft, der hat dich selbst verworfen, und allen seinen zeitlichen und ewigen Segen und kann mit nichts den Segen haben, sondern den ewigen Fluch, so verhält auch der Himmel seinen Thau und das Erdreich sein Gewächs, so wird der Himmel Eisen und die Erde Erz und ist alle Mühe und Arbeit verloren. O HErr! wir danken dir für diesen edeln, theuren Schatz, denn das Gesetz des HErrn ist ohne Wandel und erquicket die Seele, das Zeugniß des HErrn ist gewiß und machet die Albernen weise. Die Befehle des HErrn sind richtig und erfreuen das Herz, die Gebote des HErrn sind lauter und erleuchten die Augen. Die Furcht des HErrn ist rein und bleibet ewiglich, die Rechte des HErrn sind wahrhaftig, allzumal gerecht. Sie sind köstlicher denn Gold und viel feines Gold, sie sind süßer denn Honig und Honigseim. Auch wird dein Knecht durch sie erfreuet und wer sie hält, hat großen Lohn. Dein Wort, du Allerschönster, ist ein Grund der Weisheit und das ewige Gebot ist ihre Quelle, wer dasselbe ehret, der thut den rechten Gottesdienst und wer es lieb hat, den hat Gott wieder lieb. Gottes Wort spricht: Ich bin wie die Cedern auf dem Libanon und wie die Cypressen auf dem Gebirge Hermon. Ich bin aufgewachsen wie ein Palmbaum am Wasser, wie ein Rosenstock zu Jericho, wie ein schöner Oelbaum. Ich gebe einen lieblichen Geruch von mir, wie Zimmt und Myrrhen und Weihrauch. Kommt her alle, die ihr mein begehret und sättiget euch von meinen Früchten. Meine Predigt ist süßer denn Honig und meine Gabe süßer denn Honigseim. Wer von mir isset, den hungert immer nach mir und wer von mir trinket, den dürstet immer nach mir. Das ist das Buch des Bundes mit dem höchsten Gott gemacht, nämlich das Gesetz Mose, daraus ist geflossen die Weisheit, wie der Euphrat, wenn er groß ist, wie der Jordan in der Ernte, wie der Nil im Herbste. Es ist nie gewesen, der es ausgelernt hat und wird nimmermehr werden, der es auszugründen mag. Denn sein Sinn ist reicher denn sein Meer und sein Wort ist tiefer, denn sein Abgrund. Es fließen von ihm viel Bächlein in die Gärten, da werden die Bächlein zu großen Strömen und die Ströme zu großen Wassern. Denn seine Lehre leuchtet so weit als der lichte Tag und scheinet frühe. Auch schütten sie Lehre und
\quellseite{210}
Weissagung aus, die ewig bleibet. Für diesen hohen, theuerbaren Schatz danken wir dir, Gott Vater, Sohn und Heiliger Geist, hochgelobt in Ewigkeit. Amen.

\section{262. Gebet für Missionsstunden}

Ich bitte dich, mein himmlischer Vater, durch JEsum Christum, deinen lieben Sohn, unsern HErrn, du wollest dich gnädiglich erbarmen über alle Ungläubigen, wer und wo sie seien, so noch in Finsterniß wandeln, welche das Licht deines Evangelii nicht haben, die der Teufel mit Blindheit geschlagen und ihr unverständiges Herz verfinstert hat, die entfremdet sind von dem Leben, das aus Gott ist, durch die Unwissenheit, so in ihnen ist, und zu den stummen Götzen gehen, wie sie geführet werden, und deinen lieben Sohn, JEsum Christum, den Gnadenthron, aus Blindheit verfluchen und lästern. Derowegen, du getreuer Gott, nimm hinweg die Decke, welche da hänget unaufgedeckt vor den Herzen der Juden, die sich stoßen an dem Stein des Anlaufens und an dem Fels der Aergerniß. Erleuchte ihre Augen, auf daß sie erkennen den wahren Messiam, der Welt Heiland. Bringe die Heiden und alle Ungläubigen, so dein Wort für eine Thorheit achten, in den rechten Schafstall und Versammlung der Christen, in die Gemeinschaft der Heiligen, auf daß sie mit uns und allen Gläubigen dich, den Vater im Sohn und den Sohn mit dem Vater und Heiligen Geist, einen wahren Gott und HErrn ehren, preisen und anbeten. Eröffne das Verständnis allen Menschen, so ihre Seligkeit und Vertrauen nicht einzig und allein auf den rechten Grund= und Eckstein, Christum JEsum, setzen, daß sie denselben, welchen du gesandt hast, erkennen, mit rechtem Glauben und wahrer Zuversicht des Herzens für ihren Heiland und Erlöser annehmen, der mit seinem Gehorsam und Erfüllung des Gesetzes durch seinen bittern Tod eine ewige Erlösung erworben hat. Wollest auch Alle, so vom christlichen Glauben abgewichen oder sonsten in etlichen Stücken irrige und mit falscher Lehre behaftet sind, wiederbringen zu der ewigen Wahrheit deines Worts. Gott aller Gnaden, erbarme dich über die, welche des rechten Glaubens nicht sind, so da sitzen im Schatten des Todes und wandeln nach der Dunkelheit
\quellseite{211}
ihres Sinnes auf unebner Bahn. Suche das Verlorne, bringe wieder das Verirrte, erfreue die Verblendeten, eröffne die Ohren der Tauben, thue auf den Mund der Stummen, die dich nicht bekennen, richte auf die Gefallenen, hole wieder die Abgewichenen, versammle die Zerstreuten, bringe zurecht die Irrigen und Verführten um deiner Barmherzigkeit willen. Amen.

\section{263. Gebet für die Kirche}

Herr, barmherziger und gnädiger Gott, dieweil diese allein deine Bürger und Hausgenossen sind, welche erbauet werden auf den Grund der Propheten und Apostel, da JEsus Christus der Eckstein ist, und du segnest mit Gnaden allein die, so in deinem Lichte wandeln; darum bitte ich dich, du wollest dein Wort in der Welt also erschallen lassen, daß hie und da ein Häuflein sei, der die Stimme Christi höre, und dein Wort lerne, daß man überall halte die Gebote Gottes und den Glauben an unsern HErrn JEsum Christum, damit die gläubige Gemeine je länger je mehr wachse und der Wagen Gottes sei viel tausend mal tausend, welche hineingehen zu den Thoren der Gerechtigkeit, in deinem Hause zu wohnen, und deswegen von der Ungerechtigkeit abtreten, und dir allein dienen in Heiligkeit und Gerechtigkeit, die dir gefällig ist. O, HErr Jesu Christe, der du bist der Anfänger und Vollender des Glaubens, und mächtig in den Versammlungen der Heiligen, regiere durch deine innerliche Gnade der Menschen Herzen, daß sie sich zu dir bekehren: Gib deinem Donner, das ist, der Predigt deines Wortes Kraft, laß ihm gelingen, dazu du es sendest: erleuchte und stärke der Menschen Herzen, daß sie sich von deiner Gemeine nicht absondern, sondern gerne wallen mit dem Haufen zum Hause, da dein Wort und heilige Sakramente rein und unverfälscht vorgetragen und ausgespendet werden. Ob es wohl eine kleine Heerde, ein armes geringes Volk ist, über welches alle Wetter gehen und fast nirgend seinen Trost in der Welt findet, so soll sich doch die kleine Heerde nicht fürchten, dieweil es dein Wohlgefallen ist, derselben das Reich deiner Herrlichkeit zu geben, da du die Schmach deines Volks wirst hinwegnehmen, und alle diejenigen, so bei dir in deinen Anfechtungen verharrret haben, zu dir ziehen, daß sie essen und trinken sollen
\quellseite{212}
über deinem Tisch in deinem Reich, und sitzen auf Stühlen und richten die zwölf Geschlechter Israels. Erfreue den Teufel nicht, der dir und deinem Wort feind ist, und hindert, wo er kann, daß dasselbe nicht gehört und geglaubet werde. Verkläre deinen Namen, daß viel Seelen hinzugethan werden zu der Zahl der Auserwählten, die in deinem Hause wohnen und dich loben immerdar. Amen.

\section[264. Gebet für die christliche Kirche und …]{264. Gebet für die christliche Kirche und wider ihre Feinde}

Allmächtiger, ewiger, barmherziger Gott, und Vater unseres lieben Herrn und Heilandes Jesu Christi, wir sehen und fühlen es, wie es deiner Kirche in diesem Leben gehet, was sie für Glück hat, und wie sie auf so mancherlei Weise vom Teufel und von der Welt geplaget wird. Darum bitten wir dich, um desselben deines eingeborenen Sohnes willen, erstlich: du wollest unsere Herzen mit deinem Heiligen Geiste trösten und stärken, auf daß wir von so viel großer Gefahr nicht überwältiget werden, noch unterliegen. Zum andern: du wollest auch der Feinde Vornehmen und Anschläge nicht allein hindern, sondern mit deiner treuen und wunderbaren Hülfe der ganzen Welt anzeigen, erklären und beweisen, daß du für deine Kirche sorgest, sie regierest, schützest, erhaltest und errettest, der du lebest und regierest, ein ewiger Gott, Gott Vater, Gott Sohn, Gott Heiliger Geist, von Ewigkeit zu Ewigkeit. Amen.

\section[265. Gebet wider den Antichrist und Papst zu …]{265. Gebet wider den Antichrist und Papst zu Rom mit seinem Anhang}

Lieber Gott und Vater unseres Herrn Jesu Christi, wir bitten dich, du wollest uns einmal wieder heimsuchen nach allen deinen Wunderwerken, und uns zeigen den Tag der Zukunft der Herrlichkeit deines Sohnes, damit der Schalk, der Antichrist, verstöret und zerbrochen werde, der da ist der Mensch der Sünden und der Sohn des Verderbens, und mache es einmal ein Ende mit den gewaltigen Irrthümern des Teufels, durch welche leider alle Augenblick viel tausend Seelen verderbet werden und in die Hölle gerissen, von deßwegen allein, daß man die Tyrannei des
\quellseite{213}
greulichen und apostatischen, das ist, abtrünnigen Stuhls zu Rom in seinem Wesen erhalten möge, da sage alle Welt zu Amen, Amen.

\section[266. Gebet treuer Lehrer und Prediger auf …]{266. Gebet treuer Lehrer und Prediger auf Pastoral-Conferenzen}

Du gütiger und barmherziger Gott, du hast uns aus sonderlicher Gnade zum Predigt-Amt berufen, und diese geistlichen Werke: lehren, predigen, taufen, absolvieren, Sacrament reichen, aufgelegt. Nun sind wir willig und bereit, wollen es auch von Herzen gern thun. Du siehest aber, wie schwach und unvermögend wir dazu sind aus unsern Kräften; siehest auch, wie greulich sie vom Teufel und seinem Haufen werden angefochten. Darum, du lieber Gott, fördere diese Werke, gib Stärke und Kraft dazu, daß wir sie treulich vollbringen! Hindre und wehre dem Teufel und allem seinem Anhange, daß sie diese Werke nicht befudeln, daß das Gesetz und Evangelium durch falsche Lehre und Heuchelei nicht in einander gemenget, Taufe und Sacrament durch Wiedertäufer und Sacramentirer nicht geschwächet, die Absolution nicht durch falschen Ablaß und Lügen des Papstes verfälschet werde, sondern ein jegliches in seiner Kraft und Würden bleibe, und dein Heiliger Geist, der in uns wohnet, durch unsere Sünde nicht betrübt werde, sondern in allen diesen Stücken thätig und kräftig sei, auf daß wir als deine getreuen Diener allzeit erfunden werden. Amen.

\section[267. Gebet um Beständigkeit und Erhaltung bei …]{267. Gebet um Beständigkeit und Erhaltung bei dem seligmachenden Wort am Schluß von Pastoral-Conferenzen}

HErr Gott Vater aller Gnaden und Stärke, bestätige und stärke uns gnädiglich in deinem angefangenem Werke durch deinen Heiligen Geist, auf daß der Satan durch seine Gewalt noch List uns schwäche, noch müde mache, dein Wort und Reich zu verlassen. Denn es ist jetzt gefährliche Zeit, weil viel durch die Rottengeister verführet werden und viel auch abfallen, viel werden der Gnaden Gottes auch überdrüssig und kalt, daß sie, vom Satan mit List betrogen, sich dünken lassen, sie seien satt, können nun alles und haben keine Noth, werden also faul und undankbar, und bald hernach
\quellseite{214}
ärger denn vorhin. Darum laß du uns bleiben in der Brust des Glaubens, daß wir drinnen täglich zunehmen, in Christo Jesu, unserm rechten und einigen Helfer. Amen.

\section{268. Gebet in Gemeindeversammlungen}

Barmherziger, gütiger Gott, der du uns gelehret und geboten hast, daß wir zum ersten und vor allen Dingen nach deinem Reich und nach seiner Gerechtigkeit sollen trachten: ich bitte dich, du wollest Gnade geben, daß dein heiliges Wort rein und lauter in aller Welt geprediget werde, und wir unsere Vernunft unter den Gehorsam des Glaubens gefangen nehmen, auch nach demselben heilig, als die Kinder Gottes, zu deinem Wohlgefallen leben, auf daß dein Reich in und zu uns komme, und gemehret, auch ihrer viel, so noch nicht an das Wort glauben, durch einen christlichen Wandel gewonnen werden. Hilf, lieber Gott, daß, die wir aus Gnaden erlöset sein von der Gewalt der Finsterniß und versetzet in das Reich deines lieben Sohnes Jesu Christi, an welchem wir haben die Erlösung durch sein Blut, nämlich die Vergebung der Sünden, daß wir auch in seinem Reich bleiben, in heilsamer Lehre beständig verharren, und würdig, als die Kinder des Lichtes in aller Gottseligkeit und Ehrbarkeit wandeln. Nachdem aber das Reich Gottes nicht kommt mit äußerlichen Gebärden, bestehet auch nicht in Worten, sondern in der Kraft und im Geist, so verleihe Gnade, daß wir durch dein seligmachendes Wort und Heiligen Geist auf das Neue von oben herab gezeuget und wiedergeboren werden, als Miterben des Lebens, auf daß wir mit unserm Gemüth im Himmel wohnen, und stetiglich nach dem unvergänglichen, unbefleckten und unverwelklichen Erbe streben. Hilf, daß wir geistlich arm und demüthig sein, Leid tragen über unsre Sünde; gib uns Hunger und Durst, das ist, herzliches Verlangen nach der Gerechtigkeit, und daß wir auch daneben sanftmüthig, barmherzig, rein von Herzen und friedfertig sein, dazu um der Gerechtigkeit und Wahrheit willen allerlei Verfolgung und Widerwärtigkeit, böse Nachrede und unverdiente Lästerung mit Geduld leiden und überwinden. Behüte uns vor allen Aergernißen, dadurch dein heiliger Name gelästert und geschändet, dein Reich gehindert und
\quellseite{215}
geschwächet wird. Verleihe Gnade, daß wir unsern Glauben üben in allerlei Werken der Liebe und Barmherzigkeit mit Speisen, Tränken, Kleiden, Beherbergen, Trösten und Besuchen der Armen und Elenden in dieser Welt, auf daß wir an jenem Tage hören die fröhliche und gnadenreiche Stimme deines lieben Sohnes: Kommet her, ihr Gesegneten meines Vaters, ererbet das Reich, welches euch bereitet ist von Anbeginn der Welt. Amen.

\section[269. Gebet für die Kirche in …]{269. Gebet für die Kirche in Gemeindeversammlungen}

Herr und Gott, weil ja deine Heiligen abgenommen und deiner Gläubigen wenig geworden ist auf Erden, und daher zu fürchten, es möchte der Frommen so wenig werden, daß es allenthalben voll Gottloser sein dürfte: wohlan, mache dich auf, und hilf deinen Armen und Elenden, daß sie, die jetzt seufzen, nicht gar zerstöret werden. Laß deine göttliche Wahrheit getrost gelehret und ohne Heuchelei vor männiglich bekennet werden. Breite dein Wort weit und breit aus, daß wie durchläutertes Silber rein und lauter ist, und behüte uns vor dem argen Geschlechte derer, die anstatt deines Wortes unnütze Dinge vorgeben, oder sonst die Religion sich nicht lassen einen Ernst sein, und vermehre das kleine Häuflein, daß es wider die große Menge der Gottlosen fest bestehe und alle Pforten der Hölle ihm nichts anhaben mögen, um Jesu Christi willen. Amen.

\section{270. Gebet um Frieden in der Gemeinde}

Du allmächtiger, ewiger Gott, lieber himmlischer Vater, du Gott der Einigkeit und des Friedens, der du mit deinem eingebornen, geliebten Sohn, Jesu Christo, und Heiligen Geist nicht allein eins bist, sondern auch in deinem heiligen Wort verheißen hast, du wollest uns, wenn wir dich darum anrufen, ein einträchtiges Herz, Sinn und Gedanken geben: wir bitten dich innig und einmüthiglich, du wollest aller Uneinigkeit steuern und wehren, den Satan, als Stifter der Uneinigkeit, unter unsre Füße treten, daß wir keinem Eingeben nicht Raum lassen und hingegen uns sämmtlich und sonderlich ein sanftmüthiges, demüthiges und friedliebendes Herz schenken, daß wir als Brüder einträchtig bei einander leben, mit
\quellseite{216}
freundlichen Geberden und holdseligen Worten den Frieden stiften und erhalten und fleißig sein zu halten die Einigkeit im Geist durch das Band des Friedens, also daß unser aller, die wir an deinen Namen glauben und eines Leibes Gliedmaßen sind, eine Seele und ein Sinn sein und wir uns damit als deine lieben Kinder in der That erweisen mögen. O du liebreicher Gott des Friedens, der du den menschlichen Leib also wunderlich vermenget hast, daß dem dürftigen Glied am meisten Ehre gegeben würde, auf daß nicht eine Spaltung im Leibe sei, sondern alle Glieder zugleich für einander sorgen und einig sein. Verleihe uns, daß wir uns auch für aller Uneinigkeit und Widerwillen, Dräuung, Hader und Zank treulich hüten und hingegen der herzlichen brüderlichen Liebe, Fried und Einigkeit befleißigen, damit wir hier in dieser Zeit mit dem Band der Liebe zusammen verknüpfet ein friedsames, gottfeliges Leben führen, ein ruhiges Gewissen behalten, nach dieser Zeit aber mit dir, o du Gott des Friedens, und allen heiligen Engeln, Gläubigen und Auserwählten, in gewünschter Ewigkeit im himmlischen Friedensland, unauflöslicher Einigkeit und unaussprechlicher Herrlichkeit uns freuen mögen, da unser seliger Himmelsfriede nimmermehr aufhören wird. Ach lieber Herr und Gott, laß uns mit Simeon aus der Unruhe dieser Welt in Friede fahren zum ewigen Himmelsfrieden, den du uns mit deinem herben bittern Tod erworben hast. Amen.

\section{271. Gebet um Bekehrung der Verführten}

Gütiger Gott und Vater, der du willst, daß allen Menschen geholfen werde, und zur Erkenntniß der Wahrheit kommen, ich bitte dich, du wollest dich erbarmen über alle Ungläubige, so noch in Finsterniß und Schatten des Todes sitzen, deren Sinne verstockt und vom Teufel verblendet sind, daß sie das helle Licht deines Evangelii nicht sehen, damit sie von dem Irrthum ihres Weges bekehret, und ihren Seelen vom Tode geholfen werde. Oeffne die Augen der Blinden, und führe sie auf den Weg, den sie nicht wissen. Nimm hinweg die Decke, so vor ihren Augen und Herzen hänget, daß sie dich, und den du gesandt hast, Jesum Christum erkennen. O Herr, der du vor Zeiten deines Volkes Israel Elend in Egypten
\quellseite{217}
angesehen, und ihr Seufzen und Wehklagen gehöret und erhöret hast, solltest du nicht auch hören das Seufzen der Elenden, die vom Antichrist mit unerträglicher Last gedrückt und gedrängt, ge= plagt und geängstigt werden? Du siehest ja und schauest das Elend und den Jammer. Darum stehe auf und schaffe eine Hülfe, daß man getrost kehren könne. Bekehr doch das arme Volk, das von den falschen Lehrern und blinden Leitern gleichsam bezaubert wird, daß sie der Wahrheit nicht gehorchen, sondern hingehen zu den stummen Götzen, wie sie geführt werden, und in ein recht knechtisch Joch sich einsperren, sich schinden, trotzen und in das Angesicht streichen lassen. Ach wenn sie es wüßten, so würden sie freilich auch bedenken, was zu ihrem Frieden dienet. Aber dies ist die Stunde und Macht der Finsterniß, denn ihr Herz ist verstocket und verstarret, und ihre Ohren hören übel und ihre Augen schlummern. Ach HErr, erleuchte sie, daß sie ihre Herzen vor dem heiligen Evangelio nicht länger verstocken, sondern den kräftigen Irrthum, darin sie stecken, erkennen, und von Babel ausgehen, und also Gott die Ehre geben. Bekehr die Abtrünnigen, so noch zu bekehren sind, hole auch etliche unter den Fetten, hole sie aus der Tiefe des Meeres und bringe sie zu deinem Schafstall. Gib, daß sie in ihren Ohren stets die ernstliche Stimme deines Sohnes erklingen lassen: Saul, Saul, was verfolgst du mich? Es wird dir schwer sein, wider den Stachel zu löcken. Erleuchte auch, o HErr, alle diejenigen, so ihrer Vernunft mehr denn deinem Wort glauben, daß sie das Licht sehen, und erkennen, was nicht Licht ist. O Gott, du Stifter und Erhalter der Wahrheit, heilige mich in der erkannten Wahrheit, daß ich deine göttliche Ehre, die Hoheit deines Worts und die Lauterkeit des Glaubens unerschrocken vertrete, und bis ans Ende beständig bleibe, durch JEsum Christum, deinen eingebornen Sohn, in Kraft des Heiligen Geistes. Amen.

\section[272. Gebet wider die Atheisterei und …]{272. Gebet wider die Atheisterei und epikurische Sicherheit}

Hilf, HErr! Die Heiligen haben abgenommen, und der Gläubigen ist wenig unter den Menschenkindern! Es ist keine Treue, keine Liebe, kein Wort Gottes mehr im Lande, sondern eitel Sünde und epikurische Sicherheit, und recht die Zeit, von welcher
\quellseite{218}
dein liebster Sohn gesprochen, daß, wenn er werde kommen, er kaum Glauben auf Erden finden werde. Denn ja die Meisten ihres Wesens kein Hehl haben, sondern rühmen sich ihrer Sünde, wie die zu Sodom, und ist ihnen leid, daß sie es nicht ärger machen können, ja, viele Narren und Thoren sprechen gar in ihren Herzen: es ist kein Gott, sondern ohngefähr werden wir geboren, und ohngefähr fahren wir wieder dahin, daß also zu besorgen, du werdest mit deinem jüngsten Tage, welcher kommen wird wie ein Fallstrick über alle, die auf Erden wohnen, viel tausend solche gottlose Atheisten und Epikurer hinwegnehmen und sie, als ewig Verfluchte, in das höllische Feuer, das bereitet ist dem Teufel und seinen Engeln, von deinem allerheiligsten Angesicht gehen heißen. Nun, lieber Gott! weil ich weiß, an welchen ich glaube, und gewiß bin, daß du, Gott Vater, Sohn und Heiliger Geist, allein wahrer Gott bist, und zwar nicht ein Gott, dem gottlos Wesen gefällt, und wer böse ist, nicht vor dir bleibt: ach, so laß doch nimmermehr solche arge Gedanken in mein Herz kommen, es sei kein Gott, noch weniger so sicher-dahin leben, daß ich dächte bei meinen Sünden: der HErr siehet’s nicht und der Gott Jakobs achtet’s nicht. Denn wo wollte ich doch hingehen vor deinem Geist, und wo wollte ich doch hin fliehen vor deinem Angesicht, da mich nicht deine Hand daselbst führen und deine Rechte mich halten sollte; darum laß mich wohl erwägen, daß wir in dir leben, weben und sind, und wie dein Wort auf meiner Zunge, daß du, HErr, nicht alles wissest, und laß mich daher führen einen guten Wandel, und dir dienen in Heiligkeit und Gerechtigkeit, die dir gefällig ist. Hilf, daß ich verleugne das ungöttliche Wesen und die weltlichen Lüste, und züchtig, gerecht und gottselig lebe in dieser Welt. Steure aber auch durch deine unendliche Macht allem epikurischen und atheistischen Wesen, daß jedermann bedenke, wie du der große HErr seist, der die Sonne und die gesammte Welt gemacht, und alles durch dein kräftiges Wort erhältst und trägst, wie du um uns seiest, wir gehen oder liegen, und sehest alle unsere Wege. Und weil derjenige, der einen Sünder bekehret von dem Irrthum seines Weges, einer Seele vom Tode hilft, und ihm die Menge der Sünden soll bedecket werden, so gib mir Muth und Geist, daß ich nicht laufe mit den Dieben, noch Gemeinschaft habe mit den Ehebrechen
\quellseite{219}
, sondern diesen und allen andern Sündern ihre Sicherheit zu Gemüthe führe, daß sie doch sollen klug werden, und bedenken, daß, der das Auge gemacht hat, alles sehen, der das Ohr gemacht hat, alles hören müsse, auch dermaleinst ans Licht bringen werde, was im Finstern verborgen gewesen, und offenbaren den Rath aller Herzen. Stehe auch mir bei, daß ich all mein Lebenlang dich für Augen und im Herzen habe, daß ich nicht thue wider dein Gebot, sondern bei deiner letzten Zukunft würdig werde zu entfliehen diesem allem, das geschehen soll, und zu stehen vor des Menschen Sohn. Amen.

\section{273. Gebet für den Ehestand}

Heiliger Gott, barmherziger Vater, der du den Ehestand selbst eingesetzt hast, und lässest dir gefallen, daß Mann und Weib in ehelicher Liebe beisammen wohnen, die Welt erfüllen, deine Kirche mehren durch Kinderziehen, dadurch dein heiliger Name ewiglich gepriesen und gelobet werde. Ich bitte dich, weil der Satan, ein Feind aller deiner göttlichen Ordnung, mit aller List diesem Stande hart zusetzet und sich gewaltig dawider legt: du wollest allen christlichen Eheleuten deine Gnade verleihen, daß sie dich, den Stifter ihres Ordens, erkennen, und solchen Trost in ihrem Gewissen haben, daß sie in einem guten und dir angenehmen Stande sind, darin sie dir auch können dienen und selig werden, auf daß sie desto williger mit allem Wohlgefallen thun, was ihrem Beruf zustehet. Hilf, daß sie dich fürchten, und in rechter Liebe und Einigkeit christlich bei einander wohnen, im Schweiß ihres Angesichts ihr Brot essen, ihre Kinder in aller Gottseligkeit auferziehen, dazu auch das liebe Kreuz, so in diesem Stande nicht ausbleibet, in Geduld mit einander tragen. Wehre dem Eheteufel, daß er solch dein Geschöpf und Ordnung nicht verrücke noch verderbe, und nicht etwa die Ehegatten ihres Standes verdrossen und feindselig mache, und sie in sein Netz durch Unglauben ziehe, seinen Samen, allerlei Zwiespalt und Uneinigkeit, durch Verbitterung ihrer Herzen, zwischen ihnen einstreue und menge, auf daß sie nicht wider dein Gebot zu Ehebruch und Unzucht, oder einander zu verlassen, gereizet, und in andere große Sünde, Schande und
\quellseite{220}
Laster verführet und gestürzet werden, oder sonst in Zank, Hader und Zwietracht, mit Verletzung ihrer Gewissen, unglücklich leben, dadurch das Gebet und alle christliche Uebung verhindert, die Kirche geärgert und der Heilige Geist betrübet wird. Hilf, barmherziger Gott, daß alle uneinige Ehegenossen wiederum in deinem Namen versöhnet werden, einander von Herzensgrund verzeihen, ihrer ehelichen Verbündniß und Pflicht sich erinnern, und beisammen wohnen in rechtschaffener Liebe und Treue, dein heiliges Wort mit ganzem Fleiß lernen, und dasselbe in allen Ehren halten, in aller Widerwärtigkeit auf dich hoffen und bauen, dich mit ihrem Gebet und christlichem Wandel heiligen und preisen, damit sie in deinem Dienst und Wohlgefallen erfunden werden, und nach diesem vergänglichen Leben zu allen gottseeligen Eheleuten versammelt werden in dein Reich. Amen.

\section[274. Gebet, wenn man pflegt neue Obrigkeit zu …]{274. Gebet, wenn man pflegt neue Obrigkeit zu wählen}

Allmächtiger, barmherziger Gott und Vater! Du willst, daß dein Volk mit gerechten Richtern und Amtleuten versehen sei, auf daß dein Gesetz, nach alle deinem Wohlgefallen, vollstrecket werde: ich bitte dich von Herzen, gib uns Häupter und Regenten, die nach deinem Willen fragen und geneigt seien mit ganzem Herzen dein Volk also zu regieren, daß dein Lob und Ehre gefördert und gemeiner Friede erhalten werde, durch Jesum Christum. Amen.

\section{275. Gebet für die weltliche Obrigkeit}

Allmächtiger, ewiger Gott, ein König aller Könige, von dem alle Würde und Gewalt herkömmt, der du hast die weltliche Obrigkeit eingesetzt und ihr das Schwert zur Strafe der Bösen und Handhabung der Frommen befohlen, leite alle weltliche Herren und Regenten in deinen Wegen, daß sie ob dem gemeinen Frieden halten, Krieg und Aufruhr von uns abwenden, und was zur Ruhe, Ehrbarkeit und Gerechtigkeit dienet, befördern. Dazu gib ihnen, du Herr der Heerscharen, deinen göttlichen Segen, Mosis Gerechtigkeit, Davids Sanftmüthigkeit, Pinehas Eifer, Simsons Stärke und Salomo’s Weisheit. Mache und erhalte sie beständig im Guten, in Widerwärtigkeit großmüthig, in Angst getrost, in Leid
\quellseite{221}
fröhlich, in Sorgen zuversichtlich, in Versuchung sieghaftig, in Wohlfahrt unverrückt, und in allerlei Thun und Lassen weise, stark und glücklich. Laß ihre Regierung gereichen zur Unterdrückung aller Tyrannen, dieser Zeit zum Frieden, der Christenheit zum Schutz, bevor deinem heiligen Namen zum Preise und ewigen Lobe, durch JESum Christum, unsern HErrn und Heiland. Amen.

\section[276. Gebet für das Gedeihen der Früchte des …]{276. Gebet für das Gedeihen der Früchte des Landes}

HErr, allmächtiger Gott, König des Himmels und der Erde! der du durch deine überschwängliche Güte den ganzen Erdboden mit allerlei Früchten zierest und erfüllest, davon Menschen und Vieh ihre Nahrung haben. Ich bitte dich um deiner Barmherzigkeit willen, du wollest unser Land segnen und benedeien, daß es seine Früchte und sein Vermögen gebe, da wir von uns selber nicht ein Körnlein aus der Erden hervorbringen, oder demselben helfen können, wo du nicht das Gedeihen giebst. Darum verleihe den lieben Früchten ein gnädiges Gewitter, daß sie wachsen und wohl gerathen, behüte sie vor Hagel und Ungewitter, vor Verheerung, vor allem Ungeziefer und Schaden. Laß das Land nicht feiern in deinem Zorn, daß es nicht wüste liege und unfruchtbar bleibe. Schließe den Himmel nicht zu in deinem Grimm, um unserer Sünden willen, daß er nicht werde wie Eisen und die Erde wie Erz, sondern gib uns Frühregen und Spätregen und fruchtbare Zeiten, suche das Land heim und wässere es und mache es sehr reich. Gottes Brünnlein hat Wassers die Fülle. O HErr! laß unser Getreide wohl gerathen, denn also bauest du das Land, tränke unsere Furchen und feuchte unser Gepflügtes mit Regen; mache es weich und segne sein Gewächs. Kröne das Jahr mit deinem Gut, daß deine Fußstaufen triefen von Fett, so werden die Wohnungen in den Wüsten auch fett werden, daß sie triefen, und die Hügel umher werden lustig und die Anger voll Schafe sein, und die Auen werden dicke stehen mit Korn, daß man jauchzet und singet.

So laß dir, lieber Gott, in deinen gnädigen Schutz befohlen sein den lieben Samen und alle Früchte des ganzen Erdreichs, erhalte sie in Frost, Kälte, Eis und Schnee, in Wind, Hitze und Dürre, im Regen und Wetter aller Art, daß sie nicht beschädiget
\quellseite{222}
werden. Behüte uns vor Mißwachs, vor Theuerung, vor Hunger und Kummer. Siehe, wir bekennen unsere Sünden mit reuigem Herzen und schreien zu dir, unserm Gott, du wollest hören im Himmel, in dem Sitze, da du wohnest, und gnädig sein den Sünden deines Volkes, uns nicht verlassen mit leiblicher Nahrung, sondern uns versorgen, auf daß wir in allen Dingen deine göttliche Kraft und milde Hand mit Danksagung erkennen und dich preisen, der du uns reichlich ernährest an Leib und Seele, hier zeitlich und dort ewiglich. Amen.

\section{277. Fürbitte für die Wohlthäter}

O gütiger, frommer Gott! Dieweil ich aus deinem göttlichen Wort höre und erkenne, daß du an aller Undankbarkeit Mißfallen hast, und die Strafe vom Hause des Undankbaren nicht bleiben wird, sintemal eines Undankbaren Hoffnung wird wie ein Reif für der Sonnen Glanz plötzlich zergehen und wie ein unnütz Wasser verfließen. So bitte ich dich, gib mir Gnade, daß ich dir zu jeder Zeit von Herzensgrund für alle deine Wohlthaten danke, auch gegen denen, so mir jemals Gutes gethan, nicht undankbar erfunden werde, auf daß ich nicht in dein göttlich Urtheil und Gericht falle. Derohalben danke ich dir, o Herr Gott, allmächtiger Vater! denn du bist freundlich und deine Güte währet ewiglich. Du bist treu und hältst deinen Bund in tausend Glied. Dir, Gott, dem ewigen König, dem Unvergänglichen, Unsichtbaren und allein Weisen sei Ehre und Preis in Ewigkeit. Ich bitte dich für alle Menschen, die mir Gutthat erzeigt, die mich von Jugend auf erzogen und ernähret, die mich gelehret, gestrafet und zum Besten unterwiesen haben, die mir förderlich gewesen mit Worten und Werken, auch, für welche ich sonst schuldig bin zu bitten, deren Namen und Anliegen dir wohl bekannt sind, du wollest sie dir, lieber Gott, alle lassen befohlen sein, ihnen wiederum Gutes thun an Leib und Seele, hie zeitlich und dort ewiglich. Verleihe mir Gnade und deinen Segen, daß ich ihnen zur Zeit solche ihre Wohlthaten könne vergelten und an ihren Kindern wiederum verdienen. Und wie mir Gutes widerfahren ist, daß ich auch also hinfort andern Leuten räthlich und behülflich sei nach meinem Vermögen, auf daß meine
\quellseite{223}
Blätter nicht verwelken und ich nicht als ein dürre unfruchtbarer Baum ausgerottet werde. Laß mich durch Undank der Leute nicht unwillig, müde und laß werden oder aufhören, Gutes zu thun an jedermann, allermeist aber an den Glaubensgenossen, weil du, himmlischer Vater, auch gütig bist und lässest deine Sonne aufgehen über die Bösen und über die Guten und lässest regnen über die Gerechten und Ungerechten. Diesem deinen Exempel laß uns nachfolgen, auf daß wir uns beweisen als deine rechtschaffenen Kinder, auch unser Licht also leuchte vor den Leuten, daß sie unsere guten Werke sehen, dich, wahren Gott im Himmel, preisen, daß wir mit Wohlthun verstopfen die Unwissenheit der thörichten Menschen, damit dein göttlicher Name unter uns nicht entheiliget und verletzert, noch deine Ehre geschändet werde, der du lebest und regierest, ein wahrer allmächtiger Gott, von Ewigkeit zu Ewigkeit. Amen.

\section{278. Fürbitte für die Feinde}

Allmächtiger, ewiger Gott, der du uns ernstlich befohlen hast durch deinen lieben Sohn, unsern HErrn JEsum Christum, daß wir unsere Feinde lieb haben und denen, so uns beleidigen, Gutes thun und für unsere Verfolger bitten sollen: wir schreien ernstlich zu dir, daß du alle unsere Feinde heimsuchen, ihnen wahre Reue über ihre Sünden gnädiglich verleihen und ihnen mit uns, sammt der ganzen Christenheit, ein freundliches, gottfürchtiges, friedliches, einhelliges Gemüth und Herz geben wollest, durch unsern HErrn JEsum Christum. Amen.

\section[279. Herzliches Gebet und Seufzen um die …]{279. Herzliches Gebet und Seufzen um die Zukunft Christi}

O HErr JEsu Christe, du getreuer Menschenhüter, ich danke dir herzlich, daß du uns mit deiner Zukunft nicht unversehen überraschen, sondern zuvor treulich warnen wollen. Wecke mich, lieber HErr, auf vom Schlafe der Sicherheit, daß ich mich vor Sünden wider das Gewissen fleißig hüte, und als ein kluger Knecht in steter Bereitschaft stehe, und deiner Wiederkunft vom Himmel mit höchster Freude und Verlangen erwarte, als darauf in dieser erbärmlichen Zeit unsre einzige und letzte Hoffnung stehet. Komm, HErr JEsu, eile, HErr, brich ab die bösen Tage, daß nicht der arme
\quellseite{224}
Haufe gar verzage. O du himmlischer Eck- und Grundstein, schlag das Bild an seine Füße, die Thon und Eisen sind, und zermalme sie, und führe uns ein in den neuen Himmel und in die neue Erde, darin Gerechtigkeit wohnen wird. Laß doch einmal deine Heiligen das Reich einnehmen und ewig besitzen. Laß dich dazu bewegen deine göttliche Verheißung und inbrünstige Liebe gegen deine so theuer erworbene Braut, wie auch das herzsehnliche Verlangen der ganzen Natur, und das unablässige Seufzen deiner frommen Christen, und dann so vieler tausend Märtyrer Blut, so um deines Namens willen vergossen worden. Deswegen komm, lieber HErr Jesu, und eile mit deinem Gericht. Laß doch einmal deine glorreiche majestätische Stimme hören, daß Himmel und Erde erbebe. Kein Tag vergeht, wir warten dein und wollten gern bald bei dir sein. Amen, HErr Jesu, Amen.

\chapter{VIII. Buß-, Beicht- und Communion-Gebete}

\section{280. Prüfungsfragen}

\subsection{Erstes Gebot}

Du heiliger und gerechter Gott, der du Alles siehest und hörest, ich komme jetzt für dich, meinen Zustand recht zu prüfen, und zu erfahren, wie es mit mir stehe, ob ich auf dem Wege des Verderbens bin. Ich bin viel zu verderbt, mein Herz viel zu falsch und betrüglich, erforsche du mich selbst, o Gott! prüfe und erfahre mich, wie ich es meine, und gib es mir durch dein Wort und Geist zu erkennen. Habe ich auch oft an dich gedacht, du großer Schöpfer und Erhalter, du Erlöser und Heiligmacher? Habe ich auch deinen Willen und deine Gebote in meinem Thun und Lassen zur Richtschnur gehabt, daß ich immer mich selbst gefragt: Ist dieses auch recht nach Gottes Gebot? — Sind mir etwa je die abscheulichen Gedanken eingefallen: Vielleicht ist kein Gott? — Habe ich nicht vielerlei Abgötter verehret, bald mich selbst hochgehalten,
\quellseite{225}
meinem Willen allein gefolget, und nur meinen Ruhm und Vergnügen gesucht? — Habe ich nicht vor andern Menschen eine größere Furcht gehabt, als vor Gott, ihre Gegenwart mehr gescheuet, als Gottes Gegenwart, mich um ihre Gunst mehr bemühet, als um Gottes Gnade; aus Liebe gegen sie ein mehreres gethan, als ich nicht aus Liebe zu Gott thun würde? — Habe ich nicht mein Vertrauen auf große Sünder und auf Menschengunst gesetzt? — Hat mein Herz in zeitlichen Dingen, in Geld, Gut, Ehre und guten Tagen etwa seine Beruhigung funden, daß es nicht in dem ewigen Gut, von dem alles andere kömmt, seine einzige Lust und Freude gesucht? — Habe ich vielleicht der Zauberer- und Satanshülfe, statt der göttlichen gesucht, und mich vor ihrer Macht mehr, als vor Gottes Kraft gefürchtet? — Habe ich mir Gottes Macht, Allgegenwart, Gerechtigkeit und Gnade recht vorgestellet, um mich zu einer kindlichen Furcht gegen ihn zu bewegen? — Hat meine Furcht vor Gott auch rechte Früchte gehabt, als Gehorsam, Haß gegen das Böse, Reue über die Sünde, Demuth, Geduld, Wachsamkeit und dergleichen? — Habe ich getrachtet, Gott als das höchste Gut zu erkennen und über alles zu lieben? — Habe ich zu dem Ende Gottes Vollkommenheit, Majestät, Gütigkeit, Barmherzigkeit, Macht, Reichthum und Wohlthaten mit Fleiß betrachtet? — Ist meine Liebe etwa bloß um der Wohlthaten und zeitlicher Glückseligkeit willen entstanden? — Habe ich auch die Früchte der rechten, wahren Liebe Gottes bei mir verspüret, als Haß gegen die Sünde, Verachtung der Welt, Geduld, Verlangen nach mehrerer Vereinigung mit Gott, Lob, Danksagung und Freude? — Wie ist mein Vertrauen gegen Gott beschaffen gewesen; ist es ein wohlgegründet und kindliches Vertrauen? — Hab ich es auf Gottes Güte, Wahrheit, Allmacht und Verheißung gegründet? — Hat solches Vertrauen bei mir die rechten Früchte, als Hoffnung, Gelassenheit, Tapferkeit und Beständigkeit gewirket? — Wie ist also meine Erkenntniß von Gott beschaffen gewesen, ist sie eine lebendige oder todte Erkenntniß? Suche ich auch in der Erkenntniß zuzunehmen? Hüte ich mich für allen Irrthümern ernstlich? — Ist es mir in allen meinen Verlangen um Gott und seine Gnade alleine, oder doch zuwörderst zu thun?

\quellseite{226}

\subsection{Zweites Gebot}

Wie habe ich mich gegen Gottes Namen, oder gegen alles das, worinnen sich Gott in seiner Herrlichkeit offenbaret, bezeiget? — Habe ich nicht öfters Gottes und Jesu Namen ohne Ehrerbietung ausgesprochen? — Wie oft habe ich ohne rechte Andacht und Vorstellung göttlicher Majestät gebetet? — Wie ist mein Lesen und Betrachten des göttlichen Wortes beschaffen gewesen? — Habe ich nicht je zuweilen Sprüche heiliger Schrift auf eine ungebührliche Weise angeführet, mich damit sehen zu lassen, oder einen Scherz vorzubringen? — Habe ich bei der Betrachtung und bei dem Gebrauch der Creaturen ehrerbietig auf des Schöpfers Weisheit, Güte und Allmacht gesehen? — Habe ich nach Art frecher Sünder, meinen Mund zu Fluchen, Schwören und Lästern mißbraucht? — Habe ich auf Aberglauben, oder zauberische Dinge gehalten? Bin ich ein Lügner und Betrüger in geistlichen Dingen gewesen, daß ich Gottes Wort wissentlich verdrehet, oder auch durch Heuchelei andere betrogen? Wie ist mein Lob und Dank beschaffen gewesen?

\subsection{Drittes Gebot}

Habe ich den Sabbath mit geziemender Ehrerbietung angesehen, und denselben Gott gefällig zu heiligen gesucht? — Habe ich mein Gemüthe am Sonntage von zeitlichen Sorgen, und meinen Leib von ordentlicher Berufsarbeit ruhen lassen, damit Gott sein Werk in mir haben, und ich in Gott ruhen könnte? — Habe ich zu dem Ende mich früh durch ein herzliches Gebet zubereitet? Bin ich zu rechter Zeit im Hause Gottes erschienen? Habe ich mich da in Gedanken, Worten und Geberden als ein Kind Gottes verhalten, andächtig das Wort, als Gottes Wort, gehöret, herzlich mit gebetet und gesungen? Habe ich in der Kirche nach Art der Weltkinder, mit Kleidern geprahlet, von weltlichen Dingen geplaudert, mich üppig nach andern Personen umgesehen, des Predigers Vortrag nur nach menschlicher Gunst abgemessen, durch sonderbare Stellung vor den Leuten Ruhm gesucht, oder einen großen Theil der Zeit mit Schlafen zugebracht? — Habe ich zu Hause die Predigt wiederholet, oder doch auf das gedacht, was besonders meinen Zustand anging? — Habe ich den übrigen Sonntag mit heiligen
\quellseite{227}
Uebungen, als Lesen, Beten, Singen, Prüfung meiner selbst, mit Betrachtung göttlicher Werke und Wohlthaten, Erneuerung meines guten Vorsatzes u. s. f. zugebracht? — Habe ich auch andere zu unterrichten, zu erwecken und zu erbauen gesucht, gute Gespräche geführt, und armen Betrübten mit Trost und Hülfe beigestanden? Oder habe ich mit der Welt allerlei fleischliche Ergötzungen gesucht, mit sündlichem Scherz und Narrentheidungen, oder auch Spielen, die Zeit hingebracht, in Trägheit, Uebermaß Essens und Trinkens, oder üppiger Gesellschaft mein Vergnügen am Tage des HErrn gesucht? — Ist es mir ein rechter Ernst gewesen, Gott zu dienen, und recht erbaut und gebessert zu werden?

\subsection{Viertes Gebot}

Wie bin ich meinen Eltern, meinen Lehrern und meinem HErrn begegnet? Habe ich das ihnen angehängte Bild Gottes demüthig erkannt und geehrt? Habe ich ihnen alle Liebe, Treue, Gehorsam und Dankbarkeit zu erweisen mich bemühet, auch für ihre leibliche und geistliche Wohlfahrt Gott angerufen?

\subsection{Fünftes Gebot}

Habe ich nicht öfters eine Feindschaft gegen meinen Nächsten geheget, ihn in meinem Herzen gehaßet, beneidet, sein Unglück gerne gesehen und auf Rache gedacht? — Wie oft ist mein Zorn aus den Schranken gegangen? — Habe ich nicht meine Gebehrden ungebührlich in Zorn verstellet, feindselige, höhnische, kränkende Worte und wohl gar böse Wünsche herausgestoßen? — Habe ich auch vielleicht meine Feindschaft und Schalkheit verborgen, und meinem Nächsten durch verstellte Freundlichkeit und Höflichkeit meuchlings zu schaden getrachtet? — Ist nicht zuweilen meine Wuth so weit gegangen, daß ich meinem Nächsten an seinem Leibe und Gesundheit geschadet? Habe ich meinen Feind gespeiset und getränket und für ihn gebetet? — Habe ich etwa durch Unbarmherzigkeit oder Verrätherey, oder Nachlässigkeit meinen Nächsten in Leibes- und Lebensgefahr kommen lassen? — Bin ich nicht je zuweilen mit meinen Reden und Thun den Unschuldigen und Einfältigen ärgerlich worden, und habe also gewissermaßen ihre Seele
\quellseite{228}
verwundet? — Habe ich getrachtet, gegebenes Aergerniß wieder aufzuheben, und dem Nächsten für die geschehene Beleidigung möglichst genug zu thun, und ihn wieder zu besänftigen? — Habe ich meinem Nächsten in seiner Noth, Gefahr und Armuth mögliche Hülfe geleistet, oder doch mit Rath und Trost beigestanden? — Habe ich in meinem Herzen Barmherzigkeit und Liebe, in meinen Geberden Freundlichkeit, und in meinen Worten ein holdselig Wesen gehabt? Habe ich gegen meine Beleidiger ein langmüthig und zur Versöhnung williges Herze geheget, und zwischen andern Friede zu machen gesucht? — Habe ich mir meines Nächsten Bestes im Geistlichen und Leiblichen angelegen sein lassen?

\subsection{Sechstes Gebot}

Habe ich mich der Keuschheit und Reinigkeit Leibes und der Seele ernstlich beflissen, damit beide zum Dienste Gottes und des Nächsten tüchtig sein möchten? — Habe ich mich der rechten Mittel der Keuschheit gebraucht, als andächtig Gebet, Mäßigkeit, Fasten und dergleichen? — Habe ich nicht öfters böse und unreine Lüste im Herzen geheget, mich daran belustigt, auch ihnen wohl gar eine Zeitlang nachgegangen? — Sind nicht je zuweilen meine Geberden unverschämt und leichtsinnig gewesen? — Habe ich nicht öfters meine Ohren unzüchtigen Gesprächen, und meine Augen unzüchtigen Personen, Gemälden, oder andern ärgerlichen Sachen geliehen? Auch wohl selbst andere durch zweideutige ungezimemende Worte geärgert? — Habe ich mich ernstlich vor aller Unreinigkeit und Hurerei, wie sie Namen haben mag, gehütet? Müßiggang, Völlerei, Trunkenheit, böse Gesellschaft und lüderliche Kleidertracht vermieden?

\subsection{Siebentes Gebot}

Habe ich mich jederzeit der Gerechtigkeit gegen meinen Nächsten, in Ansehung seiner Güter, beflissen, so daß ich einem jeden das Seine gegeben, einem jeden das Seine gelassen, auch durch Rath, That und Abwendung alles Schadens erhalten helfen? — Habe ich das Gefundene, oder auf andere unrechte Weise erlangte Gut seinem rechten Herrn wieder zugestellt, oder auch sonst davon los zu machen getrachtet? — Habe ich nicht oft nach meines Nächssten
\quellseite{229}
Gute verlangt, durch betrügliche Worte solches gesucht, oder es sonst auf eine subtile Weise an mich gezogen? Habe ich der Obrigkeit das Ihrige treulich entrichtet, in meinem Amte und Berufe redlich und fleißig gearbeitet, die mir erzeugten Wohlthaten, oder auch gerechte Almosen dankbarlich und rechtschaffen angewandt? Habe ich meinen Nächsten im Handel und Wandel durch falsche Maaß und Gewichte, oder lüderlichen, unbilligen Preis und beißende Zinsen vervortheilet? — Habe ich klüglich hausgehalten, mich vor Schulden und Bürgschaften gehütet, oder doch das, was ich schuldig gewesen, ehrlich zu bezahlen getrachtet? — Habe ich den Müßiggang, Verschwendung und Pracht geliebet, oder mich durch unerlaubte Mittel, als Spielen, Ueppigkeit, Narrentheidung und Schmeichelei zu erhalten gesucht? — Habe ich etwa den, so mir unterthan gewesen, gedrückt und ausgesaugt, oder den Arbeitern und Gesinde ihren Lohn verstümmelt und entzogen? Wie habe ich mich gegen die Armen verhalten, habe ich eine recht aufrichtige Barmherzigkeit in der That erwiesen?

\subsection{Achtes Gebot}

Wie habe ich mich gegen meines Nächsten Ehre und guten Namen verhalten, habe ich ihn etwa durch Verrätherei, Verleumdungen, Schimpf- und Stachelreden gekränket, und in seinem Wohlstand gehindert? Oder durch falsche Complimente betrogen, und überhaupt in Beurtheilung seines Lebens meine Ehre und Vergnügen gesucht? — Habe ich mich allezeit vor bösem Argwohn, vor Falschheit, Lügen, Schmeichelei und unnützem eiteln Gewäsch ernstlich gehütet?

\subsection{Neuntes und zehntes Gebot}

Habe ich meinen aufsteigenden bösen Begierden, als: Neid, Hochmuth, Geiz und dergleichen eifrig widerstanden, und mich dagegen der Vergnügsamkeit und Demuth befliessen? — Bin ich meinem Nächsten bei aller Gelegenheit mit Gebet, gutem Wunsch, Rath und Hülfe zu statten kommen, und hat mich sein Wohlstand erfreuet? — Habe ich mich in allen Dingen, sonderlich in Essen und Trinken, Freude und Traurigkeit gemäßiget? — Habe ich den Streit wider die bösen Lüste, als einer der Christo angehöret,
\quellseite{230}
beständig fortgesetzt, und die rechten Mittel, als: Gebet, Wachsamkeit, Vermeidung der Gelegenheit, sitzliche Erneuerung des guten Vorsatzes u. s. w. gehörig angewendet? — Wohin ist mein größtes Verlangen gegangen, nach himmlischen oder irdischen Dingen? — Wie weit bin ich in der Anziehung des neuen Menschen kommen, und wie habe ich in der Erneuerung zugenommen? Ach, mein Gott, erforsche du mich selbst, prüfe und erfahre, wie ich’s meine, siehe, ob ich auf rechtem Wege bin, und leite du mich auf ewigem Wege. Du kannst allein das Herz, welches ein trotzig und verzagtes Ding ist, ergründen, und allen Betrug der Eigenliebe entdecken. Ach! wie viel Gutes habe ich unterlassen, und wie viel Böses habe ich begangen? Was ich Gutes finde, ist nur äußerlich, da ist selten eine rechte gute Absicht und guter Grund der wahren Tugend. Nun, du starker und eifriger Gott, mache, daß ich den Donner deines heiligen und unveränderlichen Gesetzes hören, vor dem Fluch und Zorn, den ich verdienet, erschrecken, und mich bußfertig zu deiner Gnade in Christo wenden möge!

\section[281. Bußgebet eines nach angestellter Prüfung …]{281. Bußgebet eines nach angestellter Prüfung erwachten Sünders}

Heiliger und gerechter, barmherziger und gnädiger Vater, Sohn und Heiliger Geist, nachdem ich durch deine Gnade die Menge, Größe und Abscheulichkeit meiner Sünden erkannt, und die Tiefe meines Elends und Verderbens eingesehen, so danke ich dir nicht nur für die hierzu verliehene Gnade und heilsame Wirkung, sondern lege auch hiermit ein aufrichtiges Bekenntniß meiner Sünden und Mißhandlungen vor dir, o Vater, ab. Ich bekenne, daß ich nicht nur durch die Erbsünde deines Ebenbildes beraubet und an Seele, Leib und allen Kräften schrecklich verderbt und daher zu allem Guten leider ganz ungeschickt und träge, und hingegen zu allem Bösen willig und geneigt bin, sondern, daß ich auch, da ich dies mein Fleisch sammt den Lüsten und Begierden kreuzigen und den alten Menschen ausziehen und tödten sollen, ich demselben in mir bisher oftmals Raum, Macht und Herrschaft gelassen und dadurch in viele, ja unzählige, wie auch große, wirkliche Sünden gefallen, eine Sünde mit der andern gehäufet, eine Last nach der
\quellseite{231}
andern auf mich geladen habe, mithin die Größe meines Elendes vergrößert und in das Verderben tiefer hinein, ja, dem Reich der Höllen immer näher gekommen bin. Ich bekenne, daß ich wider beide Tafeln deines heiligen, mit Donnern und Blitzen gegebenen Gesetzes und wider die dir und meinem Nächsten schuldige Liebe vielfältig gehandelt: dich, den ich, daß du allein wahrer Gott bist, aus deinem Wort erkennen, nach deinem Willen ehren, ja, den ich um seiner Hoheit willen über alles fürchten, um seiner Güte willen über alles lieben, und dem ich um seiner Wahrheit willen über alles vertrauen sollen, dich, o höchster Gott, sage ich, habe ich verachtet, dich nicht als meinen Herrn gefürchtet, nicht als meinen Vater geliebet noch vertrauet und in Summa, deine Herrlichkeit und Ehre, so ich dir darbringen soll, in Eitelkeit durch meinen fleischlichen und irdischen Sinn, mit unheiligem Trachten nach dem, was auf Erden ist, verwandelt: Deinen allerheiligsten Namen, vor dem alle Engel und Seraphim sich scheuen, deinen hohen Namen, den ich mit Gebet- und dankvollem Lob verehren sollen, habe ich auch wohl mit unseligem Fluchen und Schwören, mancherlei Aberglauben und Narrentheidungen, vielfältig gemißbrauchet, ohne Ehrfurcht oft genennet und ohne Andacht angerufen. Wie oft ist der, deinem Dienst besonders geweihte Tag von mir entheiliget, zu eiteln und sündlichen Verrichtungen oder Belustigungen angewendet und dein heiliges und theuerwerthes, ja, seligmachendes Wort durch meine Nachlässigkeit im Hören, durch Unglauben im Herzen und durch Widerstreben im Leben, schnöde verachtet worden. Mit Schmerzen gedenke ich an die Missethat, so ich an Eltern, Obern und andern mir von dir Vorgesetzten, anstatt des Gehorsams und Verehrung mit Verachtung, Eigenwillen und Widersetzlichkeit begangen. Und wie soll ich leugnen, daß ich an meinen, mir zur Liebe und Hilfe befohlenen Nächsten durch Lieblosigkeit, Haß, Unversöhnlichkeit, Neid, unfreundliches Bezeigen, bittere und empfindliche Worte, zugefügtes Leid und oft versagte Wohlthaten zum Mörder worden bin? Da ich, nach meiner unordentlichen Eigenliebe, durch Uebermaß in leiblicher Arbeit oder in Speise und Trank, im Zorn oder andern Gemüthsbewegungen, an mir selbst mich vergriffen und versündigt habe. Was für unreine und unzüchtige Begierden in
\quellseite{232}
meinem Herzen sich oft ereignet, welche auch wohl in ärgerliche Worte, Mienen und Thaten ausgebrochen, ist dir, o Herzenskündiger, besser, als mir bekannt. Der du auch am besten weißt, wie oft ich meines Nächsten Gut nicht nur begehrt, sondern auch wohl durch Unrecht an mich gebracht, ihm zu seiner Nahrung und leiblichem Fortkommen mehr hinderlich, als förderlich gewesen. Ja, der du alle übrige wider meinen Nächsten ausgesprochene unrichtige und unbegründete Urtheile oder andere empfindliche Verleumdungen und Lästerungen in das Buch deiner Allwissenheit aufgezeichnet hast. Und wie viel sind nicht andere Sünden von mir begangen worden, die ich nicht alle zählen noch erzählen kann? Maßen ihrer mehr sind, als des Sandes am Meer. Wie ist nicht manche Sünde bei mir zur Gewohnheit worden und wie viele sind vor meinen Augen und Herzen verborgen? Mein Leben ist bisher ein beständiges Sündigen, mein Denken, Dichten und Verlangen des Herzens ist eitel und unheilig, meine Reden unnütze und ärgerlich und mein Thun und Lassen ein beständiges Abweichen von dir und deinen Geboten gewesen. Ach, ich scheue und schäme mich billig, meine Augen aufzuheben zu dir gen Himmel, mein Gott, da meine Missethat über mein Haupt gewachsen und meine Schuld groß geworden bis an den Himmel. O, wehe mir, daß ich so oft, daß ich so muthwillig und vorsätzlich gesündiget habe, daß ich dabei deiner rufenden, leitenden und bekehrenden Gnade, o Gott, und den Wirkungen, Rührungen und Bewegungen deines Worts und Regierung deines Heiligen Geistes widerstrebet und auf meinen Sündenwegen so lange dahin gewandelt! Darum ist auch mein Herz betrübt und meine Augen sind finster worden. Ich erkenne meine Missethat und meine Sünde ist immer vor mir. An dir allein, o Gott, habe ich gesündiget und Uebel vor dir gethan. Wolltest du mit mir hadern, könnte ich dir auf tausend nicht eines antworten, wolltest du ins Gericht mit mir gehen, so wäre mein Urtheil die Verdammniß und meine Strafe der Tod und das ewige Verderben, das Verderben des Leibes und der Seelen in der untersten Hölle. Denn das ist der gerechte Lohn meiner Arbeit in Sündigen und die Frucht meines Fleißes in Uebel thun: So daß, wenn du anjetzo mein Gebet verwerfen, dein Angesicht von mir wenden und
\quellseite{233}
mich in deinem Grimm tödten wolltest, ich auszurufen und bekennen müßte: HErr, du bist gerecht und deine Gerichte sind gerecht. Wollte ich hier schweigen oder leugnen, so würde doch mein eigen Gewissen, daß du dir zur Gerichtsstube in meiner Seelen erkieset hast, mein steter Zeuge und Richter sein. O! was für Wunder fühle ich in demselben, o! was für Last liegt auf demselben! So, daß ich heute vor Unruhe meines Herzens und gar nicht weiß, wo aus noch ein. Denn vor mir sehe ich den offenen Rachen der Höllen, auf welchen ich im Sündigen zugerennet: Um mich theils meine von mir beides geängerte, und betrübte Brüder, theils die Rotte der höllischen Geister, die mir nach meiner Seele stellen: Und über mir erblicke ich deinen, als ein Feuer brennenden Zorn, dein von mir abgewandtes väterlich Angesicht und deine Hand, so mir Tag und Nacht schwer wird, fühle ich schon. Ich fühle sie, wenn ich an das Verhängniß, so mir begegnet, an die Ruthe, die ich gefühlet und an die Trübfeligkeiten, so ich um meiner Sünden willen bereits erfahre, gedenke. Hier brauset eine und da eine Tiefe. Alle deine Wasserwogen und Wellen sollen über mich gehen und alle Schalen deines Grimms über meinen Leib und Seele, als der Sünde Werkstatt und Wohnung, ausgegossen werden. Ach! HErr, die Angst meines Herzens ist groß, die Furcht vor deinen Gerichten mächtig, der Haß gegen meine Sünden aufrichtig und die Verabscheuung meiner Bosheiten ernstlich. Ich weine, wenn ich bedenke deine Liebe, ich erschrecke, wenn ich erwäge deine Wohlthaten, ich kann mich vor Betrübniß kaum lassen, wenn ich betrachte deine Treue gegen mich, die alle Morgen über mich Sünder neu und so groß gewesen, daß ich ihr allein mein Leben und Erhaltung zuzuschreiben habe. Ich schäme mich meines Undankes so sehr, als dich es reuen möchte, mich, der ich ein so großer Sünder geworden, erschaffen, erhalten und beschirmet zu haben. Die Frucht deines Blutes und Todes, HErr Jesu, ist, mich von Sünden zu erlösen, und mich dir zu erwerben, denen habe ich mich ergeben; und von dir habe ich mich losgerissen; mich ewig selig zu machen, und um die Verdammniß habe ich mit Sündigen so fleißig gearbeitet. Wie oft habe ich, o werther Heiliger Geist, vor deiner Erleuchtung meine Augen und vor deinem Wort meine Ohren verschlossen
\quellseite{234}
, vor deiner Gnadenwirkung mein Herz verhärtet, und deiner Leitung und Führung mich, zu meinem größten Schaden, widersetzet. Nun Herr, hier steh ich Armer, der Zorn verdienet hat, gib mir, o mein Erbarmer, den Anblick deiner Gnade. Das erbitte ich um deiner Barmherzigkeit, das hoffe ich, um Jesu Christi willen. Ach! Herr Gott, barmherzig und gnädig, geduldig und von großer Güte und Treue, der du vergießest Missethat, Uebertretung und Sünde. Der du mich Verlornen zu suchen und mich Verdammten selig zu machen, Jesum Christum in die Welt gesandt und ihn in seinem Blute zum Gnadenthron allen armen Sündern vorgestellet hast, stehe an, was er für mich gethan, für mich gelitten und mir erworben hat, und laß mich sein und seines vollgültigen Verdienstes jetzt und allezeit seliglich genießen. Habe ich dein Gesetz übertreten, stehe auf ihn, der es für mich erfüllet: Hab ich Fluch und Tod verdienet, stehe auf das, was er für mich gelitten: Bin ich deiner Gnade und des Heils unwürdig, so laß mir zukommen die Huld, die er mir erworben, den Segen, den er mir verdienet und die Gerechtigkeit, zu der er mir selbst gemacht worden.

So laß mich denn in seinem Blute rein, durch seinen Tod mit dir versöhnet, durch seine Auferstehung vor dir gerecht und um seiner Fürbitte willen gnädiglich erhöret werden. Ich halte mich, Vater, an deine Güte, ich halte mich, Jesu, an deine Genugthuung, ich halte mich, Heiliger Geist, an deine Wahrheit. Vergieb mir, dreieiniger Gott, die Sünde, so mein Heiland getragen: erlaß die Schuld, die mein Bürge bezahlet; halt ein mit der Strafe, die mein Versöhner erlitten: schenk mir die Gerechtigkeit, die er mir erworben. Da sein Gehorsam vor dir gilt, da sein Leiden so unschuldig, da sein Blut um Versöhnung schreit, da sein Tod so theuer geachtet ist und mir mit ihm auferwecket und sammt ihm in das himmlische Wesen versetzet sind, so laß mich gerecht in ihm und dir angenehm in diesem Geliebten sein. Das ist der Grund meiner Hoffnung zu dir, mein Gott, das ist es, was mein betendes Herz so freudig macht. Ach, so erhöre demnach mich und hilf mir, Gott meines Heils, um deines Namens Ehre willen. Errette mich und vergieb mir alle meine mir bekannte und
\quellseite{235}
unbekannte Sünde, um deines Namens willen. Das thue, o Vater, um deiner Liebe willen, das thue, HErr, weil du ein JEsus bist, das thue, Heiliger Geist, weil du mich also beten gelehret und mich zu erhören verheißen hast. Amen! Amen.

\section[282. Klage eines armen Sünders über seines …]{282. Klage eines armen Sünders über seines Herzens Unbußfertigkeit}

Ich elender betrübter Sünder komme zu dir, mein himmlischer Vater, und bringe nichts mit mir, denn eitel Sünde, darf deswegen meine Augen zu deiner hohen allerheiligsten Majestät nicht aufheben, sondern muß mich schämen, daß ich dich so oft erzürnet und deiner Stimme nicht gehorchet. Ach! es ist nichts Gutes an mir, in Sünden bin ich empfangen und geboren, an Natur dermaßen verkehrt und verderbt, daß ich in mir weder Lust noch Liebe zur Tugend und himmlischen Gütern, sondern nur zu zeitlicher Ehre und Wollust, und Zuneigung zu allerlei Bosheit fühle und empfinde. In Sünden habe ich von Kindheit auf gelebt. In Sünden lebe ich noch, so lange ich dieß sündliche Adamsfleisch an mir trage. Ich tröste mich aber, lieber HErr, deiner unermesslichen, unergründlichen, unendlichen und unaussprechlichen Gnade und Barmherzigkeit, die du allen bußfertigen Sündern in deinem Wort versprochen, und mit einem heiligen Eide bestätigt hast. Ich tröste mich des theuren Verdienstes deines eingebornen Sohnes, der auch um meiner Sünde willen dahingegeben, und um meiner Gerechtigkeit willen auferwecket ist. Verbirg, o himmlischer Vater, dein Antlitz von meinen Sünden, und siehe an das Antlitz deines Sohnes, welcher für sich selbst keine Sünde jemals gethan, ja von keiner Sünde gewußt, aber durch seinen allerheiligsten Gehorsam Blut und Tod für alle meine, und für der ganzen Welt Sünde genug gethan hat. Um dieses Heilandes willen nimm mich, o Vater, wiederum zu Gnaden an, und laß sein bitter Leiden und Sterben, und sein vollkommen theures Lösegeld, so er dir für meine Sünde bezahlt hat, an mir armen Sünder ja nicht verloren, sondern kräftig sein. So will ich dich dafür loben und preisen, hie zeitlich und dort ewiglich. Amen.

\quellseite{236}

\section{283. Gebet um Erkenntniß der Sünden}

Lieber Gott, regiere du mich, daß ich mit geistlichen Augen meine angeborne Seuche und Schwachheit erkenne und bekenne, und also zur rechten Erkenntniß Christi geführet und durch deinen Heiligen Geist regieret, gereiniget und geheiliget werde. Amen.

\section[284. Ein Gebet, wie sich ein armer Sünder vor …]{284. Ein Gebet, wie sich ein armer Sünder vor Gott demüthigen soll}

Herr Gott, ich bekenne dir meine Schwachheit und spreche mit dem Propheten: Herr, du bist unser Vater, wir sind dein Thon, du bist unser Töpfer und wir alle sind deiner Hände Werk. Weil du denn sagest, ich sei ein Sünder, will ich dir Recht lassen und gerne dieses sündlich verdammte Wesen, so in meinem Fleische und in der ganzen Natur verborgen, bekennen; auf daß du Recht behaltest in deinen Worten und gepriesen werdest, ich aber zu Schanden werde; auf daß du gerecht und das Leben seiest, ich aber sammt allen Menschen nichts denn Sünde und Tod; auf daß du das höchste Gut seiest, ich aber sammt allen Menschen das ärgste Unglück und Herzeleid. Solches bekenne ich und habe es nicht aus meiner Vernunft, sondern aus deinem Gesetz und Zusagung gelernt; denn meine Vernunft wollte gern ihre Untugend und gottloses Wesen nicht an den Tag kommen lassen, sondern decken und verbergen, ja auch schmücken; aber es ist mir mehr gelegen an dem, daß deine Ehre bestehe und vermehret werde. Amen.

\section[285. Gebet um wahre Buße und rechtschaffene …]{285. Gebet um wahre Buße und rechtschaffene Bekehrung}

O Herr Gott, himmlischer Vater, der du deine Hände den ganzen Tag gegen uns Ungehorsame ausstreckest, und allen Menschen an allen Enden und Orten ganz ernstlich geboten hast Buße zu thun, auch noch täglich uns aus deinem heiligen Worte durch die Stimme deiner Knechte und Mundboten zurufen lässest: Bekehret euch! bekehret euch! Kehret um! kehret um, die ihr abgewichen seid, ein Jeglicher von seinem bösen Wege und von seinem bösen Wesen, und dazu mit einem heiligen Eide betheuert, daß du keine Lust noch Gefallen habest am Tode des Sterbenden, sondern daß
\quellseite{237}
sich, der Gottlose bekehre und lebe, ja, daß es dein gnädiger und ernstlicher Wille sei, daß allen Menschen geholfen werde, und zur Erkenntniß der Wahrheit kommen; deswegen du auch der Menschen Sünde versiehest, daß sie sich bessern. Sintemal es aber nicht liegt an jemandes Wollen oder Laufen, sondern an deinem Erbarmen, und daher wahre Buße und Bekehrung nicht in unserm natürlichen Vermögen stehet, sondern du allein in uns wirkest, Beides, das Wollen und das Thun nach deinem Wohlgefallen: so bitte ich dich herzlich, du wollest durch deinen Heiligen Geist meine Seele und mein Gemüth erleuchten und kräftig bewegen, daß ich mich jetzt und allezeit von ganzem Herzen zu dir bekehre. Und demnach ich die leidige Sünde, sonderlich die angeborne Erbsünde, nicht sehen noch erkennen kann, ohne durch den zarten Spiegel deines heiligen Gesetzes, so erleuchte mein finsteres Herz, daß ich mich selbst in solchem Spiegel der zehn Gebote recht anschaue, und daraus nicht allein meine in Grund verderbte, böse, feinmüthige Natur, sondern auch die daraus herrührende, innerliche und äußerliche, würkliche Sünden demüthig erkenne. Zeige mir, o Herr, meine vielen und mannigfaltigen, verborgenen Fehler und Gebrechen, ja meine unzähligen Sündenschulden; und stelle mir recht vor Augen den grausamen Jammer, darein ich durch die Sünde gerathen, daß ich derentwegen von Natur sei ein Kind des Zorns und der höllischen Verdammniß, da doch der Sünden Sold anders nichts ist, als der zeitliche und ewige Fluch und Tod. Verleihe mir auch, gütiger Vater, eine rechte ernstliche Reue über meine großen und schweren Sünden, daß ich mir dieselben von Herzen lasse leid sein, und vor deinem unerträglichen Grimm und feuerbrennenden Zorne, den ich durch die Sünde auf meinen Leib und meine Seele geladen, mich herzlich fürchte und entsetze. Dein Wort ist wie ein Feuer und wie ein Hammer, der Felsen zerschmeißt. Ach, lieber Herr, zerreiße und zerschneide durch solch dein Wort hier in der Gnadenzeit mein steinernes und felsiges Herz, auf daß es fein mürbe und geartet und den Sünden ernstlich feind werde. Denn die Opfer, die dir gefallen, sind ein geängsteter Geist; ein geängstetes und zerschlagenes Herz wirst du, Gott, nicht verachten. Du wollest mir auch deinen eingebornen Sohn, als den einigen Gnadenthron,
\quellseite{238}
durchs Wort des heiligen Evangeliums recht zu erkennen geben, daß ich unter meiner schweren Sündenlast nicht verzage, sondern mich des theuern Verdienstes und vollkommenster Genugthuung meines liebsten Heilandes herzlich tröste, und also durch einen wahren, lebendigen Glauben an ihn mich wieder aufrichte. Ach Herr, mein Gott, der du Buße für die Sünde annehmen willst, lehre mich erkennen, wie viel es gekostet hat, meine Seele zu erlösen, damit ich ferner die Sünde in meinem sterblichen Leibe nicht herrschen lasse, ihr Gehorsam zu leisten in seinen Lüsten, sondern daß ich durch deine Gnade rechtschaffene Früchte der Buße thue. Amen.

\section[286. Gebet, wenn man in rechte Erkenntniß der …]{286. Gebet, wenn man in rechte Erkenntniß der Sünde kommt}

Lieber Gott, ich bekenne mich vor dir einen großen Sünder, und die zehn Gebote treiben mich und stoßen mich zur Höllen; aber das lehret mich das liebe Evangelium, daß dieses für die höchste Weisheit zu achten sei, wissen und glauben, daß du also gesonnen seist und ein solch Reich durch Christum gestiftet habest, daß du wollest gnädig sein und helfen den armen verdammten Sündern. So knüpfe ich nun an einander in ein Wort und Bekenntniß: ich bin ja ein Sünder, aber doch ist mir Gott gnädig; ich bin dein Feind, aber du bist nur mein Freund; billig würde ich verdammt, aber doch weiß ich, daß du mich nicht willst verdammen, sondern selig und zu einem Erben im Himmel haben. Ja, das willst du, das hast du mir predigen lassen und befohlen zu glauben, um deines lieben Sohnes willen, den du für mich gegeben hast. Amen.

\section[287. Demüthige Klage eines Gefallenen, aber …]{287. Demüthige Klage eines Gefallenen, aber durch Gottes Gnade Aufgemachten}

Ach du heiliger, gerechter und barmherziger Gott, ich klage und bekenne dir mit reuendem, zerbrochenem, zerschlagenem Herzen und zerknirschtem Geist meine Uebertretung und Missethat. Ach Herr! viel sind meiner Sünden, wie Sand am Meer, sie gehen über mein Haupt, und wie eine schwere Last, sind sie mir zu schwer worden, wie groß ist meine Missethat, daß sie gen Himmel steiget
\quellseite{239}
und schreiet; willst du mit mir rechten, so kann ich auf tausend nicht eins antworten. Ach, welch ein unreiner Massa und Klumpen bin ich, wie bin ich durch das tödtliche Gift der Erbsünde so jämmerlich an Seele und Leib verderbet. Siehe, ich bin aus sündlichem Samen gezeuget, meine Mutter hat mich in Sünden empfangen, vom Haupt bis auf die Fußsohlen ist nichts gesundes an mir. Ach Herr! wer will einen Reinen finden, da keiner rein ist. Ach, ich bin ein böser Zweig aus einem giftigen Baum, alle meine Kräfte sind verderbet, mein Verstand ist verfinstert, mein Wille ist dir widerspenstig, ich erkenne dich nicht recht, ich liebe dich nicht von Herzen, ich vertraue dir nicht gänzlich, meines Herzens Dichten und Trachten ist böse von Jugend auf immerdar; wie ein Brunn sein Wasser quillet, so quillet mein Herz die Sünde, da gehet heraus Verachtung, Lästerung, Hoffart, Lügen, eigene Ehre und Liebe, Ungehorsam, Feindschaft, Zorn, Nachgier, Ungeduld, Unzucht, Ungerechtigkeit, Geiz, allerlei böse Lüste, über welche du, o du gerechter Gott, das Urtheil gesprochen hast, die das thun, sollen das Reich Gottes nicht erben. Ach mein Gott, mein Schöpfer und Erlöser, ich klage dir, daß ich zwar nach deinem Bilde geschaffen, aber dasselbe verloren habe, und bin worden das Bild des Satans, das ganze Reich des Satans ist leider! in mir mit aller Bosheit und Untugend, und bin ein Werkzeug des Satans worden. Ach Herr, mein Gott, mein Elend ist größer, denn ich dir es klagen kann! Ich bin ein Greuel in allem meinem Wesen und Thun, alle meine Gerechtigkeit ist wie ein unsaftig Kleid, ich verwelke in meinen Sünden, wie die Blätter, meine Sünden werden mich auch noch dahin führen wie ein Wind. Ach Herr! ich habe alle deine Wohlthaten mißbraucht, mit allen meinen Gliedern und Kräften Leibes und der Seele dir widerstrebet, mit Leib und Seele habe ich der Sünde und dem Teufel gedienet. Ach, ich bin aus deiner Gnade gefallen in deinen ewigen Zorn, aus der Seligkeit in die ewige Verdammniß, aus dem Leben in den ewigen Tod, aus dem Himmel in die unterste Hölle. Ach Herr! ich bin der verlorne Sohn, der seine Güter so schändlich umgebracht hat, ich bin nicht werth, daß ich dein Knecht heiße; ich wäre werth, daß du mich aus deinem Hause
\quellseite{240}
tiefest, und nimmermehr wieder aufnähmest, ich bin der Knecht, der dir, meinem HErrn, zehntausend Pfund schuldig ist, und habe in Ewigkeit nicht zu bezahlen. Ich bin der Mann, der unter die Mörder gefallen, die haben mich ausgezogen, und des schönen Bildes Gottes beraubet, und haben mich an Leib und Seele so sehr verwundet, daß mich kein Mensch heilen kann; darum hilf du mir, mein Gott und HErr, schaffe mir Beistand in dieser Noth, denn Menschenhilfe ist hier kein nütze. Bekehre du mich, HErr, so werde ich bekehret, hilf mir, so ist mir geholfen, denn du bist mein Ruhm. O HErr, wer kann das böse Herz ändern, und ein neues schaffen? Du, HErr, allein, du getreuer Schöpfer in guten Werken. Wer kann die tiefen und abscheulichen Sünden-Wunden heilen. Du, HErr, und dein Wort, welches alles heilet. Nimm weg das böse steinerne Herz, und schaffe in mir, Gott, ein neues Herz, und gib mir einen neuen gewissen Geist. Verworf mich nicht von deinem Angesicht, und nimm deinen Heiligen Geist nicht von mir. Tröste mich wieder mit deiner Hülfe, und der freudige Geist enthalte mich. Ich bin das verlorne Schaf, wenn du mich nicht suchest, so irre ich ewiglich, und bleibe ewig verloren. Nimm mich auf deinen Rücken und bringe mich zu deinem himmlischen Schafstall. Gib mir eine göttliche Traurigkeit, die da wirke eine heilsame Reue zur Seligkeit, die niemand gereuet, und daß mein Herz durch den Glauben bereitet werde, fähig zu werden deines Trostes, und der hochwürdigen Vergebung der Sünden. Gib mir einen zerknirschten Geist, ein zerbrochenes, reuendes Herz, laß meine Augen Thränenquellen werden, meine Sünden herzlich zu beweinen mit Petro und mit Maria Magdalena. O mein HErr und Gott, wenn ich Tag und Nacht weinete, könte ich mein Elend nicht genug beweinen, meine Sünden nicht genug beklagen, meinen Ungehorsam nicht genug beseuften. Ach mein Gott! es reuet mich nicht so sehr, und thut mir nicht so wehe, daß ich den Tod, Hölle und Verdammniß verdienet habe (denn das habe ich billig verdienet), als daß ich dich, meinen Gott, meinen Schöpfer, meinen lieben Vater, meinen Erhalter, meinen Erlöser, meinen Seligmacher, meinen getreuen Hirten, meinen einigen, wahrhaftigen Tröster, so hoch, so oft, so vielfältig, so gröblich, so schändlich beleidiget
\quellseite{241}
und erzürnet habe, und bin dir so undankbar gewesen für deine große Liebe und Treue, Wohltat und Barmherzigkeit. Ach, wie werde ich bestehen vor deinem Angesicht, vor deinem gerechten Gericht, vor allen heiligen Engeln und Auserwählten? Ach, wo soll ich hingehen vor deinem Geist, wo soll ich hinfliehen vor deinem Angesicht? Deine rechte Hand wird mich doch an allen Orten finden und halten. Ich kann und weiß nicht mehr, denn dies einige: Vater! ich habe gesündiget in den Himmel und vor dir, und bin nicht werth, daß ich dein Kind heiße. Ach, ich habe mich der Kindschaft verlustig gemacht. Ich habe aufgehöret, dein Kind zu sein, durch meinen Ungehorsam, und habe all mein Kinderrecht verloren; aber ich glaube, du habest noch nicht aufgehöret, mein Vater zu sein, denn deine Barmherzigkeit ist groß, und währet von Ewigkeit zu Ewigkeit. Darum erbarme dich mein, und siehe nicht an meine Sünde, sondern deine unendliche, ewige Gnade und Barmherzigkeit, um deines lieben Sohnes, Jesu Christi willen. Amen.

\section{288. Gebet um wahren Glauben}

Ich danke dir, mein lieber Gott, daß ich gelernt habe, daß ich meine Sünde nicht soll angreifen mit eigener Buße, oder den Glauben ansagen mit meinen Werken und meine Sünde tilgen. Vor den Menschen dürfte ich es wohl thun, vor der Welt und dem Richter gilt es; aber vor dir, Gott, ist ein ewiger Zorn, dafür kann ich nicht genug thun, ich müßte verzagen. Darum danke ich dir, daß ein anderer für mich meine Sünde ergriffen, sie getragen, dafür bezahlet und gebüßet hat. Das wollte ich gerne glauben, es dünket mich auch fein, recht und tröstlich sein; aber ich kann mich nicht drein ergeben, ich finde in meiner Kraft nichts, das ich thun könnte, ich kanns nicht begreifen, wiewohl ich sollte. Herr, zeuch du mich, hilf du mir und schenke mir die Kraft und Gabe, daß ich glauben möge, wie David der Prophet im 51. Psalm seufzet: Schaffe in mir, Gott, ein reines Herz, und gib mir einen neuen gewissen Geist. Ein neu reines Herz vermag ich nicht zu machen, es ist dein Geschöpf und Creatur. Gleichwie ich die Sonne und Mond nicht machen kann, daß sie aufgehen und hell scheinen am Himmel, so wenig kann ich auch verschaffen, daß das
\quellseite{242}
Herz rein sei und ich einen gewissen Geist, einen starken festen Muth habe, der steif sei und nicht zappele, zweifele oder wadele an deinem Wort. Amen.

\section[289. Appellation eines bußfertigen Sünders …]{289. Appellation eines bußfertigen Sünders von Gottes Richterstuhl zu seinem Gnadenthron}

O lieber HErr, wir können nicht mit dir rechten noch vor Gericht handeln, wir wollen auch nicht von unserer Gerechtigkeit oder Sünde vor dir handeln; denn wenn du, HErr, die Sünde willst zurechnen und uns vor Gericht fragen, ob wir fromm und gerecht sind, so ist's mit uns verloren. Darum so wollen wir von solchem Gerichte zu dem Stuhle deiner Barmherzigkeit appelliren und unsere Zuflucht zu deiner Güte nehmen. Haben wir nun was Gutes gethan, so ist's aus deiner Gnade geschehen. Wende deswegen die Augen deiner göttlichen Barmherzigkeit, nicht der Gerechtigkeit deines gestrengen Gerichts, zu uns; denn wenn du die Sünde uns wirst zurechnen oder dieselbige sehen, so wird keiner selig. Vater unser! der du bist im Himmel, geheiliget werde dein Name, dein Reich komme, dein Wille geschehe, wie im Himmel, also auch auf Erden; unser tägliches Brod gib uns heute, und vergib uns unsere Schuld, als wir vergeben unsern Schuldigern, und führe uns nicht in Versuchung, sondern erlöse uns von dem Uebel, denn dein ist das Reich und die Kraft und die Herrlichkeit in Ewigkeit. Amen.

\section{290. Zuflucht zur Gnade Gottes}

O HErr, ich nehme alle deine Güte, Wohlthat und Gnade, als ein Sünder und verzweifelter Mensch, wie ich gehe und stehe, des ewigen Zorns und höllischen Feuers werth, wenn du solltest nach Recht und Verdienst mit mir handeln. Aber ich sehe nicht meine Sünde, noch was ich verdient habe, sondern dein Wort und ernstlich Gebot an, daß du heißest, vermahnest und dräuest, daß niemand kein Werk vor dich bringe, etwas zu verdienen, sondern aus väterlicher Güte Vergebung der Sünde und allerlei Wohlthat empfange und in der reinen Zuversicht deiner Gnade stehe und bleibe. Amen.

\quellseite{243}

\section{291. Macht der Gnade}

Herr, ich bin deine Sünde, du bist meine Gerechtigkeit; um deßwillen bin ich sicher. Denn meine Sünde wird deine Gerechtigkeit nicht unterdrücken, und deine Gerechtigkeit wird mich keinen Sünder bleiben lassen. Gelobet sei Gott, mein Erbarmer und Erlöser, auf dich allein vertraue ich, darum werde ich nimmermehr zu Schanden werden. Amen.

\section{292. Gebet um Gnade bei Gott}

Ich bin ein armer Sünder, o Gott, vergib mir meine Sünde. Ich will gerne meines Verdienstes geschweigen, schweige du allein deines Gerichtes. O Gott, ich will mit meinen Werken nichts vor dir verdienen, sondern sie allein dahin richten, daß ich damit dem Nächsten diene, und will mich an deine bloße Barmherzigkeit halten. Amen.

\section[293. Gebet um herzliche Erbarmung und …]{293. Gebet um herzliche Erbarmung und Vergebung der Sünden}

Ach Gott, kein Mensch noch Creatur mag mir helfen noch mich trösten, also groß ist mein Elend; denn nicht leiblich noch zeitlich ist mein Schaden. Darum, der du Gott bist, und ewig allein mir helfen kannst, erbarme dich mein; denn ohne dein Erbarmen alle Dinge mir erschrecklich und bitter sind. Nun bitte ich aber dein Erbarmen, nicht das kleine, als du dich zeitlich über die leibliche Noth erbarmest, sondern nach deiner großen Barmherzigkeit, als du dich über die Seelennoth erbarmest, erbarme dich meiner und vergib mir meine Sünde. Amen.

\section{294. Seufzer um Vergebung der Sünden}

Ich glaube und bekenne, daß ich vor Gott ein armer Sünder bin und verdammet sei, und erschrecke davor von Herzen, daß ich meinem Gott je und je ungehorsam gewesen, seine Gebote nie recht angesehen noch betrachtet, viel weniger das größte oder kleinste gehalten habe. Jedoch verzweifele ich nicht, sondern lasse mich zu Christo weisen, Gnade und Hülfe bei ihm zu suchen, und glaube
\quellseite{244}
auch fest, ich werde sie finden. Denn er ist Gottes Lamm, von Ewigkeit dazu versehen, daß er aller Welt Sünde tragen und durch seinen Tod bezahlen soll. Amen

\section{295. ZuBrünstiges Glaubensgebet}

Mein HErr Jesu Christe, du bist ja der einige Hirte und ich leider das verlorene Schaaf, das in die Irre gelaufen ist, und ist mir angst und bange, und wollte gern fromm sein und einen gnädigen Gott und Frieden im Gewissen haben; so höre ich, daß dir ja so bange ist nach mir als mir nach dir, mir ist angst und wehe, bis ich zu dir komme und mir geholfen wird; so bist du in Angst und Sorge und begehrest anders nichts, denn daß du mich wieder zu dir bringest. So komme nun zu mir, suche und finde mich, daß ich also auch komme zu dir und lobe und ehre dich ewiglich. Amen.

\section{296. Gebet um Vergebung der Sünden}

Ewiger, barmherziger Gott, ich armer, elender, betrübter Sünder, komme zu deiner göttlichen Majestät und bringe mit mir nichts anders, denn eitel Sünde, darf deswegen meine Augen nicht wohl aufheben, sondern muß mich vor dir schämen, daß ich dich, meinen Gott, meinen Schöpfer und Vater, oftmals mit so großen und so mannigfaltigen Sünden erzürnet habe, also, wenn du mit mir nach meinem Verdienst wolltest handeln, müßte ich in Ewigkeit ein Kind des Zornes und der ewigen Verdammniß sein und bleiben. Aber, lieber HErr Gott und Vater, ich tröste mich deiner großen Barmherzigkeit, unermesslichen, unergründlichen, unendlichen und unaussprechlichen Gnade, die du mir selbst verheißen, und mit einem theuren Eid und Schwur betheuret und bestätiget, da du sagst: „So wahr ich lebe, ich will nicht den Tod des Sünders, sondern daß er sich bekehre und lebe.“ Ich tröste mich deines eingebornen Sohnes, meines einigen Erlösers, Gerechtmachers, Mittlers und Seligmachers, der mit seinem ganzen Gehorsam, Leiden und Sterben mir wiederum deine Huld erworben, und mir, der ich an ihn glaube, Macht gegeben hat, daß ich dein Kind und Erbe nun und in Ewigkeit heißen und bleiben solle. Darum, so verzage
\quellseite{245}
ich nicht, sondern bekenne dir meine Sünde, innerliche und äußerliche, verborgene und unbekannte, und gebe mich vor dir schuldig als ein großer Uebelthäter und Schächer und ich weiß nichts mehr zu thun, denn daß ich auf die Zusage deines Wortes mit demüthigem Herzen und unerschrocknem Geiste in wahrem Glauben und gewissem Vertrauen auf deine versprochene Gnade, und auf das bittere Leiden deines lieben Sohnes zu deinem Gnadenthron komme, und spreche aus Grund meines Herzens: Abba, lieber Vater, HErr! erbarme dich mein! Ach Gott, sei mir armen Sünder gnädig, vergib mir meine Erb- und wirklichen Sünden, und alles, was ich jemals wider dich gedacht, geredet und gethan habe, wissentlich oder unwissentlich. Es ist ja nichts Gutes an mir. In Sünden bin ich empfangen und geboren, in Sünden habe ich gelebet, in Sünden lebe ich noch, so lange ich dies sündliche Fleisch und Blut an mir trage, und meine Sünden sind stets vor dir. Du siehst tausendmal mehr an mir, ja über tausendmal mehr und größere Sünden, denn ich selbst sehen, wissen und fühlen kann. Ach, HErr, gehe nicht ins Gericht mit deinem Knecht (Magd), denn vor dir ist kein Lebendiger gerecht. Sei mir gnädig um JEsu Christi deines Sohnes, meines HErrn und Heilandes willen, und tilge meine Sünden nach deiner großen Barmherzigkeit, welche größer ist, denn aller Menschen Sünde. Wasche mich wohl mit dem reinen Blut und besprenge mich mit dem Asopen deines Sohnes, das er am Stamm des Kreuzes vergossen hat. Laß mich dein Kind und Erbe sein und bleiben, und ein Gefäß deiner Barmherzigkeit, auch ein Brunnquell und Wohnung des Heiligen Geistes. Ach, HErr, erlöse, tröste und erquicke mich! Hilf mir wiederum auf! Erfreue meinen geängsteten Geist und mein geängstetes und zerschlagenes Herz mit deinem heiligen, freudigen und lebendigmachenden Geist. Das zerbrochene Rohr und glimmende Docht wirst du ja, HErr, nicht gar zerstoßen und auslöschen. Du wirst mich nicht verwerfen vor deinem Angesicht, das bin ich gewiß. Du wirst mich nicht lassen verzagen, sondern mich mit deinem gewissen Geist begnaden, der in mir entzünde und bekräftige wahren Glauben und Trost. Ach HErr! mein Gott und Vater, du getreuer Gott, erhöre mein Gebet, daß ich dein Kind
\quellseite{246}
und Tempel sein und bleiben möge von nun an bis in Ewigkeit. Laß mich nicht fallen in Irrthum und Laster wider mein Gewissen. Errette mich von den Blutschulden und von allem Uebel. Laß mich stets Lust haben, zu deinem Wort, daß ich dasselbe gerne höre, lese und betrachte und andern damit dienen und nütze sein könne, auf daß ich davon rede Tag und Nacht, auch deine Gerechtigkeit und deinen Ruhm stets verkündige, daher will ich dich, meinen Gott und Herrn, und deinen Sohn, meinen Heiland, Jesum Christum, zusamt deinem Heiligen Geist von nun an bis in Ewigkeit lieben, loben, preisen, danken und dienen immer und ewiglich mit allen Engeln und Menschen. Amen.

\section{297. Gebet um Versöhnung mit dem Nächsten}

Gütiger Jesu, ich habe mich durch deine Gnade bereitet und geschickt, die Vergebung meiner Sünden im Beichtstuhl zu holen, ich erinnere mich aber, daß du gesagt und ernstlich befohlen: Wann du deine Gabe auf dem Altare opferst und wirst allda eingedenk, daß dein Bruder etwas wider dich habe, so laß allda vor dem Altar deine Gabe und gehe zuvor hin und versöhne dich mit deinem Bruder, alsdann komm und opfere deine Gabe: Nun ich bin leider mit meinem Nächsten in Mißverstand und Uneinigkeit gerathen, darum so verleihe mir deine Gnade und regiere mich, daß mein Herz zur Versöhnung sich willig finden lasse. Ich erkenne meines Herzens Härtigkeit und bekenne, daß es schwer sei, Fleisch und Blut zu zwingen und sich alles Zorns und Nachgier zu entschlagen. Aber ich hoffe, ja ich bitte, du, Herr, wollest das steinerne Herz aus meinem Fleische wegnehmen und mir ein fleischern, das ist, liebreichs und versöhnlichs Herz geben und einen solchen Menschen aus mir machen, der seinen Feind liebet, die segnet, die ihm fluchen, wohl thut denen, die ihn hassen, und für die bittet, so ihn beleidigen und verfolgen. Ach! laß mich gedenken an das Gericht und die schwere Rechenschaft, die ich werde geben müssen, auf daß ich die Feindschaft, welche den Tod und Verderben suchet, fahren lasse und nach deinem Gebot mich mit meinem Widersacher, noch weil ich auf dem Wege mit ihm bin, versöhne, und die Sonne nimmermehr über meinem Zorn untergehen lasse. Und wie ich in deinem Namen
\quellseite{247}
allen denjenigen verzeihe und vergebe, die mich beleidiget haben, so laß mich auch diejenigen willfährig finden, die ich betrübet und erzürnet, daß sie mir um deinetwillen auch vergeben und verzeihen, womit ich sie beleidiget habe. O Jesu! vergib all unsere Sünde und regiere unsere Herzen, christlich und friedlich mit einander zu leben und deinen Namen zu preisen hie zeitlich, bis wir dich dort loben ewiglich. Amen.

\section[298. Gehet, wenn man zum Beichtstuhl gehen …]{298. Gehet, wenn man zum Beichtstuhl gehen und die heil. Absolution begehren will}

Herr Jesu Christe, mein Erlöser! Du hast deiner lieben Kirche hier auf Erden und deren getreuen Dienern das heilige Amt der Schlüssel hinterlassen mit angehefter Verheißung, was sie in Kraft derselben lösen oder binden werden, das soll auch im Himmel gelöset und gebunden sein, für solche gnadenreiche Mittel sage ich dir einiges Lob, Ehr und Preis, und bitte von Grund meines Herzens, weil ich als ein armer, gebundener Sünder dieses tröstlichen Löseschlüssels bedürftig, und damit ich nicht unter den Fesseln des höllischen Kerkermeisters gehalten werde, du wollest mir denselbigen durch meinen christlichen Beichtvater tröstlich widerfahren lassen, und mich von allen meinen Sünden, um deines Blutes und Todes willen, gnädig entbinden, daneben mir deinen Heiligen Geist verleihen, daß ich die heilige Absolution in rechter Buße und ungezweifelter Zuversicht, gutem Vorsatz, brüderlicher Liebe und Dankbarkeit fasse, und endlich ewig selig werde, Amen.

\section{299. Danksagung nach der Absolution}

Lieber Gott, der du mir neben deinem heiligen Worte gewisse Wahrheit gegeben hast, mich zu versichern, daß meines Herrn Christi Leben, Gnade und Himmel, darinnen er ist, meine Sünde, Tod und Hölle, mir zu gute, alle ganz und gar aufgehoben habe. Solche Verheißung wirst du mir gewiß halten, darauf bin ich so gewiß, daß die Worte, damit mich der Kirchendiener von Sünden losgesprochen hat, so fest und kräftig sein, als ob ich sie von dir, o Gott, selbst gehört hätte. Itzt nun Gottes Wort, wie es
\quellseite{248}
denn ist, so muß und wirds gewiß geschehen und auch ergehen; darauf beruhe ich und in solcher Hoffnung und Vertrauen will ich sterben. Amen.

\section{300. Danksagung nach empfangener Absolation}

Gütiger, gnädiger, barmherziger Gott und Vater, ich sage dir von Grund meines Herzens Lob und Dank, daß du mir armen Sünder jetzt abermal durch deinen Diener alle meine Sünde vergeben, das ewige Leben wieder aufs neue zugesagt und zu Gnaden angenommen hast. Ich bitte dich herzlich, gib mir deinen Heiligen Geist, und schaffe in mir ein rein Herz, damit ich festiglich glaube, daß mir alle meine Sünden durch Christum vergeben sind. Dessen zu einem gewissen Unterpfande will ich den wahren Leib und Blut deines Sohnes Jesu Christi im Brod und Wein essen und trinken zu meiner Seelen Seligkeit. Verleihe mir auch, du getreuer Gott, daß ich mich hinfort besser vor Sünden hüte, denselben feind werde und mein Leben bessere, dazu wollest du mir gnädiglich verhelfen um deines Namens Ehre willen. Amen.

\section{301. Morgengebet eines Communicanten}

Herr Jesu, ich freue mich heute, an diesem Morgen, daß ich mit andern frommen Christen soll in dein Haus gehen, auf deine große Güte, und deinen Leib und Blut mit, in und unter dem gesegneten Brod und Wein genießen. Ach! wie sehnet sich meine Seele nach den Vorhöfen des Herrn und nach deinem heiligen Altar. Wie der Hirsch schreiet nach frischem Wasser, so schreiet meine Seele nach dir, meinem Erlöser. Denn du bist mein Hirt, mir wird nichts mangeln. Ach! Herr Jesu, wie du in mir armen sündhaften Menschen das gute Werk angefangen hast, ach! so vollführe es gnädiglich, durch deines Heiligen Geistes Gnadenbeistand, zu deiner Ehr und meiner Seelen Heil und Wohlfahrt. Liebster Erlöser, ich bin zwar nicht werth, daß du unter mein Dach gehest; ach! aber deine Gnade, die alles vermag, mache mich würdig und recht wohl geschickt, damit die kostbare Speise und Trank deines wahren Leibes und Blutes in dem himmlischen Liebes- und Freudenmahl meine Seele erquicke. Bekleide du mich selbst
\quellseite{249}
mit den Kleidern des Heils, und mit dem Rocke der Gerechtigkeit des wahren Glaubens, damit ich nicht unter den Schein Gästen und Maulchristen erfunden, sondern ein würdiger Gast an deinem Tische sei, und theilhaftig werde aller deiner Gnadenschätze, die du mir erworben hast. Deine Gnade, HErr JEsu, sei und bleibe über mir, wie ich beständig auf dich hoffe. Dein Aufsehen bewahre meinen Odem, mein Leben und Weben. Gib mir auch im Uebrigen, was mir ist Noth, du getreuer Gott, in diesem und jenem Leben. Gib wahre Reue, mein Herz erneue, errette Leib und Seele; auch höre HErr! dies mein Begehr, und laß meine Bitte nicht fehlen, damit ich dir gefällig sei, hier in dieser Zeit, und dort hernach in Ewigkeit. Amen, in JEsu Namen, Amen.

\section{302. Heilige Communionandachten}

JEsu, mein HErr und mein Gott, mein Erlöser und Seligmacher, sei gelobet und gepreiset immer und ewiglich, daß du nicht allein für mich armen, verlornen Sünder den allerschmerzlichsten Tod erlitten, mich von der unzähligen Menge meiner Sünden, von der unerträglichen Last des Zornes Gottes und von der erschrecklichen Verdammniß des ewigen Todes zu erlösen, sondern auch das heilige Abendmahl eingesetzet hast, damit ich hierdurch der Gnade deines himmlischen Vaters, der Vergebung meiner Sünden und der ewigen Seligkeit versichert würde.

Du, o hochgelobter Sohn des Allerhöchsten, bist der große Stifter dieses theueren Sacraments. Hätte ein Mensch auf Erden, wenn er gleich ein Fürst, König oder Kaiser wäre, dieses Abendmahl gestiftet und angestellet, so müßte man zweifeln, ob er auch halten und geben könne, was er darin verheißen hat. Allein so bin ich genugsam versichert, daß es nicht an Einem fehlen werde, worauf du deinen Knecht vertröstet hast: denn du bist der wahrhaftige Gott, in dessen Mund niemals ein einiger Betrug erfunden worden, der allerweiseste HErr, in welchem alle Schätze der Weisheit und der Erkenntniß anzutreffen sind, der allmächtige König, dem alle Gewalt gegeben ist im Himmel und auf Erden. So will ich demnach meine Vernunft mir nichts dabider einreden lassen, sondern an deine Worte und Zusage mich einig und allein halten.

\quellseite{250}

Du hast zur Einsetzung solches Abendmahls die Nacht erwählet, in welcher du von Juda verrathen und in die Hände deiner blutdürstigen Feinde überliefert wurdest, die Nacht, da du dein bitteres Leiden um meiner Sünden willen antratest, und an dem Oelberg mit der ganzen Macht der höllischen Finsterniß bis auf den blutigen Angstschweiß kämpfen mußtest, die Nacht, welche uns den ewigen Tag und das Licht des seligen Freudenlebens im Himmel erworben hat. Was kann ich dagegen hierbei mich anderes erinnern, als deiner blutigen Passion und theueren Verdienstes? Herzliebster Heiland, du hast Brod und Wein hierzu genommen und dadurch die kräftige Wirkung dieses Sacramentes andeuten wollen, daß, gleichwie das Brod des Menschen Leib stärket, und der Wein unser Herz erfreuet, also auch meine matte und betrübte Seele im heiligen Nachtmahl neue Kraft empfangen und herrlich erquicket werden soll. Nun ich will glauben und nicht wanken, daß es sei eine Speise der Kranken, denen ihr Herz von Sünden schwer, und vor Angst ist betrübet sehr. Du hast dabei gebetet und gedanket, mich zu erinnern, daß ich mit gottgegebener Andacht allhier an dem heiligen Altar erscheinen, und der schuldigen Dankbarkeit für deine Wohlthaten nicht vergessen soll. Ei so müsse demnach mein Herz und Auge anjetzo nur auf dich andächtig gerichtet sein, und mein Mund dankgestissen seufzen: Bis willkommen, du edler Gast, der du mich nicht verschmähet hast, du kommst ins Elend her zu mir, wie soll ich immer danken dir! Du hast das Brod gebrochen, und dabei merke ich, wie mein Herz durch wahre Buße zerbrochen und zerkneischt sein müsse, wenn ich ein würdiger Gast an deinem Gnadentische sein will. Ach, so hilf durch deinen Heiligen Geist, daß ich mit einem rechten geängsteten Geist und zerschlagenen Herzen mich dabei einfinde. Du reichest uns unter dem Brod und Wein nicht ein Zeichen deines abwesenden Leibes und Blutes, sondern deinen wahren und wesentlichen Leib selbst, der für mich so viel gelitten hat, so jämmerlich zergeißelt, so grausam zermartert, und endlich gar am Kreuz getödtet ist, dein wahres und wesentliches Blut, welches für meine Missethaten so häufig vergossen ist, daß es mich reinigen möge von allen Sünden, das Blut, so aus deinen eröffneten Wunden geflossen ist, und mir eine ewige Erlösung und
\quellseite{251}
Versöhnung mit Gott erworben hat. Kein Trgeiß soll mich etwas anderes überreden; ob ich gleich mit meinen Gedanken und Sinnen nicht erreichen kann, wie es zugehe, so will ich doch deinem Wort trauen, und daran im geringsten nicht zweifeln. Du befiehlst, ich soll dieses heilige Abendmahl gebrauchen zu deinem Gedächtniß, und dabei mir zu Gemüthe führen, wie du, o heiliges Gottes-Lamm, um meiner Sünden willen dich in den Rachen der grimmigen Löwen ergeben, durch dein Blut für meine Blutschulden gebüßet, und mit deinem Tod mir das Leben und die Seligkeit erworben habest. Nun, HErr, ich gedenke an deine brünstige und unvergleichliche Liebe, die dich gezwungen hat, für mich dein Blut und Leben zu lassen. Ich gedenke an deine übergroße Marter, an deinen höchst jämmerlichen Abschied am Kreuz, und Alles, was du mir zu gute ausgestanden hast. Ich gedenke an meine obliegende Schuldigkeit, was ich dafür dir vergelten solle. JEsu, deine Liebe und Gnade ist größer, denn daß ich sie aussprechen kann. Doch will ich gleichwohl nicht stumm sein, sondern, so lange ich lebe, deinen Ruhm verkündigen und jetzt mit David sagen: Ich dankte dem HErrn von ganzem Herzen, im Rath der Frommen und in der Gemeine. Groß sind die Werke des HErrn, wer ihrer achtet, der hat eitel Lust daran. Was Er ordnet, das ist löblich und herrlich und seine Gerechtigkeit bleibet ewiglich. Er hat ein Gedächtniß gestiftet seiner Wunder, der gnädige und barmherzige HErr. Ihm sei dafür in Ewigkeit Preis und Ehre! Amen.

\section[303. Klage eines bußfertigen Communicanten …]{303. Klage eines bußfertigen Communicanten über das Verderben seiner Natur}

Ach, daß ich Wasser genug hätte in meinem Haupte, und meine Augen Thränenquellen wären, daß ich Tag und Nacht beweinen möchte das unselige Verderben, in welches mich der Fall meiner ersten Eltern gestürzt hat. Denn wo ist das herrliche Ebenbild Gottes? Wo ist die angeschaffene Weisheit, Gerechtigkeit und Heiligkeit? Wo ist die vollkommene Reinigkeit, Keuschheit und Unsterblichkeit? Ist nicht alles verloren zu meinem großen Schaden? Denn darüber bin ich nicht allein der Gleichförmigkeit mit Gott beraubet, sondern meine Natur ist noch dazu auch dermaßen
\quellseite{252}
verderbet, daß ich wohl mit Paulo seufzen muß: Ich elender Mensch! Denn was ist für große Unwissenheit in meinem Verstande und Verkehrtheit in meinem Willen! Was findet sich für unordentliches Wesen in meinen Bewegungen! Wie viel sündliche Lüste entstehen in meinem Herzen! Wie begierig ist mein Fleisch, dieselben zu vollbringen! Wie bereit sind meine Glieder zu den Werken der Ungerechtigkeit! Und was soll ich viel sagen: In meinem Fleische wohnet nichts Gutes. Wollen habe ich wohl durch deine Gnade, zu thun deinen Willen vollkommen; aber das Vollbringen des Guten finde ich nicht. Was für Unvollkommenheit ist in meinen Werken! Was für Mißtrauen und Kleinmüthigkeit spüre ich oft bei mir! Wo ist die beständige Gelassenheit im Kreuz, Verfolgung und Anfechtung! Wie kaltflinnig ist mein Gemüth zum Gebet! Wie zerstreut und unandächtig im Gebet! Wie träge ist meine Seele zu deinem Lobe! Ich bin schuldig, zu mehren allen denen, welche deinen allerheiligsten Namen mißbrauchen, aber wie oft bleibet es nach, und wie vielfältig geschiehet es ohne rechten Eifer! Wie oft wird der Feiertag von mir entheiliget! Wie viel fremde Gedanken finden sich bei dem Gottesdienst! Wie nachlässig und unbereitet ist meine Seele vielmals zum Gehör deines Wortes! Wie leicht werde ich bewegt zum Zorn und wie oft versündige ich mich durch geschwinden, unzeitigen und unmäßigen Eifer! Ach mich elenden Menschen! Wie unvorsichtig und unbehutsam sind meine äußerlichen Sinne! Wie willig zum Bösen und wie verdrossen zum Guten! O wie unbedachtsam ist meine Zunge im Reden und wie unverantwortlich im Schweigen! Wie bin ich noch so gar nicht von mir ausgegangen, denn wie viel Eigenliebe, Eigenehre, Eigenlust und Eigennutzen befindet sich leider in meinem Thun und Lassen. O wie vielen Dingen habe ich noch abzusterben, ehe ich ganz in dir leben kann. Wie viel Neigungen und Reizungen der Sünden finden sich; denn ehe noch die eine getödtet, so ist die andere schon wieder da! Wie bald wird aus einem bloßen Gedanken eine böse Begierde und aus dieser eine sündliche Ergötzung! Wie heftig setzen mir die Lüste zu, daß oft zwischen mir und dem geistlichen Tode nicht mehr als die Einwilligung und Vollbringung der Sünde ist. Daß aber diese nicht erfolget, solches habe ich in Wahrheit nicht mir zuzuschreiben
\quellseite{253}
(denn wer bin ich und was ist meine Kraft), sondern deinem Geiste, o Gott, welcher zu Hülfe kommt meiner Schwachheit. Von den Pflichten gegen Gott, gegen mich und den Nächsten mag ich nicht gedenken; denn wenn ich bedenke, wie vollkommen deine Gebote und wie unvollkommen mein Gehorsam, so entfällt mir aller Muth, denn ich finde mich zu sein, wie die Unreinen, und alle meine Gerechtigkeit wie ein unflätig Kleid. Ja ich bin nicht werth, daß ich dein Kind heiße, ich schweige, daß ich mich sollte würdig halten des Tisches, welchen du für deine Kinder bereitet hast. Weil aber meiner Seele damit nicht gerathen, wenn ich gar davon bliebe, (denn darüber würde ich je länger je elender werden), viel weniger wenn ich sollte unwürdig hinzugehen, (denn das würde zu meinem Gerichte geschehen), so komme ich, mein Gott, zu dem Thron deiner Gnade und bitte in tiefster Demuth: Gnädiger Gott, barmherziger Vater, unerträglich ist dein Zorn, welchen du bräufst allen Sündern; aber unaussprechlich ist deine Barmherzigkeit, welche du verheißest allen Bußfertigen. Darum beuge ich die Kniee meines Herzens vor dir, laß mein Flehen vor dich kommen und vergib mir, vergib mir, o Herr, alle meine Sünde um deines Namens willen, denn du bist gnädig und geduldig und von großer Güte und Treue. Gehe nicht mit mir ins Gericht, denn vor dir ist kein Lebendiger gerecht. Siehe an meine Schwachheit und Gebrechlichkeit. Ach stehe doch, wie ich durch den Fall meiner ersten Eltern verderbet und wie das Böse mich sogar besessen, daß von Natur in mir nichts Gutes wohnet. Ich erkenne meine angeborne Unheiligkeit, aus welcher, als aus einer unreinen Quelle, so viel eitle Gedanken, böse Lüste, unnütze Worte und sündliche Werke wider meinen Willen entspringen, daß Herz, Augen, Ohren, Zunge und andere Glieder des Leibes jämmerlich verunreinigt werden. Ich erkenne auch durch deine Gnade, daß mich dieses schändliche Uebel so ungestalt und häßlich macht, daß ich mir selber mißfalle; wie sollte ich denn dir, du Allerheiligster, wohlgefallen können? Aber, Herr, erbarme dich über mich elenden, schwachen und sterblichen Menschen. Ich hätte zwar keinesweges darüber zu klagen, wenn du mich gleich nach Verdienst strafen wolltest, sondern müßte vielmehr sagen: Herr, du bist gerecht, und deine
\quellseite{254}
Gerichte sind gerecht. Aber laß mir Barmherzigkeit widerfahren, mein Vater, nach deiner Verheißung. Du begehrst ja nicht mehr von einem armen Sünder, als daß er seine Sünde herzlich bereue und in wahrem Glauben Vergebung suche. Derowegen komme ich, mein himmlischer Vater, und bitte im Namen deines lieben Sohnes: Sei mir gnädig, sei mir gnädig nach deiner Güte und tilge meine Sünde nach deiner großen Barmherzigkeit. Handle mit mir nicht nach meinem Verdienst, sondern nach deiner Gnade um deiner Güte willen. Schaffe in mir, Gott, ein reines Herz und gib mir einen neuen gewissen Geist. Verw irf mich nicht von deinem Angesicht und nimm deinen Heiligen Geist nicht von mir. Laß mich nicht sterben, sondern leben, so will ich deinen Namen loben hier in der Schwachheit und dort in der Vollkommenheit. Erhöre mich, du Liebhaber des Lebens, um Jesu Christi, deines lieben Sohnes, meines einigen Mittlers und Erlösers willen. Amen.

\section[304. Gebet um heilige Begier nach dem heil. …]{304. Gebet um heilige Begier nach dem heil. Abendmahl}

Mein Herr Jesu Christe, siehe an meine Unseligkeit, Elend und Dürftigkeit; ich bin dürftig und arm, und dennoch so verdrossen zu dieser deiner Arznei, die du mir zur Vergebung der Sünden und Seligkeit geordnet hast, daß ich mich auch nach den Reichthümern deiner Gnade nicht sehne. Derhalben, o mein Herr, entzünde in mir die Begierde deiner Gnade und den Glauben deiner Zusage, damit ich dich, meinen allerfrömmsten und allergütigsten Gott nicht beleidige durch meinen verkehrten Unglauben und Faulheit, und würdiglich esse und trinke von dem Brod und Wein deines Leibes und Blutes, und durch diese heilsame Speise und Trank gestärket und erhalten werde zum ewigen Leben. Amen.

\section[305. Gebet um würdige Empfangung des heil. …]{305. Gebet um würdige Empfangung des heil. Abendmahls}

Herr Jesu Christe, ewiger Gottes Sohn, ich bin nicht werth, daß ich meinen Mund aufthue, und das hochwürdige Sacrament deines Leibes und Blutes empfange, denn ich bin ein sündiger Mensch, du aber bist der Herr, welchen aller Himmel Himmel nicht können versorgen, wie soll denn der Mensch, der Asche und Staub ist, würdig sein, deinen allerheiligsten Fronleichnam und
\quellseite{255}
dein theures Blut zu genießen. Ich weiß fast wohl und bekenne, daß meiner Sünden sehr viel ist, und ich darum ein unwürdiger Gast zu dieser himmlischen Mahlzeit bin. Wiederum aber glaube ich von Herzen und bekenne mit meinem Munde, daß du mich Unwürdigen mit deiner Gnade kannst würdig machen; denn du bist allmächtig und barmherzig, denn du allein vermagst rein und heilig zu machen, was von unreinem Samen herkommt; aus den Sündern kannst du gerechte, heilige Menschen machen, wenn du aus Gnaden alle unsre Sünde erlässest und uns mit deinem Heiligen Geist erneuerst. Derowegen bitte ich dich, durch deine Kraft und Liebe verleihe Gnade, daß ich geschickt zu deinem Tisch gehe und nicht etwa mit unwürdiger Nießung mich schuldig mache an deinem Leib und Blut, auf daß ich nicht für das Leben den Tod empfange. Gib Gnade, daß ich mich einen armen Sünder erkenne und prüfe, ein reuiges Herz über meine Missethat habe, deinen edlen, zarten Leib und theures, werthes Blut recht unterscheide, meine Vernunft, Sinn und Witz deinem Wort allezeit unterwerfe, auch mein Leben durch deine Hülfe ernstlich zu bessern gedenke, damit ich in diesem hochwürdigen Sacrament nicht allein mit dem Munde deinen Leib esse und dein Blut trinke, sondern auch mit wahrem Glauben, dich, meinen Heiland und Erlöser, annehme, in mein Herz schließe und also Leben und Seligkeit in dir habe; denn du bist das lebendige Brod, welches vom Himmel kommt und gibt den Menschen das Leben. Wer zu dir kommt, den wird nimmermehr hungern, wer an dich glaubt, den wird nimmermehr dürsten; wer dein Fleisch isset und dein Blut trinket, der bleibet in dir und du in ihm und wird nimmermehr sterben. O du süßer Herr, nach dir sehnet sich mein Geist und mein Gemüth. Wie der Hirsch schreyet nach frischem Wasser, so schreyet meine Seele zu dir; meine Seele dürstet nach Gott, nach dem lebendigen Gott; wann werde ich dahin kommen, daß ich Gottes Angesicht schaue. Erfülle mich mit deiner Gnade. Amen.

\section[306. Gebet um würdigen Genuß des heil. …]{306. Gebet um würdigen Genuß des heil. Abendmahls}

Siehe, hier komme ich, mein Heiland, auf deine holdselige Einladung. Laß mich Gnade finden vor deinen Augen, du Liebhaber
\quellseite{256}
des Lebens. Denn wer bin ich, daß du mir so freundlich rufest? Wäre ich gleich heiliger als die Engel, so wäre ich doch darum nicht würdig dieser Mahlzeit; wie sollte ich denn würdig sein, nachdem ich unreiner Lippen bin, ein Greuel und schnöde, es sei denn, daß du mich würdig machest. Derowegen komm mir zuvor mit deiner Gnade und bereite mich. Du hast ja gesagt: Wer zu mir kommt, den wollest du nicht hinausstoßen. So nimm dich doch nun meiner Seelen an nach deiner großen Barmherzigkeit. Gedenke, mein Jesu, daß ich von mir selbst zwar viel Böses, aber nichts Gutes habe, und hilf mir um deiner Güte willen. Siehe doch, wie arm und elend ich bin. Gedenke an deine Treue, du Heiland der Welt, und erfülle mein Herz mit Gnade. Wie gerne wollt ich mit brünstiger Andacht und gebührender Ehrerbietigkeit im völligen Glauben erscheinen, aber wo ist das Vollbringen, wenn ich es von dir nicht erlange, du Anfänger und Vollender meines Glaubens? Darum gib mir, was dir wohlgefällig ist. Nimm gnädiglich von mir, was dir mißfällt, oder bedecke es zum Wenigsten, das bitte ich demütigst, mit dem Rock deiner Gerechtigkeit. Wohl der Seele, die dir angenehm! Eröffne mir die Augen, daß ich sehe die Wunder an deinen Einsetzungen. Erwecke mein Herz, und entledige mich von allen fremden Gedanken, mehre meine Hoffnung, entzünde meine Begierde, mache brünstig meine Andacht, reinige mich von aller Unreinigkeit und heilige mich durch und durch, auf daß ich in freudiger Zuversicht hinzutrete, mit reinem Munde und heiligem Herzen nehme, und mit herzlicher Begierde esse und trinke zu meiner Seele Leben, Heil und Segen. Amen.

\section[307. Gebet eines Communicanten vor Empfang …]{307. Gebet eines Communicanten vor Empfang des heil. Abendmahls}

Herr Jesu Christe, Gottes und Marien Sohn, mein Erlöser und Seligmacher, weil du den Tod für mich gelitten, und von Sünden und ewiger Verdammniß mich erlöst, und selbst gesagt hast: Kommet her zu mir alle, die ihr mühselig und beladen seid, ich will euch erquicken. Ich mit vielen Schmerzen und großer Krankheit beladener Mensch komme zu dir mit zerschlagenem Herzen und demütigem Flehen, du wollest mich auch erquicken, laben,
\quellseite{257}
trösten und stärken auf meinem Lager. Weil du auch in deinem heiligen Abendmahl mit allen Gläubigen ein ewiges Gedächtnis gestiftet, und einen neuen Bund aufgerichtet, darinnen wir deinen wahren Leib und Blut empfangen, essen und trinken sollen, zur gewissen Versicherung, daß du in uns, und wir in dir sein und bleiben, auch von dir uns Niemand scheiden solle, weder Leben noch Tod, weder Teufel noch Gewalt, noch einige Creatur, so bitte ich deswegen von Grund meines Herzens, du wollest mich von meinen Sünden reinigen, meinen Leib und Seele mit diesem deinen heiligen Leibe und Blute laben und erfrischen, stärken und trösten in dieser meiner Krankheit. Allerliebster Heiland! schreibe mich an in die Zahl derer, die würdig zu ihrer Seligkeit dieses theure und kostbare Pfand empfahen, damit du in mir und ich in dir bleibe, und mich weder diese Krankheit, noch auch der Tod von dir scheide, daß hinfort nicht ich, sondern du in mir lebest, und, was ich noch lebe im Fleisch, ich in deinem Glauben lebe, und dein verbleibe, ich lebe oder sterbe. In diesem Glauben erhalte mich, mein HErr JEsu, daß ich durch dich alle Anfechtung möge überwinden, diese väterliche Züchtigung geduldig ertragen, und in gewisser Hoffnung des ewigen Lebens beständiglich ausharren. O HErr, nimm von mir, was mich wendet von dir, nimm mich mir, und ergib mich ganz eigen dir. Laß mich dein sein und bleiben bis an mein letztes Ende, und hernach in Ewigkeit, so will ich auch dich, sammt deinem Vater und Heiligen Geiste, loben und preisen, ohne Aufhören immer und ewig. Amen.

\section[308. Gebet in der Kirche vor der Empfangung …]{308. Gebet in der Kirche vor der Empfangung des heil. Abendmahls}

\subsection{Zu Gott dem Vater}

Ach HErr, allmächtiger, ewiger und barmherziger Gott und Vater unseres lieben HErrn und Heilandes JEsu Christi, ich armer, elender, sündiger Mensch sollte billig meine Augen zu dir nicht aufheben und mich keiner Gnade noch Gutes zu dir versehen, so voller Sünden stecke ich, so schändlich habe ich bisher den Bund, den du in der heiligen Taufe mit mir gemacht hast, übertreten, und sollte
\quellseite{258}
mich billig in die Zahl derjenigen mit einschließen, von welchen du aus gerechtem Zorn sagest: „Treibe sie weg von mir und laß sie hinfahren. Du hast mich verlassen und bist von mir abgefallen, darum habe ich meine Hand ausgestreckt wider dich, daß ich dich verderben will; ich bin des Erbarmens müde.“ Ja freilich habe ich dich verlassen und bin von dir abgefallen, aber, ach du frommer, getreuer, lieber Vater, du willst nicht mit uns handeln nach unsern Sünden, hast auch keine Lust an unserm Verderben, sondern willst vielmehr, daß wir uns bekehren und daß wir leben sollen. Deßhalb hat nun auch dein lieber Sohn, JEsus Christus, nach deinem väterlichen Rath und Willen das hochheilige Sacrament eingesetzt, darinnen er uns allen Gnade, Vergebung der Sünden, Leben und Seligkeit schenkt. Zu dieser deiner himmlischen Mahlzeit bin ich nun auch geladen und komme demnach jetzt als ein armes, hungriges und dürftiges Kind; speise und tränke nun meine arme Seele, du frommer himmlischer Vater, und nimm mich wieder zu Gnaden auf. Laß mich meine Sünde nicht von dir scheiden, denn ich bin deine Creatur, denn deine Hände haben mich gemacht und zum ewigen Leben geschaffen. Siehe mich mit den Augen deiner väterlichen grundlosen Güte und Barmherzigkeit an und vergib mir Alles, was ich bisher wider deine Gebote gehandelt habe, wie du es denn auch wahrhaftig schon allbereit gethan und mir im Beichtstuhl alle meine Sünden durch deinen verordneten Diener verziehen und vergeben hast, darauf du mir nun jetzt ein theures, werthes Pfand, nämlich den wahren Leib und Blut deines lieben Sohnes, JEsu Christi, geben willst. So verleihe mir nun deinen Heiligen Geist und zünde in meinem Herzen einen rechten, wahren Glauben an, damit ich dieses dein heiliges, hochwürdiges Sacrament würdiglich zu meiner Seelen Seligkeit empfange, deiner Gnade dadurch versichert und wider alle Anfechtung gestärket und erhalten werde. Hilf auch, du getreuer, lieber Vater, daß ich in solcher deiner Gnade beständig bis an mein Ende verharre und mich ja nicht wieder muthwillig durch ein neues sündiges Wesen daraus stürze, sondern daß ich forthin ein christliches und dir wohlgefälliges Leben führen, Glauben und gut Gewissen bewahren und ein frommes und gehorsames Kind sein möge, bis ich dermaleinst aus dieser
\quellseite{259}
Schwachheit zu dir anheimkomme und dir in vollkommener Heiligkeit und Gerechtigkeit dienen werde. Das wollest du thun, du frommer, himmlischer Vater, um deines lieben Sohnes, JEsu Christi, meines einigen Erlösers und Heilandes willen. Amen.

\subsection{Ein ander Gebet zu Gott dem Sohn}

Ach du frommer HErr JEsu Christe, wer kann deine grundlose Güte und Barmherzigkeit genug rühmen? und mit was Worten vermag ich doch deine Liebe und Treue gegen mich armen Sünder auszusprechen? Du hast mich nicht allein mit deinem Blut und Tod erlöset, sondern willst auch jetzt im hochwürdigen Sacrament zu mir kommen, in mein Herz einziehen und allda deinen Tempel und Wohnung haben. Ach HErr, ich bin’s ja nicht werth, daß du unter mein Dach gehen sollst, und hätte es wohl verdienet, daß du mich mit Donner und Blitz in den Abgrund der Hölle hin ein schlägest; doch weil du mich aus großer Gnade derselben werth achtest und dich mir gutwillig mit allen deinen Gnadenschätzen anbietest, wohlan, so will ich dich auch herzlich gern aufnehmen und komme demnach jetzund, dich, meinen lieben HErrn und Heiland, mit einem rechten wahren Hunger und Durst zu empfangen. So erbarme dich nun über mich, dein armes Kind, du frommer und getreuer HErr JEsu Christe, und reinige dir mein Herz durch dein heiliges rosinfarbenes Blut, damit es dein heiliger Tempel werde, wohne in mir, und laß mich ja nicht in meinen Sünden verderben. Ach HErr JEsu Christe, ich bekenne und beklage dies auch von Herzen, daß der leidige Satan, weil ich in diesem heiligen und hochwürdigen Sacrament von außen mehr nicht, als nur Brod und Wein, sehe und schmecke, daß er mir ohne Unterlaß durch meine dießfalls blinde Vernunft einraunet, dein wahrer Leib und Blut sei allda nicht gegenwärtig, sondern ich müsse mich mit meinem Glauben in den Himmel hinaufschwingen und dich allda zur Rechten Gottes empfahen. Darum so bitte ich von Herzen, hilf mir mit Gnaden, daß ich in diesem hohen und großen Geheimniß mich fest an dein heiliges Wort halte und mir dasselbe viel mehr gelten lasse, als Alles, was meine Vernunft sagt und urtheilt. Denn dein Wort ist die rechte Wahrheit, das wird mich ja nicht betrügen,
\quellseite{260}
sondern weil du dieses Sacrament testamentweise hast eingesetzt; welches man ja nicht ändern darf, so stärke meinen schwachen Glauben dadurch, auf daß ich mich dessen von Herzen tröste, Alles, was du mit deinem Blut und Tod erworben hast, sei auch mir zu gut geschehen, und hilf, daß ich alsdann solchen meinen Glauben auch mit der That, mit wahrer Gottseligkeit beweisen möge, forthin von Sünden ablasse, frommer werde und mein Leben bessere, damit ich dich nicht wieder auf’s Neue austreibe, sondern dein Tempel und Wohnung sei und bleibe, bis ich dermaleinst aus diesem Jammerthal und Schwachheit zur ewigen Vollkommenheit gelange. Der du lebest und regierest mit deinem himmlischen Vater und Heiligen Geist, wahrer Gott und Mensch, unser Mittler, hochgelobet in Ewigkeit. Amen.

\subsection{Zu Gott dem Heiligen Geist}

Du Heiliger Geist, mit Gott dem Vater und seinem lieben Sohne, JEsu Christo, gleiches Wesens, Allmacht und Herrlichkeit, unser Regierer, Leiter und Führer, dich bitt’ ich von Grund meines Herzens, du wollest mir deine göttliche Gnade verleihen, und helfen, daß ich jetzt das heilige Sacrament würdiglich zu meiner Seelen Seligkeit empfangen möge und demnach zuvörderst einen rechten wahren Glauben in mir anzünden, daß ich den Worten meines lieben HErrn und Heilandes, JEsu Christi, festiglich trauen lerne, und nicht daran zweifele, daß ich wahrhaftiglich unter Brod und Wein seinen wesentlichen Leib und Blut empfange, ungeachtet, wie heftig meine Vernunft dawider streitet, sonderlich aber damit ich mich bei diesem theuren Unterpfand von Herzen erinnere und tröste, daß mein lieber HErr Christus auch für mich sei in den Tod dahin gegeben und daß Alles, was er durch sein bitter Leiden und Sterben errungen hat, auch mir zu Gut geschehen und mich wahrhaftig für meine Person angehe. Darnach so wirke auch durch den Brauch dieser hohen himmlischen Mahlzeit eine rechtschaffene Liebe in meinem Herzen, daß ich auch meinen Nächsten mit allen Treuen meine, ihm, wo mirs möglich ist, diene und alles Gute erzeige, wie mein lieber HErr und Heiland gegen mir gethan hat, und hilf ja, du frommer und getreuer Gott, daß ich auch forthin
\quellseite{261}
mein Leben bessere, mich vor Sünden hüte, frommer werde und dich sammt meinem himmlischen Vater und seinen lieben Sohn JEsum Christum ja nicht wieder aufs Neue erzürne. Das wollest du thun, du allerhöchster Tröster, Fürsprecher und Beistand in allen Nöthen, um deines heiligen Namens Ehre und um meiner armen Seelen Heil und Seligkeit willen. Amen.

Der HErr ist mein Hirte; mir wird nichts mangeln. Er weidet mich auf einer grünen Aue, und führet mich zum frischen Wasser; er erquicket meine Seele; er führet mich auf rechter Straße, um seines Namens willen. Und ob ich schon wanderte im finstern Thal, fürchte ich kein Unglück; denn du bist bei mir, dein Stecken und Stab trösten mich. Du bereitest vor mir einen Tisch gegen meine Feinde. Du salbest mein Haupt mit Oel, und schenkest mir voll ein. Gutes und Barmherzigkeit werden mir folgen mein Lebenlang, und ich werde bleiben im Hause des HErrn immerdar.

\section{309. Gebet vor dem heil. Abendmahle}

Vater im Himmel, der du deines eigenen Sohnes nicht hast verschonet, sondern ihn für uns arme Sünder dahingegeben, erbarme dich meiner, und bereite mich zu einem würdigen Gaste des heiligen Leibes und Blutes deines allerliebsten Sohnes, meines HErrn, damit meine Seele gestärket und meine Sünden vergeben werden.

O HErr JEsu, mein Erlöser und Heiland, wahrer Gott und Mensch, der du dich auch für mich hast ans Kreuz nageln und schlagen lassen, und für meinen Ungehorsam gebüßet, erbarme dich meiner, und hilf, daß ich durch dein Kreuz und Tod, dessen Gedächtniß ich bei Empfängung dieses heiligen Sacramentes begehe, mit deinem himmlischen Vater versöhnet werde.

O Gott, Heiliger Geist, der du unsere Herzen durch dein kräftiges Wort erleuchtest, erbarme dich mein, und stärke mich im rechten wahren Glauben, daß ich in demselben meinen HErrn JEsum einwickeln und ewiglich behalten möge. Amen.

\quellseite{262}

\section[310. Gebet vor der Genießung des heil. …]{310. Gebet vor der Genießung des heil. Abendmahls}

Mein HErr Christe, ich bin gefallen, wollte wohl gerne, daß ich stark wäre; so hast du uns das Sacrament darum eingesetzt, daß wir unsern Glauben dadurch entzünden und stärken, und uns also geholfen werde. Darum bin ich da und wills empfahen. HErr, siehe, da ist das Wort, hier ist mein Gebrechen und Krankheit, so hast du selbst gesagt: Kommt her zu mir alle, die ihr mühselig und beladen seid, ich will euch erquicken, darum gehe ich herzu, und laß mir helfen.

\section{311. Ein Anderes}

Ich bin ein armer Sünder, ich bedarf Hülfe und Trost, ich will hingehen zu des Herrn Abendmahl und mich mit meines lieben Herrn Jesu Christi Leib und Blut speisen und tränken lassen. Denn du, HErr Christ, dies Sacrament drum hast eingesetzt, daß alle hungrige und durstige Seelen gespeiset und erquicket werden. Du wirst mich nicht schelten, viel weniger erwürgen, wenn ich nur in deinem Namen komme, daß ich will gesegnet sein, Hülfe und Trost haben. Ach HErr, ich bin ein großer Sünder, komme derhalben jetzt zu deinem Abendmahl, und will mit dir essen, und zweifle nicht, ich werde dir ein lieber und werther Gast sein. Amen.

\section{312. Klage über Unwürdigkeit}

HErr, wahr ist's, daß ich nicht würdig bin, daß du gehest unter mein Dach, so bin ich doch dürftig und begierig deiner Hülfe und Gnade, daß ich auch möge fromm werden. So komme ich, auf kein anders Verlassen, denn daß ich die süßen Worte gehöret habe, daß du mich mit zu deinem Tische ladest, und sagst mir Unwürdigen zu, ich soll Vergebung aller Sünden haben durch deinen Leib und Blut, so ichs esse und trinke in diesem Sacrament. Amen. Lieber HErr, dein Wort ist wahr, da zweifle ich nicht an, und darauf esse und trinke ich mit dir, mir geschehe nach deinem Willen und Worten. Amen.

\section{313. Gebet, wenn man hinzugehet}

HErr Jesu Christe, du hast das Sacrament deines Leibes und Blutes darum eingesetzt und gelassen, daß man da Vergebung
\quellseite{263}
der Sünden finden soll: so fühle ich, daß ich kein nothdürftig bin; ich bin in Sünde gefallen, und stehe in Furcht und Verzagen, bin nicht kühn, dein Wort zu bekennen, habe soviel und soviel Gebrechen; darum komme ich nun, daß du mich heilst, tröstest und stärkest. Amen.

\section[314. Kurzer Seufzer, wenn man den Leib …]{314. Kurzer Seufzer, wenn man den Leib Christi empfangen will}

HErr JESU Christe, dein heiliger Leib stärke und bewahre mich im rechten Glauben zum ewigen Leben. Amen.

\section[315. Kurzer Seufzer, wenn man das heil. Blut …]{315. Kurzer Seufzer, wenn man das heil. Blut Christi empfangen will}

HErr JESU, dein heiliges Blut stärke und bewahre mich im rechten Glauben zum ewigen Leben. Amen.

\section[316. Danksagung nach dem Genuß des heil. …]{316. Danksagung nach dem Genuß des heil. Abendmahls}

HErr JESU Christe, ich sage dir von Herzen Lob und Dank, daß du mich armen sündigen Menschen abermals von meiner Missethat gewaschen, auch zu wahrer und gewisser solcher Abwaschung und Vergebung aller meiner Sünden mit deinem wahren Leib und Blut gespeiset und getränket hast, und mich als ein besudeltes Kind nach solcher Reinigung in die väterlichen Arme deiner Gnaden und Barmherzigkeit wieder aufgenommen hast, und deinem Vater unsträflich, rein und ohne Makel weisest und fürstellest.

Ich bitte dich von Herzen und aus allen meinen Kräften, du wollest nach deiner Wohlthat noch zum Ueberfluß das dazu thun, und deine Gnade durch die Wirkung des Heiligen Geistes mir gnädiglich verleihen, daß ich solche deine große Güte und Barmherzigkeit möge erkennen, mit Dankbarkeit annehmen und dich von Herzen loben und preisen mein Lebenlang. Du wollest auch ferner durch deinen Heiligen Geist in mir wirken und schaffen, daß ich meinem Nächsten alles, so er wider mich thut, auch von Herzen vergebe, wie du, HErr, mir alle meine große und vielfältige Missethat ganz und gar verziehen und vergeben hast, ja auch ganz und gar ausgelöscht und nimmermehr gedenken willst; daß ich auch meinen
\quellseite{264}
Nächsten lieben und ihm von Herzen alles Gute erzeigen möge, wie du mir, HErr, mehr, als ich dir dafür danken kann, hast widerfahren lassen; damit du, frommer und getreuer Gott, in uns gelobet und gepriesen werdest sammt deinem Vater und Heiligen Geist immerdar. Amen.

\section{317. Eine Andere}

O allmächtiger, ewiger Gott, barmherziger Vater, ich sage deiner armlosen Güte und väterlichen Liebe von ganzem Herzen Lob, Ehr und Dank, daß du mich mit dem wahren Leib und Blut deines lieben Sohnes, meines Heilandes und Erlösers, Jesu Christi, so reichlich gespeiset und getränket hast. Nun kann ich fröhlich rühmen und sagen mit David: Der HErr ist mein Hirte, mir wird nichts mangeln. Er weidet mich auf einer grünen Aue, und führet mich zum frischen Wasser. Er erquicket meine Seele, und führet mich auf rechter Straße um seines Namens willen. Ich habe meinen Heiland genossen, darum ist meine Seele genesen. O wie soll ich dem HErrn vergelten alle seine Wohlthat, die er mir thut? Ich danke dem HErrn von ganzem Herzen, im Rathe der Frommen, und in der Gemeine. Groß sind die Werke des HErrn und wer ihrer achtet, der hat eitel Lust daran. Was er ordnet, das ist löblich und herrlich, und seine Gerechtigkeit bleibet ewiglich. Er gibt Speise denen, die ihn fürchten, er gedenket ewiglich an seinen Bund. Er hat ein Gedächtniß gestiftet seiner Wunder, der gnädige und barmherzige HErr. Ich bitte dich, lieber HErr, du wollest durch deinen Heiligen Geist in mir wirken, wie ich jetztunder das heilige Sacrament mit dem Munde empfangen, daß ich also die gnädige Vergebung der Sünden, Vereinigung mit Christo, und ewiges Leben, in und bei diesem theuren Pfand und Siegel wahrhaftiglich und festiglich begreifen und behalten möge. Hilf mir gnädiglich, daß diese meine Andacht und Gottesfurcht lang währe und nicht als ein Regenwasser und Wolke verschiebe. Behüte mich, daß ich nicht nur mit den Lippen zu dir nahe, sondern mit Leib und Seel, mit allen Vermögen und Kräften dir ganz ergeben sei und bleibe, zu einem Opfer, das da lebendig, heilig und dir wohlgefällig sei. HErr Gott Vater, durch deine Güte führe mich
\quellseite{265}
auf rechter Bahn. HErr Gott, mich heut behüte, sonst möcht ich irre gehn, halt mich im Glauben feste zu dieser argen Zeit, gib, daß ich mich stets rüste zur seligen Ewigkeit. Amen.

\section{318. Noch eine Andere}

HErr JEsu Christe, dir sei immer, hier und in Ewigkeit, Lob, Ehr und Dank, daß du mich armen elenden Sünder mit deinem wahren Leib und Blut so gnädiglich hast erquicket, und bitte von Grund meines Herzens, du wollest nun bei und in mir kräftig sein und bleiben und mich, dein armes Würmlein, mit gnädigen Augen ansehen und annehmen und aus deinen Gnaden-Händen nimmermehr fallen lassen und mit deinem Heiligen Geist hinfort in meinem Beruf und ganzem Leben also regieren, leiten und führen, daß ich nichts wider dich gedenke, fürnehme, rede, oder ins Werk setze, sondern allezeit auf dich, auf dein Wort, oder auf deinen heiligen Willen sehen möge; auch in meinem Kreuze, das du mir auflegest, geduldig und willig sei und wider dich, ja auch wider meinen Nächsten nicht murre, sondern es schlecht bleiben lasse, wie du es machst und wie es dir gefällt, und erwarte der gnädigen Erlösung und der ewigen Freude, die du geben wirst allen Gläubigen, die auf dich trauen und bauen. Hilf, HErr JEsu, daß mein Glaube in mir stark und fest sei, und bleibe, und erhalte mich in rechtem Glauben und Bekenntniß, Demuth, Geduld und Hoffnung bis an meinen letzten Seufzer, daß ich dich hier und in Ewigkeit lobe, und dir danke, der du bist, sammt dem Vater und Heiligen Geist, wahrer allmächtiger, ewiger Gott von Ewigkeit zu Ewigkeit. Amen! HErr JEsu! Amen.

\section{319. Noch eine Andere}

Ach du mein Gott, ob ich wohl ein armer Sünder bin, so bin ich dennoch kein Sünder. Ein Sünder bin ich meiner selbst halben und außerhalb Christo; aber in meinem HErrn Christo und außerhalb mir bin ich kein Sünder. Denn er hat mit seinem Blute alle meine Sünde ausgetilget, wie ich festiglich glaube, derhalben ich auch deß zum Wahrzeichen getauft, auch durch Gottes Wort absolvieret, und meiner Sünden frei, ledig und los gesprochen,
\quellseite{266}
und mit dem Sacrament des wahren Leibes und Blutes meines HErrn JEsu Christi gespeiset und getränket bin worden, als durch gewisse Gnadenzeichen, und daß ich Vergebung der Sünden empfangen habe, die mir mein lieber HErr JEsus Christus durch sein liebes Blut verdienet, erworben und erlanget hat; deß danke ich ihm in Ewigkeit. Amen.

\section[320. Eine Danksagung für die geistliche …]{320. Eine Danksagung für die geistliche Vermählung}

Lieber Gott, wie soll ich mich so hoch erheben, daß ich mich soll rühmen Gottes Braut, und Gottes Sohn meinen Bräutigam? Was kann ich armer stinkender Madensack zu der großen Ehre, welche auch denen Engeln im Himmel nicht widerfahren ist, daß sich die ewige Majestät sogar mit mir vereiniget, daß sie auch ein Leib mit mir sein will? Bin ich doch so gar, von dem Fuß an bis an die Scheitel voll Unflaths, Blattern, Grind, Ausfaß, Sünde und Stank vor Gott; wie soll ich denn der ewigen, hohen, herrlichen Majestät Braut, und mit ihr ein Leib heißen? Aber weil du es so haben willst, so sei dir Lob und Dank in Ewigkeit. Amen.

\section{321. Tröstliche Versicherung vom heil. Abendmahl}

Lieber Gott, du hast mir zugesagt, und ein gewiß Zeichen deiner Gnaden in den Sacramenten gegeben, daß Christi Leben meinen Tod in seinem Tode überwunden habe, sein Gehorsam meine Sünde in seinem Leiden vertilget, seine Liebe meine Hölle in seinem Verlassen zerstöret habe; dieses Zeichen, solches Zusagen meiner Seligkeit wird mir nicht lügen, noch trügen. Denn du hast gesagt, du kannst nicht lügen, weder mit Worten, noch Werken, da bleib ich auf, da steh ich auf. Amen.

\section{322. Eine Andere}

Gnade mir, du barmherziger Gott, ich bin ja ein armer sündiger Mensch und habe nichts denn Zorn verdienet; aber doch, ich habe gelebt, wie ich wollte, so halte ich mich hieher, daß ich weiß, und nicht zweifeln soll, daß ich getauft, und ein Christ genennet
\quellseite{267}
bin, zur Vergebung der Sünden, und daß mein HErr JEsus Christus für mich geboren, gelitten, gestorben und auferstanden ist, seinen Leib und Blut mir gegeben hat, zur Speise und Stärkung meines Glaubens. HErr JEsu Christe, ich bin in deinem Namen absolviret und entbunden von meinen Sünden, drum kann ich nicht übel fahren, noch verloren werden, so wenig, als Gottes Wort kann fehlen, oder falsch sein. Denn Gott selbst ist mir Bürge durch sein Wort. Amen.

\section{323. Klage über Unfruchtbarkeit in guten Werken}

HErr Christe, ich gehe zum Sacrament, und bleibe dennoch wie zuvor, ohne Frucht. Ich habe so großen Schatz empfangen, der bleibet da bei mir liegen und ruhen, das klage ich dir. Hast du mir den Schatz gegeben und geschenket, so gib auch, daß er Frucht und ein ander Wesen in mir schaffe, dich beweise und erzeige gegen meinen Nächsten, als ich schuldig bin. Amen.

\section[324. Trost wider Unwürdigkeit und des Teufels …]{324. Trost wider Unwürdigkeit und des Teufels Anfechtung}

HErr Gott, du hast mich durch einen Bruder absolviret, getauft und mit deinem Leib und Blut gespeiset, mache es mit deinem Knecht, wie dirs gefällt, ich will darüber nicht zürnen, noch dich lästern, sondern will alles mit Geduld ertragen, denn ich will nicht, daß dein Bund, so du mit mir in der heiligen Taufe und heiligem Abendmahl aufgerichtet, soll zu nichte werden. Ich weiß, daß ich getauft bin, daß ich den Leib deines Sohnes gegessen, und sein Blut getrunken habe, daß ich von Gottes wegen absolviret, und von meinen Sünden ledig gesprochen bin, ja, daß mir alle meine Sünden vergeben, und daß mir der Sieg verheißen wider den Teufel, Tod und Hölle. Was will ich denn nun mehr haben? Amen.

\section[325. Gebet um Gnade, Gottes Willen zu …]{325. Gebet um Gnade, Gottes Willen zu vollbringen}

Barmherziger Gott, der du uns wiederum deiner Gnaden Licht hast lassen aufgehen, durch JEsum Christum unsern HErrn, erleuchte, ermahne und stärke unsere Herzen, mit Kraft deines heiligen Geistes, im festen Glauben und feuriger Liebe, in allen Dingen
\quellseite{268}
zu thun, was dein väterlich, gnädiges Wohlgefallen ist, zu Ehren und Lob deines heiligen Evangelii, zu Trost und Nutz aller Gläubigen in Christo, dir sei Dank, Lob und Preis ewiglich. Amen.

\section[326. Gebet, daß uns Gott vor ferneren Sünden …]{326. Gebet, daß uns Gott vor ferneren Sünden behüte, und so wir sündigen, uns dieselben nicht zurechne}

Du gütiger barmherziger Gott, du lieber Vater im Himmel, du hast uns aus Gnaden und göttlicher Liebe deinen lieben Sohn geschenkt, und mit ihm alle Gnade, Heil und Seligkeit. Wir bitten dich, lieber Vater, erhalte uns solch selig Geschenk und himmlische Gabe, den freundlichen Anblick deines lieben Sohnes, Jesu Christi, daß wir ihn ja nicht durch Undankbarkeit verlieren oder sonst darum kommen. Wir sind fürwahr arme, elende und gebrechliche Menschen, fallen aus einer Sünde in die andere, sündigen jetzt mit Gedanken, jetzt mit Worten, mit Werken, und hat Mühe und Arbeit mit uns, daß wir bestehen. Da ist nimmer keine Ruhe, kein Friede; der Teufel lauert auf unsere Gedanken, schürt und bläst immerzu, die Welt lauert auf unsere Wort und Werk, Wesen und Leben, und gibt uns viel Aergerniß und Ursach zu Sünden, unser eigen Fleisch feiert auch nicht, ohne was zufällige Sünden, Laster und Untugenden sind, die uns täglich überfallen, die unser Gewissen greulich beschweren, und die Freude unseres Herzens gar zu nichte machen, und in eitel Trauer und Betrübniß verwandeln. Darum bitten wir dich, du gütiger und barmherziger Gott, ob wir vielleicht versäumlich und undankbar worden sind, und nicht so wandeln, wie wir billig wandeln sollen, so bleib du doch unser gnädiger Gott, sei uns freundlich, tröstlich, gütig und barmherzig, laß uns unsere mannigfaltigen Sünden nicht entgelten, sondern reinige durch dein Wort unser Herz und Gewissen, auf daß wir dir in Leibe und Freude dienen mögen, dich loben, ehren und preisen, zeitlich und ewiglich. Amen.

\section[327. Gebet nach dem Genuß des heil. …]{327. Gebet nach dem Genuß des heil. Abendmahls um neuen Gehorsam}

Gnädiger Gott und Vater, du hast mir aus Barmherzigkeit alle meine Sünden vergeben und mich zu Gnaden um Christi, deines
\quellseite{269}
lieben Sohnes willen angenommen. Ich habe dir auch von Herzen versprochen, mein Leben durch deine Gnade ernstlich zu bessern und frömmer zu werden. Ach Herr und Gott, weil des Menschen Thun nicht steht in seiner Gewalt, wie er wandele und seinen Gang richte, auch daß Dichten des menschlichen Herzens alle Tage böse ist; du aber dasselbe lenken kannst und in deiner Furcht leiten, wie dirs gefällig ist: so bitte ich dich, gib mir ein neues Herz, daß ich der Sünde von Herzen feind werde, laß deine Furcht mich heiligen, daß ich ein anderes und frömmeres Leben anfange und lebe in Heiligkeit und Gerechtigkeit, die dir gefällig ist. Laß mich deiner Gnade und Zusage nie vergessen, und hilf, daß ich wider mein Fleisch und Blut und wider alle bösen Lüste ernstlich streite und nicht wieder aufs Neue in Sünd und Laster falle. Ich habe vor mir Leben und Tod, laß mich ja das Leben ergreifen und nach dem Kleinod unverdrossen laufen, damit ichs ergreife und mich selbst nicht verwahre. Lehre mich thun nach deinem Wohlgefallen, auf daß ich gerecht und gottseilig lebe in dieser Welt und warte auf die selige Hoffnung und Erscheinung der Herrlichkeit des großen Gottes und unsers Heilandes Jesu Christi. Amen.

\section[328. Betrachtung der Kostbarkeit des Leidens …]{328. Betrachtung der Kostbarkeit des Leidens Christi}

Ach mein Gott und mein Herr! Laß mich erkennen und bedenken die Person, so um meinetwillen gelitten. Ist er nicht der eingeborne Sohn Gottes, das allerliebste Kind Gottes, wie ist er denn dahin gegeben in einen so schmählichen Tod? Ist er nicht der, von welchem der Vater vom Himmel gerufen: Dies ist mein lieber Sohn, an welchem ich ein Wohlgefallen habe? Wie muß er denn jetzt seinen Zorn tragen? Ist er nicht der Herr der Herrlichkeit? Wie trägt er denn jetzt solche Schmach? Ist er nicht der Allerstärkste? Wie ist er denn so schwach worden, daß er sich kreuzigen lassen? Ist er nicht ein ewiger, allmächtiger Gott? Wie leidet er denn solche unaussprechliche Noth? Ist er nicht das unbefleckte Lamm Gottes, und wird zu einer solchen blutigen Schlachtbank geführet? Ach der zarte Leib, der heilige und unbefleckte Leib, wie ist er verwundet! Ach die heiligen Hände! damit er uns segnet; ach die heiligen Füße! die auf dem Wege des Friedens gewandelt
\quellseite{270}
, wie sind sie durchgraben! Ach der holdselige Mund und rosin-farbene Lippen, wie sind sie erbleichet? Ach sein heiliges ehrwürdiges Haupt, wie ist es mit Dornen zerrissen? Ach sein liebreiches Herz, wie hat er uns dasselbe durch seine geöffnete Seiten gezeiget? Ach das heilige, zarte, reine Herzensblut, das aus seiner Seite geflossen, wasche mich von allen Sünden! Ach die freundlichen Augen, die klarer sind denn die Sonne, wie sind sie so verdunkelt? Wie sehen sie so kläglich? Wie rinnen sie mit Wasser als die Thränenquellen und die Liebebrünnlein? Ist doch alles kläglich und jämmerlich, was man an ihm siehet, ach! er ist so gar zerschlagen, daß nichts gesundes an seinem Leibe ist; das sind die Strafen und Plagen meiner Sünden, meine Missethat hat ihn so verwundet: O des großen Geheimnisses! O der großen Liebe! Daß meine Sünde an dem allerheiligsten Leibe Christi gestrafet worden. Ja durch ein solch heilig Opfer mußte ich versühnet werden, der Heilige träget meine Unheiligkeit, der Gerechte meine Ungerechtigkeit, der Fromme träget meine Bosheit, der Unschuldige trägt meine Schuld, der Herr gibt sich hin für seine Knechte, der Sohn Gottes für die Menschenkinder, der König der Ehren legt auf sich unsere Schmach, Gott gibt sich hin für sein Geschöpf, der Erlöser für die Gefangenen, der aller Welt Gott ist, für die Verlorenen, der Gesegnete, ja der Segen selbst, gibt sich hin für die Verfluchten, der Selige für die Verdammten, das Leben für die Todten. Ich habe gesündiget, was hat dies unschuldige Lämmlein gethan; ich bin gottlos gewesen, was hat dieser Gerechte mißhandelt? Ach HErr, du bist um meiner Hoffart willen gedemüthiget, um meines Ungehorsams willen gegeißelt, um meiner Wollust willen voll Schmerzen und Wunden geschlagen. Ich habe den Tod verdienet, und du mußt sterben: Ich war verkauft durch die Sünde, und du bist mein Lösegeld worden, und hast geleistet die allerhöchste Bezahlung für meine Sünde. Ach mein HErr JEsu, wie könnte eine Sünde so groß sein, die nicht mit diesem allerteuersten Schatz sollte bezahlet sein? Wie könnte eine Uebelthat so greulich sein, die nicht durch solche große Schmach, Schmerzen und Pein sollte gebüßet sein? Wie könnte der Sünden so viel sein, die nicht durch so einen großen Gehorsam, hohe Geduld, große Sanftmuth und
\quellseite{271}
tiefe Demuth dieser heiligen hohen Person sollte hingenommen sein? Ach mein Gott! wie kann dein Zorn so groß sein, daß er durch solchen Gehorsam, durch solch Trauren, Zittern und Zagen, blutigen Schweiß und bittern Tod nicht sollte versöhnet sein. Darum sei mir, o Gott! gnädig, und siehe nicht an meine Sünde, und nimm mich zu Gnaden an, und laß mich des theuren Blutvergießens deines lieben Sohnes Jesu Christi ewiglich genießen. Amen.

\section{329. Gebet über das Leiden Christi}

Ach mein Herr Jesu Christe, wenn ich dich ansehe, wie du im Garten so herzlich betrübt und vor großen Aengsten blutigen Schweiß geschwitzt, wie du nachmals gefangen, gebunden und die ganze Nacht und Stunde der Finsterniß verspottet und verhöhnet, des Morgens frühe als das unschuldige Lamm angeklagt, mit Ruthen gestrichen und endlich zum Tode verurtheilt, ans Kreuz gehängt, mit Händen und Füßen angenagelt worden bist und kein Glied an deinem heiligen Leibe unver wundet geblieben ist: — so sehe ich nichts als meine Sünde, Schuld und Missethat, welche dich dazu getrieben, und um welcher willen du dich in solche Angst, Noth, Betrübniß, Leiden und Schmerzen ergeben hast. Ach Herr, ich habe leider alles verschuldet, was du littest, ich bin des allen die Ursache. Wie wunderlich hat es doch Gott gemacht und verordnet, daß der Unschuldige bezahlen muß, was der Schuldige verschuldet, daß der Fromme entgelten muß, was der Böse verdienet. Die Schulden, welche der Knecht gemacht, muß der Herr bezahlen, was der sündige Mensch auf sich geladen, das muß Gott tragen. Ach, du Sohn des lebendigen Gottes, wie groß ist deine Liebe und Gunst gegen uns Menschen! Was hast du nicht alles gethan und ausgestanden, mich armen verlornen Sünder zu erlösen und vom Tod und der ewigen Verdammniß zu erretten. Ich habe übel gethan und du wirst gestraft; ich habe gesündigt und du mußt es erstatten. Ich bin ungehorsam gewesen und du mußt meines Ungehorsames wegen in solche Angst und Noth kommen. Ach du Herr der Ehren, was kann oder vermag ich dir zu geben für diese deine große Güte und Treue, die du an mir bewiesen! Womit kann ich armer dürftiger Mensch solches vergelten? Ach
\quellseite{272}
HErr, es bleibt unvergolten und ist mit nichts, weder im Himmel noch auf Erden, dir wieder abzuverdienen. Es ist deine Gnade und Gunst, deine Liebe und Treue, und die bleibt von mir armen Sünder ewig unvergolten. Und weil ich ja mit nichts meine Schuld abtragen kann, HErr, so will ich dir’s herzlich Dank wissen, was du an mir gethan, und will deinen Namen loben, rühmen und preisen von nun an bis in Ewigkeit. Amen.

\section{330. Danksagung für das Leiden Christi}

O du allerheiligster und holdseligster HErr JEsu Christe, wir danken dir für deine herzliche Traurigkeit, da deine Seele um unsertwillen betrübet war bis in den Tod, auf daß du die ewige Traurigkeit von uns abwendest, und uns die ewige Freude erwirbest. Für deinen demüthigen Fußfall, so du deinem himmlischen Vater thatest, da du auf die Erde fielest, auf dein heiliges Angesicht, auf daß du uns Gnade erlangtest, und unser Angesichter nicht beschämet würden. Für dein allerheiligstes Gebet, und heiligen Gehorsam, da du deinen Willen deinem himmlischen Vater ganz aufopfertest, und sprachest: Nicht was ich will, sondern was du willst, auf daß du unsern bösen Willen büßetest, heiletest, und uns in dem Willen Gottes heiligtest. Für deinen herben bittern Todeskampf, damit du dem Tode die Macht genommen, und ihn kraftlos gemacht und überwunden. Für deinen allerheiligsten blutigen Schweiß, welcher nildiglich aus deinem zarten Leichnam gedrungen, und auf die Erde gefallen, auf daß du unsern kalten Todesschweiß heiligtest, und die Angst des Todes in einen sanften Schlaf verwandeltest. O du unschuldiges, unbeflecktes Lämmlein Gottes, wir danken dir, daß du um unsertwillen bist gefangen, auf daß wir erlöset würden. Gebunden, auf daß wir von Sünden befreiet würden. Fälschlich verklaget, auf daß wir vor dem gestrengen Gerichte Gottes losgezählet würden. In dein heiliges Angesicht geschlagen, auf daß wir Friede hätten. O du allergeduldigstes und sanftmüthigstes Herz, wir danken dir, daß du um unsertwillen bist verspottet, auf daß du uns gemach würdest zur ewigen Weisheit. Verspeiet, auf daß du uns von unsern Schanden erlösetest. Gelästert, auf daß wir in dir zu Ehren gemacht
\quellseite{273}
würden. Gegeißelt, auf daß du unsern Ungehorsam büßetest. O du König der Ehren, und Herr der Herrlichkeit! wir danken dir, daß du um unsertwillen zu Hohn und Schmach mit Purpur bekleidet, auf daß du uns das hochzeitliche Ehrenkleid erwirbest. Mit Dornen gekrönet, auf daß du uns die Krone der Gerechtigkeit aufsetzest. Ein Rohr in deine rechte Hand genommen, auf daß du das schwache Rohr nicht vollends zerbrächest. Und damit lassen dein heiliges Haupt schlagen, auf daß wir unsere Häupter mit Freuden möchten aufrichten. O du allerliebster und freundlichster Herr! wir danken dir von Herzen: Für deine Vorstellung für das Volk, da Pilatus sprach: Sehet, welch ein Mensch! Auf daß dein himmlischer Vater unser Elend ansehe, und um deinetwillen sich unser erbarme: Ach Vater, siehe, welch ein Mensch! Du allerliebster, du bist von deinem Volk verworfen und verleugnet, auf daß du deiner gläubigen Kirchen zum Eckstein würdest. Du bist um unsertwillen zum Tode verurtheilet, auf daß du uns von dem Urtheil des ewigen Todes los machtest. O du allergerechtester und demüthigster Knecht Gottes, du allergehorsamster Sohn Gottes deines Vaters, wir danken dir, daß du dein Kreuz zu deinem heiligen Tode selbst hast getragen, auf daß du uns lehretest, unser Kreuz williglich auf uns zu nehmen: Bist daran mit Händen und Füßen angenagelt, auf daß du ein Opfer würdest für unsere Sünde: Bist zwischen zwei Mördern gekreuziget, und unter die Uebelthäter gerechnet: da du doch niemahl unrecht gethan, auch kein Betrug in deinem Munde erfunden worden, auf daß du uns durch deine Unschuld versöhnetest: Hast auch große grobe Lästerung und Schmach am Kreuz erlitten, auf daß du uns von der ewigen Schmach erlösetest. O du Gesegneter des Herrn, wir danken dir, daß du ein Fluch am Holz bist worden, auf daß in dir alle Völker auf Erden gesegnet würden: Du bist worden als ein Wurm, und bist doch der Schönste unter den Menschenkindern, auf daß du uns für Gott lieblich machtest: Du bist worden der allerverachtetste unter den Menschen, auf daß du uns herrlich machest: Du bist ganz trostlos am Kreuz gehangen, auf daß wir ewig getröstet würden! Du hast mit bloßem, blutigem Leibe sterben müssen, auf daß du uns mit dem Kleide des Heils, und mit dem Rock
\quellseite{274}
der Gerechtigkeit bekleidetest. O du ewiger Hoherpriester und unser einiger Mittler, wir danken dir, daß du am Kreuz für uns gebeten, auf daß du dich mit starkem Geschrei und mit Thränen deinem himmlischen Vater opfertest. Wir danken dir für das tröstliche Wort: Heute wirst du mit mir im Paradiese sein, damit hast du das Paradies den armen Sündern aufgeschlossen. Denn dies Wort ist der rechte Paradieseschlüssel. Wir danken dir für deine Angst und Noth, da du schrieest: Mein Gott! mein Gott! warum hast du mich verlassen? Auf daß wir von Gott nicht sollten ewig verlassen werden. Wir danken dir für deinen heiligen Durst am Kreuz, und den herben Essigtrunk, damit du uns vom ewigen Durst, und von der Höllen Bitterkeit erlöset hast. Wir danken dir für dein tröstliches Wort: Es ist vollbracht! das ist: Die Sünde ist nun getilget, Gott ist versöhnet, die Schrift ist erfüllet, und ist eine ewige Erlösung erfunden. Wir danken dir für deinen heiligen Tod, und für dein letztes Wort am Kreuz; denn damit ist alle unsere Sünde bezahlet, das Leben wiederbracht, und aller Gläubigen Seelen in Christo in die Hände des himmlischen Vaters überantwortet. Laß uns nun, o du allerliebreichster, gebenedeyter König, um deinen heiligen Leichnam mit Joseph von Arimathia bitten, denselben in ein rein Leinwand unsers Glaubens einwickeln, mit Myrrhen und Aloe salben, das ist, mit herzlicher Reue und Leid über unsere Sünde, in unsere Andacht aufnehmen, und in ein neues, gereinigtes Herz durch den Glauben, als in ein neues Grab legen, daß er allein, und sonst keiner mehr, darin ruhe. Und versiegle du dies Grab mit deinem Heiligen Geist, daß dich niemand, weder Welt noch Teufel, aus unserm Herzen stehle, daß wir dich nicht verlieren, sondern daß wir mit dir sterben, auferstehen, leben, gen Himmel zu dir fahren, und ewig bei dir sein und bleiben mögen. Amen.

\section[331. Gebet eines Betrübten, der sich …]{331. Gebet eines Betrübten, der sich bekümmert, als ob er das heil. Abendmahl unwürdig empfangen habe}

Himmlischer Vater, allmächtiger und barmherziger Gott, meine Zuversicht und meine Stärke, mein Heil und Hülfe in den großen Nöthen, die mich treffen! Ich rufe dich an von Herzensgrund
\quellseite{275}
im Namen Jesu Christi, deines allerliebsten Sohnes, und bitte dich, erhöre mein Gebet aus Gnaden, verbirg dich nicht vor meinem Flehen, sondern merke auf mich, wie ich so ängstlich klage und heule. Denn siehe, HErr, mein Herz ängstet sich in meinem Leibe, Furcht und Zittern ist mir angekommen über die große Traurigkeit meines Herzens, daß ich besorge, ich sei zu dem heiligen Tisch deines hochwürdigen Abendmahls nicht würdiglich gegangen und habe mich deswegen hart wider dich versündiget. Und wahr ist es, mein lieber Gott, wenn ich auf mich sehe, so bin ich vor dir in allem unwürdig und kann vor deiner heiligen Gerechtigkeit nicht bestehen. Aber es ist hingegen auch wahr, daß du deinen allerliebsten Sohn, JEsum Christum, uns zur Gerechtigkeit und Heiligung geschenkt hast, damit wir in ihm und seiner Versöhnung dir zu ewigen Gnaden gefallen und als gerecht aufgenommen werden können. Darum, mein HErr und Gott, treibe von mir durch deinen heiligen guten Geist solche und dergleichen trübselige falsche Gedanken, erleuchte mir mein Herz und meine Augen, und verleihe gnädiglich, daß ich mich von mir selbst abwende und allein nach Christo JEsu, deinem lieben Sohn, inbrünstiglich schauen möge, mich seiner hochgebendedeiten Heiligkeit wider meine unselige Unreinigkeit, seiner ewigen Gerechtigkeit wider meine verdammliche Ungerechtigkeit zu trösten und auf ihn festiglich zu verlassen. Wie ich denn weiß aus deinem heiligen Worte, daß desselben deines lieben Sohnes Gerechtigkeit in mir wohnet, sein Geist meinem Geiste Zeugniß gibt, daß ich dein Kind und du mein Herzensvater bist. Bestätige doch dasselbige in meinem Herzen und vermehre mein Vertrauen, erhalte mich in deiner Gnade und leite mich auf deinem Stege, damit ich als ein einverleibtes Glied deines lieben Sohnes ewiglich von dir ungeschieden sein und bleiben möge, um deines hochgelobten göttlichen Namens Ehre willen, hochgepreiset in alle Ewigkeit. Amen, Amen.

\quellseite{276}

\chapter{IX. Katechismusgebete}

\section{332. Gebet um wahre Erkenntniß Gottes}

Ewiger, wahrhaftiger, gerechter, lebendiger, heiliger Gott und Vater, ich klage und bekenne dir mit betrübtem und demüthigem Geist die große angeborne Blindheit und Finsterniß meines Herzens, daß ich dich von Natur nicht recht erkennen kann, auch zu deinem göttlichen Erkenntniß, in welchem doch das ewige Leben bestehet, aus natürlicher, fleischlicher Thorheit und Unwissenheit, weder Lust noch Liebe trage: Ach vergib mir solche schädliche Thorheit, gnädiger Vater, und rechne mir dieselbe nicht zu. Nimm die schwere Strafe von mir, da du durch den Propheten Esaïam das Wehe schreiest über alle, die dich nicht kennen wollen, und sprichst: Ein Ochse kennet seinen Herrn, und ein Esel die Krippe seines Herrn, aber Israel kennet es nicht, und mein Volk vernimmt es nicht. O wehe des sündigen Volks, des boshaftigen Samens. Erleuchte mich aber mit deinem Wort und Heiligen Geist, daß ich in wahrem Glauben erkenne, daß du nach deinem Wesen in drei Personen ein einiger, wahrer, wesentlicher Gott anzubeten bist, und du ewiger Vater seist die erste Person der heiligen Dreifaltigkeit, der du von Ewigkeit her deinen eingebornen Sohn aus deinem göttlichen Wesen gezeuget, und uns denselben zum Heilande geschenket hast. Ach gib mir, lieber Vater, daß ich deine Allmacht erkenne, deine Barmherzigkeit preise, deine Gerechtigkeit fürchte, deiner Wahrheit glaube, deine Weisheit rühme. Ach laß deine Allmacht sein meinen Schutz, deine Barmherzigkeit meinen Trost, deine Gerechtigkeit meine Stärke, deine Wahrheit meinen Schirm und Schild, deine Weisheit meine Regierung, deine Liebe meine Freude, deine Gnade mein Licht und Leben, dein Lob und Preis meine Herrlichkeit. Gib mir deinen lieben Sohn JEsum Christum zu erkennen, daß er sei ein ewiger, allmächtiger Gott, gleiches göttliches Wesens mit dir, und daß er in seiner angenommenen Menschheit sei die andere Person der heiligen Dreifaltigkeit, und
\quellseite{277}
sei mein HErr und mein Gott, der Glanz deiner Herrlichkeit und das Ebenbild deines Wesens, Licht von Licht, wahrhaftiger Gott vom wahrhaftigen Gott, durch welchen du alle Dinge gemacht hast, welcher ist das Leben und Licht der Menschen, welchen du uns ganz und gar geschenket hast mit allem, was er ist, Gott und Menschen, mit allen seinen ewigen Gütern und Wohlthaten, welcher ist durch seine Menschwerdung mein Bruder, durch sein Evangelium mein Prophet und Lehrer, durch seine Wunderwerke mein Arzt, durch sein Leiden und Sterben mein Erlöser, durch seine Auferstehung mein Sieg und Gerechtigkeit, durch seine Himmelfahrt meine Hoffnung, durch seine Herrlichkeit mein Preis und Ehre, durch seinen Heiligen Geist meine Salbung, durch seine Wiederkunft meine Seligkeit. Ach mein HErr JEsu, du bist der Weg, die Wahrheit und das Leben, sei mir der Weg, den ich wandeln soll, sei mir die Wahrheit, die ich glaube, sei mir das Leben, das mich erhalte und selig mache. Gib mir auch, o gütiger, gnädiger Vater, zu erkennen deinen Heiligen Geist, daß er sei die dritte Person der heiligen Dreifaltigkeit, mit dir und deinem lieben Sohn, ein ewiger, wahrer, wesentlicher Gott, von dir und deinem lieben Sohn ausgehend, und gesandt in die Herzen der Gläubigen, dieselben zu erleuchten, zu heiligen, zu trösten und in alle Wahrheit zu leiten. Ach laß denselben sein mein Licht wider alle Finsterniß und Irrthum, meinen Trost in aller Traurigkeit, meine Heiligung wider alle Unreinigkeit, meine ewige Wahrheit wider alle Lügen des Satans und aller Ketzer: Laß mich deinen Geist, als einen Geist der Wiedergeburt, erneuern, zum Tempel Gottes heiligen, daß die heilige Dreifaltigkeit zu mir komme, und Wohnung bei mir mache, daß ich durch deinen Geist eine neue Creatur werde, daß mein HErr Christus in mir lebe, und daß meine Glieder Christi Glieder sein mögen. Gib mir, o Gott, dich selbst zu erkennen, nach deinem gnädigen Willen, daß ich deine Liebe, o Gott Vater, in meinem Herzen schmecke, die Freundlichkeit und Leutseligkeit meines HErrn JEsu Christi in mir empfinde, daß ich die Frucht seiner Menschwerdung, seines Worts, seines Leidens, Auferstehung und Himmelfahrt allezeit in mir behalte, deines Heiligen Geistes Trost, Licht, Friede und Freude in meinem Gewissen empfinde, und
\quellseite{278}
also den Vorschmack und Schatz meiner Seligkeit, und mein höchstes Gut allezeit bei mir habe, und in mir trage. Dieses dein Erkenntniß, o Gott, ist das ewige Leben, dich erkennen, ist eine vollkommene Gerechtigkeit, und deine Macht wissen, ist eine Wurzel des ewigen Lebens: Durch dieses Erkenntniß werden wir gerecht und selig: Durch dieses Erkenntniß theilest du dich uns selber mit, mit allen deinen Gütern, dadurch wird unser Herz voll Gottes: Durch dieses Erkenntniß wird der Satan mit seiner Finsterniß, List und Lügen vertrieben, sein Reich zerstöret, seine Gewalt überwunden, ihm sein Palast und Harnisch genommen: Durch dieses Erkenntniß wird der Mensch in seinem ganzen Leben regieret zu Gottes Ehre und des Nächsten Nutz, und zu seinem eigenen Heil. Wo dieses Erkenntniß nicht ist, da erfüllet der Satan das Herz mit allerlei Finsterniß, Irrthum, Blindheit, Gottlosigkeit, und der Mensch wird Gottes Unehre und Schandfleck, ein unrein Gefäß des Zorns Gottes und ewiger Verdammniß. Dafür wollest du, o Gott Vater, durch dein heilig Erkenntniß mich gnädiglich behüten in Ewigkeit, durch JEsum Christum, deinen lieben Sohn, unsern HErrn. Amen.

\section{333. Gebet um wahre Gottesfurcht}

Heiliger, barmherziger, gnädiger und lieber Vater, ich klage und bekenne dir mit Reue und Seufzen meines Herzens, daß ich durch die giftige Seuche der angebornen Erbsünde also verderbet bin, daß von Natur keine Gottesfurcht vor meinen Augen ist, und daß ich wegen meiner tief verderbten Natur dich nicht kindlich fürchten, lieben noch ehren kann, auch mich leider nicht habe deine göttliche Furcht allezeit regieren lassen, sondern meine Gedanken sind mehrentheils ohne alle Gottesfurcht gewesen, und ganz eitel, wie auch meine Worte und Werke: Ach wie oft hab ich dein Wort und Gebot vergessen, und hintenan gesetzet, und mich nicht für demselben gefürchtet, wie oft hab ich den Lüsten der Welt und meines Fleisches mehr gefolget, denn deinem Wort und Befehl: Ach vergib mir, lieber Vater, solche Sicherheit, und Verachtung deiner Gebote und Drohungen, und wende die schwere Strafe von mir ab, die du drohest allen denen, die dich nicht fürchten, daß du dieselben
\quellseite{279}
verwerfen wollest, weil sie dein Wort verwerfen. Dieweil aber die Furcht Gottes die allerschönste Weisheit ist, dadurch ein Mensch Gott gefällt, so bitte ich dich von Herzen, du wollest-erstlich die angeborne Sicherheit und Vermessenheit, auch Verachtung deiner Gebote und göttlichen Drohungen, in meinem Herzen dämpfen und austilgen, und dagegen durch deinen Heiligen Geist, weil er ist ein Geist der Furcht des HErrn, in mir deine göttlich-kindliche Furcht erwecken, daß ich mich stets und an allen Orten, in allen meinen Gedanken, Worten und Werken, fürchten lerne für deinem Zorn und Ungnade, daß ich auch über alle und jede Sünden stetige wahre Reue und Leid in meinem Herzen habe, und daß mir viel mehr leid sei, und schmerzlich weh thue, daß ich dich beleidiget habe, denn daß ich gestrafet werde. O mein Gott und HErr, laß mich erkennen, daß ich dich unbillig beleidiget habe, ich aber billig gestrafet werde, darum verleihe mir deine Gnade, daß ich allen Sünden von Herzen feind werde, als den Werken des Teufels, daß ich mich auch in allen meinen Schwachheiten und Gebrechlichkeiten, im Kreuz und Anfechtung trösten möge deiner väterlichen Gnade und Barmherzigkeit, und dich nicht als meinen Feind, sondern als meinen Vater fürchte und liebe. Bewahre auch meine Seele und innerlichen Menschen durch deine göttliche Furcht, daß ich nichts wider deinen heiligen Willen denke, rede und thue, sondern daß ich alles gedenke, rede und thue, als vor deinen heiligen Augen und Angesicht, daß auch mein innerliches Auge allein auf dich gerichtet und gewandt sei, also, daß ich alle meine Worte und Werke in deiner Furcht zuvor wohl bedenke, und in allen Dingen deine göttliche Weisheit, Allmacht und Hülfe zuvor demüthig anrufe, daß ich mich auch kein zeitlich Ding, Ehre, Reichthum, weltliche Freude, Lust oder Menschenfurcht von deiner göttlichen Furcht lasse abwenden, sondern daß ich stets eingedenk sei der gnädigen Verheißung, da du den Gottesfürchtigen zugesagt hast Errettung, Erbarmung, Segen, Gnade, Hülfe, Weisheit, Erhaltung, Heil und Trost: Und daß mich deine göttliche Furcht, Gnade und Barmherzigkeit allezeit, in meinem ganzen Leben, regieren, leiten und führen möge, durch Jesum Christum, deinen lieben Sohn, unsern HErrn. Amen.

\quellseite{280}

\section{334. Gebet um herzliche brünstige Liebe}

O du liebreicher Gott, freundlicher und gnädiger Vater, du Ursprung und Brunnen aller Liebe, Freundlichkeit, Güte, Gnade und Barmherzigkeit! Ich klage und bekenne dir schmerzlich, daß in meinem Herzen alle deine göttliche Liebe, damit ich dich über alle Dinge lieben sollte, durch die Erbsünde so gar erloschen und erstorben ist, daß ich von Natur mehr geneigt bin zu mir selbst, zu meiner eigenen Liebe, und zu der Liebe der Creaturen, als zu dir, meinem lieben Gott und Vater; und habe demnach dich, meinen allerliebsten Vater, und meinen Erlöser, Jesum Christum, und den Heiligen Geist, meinen wahren Tröster, nie rechtschaffen und über alle Dinge geliebet, da du das höchste und ewige Gut selber bist, welches über alle Dinge sollte geliebet werden. Ach verzeih mir solche meine schwere Sünde und große Thorheit, und wende die schwere Strafe von mir ab, da geschrieben stehet: Wer Jesum Christum nicht lieb hat, der sei verflucht. Tilge aber in mir aus alle unordentliche Liebe, die Welt- und Creaturliebe, und eigene Ehre, die Augenlust, die Fleischeslust, das hoffärtige Leben, welches die Herzen der Menschen von dir abreißet. Zünde aber dagegen in mir an durch deinen Heiligen Geist die reine unbefleckte Flamme deiner Liebe, daß ich dich um dein selbst willen, als das höchste Gut, als die holdseligste Lieblichkeit, als die überschwänglichste Gütigkeit, als die wesentliche Heiligkeit, als die lauterste Wahrheit, Gerechtigkeit und Weisheit, und in Summa, als alles Gut, und einen ewigen Ursprung alles Guten, von Herzen lieb habe, ohne alle Betrachtung einigerlei Belohnung, eigenes Nutzens oder Vergeltung, sondern allein um dein selbst willen, daß ich dir auch mein Herz, Willen und Verstand unterwerfe, und mit Freuden deine Gebote halte, und deinen Willen gerne thue: denn darin stehet die wahre Liebe gegen dich. O Herr Jesu! die reine Liebe deines reinen Herzens entzünde mein kaltes Herz. Die zarte Liebe deiner unbefleckten Seele erleuchte meine Seele. Die Liebe deines edlen Gemüths erfülle meine Sinne und Gemüth. Die Liebe deiner göttlichen Kraft stärke unsere Leibes- und Seelenkräfte in deiner Liebe, daß ich auch um deiner Liebe willen Alles unterlasse, was dir mißfället; und Alles thue und leide um deiner Liebe willen, was dir
\quellseite{281}
gefället, und darin verharre bis ans Ende: denn dich also lieben ist die allerschönste Weisheit, und wer sie siehet, der liebet sie, denn er siehet, wie große Wunder sie thut; daß mich auch deine Liebe zu dir ziehe, mit dir vereinige, und mit dir, meinem HErrn, Einen Geist, Einen Leib und Seele mit dir mache; daß ich allezeit an dich gedenke, von dir rede, nach dir hungere und dürste, und in deiner Liebe gesättiget werde, in dir bleibe und du in mir; daß ich auch durch solche Liebe alle Menschen in dir, und um deinetwillen, als mich selbst liebe; daß ich auch aus erbarmender Liebe meinen Feinden vergebe, und dieselben liebe, und für sie bitte, und denselben wohlthue, und sie mit Güte überwinde. Und das Alles um der großen Liebe willen meines Vaters im Himmel, meines Erbarmers, um der vollkommenen Liebe willen Gottes des Sohnes, meines Erlösers, und um der inbrünstigen feurigen Liebe willen Gottes des Heiligen Geistes, meines einigen und wahren Trösters. Amen.

\section{335. Gebet um wahre beständige Hoffnung}

O Gott, du ewige Wahrheit, der du Glauben hältst ewiglich, an welchem nicht zu Schanden werden, die auf dich hoffen! Ich klage und bekenne dir, daß mein verderbtes Fleisch und Blut sehr am Zeitlichen hanget, immer siehet nach leiblichem und zeitlichem Trost, und vergisset oft deiner Verheißung, Allmacht und Barmherzigkeit. Ach vergieb mir diese schwere Sünde, und wende die harte Strafe von mir, da du dräuest, daß, die ihre Hoffnung aufs Zeitliche setzen, nicht sehen sollen den zukünftigen Trost. Tilge in mir aus alle falsche betrügliche Hoffnung, die wie ein Rauch und dünner Reif verschwindet. Benimm mich aller vergeblichen und unnützen Sorgen. Laß mich von Herzen glauben, daß du für mich sorgest und wachest. Laß meines Herzens Vertrauen und Hoffnung allezeit in dich gerichtet sein, daß dein liebreicher und süßer Einfluß deiner Gnade und Güte in mir nicht verhindert, und meines Herzens Einkehren und Ruhe in dir nicht zerstöret werde. Denn alle Hoffnung und Begierde, die in dir nicht ruhet, muß in ewiger Unruhe bleiben. Laß meine Hoffnung in dir allein, als einen Anker in einem festen Grunde, haften, und deiner Hülfe ungezweifelt erwarten. Laß mein Herz empfinden, daß deine Barmherzigkeit grundlos, deine Güte
\quellseite{282}
unendlich, deine Verheißung wahrhaftig sei: daß dadurch meine Hoffnung befestiget werde, daß sie nicht wanke; mein Gebet versichert, daß es gewiß erhöret sei; meine Zuversicht versiegelt, daß ich in deinem Schutz und Schirm behütet, und sicher sei, als in einer Festung. Der Grund meiner Hoffnung, o Vater, ist die holdselige Menschwerdung deines lieben Sohnes, und sein heiliges theures Verdienst, seine Auferstehung und Himmelfahrt, dadurch du uns zu einer lebendigen Hoffnung neu geboren, die mich auch gewißlich nicht wird lassen zu Schanden werden. Denn in Christo bin ich allbereit selig, mit ihm und in ihm bin ich allbereit auferstanden und gen Himmel gefahren, und in das himmlische Wesen gesetzet. Darum habe ich in Christo Jesu, meinem HErrn, allbereit das ewige Leben, und warte nur auf die Offenbarung der zukünftigen Herrlichkeit. Darum laß mich, o Gott, in allen Dingen, in allem Thun und Lassen auf dich allein hoffen, daß du mir Alles allein seiest, was mein Herz wünschet. Laß mich auch in Trübsal über den Verzug deiner Hülfe nicht allzusehr betrübet werden. Denn je länger du verziehest, je herrlicher du hernach hilfest. Lindere mir aber mein Kreuz, auf daß ich nicht zu müde werde, und stärke mich allezeit mit deinem Trost, daß ich mit denen, die auf den HErrn hoffen, neue Kraft bekomme, und nicht falle, sondern ewiglich bleibe, wie der Berg Zion, durch Jesum Christum, unsern HErrn. Amen.

\section{336. Gebet um wahre Demuth}

O HErr Jesu, du allerdemüthigstes Herz, ich bekenne und klage dir, daß ich von Natur zu eigener Ehre sehr geneigt bin, und der böse Geist mein Herz durch die Hoffart, welche ist ein Anfang aller Sünde, vergiftet hat: also, daß ich mich oft mit eigener angemaßter Ehre an dir versündiget, und nicht bedacht, daß alle Ehre dir allein gebühre, und keiner Creatur, meinen Nächsten oft verachtet, meinen Kräften zu viel vertraut, und auf mich selbst zum öftern mich verlassen habe. Ach HErr, der du den Demüthigen Gnade gibst, und widerstrebest den Hoffärtigen, rechne mir diese Sünde nicht zu, und wende die schwere Strafe von mir, so du den Hoffärtigen gedrohet hast, daß du sie stürzen und zerstreuen
\quellseite{283}
wollest, und laß meine Ehrfurcht, durch deine heilige Demuth, getilget und bezahlet sein. Lehre mich aber mein Elend erkennen, und die hohe Majestät Gottes fürchten, denn was bin ich, denn eine Hand voll Erde und Asche, ein faules Aas, ein Gefäß voller Unsauberkeit, ein elender Wurm, eine sündliche Geburt, ein Kind des Zorns von Natur, in Sünden bin ich empfangen und geboren, in Mühseligkeit und Eitelkeit lebe ich, mit Schmerzen und erbärmlicher Gestalt sterbe ich, wer weiß, wie und wo mein Ende sein wird? Ach ich bin ein Abgrund des Elendes, blind in deinem Erkenntniß, stumm zu deinem Lob, taub zu deinem Wort, lahm in deinen Wegen. Ich habe von mir nichts denn Sünde, Tod und Verdammniß, was ich aber gutes habe, das ist dein und nicht mein, denn was hat ein Mensch, das er nicht empfangen hat, darum ist dein allein die Ehre und nicht mein. Ich habe dir aber oft deine Ehre geraubet, und dieselbe mir durch Hoffart zugeeignet, mit fremdem Gut stolziret und gepranget, als ein ungetreuer Knecht, und ungerechter Haushalter. Ach fordere mich, lieber Herr, nicht zur Rechnung, ich kann nicht bestehen! Gib mir aber, daß ich demüthig sei im Herzen, einfältig in Worten, niedrig in Werken, daß ich mich nicht hoch achte, sondern gering sei in meinen Augen, pflanze in mein Herz wahre Demuth, daß ich dir in allen Dingen gehorsam sei, wozu du mich gebrauchen willst, gib, daß ich die Werke meines Berufes in Einfalt des Herzens verrichten möge, und nicht auf meine, sondern auf deine Kraft und Hülfe sehe, und derselben mit Geduld erwarte, hilf mir, daß ich das Meine, was du mir auferleget hast, getreulich ausrichte, und nicht trachte nach dem, dazu ich nicht berufen bin. Laß mich nicht in solche Blind- und Thorheit gerathen, daß ich große Dinge durch mein Vermögen getraue zu verrichten, laß mich bedenken, daß ich meinen Schatz in irdischen Gefäßen trage, laß nicht in meinem Herzen aufsteigen eine Verachtung meines Nächsten, auch nicht des geringsten Menschen, denn durch die Demüthigen verrichtest du dein Geschäft auf Erden, und große Dinge, gibst denselben deine Gnade, daß sie deine heilsamen Werkzeuge seien. Lehre mich durch Demuth Friede, und Einigkeit erhalten, o Gott des Friedens. Lehre mich, o demüthiger Jesu, deinem Exempel folgen, der du dich unter Gott, Engel
\quellseite{284}
und Menschen und alle Creaturen gedemüthiget und erniedriget hast, und dich einen Wurm genennet und gleich geachtet, weil ein Wurm das geringste ist unter den Creaturen. Ach laß dies Wort in meinem Herzen Frucht bringen, da du sagest: Lernet von mir, ich bin sanftmüthig, und von Herzen demüthig, nicht auswendig im Schein, sondern von Herzen. Ach, du HErr der Herrlichkeit, hast dich gedemüthiget, und der elende Wurm des menschlichen Herzens bläset sich auf. Ach du demüthiges Herz, du hast alle Ehre dieser Welt geflohen, und die deine Kinder sein wollen, können der weltlichen Ehre nicht satt werden, wann werde ich doch der Welt-Ehre absterben, daß ich mich keiner Ehre würdig achte, von niemand Ehre begehre, mich lerne für nichts achten, mir nicht selbst gefalle, sondern mir selbst mißfalle, mich selbst strafe, alles mein Thun für Koth achte, und für ein unreines Tuch vor deinem Augen. Gib mir auch, daß ich mich den lobenden Mund der Heuchler nicht lasse betrügen, sondern bedenke, daß mir tausend mal mehr mangelt. Gib mir die Niedrigkeit meines Herzens zu einem Schatz und zum Fundament aller Tugend, zur Nachfolge deiner heiligen Demuth. Laß mich bedenken das Wort, das du sagest: Wer sich erniedriget wie dies Kind, der ist der größte im Himmelreich. Item: Ich sehe an den, der zerbrochenen Geistes ist. Item: Daß der Höchste auf das Niedrige sehe im Himmel und auf Erden. Laß mich nicht ein Greuel werden vor deinen Augen durch Hoffart, einen Gottes-Dieb durch Ehrgeiz, einen Nachfolger und Gesellen des Lucifers durch Hochmuth, wende von mir alle hoffärtige Gedanken, stolze Gebeerden, prächtige Worte, und schließe mein Herz in dein demüthiges Herz, daß meine Seele in demselben ihre Ruhe und Wohnung habe ewiglich. Amen.

\section{337. Gebet um wahre Geduld}

Ach HErr Gott, lieber Vater, barmherzig und gnädig, geduldig und von großer Güte und Treue: Und o HErr Jesu Christe, du geduldiges Lämmlein Gottes: O Gott heiliger Geist, ein Geist des Trostes und Friedens: Ich klage und bekenne dir meines Herzens angeborne Ungeduld und Ungehorsam, daß ich oft wider deinen heiligen Willen gemurret und gezürnet habe: Ach
\quellseite{285}
Vergib mir solche schwere Sünde und Ungehorsam, und rechne mir dieselbe nicht zu, wende die harte Strafe von mir ab, die du dräuest denen, die dein Joch nicht tragen wollen; denn dieselben können nimmermehr Ruhe finden für ihre Seele. Gib mir aber deine Gnade, daß ich im Kreuz dir allezeit möge gehorsam sein, und ja nicht wider dich murre, noch mit dir zürne, sondern erkenne, daß solche meine Trübsal durch deinen göttlichen Rath mir auferleget worden, weil mir ohne deinen Willen kein Härlein kann vom Haupte fallen, und daß ich demnach mit gütigem, geduldigem, dankbarem Herzen alles Kreuz und Leiden von der Hand deiner väterlichen Vorsehung aufnehmen möge, ja mich noch größerer Strafe würdig achte, weil mein Kreuz und Leiden viel geringer ist, denn meine Sünde, ja daß ich des Kreuzes nicht allein nicht müde werde, sondern mehr begehre zu leiden, was dein Wille ist, und was zu deinen Ehren und mir zu Nutz gereichen soll. Ich laß mich erkennen, daß du aus großer Liebe mir mein Kreuz zugeschicket hast, auf daß du mich demüthigest, mein Fleisch kreuzigest, meinen Glauben, Liebe und Hoffnung prüfest, mein Gebet durch Geduld und Beständigkeit bewahrest, mich lehrest, tröstest, zum Reich Gottes bereitest, und im Himmel durchs Kreuz groß machest, und in deine Herrlichkeit einführest. Darum gib mir ein solch Herz, daß ich mein Kreuz lieb habe, dir dafür danke, mich dessen erfreue, daß ich dadurch dem Ebenbilde meines HErrn Jesu Christi gleich werde. O HErr Jesu, gib mir ein solch Herz und Sinn, daß ich in allen meinen Leiden allezeit dein Bild anschauen möge, wie du mit großer Geduld dein Kreuz getragen, wie dein heiliges Haupt geschlagen, mit Dornen gekrönet und zerrissen, dein Angesicht verspeiet und verspottet, dein heiliger Leib gegeißelt, verwundet und getödtet ist; ach wie hast du deinen Willen deinem himmlischen Vater in so großer Geduld aufgeopfert: Ich laß mich dir auch meinen Willen ganz und gar aufopfern, denselben verleugnen, mich dir ganz ergeben, und mich erfreuen, daß dein Wille, der allezeit gut ist, und alles gut machet, an mir möge erfüllet werden. Vertreib aus meinem Herzen alle Ungeduld, Traurigkeit, Angst, Furcht und Kleinmüthigkeit, daß ich durch Geduld in dir stark werde, alles zu überwinden, ja daß ich solche Schmerzen und Traurigkeit, nach deinem
\quellseite{286}
Willen, und durch deinen Willen in dir, mit dir, durch dich geduldig trage, und leide, so lange dirs gefällt, daß ich es halte für meines Herrn Jesu Christi Schmerzen, ihn dafür liebe und lobe jetzt und in Ewigkeit. Gib mir, Herr Jesu, Geduld, alles zu leiden, was du willst, gib mir deine Sanftmuth, daß ich nicht wider dich murre, deine Demuth, daß ich mich aller Strafe würdig achte, deine Gnade, alles zu erdulden, durch wahren Glauben alle meine Sorgen auf dich zu werfen, deinen wahrhaftigen Verheißungen zu trauen, die Zuversicht, daß du auch mitten im Kreuz mich lieb habest, die Hoffnung, daß du mir mein Kreuz werdest helfen tragen und lindern, tröste mich auch mit deinem Heiligen Geist, mit dem Vorgeschmack und Anblick des ewigen Lebens, daß ich mich mehr umsehen möge nach innerlichem himmlischen, denn nach äußerlichem Trost, gib, daß ich mich nicht fürchte vor denen, die den Leib tödten können, die Seele aber nicht; verleihe mir, daß ich mitten in Traurigkeit und Widerwärtigkeit, dennoch ein ruhiges und stilles Herz behalten möge, auch die Sanftmuth gegen meine Feinde, daß ich mich nicht räche weder mit Worten noch Werken, noch Begierden, auch nicht mit Gebehrden. Gib mir wahre Beständigkeit im Glauben, daß ich beharre bis ans Ende, und selig werde, und weil keines Christen Leben, Zeit und Beruf ohne Kreuz sein kann und muß: Du aber, o mein himmlischer Vater, deines armen Kindes Schwachheit weißt, so lege mir auf, was ich tragen kann, und mache mir es nicht zu schwer oder zu lange, laß mich deine Gnade mit meinem Kreuz oder mit meiner Last heben und tragen. Gib mir Gnade, die Trübsal so lange zu dulden, bis dir es gefällt, mich zu erlösen, auf daß ich wohl bewähret, geläutert und gereiniget werde. Laß meinen Herrn Christum mit seiner Sanftmuth, Demuth und Geduld in mir leben, daß nicht ich, sondern derselbe in mir lebe, der mein Leben ist. Gib mir auch gegen alle meine Beleidiger einen sanften Muth, eine erbarmende Liebe, daß ich mit dir sagen kann: Vater, vergib ihnen; einen gelinden, holdseligen Mund, eine milde und hülfreiche Hand, daß ich mit dir und deiner Liebe ewiglich verbunden und vereiniget bleiben möge. Amen.

\quellseite{287}

\section{338. Gebet um Heiligung des Namens Gottes}

Ach HErr und Gott, ich bekenne und klage dir von Grund meines Herzens, daß ich deinen heiligen, hohen und werthen Namen nicht allezeit geheiliget, und recht gebraucht, sondern manchmal denselben mit Gedanken, Worten und Werken innerlich und äußerlich entheiliget, unnützlich geführt und mißbrauchet. Ach lieber himmlischer Vater, verzeihe mir diese schwere Sünde, und wende von mir ab die Strafe, so du gedräuet hast denen, so deinen heiligen Namen unnützlich führen, gib auch Gnade, daß ich forthin mit höchstem Fleiß vor aller Entheiligung deines Namens mich hüten möge. O HErr Jesu Christe, der du mich gelehret hast, täglich zu beten: Vater unser, der du bist im Himmel, geheiliget werde dein Name, gib Gnade, daß solche Heiligung des Namens Gottes auch bei mir statt habe und von mir geschehen möge. Du allerheiligster Sohn des Allerhöchsten hast den Namen deines himmlischen Vaters vollkommen geheiliget, und uns Menschen geoffenbaret, um dieses deines vollkommenen Gehorsams willen wollest du mir vergeben, was jemals von mir wider die Heiligung des göttlichen Namens geschehen, und Gnade geben, daß ich forthin denselben mit Beten, Loben und Danken möge rechtschaffen heiligen. O Heiliger Geist, der du bist ein Geist, der da heiliget, du wollest auch mich durch und durch heiligen, daß ich in der Heiligung des Namens Gottes immerdar fortfahre und täglich völliger werde, und mein Geist sammt der Seele und Leib behalten werde auf die Zukunft meines HErrn Jesu Christi. Amen.

\section{339. Gebet um die Gnade, Gott zu loben}

O allmächtiger Gott, barmherziger Vater, ich bekenne und klage dir von Grund meines Herzens, daß ich deine vielfältigen Wohlthaten, so du mir und Andern erzeigest, nicht so inniglich betrachtet, und so herzlich dafür gedanket, als es billig sein sollte, daß ich dein Lob und Preis nicht mit solcher Andacht täglich verkündiget, als du es von mir erforderst. Ach, lieber Vater, verzeihe mir diese Sünde und wende von mir die Strafe, so du allen Undankbaren gedräuet hast; gib mir aber deine Gnade, daß ich hinfort dir mit Mund und Herzen täglich lobsinge, und für
\quellseite{288}
Alle deine Wohlthaten dir danke. Gib, daß ich dich allezeit lobe für deine große Gnade, und dir allezeit Preis und Lob opfere. Dich, HErr, loben Himmel, Erde und Meer, alle Wasserbrunnen und alle Creaturen, und was darinnen ist. Hilf, daß ich vielmehr, als deine vernünftige Creatur, dich täglich möge loben. Dich, HErr, loben die heiligen Engel, die starken Helden, deine Heerscharen und Diener, hilf, daß ich auch mit denselben meine Stimme erhebe und dir danke. O HErr JEsu Christe, der du in den Tagen deines Fleisches deinem himmlischen Vater so oft und so herzlich gedanket, durch dieses dein heiliges Exempel auch mich zum Lob und Preis Gottes und zur Danksagung zu bewegen, gib mir deine Gnade und Geist, daß ich auch möge jederzeit für alle Wohlthaten dankbar sein, und Gottes Lob verkündigen. Du hast vor deinem himmlischen Vater mich zum geistlichen Priester gemacht, hilf, daß ich durch dich das Lobopfer Gott allezeit opfern möge. O Heiliger Geist, der du bist ein Geist der Wiedergeburt und Erneuerung, du wollest auch mir eine neue Zunge geben, damit ich Gott, den HErrn, könne loben. Du himmlisches Freudenöl, erfreue und erquicke mein Herz mit geistlicher Freude, daß ich jederzeit mit fröhlichem Munde Gott, den HErrn, lobe, daß ich singe und spiele dem HErrn in meinem Herzen, und sage Dank allezeit für Alles, Gott und dem Vater, im Namen unsers HErrn JEsu Christi. Amen.

\section[340. Gebet um Liebe und Hochhaltung des …]{340. Gebet um Liebe und Hochhaltung des göttlichen Worts}

Allmächtiger, ewiger, gütiger Gott, der du uns hast deinen lieben Sohn, JEsum Christum, gegeben, welcher dein ewiges und wesentliches Wort und Ebenbild ist und uns dein Wesen und Willen in deinem Wort geoffenbaret hat: wir bitten dich, gib uns deinen Heiligen Geist, daß wir zu deinem Wort Lust und Liebe allezeit haben und alle rechte Lehre, Trost, Freude und Erquickung darin suchen und finden, und unser Leben, Amt und Beruf darnach anstellen und richten, auch andere damit unterweisen, und durch dasselbige Wort alle Widermärtigkeit und den Teufel selbst, die Sünde, den Tod und alles Unglück überwinden. Ach, HErr, gib, daß wir ja dein Wort für den größten, herrlichsten Schatz halten,
\quellseite{289}
höher und größer, denn die ganze Welt, mit allem, was sie ist, hat und vermag. Laß dein Wort eine brennende Leuchte sein in unsern Herzen; eine Richtschnur, darnach wir glauben; einen geraden Weg, auf dem wir in unserm Leben und Beruf wandeln; einen Trost im Kreuz, eine Labung in aller Noth, eine feste Burg wider alle Feinde, einen Trotz wider den Teufel, einen Schutz und Schirm, und Freude unseres Herzens und Gewissens. Erhalte uns, HErr, bei deinem Wort, und wehre allen Feinden und Verfolgern deines Worts; daß wir dich recht erkennen und loben und preisen von nun an bis in Ewigkeit; der du bist ein einiger, ewiger, allmächtiger, wahrer Gott, Vater unseres HErrn JEsu Christi, mit sammt deinem Sohne, unserm HErrn JEsu Christo, und dem Heiligen Geiste. Amen.

\section[341. Gebet um Lust und Liebe zum göttlichen …]{341. Gebet um Lust und Liebe zum göttlichen Worte}

Ach barmherziger Gott und Vater, du siehest, wie ich wegen meiner verderbten Natur keine Lust zu deinem Worte habe, und wie leicht ich mich vom Teufel, meinem Fleisch und Blut, von den Weltkindern, auch falschen Lehrern und Predigern, ja andern gar lüderlichen Ursachen davon so gar leicht lasse abhalten, unfleißig, faul und träge bin, dein Wort zu hören und zu bewahren. Darum bitte ich nun, o ewiger Gott, verleihe mir diese meine angeborne Faulheit und Trägheit und gib mir doch ein williges und munteres Herz, dein Wort zu hören und zu betrachten; erwecke in mir eine solche herzliche Begierde, daß ich begierig sei nach der vernünftigen lautern Milch göttlichen Worts, als die jetzt gebornen Kindlein. Hilf, daß ich meine beste Lust und Ergötzlichkeit an deinem Worte habe, daß in dieser Welt mir nichts lieber als dein Wort sei, daß ich dein Wort liebe über Gold und viel feines Gold und dasselbe allzeit für meinen besten Schatz achte. Dieweil ich auch leider die Zeit erlebt habe, von welcher dein lieber Sohn selbst geweissaget hat, daß falsche Christi aufstehen und Wunder thun werden, daß verführet würden in den Irrthum, wo es möglich wäre, auch die Auserwählten: so bitte ich dich, du wollest mich ja vor Irrthum und falscher Lehre gnädiglich behüten und bewahren. Erhalte mich bei deiner Wahrheit, denn dein Wort ist die
\quellseite{290}
Wahrheit, daß ich demselben als der himmlischen Wahrheit anhangen, auch darin bis an mein letztes Ende standhaftig verharren möge, und dieses wollest du thun um deines allerheiligsten und hochgelobten Namens Ehre willen. Amen.

\section{342. Gebet, wenn man in der Bibel lesen will}

Ewiger Gott und Vater unsers Herrn Jesu Christi, verleihe mir deine Gnade, daß wir die heilige Schrift wohl und fleißig studiren und Christum darin suchen und finden und durch ihn das ewige Leben haben. Das hilf uns, lieber Gott, mit Gnaden. Amen.

\section{343. Gebet vor dem Bibellesen}

Allerliebester Jesu! der du mit freundlichen Lippen gesprochen: Suchet in der Schrift, denn ihr meinet, ihr habt das ewige Leben darinnen, und sie ist's, die von mir zeuget; laß dir gefallen die Begierde meines Herzens zu deinem Wort und segne selbst mein jetziges Beginnen. Ohne dich kann ich freilich nichts thun, so verrimmt auch der natürliche Mensch nichts vom Geiste Gottes, es ist ihm eine Thorheit, und kann es nicht erkennen: darum, so öffne du mir selbst das Verständniß, daß ich die Schrift verstehe, und deine Wunder sehen möge in deinem Gesetz. Gib mir den Geist der Weisheit und der Offenbarung zu deiner Selbsterkenntniß, und erleuchtete Augen meines Verständnisses, damit ich möge erkennen die Hoffnung unsers Berufs, und welches sei der Reichthum deines herrlichen Erbes an deinen Heiligen. Laß mich gefangen nehmen alle Vernunft unter den Gehorsam Christi, und hingegen glauben alle dem, was hierinnen die heiligen Menschen Gottes haben geredet, getrieben von dem Heiligen Geist; damit ich wachsen möge in der Erkenntniß Gottes und gestärkt werde mit aller Kraft nach deiner herrlichen Macht: bis daß wir alle hinankommen zu einerlei Glauben und deiner Erkenntniß, und ein jeder ein vollkommener Mann werde, der da sei in der Maße deines vollkommenen Alters, um dein selbst willen. Amen.

\quellseite{291}

\section{344. Gebet nach dem Bibellesen}

Ich danke dir, herzliebster Jesu, du Anfänger und Vollender meines Glaubens, daß du in mir beständige Lust zu deinem Wort und also beide, das Wollen und das Vollbringen nach deinem Wohlgefallen gewirket hast. Nun erfahre ich, daß es wahr sei, was der Geist durch David gesprochen: Wenn dein Wort offenbar wird, so erfreuet es, und macht klug die Einfältigen. Versiegle du nun in mir alle diese Worte der Weissagung in diesem Buch, und nimm nicht von meinem Munde das Wort der Wahrheit, denn ich hoffe auf deine Rechte. Es ist ewiglich mein Schatz und meinem Munde süßer denn Honig, es ist meines Fußes Leuchte und ein Licht auf meinen Wegen. Ja, erhalte uns, Herr, dies dein Wort; denn es ist unsers Herzens Freude und Trost. Amen.

\section{345. Gebet um die Ehre des heil. Predigtamtes}

O allmächtiger, ewiger, lieber himmlischer Vater, ich danke dir von Grund meines Herzens, daß du uns treue Lehrer und Prediger gegeben, welche uns dein heiliges, göttliches Wort rein und unverfälscht vortragen. Ich bekenne dir aber auch herzlich, daß ich dieses theuren werthen Schatzes nicht allezeit so hoch geachtet, als ich thun sollen, daß ich meine vorgesetzten Lehrer nicht für deine Diener und Botschafter jederzeit geehret. Ach, lieber himmlischer Vater, verzeihe mir diese Sünde, und gib Gnade, daß ich forthin alle schuldige Ehre, Gehorsam und Gutthätigkeit ihnen erweisen möge. O Herr Jesu Christe, der du als der einige Hohepriester des Neuen Testaments und Haupt deiner Kirche in die Höhe aufgefahren, dich zur Rechten deines himmlischen Vaters gesetzt, und gibst noch heutiges Tages deiner Kirche Hirten und Lehrer, daß durch dieselben dein geistlicher Leib erbauet werde, ich bitte dich herzlich, du wollest auch diejenigen, so du aus Gnaden gesendet und gegeben, mit Kraft aus der Höhe anthun, und mit deinem Heiligen Geist begaben, daß sie dein Wort mit aller Freudigkeit predigen mögen. Lege dein Wort in ihren Mund, und bedecke sie unter dem Schatten deiner Hände, daß sie den Himmel pflanzen und die Erde zurichten. Schmücke sie mit reichem Segen, daß sie einen Sieg nach dem andern wider den Teufel und sein Reich erhalten, daß
\quellseite{292}
man sehen möge, der rechte Gott sei zu Zion. O Heiliger Geist, der du in Gestalt feuriger Zungen über die Apostel am Pfingsttage dich ausgegossen, erfülle auch unsere Lehrer mit deinem heiligen Licht, daß sie auch mit feurigen Zungen reden, damit es den Zuhörern durchs Herz gehe. Du himmlischer Lehrer der Wahrheit, gib zu ihrem Pflanzen und Begießen dein Gedeihen. Gib uns aufmerksame Ohren und willige Herzen, daß wir demjenigen, was uns aus Gottes Wort wird vorgetragen, im Gehorsam williglich folgen mögen, und dadurch zur ewigen Seligkeit geheiliget, geleitet und bewahret werden. Amen.

\section{346. Gebet um die Ehre der Eltern}

O allmächtiger, ewiger Gott, himmlischer Vater, der du mit großem Ernst geboten, die Eltern zu ehren, ich danke dir von Grund meines Herzens, daß du mich hast lassen von christlichen Eltern geboren werden, und dieselben bisher mir erhalten; daneben aber bekenne und klage ich dir, daß ich meine Eltern nicht allezeit in solchen Ehren gehalten, als ich zu thun schuldig gewesen, sie nicht also geliebet, gefürchtet und alles Gutes ihnen erzeiget, als ich billig thun sollen, sondern oftmals mit Worten und Werken mich ihnen widersetzt. Ach, lieber Gott und Vater, verzeihe mir diese Sünde meiner Jugend und meine Uebertretung, gedenke aber meiner nach deiner großen Barmherzigkeit um deiner Güte willen. Gib mir auch forthin ein gehorsames, liebreiches Herz, daß ich meine lieben Eltern von Herzen mit Worten und Werken ehre. O Herr Jesu Christe, der du in den Tagen deines Fleisches deinen Eltern unterthänig gewesen und ein Exempel der Nachfolge auch hierinnen mir gelassen, daß ich deinen Fußstapfen folgen soll, gib mir auch ein solches dankbares, unterthäniges, williges und gehorsames Herz gegen meine Eltern, daß ich mit Gehorsam, Liebe und Furcht, Geduld und Demuth sie ehre. Siehe auch meine lieben Eltern an mit den Augen deiner Barmherzigkeit, laß sie zu einem geruhigen Alter kommen, segne ihre Nahrung, lindere ihr Kreuz, verleihe ihnen zeitlichen und ewigen Segen. O Heiliger Geist, der du bist ein Geist der Liebe und der Zucht, regiere und führe mich also, daß ich meine Eltern herzlich liebe, und ihre Züchtigung
\quellseite{293}
mit Geduld annehme. Erleuchte mein Herz durch deine himmlische Gnade und Licht, daß ich mit dankbarem Herzen erkenne, was für große Wohlthaten durch meine lieben Eltern mir von Gott erzeiget, damit ich ja nicht undankbar und ungehorsam erfunden werde, sondern gegen Gott und meine Eltern alle schuldige Unterthänigkeit und Gehorsam erweisen möge. Amen.

\section{347. Gebet um christliche Freundlichkeit}

Ach allerliebster Jesu, auserwählter Seelenfreund! wie freundlich hast du dich in den Tagen deines Fleisches gegen jedermann erzeiget? Was von dir längst geweissaget worden, daß du nicht würdest mürrisch und greulich sein, das hast du in der That erfüllet. Wie liebreich war dein Herz, wie schön dein Angesicht? Wie holdselig sind deine Lippen, wie gütig deine Augen, wie anmuthig deine Geberden, wie hülfreich deine Hände gewesen? Damit hast du mir ein Vorbild gelassen, daß ich mich gleichermaßen gegen meinen Nächsten bezeigen soll. So gib mir nun die Gnade, daß ich von dir solche Tugend lerne, und deinen Fußstapfen folgen möge. Laß aus meinen Augen leuchten deiner Augen Lieblichkeit, aus meinem Munde deiner Lippen Holdseligkeit, aus meinem Angesicht deines Angesichts Freundlichkeit. Mache mich deinem Bilde, allerschönster Heiland, ähnlich, daß ich mit Liebe, Gütigkeit und Gelindigkeit jedermann begegne, mit Rath erscheine denen Elenden, mit Trost den Betrübten, mit Hülffe den Dürftigen; daß ich niemand betrübe mit meinen Lippen, niemand schade mit meinen Worten, niemand erschrecklich sei mit meinen Geberden. Laß mich nicht durch Zorn und Grimm dem grausamen Teufel und scheußlichen Drachen aus der Hölle gleich werden, sondern mein Gesicht allezeit holdselig und freundlich wie eines Engels Angesicht erscheinen, und mich an Gnade und Gunst bei Gott und Menschen von Tag zu Tag zunehmen. Amen, o freundlichster und holdseligster Jesu, Amen.

\section{348. Gebet um Keuschheit des Herzens}

Ach du edles, heiliges, reines, keusches, unbeflecktes, züchtiges Herz, Herr Jesu Christe, du Liebhaber der Reinigkeit, du Krone aller Ehre und Tugend! Ich klage und bekenne dir meines Herzens
\quellseite{294}
angeborne Unreinigkeit, dadurch ich meinen Leib und Seele oft beflecket habe durch unreine Gedanken, Worte und Werke. Ach vergib mir, du reines, mildes, gütiges Herz, diese meine große Sünde und wende die schwere Strafe von mir ab, so du den Unreinen bräufst: denn so die selig sind, die reines Herzens sind, und Gott schauen werden, so werden die ohne Zweifel unselig sein, die unreines Herzens sind, und werden Gott nicht schauen. Darum schaffe in mir, Gott, ein reines Herz, und verwirf mich nicht von deinem Angesicht um meiner Unreinigkeit willen. Und weil ich erkenne, daß ich nicht anders kann züchtig leben, du gebest mir's denn, und solches Erkenntniß ist auch eine große Gnade, so bitte ich dich demüthiglich, heilige und reinige mein Herz durch den Glauben, durch den Heiligen Geist, durch die Buße und durch die neue Geburt, und stärke mich, daß ich den unreinen Geist in mir nicht lasse herrschen, oder mich einnehmen und besitzen lasse, wie ein unreines Haus, auf daß er nicht meine Seele beflecke, meine Gedanken vergifte, meinen Leib verunreinige. Lösche in mir aus die Flammen der Unzucht, umgürte meine Lenden und Nieren mit dem Gurt der Keuschheit, du keuscher, züchtiger und edler Bräutigam meiner Seelen! Umfasse mein Herz mit deiner reinen Liebe, reinige und vermähle meine Seele mit deinem keuschen Herzen, erfülle mein Herz mit reinen und heiligen Gedanken, daß ich von deinem Leibe ein reines und unbeflecktes Glied sein und allezeit bleiben möge, damit ich nicht ein unreines Gefäß und Glied des Satans werde, ein Gefäß der Unkehren und Schanden, sondern ein Gefäß der Gnaden und Ehren, und meine Gaben, die du in mich, als in ein Gefäß der Barmherzigkeit, geleget hast, nicht verschütte und gar verliere; daß ich mich durch Zucht und Reinigkeit absondere von den unsaubern Geistern, den Teufeln, und von allen unreinen Heiden, daß ich nicht ausgestoßen werde aus dem neuen himmlischen Jerusalem, sondern mit dir vereinigt bleibe und ein Geist, ein Herz, ein Leib mit dir sei, gleichwie du mich dazu in der heiligen Taufe gewaschen, geweihet, und mit dem Heiligen Geist gesalbet und geheiliget hast zu deinem heiligen Tempel und Wohnung. Ach, lehre mich bedenken, daß mein Leib ein Tempel und Wohnung sei des Heiligen Geistes, daß ich Gottes Tempel nicht zerbreche,
\quellseite{295}
auf daß du mich nicht wieder zerbrechest, und daß ich ein Glied Christi sei, und aus mir selbst kein unehrliches Glied mache, und an meinem eigenen Leibe sündige. O du edler, keuscher, züchtiger Himmelsbräutigam, der du unter den Rosen der Reinigkeit weidest, weide meine Seele mit deinem Erkenntniß und reiner Liebe und treibe von mir alle unreinen Gedanken, daß du mit deinem heiligen Geist in mir mögest wohnen, und die heiligen Engel bei mir bleiben mögen allezeit. Amen.

\section{349. Gebet um Gerechtigkeit im Handel und Wandel}

Allmächtiger, ewiger, barmherziger und gerechter Gott, ich klage dir von Grund meines Herzens, daß ich manchmal mich habe des Nächsten gelüsten lassen und habe ihn im Handel und Wandel übervortheilt. Dadurch habe ich die schwere Strafe auf mich geladen, welche du allen Ungerechten gedrohet hast. Ich bitte aber herzlich, du wollest um Christi, deines lieben Sohnes, des Gerechten willen, meine Sünde und Ungerechtigkeit mir verzeihen, und Gnade geben, daß ich forthin in allem meinem Thun und Wandel der Gerechtigkeit mich befleißige, dem Nächsten ja nicht das Seine durch List oder Gewalt mit oder ohne den Schein des Rechten entwende, sondern vielmehr dasselbe zu behalten ihm förderlich und dienstlich sei. Behüte mich gnädiglich, lieber himmlischer Vater, daß ich nicht fehle oder an deinem heiligen Namen mich vergreife. O Herr Jesu Christe, der du in den Tagen deines Fleisches in allem Handel und Wandel auf das Gesetz der Natur uns gewiesen: Alles, das ihr wollt, daß euch die Leute thun sollen, das thut ihr ihnen; verleihe mir, daß ich demselben treulich nachlebe. Du hast das ewige, unwandelbare Gebot Gottes, daß man nicht stehlen soll, zum öfteren wiederholet; hilf, daß ich dasselbe nimmermehr aus meinem Herzen lasse kommen, sondern stets daran gedenke und mich für aller Ungerechtigkeit hüte. O Heiliger Geist, dämpfe und tödte in meinem Herzen die unordentliche Lust und Begierde, reich zu werden, daraus Diebstahl und Ungerechtigkeit erfolget, bekräftige und vermehre in mir die Liebe des Nächsten, daß dieselbe eine Meisterin und Regiererin sei in allem meinem
\quellseite{296}
Handel und Wandel, damit ich meine Sache also ausrichte, daß ich niemand Unrecht thue, sondern in der Gnade Gottes und Liebe des Nächsten zum ewigen Leben erhalten werde. Amen.

\section{350. Gebet wider den Geiz}

O allmächtiger, ewiger Gott, himmlischer Vater! ich bekenne und klage dir von Grund meines Herzens, daß ich von Natur zum Geiz geneigt bin und manchmal demselben in meinem Herzen Raum und Statt gegeben habe. Ach lieber himmlischer Vater, verzeihe mir diese Sünde und gib Gnade, daß ich forthin dem Geiz von Herzen feind werde, weil ja derselbe eine Wurzel ist alles Uebels, daraus Ungerechtigkeit, Unbarmherzigkeit gegen den Nächsten, Vertrauen aufs Zeitliche und viel böser schädlicher Früchte erwachsen. Reinige mein Herz, o himmlischer Vater, von diesem schädlichen Gift und pflanze in mich die edle Tugend der Genügsamkeit, daß ich mir begnügen lasse an dem, das da ist, in gewisser Hoffnung, du werdest mich nicht verlassen noch versäumen. O Herr Jesu Christe, der du so treulich und brüderlich uns gewarnet hast: Sehet zu, hütet euch vor dem Geiz, denn niemand lebt davon, daß er viel Güter hat, ach, gib Gnade und Segen, daß ich dieser Warnung möge folgen. Du hast gelehret: Ihr sollt euch nicht Schätze sammeln auf Erden, da sie die Motten und der Rost fressen und da die Diebe nicht nachgraben und stehlen, ach hilf, daß ich diese Lehre zu Herzen nehme und mein Leben darnach richte. Es ist ja der Geiz eine schändliche Abgötterei; dieselbe wollest du, o mein getreuer Heiland, aus meinem Herzen tilgen und austreiben. Laß mich, o lieber himmlischer Bräutigam, mit reiner Liebe dir anhangen, daß ich nicht Fleischeslust, Augenlust und hoffärtiges Leben in mir lasse herrschen. O Heiliger Geist, dämpfe und tödte in mir den Geiz, welcher ist ein schädliches, verdammliches Werk des Fleisches. Erneuere mein Herz, daß ich den alten Menschen mit seinen Werken ausziehe und meine Glieder, die auf Erden sind, böse Lust und den Geiz tödte. Behüte mich gnädiglich, daß der betrügliche Reichthum und Sorge der Nahrung als stachlichte Dornen mein Herz nicht mögen ritzen und den Samen des göttlichen Worts in mir ersticken. Hilf, daß ich nicht meine Seele mit Gift des unrechten
\quellseite{297}
Guts, sondern mit dem himmlischen Manna des göttlichen Worts speise, denn weil sie ewig unsterblich ist, muß sie eine ewige und unsterbliche Speise haben. Gib mir die rechte Ruhe der Seele, daß ich durch Genügsamkeit in Gott ruhen möge und endlich in die ewige Ruhe eingehe. Amen.

\section{351. Gebet um Gutthätigkeit}

Allmächtiger, ewiger Gott, himmlischer Vater, ich bekenne und klage dir von Grund meines Herzens, daß ich manchmal verdrossen und träge gewesen bin, Almosen zu geben. Ich habe leider oft mein Herz vor dem dürftigen Bruder zugeschlossen und die hilfreiche Hand ihm nicht geboten. Ach, lieber himmlischer Vater, verzeihe mir diese schwere Sünde und gib Gnade, daß ich forthin mitleidig sei gegen den armen Lazarus, ihm gerne gebe und mittheile. Was bin ich, daß ich sollte vermögen Kraft, freiwillig zu geben? Du wollest selber durch deinen Heiligen Geist mir ein mildreiches, freiwiliges Herz geben und ewiglich solchen Sinn in meinem Herzen bewahren und mein Herz zu dir schicken. O HErr JEsu Christe, der du uns so treulich ermahnt hast: „Verkaufet, was ihr habt und gebt Almosen!“ „Gib dem, der dich bittet, und wende dich nicht von dem, der dir abborgen will.“ Gib mir deine Gnade und Geist, daß ich diese treuen Vermahnungen zu Herzen nehme und derselben nimmermehr vergessen möge. Ach, mein treuer Heiland, regiere und leite mich durch deinen guten Geist, daß ich in diesem Leben also möge wandeln, daß ich dermaleinst am jüngsten Tage zu deiner Rechten gestellt werde und die fröhliche Stimme höre: „Kommt her, ihr Gesegnete meines Vaters, ererbet das Reich, das euch bereitet ist von Anbeginn der Welt. Denn ich bin hungrig gewesen, und ihr habt mich gespeiset. Ich bin durstig gewesen, und ihr habt mich getränket.“ O Heiliger Geist, dessen Frucht ist die Gütigkeit und Gutthätigkeit, ändere und bessere mein hartes, unbarmherziges Herz, nimm von mir das steinerne Herz und gib mir ein fleischernes Herz, daß ich des Nächsten Noth und Dürftigkeit mir lasse zu Herzen gehen, hilf, daß wir als Glieder eines geistlichen Leibes untereinander mitleidig und behülflich sein, damit wir uns hierdurch als Kinder Gottes, Christi Glieder und Erben des ewigen Himmelreichs erweisen mögen. Amen.

\quellseite{298}

\section{352. Gebet um Wahrhaftigkeit}

Allmächtiger, ewiger Gott, himmlischer Vater, ich bekenne und klage dir von Grund meines Herzens, daß ich nicht allezeit die Wahrheit geliebt und von Herzen mit meinem Nächsten die Wahrheit geredet habe, sondern habe manchmal ihn belogen und betrogen, ihn verleumdet und verlästert, seinen guten Namen geschändet und bösen Leumund gemacht. Ach, lieber himmlischer Vater, vergib mir diese Sünde und wende von mir ab die schwere Strafe, die du allen Lügnern, Lästerern und Verleumdern gedroht hast. Gib mir aber deine Gnade, daß ich forthin in der Wahrheit wandle, vor allen Lügen und Falschheit mich fleißig hüte. Abgötterei und Lügen laß, o lieber Gott und Vater, ferne von mir sein. O HErr JEsu Christe, der du dich selber die Wahrheit nennest und in allen deinen Thun und Worten dich wahrhaftig erwiesen hast, regiere mein Herz also, daß ich der Lügen Feind sei und der Wahrheit beständig anhange. Behüte mich auch gnädiglich vor den Lügnern, Lästerern und Verleumdern, daß sie mich nicht berücken noch zu Fall bringen mögen. Laß sie zu Schanden werden, die mich mit Lügen wollen niederdrücken. Hilf mir wider die, so mich mit Lügen verfolgen. Laß mich nicht unter die Lästerer gerathen und laß mich nicht unter ihnen verderben. O Heiliger Geist, der du bist ein Geist der Wahrheit, gib mir Stärke und Kraft, daß ich die Wahrheit frei heraus jederzeit bekenne und vertheidige. Dämpfe und tödte in mir mein sündliches Fleisch, welches mich immerdar zu Lügen und Falschheit, zur Heuchelei und Betrug anreizet, gib Gnade, daß ich seiner Anreizung nicht folge, sondern mich jederzeit bestleißige, daß Mund und Herz zusammenstimmen, und ich auf dem Wege der Wahrheit zum Reich der Wahrheit allezeit wandeln möge. Amen.

\section{353. Gebet um Regierung der Zunge}

Allmächtiger, ewiger Gott, himmlischer Vater, ich bekenne und klage dir, daß ich oft und viel mit meinem Munde wider dich, meinen lieben Gott, und wider meinen Nächsten gesündiget, daß ich meine Zunge nicht so fleißig im Zaum gehalten, als ich billig hätte
\quellseite{299}
thun sollen, dieselbe nicht allezeit zu deines Namens Ehre und des Nächsten Besserung gebraucht, manchmal unnütze Reden geführt habe. Ach, lieber Gott, verzeihe mir diese Sünde und wende von mir ab die schwere Strafe, die du drohest allen denen, welche ihre Zunge mißbrauchen dir zur Unehre und dem Nächsten zum Schaden. Gib mir aber deine Gnade, daß ich forthin meine Zunge im Zaum halte, lehre mich, daß ich ein Schloß an meinen Mund lege und ein festes Siegel auf mein Maul drücke, daß ich dadurch nicht zu Fall komme. O Herr Jesu Christe, du zartes Lämmelein Gottes, in dessen Mund kein Betrug erfunden worden ist, der du nicht wieder gescholten hast, da du gescholten worden, noch gedrohet hast, da du littest, du hast deinen Mund nicht aufgethan, wie ein Lamm, das zur Schlachtbank geführet wird, und wie ein Schaf, das verstummet vor seinem Scherer und seinen Mund nicht aufthut. Dadurch hast du nicht allein für unserer ersten Eltern und meine unnützen vergeblichen Worte gebüßt, sondern auch die Verschwiegenheit und geduldiges Stillschweigen mich lehren wollen. Ach, lieber Heiland, gib mir einen stillen, geduldigen Geist und einen verschwiegenen Mund, daß ich mit Worten mich nicht vergreife. Verleihe mir die Gnade, daß ich meine Gedanken und Zunge könne im Zaum halten und mein Herz mit Gottes Wort züchtigen und mein nicht schonen, wo ich fehle, auf daß ich nicht Sünde anrichte und großen Irrthum stifte und viel Uebels begehe. O Heiliger Geist, der du dich in feurigen Zungen über die Apostel ausgegossen und ihre Zungen also regieret hast, daß sie die großen Thaten Gottes verkündigten, erfülle auch meine Zunge mit dem himmlischen Feuer der Weisheit, daß ich dieselbe nicht mißbrauche zu unnützen Reden, sondern nach der Vermahnung, welche du mir durch dein theures Werkzeug, den Apostel Paulum, gegeben hast, allein dasjenige rede und fürbringe, was nützlich zur Besserung ist, da es noth thut. Thue meine Lippen auf, daß ich Gottes Ruhm verkündige und des Nächsten zeitliches und ewiges Heil treulich befördere. Das wollest du, o treuer Gott, mir aus Gnaden um Christi willen geben und verleihen. Amen.

\quellseite{300}

\section{354. Klage über das angeerbte Verderben}

Heiliger Gott, du gerechter Richter, ich weiß, daß ich in Sünden empfangen und geboren, ich weiß, daß ich aus unreinem Samen in Mutterleibe gebildet bin. Dieses Sündengift hat meine ganze Natur so gar verderbet und zerstört, daß keine Kraft meiner Seele von seiner Ansteckung frei geblieben ist. Die heilige Beilage des göttlichen Ebenbildes, welche in unserem Stammvater auch mir nit anvertrauet war, ist in mir verloren gegangen. Es ist keine Kraft mehr vorhanden, einen Anfang in deiner Erkenntniß, Furcht, Liebe und Vertrauen zu machen, kein Vermögen mehr, deinen Geboten Gehorsam zu leisten. Mein Wille ist vom Gesetz abgewendet, und der Sünde Gesetz in meinen Gliedern, das da widerstreitet dem Gesetz in meinem Gemüthe, macht, daß meine ganze Natur verdorben und verkehrt ist. Ich armer, elender Mensch fühle den Anlauf der Sünde, die in meinen Gliedern wohnet und mir immer anklebet, ich fühle das Joch der mich schwer drückenden bösen Lust. Denn obwohl im Bad der Taufe durch den Geist der Gnade wiedergeboren und erneuert, bin ich doch nicht gänzlich frei von dem Joch und der Gefangenschaft der Sünde, sintemal die in mir verborgene liegende bittere Wurzel immer wieder hervorzuwachsen trachtet. Der Sünde Gesetz wütet in meinem Fleische und will mich gefangen nehmen. Ich bin voller Zweifel, voll Mißtrauen, voller Sucht nach eigener Ehre; aus dem Herzen kommen arge Gedanken, die mich vor dir gänzlich verunreinigen; aus vergifteter Quelle fließen giftige Bäche hervor. Siehe, Herr, darum nicht ins Gericht mit deinem Knecht, sondern sei mir gnädig nach deiner großen Barmherzigkeit. Der Abgrund meines Elendes rufe den Abgrund deines Erbarmens an. Für diesen Gräuel meiner unreinen Natur bringe ich dir dar die allerheiligste Empfängniß deines Sohnes. Mir ist er geborn, mir also auch empfangen; mir ist er gemacht zur Heiligung und zur Gerechtigkeit, also ist er auch mir gemacht zur Reinigung und zur Unbeflecktheit. Durch diesen deinen Sohn und um seinetwillen erbarm dich mein, du Allerhöchster, und stelle das verborgene Uebel des meiner Natur anklebenden Verderbens nicht in das Licht vor
\quellseite{301}
deinem Angesicht, sondern siehe an deinen lieben Sohn, meinen Mittler: seine allerheiligste und unbefleckte Empfängniß komme zu Hilfe meinem Elende. Amen.

\section{355. Gebet um wahre Bekehrung}

Allmächtiger Gott, barmherziger Vater! Es erschallet zwar täglich in meinen Ohren und Herzen deine Buß- und Gnaden-Stimme, die mich und alle sündige Seelen zur Bekehrung aufmuntert: Kehre wieder, du Abtrünnige, so will ich mein Angesicht nicht gegen dir verstellen, denn ich bin barmherzig, und will nicht ewiglich zürnen, allein erkenne deine Missethat, daß du wider den Herrn, deinen Gott, gesündiget hast. Aber ach! ich bin ja von mir selber, als von mir selber, nicht tüchtig, mich zu dir zu bekehren. So wenig ein Mohr seine Haut verwandeln, und ein Pardel seine Flecken ablegen kann, so wenig kann ich aus meinen natürlichen Kräften von Sünden ablassen, und Gutes verrichten, weil das Dichten und Trachten meines Herzens zum Bösen geneigt ist von Jugend auf und immerdar. Hat gleich deine Gnade so viel bei mir gewirket, daß ich mir vorgesetzt: ich wolle mich aufmachen, mich wegen meiner Sünden wieder mit dir versöhnen, und dir einen geängsteten Geist, ein zerschürftes und zerschlagenes Herz aufopfern, ach! so wirft mir doch der verfluchte Satan tausenderlei Hindernisse in den Weg. Mir widerstrebet mein eigenes sündliches Fleisch und Blut, daß ich das Gute, so ich will, nicht vollbringen, sondern in Sünden beharren, und darinnen endlich verderben soll. Darum, Abba, lieber Vater! beuge und neige ich vor dir meine Kniee in tiefster Demuth, und bitte um deiner unendlichen Barmherzigkeit willen: Bekehrte du mich, HErr! so werde ich bekehret, hilf du mir, so ist mir geholfen. Laß deine Gnade, mein Gott! an mir nicht vergeblich sein, sondern wie du das Wollen gegeben hast, also verleihe auch das Vermögen zum Vollbringen. Steure dem Muthwillen des höllischen Bösewichts, wehre dem Eingeben meines verkehrten Herzens, daß dadurch mein heiliges Vorhaben nicht gehindert und rückgängig gemacht werde. Thue meine Augen auf, daß ich meine Sünden herzlich erkenne, schmerzlich beklage und bitterlich beweine. Gib mir eine göttliche Traurigkeit
\quellseite{302}
über meine Missethaten, die da wirke eine Reue zur Seligkeit, welche mich nimmermehr gereue. Laß mich aber, o Gott! nicht unter der Sündenlast verderben, noch mit Cain und Judas verzagen, sondern wende dich zu mir, und sei mir gnädig, siehe an meinen Jammer und Elend, und vergib mir alle meine Sünden. Zeige mir doch dein liebreiches und erbarmendes Vater-Herz, welches nicht will den Tod des Sünders, sondern daß er sich bekehre, und lebe; sprich meiner Seelen tröstlich zu, daß du wollest Gnade für Recht ergehen lassen, und aller meiner Uebertretungen nicht mehr gedenken. Weise mir mit deinem göttlichen Gnadenfinger den gekreuzigten JESum, wie er um meiner Sünden willen gestorben, und um meiner Gerechtigkeit willen auferwecket worden. Oeffne mir zur sichern Zuflucht seine blutige Wunden, in welchen so mancher armer Sünder Heil und Trost gefunden hat. Schaffe auch in mir ein reines Herz, und gib mir einen neuen gewissen Geist; verwirf mich nicht von deinem Angesicht, und nimm deinen heiligen Geist nicht von mir; tröste mich wieder mit deiner Hülfe, und der freudige Geist enthalte mich in einem beständigen Vorsatz, mein sündliches Leben zu bessern, und mein Lebenlang dir in Heiligkeit und Gerechtigkeit zu dienen, wie es dir gefällig ist. Ach HErr! höre, ach HErr! sei gnädig, ach HErr! merke auf, und thue es, um Jesu Christi willen. Amen!

\section{356. Gebet um den rechten Glauben}

HErr, allmächtiger Gott, ein Vater des Lichts, bei welchem keine Veränderung noch Wechsel des Lichts und der Finsterniß ist, von dem alle guten und vollkommenen Gaben zu uns herabkommen, ich bitte dich, weil der Glaube nicht Jedermanns Ding ist, du wollest durch deinen Heiligen Geist rechtschaffene Erkenntniß deines Sohnes JESu Christi in meinem Herzen pflanzen und erhalten, auch von Tag zu Tag mehren, daß ich erfüllet werde mit Erkenntniß deines Willens in allerlei geistlicher Weisheit und Verstand, auf daß ich würdig wandele dir zu allem Gefallen, und fruchtbar sei in allen guten Werken, und wachse in denselben nach deiner herrlichen Macht in aller Geduld und Langmuthigkeit mit Freuden. Gib mir Kraft, nach dem Reichthum deiner Herrlichkeit stark zu werden an dem inwendigen
\quellseite{303}
Menschen durch deinen Geist und Christum zu wohnen durch den Glauben in meinem Herzen. O lieber Gott! dieweil niemand deinen Sohn kennt denn nur du, Vater, und niemand kennet dich, Vater, denn nur derselbige dein Sohn, und wem er es will offenbaren; so bitte ich dich, ziehe du mich zu ihm, und er ziehe mich zu dir. Gib mir Erkenntniß des Heils, die da ist in Vergebung der Sünden. Komme zu Hülfe meinem schwachen Glauben, der so klein ist wie ein Senfkorn, auf daß er zunehme, und ich in dir eingewurzelt und erbauet, dazu fest und unbeweglich sei. Frommer Gott, der du das Fünklein des Glaubens in uns angezündet und das gute Werk in uns angefangen hast, ich schreie zu dir, du wollest es auch vollführen bis an das Ende, daß wir je mehr und mehr reich werden in allerlei Erkenntniß und Erfahrung, auf daß wir seien lauter und unanstoßig bis auf den Tag Christi, erfüllet mit Früchten der Gerechtigkeit, die geschehen durch Christum Jesum zu Ehren und Lob Gottes. Erhalte, was du in uns gewirket hast, auf daß wir eine gute Ritterschaft üben, Glauben und ein gut Gewissen bewahren und nicht etwa in Kreuz und Anfechtung umgestoßen, am Glauben Schiffbruch möchten erleiden. Derhalben behüte mich, mein Gott, daß ich unter so viel Irrthümern, Spaltungen und Ketzereien dieser Welt nicht verführet werde. Bewahre mich vor Aberglauben und allem falschen Gottesdienst, und daß ich in keinem Artikel irre und zweifele. Dazu so gib auch Gnade, daß derselbe mein Glaube nicht todt, werkloß oder unkräftig sei, sondern thätig und geschäftig, der durch die Liebe ausbreche, auf daß ich das Ende des Glaubens davon bringe, nämlich der Seelen Seligkeit. Amen.

\section{357. Gebet um Erlösung vom Fluch des Gesetzes}

Herr Gott, wer kann deine Gebote halten? Je mehr du gebietest, je weniger man thut. Wir sollten dir trauen und deine Gebote halten; das thun wir nicht, finden nichts mehr am Gesetz, denn daß nichts Gutes an uns ist. Darum hat es Moses gegeben, daß es den Fluch offenbaren soll, und wenn wir denn solches sehen und fühlen, daß wir dann anheben und sagen: Nun komme, Herr, und gib uns den Segen, erlöse uns von diesem Fluch! Das
\quellseite{304}
Gesetz hilft nichts zum Gewissen, der Segen muß es allein thun. Hier im Gesetz ist eitel Tod und Gottes Zorn; nun ist das Evangelium dein Wort. Darum halte, was du gesaget hast, daß du uns den Segen gebest durch den gebenedeiten Samen, der kommen ist, der uns solches gebe und helfe vom Tode durchs Leben, von Sünde durch Gerechtigkeit. Amen.

\section{358. Danksagung für die Schöpfung}

O allmächtiger Gott, himmlischer Vater, der du nicht geschaffen, nicht geboren, sondern vor allem Anfang in Ewigkeit bist, ich anbete dich, ich ehre, lobe und preise dich, und sage dir Dank von Grund meines Herzens für alle deine Wohlthaten und sonderlich, daß du Himmel und Erde, Sonne, Mond und alle Creaturen, dazu alles, was da lebet und ist, durch deine Allmacht und göttliches Wort geschaffen hast, durch deine Weisheit regierest und erhältst. Auch danke ich dir, daß du den ganzen Erdboden sammt allen erschaffenen Creaturen durch deine Güte und Barmherzigkeit dem menschlichen Geschlechte dienstbar untergeben, und zu gut ausgetheilt hast. Insbesondere aber lobe ich dich, meinen Schöpfer und HErrn, daß du mich auch, deine Creatur, zu einem vernünftigen Menschen nach deinem Bildniß formieret und gemacht hast; mir Leib und Seele, und alle Glieder, Vernunft, Verstand und alle Sinne gegeben hast und noch bewahrest. Groß und wunderbar ist deine Güte an mir, der du mich, als ich noch im Mutterleibe verschlossen war, erhalten und errettet hast, auch von dannen herausgezogen; mich nicht blind, taub, stumm, lahm, oder gebrechlich hast lassen geboren werden. Viel und mannigfaltig ist deine Barmherzigkeit an mir, daß du mich nicht zu einem Viehe, welches ohne Verstand und Vernunft ist, gemacht und erschaffen hast. Wer kann die großen Thaten des HErrn ausreden, und alle seine löblichen Werke preisen? Wer kann alle sein Lob erzählen? Niemand ist, der alle seine Wohlthaten aussprechen mag. Und ob ich armer Sünder dich nicht loben kann, so viel und hoch als du würdig und ich schuldig bin, so will ich doch darum nicht schweigen, sondern deinen heiligen Namen ohne Unterlaß rühmen, so viel ich kann und mir möglich ist. Ich will deine Gerechtigkeit, Gnade
\quellseite{305}
und große Barmherzigkeit rühmen und preisen, so lange ich lebe; und weil ein lebendiger Athem in mir ist, will ich nicht vergessen, was du Gutes an mir erzeiget hast. Mein Mund soll verkündigen deine Gerechtigkeit, täglich dein Heil, die ich nicht alle zählen kann. Meine Zunge, die du mir erschaffen hast, soll deiner Ehren und deines Preises voll sein täglich und immer sagen: Ich danke dem HErrn in seinem Heiligthum, ich lobe ihn in der Veste seiner Macht, ich preise ihn in seinen Thaten, ich rühme ihn in seiner großen Herrlichkeit. Alles, was Odem hat, danke dem HErrn, Halleluja.

\section[359. Danksagung für alle von Mutterleib an …]{359. Danksagung für alle von Mutterleib an erzeigten Wohlthaten}

HErr, ehe ich noch war, lebete, webete, und etwas thun konnte, warst du über mir im Mutterleibe, nahmst dich meiner, als deines Geschöpfs, gnädiglich an, sorgetst herzlich für mich und erhieltest mich wunderbarer Weise. Vielmehr thust du solches, treuer Menschenhüter, an mir, der ich nun ein Mensch zur Welt geboren, lebe, gehe, stehe, schaffe und durch dein Wort dich kenne, ob es wohl vor Menschenaugen viel anders scheinet, und mein alter Adam, der mir am Hals bis in die Grube hänget, das Widerspiel fühlet. Es scheine aber und fühle sich, wie es wolle, ich kehre mich nicht dran, lasse mich's auch nicht irren, sondern halte mich an dein Wort, daß du mein HErr von Mutterleib an bist; das trägt und fehlet nicht. Darauf verlasse ich mich, erwecke und stärke dadurch meinen Glauben, welcher nicht aufs Sichtbare, das zugegen ist, stehet, sondern das, welches unsichtbar ist, durch Hoffnung in Geduld erwartet. Gelobet seist du, mein HErr und mein Gott, in Ewigkeit. Amen.

\section{360. Danksagung für die Erlösung}

O HErr Jesu Christe, des lebendigen Gottes Sohn, der du bist das Ebenbild und der Glanz deines ewigen Vaters, und das Ebenbild seines Wesens, nicht gemacht noch geschaffen, sondern von ihm vor aller Ewigkeit geboren, in einem göttlichen unzertrennlichen
\quellseite{306}
Wesen, gleicher Majestät und Herrlichkeit, ich anbete dich, ich ehre, lobe und preise dich, und sage dir Dank für alle deine Wohlthaten, sonderlich aber, daß du zur bestimmten Zeit nach deinem göttlichen Willen und Wohlgefallen von Maria, der reinen keuschen Jungfrau, wahre menschliche Natur an dich hast genommen, mich armen verlornen und verdammten Menschen erlöset, erworben, gewonnen von allen Sünden, vom Tod und von der Gewalt des Teufels, nicht mit Gold oder Silber, sondern mit deinem heiligen theuren Blut, und mit deinem unschuldigen Leiden und Sterben, und solches aus lauter Gnade und Barmherzigkeit, ohne alle meine Werke, Verdienst und Würdigkeit. O lieber HErr Gott, wie groß ist deine Liebe gegen das menschliche Geschlecht, sintemal du Gerechter für die Ungerechten und Gottlosen gestorben bist. Wie kann doch größere Liebe immer sein, denn so jemand sein Leben lässet für seine Freunde? Nun aber beweisest du deine allerhöchste überschwängliche Liebe an uns, daß du für uns gestorben bist, da wir noch Sünder und deine Feinde waren, und hast uns durch deinen Tod mit deinem Vater versühnet, seinen gerechten Zorn gestillet und ausgetilget die Handschrift, so wider uns war, welche durch Satzungen entstanden und uns entgegen war, und hast sie aus dem Mittel gethan und an das Kreuz geheftet. Du hast ausgezogen die Fürstenthümer und die Gewaltigen der Finsterniß, sie Schau getragen öffentlich, und einen Triumph aus ihnen gemacht durch dich selbst. Dazu mit deiner heilmärigen Auferstehung von den Todten hast du wiederbracht Unschuld und ewiges Leben, wahre Gerechtigkeit und Seligkeit, mir und allen, die solches von Herzen glauben. O du treuer Heiland, laß dein heilig bitter Leiden an mir armen Sünder nicht übel angewendet sein. Hilf, daß ich hinfort nicht mir selbst lebe, sondern dir, der du für mich gestorben und wieder auferstanden bist, daß ich von Sünden aufstehe und lebendig werde in dir und dein eigen sei, in deinem Reich unter dir lebe, dir diene in ewiger Gerechtigkeit, der du bist aufgefahren gen Himmel, sitzest zur Rechten deines Vaters, immer und ewiglich. Amen.

\quellseite{307}

\section[361. Gebet zu Christo Jesu um seine Zukunft …]{361. Gebet zu Christo Jesu um seine Zukunft zum jüngsten Tage}

Herr Jesu Christe, ob zwar von der Stunde deiner Zukunft niemand weiß, auch die Engel im Himmel nicht, sondern allein der Vater, der es seiner Macht vorbehalten hat: so wird es doch mit dieser Welt ein Ende nehmen und ihr Wesen vergehen; du wirst kommen mit Feuerflammen, Rache zu geben über die, so dich, Gott, nicht erkennen und deinem Evangelio nicht gehorsam sind. Und damit wir nicht zweifeln, so hast du die Zeichen deiner Zukunft aus treuem Herzen uns wissen lassen und namhaft gemacht; wie sie denn jetund im vollen Schwange gehen und also gewiß beweisen, daß das Ende aller Dinge nahe ist. Denn es geschehen große Zeichen an Sonne, Mond und Sternen, die fallen vom Himmel und verlieren ihren Schein. Man hört von Krieg und Kriegsgeschrei. Ein Volk empört sich über das andere und ein Königreich über das andere. Es geschehen Erdbeben hin und wieder. Es ist theure Zeit und Hungersnoth. Die Ungerechtigkeit nimmt, je länger, je mehr überhand. Es erkaltet in vielen die brüderliche Liebe. Es sind so schreckliche böse Zeiten, daß den Leuten angst und bange ist. Sie verschmachten und vergehen schier vor Trübsal und Jammer, der in der Welt ist. In der Kirche sind viele Männer aufgestanden, die verkehrte Lehre geredet und dein Wort verfälscht haben, und dein Wort muß noch bei vielen Ketzerei heißen. Im weltlichen Regiment geht oft Gewalt für Recht. Das Recht wird in Galle und die Frucht der Gerechtigkeit in Wermuth verwandelt. Das Böse muß gut und das Gute bös und schwarz muß weiß heißen. Im Hausregiment ist große Untreue, Ungehorsam, Uneinigkeit, Zwietracht, Zank und Hader, und wird insgemein ein gottloses Leben geführt, das man doch nicht für Sünde achtet. Welches alles sind Zeichen des zunehmenden jüngsten Tages. Weil nun dieselben alle vor Augen schweben, so hilf, o Herr, gnädiglich, daß wir es zu Herzen nehmen, und nicht sicher sein, und etwa mit den thörichten Jungfrauen übereilt und zerschmettert werden, sondern daß wir allezeit wacker sein und beten, Gutes thun und nicht müde werden, damit wir deinem gestrengen Urtheil und Gericht entfliehen, und würdig sein mögen, zu stehen
\quellseite{308}
vor deinem heiligen Angesicht, wenn du kommen wirst in den Wolken mit großer Kraft und Herrlichkeit und deine Engel senden und versammeln deine Auserwählten von den vier Winden und von dem Ende der Erde bis zum Ende der Himmel. HErr, wir warten täglich auf dein Heil. So komme nun, o HErr JEsu, und mache alles unsers Jammers ein Ende und nimm uns sammt dem gläubigen Schächer in das Paradies auf. Komm, HErr JEsu, komm. Amen.

\section{362. Gebet um baldige Zukunft zum Gericht}

Lieber HErr JEsu Christe, stärke und vollbringe dein Werk, das du in uns angefangen hast, und eile ja herzu mit dem herrlichen Tage unserer Erlösung, den wir von Gottes Gnaden herzlich begehren, drum seufzen und drauf warten, in einem rechten Glauben und gutem Gewissen, damit wir gedienet haben der undankbaren Welt, an welcher keine Besserung zu hoffen ist, sondern ist ein Feind zugleich ihrer eigenen und unserer Seligkeit. Komm, lieber HErr JEsu, und wer dich liebet, spreche: Komm, lieber HErr JEsu. Amen.

\section{363. Sehnliches Seufzen nach dem jüngsten Tage}

HErr JEsu Christe, eile doch und bringe herzu den seligen Tag, da die Hoffnung unserer herrlichen Erlösung soll erfüllet werden, denn eben darum hast du uns heißen bitten im Vaterunser: Dein Reich komme! Weil du uns denn solches zu bitten befohlen hast, so gib auch Gnade und hilf, daß wirs thun und daneben festiglich glauben, daß wir endlich zu solcher Herrlichkeit kommen werden. Gib auch, daß derselbe fröhliche, selige Tag unserer Erlösung und Herrlichkeit bald komme, und wir solches alles erfahren, wie wirs jetzt im Wort hören und glauben. Amen.

\section{364. Betrachtung des jüngsten Gerichts}

HErr JEsu, wer will den Tag deiner Wiederkunft zum Gerichte erleiden, da du nicht mehr wie in deiner ersten Zukunft ins Fleisch, in armer, verachteter Knechtsgestalt, sondern in den Wolken des Himmels mit großer Kraft, in der Herrlichkeit deines Vaters mit allen heiligen Engeln und Erzengeln erscheinen und
\quellseite{309}
auf dem Stuhl deiner Majestät zur Rechten der Kraft sitzen wirst. O wie schrecklich wird es alsdann zugehen! Denn da wirst du deine Engel senden mit hellen Posaunen, daß sie deine Auserwählten sammeln von den vier Winden, von einem Ende des Himmels zum andern, und die Bösen von ihnen scheiden und sie in den Feuerofen werfen, da sein wird Heulen und Zähneklappen. O HErr, wer wird bestehen, wenn du wirst erscheinen? Denn da wirst du deinen Eifer nehmen zum Harnisch, und die Kreatur rüsten zur Rache über deine Feinde. Da wirst du Gerechtigkeit anziehen zum Krebs und das ernste Gericht aufsetzen zum Helm. Heiligkeit wirst du nehmen zum unüberwindlichen Schild, und den strengen Zorn wetzen zum Schwert. Da wirst du den ganzen Erdkreis richten mit Gerechtigkeit, und einem jeden geben, nach dem er gehandelt hat bei Leibes-Leben, es sei gut oder böse. Da wird alles haarfein herfürkommen, was ein jeder Mensch die ganze Zeit seines Lebens gedacht, geredet oder gethan hat, da werden einem jeglichen seine Sünden, die er in den innersten Winkeln begangen, an der Stirn geschrieben stehen. Denn da wirst du alles an das helle Tageslicht bringen, was im Finstern verborgen ist, und den Rath der Herzen offenbaren, also daß wir auch von einem jeglichen unnützen Worte, so aus unserm Munde gangen, Rechenschaft werden geben müssen. O ich elender Mensch, wie will ich in solchem ernsten Gericht bestehen? wie will ich alle meine bösen Gedanken, sündliche Gebehrden, unnütze Reden und untüchtige Werke, deren unzählich viel sind, verantworten? Ach, wie viel sind der Sünden, die ich nicht weiß, ja die ich jetzt nicht für Sünde, sondern für gute Werke achte, welche vor deinem Gericht recht offenbar und kohlschwarz erscheinen werden? Ach, wehe mir, wie übel habe ich gethan? Wo soll ich fliehen hin, da ich möge bleiben? Verbergen kann ich mich nicht: denn führe ich gen Himmel, so bist du da, nähme ich Flügel der Morgenröthe und bliebe am äußersten Meer, so würde mich doch deine Hand daselbst führen und deine Rechte mich halten. Ja, wenn ich mich gleich in die Hölle vergrübe würde mich doch deine Hand von dannen herholen. Entschuldigen kann ich mich viel weniger. Ach, wo soll ich deinen Rath und Hülfe suchen? Bei dir, HErr Jesu Christ allein, der du nicht allein
\quellseite{310}
mein Richter, sondern auch mein Advokat und Fürsprecher sein wirst, weil du mein Bruder und Heiland bist, und um meinetwillen dich richten lassen, auf daß ich im jüngsten Gericht mit Freuden bestehen möge. Lob sei dir, o mein liebster Immanuel, Hosianna, gelobet sei, der da kommt im Namen des HErrn, mein König und mein Bräutigam. Komm, HErr JEsu, komm bald. Amen.

\section{365. Danksagung für die Heiligung}

O Gott Heiliger Geist, der du ausgehet vom Vater und gesandt wirst von dem Sohn, bist mit ihnen gleichförmig, mit ewig in einem göttlichen und ungetheilten Wesen, ich anbete dich, ich ehre, lobe und preise dich, und danke dir aus Grund meines Herzens für alle deine Wohlthaten, und sonderlich, daß du mich durch deine Gnade zum heiligen christlichen Glauben berufen und gebracht, dich selbst in meiner Taufe eingegossen hast und ohne Aufhören in mir wirkest deine Werke. Ich bitte dich, weil ich nicht aus eigener Vernunft noch Kraft an JEsum Christum, meinen HErrn, glauben oder zu ihm kommen kann, du wollest mich lehren, in alle Wahrheit leiten und führen, den wahren Glauben in mir mehren und ernähren, mich mit rechtem Verstand deines Wortes begaben, in wahrer Anrufung und rechtem Vertrauen heiligen, innerlichen Frieden des Gewissens, mit ewigwährendem Trost neben aller Freudigkeit und Beständigkeit geben und schenken. Schreibe du JEsum Christum in mein Herz, der du bist der lebendige Finger Gottes, vergewissere mich mit deinem Zeugniß, tröste mich wider alle Furcht und Blödigkeit des Fleisches. Sei du mein treuer Beistand wider des bösen Feindes Anklagung; stärke mich wider die Verzagung, auf daß ich also im Glauben durch deine Gnade der Verheißung versiegelt werde; denn du bist das Pfand meines Erbes zu meiner Erlösung. Verleihe mir auch rechte Andacht und Seufzen zum Gebet, daß ich mit wahrer Zuversicht und kindlichem Vertrauen schreie und spreche: Abba, lieber Vater! Gib mir rechte Einfalt des Herzens, ohne alle Galle und Bitterkeit, der du in der Gestalt einer Taube über unserm HErrn JEsu Christo am Jordan in seiner heiligen Taufe erschienen bist. Entzünde in mir das Feuer der inbrünstigen Liebe, der du dich in feurigen Zungen
\quellseite{311}
über die heiligen Apostel schwebend hast sehen lassen, damit ich deine Salbung empfange. Erneuere mich an dem inwendigen Menschen, daß ich durch dich, der du im Wort und Sacrament kräftig bist, möge wiedergeboren werden. Behüte mich vor Abgötterei und Lügen, vor allem Irrthum und Ketzerei. Nimm hinweg die Unreinigkeit meines Gemüths. Lösche und dämpfe in mir alle fleischliche Begierde. Laß die sündlichen Lüste nicht kräftig werden in meinem Leibe, sondern mache du deine Wohnung in mir und herrsche in allen meinen Gliedern, begleite mich in das rechte Vaterland. Amen.

\section{366. Dankstagung für die Auferstehung der Todten}

O du Herzog unserer Seligkeit, Christe Jesu, mein Herr und mein Gott, der du selbst gesprochen: Ich bin die Auferstehung und das Leben; wer an mich glaubet, der wird leben, ob er gleich stürbe, stärke meinen Glauben in diesem Geheimniß von Tag zu Tage, daß ich weder mit Worten noch Gedanken erfinden werde unter denen, die da halten, es sei keine Auferstehung, sondern festiglich glaube, daß, wie du von den Todten auferstanden und der Erstling worden unter denen, die da schlafen; also auch ich im Grabe nicht bleiben werde, sondern daß du mit dem himmlischen Vater und Heiligen Geist dermaleins wirst verklären meinen nichtigen Leib, daß er ähnlich werde deinem verklärten Leibe nach der Wirkung, damit du kannst auch alle Dinge dir unterthänig machen. Ach, mein trautester Jesu, hoffeten wir allein in diesem Leben auf dich und glaubeten nicht, daß nach diesem Leben ein besseres und ewiges Leben zu erwarten wäre, so wären wir die elendesten unter allen Menschen. Nun du aber durch deine Auferstehung das Leben und ein unvergänglich Wesen ans Licht bracht, auch überdieß klärlich gesagt: Ich lebe und ihr sollt auch leben, sowohl dem Leibe als der Seele nach, so freuet sich mein Herz und meine Ehre ist fröhlich, denn ich weiß, daß mein Erlöser lebet, und er wird mich hernach aus der Erden auferwecken und werde darnach mit dieser meiner Haut umgeben werden und werde in meinem Fleisch Gott sehen; denselben werde ich mir sehen und meine Augen werden ihn schauen und kein Fremder. Will schon zu Zeiten mein
\quellseite{312}
Glaube sinken, und anheben an diesem Geheimniß zu zweifeln, so hilf mir selber und vertritt mich bei Gott, daß mein Glaube nicht aufhöre. Gib nicht zu, daß das glimmende Töchtlein in mir vollends werde ausgelöscht, sondern stärke mir den Glauben, daß ich sowohl deiner Allmacht, welcher kein Ding unmöglich, traue, als auch immer dir mit gläubigem Herzen danke und sage: Wir haben in dir einen Gott, der da hilft, und einen HErrn HErrn, der, kraft seiner siegreichen Auferstehung, vom Tode errettet, dir sei Ehre von Ewigkeit zu Ewigkeit. Amen.

\section[367. Gebet nach den sieben Bitten des heil. …]{367. Gebet nach den sieben Bitten des heil. Vaterunsers}

O lieber himmlischer Vater, du dreieiniger Gott, wir, deine Kinder, bitten dich von Herzen, du wollest uns vor falscher Lehre und sündlichem Leben behüten, getreue Lehrer geben und erhalten, daß sie dein liebes Wort rein und lauter predigen, und deine heiligen hochwürdigen Sacramente recht verwalten, und wir auch heilig und gottseelig darnach leben mögen in deiner Furcht. Schenke uns deinen Heiligen Geist, Gerechtigkeit, Friede und Freude des Herzens, und bewahre in uns die Beilage bis auf den Tag Jesu Christi. Regiere uns, daß wir thun, was du in deinem Worte von uns haben willst, daß wir an Christum glauben, unser Leben bessern, geduldig sein und beständig verbleiben. Hindere und brich allen bösen Rath und Willen, als des Teufels, der Welt und unseres Fleisches, auf daß dein Wille von uns auf Erden, wie von den lieben Engeln im Himmel unverhindert geschehe. Tröste und stärke alle Kranken, Betrübten und Angefochtenen, und erhalte sie in christlicher Geduld und wahrem Glauben bis ans Ende. Verleihe auch unserer Obrigkeit Gnade, Geist und Weisheit, die Unterthanen nach deinem Wohlgefallen recht zu regieren. Behüte uns vor Aufruhr, Krieg und Blutvergießen, vor Feuer und Wassersnoth, Hagel, Ungewitter, Mißwachs, Pestilenz und theurer Zeit. Laß dir auch den lieben Hausstand, und darinnen Mann und Weib, Eltern und Kinder, Herren und Gesinde, arme Wittwen und Waisen, Junge und Alte treulich empfohlen sein, und laß uns genießen mit Gebet und Danksagung unser täglich Brod. Ach barmherziger Gott und Vater, siehe nicht
\quellseite{313}
an unsere Sünde und Missethat, vergib uns dieselbe und erlasse uns auch die wohlverdiente Strafe, erneuere unsere Herzen und Sinne, daß wir den Sünden je mehr und mehr absterben und frömmer werden, daß wir uns unter einander auch gerne vergeben, lieben und in Frieden leben. Behüte uns auch vor den schweren Anfechtungen des Teufels, der Welt und unseres Fleisches und stärke uns durch deines Geistes Kraft, daß wir solchen unsern geistlichen Feinden Widerstand thun, Alles wohl ausrichten und den Sieg erlangen mögen. Dieweil auch, leider! unser Jammer allhier mancherlei und groß ist, so bitten wir dich, o treuer Gott und Vater, du wollest dich unser annehmen, und in keiner Noth verlassen, sondern unsern Jammer und Elend lindern und mindern, uns darinnen trösten und helfen, ja von allem Uebel uns endlich erlösen, ein seliges Sterbestündlein bescheren und alsdann im rechten Glauben mit Gnaden aus diesem Jammerthal zu dir in die ewige Freude und Seligkeit aufnehmen. Amen, in Jesu Namen Amen.

\section[368. Gebet um Beständigkeit im Bekenntniß der …]{368. Gebet um Beständigkeit im Bekenntniß der Wahrheit in Verfolgungszeit}

Halte was du hast, daß niemand deine Krone nehme. Offenb. 3, 11.

Allerliebster JESU, mein Herr und mein Gott, du hast selbst in den Tagen deines Fleisches gesprochen: „Wer mich bekennet vor den Menschen, den will ich bekennen vor meinem himmlischen Vater. Wer mich aber verleugnet vor den Menschen, den will ich auch verleugnen vor meinem himmlischen Vater.“ Du weißt alle Dinge, du weißt, daß ich dich lieb habe, du weißt auch meine Trübsal und die Lästerung von denen, die da sagen, sie sind Christen und sind’s nicht, sondern sind des Satans Schule. Ich bitte dich herzlich, heilige mich in deiner Wahrheit, dein Wort ist die Wahrheit, und gib, daß ich mich nicht lasse wägen noch wiegen von allerlei Wind der Lehre, sondern daß ich stehe in deiner Wahrheit, die ich erkannt, und bleibe in dem, das ich gelernet habe. Mache mich stark durch deine Gnade, daß ich deinen Namen trage vor den Völkern, Fürsten und Königen, und ja nicht sei ein Rohr, das der Wind hin und her wehet. Denn weil ich wohl weiß, daß keiner, außer der beharret bis ans Ende, selig werden soll, so erhalte mein
\quellseite{314}
Herz bei dem einigen, und wende von mir allen falschen Weg, daß sich mein Herz nicht von dir abkehre, noch sich heimlich bereden lasse, und damit eine Wurzel in mir werde, die Galle und Wermuth trage. Denn was würde mir's doch helfen, wenn ich schon die ganze Welt gewönne, und nähme doch Schaden an meiner Seele? Wollen meine Widersacher und Verfolger gar an mich, mein Fleisch zu fressen und mich zu einer widrigen Religion zu zwingen, so gib selbst Kraft und Muth, daß ich mich nicht fürchte vor denen, die den Leib tödten, und die Seele nicht mögen tödten, sondern vielmehr vor dem, der Leib und Seele verderben kann in die Hölle, und daher mein Leben nicht lieber bis in den Tod. Und weil der Geist willig, aber das Fleisch schwach ist, so vollführe du selbst das gute Werk, so du in mir angefangen hast, und laß mich dir getreu verbleiben bis in den Tod. Sei in mir Schwachen mächtig, daß ich weder Leben noch Tod, noch sonst einige Gewalt mich scheiden lasse von der Liebe, die da ist in dir, und bringe mich endlich dahin, da ich befreiet von aller Verfolgung, mit Herz und Mund singen kann: Es hat überwunden der Löwe, der da ist vom Geschlecht Juda. Dem, der auf dem Stuhl sitzt, und dem Lamm sei Lob und Ehre und Preis und Gewalt von Ewigkeit zu Ewigkeit. Amen.

\section[369. Ernstes Gebet um Bewahrung vor …]{369. Ernstes Gebet um Bewahrung vor Verstockung und Abfall}

O Herr Gott, so wir ja sündigen sollen, so laß uns andere Sünden thun denn solche, da man hinfähret und scheucht den Segen als einen Fluch, und ein verstockt, verblendet und verhärtetes Herz bekommt, das weder siehet noch höret, und ihm schlechtes nichts sagen läßt, und meinet, es gehet im Segen und nicht im Fluch. Lieber Gott, laß uns in die Sünde nicht fallen, so die offenbare Wahrheit nicht leiden will; denn da ist weder Rath, noch Hülfe, noch Entschuldigung, und der Zorn endlich angegangen. Davor behüte uns, lieber treuer Gott. Amen.

\section{370. Gebet um ein seliges Ende}

O gütiger Gott, du hast dem Menschen ein Ziel gesetzet zu leben, welches er nicht kann übergehen. Denn er hat seine bestimmte
\quellseite{315}
Zeit, die Zahl seiner Monden stehet bei dir, alle unsere Tage hast du gezählet, welche doch schnell dahin fahren wie ein Strom, nicht anders, als flögen wir dahin. Alle unsere Jahre sind wie ein Rauch oder Schatten, der plötzlich vergehet. Der Mensch ist doch wie Gras, welches bald verdorret, und wie eine Blume auf dem Felde verwelket, so lehre mich nun, o gütiger Gott, erkennen und zu Gemüthe führen, daß es ein Ende mit mir haben muß und mein Leben ein Ziel hat, und ich davon muß. Siehe, meine Tage sind wie eine Hand breit bei dir und mein Leben ist wie nichts vor dir. Wie gar nichts sind alle Menschen, die doch so sicher leben. Herr, lehre mich bedenken, daß ich sterben muß und allhier in dieser Pilgerschaft keine bleibende Statt habe. Thue mir kund mein kurzes und vergängliches Wesen, daß ich oft und viel an mein Ende gedenke, auf daß ich in dieser Welt nicht mir selbst, sondern dir lebe und sterbe, damit ich im Glauben wacker und fröhlich erwarte den Tag meiner Heimfahrt, und der Erscheinung deines lieben Sohnes Jesu Christi und zu derselbigen mit heiligem Wandel und gottseeligem Wesen eile. Begnadige mich, mein Gott, mit einem seligen Abschiede, wenn mein Stündlein herzunah et, daß ich selig sterbe, ein vernünftiges Ende nehme in wahrer Erkenntniß, daß mein Verstand und Sinn nicht verrücket werden und ich nicht aberwitzige Rede oder Lästerworte wider dich, meinen Herrn, oder wider meine Seligkeit führe. Behüte mich vor einem bösen schnellen Tod und vor der ewigen Verdammniß. Laß mich nicht plötzlich und unversehens mit meinem letzten Stündlein überfallen werden, sondern daß ich mich zuvor mit wahrer Buße und rechtem Glauben bereite. Und wenn dasselbe kommt, so mache mich freudig und unverzagt zu dem zeitlichen Tode, der mir nur die Thür aufthut zum ewigen Leben, und laß mich, deinen Diener, alsdann im Frieden fahren, denn meine Augen haben deinen Heiland gesehen, welchen du bereitet hast vor allen Völkern, ein Licht, zu erleuchten die Heiden und zum Preis deines Volkes Israel. Gib, daß mein letztes Wort sei, welches dein lieber Sohn am Kreuz gesprochen hat: Vater, in deine Hände befehl ich meinen Geist. Und wann ich nimmer reden kann, so erhöre doch mein letztes Seufzen durch Jesum Christum. Amen.

\quellseite{316}

\section[371. Gebet um Bewahrung vor einem bösen …]{371. Gebet um Bewahrung vor einem bösen schnellen Tode}

O allerliebster Jesu, o mein HErr und Gott! durch alle deine Arbeit und Schmerzen, durch dein kostbarliches Blut, durch deine letzten Worte, die du, o mein allersüßester Jesu, am Kreuze gesprochen: Mein Gott, mein Gott, warum hast du mich verlassen! durch den starken Schrei, da du zu deinem himmlischen Vater geschrieen: Vater, in deine Hände befehle ich meinen Geist! durch- dieses Alles bitte ich dich inbrünstiglich, du wollest mich vor einem bösen, geschwinden und schnellen Ende gnädiglich behüten. Deine Hände, o mein Erlöser, haben mich geschaffen und gemacht, sammt Allem, was an und um mich ist. Ach lieber HErr, stürze mich nicht so jählings. Gib und verleihe, bitte ich, geraume Zeit zur Buße. Gib mir ein glückseliges Sterbestündlein, daß ich in deiner Gnade hinscheide, damit ich dich von ganzem Herzen lieben, loben und benedeien möge in Ewigkeit. Jedoch stehet es Alles, o gütiger HErr Jesu, in deiner Hand, und ist niemand, der deinem Willen mag widerstehen. Mein Leben hänget an deinem Willen, das wird sich enden, wann du selber willst. So will auch ich anders nichts, du mein lieber Gott, denn nur allein dieses, was du willst, wo du willst, wann du willst und wie du willst, daß ich sterbe, also geschehe es, dieses Alles ergebe ich deiner Fürsichtigkeit. Keinen Ort, keine Zeit, keinen Tod, er sei so schmählich und so schmerzlich, als er immer wolle, begehre ich auszuschließen. Nur dies allein bitte ich, o Christe, mein Gott, daß ich keines jähen, unversehenen Todes sterbe. Jedoch nicht mein, sondern dein Wille geschehe. Ist es dann vor deinem Angesicht also gefällig, daß ich jählings dahinstürbe, so bin ich auch nicht dahinder. Dein Wille, o mein Gott, geschehe in allen Dingen. Denn ich bin der gänzlichen Hoffnung und Zuversicht, ich werde durch deine große Barmherzigkeit, durch welche ich dieses Alles inniglich bitte, alsdann in deiner Gnade und Liebe dahin sterben, und wenn ich darin sterbe, wird mich kein jäher Tod von dir mehr scheiden können. Denn der Gerechte, ob er gleich zeitlich stirbet, ist er doch in der Ruhe. Dieses ist kein unversehener Tod, dessen Leben behutsam und vorsichtig gewesen. Wenn ich alsdann ja
\quellseite{317}
nicht soll Raum und Zeit haben (welches, o mein Gott, dir allein bewußt), daß ich mich dir könnte befehlen: siehe, so thue ich solches jetzt, und mit möglichster Demüthigkeit und inbrünstigem Herzen rufe ich aus der Tiefe zu dir gen Himmel: Erbarme dich mein, o Gott, nach deiner großen Barmherzigkeit. Dein Wille geschehe, o Herr, wie im Himmel, also auch auf Erden. In deine Hände befehle ich meinen Geist. Du hast mich erlöset, Herr, du getreuer Gott. Es hendrein und loben dich, o Gott, alle Geschöpfe. In dich, o Herr, hoffe ich, du wirst mich nicht lassen zu Schanden werden in Ewigkeit. Amen.

\section{372. Gebet von der heil. Taufe}

Ach Herr Jesu Christe, der du bist der Anfang, Mittel und Ende unsrer Seligkeit, und das heilige Sacrament der heiligen Taufe gestiftet hast und dasselbe bezeugt in dem Werk der Erlösung, da aus deiner eröffneten Seite Blut und Wasser geflossen, daher St. Johannes sagt: Drei sind auf Erden, die da zeugen: der Geist, das Wasser und das Blut und die drei sind beisammen; hast auch dies Sacrament durch deine eigene Taufe im Jordan als durch dein selbst Exempel bestätiget und geheiliget: ich danke dir herzlich, daß du mich durch dies Sacrament deiner heiligen christlichen Kirche einverleibet hast und mich dadurch in die Gemeinschaft aller deiner himmlischen und ewigen Güter gesetzt, so du durch dein bitter Leiden und Sterben erworben hast, hast mich dadurch von dem greulichen Gift und Aussatz der Erbsünde gereiniget und abgewaschen, alle meine Sünde, Schuld und Missethat aus Gnaden vergeben und hast mich gereiniget durch das Wasserbad im Wort und alle meine Uebertretung getilget, daß nicht eine ist übergeblieben, die mich könnte verdammen, daß nun nichts Verdammliches ist an denen, die in Christo Jesu sind, hast mich auch, weil die Sünde vergeben, erlöset vom ewigen Tod und von der Angst, Furcht und Schrecken des zeitlichen Todes, auch von der Gewalt des Teufels, weil du mich dir einverleibet hast, zu einem Gliede deines geistlichen Leibes gemacht, mich in dein Gnadenreich versetzt, mich mit dir vermählet und in Ewigkeit verlobet. Und weil ich durch die Taufe dich angezogen, so bin ich mit deinem ganzen
\quellseite{318}
heiligen Gehorsam, Verdienst, Gerechtigkeit, Heiligkeit und Unschuld, als mit dem Kleide des Heiles und mit dem Rock der Gerechtigkeit gekleidet, hast mich durch das Bad der Wiedergeburt und Erneuerung im Heiligen Geist neu geboren, aus einem Sünder einen Gerechten gemacht, indem ich in deinen Tod getauft und der Frucht deines Todes bin theilhaftig worden und deinem Gehorsam und Verdienst einverleibet und demnach aus einem Kinde des Zornes ein Kind der Gnaden worden, hast mich aus einem armen, elenden, verlornen Menschenkind zu Gottes Kinde gemacht, mir den Heiligen Geist der Kindschaft gegeben, durch welchen wir rufen: Abba, lieber Vater, hast mich dadurch in deinen ewigen Gnadenbund aufgenommen, und dich mit mir verbunden durch deine Verheißung, mein Vater zu sein in Ewigkeit, hast mir auch geschenket das Kindesrecht: nämlich das ewige himmlische Erbe und Seligkeit. Summa: Du hast mir in der heiligen Taufe das höchste Gut geschenkt, deinen lieben Sohn mit allen seinen Wohlthaten; Christus ist mein mit allem, was er ist und hat. Was kann mir thun die Sünde? Ist doch Christus meine Gerechtigkeit; was kann mir thun der Tod? Ist doch Christus mein Leben; was kann mir der Teufel thun? Christus ist meine Stärke und Sieg; was kann mir die Welt thun? Christus hat sie überwunden. Ja, mein HErr Christus hat mich schon selig gemacht und mir alle Seligkeit geschenkt in der heiligen Taufe. Darum warte ich in Geduld der künftigen Herrlichkeit. Was kann mir Armuth, Elend, Kreuz, Verfolgung, Verachtung schaden? Bin ich doch schon selig. Die Güter der Gnaden habe ich alle in und mit Christo empfangen und warte auf die Güter der Herrlichkeit; weil ich aber, o mein HErr Jesu Christe, durch die Taufe mit dir gestorben und begraben bin, so hilf mir, weil ich noch im Fleisch lebe, daß ich nicht nach dem Fleisch lebe, sondern daß meine Taufe täglich in mir fruchtbar sei, und wirke die Tödtung des Fleisches, daß ich täglich mit dir sterbe durch herzliche Reue und Leid; denn wer täglich in ihm selbst stirbet, der hat allezeit einen neuen Anfang seines Lebens in dir. Und weil ich dir, mein HErr und Gott, einverleibet bin als ein Pfropfreislein dem Baum des Lebens, so laß mich in dir als eine Rebe am lebendigen Weinstock Frucht bringen, nicht Früchte
\quellseite{319}
des alten Menschen, sondern des neuen Menschen, der neuen Creatur, Frucht des Geistes, und daß ich täglich bedenke, daß ich zu einem neuen Leben getauft bin, daß ich dich habe angezogen als ein Geschenk der ewigen Gerechtigkeit und Seligkeit und als ein neues Leben, dadurch du in mir und ich in dir leben und bleiben möge ewiglich; ja daß ich nimmermehr vergesse des Bundes, den du mit mir gemacht hast, denselben nicht breche, denselben nicht verleugne, weder mit Herzen noch mit Mund, noch mit einem gottlosen Leben, sondern mich dessen jederzeit freue und tröste, und darauf wider alle Anfechtung mich sicherlich verlasse und in diesem Bund, welcher ewig fest und gewiß ist, auch ewig bleiben und selig werden möge. Amen.

\section[373. Danksagung für das Amt der Schlüssel, …]{373. Danksagung für das Amt der Schlüssel, vor der Beichte zu beten}

O liebster HErr JEsu, du Sohn des lebendigen Gottes, wie hast du die Leute so lieb, daß du nicht allein um unseres Heils und ewiger Seligkeit willen unseres Fleisches und Blutes bist theilhaftig worden, und durch dein bitter Leiden und Sterben uns von Sünden erlöset, auch durch deine siegreiche Auferstehung die vor Gott geltende Gerechtigkeit erworben und dann durch deine glorreiche majestätische Himmelfahrt den weiland durch die Sünde versperrten Himmelspaß zu deinem Vater uns wiederum eröffnet, sondern auch vor deiner Auffahrt die Schlüssel des Himmelreichs deinen lieben Aposteln, im Namen und anstatt deiner Hauslehre, der ganzen Kirche, hinterlassen hast, auf daß allen denjenigen, welche nach empfangener Taufe durch falsche Lehre oder sündliches Leben sich selbst aus dem Himmelreich verstoßen, und von der Gnade gefallen, wenn sie rechtschaffene Buße thun, der Himmel durch den Lößschlüssel wiederum aufgesperret und also der Zugang zu Gott, deinem himmlischen Vater, in aller Freudigkeit und Zuversicht durch den Glauben an dich eröffnet würde, laut deiner wahrhaftigen Verheißung: „Was ihr (meine Jünger, und eure Nachfolger im öffentlichen Predigtamt) auf Erden lösen werdet, soll auch im Himmel los sein.“ Und wiederum: „Welchen ihr die Sünden erlasset, denen sind sie erlassen.“ Auf solch dein Wort will ich nun in deinem Namen mit zerknirschtem, leidmüthigem Herzen,
\quellseite{320}
in kindlichem Vertrauen zu deiner Gnade, zu deinem ordentlichen Diener und Haushalter über deine Geheimnisse hingegeben, allda zur Stärkung meines Glaubens die sonderbare Absol[u]tion oder Vergebung der Sünden aus seinem Munde an deiner Statt zu empfangen. Denn ich glaube festiglich, was der von dir verordnete Diener an deiner Statt mit mir redet und handelt, sonderlich wenn er mich als einen bußfertigen Sünder auf deinen göttlichen Befehl und theures Verdienst von meinen Sünden absolvirt, quitt, ledig und los zählet, daß solch seine Entbindung also kräftig und gewiß sei, als wenn sie von dir selbst mit eigenem Munde vom Himmel herab geschähe, und also die Stimme der Absolution wahrhaftig dein selbsteigen Wort sei. O getreuer Heiland, erhalte solchen höchsten, werthesten Schatz bis ans Ende der Welt, und regiere mein Herz, daß ich dies Gnadenwerk hoch und werth halte und den Trost, welchen ich aus dem Munde deines Dieners hören werde, durch starken Glauben fasse, annehme und behalte. Verleihe mir auch Gnade, daß ich meinen Nächsten, auch meinen Feinden von Herzen vergebe und gerne wohlthue denen, die sich wider mich versündigen. Hierauf will ich gehen in dein Haus auf deine große Güte und will anbeten in deiner Furcht. Herr, leite mich in deiner Gerechtigkeit, und richte deinen Weg vor mir her. Herr, ich rufe zu dir; eile zu mir, vernimm meine Stimme, wenn ich dich anrufe. Mein Gebet müsse vor dir taugen wie ein Rauchopfer, meiner Hände Aufheben wie ein Abendopfer. Laß mich in der heiligen Absolution hörien Worte des Lebens, dadurch ich selig werde. Ja laß mich hören Freude und Wonne, daß die Gebeine fröhlich werden, die du zerschlagen hast. Nun gehe ich; ach laß mich mit dem bußfertigen Zöllner gerechtfertigt wieder heimkommen, und hilf, daß ich ferner in ein neues Leben trete. Amen.

\section[374. Danksagung für die Einsetzung des heil. …]{374. Danksagung für die Einsetzung des heil. Abendmahls}

Allerliebster HErr JEsu Christ, was ist doch der arme Mensch, daß du seiner so inständig gedenkest und was ist des Menschen Kind, daß du dich sein so eifrig annimmst? Allerliebster HErr JEsu, in der Nacht, da du verrathen worden, hast du mehr für unser Heil, als für deine eigene Wohlfahrt gesorget. Du hast
\quellseite{321}
eingesetzt das hochwürdige Abendmahl, auf daß wir dadurch dir einverleibet, im Glauben gestärket und zum ewigen Leben ernähret würden. Lob und Dank sei dir gesagt für solche unaussprechliche Wohltat. Groß sind, Herr, deine Werke; wer ihrer achtet, der hat eitel Lust daran. Was du ordnest, das ist löblich und herrlich. Du hast ein Gedächtnis gestiftet deiner Wunder, du gnädiger und barmherziger Herr. Du gibst Speise denen, die dich fürchten, und gedenkest ewiglich an deinen Bund. Herr, du hast zu deinem heiligen Abendmahl Brod und Wein genommen. Ach hilf, daß, wie Brod und Wein den Leib nähren und erhalten, also auch unsere Seelen durch solche Wunderzeichen kräftiglich gespeiset und getränket werden. Du hast über dem Brod und Wein gedanket und gebetet: erwecke in uns auch jederzeit wahre Andacht, daß wir wie bei anderen Sachen, also vornehmlich in dieser heiligen Handlung, dich inbrünstig anrufen und dir herzlich danken. Du hast das Brod gebrochen: ach, laß mein Herz auch durch wahre Buße und Reue gebrochen und zerknirscht werden. Du hast es ausgetheilet: ach, laß mich auch von dem Segen, den du mir gibst, williglich austheilen. Du hast gesagt: „Nehmet hin und esset!“ auf solche Worte erwecke in mir eine herzliche Begierde, daß ich mit großem Verlangen hinzukomme und dein Gebot erfülle. Du hast nicht gesagt: „Das bedeutet,“ sondern das ist; nicht: Ein Zeichen meines Leibes, sondern mein Leib und zwar der Leib, der für euch gegeben wird. Dieses sind deine Worte, die hast du ja deswegen geredet, daß wir daraus deinen Willen und Meinung vernehmen sollen. Derowegen so glaube ich denselben, wie sie lauten, und lasse mich von dem klaren Buchstaben nicht abwendig machen. Du bist der Allweise und weißt wohl, wie und was du reden sollst. Du bist der Allermächtigste und kannst wohl erfüllen, was du geredet hast. Du bist der Allwahrhaftigste und wirst auch wohl leisten, was du hast zugesagt. O Herr, bekräftige in uns solchen Glauben. Heilige uns und erhalte uns in deiner Wahrheit. Dein Wort ist die Wahrheit. Bringe zurecht alle Irrende und erleuchte derer Augen, die sich durch die unvollkommene Vernunft und übel fundene Deuteleien verblenden lassen. Herr Jesu, du hast auch gesagt: „Solches thut zu meinem
\quellseite{322}
Gedächtnis.“ So gib mir nun, HErr, gnädiglich ein frisches und stetiges Gedächtnis aller deiner Wohlthaten und sonderlich deiner Bande und Wunden, deiner Schläge und Striemen, deines Leidens und Sterbens. Du hast ja in Gnaden an mich gedacht vor der Zeit, da du mich zu einem lebendigen Glied deiner Kirche und Mitbürger deines himmlischen Reichs hast auserkoren. Du hast in Gnaden an mich gedacht in der Zeit, da du für mich bist Mensch worden, meine Schuld bezahlet und noch in deinem letzten Seufzer am Kreuz deinen himmlischen Vater für mich gebeten hast. Du wirst in Gnaden an mich gedenken nach der Zeit, wenn du mich ins himmlische Paradies wirst einführen. Sollte ich denn nicht schuldig, ja nicht willig sein, hinwiederum an dich zu gedenken und zumal bei Genießung deines heiligen Abendmahls deine Liebe und Treue, deine Gnade und Gunst, deine Holdseligkeit und Barmherzigkeit zu verkündigen! Du hast es nicht lassen genug sein, für uns zu leiden und zu sterben, sondern du hast auch deinen eigenen Leib und dein eigenes Blut unter dem Brod und Wein mit dem Munde zu empfahen uns verordnet und eingesetzt. Als die Mutter des HErrn, Maria, zu ihrer Gefreundtin Elisabeth kam, sagte dieselbe mit Verwunderung: „Woher kommt mir das, daß die Mutter meines HErrn zu mir kommt?“ Alhie ist nicht die Mutter des HErrn, sondern der HErr selber. Soll ich denn nicht auch mit Verwunderung sagen: „Woher kommt mir das, daß mein HErr zu mir kommt?“ Als der HErr sich gegen den Hauptmann zu Capernaum erböt, er wollte zu ihm kommen, erkannte der redliche Soldat alsobald seine Unwürdigkeit und sprach: „HErr, ich bin nicht werth, daß du unter mein Dach gehest.“ O HErr, wie vielmehr soll ich mich demüthigen, wenn du in meinen sündhaften, unreinen Mund eingehen und mir deinen Leib zu essen und dein Blut zu trinken geben willst! Wie vielmehr sollen meine Knie sich beugen, mein Herz sich schmiegen, meine Zunge sich erheben und sagen: HErr, ich bin nicht werth, daß du in meinen Mund gehest. Aber weil es dir also gefallen, mir deinen Leib zu einer Speise und dein Blut zu einem Trank zu ordnen, sei dir Lob und Dank dafür gesagt und deine überschwängliche Barmherzigkeit müsse gerühmet und gepriesen werden immer und ewiglich. Amen, HErr JEsu, Amen, Amen.

\quellseite{323}

\section{375. Gebet wider die Sacramentirer}

Mein lieber Herr Jesu Christe, es hat sich ein Hader über deine Worte im Abendmahl erhoben. Etliche wollen, daß sie anders sollen verstanden werden, denn sie lauten. Aber dieweil sie mich nichts Gewisses lehren, sondern allein verwirren und ungewiß machen, und ihren Text in keinem Wege können noch wollen beweisen, so bin ich blieben auf deinem Text, wie die Worte lauten. Ist etwas finster drinnen, so hast du es wollen so finster haben, denn du hast keine andere Erklärung darüber gegeben, noch zu geben befohlen. So findet man in seiner Schrift noch Sprachen, daß ist, deutet, oder mein Leib, Leibzeichen heißen. Wäre nun eine Finsterniß drinnen, so wirst du’s mir wohl zu gute halten, daß ich’s nicht treffe, wie du denn deinen Aposteln zu gute hieltest, daß sie dich nicht verstanden in vielen Stücken, als du von deinem Leiden und Auferstehung verkündigtest, und sie doch die Worte, wie sie lauteten, behielten, und nicht anders macheten, wie auch deine liebe Mutter nicht verstand, da du zu ihr sagtest: Ich muß sein in dem, das meines Vaters ist, und sie doch einfältiglich die Worte in ihrem Herzen behielt, und nicht andere daraus machte. Also bin ich auch an diesen deinen Worten blieben: Das ist mein Leib, und habe mir kein anderes daraus machen wollen, noch machen lassen, sondern dir befohlen und heim gestellet, ob etwas finster drinnen wäre, und sie behalten, wie sie lauten, sonderlich weil ich nicht finde, daß sie wider einige Artikel des Glaubens streben. Amen.

\chapter{X. Wettergebete}

\section{376. Gebet zur Zeit der Dürre}

O getreuer, wahrhaftiger und gerechter Gott! der du in einem wahrhaften Worte auf das allertrostlichste hast angezeigt, wie daß ein recht ernstliches, eifriges, glaubreiches und beharrliches Gebet durch die Wolken des Himmels kräftiglich dringe und nicht allein gewiß vor dein Angesicht komme, sondern auch gar viel vermöge
\quellseite{324}
und ausrichte; wir Armen, deine Kinder, kommen hier auf dein tröstliches Anzeigen, in dieser ungewöhnlichen und über die Maßen lang währenden hitzigen und trocknen Zeit, mit unserm gläubigen, herzlichen und recht ernstlichen Gebete vor das Angesicht deiner hohen göttlichen Majestät und bitten dich demüthiglich, du wollest uns deiner Zusage nach gnädiglich erhören, deinen großen gefaßten Zorn von uns mildiglich abwenden, die tief ausgedorrte Erde wiederum nach Nothdurft befeuchten, und einen gebührlichen Regen vom Himmel herunter auf sie fallen, tropfen und regnen lassen, damit sie auch wiederum ihre befohlene Früchte zu nothwendigem Aufenthalt unsers zeitlichen Lebens, dazu du sie in der ersten Schöpfung väterlich verordnet hast, hervorbringen und geben könne. Wie du nun vor Zeiten den geschlossenen Himmel auf das ernstliche Gebet deines getreuen Knechts und Propheten Elia wiederum eröffnet und auf die durstige Erde hast regnen lassen; also versehen wir uns auch jetzt zu dir, du werdest in der großen dürren Hitze mit uns allergnädigst handeln und uns dieser ernstlichen Bitte zu deinem Lob und unserm Heil väterlich gewähren, durch Jesum Christum, deinen lieben Sohn, unsern einigen Erlöser und Seligmacher. Amen.

\section{377. Danksagung nach erlangtem fruchtbaren Regen}

O treuer Gott, ob wir zwar als arme Sünder deinen Zorn erwecket und nichts denn Strafen verdienet haben; wie du uns denn eine Zeit hero alle heimgesuchet; da die Sonne aufging mit der Hitze, daß das Gras verwelkte, die Blüthe abfiel und ihre schöne Gestalt verdarb; da du den Wolken gebotest, daß sie nicht darauf regneten, und plagest uns mit dürrer Zeit: so hast du doch nun den Himmel mit Wolken bedecket und gibst Regen auf Erden, lässest Gras wachsen auf den Bergen, gibst dem Vieh sein Futter und erquickest also das dürre Land wiederum. Wir danken dir, HErr, daß du so zornig bist gewesen über uns und dein Zorn sich gewendet hat, und tröstest uns mit deiner Hülfe. Du lässest dich doch nicht unbezeuget, hast uns viel Gutes gethan und vom Himmel wiederum Regen und fruchtbare Zeiten gegeben, daß du unser Herz erfüllest mit Speise und Freuden. — So jauchzet Gott, alle Lande, lobsinget
\quellseite{325}
zu Ehren seinem Namen, rühmet ihn herrlich, und sprechet zu Gott: Wie wunderlich sind deine Werke! Alle Lande beten dich an und lobsingen dir und deinem heiligen Namen! Hilf aber auch, du getreuer Gott, daß wir uns zu dir, dem Herrn, bekehren und unser sündliches Leben lassen, vom Bösen aufhören und anfangen Gutes zu thun und darinnen beständig verharren, damit wir dich nicht auf ein neues erzürnen und durch vorsätzliche Sünden entweder die vorigen oder noch größere Strafen auf uns laden. Regiere uns mit deinem Heiligen Geist, daß wir dir mit Freude und Lust unseres Herzens dienen, damit alle deine Segen weiter über uns kommen in der Stadt, auf dem Lande, auf dem Acker, im Hause und in Allem, was wir haben. Ja, daß wir gesegnet seyn, wenn wir aus- und eingehen, bis wir dermaleins den reichen geistlichen Segen in himmlischen Gütern durch Christum vollkommenlich empfangen; welchem sammt dir, o himmlischer Vater, und dem Heiligen Geiste, ewigem wahren Gott, sei Lob, Ehr und Preis in alle Ewigkeit. Amen.

\section[378. Gebet um einen gnädigen warmen …]{378. Gebet um einen gnädigen warmen Sonnenschein zur Zeit langwieriger Nässe}

Herr, allmächtiger Gott! der du den Himmel mit Dunkel bekleidest, und seine Decke als einen Sack machest; du hast die unbußfertige erste Welt um ihrer großen Bosheit und Frevels willen mit dem Wasser der Sündfluth in deinem Zorn gestrafet, und alles, was im Trocknen auf Erden einen lebendigen Odem, Geist und Wesen gehabt, verderbet! Wir werden auch zu dieser unserer Zeit um unserer Uebertretung willen, die wir allenthalben wider dich gethan haben, nach deiner strengen Gerechtigkeit mit ungewöhnlichem Regen, steter und langwieriger Nässe heimgesuchet, dadurch dem Lande und seinem Gewächse großer Schade geschieht. Barmherziger Vater! der du die Brunnen der Tiefe sammt den Fenstern des Himmels in der Sündfluth verstopfest, und dem Regen vom Himmel gewehrest; nimm von uns alle unsere Sünde und Missethat, und wende die Strafe und schädliche Regen und große Gewässer von uns. Denn das Land ist betrübt, und der Himmel droben traurig. Siehe den Zeugen in den Wolken, den
\quellseite{326}
Regenbogen an, und erinnere dich deiner Zusage, gib gutes und fruchtbares Wetter, damit die Früchte des Feldes und in Gärten ihr Gedeyen und Fortkommen haben, und durch die vielfältige und übrige Nässe nicht verderbet werden, und Theurung folge. HErr, du wahrhaftiger Gott! der du nicht lügen kannst, du hast gesagt: Es soll hinfort keine Sündfluth mehr kommen, die die Erde verdirbt! Darum sei gnädig, HErr! und thue deine Barmherzigkeit nicht vor Zorn verschließen, nimm deine Gnade und Barmherzigkeit nicht von uns. Gib uns nicht müßige Zähne und Mangel am Brode; segne unser Brod, daß wir zu essen haben, und deinen Namen preisen; so werden wir merken, daß wir deine rechten Kinder sind, für die du treulich sorgest, und dich als unsern himmlischen Vater ehren und preisen all unser Lebenlang. Amen.

\section{379. Gebet für die lieben Früchte}

Lieber HErr Gott, behüte gnädiglich die lieben Früchte auf dem Felde, reinige die Luft, gib seligen Regen und gut Gewitter, daß die Früchte wohl gerathen, daß sie auch nicht vergiftet werden und wir mit dem Vieh daran essen möchten die Pestilenz, Fieber und andere Krankheiten; denn solche Plagen kommen daher, daß die bösen Geister die Luft vergiften und darnach die Früchte, Wein und Korn, daß wir durch dein Verhängniß den Tod und Plagen essen und trinken müssen an unsern eigenen Gütern. Darum laß sie, lieber Gott, gesegnet sein, daß sie uns gesund und selig gedeihen mögen, wir auch derselben nicht mißbrauchen zum Schaden der Seelen oder Vermehrung der Sünden, Völlerei und Müßiggangs, daraus Unkeuschheit, Ehebruch, Fluchen, Schwören, Morden, Krieg und alles Unglück folget; sondern gib uns Gnade, deiner Gaben zu brauchen zur Seelen Seligkeit und Besserung unsers Lebens, und daß die Früchte eine Ursache seien, des Leibes und der Seelen Gesundheit zu behalten und zu mehren. Amen.

\section[380. Gebet, wenn durch Gottes Verhängniß den …]{380. Gebet, wenn durch Gottes Verhängniß den Feldfrüchten vom Mehlthau und allerlei Ungeziefer Schaden geschieht}

Wahrhaftig und groß sind alle deine Gerichte, wunderbar und unerforschlich sind alle deine Werke, du gewaltiger und allmächtiger
\quellseite{327}
O Herr Himmels und der Erde. Denn wie du verlangst in deinem heiligen Wort über Alle, die sich deinen Geboten muthwillig widersetzen, erschreckliche Strafen und Plagen gedroht und unter andern dich auch hast vernehmen lassen, wider sie allerhand Ungeziefer zu wappnen und zu führen; also erfahren und sehen wir, als die Unbußfertigen, auf diesen heutigen Tag mit unserm und unserer Nahrung äußersten Schaden und Nachtheil, daß solche gemeldeten Strafen und Plagen um unsrer Sünde willen, uns von unserm gottlosen Wesen abzuschrecken, auch über uns gekommen sind. Denn was wir jährlich mit den Unsrigen zur Erhaltung unseres zeitlichen Lebens einsammeln, einernten, einschneiden und folgend genießen und gebrauchen sollen, das fressen, verzehren und verderben vor unsern Augen die schädlichen wilden Thiere, die unzähligen Raupen, die scheußlichen Heuschrecken und Käfer, ja der beschwerliche Mehlthau und andere beschwerliche Würmer und Ungeziefer. Wir bitten dich, o Gott, in diesem gegenwärtigen Jammer und Elend durch deinen geliebten Sohn JEsum Christum, strafe uns nicht in deinem grimmigen Zorn, wie du etwa die verruchten Aegyptier gleichfalls gestraft hast, sondern nimm gegenwärtige Beschwerniß von uns gnädiglich hinweg; beschütze und beschirme Alles, was wir zu nothwendiger Unterhaltung des Leibes für uns und die Unsrigen täglich bedürfen und brauchen. Insbesondere aber, lieber Vater, behüte uns vor dem höllischen Ungeziefer und vor dem bösen, unsterblichen Wurm eines bösen, unruhigen, für und für nagenden und beißenden Gewissens, damit wir also durch dich hier zeitlich und dort ewiglich in Gnaden erhalten werden, der du mit Gott, deinem eingebornen Sohne, und Gott, dem Heiligen Geiste, wahrer Gott, hochgelobet, lebest und regierest in alle Ewigkeit. Amen.

\section{381. Gebet bei Ungewitter}

Ach HErr Gott vom Himmel! großer und schrecklicher Gott! der du herrschest in allen Königreichen auf Erden, und mächtig, schrecklich, auch wunderthätig bist, deine Stimme gehet auf den großen Wassern; du Gott der Ehren donnerst; deine Stimme gehet mit Macht; deine Stimme, HErr, gehet herrlich, deren Grimm ist ein verzehrend Feuer; und es ist schrecklich, in deine Zornhände
\quellseite{328}
fallen. Wenn du ein Land anrührest, so zerschmelzet es, daß alle Einwohner trauern müssen. Unsere Missethaten haben es verdient, daß du uns strafest; hilf uns aber in dieser Noth um deines Namens willen. Laß deinen Zorn nicht aufwachsen über deine brünstige und herzliche Barmherzigkeit. Du bist ja unser Vater, von Alters her ist das dein Name. Du bist unsere Hoffnung; und hast uns versprochen, wenn wir uns würden zu dir kehren, daß dich des Unglückes gereuen soll, das du gedachtst uns zu thun. So vertilge denn unsere Missethat wie eine Wolke, und unsere Sünde wie den Nebel. Ach Herr! laß uns nicht und ziehe deine Hand nicht von uns ab, Gott unser Heil! nimm die Zahlung für unsere Sünde von dem blutigen Verdienste deines Sohnes Jesu, damit er uns erworben, und eine ewige Erlösung erfunden hat, und daß uns reiniget von allen unsern Sünden. Er ist ja die Versöhnung für unsere Sünde, und durch ihn haben wir Friede mit dir. Kehre dich jetzt von dem Grimm deines Zornes, und sei gnädig der Missethat deines Volkes, nach deiner großen Barmherzigkeit. Verderbe nicht dein Volk und Erbtheil, das du durch deine große Kraft erlöset hast, und durch deines lieben Sohnes eigenes Blut erkauft hast; behüte unsern Leib und Leben wie deinen Augapfel. Bewahre auch in diesem großen Ungewitter unser Haus, Hof, Güter, Nahrung, die dein Segen uns bescheret hat. Wende von uns väterlich ab alles wohl verdiente Uebel, und hilf uns, Gott unseres Heils, um deines großen Namens Ehre, um aller deiner Barmherzigkeit wegen, und um des Blutes Jesu, unseres Heilandes willen. Amen.

\section{382. Gebet, wenn das Gewitter lange anhält}

Allmächtiger, ewiger, gewaltiger und schrecklicher Gott, wir sehen aus gegenwärtigem Wetter, wie heftig du über uns erzürnet bist, daß du auch deinem Firmamente schon befohlen hast, sich wider uns zu setzen und um unserer Sünde willen uns zu strafen, weil ein grausamer Blitz und Donnerschlag auf den andern folget, also daß wir in Gefahr Leibes und Lebens stehen, wo du uns nicht väterlich beschützest. Darum, o Herr, sei nicht ferne von uns, denn Angst ist nahe, und ist sonst hier kein Helfer. Ach, lieber Herr
\quellseite{329}
und Gott! verschone unser, behüte uns und die lieben Früchte auf dem Felde, damit wir unsern leiblichen Unterhalt haben. Sei uns gnädig und erbarme dich unser, verleihe uns allezeit hier zeitlich deine Gnade und Barmherzigkeit, und nach diesem Leben die ewige, immerwährende Seligkeit. Amen.

\section{383. Gebet bei schrecklichem Sturm}

Ach Herr, Herr, du großer und erschrecklicher Gott, wir, deine elenden Kinder, hören mit Furcht und Zittern, wie du einen schrecklichen Sturmwind, Brausen und Toben erweckest, und der gottlosen, bösen Welt deinen gerechten Zorn und wohlverdiente Strafe dadurch anzeigen lässest. Ach Herr, stehe nicht an unsre Sünde groß, sprich uns derselben in Gnaden los, mach uns von allen Plagen frei, steh uns in unserm Elend bei, bedrohe den großen Sturmwind, damit sich lege sein schreckliches Reisen und Wüthen, steure dem Satan, daß er keine Macht an unserm Leib und Leben, Haus und Hof, Hab und Gut, finden möge, wende ab gefährliche Feuersbrunst, und allen Schaden an Häusern und Gebäuden, an Menschen und Vieh, zu Wasser und zu Lande. Laß uns, lieber Vater, in deine allmächtige Hand eingeschlossen sein und bleiben, so wird uns weder dieser Wind noch einige andre Creatur im Himmel und auf Erden schaden, oder von dir reißen können in Ewigkeit. Hilf uns, Herr Gott, so ist uns geholfen, auf daß von Herzen können wir nachmals mit Freuden danken dir, gehorsam sein nach deinem Wort, dich allzeit preisen hier und dort, um deines lieben Sohnes JEsu Christi, unsers Erlösers und Seligmachers, willen. Amen.

\section{384. Gebet bei Ungewitter}

Großmächtiger und starker Gott im Himmel und auf Erden! aus deinem jetzigen Donner und Blitzen werden wir erinnert, daß du über uns sündige Menschen sehr zornig bist. Du rufest uns durch solchen Schall, vor welchem alle Creaturen zittern und erschrecken müssen, als durch die rechte Bußglocke, zur Erkenntniß unserer Sünden und zur wahren Buße, damit wir nicht mit der gottlosen bösen Welt mögen verdammet werden. Du bist ein so mächtiger und gewaltiger Herr, daß du auch, wo du nach deiner strengen
\quellseite{330}
Gerechtigkeit verfahren wolltest, mit einem einzigen Donnerschlage und Blitze die ganze Welt im Augenblicke verderben könntest. Aber wir fallen auf die Knie unseres Herzens und bitten demüthiglich, laß ab von deinem grimmigen Zorne, der wider die Sünde wie ein Feuer brennet. Erbarme dich unser, verschone unser, erhebe über uns dein gnädiges Angesicht und sei uns barmherzig. Bewahre unser Leib und Leben, Haus und Hof, Stadt und Land, unser Hab und Gut und Alles, was wir haben, auch die Früchte auf dem Felde vor Hagel, Schlossen, Anzündung des Wetters, vor großen Wasserfluthen und vor allem Schaden. Behüte uns auch vor einem bösen, schnellen Tod und verleihe gnädiglich, daß viel Sünder durch dein Wort, Donner und Blitzen vom gottlosen Wesen und Leben abgeschrecket, sich von Herzen zu dir wenden und bekehren mögen. Laß uns alle Stunden und Augenblicke an deine herrliche Zukunft zum Gericht gedenken, und bereit sein, dich mit Freuden zu empfahen und mit dir zur ewigen Freude einzugehen. Amen. Ach HErr! sei uns armen Sündern an Leib und Seel, hier in diesem und dort im zukünftigen Leben, gnädig und barmherzig. Amen.

\section[385. Danksagung, wenn das Gewitter glücklich …]{385. Danksagung, wenn das Gewitter glücklich vorüber gegangen ist}

Allmächtiger Gott, himmlischer Vater, der du uns gesagt hast: Rufe mich an in der Zeit der Noth, so will ich dich erretten, so sollst du mich preisen: wir sagen dir von Grund unsers Herzens Lob und Dank, daß du unser Gebet gnädig erhöret und das zornige Wetter also hast vergehen lassen, daß uns dadurch an Leib und Gut kein Schaden geschehen ist; damit du abermals dein getreues Vaterherz zu erkennen gegeben hast, daß du nicht mit uns handeln wollest nach unsern Sünden, und uns nicht vergelten nach unsern Missethaten. Verleihe uns, barmherziger Vater, um deines eingebornen Sohnes JEsu Christi willen, daß wir uns ob solchen deinen ernstlichen Warnungen beständiglich bessern, in deiner Furcht hinfür leben, und auf die Zukunft deines lieben Sohnes uns bereit machen, damit wir ihm mit Freuden entgegen gehen, den neuen Himmel, in welchem Gerechtigkeit wohnet, mit Jauchzen
\quellseite{331}
und Freuden einnehmen, und mit dir in deiner seligen Anschauung ewiglich besitzen mögen, durch denselben deinen lieben Sohn, Jesum Christum, unsern Herrn. Amen.

\section[386. Gebet, wenn ein Comet oder andere …]{386. Gebet, wenn ein Comet oder andere wunderliche Zeichen am Himmel sichtbar werden}

Großer Gott, du thust Wunder oben im Himmel und unten auf Erden, du zählest die Sterne und nennest sie alle mit Namen, und thust große Dinge, die nicht zu forschen sind. O verborgener Gott, der du das Auge gemacht hast, öffne uns die Augen, daß wir das gegenwärtige Wunderzeichen am Himmel mit Verwunderung ansehen, und bedenken, daß es nicht ohngefähr, sondern durch deine göttliche Vorsehung geschehe, die Gottlosen zu schrecken und die Frommen zu trösten. O Herr Jesu Christe, hierdurch werden wir deiner Zukunft erinnert und unserer Erlösung vertröstet, daß der jüngste Tag vor der Thür sei. O Gott heiliger Geist, weil die Zeichen des Himmels unsere Bußprediger sein, so hilf, daß wir uns unserer Sünden dadurch erinnern und nach solchen Zeichen uns bessern, durch Jesum Christum. Amen.

\chapter{XI. Kreuz- und Nothgebete}

\section{387. Gebet in großer Schwachheit des Glaubens}

Ach Herr, ich erfahre es nun mit der Wahrheit, daß der Glaube nicht Jedermanns Ding sei. Ich glaube, lieber Herr, aber hilf meinem Unglauben! Du wollest das zerstoßene Rohr nicht zerbrechen, und das glimmende Döchtlein nicht auslöschen, o Jesu, der du sitzest zur Rechten Gottes, vertritt und bitte für mich, daß mein Glaube nicht aufhöre. Sei der Anfänger und Vollender des Glaubens, damit ich auslösche alle feurigen Pfeile des Bösewichts. Laß mich glauben, ob ich gleich nicht sehe, und also selig sein. Amen.

\quellseite{332}

\section{388. Klage wegen Mangels an Andacht im Gebet}

Mein Gott und Herr, ich klage dir mit Seufzen, daß ich dich gern mit inbrünstiger Andacht anrufen und das Opfer eines eifrigen Gebets bringen wollte, aber mein Herz ist zu schwer und zu kalt, es will sich nicht recht und feurig von dem Irdischen zu dem Himmlischen erheben. Ich befinde mich ungeschickt und unwürdig zum Gebet und muß oft mein Fleisch und Blut zur Andacht mit aller Macht zwingen, darüber ich mich betrübe und mir einbilde, du seist mir ungnädig geworden, weil ich keine rechte Bewegung des Heiligen Geistes, mit dir zu reden, empfinde. Lieber Vater, habe ich solches mit unbedächtigem Wandel und Nachlässigkeit verschuldet, habe ich dich nicht in deinem Wort, wenn du mich zur Buße gerufen, fleißig gehöret, habe ich deines Heiligen Geistes Trieb nicht gehorsam gefolget, habe ich etwa meine Gedanken zu sehr auf das Irdische gewendet, so vergib mir, um deines lieben Sohnes Jesu Christi willen. Reiße von mir alle Hindernisse, die mich von seliger Andacht abhalten, und nimm mit meiner Schwachheit, die mir sehr leid ist, väterlich vorlieb; nimm alle geistliche Unlust aus meinem Herzen und laß mich willig und bereit sein, dir mit andächtigem und von der Welt abgewandtem Gemüth zu dienen. Ich sehne mich nach deiner himmlischen Gnade, nimm dich meiner Seele an und laß mich ein kindliches Gespräch mit dir, als meinem himmlischen Vater, erfreulich halten, daß der Heilige Geist in mir täglich zu dir rufe: Abba, lieber Vater! Hierüber will ich mich von Herzen freuen und mich dessen trösten, daß dir mein armes Gebet, als ein williges Opfer, angenehm und gefällig sei und ich bei dir ewig in Gnaden sein und bleiben werde. Amen.

\section{389. Gebet in Angst und Bangigkeit des Herzens}

Du Gott alles Trostes, der du reichlich erquickest Alle, die in ihrer Angst und Trauern ihre kindliche Zuflucht zu dir haben, ich armer, hochgeängsteter Mensch weiß nirgends hinzufliehen, als zu deiner Barmherzigkeit; höre doch auf mein Geschrei und merke auf mein Gebet! aus Betrübniß rufe ich zu dir. Ich ängstige mich täglich in meinem Herzen, Angst der Hölle hat mich getroffen,
\quellseite{333}
und um Trost ist mir sehr bange; Furcht und Zittern ist mich angekommen. HErr, du lässest mich erfahren viel und große Angst, so mache mich doch auch wieder lebendig und führe mich wieder aus der tiefen Erde herauf. Du hast mich in eine harte Angstschule nach deinem väterlichen Willen gerathen lassen, HErr; Angst ist nahe, und ist hier kein Helfer, ohne du allein. Hilf mir diesen hohen Angstberg durch deine Gnade überwinden und gib mir für Angst Jauchzen, für Bangigkeit Fröhlichkeit, für Zittern Herz und Muth, für Zagen Trost und Leben, daß ich der vorigen Angst vergesse und in deiner göttlichen Erlösung mich ergötze, und meine Gebeine sich freuen, die jetzt sehr erschrocken sind. HErr, ich weiß, daß du deine lieben Kinder nicht von Herzen betrübest, sondern nur zu ihrer Seligkeit und ihrem Besten dein freundliches Angesicht eine Zeit lang verbirgest; du hast mir ja vorhin und andern deinen Gläubigen so oft aus Angst und Noth ganz wunderbar geholfen. Beweise auch jetzt deine Gnade an mir, ich will alle Stunden auf deine Hülfe hoffen; ich weiß, du wirst mich nicht ohne Trost lassen, weil ich sein so herzlich begehre. Wende dich zu mir und sei mir gnädig; denn ich bin einsam und elend, die Angst meines Herzens ist groß, führe mich aus meinen Nöthen. Siehe an meinen Jammer und Elend, und vergib mir alle meine Sünde um deiner Barmherzigkeit willen. Amen.

\section{390. Gebet in großer Sündenangst}

Ach HErr, du gnädiger und barmherziger Gott, strafe mich nicht in deinem Zorn, und züchtige mich nicht in deinem Grimm; denn deine Pfeile stecken in mir und deine Hand drücket mich; es ist nichts Gesundes an meinem Leibe vor deinem Drohen, und ist kein Friede in meinen Gebeinen vor meiner Sünde, mein Herz bebet, meine Kraft hat mich verlassen, und das Licht meiner Augen ist nicht bei mir; es hat mich Leiden umgeben ohne Zahl, es haben mich meine Sünden ergriffen, daß ich nicht sehen kann, ihr ist mehr denn Haar auf meinem Haupte, und mein Herz hat mich verlassen. O gütiger Vater, weil ich dich mit meinen vielfältigen Sünden so hoch erzürnet habe, leide ich daher in meinem Gewissen sonderlich die schwere und schreckliche Anfechtung, daß ich nicht an deiner
\quellseite{334}
Gnade zweifeln, sondern vielmehr verzweifeln soll, als ob sie für mich ganz und gar aus sei, und ein Ende hätte. O Gott, wie wird mein Herz geängstiget und gequälet! Ach stärke und erhalte mich in diesen meinen großen Aengsten und zagen durch deinen Heiligen Geist! O Herr, errette mich, eile mir zu helfen, führe meine Seele aus der Höllenangst, vergib mir alle meine Sünde, aus welcher der Teufel eitel hohe, große Berge macht, hinter welchen ich deine Gnade nicht sehen soll. Laß mir doch einen Anblick werden deiner Barmherzigkeit, und beraube mich nicht deiner göttlichen Hülfe. Tröste mich wieder, mein Gott, auf daß mein Herz gestillet und fröhlich werde. Du bist ja größer denn unser Gewissen, welches, ob es uns gleich anklaget, kannst du es doch durch dein Evangelium und fröhliche Botschaft wohl schweigen. Und obwohl meine Sünden mächtig und blutroth sind, so ist doch deine Gnade noch viel mächtiger und kannst sie schneeweiß machen. So sei nun wieder zufrieden, meine Seele, was betrübest du dich, und bist so unruhig in mir? Harre auf Gott, der meines Angesichtes Hülfe und mein Gott ist. Er wird dich aus der Angst reißen, denn der Herr thut dir Gutes, und errettet dich aus allen deinen Sünden durch den unschuldigen Tod Jesu Christi, welcher das heilige, unbefleckte Opfer für unsere Missethat geworden ist. So betrübe dich doch nicht also, meine Seele, und sei nicht unruhig in mir. Harre nur auf Gott, denn ich werde ihm noch danken, daß er meines Angesichtes Hülfe und mein Gott ist. Dieses Alles, mein Gott, schreibe in mein Herz, und stärke damit meinen Glauben, daß ich in aller Anfechtung ritterlich ringe, und durch deine kräftige Gnade Alles wohl ausrichten und das Feld behalten möge. Amen.

\section{391. Gebet nach einem tiefen Sündenfall}

Ach mein herzallerliebster himmlischer Vater, o du gerechter Richter der Welt, wie groß habe ich armer Mensch mich an dir versündiget! Ach, was habe ich gethan, wie habe ich deiner heiligen Gebote, deiner Gutthaten gegen mich, deiner schweren Drohungen wider die Sünde, ja deines offenbaren schrecklichen Gerichts so schändlich, so gar lüderlich vergessen, daß ich diese große Sünde,
\quellseite{335}
die mich nun in meinem Herzen Tag und Nacht ängstet, begangen habe! Ach Gott, wie bange ist meiner Seele, wie bebet mein Herz, wie muß ich mich nun schämen in meinem eigenen Gewissen vor dir! Welch einen tiefen Abgrund der Sünden und der Verdammniß sehe ich vor meinen Augen! Nun, HErr, mein Gott, ach lieber Vater, in dieser meiner Angst nehme ich meine Zuflucht in kindlichem Vertrauen zu deiner grundlosen Barmherzigkeit und bitte dich durch JEsum Christum um Gnade. Ach, wie sind meiner Sünden so viel, wie ist meine Schuld so sehr groß und ich kann in Ewigkeit nimmermehr bezahlen! Habe doch Geduld mit mir armen Sünder, du lieber Gott, und tilge meine Missethat nach deiner übergroßen Barmherzigkeit um Christi willen, der auch für mich bezahlet hat. Ich habe ja mißgehandelt und sehr grob gesündiget. Es ist mir aber ja doch von Herzen leid, und kränke mich deswegen täglich, wie du, o Herzenserkündiger, an mir siehst. Ei nun, lieber Gott, ein solch geängstigt und zerschlagen Herz wirst du ja nicht verstoßen; du bist ja nahe denen, die zerschlagene Gemüther haben, und hilfst denen, die zerbrochenen Herzens sind. Ach, so verstoße doch auch mich nicht! Heile, HErr, mein zerschlagenes Herz, mein verwundetes Gewissen. Hilf mir armen Sünder um Christi willen und verbinde meine Schmerzen. Du bist ja mein Vater, mein Schöpfer, und weißt, was für ein Gemächt ich bin, und erkennst, daß ich Staub bin. Dein Sohn JEsus Christus ist ja auch in die Welt gekommen, die armen Sünder selig zu machen, und sein Blut reiniget uns von allen Sünden. Ach Gott, so laß doch das theuere Verdienst JEsu Christi an mir armen, betrübten, bußfertigen Sünder nicht verloren sein. Du hast ja verheißen, du wollest nicht den Tod des Sünders, sondern daß er sich bekehre und lebe. Nun, lieber Gott, ich begehre von Herzen Buße zu thun durch deine Gnade. Du erkennest mein Herz und weißt, daß mir meine begangene Sünde recht herzlich leid ist, und begehre mich zu bessern. O schenke mir das Leben und bekehre mich, HErr, so werde ich bekehret. Gehe doch nicht mit mir ins Gericht und laß mich nicht verzagen in meinen Anfechtungen. Ach Gott, wie manchmal bin ich so betrübt wegen meiner Sünden, daß ich nicht weiß wohin, und was ich thun soll. Ach HErr, verwirf mich doch nicht von
\quellseite{336}
deinem Angesicht und nimm deinen Heiligen Geist nicht von mir. Reuch meinen Fuß aus dem Netze. Verleih mir doch meine Missethat und laß mich durch die Kraft des Heiligen Geistes empfinden Freude und Wonne, daß die Gebeine fröhlich werden, die du zerschlagen hast. Ich finde nun, daß, so leicht der Satan vorhin mir die Sünde gemacht hat, mich darein zu stürzen, so grausam und schwer macht er sie nun in meinem Gewissen, mich zu der Verzweiflung zu bringen. Ach mein Gott, verlaß mich nicht in solchem harten Kampf, sondern laß deine Kraft in meiner Schwachheit mächtig sein. Wappne du mich durch die Kraft deiner Stärke, damit ich bestehen möge wider die listigen Anläufe des Teufels. Erleuchte auch meine Augen, daß ich hinfüro fürsichtig wandle und vor dergleichen Sünden mich hüte. Ach Gott, wo du mein Herz erfreuest mit deiner Gnade, wie herzlich will ich dir danken, wie fleißig will ich durch deine Gnade mein Leben bessern, wie gern will ich andere Uebertreter deine Wege lehren, daß sich die Sünder zu dir bekehren. Es soll mir, o lieber Gott, eine Warnung sein, und ich werde mich hüten vor solcher Plage des Gewissens mein Lebenlang. Ich hoffe darauf, mein Gott, daß du so gnädig bist, und mein Herz freuet sich, daß du so gern hilfst und keinen hinausstößest, der zu dir kommt. Darum komme ich nun und bitte dich nochmals, sei mir armen Sünder gnädig um deines allerliebsten Sohnes, meines einigen Heilandes und Fürsprechers JEsu Christi willen, welcher verheißen hat, was wir dich in seinem Namen bitten werden, das werdest du uns geben, und hat uns darauf befohlen, also zu beten: Vater unser, der du bist im Himmel, geheiliget werde dein Name. Dein Reich komme. Dein Wille geschehe auf Erden, wie im Himmel. Unser täglich Brod gib uns heute. Und vergib uns unsre Schuld, als wir vergeben unsern Schuldigern. Und führe uns nicht in Versuchung, sondern erlöse uns von dem Uebel. Denn dein ist das Reich und die Kraft und die Herrlichkeit in Ewigkeit. Amen.

\section{392. Gebet in Anfechtung zur Verzweiflung}

Barmherziger Gott, nachdem ich oft muthwillig deine heiligen Gebote übertreten, dich, meinen Gott und Herrn, verachtet, erzürnt
\quellseite{337}
und heftig beleidiget habe, daher mir mein Gewissen hart beschweret und verwundet ist, daß ich darob fast kleinmüthig und zaghaftig bin. Wiewohl mir dein heiliges Wort Vergebung meiner Sünden aus lauter Gnade unwiderruflich zusagt, so ist doch mein Glaube schwach und der Teufel stark, der mir gern allen Trost stehlen und aus dem Herzen reißen möchte. Darum rufe ich zu dir, o Heiliger Vater, laß mich an deiner göttlichen Gnade nun und nimmermehr verzagen, daß ich nicht in die allerggrößte Sünde des Unglaubens und der Verzweiflung falle oder darein willige. Stärke mich, daß ich mitten im Tode auf dich, mein Leben, hoffe und an deiner Barmherzigkeit und Hülfe nicht verzage, auf daß ich nicht wie der gottlose Cain meine Sünden größer achte, denn daß sie mir könnten vergeben werden. O Christe, Gottes Sohn, du lebendiger Brunn aller Gnaden, der du überfließest mit eitel Quellen der Barmherzigkeit, zu dir rufe ich von ganzem Gemüthe, mehre mir den Glauben an dein heiliges Leiden und Sterben; denn das ist ja gewißlich und unwidersprechlich wahr, daß ein einiges Tröpflein deines allerheiligsten Blutes, für mich vergossen, viel kräftiger und mächtiger ist, als alle meine größten und mächtigsten Sünden. Siehe mich an mit den Augen deiner Barmherzigkeit, wie du angesehen hast den lieben Petrus, nachdem er dich verleugnet und sich verflucht hatte, auf daß ich nicht wie Judas, der Verräther, verzweifele und die Sünde wider den Heiligen Geist begehe. O Gott Heiliger Geist, du milder Schatz, stehe mir bei in meiner letzten Noth, wenn mich der böse Feind anklagt und mein Gewissen mich beschuldiget, wenn mich erschreckt der Hölle Anblick und ich mit eitel Todesnöthen und greulichen Anfechtungen umfangen bin. Wenn mich die ganze Welt verläßt und alles wider mich steht, so tröste mich, daß meine Hoffnung mir nicht entfalle. Bekräftige mein Herz mit deinem Zeugniß und deiner Versiegelung, daß ich fest glaube eine Vergebung der Sünden, die mir und allen, so der Verheißung Gottes trauen, widerfahren wird. Laß mich des Bundes meiner heiligen Taufe eingedenk sein und mich dabei der Zusage: Wer glaubet und getauft wird, der wird selig werden, von Herzensgrunde annehmen und trösten. Amen.

\quellseite{338}

\section{393. Zuflucht zu Christo in Sünden noth}

Lieber Herr Jesu Christe, ich fühle meine Sünden, sie beißen, jagen und erschrecken mich, wo soll ich hin? Ich sehe dich, Herr Jesu Christe, an und glaube, wiewohl schwächlich, an dich; doch halte ich mich an dich und bin gewiß, du hast gesprochen: Wer an mich gläubet, soll haben das ewige Leben. Ob nun gleich mein Gewissen beschweret ist, und die Sünden mich erschrecken und das Herze zittern machen, so heißt es doch: Sei getrost, mein Sohn, deine Sünden sind dir vergeben und du sollst das ewige Leben haben, und ich will dich auferwecken am jüngsten Tage. Amen.

\section{394. Ein anderes Gebet}

Ich glaube an Christum, der von der Jungfrau Maria geboren ist, gelitten und gestorben, und verlasse mich darauf, daß er selber sagt: wer zu ihm kommt, den will er nicht ausstoßen. Auf diese Worte verlasse ich mich und komme darauf zu dir, lieber Herr Christe, denn das ist dein Wille und Herz, auch dein Mund. Deine Worte sind mir genug und gewiß, ich weiß wohl, daß du mir nicht lügst, die Worte werden mir nicht fehlen: du willst die nicht verstoßen, die zu dir kommen. Ob ich schon ein Bube bin und nicht genug heilig und fromm, daß ich bestehen könnte, so bist du dennoch wahrhaftig und wirst, daß ich am jüngsten Tage soll auferwecket werden. Ob ich nun nicht kann bestehen, so willst du doch, lieber Herr Jesu Christe, wohl stehen und mich nicht verwerfen. Amen.

\section[395. In großer Melancholie und Traurigkeit …]{395. In großer Melancholie und Traurigkeit des Gemüthes}

Allmächtiger, gerechter, barmherziger Vater im Himmel, du weißt, daß wir armen Menschen durch die Erbsünde also sind verderbet, daß wir nicht allein mit vielen Sünden deinen gerechten Zorn verdienen, denselben auch in unsern Herzen bisweilen schmerzlich fühlen, sondern auch mit vieler trübseliger Zeit und Widerwärtigkeit beleget werden, und oft in große Melancholie des Gemüths gerathen, welcher wir uns aus menschlicher Schwachheit nicht entschlagen können ohne deinen heiligen göttlichen Trost, welcher die betrübten Herzen zur Ruhe stellen, sie erfreuen und zum ewigen
\quellseite{339}
Leben erhalten kann. Darum, o lieber Vater, siehe mich auch mit Gnaden an, der ich jetzt von Herzen betrübt bin, und fast den ganzen Tag traurig gehe. Herr, meine Seele ist betrübt bis in den Tod, daß ich auch vor großer Traurigkeit und Trägheit meines Herzens fast zu allen Werken meines Berufes untüchtig werde und darüber vergesse mein Brod zu essen. Ach laß mich, o treuer Gott, in solcher Melancholie und Betrübniß meiner Seelen nicht verharren, an welcher der leidige Teufel seine Lust und Gefallen hat, sondern wende von mir die Traurigkeit, in welcher ich mich so sehr ängstige. Ach HErr, verwirf mich nicht von deinem Angesicht, verbirg dein Antlitz nicht von mir. Wie lange soll ich sorgen in meiner Seele und mich ängsten in meinem Herzen? Behüte mich vor Anfechtung und Ungeduld, daß ich aus meiner Noth errettet, dich ganz fröhlich rühmen könne. Du willst ja den glimmenden Docht nicht auslöschen und das zerstoßene Rohr nicht gar zerbrechen; so weiß ich auch, daß, wer sich von Herzen zu dir hält, von dir nimmermehr verstoßen wird. Und ob ich zwar in diesem meinen Elende nicht eine Stunde oder einen Tag mich meiner Traurigkeit entschlagen kann, so will ich dir doch, o lieber Vater, geduldig stille halten und mich tröstlich erinnern deiner Hilfe, die du mir vorhin oft väterlich geleistet, und mich aus mancher betrübten Stunde gnädiglich erlöset hast; warte derwegen alle Augenblick auf deine göttliche Erquickung des Heiligen Geistes und Aenderung meines Gemüthes. Darum sei dieses meine Freude, daß ich mich zu Gott halte in dieser meiner Traurigkeit; ich weiß, daß ich Gott zum Trost habe, und wenn ich nur, HErr, dich habe, so frage ich nichts nach Himmel und Erde. Wie der Hirsch schreiet nach frischem Wasser, so schreiet meine Seele, Gott, zu dir; meine Seele dürstet nach Gott, nach dem lebendigen Gott. Ich will mich freuen und fröhlich sein über deine Güte, daß du mein Elend ansiehest, und erkennest meine Seele in der Noth. Ich weiß ja, daß dem Gerechten das Licht immer muß wieder aufgehen und Freude den frommen Herzen, und daß ich um deines lieben Sohnes JEsu Christi willen gewiß ein Erbe des ewigen Lebens bin und bleiben werde, ob ich gleich noch so vielen Elende in dieser Welt unterworfen wäre, wenn ich gleich noch mit so viel traurigen
\quellseite{340}
Stunden geplaget würde. Derohalben hoffe ich auch, du werdest nach dieser meiner großen Betrübniß mein trauriges Herz durch deinen Heiligen Geist wieder erquicken, daß ich mich deiner Gnade von Herzen freuen möge. Und dafür will ich dem HErrn singen mein Lebenlang, und meinen Gott loben, so lang ich hier bin, denn du lässest mich erfahren viel und große Angst, und machest mich wiederum lebendig, dein Lob soll immerdar in meinem Munde sein, meine Lippen und meine Seele, die du erlöset hast, sollen fröhlich sein und dir Lob singen. Warum betrübest du dich, meine Seele, und bist so unruhig in mir? Harre auf Gott, denn ich werde ihm noch danken, daß er meines Angesichts Hülfe und mein Gott ist. Ich hoffe darauf, daß du, o Gott, so gnädig bist, und mein Herz freuet sich, daß du so gerne hilfest, ich will dem HErrn singen, daß er so wohl an mir thut. Amen.

\section[396. Gebet, wenn einer wegen der ewigen …]{396. Gebet, wenn einer wegen der ewigen Versehung angefochten wird}

Ewiger, gnädiger und gütiger Gott, der du in deinem Sohn JEsu Christo, ehe noch der Welt Grund geleget war, uns zum ewigen Leben erwählet und geschaffen, auch nach dem Fall unserer ersten Eltern denselben deinen lieben Sohn ohne einigen Unterschied allen Menschen zu einem Heiland gesandt hast, solchen deinen gnädigen Willen durch dein heiliges Wort der Welt offenbaren und durch die heiligen Sacramente, als gewisse und unfehlbare Gründe und Zeichen, bekräftigen lässest: ich armer betrübter Mensch klage dir, daß der listige und böse Geist mich mit diesen geschwinden und gefährlichen Anfechtungen und Gedanken quälet und plaget, daß ich nicht in die Zahl der Auserwählten, sondern in die Zahl der Verdammten gehöre, und nimmermehr könne selig werden, mich dadurch von deinem geoffenbarten Worte abzureißen, dich, wie andere Gottlose, zu lästern und zu schmähen, deinen Sohn, den du der ganzen Welt zu einem Heiland gegeben hast, nicht für meinen Heiland zu erkennen, sondern sein heiliges theures Blut, das er für mich vergossen hat, mit Füßen zu treten und mich in ewige Verzweiflung zu stürzen. Ach gütiger, barmherziger Gott, laß mich nicht unter diesen gefährlichen und feurigen Pfeilen des Bösewichts stecken,
\quellseite{341}
daß ich nicht darin ewiglich sterbe und verderbe, sondern sende den Tröster, den Heiligen Geist, zu mir, daß ich durch seine Hülfe in einfältigem Glauben mich an dein geoffenbartes Wort halte, deinen Sohn, den du, wie der ganzen Welt, also auch mir, zu einem Heiland gegeben hast, erkenne, sein nimmermehr vergesse, sondern stets mich erinnere, tröste und freue alles dessen, was er um der ganzen Welt willen gethan, auch zu meiner Seligkeit vollbracht hat. Gib, daß ich meine Taufe, Absolution und Abendmahl deines lieben Sohnes, als einen bewährten Schild, allen giftigen Pfeilen des leidigen Satans möge vorwerfen, und in großer Freudigkeit pochen, darauf trogen und mich gänzlich darauf verlassen, damit ich aus solchen feurigen Pfeilen des Teufels errettet, dich mit andern Auserwählten, die in gleichen Anfechtungen gestecket, in Christo JEsu, meinem Heilande, durch Kraft des Heiligen Geistes hier und dort mit fröhlichem Herzen und Munde loben und preisen möge, Amen, und abermals Amen in Christo JEsu mit dem Vater und dem Heiligen Geist hochgelobet. Amen.

\section{397. Gebet in satanischen Anfechtungen}

Ach HErr, wie sind meiner geistlichen Feinde so viel, die sich wider mich setzen! Sie sagen zu meiner Seele, sie habe keine Hülfe bei dir; errette du mich, so will ich mich nicht fürchten vor viel hundert tausend bösen Geistern, die meine Seele ängsten wollen. Auf dich, HErr, traue ich; laß mich nicht zu Schanden werden, und errette mich, daß der Teufel, mein Verfolger, zurückkehren müsse, und nicht, wie ein Löwe, meine Seele erhasche und zerreiße, weil kein Erretter da ist. HErr Christe, erhebe dich über den Grimm meines Feindes, der mir eine Grube gegraben, darinnen ich verderben soll; siehe an mein Elend, und erhebe mich aus den Thoren des Todes, daß ich deinen Preis erzähle, daß die alte Schlange in die Hölle gefehret werde. HErr, verbirg dich nicht vor mir, denn der böse Geist treibet Uebermuth über mich; er lauert auf mich, wie ein Löwe in der Höhle, und erdenket böse Tücke, auf daß er mich erhaschen möge; er spannet seinen Bogen und schießt seine feurigen Pfeile heimlich in mein elendes Herz. Erleuchte meine Augen, daß ich nicht vom höllischen Geiste verblendet werde, noch in meinen
\quellseite{342}
Sünden sterbe; gib ja nicht zu, daß sich dieser mein Feind rühme, er sei mein mächtig worden, und mein Widersacher sich nicht rühme, daß ich darniederliege. Er begehret meiner Seele, wie ein Löwe des Raubes; er stehet wider mich und ängstiget mich; darum, o Gott, behüte mich unter dem Schatten deiner Flügel, laß mich die alte listige Schlange nicht betrügen. Bewahre meine Seele und errette mich; wenn die bösen Geister an mich wollen, mich zu fressen, so laß sie anlaufen und fallen; verbirg mich in deinem Gezelt wider ihren Zorn. Weil ich so sehnlich zu dir, HErr, schreie, so schweige nicht, auf daß ich nicht gleich werde denen, die in die Hölle fahren; erkenne meine Seele in dieser großen Noth, streite wider meine Bestreiter, sprich zu meiner Seele: Ich bin deine Hülfe. Laß den Satan werden wie Spreu vor dem Winde; der Engel des HErrn stoße ihn weg, weil er mich gedenket in Herzeleid zu bringen; er müsse sich schämen und an mir zu Schanden werden. Denn du bist mein Gott und Stärke; wenn gleich die ganze Welt voller böser Geister wäre, mich zu verschlingen, so soll es doch, weil du bei mir stehest, ihnen nicht gelingen. HErr, schaffe meiner Seele Ruhe, weil der Fürst dieser Welt wider mich stehet. Gott sei mir gnädig, denn der böse Geist will mich versenken und streitet täglich wider mich; lege du aber Ehre ein und zerstreue meine leiblichen und geistlichen Feinde; schilt den Teufel, daß er von mir fliehe. Du guter und getreuer Hirte, Christe, der du mich von der Gewalt des Teufels mit deinem theuern Blut erlöset hast, errette meine Seele aus dem Rachen des höllischen Wolfes. Laß deine Gnade mein Trost sein; führe mich aus der Angst meines Herzens; dein guter Geist erfreue mich, daß ich deine Wunder preise und dir für deine Hülfe danken möge. Amen.

\section{398. Kurzes Gebet wider des Satans Mordpfeile}

Ich danke dir, mein lieber HErr JEsu Christe, daß du als ein Stärkerer über den starken Gewappneten kommen, ihn überwunden, seine Wehr und Harnisch, darauf er sich verließ, genommen, ihm sein Reich zerstöret und den Fürsten der Welt ausgestoßen, ihn auch sammt allen meinen Feinden öffentlich durch dich selbst Schau getragen hast, und bitte dich, ach mein herzallerliebster
\quellseite{343}
Herr Jesu Christe, du wollest dem verdammten Schaununstath kräftiglich steuren und wehren, daß er mich nicht in Vermessenheit und Verzweifelung stürze, mich auch nicht mit seinen Mordpfeilen der Versehung, außer Gottes geoffenbartem Wort, verwunde, und mich wegen der Schwachheit meines Glaubens oder anderer Ursachen wegen kleinmüthig und verzagt mache. Ach liebster Heiland, verleihe mir Kraft und Stärke, daß ich dadurch diese und alle seine feurige Mordpfeile überwinde und ausschlage, und endlich durch dich die Krone des Lebens, und den unvergänglichen Ehrenkranz im ewigen Leben erlange. Amen.

\section{399. Danksagung nach überwundener Anfechtung}

O ewiger Gott, der du Niemand lässest versucht werden über sein Vermögen, sondern schaffest, daß die Versuchung so ein Ende gewinne, daß er es kann ertragen, und hast zugesagt: Ich habe dich einen kleinen Augenblick verlassen, aber mit großer Barmherzigkeit will ich dich wieder sammeln. Ich habe mein Angesicht im Augenblick meines Zornes ein Wenig von dir gewendet, aber mit ewiger Gnade will ich mich deiner erbarmen. Ich hebe meine Hände zu dir auf, deinen Namen will ich allezeit erhöhen, dein Lob soll immerdar in meinem Munde sein; denn du weißt, in was für Jammer und Noth, Seelenangst und geistlicher Anfechtung ich eine Zeit lang gestecket, wie mein Geist sich in mir geängstet und mein Herz mir in meinem Leibe vor Unruhe und teuflischer Versuchung fast verzehret gewesen; aber da ich zu dir rief mit meinem Munde, so merktest du auf mein Flehen; denn du bist der Gott, der der Elenden Seufzen nicht verschmähet, und der betrübten Herzen Verlangen nicht verachtet. Du wendest dich zum Gebet der Verlassenen, und verschmähest ihr Schreien nicht, und machest los die Kinder des Todes. Darum so sei nun deine göttliche Kraft und Allmacht gelobet, deine grundlose Güte und Barmherzigkeit sei gebenedeiet, deine ewige Wahrheit sei gepriesen. Wo du, Herr, nicht hättest geholfen, so läge meine Seele schier in der Hölle, aber deine Gnade, Herr, hielt mich; ich hatte viel Bekümmerniß in meinem Elende, aber deine Tröstungen ergötzten meine Seele. Du hast mich lassen erfahren viel und große
\quellseite{344}
Angst, und hast mich wieder getröstet und lebendig gemacht, und aus der Tiefe gehollet. Darum danke ich dir auch für deine große Treue, mein Gott! Meine Lippen und meine Seele, die du erlöset hast, sind fröhlich und lobsingen dir; du bist freundlich und deine Güte währet ewiglich, und deine Wahrheit für und für. Ach treuer Gott, behüte mich durch deine große Güte forthin vor dergleichen Versuchung! Und da der leidige Satan (der dürre Stätte durchwandelt, und gern in dem Hause seine Wohnung aufschlagen wollte, da er zuvor ausgetrieben) mich ferner wollte anfechten, so stehe mir bei, du starker Gott, HErr Zebaoth; gib mir des Geistes Waffen gegen des Teufels listige Pfeile, daß ich ihn ritterlich überwinden möge; verleihe mir deinen guten Geist, der des bösen Geistes Werk zerstöre, und in mir Alles, das wider deine göttliche Ehre und meine Seligkeit ist, dämpfe. Zu dem gib mir Gnade, daß ich allen Untugenden von Herzen feind werde, und mich dir ganz allein ergebe, die ganze Zeit meines Lebens in wahrer Gottesfurcht und andächtigem Gebet und christlichen Tugenden zubringe, eine gute Ritterschaft übe, rechten Glauben und ein gut Gewissen bewahre, auf daß auch mir beigeleget sei und ich erlange die Krone der Gerechtigkeit, die Jesus Christus geben wird Allen, die seine Erscheinung lieb haben. Solches verleihe mir, du heilige Dreifaltigkeit, Gott Vater, Sohn und Heiliger Geist, hochgelobet in Ewigkeit. Amen.

\section{400. Gebet bei Gefahr von Verführung}

Barmherziger Gott, der du uns gewarnet hast, daß wir uns hüten sollen vor den falschen Propheten, welche in Schafskleidern zu uns kommen, inwendig aber reißende Wölfe sind; und weil in den letzten Tagen, darinnen wir jetzt leben, greuliche Zeiten kommen werden, da sich der Satan in einen Engel des Lichtes und die falschen Lehrer und trüglichen Arbeiter zu Christi Aposteln sich verstellen, welche den Schein eines gottseeligen Wesens haben, aber seine Kraft verleugnen sie; dazu nachdem auch der Abfall von der reinen Lehre und der Mensch der Sünden und das Kind des Verderbens, der da ist der Widerwärtige, offenbar werden soll, so hilf, ewiger Gott, daß wir die Liebe zur Wahrheit haben, alle Lugengeister
\quellseite{345}
sammt allen Irrthümern und allen falschen Schein vermeiden und fliehen, auf daß wir uns vom rechten Glauben nicht bewegen, noch Jemand das Ziel verrücken lassen, sondern standhaftig bei deinem Wort bis ans Ende verharren, und uns Niemand in keinerlei Weise lassen verführen, weder durch lügenhafte Kräfte, Zeichen und Wunder, noch mit irgend einer Verführung zur Ungerechtigkeit, daß wir uns nicht mögen oder wiegen lassen von allerlei Wind der Lehre, durch Schalkheit der Menschen und Täuscherei, damit sie an uns hantieren und gedenken uns zu erschleichen. Behüte, lieber Gott, deine Auserwählten in dieser armen Welt, damit sie nicht in Irrthum verführet werden. Verkürze die bösen Tage in diesen gefährlichen Zeiten. Bringe um den Widerchrist, das boshafte Kind aller Verführung, mit dem Geist deines Mundes. Mache sein ein Ende durch die Erscheinung deines lieben Sohnes. Bewahre uns auch vor allen Rottengeistern und Spaltungen, vor den greulichen Wölfen, die der Heerde nicht verschonen, vor Männern, die verkehrte Lehre reden und die Jünger an sich ziehen, vor Menschen, die zerrüttete Sinne haben und Zertrennung oder Aergerniß anrichten, vor dem Unkraut, welches der Feind aussäet, vor Dieben und Mördern der Seele. Laß uns deine Stimme allein hören, derselbigen von Herzen in wahrem Glauben und einfältigem Gehorsam folgen, auf daß dein Wort unser einig Richtschnur sei, darnach wir alle Lehre urtheilen, und die fremde fliehen, und keine Lust an der Ungerechtigkeit haben, so werden wir Freudigkeit haben, und nicht zu Schanden werden an jenem Tage, da erscheinen wird JEsus Christus, unser Heiland. Amen.

\section{401. Gebet für unschuldig (Gefangene) Verfolgte}

HErr! ich erhebe meine Stimme, HErr! dich suche ich, HErr! ich rufe dich an, Kreuz und Noth hat mich betroffen, ich harre dein, mein Gott! darum wollest du deine Ohren gnädig zu mir neigen, und auf mein Schreien hören. Ach, wie verfolgen mich die Gottlosen, und treten mich mit Füßen. Sie haben mir ohne Ursach ihr Netz gestellet, und meiner Seele eine Grube zugerichtet zum Verderben. Sie thun mir Arges um Gutes, und beißen ihre Zähne zusammen über mich, mich in Herzeleid zu bringen
\quellseite{346}
. Sie freuen sich über meinen Schaden, und hören nicht auf, mich unbillig zu hassen und zu ängstigen. Wie groß ist mein Kreuz, wie häufet sich meine Noth, wie mehret sich meine Trübsal! Mein Gott! betrübt ist meine Seele in mir, betrübt ist mein Herz, betrübt sind meine Gedanken, betrübt ist alles, was in mir und an mir ist, und ich weiß nicht, wo ich in dieser Welt bleiben, wo ich Hilfe und Trost erlangen soll; muß mich derowegen in die Höhe schwingen zu dir, mein Gott, meine Burg, meine Hilfe, mein Erretter, auf den ich hoffe! Ich klopfe an die Thür deiner Gnaden, und suche bei dir väterliche Hilfe und starke Errettung, darnach ich lange geseufzet und viel Thränen vergossen. Denn du bist ja ein treuer Gott, der mich und keinen Menschen mit mehr Trübsal beladen und beschweren wird, als wir ertragen können, der allmächtige Gott, dessen starke Hand unverkürzt ist, und aus aller Noth erretten kann. Abba, lieber Vater! ich verlasse mich auf dein väterliches Herz, damit du mir noch gewiß zugethan bist, ach, siehe doch mein großes Kreuz und Elend an, höre, wie ich zu dir so sehnlich schreie, und wende deine Barmherzigkeit nicht von mir, laß mich deine Güte und Treue allezeit behüten, und erlöse mich aus der grausamen Grube, und von der Hand derer, die mich hassen, und mir trachten Schaden zu thun. Herr, streite wider meine Bestreiter, und mache dich auf, mir zu helfen, wie du vormals vielen deine Hilfe erzeigt,- und dich einen starken Erretter bewiesen hast, derer Exempel ich armer und wohl geplagter Mensch mir zu Trost ansehe, und gedenke, daß auch mein Kreuz dir nicht so groß sei, daß du mir nicht helfen könntest. Mein Herr und mein Gott! laß meine Feinde und Verfolger mich nicht überwältigen und unterdrücken, laß sie nicht über mich sich freuen und sagen: Da, da, das wollen wir! wir haben ihn verschlungen; sondern erwecke dich, und wache auf zu meinem Rechte, führe meine Sache, gib mir christliche Geduld, und sprich zu meiner Seele: Ich bin deine Hilfe. Herr! behalte meinen Feinden ihre Sünden nicht, sondern bekehre sie, und vergib ihnen, was sie mir Uebels gethan. Laß sie mit Furcht und Zittern sehen, daß du mir beistehst, und mir Recht schaffest, mich aber freuen und fröhlich sein ewiglich, daß du mein Elend ansiehst, und erkennest meine Seele in der Noth. Amen, Amen.

\quellseite{347}

\section{402. Gebet Vertriebener und Verfolgter}

Ach du getreuer Gott, du einige Zuflucht der Verlassenen, erhöre uns nach deiner wunderlichen Gerechtigkeit, Gott unser Heil, der du bist die Zuversicht aller auf Erden, und ferne am Meer! Du bist unsere Zuflucht für und für; ehe denn die Berge worden, und die Erde und die Welt geschaffen worden, bist du, Gott, von Ewigkeit zu Ewigkeit. Erbarme dich unsers Elendes! Siehe, die Welt will uns nirgend leiden, wir werden verfolgt und haben nirgend eine bleibende Stätte, und wird an uns erfüllet, was unser Erlöser Christus Jesus uns zuvor gesagt hat: Siehe, ich sende euch wie Schafe mitten unter die Wölfe; ihr müßt gehasset werden von Jedermann um meines Namens willen; wer aber beharret bis ans Ende, der wird selig. Ach laß uns, lieber Vater, in wahrer Geduld und Beständigkeit ausharren! und weil du gesagt hast: Wenn sie euch in einer Stadt verfolgen, so fliehet in eine andere; so haben wir Hoffnung und Trost, du habest uns mit diesen Worten die Herberge bestellet, und werdest bei uns sein in unserer Flucht, wie bei dem Kindlein Jesu, und wie bei dem Erzvater Jakob, da ihm das Heer Gottes begegnete. Ach Herr Jesu Christe, laß dein Exempel unsern Trost sein! da du sprichst: So euch die Welt hasset, so wisset, daß sie mich vor euch gehasset hat; der Jünger ist nicht über seinen Meister. Ich habe euch von der Welt erwählet, darum hasset euch die Welt. Darum gib Geduld, weil es nicht anders sein kann, und weil alle die, so in dir, o Christe, wollen gottselig leben, Verfolgung leiden müssen, so gib Gnade, daß wir das Kreuz geduldig auf uns nehmen, und bedenken die Exempel der heiligen Apostel, wie davon St. Paulus sagt: Ich halte, Gott habe uns Apostel für die Allergeringsten dargestellet, als dem Tode übergeben; denn wir sind ein Schauspiel worden der Welt und den Engeln und den Menschen; bis auf diese Stunde leiden wir Hunger und Durst, und sind nackend, und werden geschlagen, haben nirgend eine gewisse Stätte; man schilt uns, so segnen wir, man verfolgt uns, so dulden wirs, man lästert uns, so flehen wir; wir sind stets als ein Fluch der Welt und Fegopfer aller Leute. Laß, o Herr Christe, dein Wort unsern Trost sein,
\quellseite{348}
da du sprichst: Selig sind, die um der Gerechtigkeit willen verfolgt werden, denn das Himmelreich ist ihr. Selig seid ihr, wenn euch die Menschen um meinetwillen schmähen und verfolgen, und reden allerlei Uebels wider euch, so sie daran lügen. Seid fröhlich und getrost, es wird euch im Himmel wohl belohnet werden; denn also haben sie verfolgt die Propheten, die vor euch gewesen sind. Ach HErr Gott, laß das Wort des heiligen Apostels Pauli auch unsern Trost sein: Wir werden verfolgt, aber wir werden nicht verlassen. Und abermal: Wir müssen durch viel Trübsal ins Reich Gottes eingehen. Laß uns auch an die freundlichen Worte St. Petri gedenken, da er spricht: Freuet euch, daß ihr mit Christo leidet, auf daß ihr auch zur Zeit der Offenbarung seiner Herrlichkeit Freude und Wonne haben möget. Selig seid ihr, wenn ihr geschmähet werdet über dem Namen Christi; denn der Geist der Herrlichkeit Gottes ruhet auf euch; bei ihnen ist er verlästert, aber bei euch ist er gepriesen. Ach Gott, wie ist allen deinen freudigen Bekennern dieses ein großer Trost! da du sagst: Wer euch antastet, der tastet meinen Augapfel an. Darum hoffen wir alle deine Pilgrime, und die wir ins Elend getrieben werden, du werdest uns behüten, wie einen Augapfel im Auge, und beschirmen unter dem Schatten deiner Flügel, und an allen Orten bei uns sein, uns begleiten und versorgen; denn die Erde ist des HErrn, und was darinnen ist, der Erdboden und was darauf wohnet. HErr, deine Gnade gehet, so weit der Himmel ist, und deine Wahrheit, so weit die Wolken gehen. Zähle unsre Flucht, fasse unsre Thränen in deinen Sack, ohne Zweifel, du zählest sie. Wohl dem, deß Hilfe der Gott Jakob ist, deß Hoffnung auf den HErrn, seinen Gott, stehet, der Himmel, Erde, Meer und Alles, was darinnen ist, gemacht hat, der Glauben hält ewiglich, der Recht schaffet denen, so Gewalt leiden, der die Hungrigen speiset. Der HErr löset die Gefangenen, der HErr machet die Blinden sehend, der HErr richtet auf, die niedergeschlagen sind, der HErr liebet die Gerechten, der HErr behütet die Fremdlinge und Waisen und erhält die Wittwen, und kehret zurück den Weg der Gottlosen; der HErr ist König ewiglich, dein Gott Zion für und für. Halleluja.

\quellseite{349}

\section[403. Gebet eines Christen, der unter …]{403. Gebet eines Christen, der unter Irrgläubigen leben muß}

Ach Gott! du siehest, wie ich leider unter solchen Leuten mich aufhalten muß, bei welchen weder dein seligmachendes Wort unverfälscht gelehret wird, noch die heiligen Sacramente ungeändert ausgetheilet werden. Dieweil denn mein Fleisch und Blut blöde ist, und sich leichtlich entweder aus Furcht des bittern Hasses, oder aus Begierde der schnöden Ehre könnte verleiten lassen; so bitte ich dich inniglich, daß du mich mit deinem Heiligen Geist stärkest, auf daß ich die erkannte evangelische Wahrheit höher denn aller Welt Gut achte, und mich davon weder Verheißung noch Bedrohung lasse abwendig machen. Verleihe mir auch eine christliche Vorsichtigkeit, auf daß ich also handle und wandle, damit man außer eines erforderten freudigen Bekenntnisses und Rechenschaft der Hoffnung, die in mir ist, keine Ursache an mir finden möge, mich zu bestricken oder zu bestrafen, und hilf mir zu rechter Zeit wiederum zu meinen Glaubensgenossen und zu dem rechten Gottesdienst, der allein auf deinem uns geoffenbarten, trostreichen Worte bestehet, um meines liebsten Herrn Jesu Christi willen. Amen, Amen.

\section[404. Gebet wider die Furcht bei unheimlichem …]{404. Gebet wider die Furcht bei unheimlichem Gepolter}

O Herr Jesu, du treuer Mittler zwischen Gott und den Menschen, der du nicht wolltest, daß deine Jünger sich sollten fürchten, da du durch verschlossene Thüren zu ihnen kamest, sondern ihnen alle Furcht benahmest und sprachest: Was seid ihr so erschrocken, und warum kommen solche Gedanken in eure Herzen? der du auch selbst in der einsamen Wüste mir zum Trost dich vom Satan versuchen lassen und durch dein kräftiges Wort denselben verjaget und der Engel Dienst empfunden hast: ich bitte dich von Herzen, laß mich keine plötzliche Furcht erschrecken, daß ich nicht weiß, wo ich hinaus soll. Du weißt, daß es um das menschliche Herz ein verzagtes Ding ist und durch geschwindes Schrecken auch manche in Noth und Tod gefallen sind; darum decke mich selbst mit deinen Fittigen, daß ich nicht erschrecken müsse vor dem Grauen des Nachts, noch vor den Pfeilen, die des Tages fliegen. Will mich Satan mit allem
\quellseite{350}
seinen Heer betrüben und furchtsam machen, so erhebe deine göttliche Stimme über mich und sprich: Hebe dich weg, Satan; und befehl deinen himmlischen Geistern, daß sie mich, Armen, vor diesem Schrecken schützen und befreien. Gib mir zu erkennen, daß ohne deinen Willen kein Haar von meinem Haupte fallen könne, und rufe mir in mein Herz: Fürchte dich nicht, denn ich bin mit dir, weiche nicht, denn ich bin dein Gott, ich stärke dich, ich helfe dir auch, ich erhalte dich durch die rechte Hand meiner Gerechtigkeit. Es schrecken mich zwar meine Sünden, weil ich an dir gesündiget und übel vor dir gethan habe, aber dein Blut machet mich auch rein von allen meinen Sünden. Darum erquicke mich allezeit in meiner Angst und sprich zu meiner Seele: Ich bin deine Hilfe. Stehe auf und laß alle deine Feinde zerstreuet werden, und die dich hassen, vor dir fliehen, damit ich still und sicher wohne und niemand sei, der mir leid thue. Amen.

\section{405. Gebet wider die Verleumder}

Mein Gott und Vater im Himmel! dir klage ich mit betrübtem und wehmüthigem Herzen die Schmach, welche mir jetzt von feindseligen Leuten durch schändlichen Argwohn und verleumderische Nachrede widerfährt. Denn du prüfest Herzen und Nieren und weißt am besten, wie unrecht sie an mir thun; du bist der allmächtige Gott, der meine Sache am besten ausführen und mir von den Lügenmäulern helfen kann. Zwar geschieht es nicht ohne dein gerechtes Verhängniß, daß meine Feinde mich also schmähen und das Gift ihrer Zungen wider mich ausgießen; und ob ich gleich diesesfalls vor den Menschen unschuldig bin, so habe ich doch mit unzähligen Sünden, die ich wider dich begangen, solches wohl verdienet. Aber ich verlasse mich mit kindlicher Zuversicht auf deine große Barmherzigkeit und hoffe, du wirst mein Elend gnädig ansehen, mein Weinen und Flehen erhören, und meiner Sünden um Christi willen nicht gedenken, sondern meine Unschuld zu rechter Zeit vor der Welt offenbar machen, daß es sehen, die mir Übles wollen, und sich schämen müssen. Jesu, mein Heiland! ich bedenke jetzt, welche Verleumdungen und Schmachreden du von deinen Feinden hören müssen, da du doch keine Sünde gethan hattest; und damit
\quellseite{351}
tröste ich mich bei allen unverdienten Lästerungen, die ich leiden muß. Du hast um meiner Sünde willen dich müssen verschmähen und verspeien lassen, und bist der Allerverachteteste und Unwerteste, ein Spott der Leute und eine Verachtung des Volks gewesen; so will ich demnach um deinetwillen und dir zu Ehren diese Schmach gern tragen. Aber, mein Seelenbräutigam! erbarme dich meiner nach deiner großen Güte, und errette meine Seele wieder von den bösen Menschen und frevelen Leuten, die da Böses gedenken in ihrem Herzen, und ihre Zungen schärfen wie Schlangen. Jedoch vergilt ihnen nicht nach ihren bösen Werken, und rechne ihnen diese Sünde nicht zur Verdammniß zu, sondern sei ihnen um deines Verdienstes willen gnädig, und hilf, daß sie sich bekehren und leben. Dir befehle und überlasse ich meine Sache, du wirst endlich meine Gerechtigkeit hervorbringen wie das Licht, und mein Recht wie den Mittag. Herr Gott Heiliger Geist! Tröster aller Betrübten! Du hast mich in deinem Wort versichert, daß denen, die Gott lieben, alle Dinge zum Besten dienen müssen. Darum bitte ich dich demüthigst, du wollest mich solches auch in gegenwärtiger meiner Noth erfahren lassen. Gib, daß alle grausamen Lügen und Spottreden meiner Widerwärtigen zu meinem Besten ausschlagen hier in der Zeit, damit ich desto vorsichtiger wandle, und allezeit züchtig und gerecht und gottselig lebe; zu meinem Besten gereichen dort in der Ewigkeit, da nach der Verheißung meines JESU einen solchen Gnadenlohn zu gewarten haben sollen Alle, die unschuldigerweise von Menschen geschmähet werden. Indessen aber verleihe mir christliche Geduld, daß ich ohne Murren dieses Kreuz trage, und nicht Böses mit Bösem, noch Scheltwort mit Scheltwort vergelte, sondern segne, die mir fluchen, liebe, die mich hassen, wohlthue denen, die mich beleidigen. Ach Herr, du dreieiniger Gott! thue solches in großen Gnaden, so will ich deinen Namen ehren und loben in Zeit und Ewigkeit. Amen.

\section{406. Gebet, wenn Kinder krank sind}

Lieber Vater, scharf ist deine Ruthe; aber Vater bleibst du, das weiß ich fürwahr. Und du, lieber Herr und Heiland, JESU Christe, der du ein Vorbild alles unseres Leidens gewesen, tröste
\quellseite{352}
und drücke dich selbst in unser Herz, auf daß ich dies Opfer dieses betrübten Geistes vollbringen, und dir unseren Isaac mit willigem Geist übergeben möge. Amen.

\section[407. Gebet eines armen gefangenen Sünders um …]{407. Gebet eines armen gefangenen Sünders um Vergebung der Sünden}

HErr, du gerechter und barmherziger Gott! deine Gerichte sind Gerechtigkeit und Wahrheit, du hast mir deine Gnade aus gerechtem Gerichte entzogen, und weil ich nach dem schändlichen Willen des Teufels, der bösen Gemeinschaft, und meines verderbten Fleisches, dein heilig Wort und das liebe Gebet muthwillig und verächtlich übergangen, mich deswegen in grobe und gräuliche Sünden, und also in Gefahr des Leibes und der Seelen sinken und fallen lassen, daß mir nun widerfährt, was ich werth bin, und wie ich verdient habe, indem ich nicht nur leiblich mit Ketten und Banden umgeben bin, und des Gerichts zum Tode, ja den gewaltsamen Tod selbst alle Augenblick zu gewarten habe, sondern auch in meiner Seele um begangener Sünden willen gar gräulich geängstiget, und mit großer Angst des nagenden Gewissenswurms, der Hölle, des Teufels, deines Zornes und des Todes geplaget werde. In dieser Noth und Trübsal aber erinnere ich mich aus deinem Wort, daß deine Barmherzigkeit und Liebe viel größer und mächtiger ist, als die große Last meiner Sünden; daß du allen armen bußfertigen Sündern, ohne Unterschied, mit einem Eide Gnade und Barmherzigkeit, auch Vergebung zugesagt hast; daß dein lieber Sohn, mein Erlöser, mit seinen Leiden und Tode alle arme Sünder, und also auch mich, den größten Uebelthäter, selig zu machen, in die Welt kommen ist; daß du mich in der heiligen Taufe zu einem Kinde und Erben des ewigen Lebens angenommen, durch deinen guten Geist in wahrem Glauben geheiliget, und mir die Frucht des Leidens und Sterbens Jesu Christi, ja Christum selbst mit seinem ganzen Verdienste versprochen, zugeeignet und geschenket hast. Worauf ich denn mit getroster Zuversicht zu dir schreie: Gott, sei mir armen Sünder gnädig, und handle nicht mit mir nach meinen Sünden; siehe an deine grundlose Gnade und Barmherzigkeit, und rechne nicht mit mir nach deiner gestrengen Gerechtigkeit. Gedenke
\quellseite{353}
doch an die Wunden und Tod deines lieben Sohnes Jesu Christi, meines werthesten Erlösers, an deine gnadenreiche Verheißung, und an meinen Taufbund, und laß mich, dein so theuer erkauftes Geschöpfe, und dein so hochgeheiligtes Kind, nicht ewig unter die Gewalt des Teufels und der Hölle gerathen. Tröste mich wieder mit deiner Hilfe, und stärke meine Seele in wahrem Glauben, und rechtschaffener Bußfertigkeit wider alle Anfechtung, Angst und Beschwerlichkeit. Ach! erbarme dich meiner, Herr, dreieiniger Gott! und laß dir mein Elend zu Herzen gehen, daß ich dort ewig bei dir leben, und dich loben möge. Amen.

\section[408. Gebet eines Verbrechers, der zum Tode …]{408. Gebet eines Verbrechers, der zum Tode verurtheilt worden ist}

O Herr, gerechter Gott! Ich armer Sünder klage und bekenne dir an diesem Tage meine von mir wohlverschuldete Betrübniß, und in dem jetzt herannahenden bösen Stündlein, wie mich der Satan von deinem Worte abgehalten, mir alle Zeichen deines gerechten Zornes aus den Augen gerissen, und es dahin gebracht, daß ich nicht geglaubt, daß du, als ein Feind der Sünden, mich strafen werdest. Ich habe meiner Lehrer Rath verachtet, und ihre Vermahnung in den Wind geschlagen, dich im Gebet um Bewahrung vor Sünden nicht ersuchet (mich zu böser Gesellschaft gehalten, Lügen, Müßiggang, Unzucht und außer ruchloses Leben geliebet, und gethan, was mir gelüstet). Darum hast du deine Hand von mir abgezogen, und mich anlaufen lassen, daß ich, weil du meinem Muthwillen nicht länger zusehen mögen, nun Andern zum Abscheu mit einem geschwinden schmählichen Tod belegen werde und ohne alle Gnade sterben muß. Ach Herr, mein Gott! ich gehe heute mit dem verlorenen Sohn in mein Gewissen, und beichte dir aus rechtschaffener Reue alle meine Mißhandlung; verzeihe sie mir um Jesu Christi willen. Herr, du bist gerecht, und alle deine Gerichte sind unsträflich, darum entsage ich jetzt, was meine Thaten werth sind. Barmherziger Gott, du hast eine lange Zeit zur Buße gegeben, wiewohl dir nicht unbewußt war, daß ich aus angeborner Bosheit meine Gedanken und Thun nicht ändern würde, aber ich bin mehr ärger als frömmer worden. Daher ist jetzt meine
\quellseite{354}
Seele hoch bekümmert, Stricke des Todes haben mich umgeben, und Höllenangst hat mich troffen. Ich bin in Jammer und Noth kommen. Laß mich, o lieber Gott! darinnen nicht verzagen, sondern sei mir großen Sünder gnädig. Du hast mich einen kleinen Augenblick des Zorns verlassen, und dein Angesicht vor mir verborgen; aber mit großer Barmherzigkeit wirst du dich mein erbarmen, und mich zum ewigen Leben wieder sammeln, um JEsu Christi willen. Amen.

\section[409. Drei Gebete eines zum Tode …]{409. Drei Gebete eines zum Tode Verwurtheilten in der letzten Stunde}

\subsection{Das erste zu Gott dem Vater}

Ach HErr, mein Gott, mein allerliebster Vater, die letzte Stunde meines Abschieds aus diesem Jammerthal ist nun vorhanden, da kann mir niemand denn du allein helfen. Ach HErr! der du niemals Jemand verschmähest, verstoße auch mich Armen zu dieser Stunde nicht, sondern laß die Gnadenthür des himmlischen Paradieses meiner armen Seele offen stehen; verleihe mir gnädiglich, daß sich mit meiner letzten Stunde alle meine andere, zeitliche und ewige Noth, Sünde und begangene Schande und Laster ende; ja verleihe, daß diese meine letzte Stunde auf Erden die erste sei bei dir in der ewigen Seligkeit, durch JEsum Christum. Amen.

\subsection{Das andere zu Gott dem Sohn}

O HErr JEsu Christe! ich bitte dich um der großen Bitterkeit willen, die deine edle und unschuldige Seele, da sie ist von deinem heiligen Leichnam abgeschieden, gehabt hat; erbarme dich, lieber HErr und Gott! über meine arme Seele in dieser Stunde, und in dem Augenblicke, wenn sie wird von meinem sündhaften Leibe abscheiden. HErr, erhöre mich, Christe erhöre, o du Gottes Lamm, das der Welt Sünde getragen, erhöre! Amen.

\subsection{Das dritte zu Gott dem Heiligen Geiste}

O Heiliger Geist! der du von dem allmächtigen, ewigen Gott gesendet wirst zu einem Tröster und Beistand aller Betrübten. Ach, thue auch in dieser Stunde, in diesen meinen letzten Nöthen,
\quellseite{355}
einen treuen Beistand meinem armen schwachen Herzen und zaghaften, blöden Gewissen, daß ich an meinem Erlöser und Seligmacher Christo beständig bis an mein Ende verharre, wenn auch mein Mund nimmer reden würde, ach! so sprich du für mich und meine arme Seele: Vater, in deine Hände befehle ich dir meinen Geist, du hast mich erlöset, o du HErr Gott! Amen.

\section{410. Gebet eines verlassenen armen Waisenkindes}

Ach allerliebster HErr JEsu Christe, der du gesagt hast: „Ich will euch nicht Waisen lassen!“ erfülle an mir deine gnädigste Verheißung und wahrhaftiges Wort, der ich eine arme Waise bin. Ach lieber HErr, nimm du mich auf; du bist barmherzig und gnädig, geduldig und von großer Güte und Treue, der du beweisest Gnade und vergibest Missethat, Uebertretung und Sünde. Du siehest ja, denn du schauest das Elend und Jammer. Es stehet in deinen Händen. Die Armen befehlen es dir. Du bist der Waisen Helfer. Das Verlangen der Elenden hörest du, HErr, ihr Herz ist gewiß, daß dein Ohr darauf merket, daß du Recht schaffest den Armen und Waisen. Ach! vergib mir meine Unwissenheit, Sünde und Thorheit, damit ich dir und meinen lieben Eltern jemals zuwider gewesen und diese und andere Strafen, Jammer und Herzeleid verdienet habe. Gedenke, HErr, an deine Barmherzigkeit und an deine Güte, die von der Welt her gewesen ist. Du bist recht ein Vater der Waisen, sei du auch mein Vater, beschütze mich mit deiner Barmherzigkeit, bewahre mich mit deiner Allmacht, und laß dein Wort eine starke Mauer um mich sein. Ach, herzliebster HErr JEsu, gib mir Nahrung und Kleider, beschere mir gute Freunde, treue Vormünder und Förderer, die mir mit Rath und That behülflich sein. Regiere mich auch sonderlich durch deinen Heiligen Geist, daß ich unsträflich wandele, und mich halte nach deinem Wort, fromm und gottesfürchtig sei, wie der verlassene Joseph und der fromme Daniel, wie die gottesfürchtige Esther; denn die Gottseligkeit ist doch zu allen Dingen nütze und hat die Verheißung dieses und des zukünftigen Lebens. Gib, daß ich dich mein Lebenlang vor Augen habe, und mich hüte, daß ich in keine Sünde willige und wider deine Gebote thue. Hilf, daß ich mich
\quellseite{356}
enthalte von fleischlichen Lüsten, welche wider die Seele streiten, alle Hurerei vermeide, und meinen Leib und meine Seele bewahre in Heiligung und Ehren; daß ich auch gern sei bei den Alten, nicht verachte, was sie reden, und mich richte nach ihren Sprüchen. Laß mich still und verschwiegen sein, daß ich mich schäme, nachzusagen Alles, was ich höre, und zu offenbaren heimliche, vertraute Reden. Bescherest du mir nun einen Freund und Förderer, so hilf, daß ich ihn auch wie ein Auge halte, ihm treu und verschwiegen, auch sonst schamhaftig, züchtig und ehrbar sei, nur rede, was nötig ist; denn durch Scham, Zucht, Verschwiegenheit, Treue und Redlichkeit ist manch armes Kind und elendes Weiblein befördert worden, zu Gut und Ehre gelanget. Behüte mich auch vor böser Gesellschaft und leichtfertigen Leuten, daß ich nicht ihre Wege lerne und meine Seele Ärgernis empfange, von ihnen Hoffart lerne, daß sie mich in Unglück bringen, und ich um ihrer Thorheit willen Schaden leiden müsse. Leite mich in deiner Gerechtigkeit, lehre mich thun nach deinem Wohlgefallen, dein guter Geist führe mich auf ebner Bahn, und erhalte mich zum ewigen Leben um deiner Barmherzigkeit willen. Amen.

\section{411. Gebet in Armuth}

Barmherziger, gnädiger Gott, lieber Vater, du hast mir das Kreuz der leiblichen Armuth auferleget, ohne Zweifel mir zum Besten; gib mir Geduld, dieses Kreuz recht zu tragen, und deinem gnädigen Willen mich gehorsam zu unterwerfen; denn es kommet Alles von dir: Glück und Unglück, Armuth und Reichthum, Leben und Tod. Du machest arm und machest reich, du erniedrigest und erhöhest. Zeige mir aber, du getreuer Gott, Mittel und Wege, wie ich mich und die Meinen ehrlich ernähren möge, und segne meinen Beruf und Arbeit. Du hast ja gesagt, daß sich alle Armen ihrer Handarbeit nähren sollen; laß diesen Segen auch über mich kommen, laß mich aus deiner milden Vaterhand auch essen, und mit den Meinigen gespeiset werden, wie die Vögel unter dem Himmel, die weder säen noch ernten, und du himmlischer Vater nährst sie doch. Du kleidest und schmückest jährlich die Lilien-Blume auf dem Felde mit neuem Röcklein, ja du gibst einem jeden jährlich
\quellseite{357}
einen neuen Rock nach seiner Art; du wirst ja meiner und der Meinen nicht vergessen, denn du, himmlischer Vater, weißest, daß wir das Alles bedürfen. Danach lehre mich zuvörderst suchen dein Reich und deine Gerechtigkeit, so wirst du mir nach deiner Verheißung auch die zeitliche Nothdurft zuwerfen. Und weil ich auf Erden keinen Trost habe — denn dem Armen ist Niemand Freund — so sei du mein Trost, der du der Armen und Geringen Zuflucht und Stärke genannt wirst. Laß mir meine Armuth nicht Anlei-tung geben zum Bösen und zu verbotenen Mitteln, sondern zu desto stärkerem Glauben und Vertrauen auf dich und zu allen christlichen Tugenden, zur Demuth, Sanftmuth, Geduld, Gebet, Hoffnung und Beständigkeit. Gib, daß ich Niemand mit der Armuth beschwerlich sei, sondern erwecke mir freiwillige Wohlthäter, die du lieb hast, und verleihe, daß allerlei Gnade bei ihnen wohne. Ach lieber Gott, ein Armer ist wie ein Fremdling auf Erden, welchen Niemand kennen will, deß sich Niemand annimmt. Das ist aber mein Trost, daß du gesagt hast: Ihr sollt die Waisen und Wittwen und Fremdlinge lieb' haben, und ihnen Speise und Kleider geben, und euch fürchten vor dem HErrn euerm Gott. So ist auch das mein Trost, daß geschrieben stehet: der HErr hebt auf den Dürftigen aus dem Staube, und erhöhet den Armen aus dem Koth, daß er ihn setze unter die Fürsten, und den Stuhl der Ehren erben lasse. Du, lieber Vater, hast es also geordnet, daß Reiche und Arme müssen unter einander sein; du aber, HErr, hast sie alle gemacht. Darum, lieber Vater, nimm dich auch meiner an, und laß mir nicht Unrecht noch Gewalt geschehen, daß ich nicht unterdrückt werde. Du fragst ja: Es ist besser ein Armer, der in seiner Frömmigkeit gehet, denn ein Reicher, der in verkehrten Wegen gehet. Laß mich eingedenk sein, was der alte Tobias seinem Sohne sagt: Wir sind wohl arm, aber wir werden viel Gutes haben, so wir Gott fürchten, die Sünde meiden und Gutes thun. Und was David sagt: Das Wenige, das ein Gerechter hat, ist besser, denn das Gut vieler Gottlosen. Ich bin jung gewesen und alt worden, und habe noch nie gesehen den Gerechten verlassen oder seinen Samen nach Brod gehen. Dieser Verheißungen tröste ich mich, und lasse mich wohl begnügen; denn es ist besser Wenig mit
\quellseite{358}
Gerechtigkeit, denn viel Einkommens mit Unrecht; denn wir haben nichts in die Welt gebracht, werden auch nichts mit hinaus nehmen. Darum bitte ich dich um ein solches Herz, welches sich mehr um den ewigen Reichthum bekümmert als um zeitlich Gut — du wirst mir mein bescheiden Theil wohl geben. Laß mich folgen der schönen Lehre des weisen Hauspredigers, da er sagt: Vertraue Gott und bleibe in deinem Beruf; denn es ist dem HErrn gar leicht, einen arm oder reich zu machen. Laß mich anschauen das Exempel meines Erlösers JEsu Christi, da er spricht: Die Vögel unter dem Himmel haben ihre Nester, und die Füchse ihre Löcher, aber des Menschen Sohn hat nicht so viel, da er sein Haupt hinlegt. Du, HErr, bist mein Gut und mein Theil, du erhältst mein Erbtheil, du erfreuest mein Herz, ob Jene gleich viel Wein und Korn haben; ich bin arm und elend, der HErr aber sorget für mich; ich hoffe auch, daß ich sehen werde das Gute des HErrn im Lande der Lebendigen. Seid getrost und unverzaget alle, die ihr des HErrn harret.

\section{412. Trost in Armuth und Mangel}

Ob ich gleich hier Armuth leide, schadet nicht; dennoch weiß ich, daß du, mein lieber Gott, mich wirst nicht lassen Noth leiden, denn du hast mir Christum gegeben und alle Seligkeit in ihm. Du wirst mir auch so viel zuwerfen, daß der Leib die kurze Zeit seines Lebens seine Nothdurft haben wird. Amen.

\section{413. Gebet um Ueberwindung in aller Noth}

Ach Vater und Gott alles Trostes, verleihe uns durch dein heiliges Wort und Geist festen, fröhlichen und dankbaren Glauben, damit wir diese und alle Noth mögen selig überwinden, und endlich schmecken und erfahren, daß es die Wahrheit sei, da dein lieber Sohn Christus selbst spricht: Seid getrost, ich habe die Welt überwunden. Amen.

\section[414. Gebet eines Unterdrückten um der …]{414. Gebet eines Unterdrückten um der Wahrheit willen}

Ich leide viel, und gehet mir übel; aber meinen Feinden gehet es wohl. Sie leben, ich sterbe ohne Unterlaß. Sie sind
\quellseite{359}
mächtig und stark, ich werde ohne Unterlaß niedergedrückt. Sie sind in Ehren, ich in Schmach. Sie im Frieden, ich im Unfrieden. Sie mehren sich und haben ihrer viel, die ihnen günstig sind, die sie loben, die es mit ihnen halten; ich bin allein, verlassen, und niemand hält's mit mir, auch ist mir niemand günstig, ich bin ein Einsamer, von allen verlassen und verachtet. Darum, lieber HErr Gott, nimm mich auf, und verlaß mich nicht, eile mir zu helfen, denn alle andere helfen mir zum Verderben. Ich suche kein Heil noch Seligkeit, weder in mir selbst, noch in jemand anders, denn bei dir allein. Amen.

\section[415. Gebet eines in einen Prozeß verwickelten …]{415. Gebet eines in einen Prozeß verwickelten Christen}

O HErr, du bist ja ein gerechter Gott und hast Gerechtigkeit lieb; aber in der Welt gehet es oft gar anders, denn recht und der Gottlose übervortheilt den Gerechten. Ungerechtigkeit hat leider! überhand genommen; viel sind, die das Gute hassen und das Arge lieben, die das Recht in Wermuth und Galle verkehren und die Gerechtigkeit zu Boden stoßen. Wie viel findet man, so die Sache der Armen beugen und den richtigen Weg der Elenden hindern, ja die Gerechten um Geld und die Armen um ein Paar Schuhe verkaufen, daß die Wittwen ihr Raub und die Waisen ihr Blut sein müssen! Was der Fürst will, das spricht der Richter, daß er ihm wieder einen Dienst thun soll, wie dein heiliges Wort darüber klaget. Die Gewaltigen rathen nach ihrem Muthwillen und drehen's, wie sie wollen. Es gehet Gewalt über Recht, das klaget manch armer Knecht. HErr, du weißt ihr Uebertreten, wie sie die Gerechten dringen und Geld nehmen und die Armen im Thor und Gerichte unterdrücken, denn es ist eine böse Zeit, darum wollest du, du gerechter Richter, hören im Himmel und Recht schaffen den Elenden. Da ich aber ja in dieser Welt zu meinem Rechte nicht gelangen kann, sondern der Gewaltigen Trutz fortgehet, so verleihe mir Geduld, daß ich mich nicht entrüste, sondern mir vielmehr Unrecht thun und mich verurtheilen lasse, als daß ich dem Uebel widerstreben und Böses fürnehmen wollte. Regiere mich auch mit deinem Heiligen Geist, daß ich vor Gericht nicht zänkisch sei, sondern den Richter in Ehren halte und mich schäme, für meine
\quellseite{360}
Seele das Recht zu bekennen, und da ich etwa in dieser Sache gefehlet, nicht wider die Wahrheit rede. Da ich aber bei meiner gerechten Sache unterliegen müßte, ach HErr, so stärke mich, daß ich meine Seele in Geduld fasse und dir meine Sache sammt der Rache in kindlichem Gehorsam befehlen möge. Amen.

\section{416. Gebet im Verlust zeitlicher Habe}

Lieber Gott, du siehest und weißest alle meine Noth, Anliegen, Verlust und Schaden an Habe, Gut und Ehre, du stehest ja und schauest mein Elend und Jammer, dir sei es allein heimgestellt und anbefohlen; hilf du mir, so ist mir geholfen, gib Geduld und Trost in dieser großen Angst und Drangsal. Siehe, hie bin ich, mach’s mit mir, wie es dir wohlgefället, laß mich nach ausgestandener Hiobslage auch wieder deines getreuen Knechtes, des geduldigen Hiobs, Trost und Freude empfinden hier zeitlich und dort ewiglich. Amen.

\section{417. Gebet in Feuersgefahr}

Ach HErr, HErr, gnädig und barmherzig, geduldig und von großer Güte und Treue! wir erkennen mit herzlicher Reue und Buße, daß wir die vor Augen schwebende Feuersbrunst wohl verdienet, welche leichtlich überhand nehmen, und uns wie Adama und Zebdim zurichten könnte. Dieweil du aber nicht Lust hast an unserm Verderben, sondern vielmehr willst, daß sich der Sünder bekehre und lebe: so bitten wir dich herzlich, du wollest auch diese wohlverdiente Strafe unsrer Sünden sammt allem besorglichen, reißenden und tobenden Wüthen gnädiglich abwenden, mitten im Zorn deiner Gnade und Barmherzigkeit eingedenk sein, und dem grausamen Feuer gebieten und sagen: Es ist genug! um deines lieben Sohnes Jesu Christi willen, der uns von deinem gerechten Zorn und Strafe hat befreit, all unsre Sünden getragen, völlig gebüßet und für uns bezahlet. Ach HErr, höre! ach HErr, sei gnädig! ach HErr, merke auf und thue es, und verzeuch nicht, um dein selbst willen, HErr unser Gott. HErr Gott Vater im Himmel, erbarme dich über uns! HErr Gott Sohn, der Welt Heiland, erbarme dich über uns! HErr Gott, Heiliger Geist, erbarme dich
\quellseite{361}
über uns! Sei uns gnädig, verschon uns, lieber HErr Gott. Sei uns gnädig, hilf uns, lieber HErr Gott. In dieser großen Feuersnoth behüte uns, lieber HErr Gott. Christe, du Lamm Gottes, der du trägst die Sünde der Welt, erbarme dich unser! O JEsu Christe, Gottes Sohn, erbarme dich über uns, so wollen wir, dein Volk und Schafe deiner Weide, dir nächst schuldigen Gehorsam und ernster Besserung unseres sündhaften Lebens danken ewiglich und deinen Ruhm verkündigen für und für. Amen.

\section{418. Gebet in großer Wassersnoth}

HErr Gott, du bist unsre Zuflucht für und für, eine Hülfe in den großen Nöthen, die uns getroffen haben. Der Name des HErrn ist ein festes Schloß, und der Gerechte läuft dahin und wird beschirmet, darum kömmt auch zu dir alles Fleisch. Ach, schaue von deiner heiligen Höhe und siehe an unsern Jammer und Elend! Denn deine Wassermogen und Wellen gehen über uns. Wir sind im tiefen Wasser, und die Fluth will uns ersäufen, und Niemand ist, der uns erretten kann, als du, HErr, alleine. Darum hilf du uns, Gott, unser Helfer, um deines Namens Ehre willen! Unsere Sünden sind zwar nicht geringer, denn der ersten Welt, die du auch mit Wasser überschwemmet und heimgesucht hast; aber gedenke an deine Barmherzigkeit und an deine Güte, die von der Welt her gewesen ist. Erfülle an uns die Verheißung, durch David gesprochen, daß, wenn große Wasserfluthen kommen, sie nicht an die gelangen sollen, welchen die Uebertretungen vergeben sind, denen die Sünde bedeckt ist, und der HErr die Missethat nicht zurechnet. Ach, HErr, die Wasserströme erheben sich, die Wasserströme erheben ihr Brausen, und Wassermogen sind groß und brausen greulich; aber sei du, HErr, noch größer in der Höhe. Brich ihnen den Lauf mit deinem Damm, und setze ihnen Riegel und Thür. Sprich selbst zu dem Wasser: Bis hieher sollst du kommen und nicht weiter, hier sollen sich legen deine stolzen Wellen. Verstopfe die Brunnen der Tiefen, sammle den Fenstern des Himmels, daß die Gewässer sich verlaufen von der Erden, und der Erdboden wieder trocken werde. Vergiß nicht, was du sonsten deinem Diener Israel hast zugesagt, daß, wenn sie durchs Wasser gehen, du bei ihnen sein wollest, daß sie
\quellseite{362}
die Ströme nicht sollen ersäufen. Darum erbarme dich aller derer, die wirklich in solcher Gefahr sind, und denen das Wasser gehet bis an die Seele. Rette sie aus dem Wasser, wie die Tochter Pharao den Mose, oder behalte sie, wie den Noah mit den Seinigen, mitten in solcher Wassergefahr. Sollten aber nach deinem gerechten Gericht diese deine Zornesfluthen Jemanden ersäufen, so sei doch deinem Volke gnädig, behalte ihnen ihre Sünde nicht, sondern nimm auf ihren Geist, und laß es ihrer Seele nicht schaden. Ja, sollte auch mich das Unglück treffen, so geschehe dein Wille, ich befehle dir Leib und Seele und Alles, was mir angehöret, in deine Hände. Denn nackend bin ich von Mutterleibe gekommen, nackend werde ich auch wieder dahinfahren. Nun, Gott, sei nicht ferne von mir; mein Gott, eile mir beizustehen und zu helfen, um Christi willen. Amen.

\section{419. Gebet, wenn einer zum Krüppel geworden ist}

Allerliebster Herr JESU, der du der Lahmen, Krüppel und Gebrechlichen in deinem Wort nicht vergessen, und dich ihrer in den Tagen deines Fleisches auf das allergetreulichste angenommen hast, ach, erbarme dich auch über mich elenden und gebrechlichen Menschen! Erwecke bei ehrlichen Leuten christliches Mitleiden, daß sie mir nothdürftige Handreichung thun. Gib mir vor Allem eine gesunde Seele, erhalte mich in einem christlichen Leben und laß mich wohl bedenken, daß ich vielleicht mit diesen meinen Gliedern schwer würde sündigen, wenn sie noch alle gesund und stark wären; nun mir aber dieselben verdorben sind, suchest du dadurch meiner Seelen Wohlfahrt, und ist mir freilich besser, daß ich als ein Krüppel zu dem Leben eingehe, denn daß ich als ein Gesunder sollte in das höllische Feuer geworfen werden. So weiß ich auch, daß dieser mein ungesunder Leib alle seine Gebrechlichkeit wird dermaleins im Grabe lassen, und ich werde ganz gesund und frisch in großer Kraft zum ewigen Leben auferwecket werden. Solchen Trost bewahre, du getreuer Gott, um des bittern Leidens und Sterbens JESU Christi willen, durch die Gnade deines Heiligen Geistes in meinem Herzen, und hilf, daß, wenn gottlose Leute meiner spotten, ich ja nicht ungeduldig werde, noch in meinem Elend verzage, oder jemanden seine Gesundheit
\quellseite{363}
mißgönne, sondern fleißig bete, meine Hoffnung auf dich setze und sicherlich glaube, du werdest mich ernähren, kleiden, speisen, tragen und heben bis ins ewige Leben. Amen.

\section{420. Gebet in großer Noth und Gefahr}

O allmächtiger, hilfreicher Gott und Vater, du siehest und erkennest, daß wir in sehr großer Noth und Gefahr leben, weil wir weder Rath noch Hülfe wissen, und in uns keine Kraft noch Macht ist, solcher großen Noth zu entgehen, so sehen unsere Augen nach dir, der du unser einiger Trost, Zuflucht und Helfer bist und gesagt hast: Rufe mich an in der Noth, so will ich dich erretten, und du sollst mich preisen. Unser Herz hält dir vor das Wort deiner Verheißung, und suchet nun in kindlicher Demuth und Zuversicht dein Antlitz. Ach, HErr! verbirg das Antlitz deiner Gnade nicht vor uns, und verstoße nicht im Zorn deine Kinder; verlaß uns nicht, und zeuch deine Hand nicht von uns ab, Gott, unser Heil, der du alle Hülfe thust, so auf Erden geschieht. So laß uns darum deine Hülfe widerfahren, und sei uns auch gnädig. Auf dich, HErr, trauet unsere Seele, und unter dem Schatten deiner Flügel haben wir Zuflucht, bis das Unglück vorübergehe. Kehre dich doch wieder zu uns, und fülle uns frühe mit deiner Gnade, so wollen wir dir ein Freudenopfer thun, und deinem Namen danken, daß er so tröstlich ist. Erfreue uns nun wieder, nachdem du uns so lange geplaget, nachdem wir so lange Unglück leiden. Schaffe uns Beistand in der Noth, so wollen wir dich rühmen und fröhlich sein über deiner Hülfe unser Lebenlang. Amen.

\section{421. Ein Anderes}

Ich weiß gewiß, daß mich dennoch unser HErr Gott herzlich lieb hat, ob ich gleich jetzt in dieser großen Noth stecke und nicht sehe, wie mir geholfen könnte werden. Ich befehle es aber meinem lieben Gott, der jetzt in diesem Jammer auf mich siehet wie eine Mutter auf ihr Kindlein, das sie unter ihrem Herzen getragen hat, der wird es wohl machen, den will ich auch darum bitten und gewißlich glauben, daß er mich hören und erretten wird. Denn wenn die Gerechten schreien, so erhöret der HErr und errettet sie aus ihrer Noth. Amen.

\quellseite{364}

\section{422. Noch ein Anderes}

Ach lieber Vater, du bist ja mein lieber Vater, denn du hast ja deinen einigen lieben Sohn für mich gegeben. Darum wirst du ja nicht mit mir zürnen, noch mich verstoßen, du siehest meine Noth und Schwachheit. Darum wollest du mir helfen und mich erretten. Amen.

\section{423. Noch ein Anderes}

Ach Vater aller Barmherzigkeit und Gott alles Trostes, stärke und kräftige mich durch deinen Geist, bis erscheine und komme dieses Werk, darauf du uns heißest harren in allen Trübsalen, denn du nicht von Herzen die Menschen plagest und betrübest, und du lässest nichts Böses geschehen, wenn du nicht etwas Gutes daraus könnest schaffen. Amen.

\section{424. Gebet in großer Widerwärtigkeit}

Ach Gott, strafe mich nicht im Zorn, laß in Gnaden sein und zeitlich, sei Vater und nicht Richter. Ach Gott, zürne hier, haue hier und schlage hier, und schone unser dort. Amen.

\section{425. Ein Anderes}

Lieber Vater, schlage und streiche getrost zu, ich habs leider wohl verschuldet; doch laß es eine Vater-ruthe sein. Wie du denn alle deine Kinder, so du lieb hast, stäupest, denn welche du um ihrer Sünde willen nicht strafest, sind nicht Kinder, sondern Bastarde: darum streiche, peitsche, dresche flugs auf uns, gerechter Richter, doch auch barmherziger Vater, also aber, daß du dein göttlich Vaterherz nicht von uns wendest, auf daß wir dich hier und dort ewiglich loben und preisen mögen. Amen.

\section{426. Gebet in theurer Zeit}

Du gnädiger und barmherziger Gott, lieber Vater! wir erkennen und bekennen mit reuigem Herzen, daß wir die Zeit unsers Lebens sehr übel angewendet und gottlos gelebt haben, wir bitten aber, du wollest durch das hoch theure Verdienst deines allerliebsten Sohnes Jesu Christi, unseres Heilandes, uns armen Sündern
\quellseite{365}
gnädig sein, und unsere Missethat verzeihen und austilgen, und ferner verleihen, daß wir dies dein Kreuz und Heimsuchung mit Geduld ertragen und leiden. Unter anderem Kreuz aber, welches du auflegest, ist unserm Fleisch und Blut kein unerträglicheres, denn die bittere Hungersnoth und der hochbeschwerliche, schmähliche Mangel täglicher Leibesnahrung. Solche Strafe haben wir gar wohl verdienet durch großen Ueberfluß, schändliche Pracht und Mißbrauch deiner vielfältigen Gaben, dafür wir dir auch leider wenig Dank gesagt. Weil du uns aber in deinem Wort deine milde Hand und leibliche Unterhaltung als ein gütiger Vater zugesagt hast, denn also spricht dein lieber Sohn Matth. 6.: Ihr sollt nicht sorgen und sagen, was werden wir essen, was werden wir trinken? Euer Vater weiß, daß ihr das alles bedürfet; du als ein allwissender Gott weißt am besten, was uns, als deinen armen Kindern, nützlich und gut ist: so bitten wir dich von Grund unsers Herzens, weil großer Hunger und vielfältiger Mangel allenthalben uns drücket und ängstiget, du wollest deiner väterlichen Verheißung eingedenk sein und uns unser liebes täglich Brod väterlich mittheilen. Deiner wahrhaftigen Zusage trauen wir gewißlich, denn was du zusagst, das hältst du gewiß; denen, so dich im Glauben beständig suchen und lieben, soll nichts an irgend einem Gute mangeln. Du thust deine milde Hand auf, und sättigest alles, was lebet, mit Wohlgefallen. Die Elenden sollen essen, daß sie satt werden, und die nach dem HErrn fragen, werden ihn preisen. Der HErr wird euch im Trübsal Brod, und in Aengsten Wasser geben; der HErr wird deinem Samen, den du auf den Acker säest, Regen geben, und Brod von des Ackers Einkommen, und desselben volle Genüge. Du hast vorzeiten Elian, den Propheten, durch einen Engel und Raben erhalten; du hast den Propheten Elisa durch eine Wittfrau gespeiset. Du hast Daniel in der Löwengrube gesättiget. Du hast durch deinen lieben Sohn mit einem kleinen und sehr geringen Vorrath, fünf Gerstenbroden und gar wenig Fischen, fünftausend Mann nicht allein reichlich gesättiget, sondern auch zwölf Körbe voll der Brocken aufsammeln lassen. Darum sind wir getroster Zuversicht und glaubensreicher Hoffnung zu dieser theuern Zeit, du werdest uns auch
\quellseite{366}
in dieser schweren Zeit wohl wissen und können väterlich ernähren, bis daß wir diesem elenden Leben durch den natürlichen Tod Urlaub geben und endlich mit dir und allen Auserwählten in deinem ewigen Reich über deinem Tisch essen und trinken die himmlische, ewige und unvergängliche Speise und Trank der ewigen, von dir hochversprochenen und durch Christum theuer erworbenen Freude und Seligkeit. Denn du kannst und willst deine Kinder nicht verlassen, dieweil wir in diesem Elendsthal wandern, du hast ein Vater- und Mutterherz. Weil du uns geistlich speisest, wirst du uns leiblich Essen und Trinken zur Nothdurft auch nicht versagen. Du gibst dem Vieh sein Futter und speisest wunderbarer Weise die jungen Raben, viel weniger wirst du unser vergessen, die du mit deines lieben Sohnes theurem Blut gewaschen und gereiniget hast. Solches wollest du thun, o ewiger Vater! um deines geliebten Sohnes, unsers einigen Erlösers und Seligmachers Jesu Christi willen. Amen.

\section[427. Gebet in ansteckenden Seuchen und …]{427. Gebet in ansteckenden Seuchen und Sterbensläuften}

Allmächtiger, ewiger Gott und Vater unsers HErrn Jesu Christi, ein HErr Himmels und der Erde, wir arme, elende Sünder müssen bekennen, daß wir dich, o Gott, mit unsrem sündlichen Leben und Wesen auf’s greulichste erzürnet haben, darum du auch billig über uns deinen Zorn ausgeußt und uns mit mancherlei Plagen, Seuchen und Krankheiten angreifst. Wie sollen wir nun thun? Sollen wir verzagen? Das sei ferne! Wir erkennen zwar, daß wir nicht allein diese jetzt regierende Seuche, als eine Strafe unserer vielfältigen Sünden, sondern auch wohl größere und ärgere verdienet haben. Wo sollen wir denn hinfliehen, wo sollen wir uns hinwenden und kehren, daß wir mögen vor dieser und anderen Seuchen und Plagen sicher sein? Allein zu dir, HErr Jesu Christ, wir haben sonst keinen Trost weder im Himmel noch auf Erden, ohne dich, der du uns erlöset hast. Du wirst ja dein Geschöpf nicht verwerfen; darum rufen, seufzen und schreien wir zu dir demüthiglich von ganzem Herzen und sagen: Gott sei uns gnädig und tilge alle unsere Sünde nach deiner übergroßen Gnade, Güte und Barmherzigkeit. Laß ab von deiner Ungnade, Grimm
\quellseite{367}
und Zorn über uns. Erzeige uns wiederum deine Gnade und verschone uns mit dieser jetzt regierenden Seuche und abscheulichen Krankheit gnädiglich. Laß dich erbitten, o Herr, laß dich erbitten und verschone uns. Beschütze und beschirme uns väterlich, daß uns diese Seuche nicht beschädige oder treffe, hinraffe oder wegnehme. Ist es aber ja dein göttlicher Wille, daß wir in solcher Seuche unser Leben beschließen, und von dieser Welt abscheiden sollen, so geschehe derselbe dein gnädiger Wille, welcher allezeit der beste ist. Darauf so befehlen wir uns nun mit Leib und Seele, mit Weib und Kind und ganzem Hausgesind in deine göttliche Gnade und väterliche Hand, und bitten von Herzen demüthiglich, wenn wir also mit einem solchen Stündlein übereilet werden und nun an dem ist, daß unser Leib und Seele sich scheiden sollen, du wollest uns gnädiglich bei guter Vernunft erhalten, daß wir dir unsre Seele fein mit Verstand befehlen können, und verleihe uns dazu ein seliges Ende, auf daß wir also durch den zeitlichen Tod, welcher ein Ende ist alles Jammers und Elendes und uns die Thür zum ewigen Leben aufmacht, in das ewige Leben dringen und desto eher zu unserm Erlöser und Seligmacher kommen und uns mit allen Auserwählten Gottes im Himmel ewig freuen mögen. Solches wollest du uns, o treuer Gott und Vater, gnädiglich widerfahren lassen durch JEsum Christum unsern HErrn, welcher mit dir und dem Heiligen Geist lebet und regieret, gleicher Gott, hochgelobet in Ewigkeit. Amen.

\section[428. Gebet Eines, der Amts halber zur Bett …]{428. Gebet Eines, der Amts halber zur Bett der Pestilenz nicht fliehen kann}

Hebe dich, Teufel mit dem Schrecken! Weil dichs verdrießt, so will ich dir zu Trutz um desto eher hinzugehen zu meinem kranken Nächsten, ihm zu helfen, und will dich nicht ansehen, und will auf zwei Stücke gegen dich pochen. Das erste ist, daß ich fürwahr weiß, daß dies Werk Gott und allen Engeln wohlgefällt, und wo ichs thue, daß ich in seinem Willen und rechtem Gottesdienst und Gehorsam gehe; sonderlich weil es dir so übel gefällt und du dich so hart darüber setzest, so muß es freilich insonderheit Gott gefallen. Wie willig und fröhlich wollte ichs thun, wenns nur einem
\quellseite{368}
Engel wohlgefiele, der mir zusähe und sich mein darüber freuete; nun es aber meinem HErrn JEsu Christo und dem ganzen himmlischen Heere wohlgefällt, und ist Gottes, meines Vaters, Wille und Gebot: was sollte mich dein Schrecken denn bewegen, daß ich solche Freude im Himmel und Lust meines HErrn sollte hindern, und dir mit deinen Teufeln in der Hölle ein Gelächter und Gespött über mich anrichten und höhnen? Nicht also, du sollst es nicht enden. Hat Christus sein Blut für mich vergossen und sich um meinetwillen in den Tod gegeben, warum sollte ich nicht auch um seinetwillen mich in eine kleine Gefahr geben, und eine ohnmächtige Pestilenz nicht dürfen ansehen? Kannst du schrecken, so kann mein Christus stärken; kannst du tödten, so kann Christus Leben geben; hast du Gift im Maul, Christus hat noch viel mehr Arznei. Sollte mein lieber Christus mit seinem Gebet, mit seiner Wohltat und allem Trost nicht mehr gelten in meinem Geist, denn du leidiger Teufel mit deinem falschen Schrecken in meinem schwachen Fleische? Das wollte Gott nimmermehr! Heb dich, Teufel, hinter mich. Hier ist Christus, und ich sein Diener in diesem Werk. Der soll's walten. Amen.

\section[429. Wenn man nicht gebunden ist, sondern zu …]{429. Wenn man nicht gebunden ist, sondern zu Pest- und Sterbenszeiten weicht und fleucht}

HErr Gott, ich bin schwach und furchtsam; darum fliehe ich das Uebel; und thue so viel dazu, als ich kann, daß ich mich davor hüte; aber ich bin gleichwohl in deiner Hand in diesem und allerlei Uebel, so mir begegnen möge; dein Wille geschehe, denn meine Flucht wirds nicht thun. Denn der Teufel feiert und schläft nicht, welcher ist ein Mörder von Anfang, und sucht allenthalben eitel Mord und Unglück anzurichten.

\section{430. Gebet in Kriegsnoth}

Ach, HErr Zebaoth, wir bekennen vor deinem heiligen Angesichte, daß wir dies gar schreckliche Kriegsfeuer und grausame Landesverwüstung durch unsere vielfältigen, großen und schweren Sünden angezündet und verursacht haben. Denn weil wir deine
\quellseite{369}
mannigfaltigen Wohlthaten vergessen und des lieben edlen Friedens zu allerhand Sünden und Ungerechtigkeit gemißbraucht haben und wacker gewesen sind, Uebels zu thun: so bist du nach deiner Gerechtigkeit auch wacker gewesen mit diesem Unglück über uns her, daß du über uns hast gehen lassen. Aber hilf uns und errette uns, HErr, nach deiner großen Güte, laß des Raubens, Würgens, Sengens und Brennens ein Ende werden. Laß dich das Elend so vieler tausend Seelen, die du dir durch das Blut deines Sohnes zum Eigenthum erkauft hast, jammern und schaffe unsern Grenzen wieder Friede und Ruhe. Du Friedefürst, HErr Jesu Christe, wahrer Mensch und wahrer Gott, der du uns zum Frieden berufen hast, mache fest die Riegel unserer Thore: Und weil du aller Menschen Herzen in deiner Hand hast und leitest sie wie die Wasserbäche: so laß alle Mächtigen auf Erden Friedensgedanken haben, damit nicht durch Krieg und Blutvergießen das Land von Frevel und Ungerechtigkeit angefüllt und der Gottesdienst, Zucht und Ehrbarkeit gehindert werde, sondern wir forthin ein geruhiges und stilles Leben führen und dich für deine Gnade und wunderliche Güte rühmen und preisen mögen, so lange wir hier sind und leben. Amen.

\section{431. Ein Anderes}

Himmlischer Vater, wir haben ja wohl verdienet, daß du uns strafest, aber strafe du uns selbst nach deiner Gnade, und nicht in deinem Grimm. Es ist uns besser, in deiner Hände Staube uns geben, denn in der Menschen oder des Feindes Hände, wie David auch bat (2 Sam. 2.), denn groß ist deine Barmherzigkeit. Wir haben dir gesündiget und deine Gebote nicht gehalten, aber du weißt, allmächtiger Gott Vater, daß wir dem Teufel, Papst und Türken nicht gesündiget haben, sie auch kein Recht noch Macht haben, uns zu strafen, sondern du kannst und magst ihrer brauchen als deiner grimmigen Ruthen wider uns, die wir an dir gesündiget und alles Unglück verdienet haben. Ja, lieber Gott, himmlischer Vater, wir haben keine Sünde wider sie gethan, darum sie Recht hätten, uns zu strafen, sondern viel lieber wollten sie, daß wir sammt ihnen aufs greulichste wider dich sündigten. Denn sie fragen
\quellseite{370}
nicht darnach, ob wir dir ungehorsam wären, dich lästerten, allerlei Abgötterei trieben, wie sie thun, mit falscher Lehre, Glauben und Lügen umgingen, Ehebruch, Unzucht, Mord, Diebstahl, Räuberei, Zauberei und alles Uebel wider dich thäten, da fragen sie nicht nach: sondern das ist unsere Sünde wider sie, daß wir dich, Gott Vater, den rechten einigen Gott, und deinen lieben Sohn, unsern Herrn Jesum Christum, und den Heiligen Geist, einen ewigen Gott predigen, glauben und bekennen, ja das ist die Sünde, die wir wider sie thun. Aber wo wir dich verleugnen, würde uns der Teufel, Welt, Papst und Türke wohl zufrieden lassen, wie dein lieber Sohn spricht: Wäre't ihr von der Welt, so hätte die Welt das ihre lieb. Hier siehe nun drein, du barmherziger Vater über uns, und ernster Richter über unsere Feinde, denn sie sind deine Feinde mehr, denn unsere Feinde, und wenn sie uns verfolgen und schlagen, so verfolgen und schlagen sie dich selber, denn das Wort, so wir predigen, glauben und bekennen, ist dein, nicht unser, alles deines Heiligen Geistes Werk in uns. Der Teufel will solches nicht leiden, sondern an deiner Statt unser Gott sein, an deines Wortes Statt Lügen in uns stiften. Der Türke will seinen Mahomet an deines lieben Sohnes, Jesu Christi, Statt setzen, denn er lästert ihn und spricht, er sei kein rechter Gott, sein Mahomet sei höher und besser, denn er ist. Ist's nun Sünde, daß wir dich den Vater und deinen Sohn und den Heiligen Geist für den rechten einigen Gott halten, bekennen und rühmen, so bist du selbst der Sünder, der du solches in uns wirkest, heißest und haben willst. Darum so hassen, schlagen und strafen sie dich selbst, wenn sie uns um solcher Sache willen hassen, schlagen und strafen. Darum wache auf, lieber Herr Gott, und heilige deinen Namen, den sie schänden; stärke dein Reich, das sie in uns zerstören, und schaffe deinem Willen, den sie in uns dämpfen wollen, und lasse dich nicht um unserer Sünde willen also mit Füßen treten von denen, die nicht unsere Sünde in uns strafen, sondern dein heiliges Wort, Namen und Werk in uns tilgen wollen, daß du kein Gott sein solltest und kein Volk haben, das dich predige, glaube und bekenne. Amen.

\quellseite{371}

\section{432. Gebet um Frieden in Kriegszeiten}

Herr Zebaoth, laß dich doch erbarmen, daß alle anderen gottlosen Könige und Lande sitzen im Frieden, aber allein dein eigen Volk muß Unfrieden haben; so doch billiger wäre, daß jene Unfrieden und dein Volk Frieden hätte. Weil du denn gerecht bist und wahrhaftig in deinen Worten, so wollest du doch dich selbst ansehen und an deine Gerechtigkeit und Wahrheit gedenken, wenn du ja uns nicht willst ansehen, und uns wieder Frieden geben. Amen.

\section[433. Gebet bei allgemeinen Strafen und …]{433. Gebet bei allgemeinen Strafen und Landplagen}

O du gütiger Gott, barmherziger Herr, hast du dich von uns durch deinen grimmigen Zorn um der Sünde willen abgewendet und unser Lebenlang mit Angst und Trübsal, auch mit der Hölle und ewigem Tode gedräuet und aufs Höchste geängstiget, so laß nun ab, lieber Vater, kehre dich wieder zu uns, laß dich erbitten und sei deinen Knechten gnädig, fülle uns frühe mit deiner Gnade, so wollen wir rühmen und fröhlich sein unser Lebenlang. Amen.

\section{434. Ein Anderes}

O du gütiger barmherziger Gott, du hast uns wie das Gras, das frühe blühet, durch deinen Zorn und Grimm lassen welk und dürr werden, unser Leben verkürzet und deine Ungnade über uns ausgeschüttet. Wir bitten dich, lieber barmherziger Gott, kehre dich wieder zu uns, fülle uns frühe mit deiner Gnade, nicht mit einer schlechten, geringen und einzelnen Gnade, dadurch die Haushaltung, das weltliche Regiment, oder Predigtamt, oder Gesundheit des Leibes erhalten werde, sondern gib uns eine reiche, überschwängliche, ewige Gnade, durch welche wir ewig erhalten und vom Teufel, Tod, Sünde und Hölle errettet werden. Amen.

\section{435. Gebet um Geduld und Ueberwindung}

Lieber Gott, verleihe uns auch ein friedlich Herz und guten Muth im Kampf und Unruhe wider den Teufel, auf daß wir nicht allein erdulden und endlich obsiegen, sondern auch mitten im Kampf
\quellseite{372}
und Unruhe Frieden haben mögen, dich loben, und danken und nicht murren, noch ungeduldig werden wider deinen göttlichen Willen, damit der Friede in unserm Herzen den Sieg behalte, daß wir nichts wider dich, unsern Gott, noch Menschen durch Ungeduld vornehmen, sondern beide, inwendig und auswendig, gegen Gott und den Menschen still und friedlich bleiben, bis der endliche und ewige Friede komme. Amen.

\section[436. Die Litanei, oder demüthige Fürbitte für …]{436. Die Litanei, oder demüthige Fürbitte für die allgemeine Noth der ganzen Christenheit}

Wer eigentlich der Autor derselben sei und sie also zusammengetragen, kann man so eben nicht wissen; ist schon zu den Zeiten des alten Lehrers Basilii, ungefähr nach Christi Geburt im 370. Jahr, zu Neo-Cäsarien und mehreren Orten in Griechenland, und zu Ambrosii, Bischofs zu Mailand, Zeiten, im Jahr Christi 380, in Weltschland und andern Orten gemein gewesen. Weil sie aber hernach im Papstthum mit Anrufung der verstorbenen Heiligen verderbet worden, ist der selige Herr Lutherus verursachet worden, neben andern eingerissenen Irrthümern des Papstthums auch die Litanei zu reformiren und nach Inhalt göttlichen Worts also zu stellen, wie sie noch heutiges Tages in allen evangelischen Kirchen gesungen und gebraucht wird. Diese teutsche Litanei ist von dem nunmehr seligen und um die ganze evangelische Kirche hochverdienten Theologen, Herrn Johann Gerhard, der heiligen Schrift Doctor und bei der hochlöblichen fürstlichen sächsischen Universität Jena weitberühmtem Professor, mit schönen andächtigen Herzensseufzern erkläret worden, wie folget:

Kyrie Eleison! HErr, unser Herrscher, * ein HErr aller Herren, † der du bist der rechte Vater über Alles, das Kinder heißt im Himmel und auf Erden, erbarme dich unser. * Ps. 8, 2. * 1 Tim. 6, 15. † Eph. 3, 15.

Christe Eleison! Du Gesalbter des HErrn, * unser König, und † ewiger Hoherpriester, ‡ unser einiger Mittler, Heiland und Erlöser, erbarme dich unser. * Ps. 2, 2. v. 6. † Hebr. 5, 6. ‡ 1 Tim. 2, 5.

Kyrie Eleison! O heiliger Geist, der du mit dem Vater und Sohn bist * ein HErr und † Gott, ‡ unser Tröster, erbarme dich unser. * 1 Cor. 12, 5. † v. 6. ‡ Joh. 15, 26.

Christe, erhöre uns! O HErr Jesu, unser * einiger Fürsprecher bei dem Vater, † der du bist zur Rechten Gottes und vertrittest uns, bitte für uns und erhöre uns. * 1 Joh. 2, 2. † Röm. 8, 34.

\quellseite{373}

HErr Gott Vater im Himmel, erbarme dich über uns! Der du uns erschaffen und * uns also geliebet, daß du uns deinen eingebornen Sohn gegeben, † o HErr, segne und behüte uns.

HErr Gott Sohn, der Welt Heiland, erbarme dich über uns! Der du uns also geliebet hast, daß du dich selbst für uns dargegeben, uns zu erlösen, * o HErr, laß dein Angesicht leuchten über uns und sei uns gnädig.

HErr Gott Heiliger Geist, erbarme dich über uns! Der du uns wiedergebierest, erneuerst und * heiligest, der du uns † vertrittst mit unaussprechlichem Seufzen, † o HErr, hebe dein Angesicht über uns und gib uns deinen Frieden.

Sei uns gnädig, verschone unser, lieber HErr Gott! Ach HErr, * handele nicht mit uns nach unsern Sünden und vergilt uns nicht nach unserer Missethat, † strafe uns nicht in deinem Zorn und züchtige uns nicht in deinem Grimm.

Sei uns gnädig, hilf uns, lieber HErr Gott! Hilf du uns, Gott, unser Helfer, um deines Namens Ehre willen, errette uns und vergib uns unsere Sünde um deines Namens willen.

Vor allen Sünden behüte uns, lieber HErr Gott! Behüte uns, daß wir nicht wider dich sündigen. * Ach HErr, lehre uns thun nach deinem Wohlgefallen, und dein guter Geist führe uns auf ebener Bahn.

Vor allem Irrsal behüte uns, lieber HErr Gott! Wende unsere Augen ab, daß sie nicht sehen nach unnützer Lehre; laß dein Wort sein * unseres Fußes Leuchte und ein Licht auf unserm Wegen.

Vor allem Uebel behüte uns, lieber HErre Gott! Ach, himmlischer Vater, wie du uns durch Christum vom Sündenübel erlöset hast: * also erlöse uns noch von allerlei Strafübel † und hilf uns aus zu deinem himmlischen Reich.

Vor des Teufels Trug und List behüte uns, lieber HErre Gott! Rüste uns aus mit der Macht deiner Stärke, * daß wir bestehen mögen gegen die listigen Anläufe des Teufels.

\quellseite{374}

Vor bösem, schnellem Tod behüte uns, lieber HErr Gott! Ach HErr, * raffe unsere Seelen nicht hin mit den Sündern; noch unser Leben mit den Blutdürstigen. † Lehre uns bedenken, daß wir sterben müssen, und hilf, daß wir nicht unversehens mit unserm letzten Stündlein übereilet werden. * Ps. 26, 9. † Ps. 90, 12.

Vor Pestilenz und theurer Zeit behüte uns, lieber HErr Gott! Laß uns unter deinem Schirm und deinem Schatten bleiben, * daß wir bewahret werden vor der schädlichen Pestilenz; † nimm auch nicht hinweg den Vorrath des Brodes, † laß keine Theuerung über uns kommen. Ps. 91, 1. * v. 8. † Ezech. 14, 13. ‡ c. 36, 29.

Vor Krieg und Blutvergießen behüte uns, lieber HErr Gott! Zerstreue die Völker, die gern kriegen; * steuere den Kriegen in aller Welt; † mache fest die Riegel unsrer Thore und schaffe unsern Grenzen Friede. Ps. 68, 31. * Ps. 46, 10. † Ps. 147, 13.

Vor Aufruhr und Zwietracht behüte uns, lieber HErr Gott! O du * HErr des Friedens, regiere alle Unterthanen, daß sie † nicht Aufruhr im Lande anrichten, sondern friedlich und ‡ einträchtig beisammen wohnen. * 2 Thess. 3, 16. † Sir. 7, 7. ‡ Ps. 133, 2.

Vor Hagel und Ungewitter behüte uns, lieber HErr Gott! Wende ab schädliche Hagel im Regen und schreckliche Ungewitter, * beschirme uns zur Zeit des Donners, † mit dem Schatten deiner Flügel. Ps. 105, 32. * Jer. 30, 23. † Ps. 17, 8.

Vor dem ewigen Tode behüte uns, lieber HErr Gott! Erlöse uns aus der Hölle und bewahre uns vor dem ewigen Tode; * erhalte uns aus deiner Macht durch den Glauben zur Seligkeit. Hos. 13, 14. * 1 Pet. 1, 5.

Durch deine heilige Geburt hilf uns, lieber HErr Gott! Unsere * sündliche Empfängniß und unreine Geburt wollest du, o HErr Christe! durch deine allerdings reine Empfängniß und Geburt heiligen. Ps. 51, 7.

Durch deinen Todeskampf und blutigen Schweiß hilf uns, lieber HErr Gott! Du hast * die Kelter des göttlichen Zorns wider unsere Sünde allein getreten und darüber Blut geschwitzet, um deß willen wollest du uns von dem † zukünftigen Zorn erlösen. * Jes. 63, c. † Offb. 19, 15. ‡ 1 Thess. 1, 10.

\quellseite{375}

Durch dein Kreuz und Tod hilf uns, lieber HErr Gott! Du hast unsere Sünde selbst geopfert an deinem Leibe auf dem Holz, laß uns unserer Sünde halben nimmermehr verzagen; * du hast durch deinen Tod dem Tode die Macht genommen, ach, stehe uns bei in Noth und Tod. 1 Pet. 2, 24. *2 Tim. 1, 10.

Durch dein heilig Auferstehen und Himmelfahrt hilf uns, lieber HErr Gott! Hilf, * daß wir der Sünde täglich absterben und in einem neuen Leben wandeln, † daß wir trachten nach dem, das droben ist, und nicht nach dem, das auf Erden ist. *Röm. 6, 2. †Col. 3, 1.

In unserer letzten Noth hilf uns, lieber HErr Gott! Sei unser Licht und Führer, * wenn wir durchs finstere Thal des Todes wandern sollen; laß uns in der Finsterniß des Todes † deine Herrlichkeit sehen und ‡ nimm unsere Seele auf in deine Hände. *Ps. 23, 4. †Apost. 7, 55. ‡Ps. 31, 6. Apok. 7, 59.

Am jüngsten Gericht hilf uns, lieber HErr Gott! (Sib, o HErr, * daß wir Barmherzigkeit finden bei dem HErrn am letzten Gerichtstage; laß uns hören die fröhliche Stimme: †Kommt her, ihr Gesegneten, ererbet das Reich, so euch bereitet ist von Anbeginn der Welt. *2 Tim. 1, 17. †Matth. 5, 34.

Wir arme Sünder bitten, du wollest um unserer Sünde willen * deine Barmherzigkeit nicht vor Zorn verschließen; †laß unser Gebet vor dich kommen und neige deine Ohren zu unserm Geschrei. *Ps. 77, 10. †Ps. 88, 3.

Du wollest uns erhören, lieber HErr Gott! Und deine heilige christliche Kirche regieren und führen. Hilf deinem Volk, o HErr, und segne dein Erbe, und * weide sie, und erhöhe sie ewiglich; regiere und führe deine heilige Gemeine, †welche dein Sohn durch sein eigen Blut erworben hat; ‡ suche heim deinen Weinstock, den deine Rechte gepflanzet hat. *Ps. 28, 8. †Ps. 48, 15. ‡Apost. 20, 28. **Ps. 80, 15.

Erhöre uns, lieber HErr Gott! Alle Bischöfe, Pfarrherren und Kirchendiener in heilsamem Wort und heiligem Leben erhalten. Hilf, o HErr, daß Lehrer und Prediger *das Geheimniß des Glaubens in reinem Gewissen
\quellseite{376}
bewahren, † daß sie halten ob dem Wort, das gewiß ist, und ‡ die Lehre zieren in allen Stücken. * 1 Tim. 3, 9. † Tit. 1, 9. ‡ Tit. 2, 10.

Erhöre uns, lieber HErr Gott!

Allen Rotten und Aergernißen wehren. Gib Gnade, * daß wir fleißig seien zu halten die Einigkeit im Geist durch das Band des Friedens. † Steuere allen Irrgeistern, welche Zertrennung und Aergerniße anrichten. * Eph. 4, 3. † Mich. 2, 11.

Erhöre uns, lieber HErr Gott!

Alle Irrigen und Verführten wiederbringen. Du treuer * Hirte unserer Seele, wollest † die verlorenen und verirrten Schäflein suchen und wiederbringen, ‡ auch ihre Füße richten auf den Weg des Friedens. * Ps. 23, 1. † Ezech. 34, 16. ‡ Luc. 1, 79.

Erhöre uns, lieber HErr Gott!

Den Satan unter unsere Füße treten. O HErr, mache uns theilhaftig * des Sieges, welchen dein lieber Sohn wider den Satan erhalten, daß wir in seiner Kraft † treten auf den häßlichen Löwen und Drachen; ‡ zertritt den Satan unter unsere Füße in Kurzem. * Matth. 4, 10. † Ps. 91, 13. ‡ Röm. 16, 20.

Erhöre uns, lieber HErr Gott!

Treue Arbeiter in deine Ernte senden. Gib treue * Hirten und Lehrer, gib reichlich dein heiliges Wort, † und mit großen Schaaren Evangelisten. ‡ Gib uns Lehrer nach deinem Herzen, die uns weiden mit Lehre und Weisheit. * Eph. 4, 11. † Ps. 68, 12. ‡ Jer. 3, 15.

Erhöre uns, lieber HErr Gott!

Deinen Geist und Kraft zum Wort geben. Gib * zu dem Pflanzen und Begießen der Lehrer das Gedeihen. Befeuchte † den Samen deines Worts ‡ mit dem Regen deines göttlichen Segens. ** Gib deinem Donner Kraft, daß dein Wort †† wachse und ‡‡ laufe. * 1 Cor. 3, 6. † Luc. 8, 11. ‡ Jes. 55, 10. ** Ps. 68, 34. †† Apost. 12, 24. ‡‡ 2 Thess. 3, 1.

Erhöre uns, lieber HErr Gott!

Allen Betrübten und Blöden helfen und trösten. O Vater der Barmherzigkeit und Gott alles Trostes, * erquicke die müden Seelen, gib den Müden Kraft, und Stärke genug den Unvermögenden,
\quellseite{377}
† requide den Geist der Gedemüthigten und das Herz der Zerschlagenen. * 2 Cor. 1, 3. † Jes. 40, 29. ‡ Jes. 57, 15.

Erhöre uns, lieber HErr Gott!

Allen Königen und Fürsten Friede und Eintracht geben. Wir bitten dich * für alle Obrigkeit, daß wir unter ihnen ein geruhliches und stilles Leben führen mögen in aller Gottseligkeit und Ehrbarkeit; † leite die Herzen aller hohen Potentaten zum Frieden, damit ‡ Friede und Treue sei, weil wir leben. * 1 Tim. 2, 2. † Ps. 33, 15. ‡ Jes. 39, 8.

Erhöre uns, lieber HErr Gott!

Unserm Kaiser stetig Sieg wider seine Feinde gönnen. Ach Gott, * es ist ja des Königs Herz in deiner Hand, wie die Wasserbäche. Gib unserm Kaiser ein gütig Herz gegen dein Wort und Kirche. Steuere allen denen, so ihre Hände an deine Friedensamen legen und deinen Bund zu entheiligen begehren. † Dämpfe die Feinde deiner Kirche und wende deine Hand über ihre Widerwärtigen. * Spr. 21, 1. † Ps. 55, 21. ‡ Ps. 81, 15.

Erhöre uns, lieber HErr Gott!

(Unsern gnädigen Landesherrn mit allen seinen Gewaltigen leiten und schützen.) (Gib Gnade, daß unser gnädiger Landesherr sei: * ein Pfleger und Säugamme deiner Kirche; † sei du, HErr, seines Gefallten Stärke, und erfülle alle seine christlichen Anschläge; ‡ überschütte ihn mit gutem Segen, ** und durch deine Güte laß ihn fest bleiben.) * Jes. 49, 23. † Ps. 28, 8. ‡ Ps. 20, 5. ** Ps. 21, 4. 8.

Erhöre uns, lieber HErr Gott!

Unsern Rath und Gemeine segnen und behüten. Gib allen Rathsherren * den Geist des Raths und der Weisheit, † daß sie nach deinem Wort das Recht setzen; ‡ segne unsere ganze Gemeine, segne Alle, die dich fürchten, beide Kleine und Große; ** segne uns je mehr und mehr, uns und unsere Kinder. * Jes. 11, 2. † Spr. 8, 15. ‡ Ps. 115, 3. ** v. 4.

Erhöre uns, lieber HErr Gott!

Allen, so in Noth und Gefahr sind, mit Hülfe erscheinen. Sei du, o HErr, * der Armen Schutz, ein Schutz in der Noth,
\quellseite{378}
† errette die Elenden aus ihrem Elend † und reiße sie aus dem weiten Rachen der Angst, die keinen Boden hat; ** hilf ihnen aus allen ihren Nöthen.\footnote{* Ps. 9, 11. † Hiob 36, 15. ‡ v. 16. ** Ps. 34, 7.}

Erhöre uns, lieber Herre Gott!

Allen Schwangeren und Säugern fröhliche Frucht und Gedeihen geben. Ach, Herr * Gott über alles lebendige Fleisch, der du die Menschenkinder im Mutterleibe † wunderbarlich bildest, stärke alle Schwangeren, behüte sie vor allem Unfall, und † wenn die Kinder bis an die Geburt kommen, so gib Kraft zu gebären.\footnote{* 4 Mos. 27, 16. † Ps. 139, 14. ‡ Jes. 37, 3.}

Erhöre uns, lieber Herre Gott!

Aller Kinder und Kranken pflegen und warten. Segne, Herr! alle Kinder, * daß sie an Weisheit, Alter und Gnade zunehmen. Du, Herr, † unser Arzt, wollest allen Kranken, nach deinem gnädigen Willen, ‡ von ihrer Krankheit helfen, auch Gnade geben, ** daß sie in ihrer Krankheit dich, den Herrn, suchen.\footnote{Ps. 147, 13. * Luc. 2, 52. † 2 Mos. 15, 26. ‡ Ps. 41, 4. ** 2 Chron. 16, 12.}

Erhöre uns, lieber Herre Gott!

Alle Gefangenen los und ledig lassen. Hilf, o Herr! aus ihren Aengsten allen denen, * welche sitzen müssen im Finsterniß und Dunkel, gefangen im Zwang und Eisen; † führe sie aus dem Finsterniß und Dunkel, ‡ verachte nicht deine Gefangenen, sondern ** führe sie aus zu rechter Zeit.\footnote{Ps. 107, 10. * v. 13. † v. 14. ‡ Ps. 69, 85. ** Ps. 68, 7.}

Erhöre uns, lieber Herre Gott!

Alle Wittwen und Waisen vertheidigen und versorgen. Du bist ja * ein Vater der Waisen und ein Richter der Wittwen; † behüte alle Waisen und Wittwen, ‡ laß die Thränen der Wittwen vor deinen Thron kommen und ** schaffe Recht den Waisen.\footnote{* Ps. 68, 6. † Ps. 146, 9. ‡ Sir. 35, 21. ** Ps. 10, 8.}

Erhöre uns, lieber Herre Gott!

Aller Menschen dich erbarmen. Du, Herr, bist ja * ein Heiland aller Menschen, † dein unvergänglicher Geist ist in allen, und dein lieber Sohn hat sich ‡ für alle Menschen dahingegeben, sie zu erlösen; darum wollest du auch dich aller erbarmen.\footnote{* 1 Tim. 4, 10. † Weish. 12, 1. ‡ 1 Tim. 2, 6.}
\quellseite{379}
Erhöre uns, lieber Herre Gott!

Unsern Feinden, Verfolgern und Lästerern vergeben und sie belehren. Vergib, Herr, unsern Feinden und Verfolgern, * denn sie wissen nicht, was sie thun; † behalte ihnen diese Sünde nicht, daß sie uns ‡ ohne Ursache hassen und verfolgen.\footnote{* Luc. 23, 34. † Apost. 7, 60. ‡ Ps. 119, 161.}

Erhöre uns, lieber Herre Gott!

Die Früchte auf dem Lande geben und bewahren. Suche, Herr, unser Land heim, und wässere es, und mache es sehr reich; * laß unser Getreide wohl gerathen, † mache das Land voll Früchte, die du schaffest, ‡ gib uns Frühregen und Spätregen zu rechter Zeit.\footnote{* Ps. 65, 10. † Ps. 104, 13. ‡ Jer. 5, 24.}

Erhöre uns, lieber Herre Gott!

Und uns gnädiglich erhören. Neige, * Herr, deine Ohren zu unserm Gebet, † höre das Verlangen der Elenden; ‡ ach Herr, höre, ach Herr, sei gnädig, ach Herr, merke auf und thue es, und verzeuch nicht um dein selbst willen.\footnote{* Ps. 17, 6. † Ps. 10, 17. ‡ Dan. 9, 19.}

Erhöre uns, lieber Herre Gott!

O Jesu Christe, Gottes Sohn, o Herr Jesu, unser Heiland, * hilf uns um deines Namens willen, † du Gesalbter des Herrn, mache uns theilhaftig deiner geistlichen ‡ Salbung und des Freudenöls des Heiligen Geistes**, du eingeborner Sohn Gottes††, mache uns zu Kindern deines himmlischen Vaters.\footnote{* 1 Chron. 17, 35. † Joh. 1, 41. ‡ 1 Joh. 2, 27. ** Ps. 45, 8. †† Joh. 1, 12.}

Erhöre uns, lieber Herre Gott!

O du Gotteslamm, das der Welt Sünde trägt, o du *unschuldiges und unbeflecktes Lämmlein Gottes, durch dessen theures Blut wir von unsern Sünden gereiniget sind.\footnote{* 1 Pet. 1, 19.}

Erbarme dich über uns!

O du Gotteslamm, das der Welt Sünde trägt, da du um unsertwillen gemartert, * hast du deinen Mund nicht aufgethan, wie ein Lamm, das zur Schlachtbank geführet, und wie ein Schaf, das verstummet vor seinem Scherer und seinen Mund nicht aufthut.\footnote{* Jes. 53, 7}
\quellseite{380}
Erbarme dich über uns! O du Gotteslamm, das der Welt Sünde trägt, du bist für uns ein Fluch worden, * daß der Segen über uns käme; † die Strafe unserer Sünde ist auf dich geleget, auf daß wir Friede hätten. * Gal. 3, 13. † Jes. 53, 5.

Verleihe uns steten Frieden! Christe, erhöre uns! JEsu, unser einiger * Fürsprecher bei dem Vater, † der du bist zur Rechten Gottes und vertrittst uns. *1 Joh. 2, 2. † Röm. 8, 34.

Kyrie Eleison! HErr, unser Herrscher, * ein HErr aller Herren, † der du bist der rechte Vater über Alles, das Kinder heißt im Himmel und auf Erden, erbarme dich unser! Pf. 8, 2. *1 Tim. 6, 15. † Eph. 3, 15.

Christe Eleison! Du Gesalbter des HErrn, * unser König und † ewiger Hoherpriester, unser ‡ einiger Mittler und einiger Heiland, erbarme dich unser! * Ps. 2, 2. † v. 6. ‡ Heb. 5, 6.

Kyrie Eleison! Gott Heiliger Geist, der du mit dem Vater und Sohn bist * ein HErr und † ein Gott, unser ‡ Tröster, erbarme dich unser! Amen. *2 Tim. 2, 5. †1 Cor. 12, 6. ‡ Joh. 15, 26.

\chapter{XII. Dank- und Lobgebete}

\section{437. Gebet um christliche Dankbarkeit}

Ach du gnädiger, gütiger, übermilder Gott und Vater! Wie groß ist deine Gnade! Wie gütig ist dein mildes Vaterherz! Wie groß sind deine Wohlthaten, Liebe, Güte und Barmherzigkeit gegen uns! Ich klage und bekenne dir, daß mein Herz von Natur so unverständig, so viehisch und grob ist, daß ich nie von Herzen erkannt habe deine Wohlthaten der Schöpfung und Erhaltung, der Erlösung und Heiligung. Ich habe dir nie von Herzen dafür gedanket, dir nie deine gebührende Ehre dafür gegeben. Ach ich erkenne und bekenne, daß ich viel zu gering bin aller deiner Barmherzigkeit, die du an mir von Mutterleibe gethan hast. Ich erkenne, daß ich
\quellseite{381}
deiner geringsten Wohlthat nicht werth bin, sondern vielmehr schuldig und werth bin deines Zorns und Ungnade. Dennoch hast du mir Unwürdigem so große Barmherzigkeit erzeiget, aus lauter Gnade und Güte. Ich habe es nicht verdienet, kanns auch nicht verdienen, und werde es auch in Ewigkeit nicht verdienen können. Denn deine Gnade ist's, was ich um und um bin. Du hättest wohl Macht, wegen meiner Undankbarkeit, alle deine Gaben, leibliche und geistliche, wieder zu nehmen, denn es ist dein Gut! Ach, vergib mir solche große Undankbarkeit, und wende von mir die Strafe, die du dräuest, daß das Böse vom Hause des Undankbaren nicht lassen soll, und gib mir ein verständiges, dankbares Herz, daß ich erkennen möge, du seiest der Brunnen und Ursprung aller guten Gaben, und daß ich ohne dich nichts bin, denn ein todter, lebloser Schatten in allem meinem Thun. Ach, wie herzlich haben dir alle deine Heiligen gedanket, und gesaget: Das ist ein köstlich Ding, dem HErrn danken, und deinen Namen loben, du Allerhöchster. Item: So will ich dir ein Freudenopfer thun, und deinem Namen danken, daß er so tröstlich ist. Lobe den HErrn, meine Seele, und vergiß nicht, was er dir Gutes gethan hat. Gib, daß ich alles mit dankbarem Herzen von deiner Hand empfange, und deine Gaben aus den Schätzen deiner Gnade und Barmherzigkeit aufnehme, daß ich auch erkennen lerne, daß du deine Güter und Gaben mir allein erhältst, und ich nicht selbst, daß ich dich auch darum bitte, liebe, ehre und preise, daß ich auch alles, was du mir gibst, zu deines Namens Lob und Ehre gebrauche und anlege, und nicht mir, sondern dir die Ehre in allen Dingen gebe, denn dies ist die Dankbarkeit und Gerechtigkeit, die ich dir schuldig bin, weil alles dein ist, und das ist die Wahrheit, daß ich dasselbe erkenne und preise, und so du etwas Gutes durch mich wirkest, daß ichs dir und nicht mir zuschreibe, sondern wenn ich alles gethan habe, spreche: Ich bin ein unnützer Knecht, ein unwürdiges Werkzeug deiner Gnade, nicht ich, sondern deine Gnade, die in mir ist, thut alles durch mich. Gib mir auch, daß ich dankbar bin gegen die, durch welche du mir Gutes thust, und um deinetwillen dieselben liebe und ehre, und ihnen wieder durch deine Gnade diene, und von deinen Gütern Gutes thue, und für sie bitte! Ja, daß ich auch um deinetwillen
\quellseite{382}
meine Feinde liebe, und ihnen Gutes thue. Laß die Undankbarkeit, welche das allerschändlichste Laster ist, in mir nicht einwurzeln, daß mich nicht der Fluch treffe, sondern laß diese edle Tugend welche eine Mutter ist vieles Segens, allezeit bei mir bleiben, daß ich mit fröhlichem Herzen und Gewissen dich, mit allen heiligen Engeln, in Ewigkeit für alle deine Wohlthaten lobe und preise, durch Christum unsern Herrn. Amen.

\section{438. Danksagung für die Liebe Gottes}

Ach du barmherziger, gnädiger, leutseliger und freundlicher Gott, du Liebhaber der Menschen! Ich danke dir von Grund meines Herzens für deine große, unaussprechliche, väterliche, herzliche Liebe, damit du, liebreicher Gott und Vater, mich allezeit geliebet hast. Du hast mich in deiner Liebe als in deinem Herzen getragen. Du denkst allezeit an mich und vergissest mein nicht. Du sorgest für mich, hörest mein Gebet, zählest meine Thränen, zählest meine Seufzer, weißest alle meine Trübsal, erkennest mein Herz, und hast mir so theure werth(e) Verheißung deiner Gnade, Hülfe, Trostes, Vergebung der Sünden und des ewigen Lebens gegeben und geoffenbaret. Du hast mir deinen lieben Sohn geschenket und mich durch seine Menschwerdung lassen versichern deiner Huld und Gnade, durch sein heiliges Evangelium mich lassen lehren und trösten, durch sein Exempel mir einen Weg und Richtschnur eines heiligen Lebens gezeiget, durch sein Leiden und Sterben mich vom ewigen Tode erlöset, durch seine Auferstehung aus dem Tode gerissen, durch seine Himmelfahrt mir das Paradies eröffnet und eine Stätte im Himmel bereitet. Du hast mich durch deinen Heiligen Geist erleuchtet, geheiliget, getröstet, gestärket, gelehret und erquickt und mich durch ihn der Kindschaft Gottes und ewigen Erbes versichert. Du hast mir nicht allein große Wohlthaten erzeiget, sondern dich mir selbst gegeben mit deinem lieben Sohn und Heiligen Geist. Wie kann ich dir für diese große Liebe genugsam danken? Sie ist größer denn Himmel und Erde, sie ist ewig und höret nimmermehr auf, sie ist unendlich, höher denn der Himmel, breiter denn die Erde, tiefer denn der Abgrund, länger denn der Aufgang vom Niedergang. So hoch der Himmel über der Erde ist, lässest du deine Gnade walten über alle,
\quellseite{383}
die dich fürchten. So fern der Abend ist vom Morgen, lässest du unsere Uebertretung von uns sein und erbarmest dich über uns, wie sich ein Vater über seine Kinder erbarmet. Ach, gib mir ein solch Herz, daß ich deiner Liebe nimmer vergesse, darin gehe, liege, schlafe und wache, stehe und wandele, lebe und sterbe, und in deiner Liebe wieder aufstehe, und darin ewig lebe. Denn deine Liebe ist das ewige Leben und die ewige Ruhe. Und das ewige Leben ist nichts anderes, denn deine ewige Liebe. Ach, laß dieselbe nimmermehr in meinem Herzen verlöschen, sondern immer größer werden, daß in meiner Seele nichts mehr sei, leuchte, lebe, scheine, denn deine Liebe; daß ich vor Liebe meine Thränen vergieße, wie Maria Magdalena, und dich mit Freuden aufnehme in das Haus meines Herzens, wie Zachäus. Ach, du hast niemand jemals verschmähet, du Liebhaber der Menschen, verschmähe mich auch nicht. Ich bin zwar nicht werth, daß ich dich liebe oder von dir geliebet werde, ich elender Wurm, stinkender Koth, todter Hund! Mache du mich aber durch deine Liebe würdig. Neige mein Herz von aller Creaturliebe zu dir, auf daß du meine Liebe allein seiest, meine Freude, meine Hoffnung, meine Stärke, mein Licht, mein Heil, mein Leben, mein Arzt, meine Speise, mein Traut, mein Hirte, mein Schutz, meine Kraft, Schatz und ewiger Reichthum. O selig bin ich, wenn ich in deiner Liebe bleibe, damit du mich liebest! O heilig bin ich, wenn ich dich allezeit liebe! O unselig bin ich, wenn ich aus deiner Liebe falle: so verliere ich mein höchstes Gut und komme von dem Licht in die Finsterniß, vom Leben in den ewigen Tod! O du reine, göttliche Liebe, entzünde mein Herz, erleuchte meinen Verstand, heilige meinen Willen, erfreue mein Gedächtniß und vereinige mich mit dir ewiglich. Amen.

\section[439. Danksagung für die ewige Gnadenwahl in …]{439. Danksagung für die ewige Gnadenwahl in Christo}

Ach du barmherziger, gnädiger, gütiger Gott! Wie hast du den Reichthum deiner Gnade in deinem lieben Sohn, Jesu Christo, so mildiglich, so reichlich, so überflüssig über uns arme, unwürdige, verlorne und verdammte Menschen ausgetheilt und bist alle unsern Verdienst zuvorkommen. Du hast uns geliebet in deinem Geliebten, ehe der Welt Grund geleget ist. Deine Barmherzigkeit gehet über
\quellseite{384}
alle Menschen. Du willst nicht, daß jemand verloren werde, sondern daß sich jedermann zur Buße kehre und lebe. Du willst, daß allen Menschen geholfen werde und sie zur Erkenntniß der Wahrheit kommen. Darum hast du deinen lieben Sohn der Welt geschenket und geoffenbaret und in deiner ewigen Gnadenwahl verordnet, alle die selig zu machen, die an ihn glauben. Derhalben hast du ihn für die Sünde der Welt dahingegeben und auf ihn, als das unschuldige Lämmlein Gottes und einiges Versöhnopfer, alle unsere Sünde geleget, dieselbe zu tragen, hinwegzunehmen und die Welt zu versöhnen. Darum ist er eine Versöhnung für der ganzen Welt Sünde. Er ist kommen zu suchen und selig zu machen, das verloren ist. Nun waren alle Menschen verloren. Darum ist er kommen, alle Menschen selig zu machen. Derowegen hast du, getreuer Gott, das Evangelium allen Völkern verkündigen lassen, und dein lieber Sohn hat alle zu sich gerufen, die mühselig und beladen seyn, daß er sie erquicke. Du hast auch zugesaget, durch dein Wort in uns den Glauben zu erwecken. Denn dein Wort soll nicht leer wieder zu dir kommen, sondern alles ausrichten, dazu du es gesandt hast; und hast's gemacht zur Kraft Gottes, gerecht und selig zu machen, die daran glauben; hast auch zugesagt, die Gläubigen zu heiligen in der Liebe, auch dieselbigen zu erhalten im Glauben bis ans Ende, also, daß deine Schäflein niemand soll aus deiner Hand reißen, und daß du das gute Werk, so du in uns hast angefangen, vollführen wirst bis auf den Tag Jesu Christi; hast auch deinen Gläubigen gnädigen Schutz verheißen wider den Teufel und die Welt, und so tröstliche Siegel deiner Gnade an dein Wort gehänget, die heilige Taufe und Nachtmahl, daß wir an deiner Gnade nicht zweifeln sollen. Für die große Liebe, allgemeine Erlösung und Bezahlung unserer Sünde, allgemeinen Beruf, und für die Kraft deines göttlichen Worts und den Heiligen Geist, der in uns den Glauben erwecket, weil der Glaube dein Werk ist; auch für die Erhaltung im Glauben, für die tröstliche Verheißung deiner ewigen Gnade, die nicht von uns weichen soll, ob gleich Berge und Hügel hinfallen; auch für die tröstlichen Siegel der heiligen hochwürdigen Sacramente danke ich dir von Herzen. Und bitte, du wollest mich dabei erhalten, mich vor schweren Anfechtungen bewahren
\quellseite{385}
und mein Herz sich gründen lassen auf deine unaussprechliche allgemeine Liebe, auf das allgemeine Verdienst meines Erlösers, auf die allgemeinen Verheißungen, Beruf und Trost der hochwürdigen Sacramente; welcher Wohlthaten du mich aus Gnaden auch hast theilhaftig gemacht, daran ich erkenne, daß ich dein auserwähltes Kind bin, in Christo vor der Welt aus lauter Gnade erwählet und zur Kindschaft verordnet durch meinen HErrn JEsum Christum, in welchem du mich geliebet und mich mit deinem Heiligen Geist versiegelt hast. Gib auch, daß ich heilig sei und unsträflich in der Liebe, zu Lob deiner herrlichen Gnade; und versichere mein Herz nach deiner Verheißung, daß ich gewiß sei, daß mich weder Tod noch Leben noch einige Creatur scheiden kann von deiner Liebe in Christo JEsu, meinem HErrn. Amen.

\section{440. Danksagung für die heil. christliche Kirche}

Ach du gütiger, gnädiger, barmherziger Vater, du Liebhaber der Menschen! Ich danke dir, daß du von Anfang der Welt aus dem menschlichen Geschlecht dir allezeit eine heilige christliche Kirche gesammlet und berufen hast, der du dich in deinem heiligen Wort offenbaret hast, welche dich, den Vater und den Sohn und den Heiligen Geist recht erkannt, an dich gegläubet, dich angerufen, gehret, gelobet und gepreiset hat, und daß du mit derselben einen ewigen Gnadenbund gemacht hast, daß du ihr gnädiger Gott sein wollest, und ihre Sünde vergeben, um des theuren Verdienstes Christi, deines lieben Sohnes, willen, der du auch deinen Heiligen Geist verheißen, der sie erleuchte, lehre, stärke, tröste, heilige, in alle Wahrheit leite und erhalte, der sie erneuere, daß sie dir im heiligen Leben, in Weisheit und Gerechtigkeit diene, die du auch verheißen hast aus Gnaden selig und heilig zu machen, zu welcher du gesagt hast durch deinen lieben Sohn: fürchte dich nicht, du kleine Heerde, es ist des Vaters Wille, dir das Reich zu bescheiden. Welcher du viel herrliche, gnadenreiche Verheißungen gegeben der Vereinigung mit dir: Ich will mich mit dir verloben in Ewigkeit, ich will mich mit dir vertrauen in Gericht und Gerechtigkeit, ja im Glauben will ich mich mit dir verloben, und du sollst den HErrn erkennen. Item: Auf diesen Fels will ich bauen meine Gemeine,
\quellseite{386}
und die Pforten der Hölle sollen sie nicht überwältigen; welche du dir gereiniget hast durchs Wasserbad im Wort, und dir zugerichtet eine Gemeine, die herrlich sey, die nicht habe einen Flecken oder Runzel oder der etwas, sondern daß sie heilig sey und unsträflich; welche du berufen hast, und sie verordnet, daß sie ähnlich soll sein dem Ebenbilde deines lieben Sohnes, welche du auch hast gerecht gemacht und auch herrlich machen wirst. Welche du auch mit so vielen holdseligen Namen nennest, daß sie sei eine Braut deines lieben Sohnes Jesu Christi, und er ihr ewiger Bräutigam; sie sein geistlicher Leib, und er ihr einiges, ewiges Haupt, von welchem sie alle Gaben und Fülle empfähet, Licht, Leben, Trost, Stärke, Kraft, Sieg. Sein königlich Priesterthum, ein Volk des Eigenthums, das auserwählte Geschlecht, sein schönes Erbtheil, und seine Liebliche, auf welche ihm das Loos gefallen, in welcher er allein König und Hoherpriester ist: sein Schafstall, da er allein Hirte ist, eine schöne Stadt Gottes und himmlisches Jerusalem, ein Pfeiler und Grundfeste der Wahrheit, auf den Eckstein Jesum Christum gegründet; ein Weinberg, darin Christus der lebendige Weinstock, und wir seine Reben sein, welchen du deine Gegenwart verheißen hast bis ans Ende der Welt. Dafür dank ich dir von Herzen, daß du mich auch zur Gemeinschaft dieser deiner heiligen Kirche berufen hast, daß ich derselben Glied und Bürger bin des himmlischen Jerusalems, in der Gemeinschaft vieler tausend Heiligen, in welcher ich alle himmlische Güter und Gnadenschätze in Christo mit allen Heiligen gemein habe, nemlich einen Erlöser und Seligmacher, ein Haupt, einen Hirten, eine Taufe, ein Nachtmahl, einen Glauben, einen Gott und Vater unser aller, der da ist in uns allen, durch uns alle, und über uns alle, da alle Glieder der Gläubigen wahrhaftigen Trost, Licht, Leben, Erquickung, Weide und Schutz haben, und wenn gleich die Welt unterginge, und die Berge mitten ins Meer sänken, dennoch soll die Stadt Gottes sein lustig bleiben mit ihren Brünnlein, da die heiligen Wohnungen des Höchsten sind, Gott ist bei ihr drinnen, und hilft ihr frühe, darum wird sie wohl bleiben. Hilf uns nun, und stärke uns, lieber Gott, daß wir in Kreuz und Verfolgung uns deiner Gegenwart trösten, deinen Namen freudig bekennen, für deine Ehre ritterlich streiten,
\quellseite{387}
deinen Namen und Wort, so es sein soll, mit unserm Blute bezeugen, und daß wir auch endlich aus dieser streitenden Kirche aufgenommen werden in die siegende Kirche, da aus Engeln und Menschen eine Kirche wird, die dich in Ewigkeit wird loben und preisen. O selig sind, die in deinem Hause wohnen, immer und ewiglich. Amen.

\section[441. Danksagung für den Schutz der heiligen …]{441. Danksagung für den Schutz der heiligen Engel}

Ach du getreuer Gott, du Liebhaber, Beschirmer und Erhalter des menschlichen Geschlechts, du HErr der Heerschaaren, vor welchem stehen tausendmal tausend, und zehenmal hundert tausend dienen dir, durch welchen alles erschaffen, beides das Sichtbare und das Unsichtbare, die Thronen und die Herrschaften, wie hast du die Menschen so lieb, daß du ihnen von Kindheit auf zugeordnet hast deine heiligen Engel, die beständigen, demüthigen, freundlichen, Gott lobenden Engel, die gehorsamen, keuschen, lieblichen, wahrhaftigen, guten Geisterlein, die schönen glänzenden himmlischen Feuer-Flämmlein, die starken Helden, das Heer Gottes, die heiligen Mächter, die sich um uns her lagern, die mir deinen Namen fürchten, welche unsere Engel allezeit sehen das Angesicht unsers Vaters im Himmel, die für deinem Thron stehen, derer etlicher Name du uns geoffenbaret hast; als des Erzengels Gabrielä, das heißt, Gottes Kraft; Michael, das ist, wer ist wie Gott? Raphael, Gottes Arzt. Mit welchen heiligen Engeln du die drei Stände, so du selbst geordnet hast auf Erden, beschützest, wie denn der Erzengel Gabriel, der vor Gott stehet, der Jungfrau Maria den Gruß brachte, dem heiligen Priester Zacharias die Geburt Johannis des Täufers verkündigte, und dem heiligen Propheten Daniel die Zeit der Zukunft des Messias offenbarte, als ein herrlicher Kirchen-Engel. Der Großfürst Michael, als ein Regenten- und Schutz-Engel, der die Obrigkeit und Lande schützet, auch für Gottes Volk streitet, und der Engel Raphael, welcher zum Haus-Engel verordnet ist, den Asmodi und Haus-Teufel in die Wüste zu binden. Ach lieber, getreuer Gott, was ist der Mensch, daß du ihn so hoch achtest, dir sage ich ewig Lob und Dank für diese Wohlthat, daß du diese dienstbare Geisterlein ausgesandt
\quellseite{388}
hast zu dienen denen, so die Seligkeit ererben sollen. Und bitte dich, du wollest deinen Engeln über mir Befehl thun, daß sie mich behüten auf allen meinen Wegen; daß sie mich auf den Händen tragen, und ich meinen Fuß nicht an einen Stein stoße, daß ich möge gehen auf den Löwen und Ottern, und treten auf den jungen Löwen und Drachen. Treibe von mir alle bösen Geister, welche sind Lügner und Mörder von Anfang, behüte mich vor ihrem Grimm und Wüthen, vor ihren Lügen und Lästerungen, vor ihrer List und Betrug, daß sie nicht ihr Unkraut unter den Weizen säen in mein Herz. Wehre den Lügen-Geistern in aller falschen Propheten Munde, dem Mordgeist in allen Tyrannen, dem Hoffarts- und Geiz-Teufel im Hausstande. Laß mich deine heiligen Engel allezeit begleiten, wie den Jacob; schützen, wie den Propheten Elisaum, mit feurigen Wagen und Roß, auch Brod und Wasser bringen, wie dem Elia, das ist, Rath und Trost. Laß sie bei mir sein in meinem Kreuz, wie bei den drei Männern im feurigen Ofen, und bei dem Propheten Daniel in der Löwengrube, aus allen meinen Nöthen mich erretten und ausführen, wie den Loth aus dem Feuer zu Sodom, wie Petrum aus dem Gefängniß, Paulum aus dem Schiffbruch. Laß mein Haus und Hof, Kinder, und alles was ich habe, durch deine heiligen Engel bewahret werden, wie das Haus Hiob, daß der Feind keinen Eingriff thun kann. Laß mich in deiner Furcht leben, dein Wort und Evangelium lieb haben, welches die Engel gelüstet anzuschauen. Gib wahre Buße in mein Herz, auf daß sich die Engel über mich freuen im Himmel, zünde in mir an einbrünstiges Gebet und Lob deines Namens, daß ich eines Engels Amt möge verrichten, und mit ihnen singen: Heilig, Heilig, Heilig ist Gott, der HErr Zebaoth, und endlich laß auch meine Seele von den Engeln in Abrahams Schoß getragen werden, und mache mich in der Auferstehung an jenem Tage deinen heiligen Engeln gleich, daß ich derselben Gesellschaft ewig beiwohnen möge. Amen.

\section[442. Danksagung für die gnädige Erhaltung und …]{442. Danksagung für die gnädige Erhaltung und Behütung in so mancherlei Gefahr}

Lobe den HErrn, meine Seele, und alles, was in mir ist, seinen heiligen Namen. Lobe den HErrn, meine Seele, und vergib
\quellseite{389}
nicht, was er dir Gutes gethan hat, der dir alle deine Sünde vergibt, und heilet alle deine Gebrechen, der dein Leben vom Verderben errettet, der dich krönet mit Gnade und Barmherzigkeit. Mein Gott, ich danke dir, daß du meinen Gang erhalten hast auf deinen Fußsteigen, daß meine Tritte nicht haben gegleitet. Du hast deine wunderliche Güte an mir bewiesen wider die, so sich wider deine rechte Hand gesetzet haben, du Heiland derer, die dir vertrauen. Du hast mich behütet wie einen Augapfel im Auge. Du hast mich beschirmet unter dem Schatten deiner Flügel vor den Gottlosen, die mich verstören; vor meinen Feinden, die um und um nach meiner Seele trachten. Darum will ich schauen dein Antlitz in Gerechtigkeit. Ich will satt werden, wenn ich erwache nach deinem Bilde. Du, HErr, erleuchtest meine Leuchte; der HErr, mein Gott, machet meine Finsterniß licht. Ach mein Gott, du hast mich in meiner Noth oft erhöret. Dein Name hat mich oft geschützet. Du hast mir oft Hülfe gesandt vom Heiligthum, und mich aus Zion gestärket. Darum rühme ich, daß du mir hilfest, und in deinem Namen werfe ich Panier auf. Du hast mich mit deinem Segen überschüttet, und erfreuet mit Freuden deines Antlitzes. Du hast nicht verschmähet das Elend des Armen, und dein Antlitz vor mir nicht verborgen; und da ich zu dir schrie, hast du mir ausgeholfen. Du, getreuer Hirte, hast mir, deinem armen Schäflein, nie etwas mangeln lassen, hast mich allezeit auf grüner Aue deines göttlichen Worts geweidet, und zum frischen Trost-Brünnlein geführt. Du hast allezeit meine Seele erquicket, und mich auf rechter Straße geführet, um deines Namens willen. Und, ob ich gleich im finstern Thal des Kreuzes und mancherlei Gefahr oft gewandelt habe; so habe ich doch kein Unglück gefürchtet: denn du bist allezeit bei mir gewesen. Dein Stecken und Stab hat mich allezeit getröstet. Du hast vor mir einen Tisch bereitet gegen meine Feinde. Du hast mein Haupt mit dem himmlischen Freudenöl deines Heiligen Geistes gesalbet, und mir voll eingeschenket deines göttlichen Trostes und Rathes. Du hast mir Gutes und Barmherzigkeit mein Lebenlang folgen lassen, und hoffe, ich werde auch bleiben in deinem Hause immer und ewiglich. Ach HErr, du hast meine Seele aus der Hölle geführet; du hast mich lebendig behalten,
\quellseite{390}
da andere in die Hölle fuhren. Du hast mich unterwiesen, und mir den Weg gezeigt, den ich habe wandeln sollen; du hast mich mit deinen Augen geleitet. Ach Gott, wie theuer ist deine Güte, daß Menschenkinder unter dem Schatten deiner Flügel trauen! Sie werden trunken von den reichen Gütern deines Hauses. Du tränkest sie mit Wollust, als mit einem Strom. Denn bei dir ist die lebendige Quelle, und in deinem Licht sehen wir das Licht. Mein Gott, du hast mich erfahren lassen viel und große Angst, und hast mich wieder lebendig gemacht, und aus der Tiefe der Erde wieder herausgeholet. Darum sei nun wieder zufrieden, meine Seele, der HErr thut dir Gutes. Denn du hast meine Seele aus dem Tode gerissen, meine Augen von den Thränen, meine Füße vom Gleiten. Ich will wandeln vor dem HErrn im Lande der Lebendigen. Wie soll ich dem HErrn vergelten alle seine Wohlthat, die er an mir thut? Ich will den heilsamen Kelch nehmen, und des HErrn Namen predigen. Ach mein Gott, wie oft hast du mich vor dem brüllenden Löwen, dem Teufel, bewahret, daß er meine Seele nicht verschlungen! Wie oft hat der Satan mein begehrt, daß er mich sichten möchte wie den Weizen, du aber, mein HErr JEsu Christe, mein Fürsprecher, hast für mich gebeten, daß mein Glaube nicht aufhörte! Wie oft hast du mich von den Lügenmäulern und von den zänkischen Zungen errettet, die ihre Zungen schärfen wie eine Otter! Wie oft hast du mich aus Noth und Gefahr errettet, und mich bedecket in deiner Hütte zur bösen Zeit, und mich heimlich verborgen in deinem Gezelt vor jedermanns Trog! Wie oft hast du mich vor Zorn, Nachgier und andern fleischlichen Lüsten behütet, daß ich nicht schrecklich gefallen bin! Ja wie oft bin ich gefallen, und du hast mich nicht weggeworfen, sondern mich bei meiner Hand gehalten, und mich wieder aufgerichtet. Ach Gott, wie kann ich dir genugsam für deine gnädige Hülfe, Errettung, Stärke und Trost danken! Wie oft hast du deinen Engeln über mir Befehl gethan, daß sie mich behütet haben auf allen meinen Wegen, daß sie mich haben müssen auf den Händen tragen, und ich meinen Fuß nicht habe an einen Stein gestoßen! Gelobet seist du, o mein HErr und Gott, der du allein Wunder thust, und gelobet sei dein herrlicher Name ewiglich, und alle Lande müssen deiner Ehre voll werden. Amen, Amen.

\quellseite{391}

\section{443. Danksagung nach der Arbeit}

Allmächtiger, ewiger, gütiger Gott! Ich sage dir herzlich Lob und Dank um deine allerheiligste Hülfe und Beistand, die du mir heute diesen Tag erwiesen hast, denn ohne dich hätte ich wahrlich nichts vermocht, verobhalben dir billig allein das Lob gebühret. Ich bitte dich, du wollest dir solch Werk gefallen, und es mir und meinem Nächsten zur Wohlfahrt Leibes und der Seelen gedeihen lassen, durch Jesum Christum, unsern Herrn, welcher mit dir und dem Heiligen Geist, wahrer Gott, lebet und regieret von Ewigkeit zu Ewigkeit. Amen.

\section[444. Danksagung nach Vollendung eines neuen …]{444. Danksagung nach Vollendung eines neuen Wohnhauses}

Gnädiger und barmherziger Gott, ich bin nicht werth, daß du unter mein Dach gehest, doch hast du verheißen, zu den Deinen zu kommen und sie zu segnen. Mit deinem Rath habe ich mir, nach meinem Vermögen, meine häusliche Nahrung zu bessern, diesen Bau vorgenommen und durch deine allmächtige Hülfe ist derselbe nun vollendet, dafür sei dir auch von Grund meines Herzens Lob und Dank gesagt. O du Hüter Israel, bewahre dies mein Gebäude vor Feuer, Hagel und Sturmwind, durch deinen heiligen Wächter. Laß unsre Kindeskinder drinnen wohnen, und hilf, daß es erbe auf unsere Nachkommen. O Gott Vater, Sohn und Heiliger Geist! du heilige Dreifaltigkeit, komme samt deinen heiligen Engeln zu mir und den Meinen, wohne unter uns, daß wir und dies Haus deine Wohnung und Tempel sein und bleiben. Lehre uns aber bedenken, daß wir Fremdlinge und Gäste vor dir sein. Laß uns Friede und Einigkeit lieben, allzeit in die ewigen Hütten und Wohnungen durch das Verdienst und Blut Jesu Christi mit allen heiligen Engeln und Seligen eingehen, und uns immer und ewiglich freuen. Amen.

\section{445. Dankgebet nach glücklich vollbrachter Reise}

Gnädiger Gott, himmlischer Vater, ich danke dir von Herzensgrund, daß du mich meine Reise glücklich lassen vollbringen und durch deiner lieben Engel Schutz mit Freuden wieder zu den Meinigen gebracht, auch für allem Uebel so väterlich bewahret, daß ich
\quellseite{392}
nicht von Räubern und Mördern bin umgebracht, oder von wilden Thieren zerrissen, oder sonst in Wassernoth und andern Gefährlichkeiten verdorben oder beschädiget bin; sondern daß ich, wie der junge Tobias von dem Engel Raphael, frisch und gesund hin und her geleitet worden; dieses alles habe ich deiner väterlichen Vorsorge und allmächtigem Schutz zu danken, und bitte dich von Herzen, du wollest weiter mich und die Meinigen gnädiglich behüten, in deinen gewaltigen Schutz nehmen und an Leib und Seele zum ewigen Leben mächtig erhalten, um Jesu Christi willen. Amen.

\section[446. Dankgebet nach Zurücklegung einer …]{446. Dankgebet nach Zurücklegung einer gefährlichen Seereise (besonders für Einwanderer)}

O du gewaltiger Herrscher Himmels und der Erden! laß doch unsern Mund deines Ruhmes und Preises voll sein täglich, denn dich wollen wir loben alle Zeit, dein Lob soll immerdar in unserm Munde sein; wir danken dem Herrn von ganzem Herzen im Rathe der Frommen und in der Gemeine: Groß sind die Werke des Herrn, und wer ihrer achtet, hat eitel Lust daran. Aber wie sollen wir dem Herrn vergelten alle seine Wohlthaten, die er uns bewiesen! Wir sind viel zu gering aller Barmherzigkeit und Treue, die du an deinen Knechten und Mägden gethan hast. Du hast uns auf dieser Reise erfahren lassen viel und große Angst; du hast uns sehen und erfahren lassen die Gefährlichkeit d. a. Meeres, wie die Wasserströme sich erhoben, greulich gesauset und gebrauset, daß das Schiff fast mit Wellen bedeckt ward, welches bald in den Himmel, bald in den Abgrund fuhr, davon wir wie die Trunkenen taumelten; du warfest uns in die Tiefe des Meeres, daß Wasserfluthen uns umgaben, alle deine Wassermogen und Wellen gingen über uns, daß wir auch gedachten, der Schlund des grausamen Meeres würde uns verschlingen. Ich sprach allbereit: Nun muß ich in mein Grab fahren; nun muß ich nicht mehr schauen die Menschen. Ich dachte oft: O möchte ich bis morgen leben; o wie winselte ich unten im Schiff, wie ein Kranich und Schwalbe, und girrete wie eine Taube, und ward mir um Trost sehr bange. Aber du Menschenhüter und Liebhaber des Lebens! du hast dich meiner Seele und der Meinigen herzlich angenommen, daß sie nicht
\quellseite{393}
ertrunken; denn ich rief zu dem HErrn in meiner Angst, und er antwortete mir; ich schrie aus dem Bauche des Meeres, und du erhörtest meine Stimme: Du stilletest das Brausen des Meeres, das Brausen seiner Wellen, daß sich die stolzen Wellen mußten legen, du hast mein Leben aus dem Verderben geführet, du warst bei mir in der Noth; du hast mich aus derselbigen gerissen, und hast mich auch wieder lebendig gemacht, und aus der Erde geholet; du hast mich sehr groß gemacht und wiederum getröstet. Deshalb lobet man dich in der Stille, und dir bezahlet man Gelübde; du erhörst Gebet. Gleichwie du mich mit Noth nach meines Herzens Wunsch wiederum zu Lande gebracht, daß ich nunmehr nebst den Meinen den Erdboden mit Gesundheit wiederum betreten mag; also komme ich auch jetzt, und schlage einen Dankaltar auf, die Farren meiner Lippen zu opfern, von Herzen sprechend: Nun danket alle Gott, der große Dinge thut an allen Enden; der uns von Mutterleibe an lebendig erhält, und thut uns alles Gutes. Lobe den HErrn, meine Seele, und vergiß nicht, was er dir Gutes gethan. Jauchzet Gott, alle Welt, lobsinget zu Ehren seinem Namen, rühmet ihn herrlich, und sprechet zu Gott: Wie wunderbarlich sind deine Werke. Lobet, ihr Völker, unsern Gott, laßt seinen Ruhm weit erschallen, der unsere Seelen beim Leben erhält, und läßt unsre Füße nicht gleiten; denn Gott, du hast uns versucht und geläutert, wie das Silber geläutert wird; wir sind in Feuer und Wasser kommen, aber du hast uns ausgeführet und erquicket; darum will ich mit Brandopfer gehen in dein Haus, und dir meine Gelübde bezahlen, wie ich meine Lippen habe aufgethan und mein Mund geredet hat in meiner Noth. Und dieses will ich nicht allein in der Gemeinde rühmen, sondern auch unsern Kindern und Nachkommen erzählen, auf daß es unsre Nachkommen lernen, und die Kinder, die noch sollen geboren werden, wenn sie aufkommen, daß sie auch ihren Kindern verkündigen, daß sie setzen auf Gott ihre Hoffnung und nicht vergessen der Thaten Gottes. Ich bitte dich, o gütiger Gott, ferner herzlich und inbrünstig, du wollest ja nicht aufhören, auch künftig mir und den Meinen mehr und mehr Gutes zu thun. Nimm uns allezeit in deinen väterlichen Schutz, bewahre uns vor allem Unfall; nimm
\quellseite{394}
in Gnaden weg von uns, was schädlich ist, gib uns aber, was uns an Leib und Seele ersprießlich sein mag. Laß uns unter deinem Schutz allzeit sicher sein. Gib Glück, Gedeihen und Segen zu eines jeglichen Vornehmen in seinem Amte und Stande; befördere das Werk unserer Hände bei uns, ja das Werk unserer Hände wollest du fördern. Richte all’ unser Vorhaben, Thun und Lassen zu deines Namens Lob, Ehr’ und Preis deiner Kirche, zur Erbauung unseres Nächsten.

Ach HErr! ich trete nunmehr (mit den lieben Meinigen) auf ein neues Land, erneuere uns, und mache an uns Alles neu; schaffe in uns ein reines Herz, und gib uns einen neuen, gewissen Geist. O HErr, hilf! O HErr, laß wohlgelingen! Dein ist die Erde, und du wohnest allenthalben. Sei und bleibe du nur in diesen fremden Oertern unser gnädiger Gott und Vater, unser Versorger und Beförderer; verleihe, daß wir Alles also vor uns finden, wie wir es selbst wünschen, oder uns an Leib und Seele ersprießlich sein mag. Gib auch, du Friedensfürst, unserm Lande den lange gewünschten Frieden beständig; behüte und bewahre eines jeden Landes Obrigkeit; erhalte sie in glückseliger, friedlicher Regierung, Gesundheit und langem Leben. Laß dir auch befohlen sein unsere hinterlassenen, vielgeliebten Freunde, von welchen wir mit so vielen Zähren und Thränen geschieden. Tröste, behüte und bewahre sie vor allem Unglück und Unfall Leibes und der Seele. Segne sie in ihrer Haushaltung und Nahrung, erfreue sie und uns, daß wir allezeit fröhliche und gewünschte Nachricht eines von dem andern haben mögen. Nun segne uns Gott, er gebe, was unser Mund bittet, und das Herz wünschet. Er lasse ihm alle unser Thun und Lassen wohlgefallen. Der HErr segne uns; der HErr behüte uns; der HErr lasse sein Antlitz leuchten über uns, und sei uns gnädig; der HErr erhebe sein Angesicht auf uns, und gebe uns seinen Frieden. Amen.

\section{447. Danksagung für Erlösung aus Trübsal}

Ich danke dir, HErr, mein Gott! von ganzem Herzen und ehre deinen Namen ewiglich. Ich danke dir ewiglich, denn du kannst es wohl machen. Siehe, um Trost war mir sehr bange, du
\quellseite{395}
aber hast dich meiner Seele herzlich angenommen, daß sie nicht verdürbe. Ich schrie zu dir in meiner Angst, da erhörtest du meine Stimme, und mein Geschrei kam vor dich zu deinen Ohren, und du tröstetest mich. Du hast mir treulich beigestanden, meine Seele gnädig erhalten, und mich endlich aus meiner großen Trübsal väterlich errettet. O mein HErr und mein Gott! wie soll ich dir vergelten alle deine Wohlthaten, die du täglich an mir thust? Wohlan! ich will deinen Namen preisen, so lange ich lebe, ich will dich rühmen vor allem deinem Volk, und dein Lob soll immerdar in meinem Munde sein. Hilf mir, o getreuer Gott! durch deinen Geist, daß ich diese meine Gelübde bezahle und zu seiner Zeit vergelte, was du an meiner Seele gethan hast. Daß ich nicht nur in dieser Sterblichkeit dein Lob ausbreite, sondern auch dort in der Ewigkeit mit den Außerwählten mein ewiges Halleluja dir zu Ehren erschallen lasse. Amen.

\section{448. Danksagung nach erlangtem Frieden}

Wir danken dir, gütiger Gott, daß du uns von unsern Feinden Ruhe umher gegeben, und unsere Schmerzen in Freude, und unser Leid in gute Tage verkehret hast. Wir danken dir, daß du uns wie einen Brand aus dem Feuer gerissen hast, daß wir unser Leben als eine Beute davon gebracht. O barmherziger Vater, wie können wir diese große Wohlthat genugsam vergelten, daß du uns, die wir nicht werth sind, unsre Augen zu dir gen Himmel aufzuheben, so gnädig angesehen, unser Gebet erhört und gethan hast, was wir von dir begehrt und gebeten haben. O HErr, du hast ja freilich nach all deiner Barmherzigkeit mit uns gehandelt, daß unser Land (Stadt), so zuvor voll Furcht und Schrecken gewesen, jetzt voll Freude und Wonne worden. Wir bitten demüthiglich, du wollest uns, die wir vom Schwert errettet und übergeblieben, auch forthin Gnade finden, und unser Horn wiederum wachsen lassen, daß wir wieder unter uns wurzeln und über uns Frucht tragen. Baue uns, HErr, daß wir sicher wohnen, und gib Gnade, daß die Städte, so jetzt zerstöret, öde und zerrissen sind, wiederum besezt und festgebauet, und das verwüstete Land wieder gepflüget werde, auf daß alle Völker erfahren, daß du der
\quellseite{396}
Herr bist, der da bauet, was zerrissen ist, und pflanzet, was verheeret war.

O du Gott und Liebhaber des Friedens, verleihe uns forthin stäten Frieden, und laß uns unsere Grenze und Häuser mit Ruhe und Frieden besitzen, daß der Kriegsboten Stimme uns nicht mehr erschrecke, noch der Kriegsmann beschädige. Setze uns in unser Land, und pflanze uns, daß wir aus demselben nicht vertrieben werden, in der Irre gehen, noch die Kinder der Bosheit uns schmähen und drängen mögen. Ach Herr! laß uns nicht mehr in solchen Jammer kommen, sondern ob wir uns forthin an dir versündigen würden, so strafe du uns gnädiglich. Gib uns Frieden zu unsern Zeiten, denn es ist doch kein anderer nicht, der für uns könnte streiten, denn du, unser Gott, allein.

\section[449. Ein Lob der Allmacht, Barmherzigkeit und …]{449. Ein Lob der Allmacht, Barmherzigkeit und Gerechtigkeit Gottes}

Ich will dich erhöhen, mein Gott, du König, und deinen Namen loben immer und ewiglich. Der Herr ist groß, und sehr löblich, und seine Größe ist unaussprechlich. Kindesfinder werden deine Werke preisen, und von deiner Gewalt sagen. Ich will reden von deiner herrlichen schönen Pracht, und von deinen Wundern, daß man soll reden von deinen herrlichen Thaten, und daß man erzähle deine Herrlichkeit, daß man preise deine große Güte, und deine Gerechtigkeit rühme. Gnädig und barmherzig ist der Herr, geduldig und von großer Güte. Der Herr ist allen gütig, und erbarmet sich aller seiner Werke. Es sollen dir danken, Herr, alle deine Werke, und deine Heiligen dich loben, und die Ehre deines Königreichs rühmen, und von deiner Gewalt reden, daß den Menschenkindern deine Gewalt kund werde, und die herrliche Pracht deines Königreichs. Dein Reich ist ein ewig Reich, und deine Herrschaft währet für und für. Der Herr erhält alle, die da fallen; und richtet auf, die da niedergeschlagen sind. Aller Augen warten auf dich, und du gibest ihnen ihre Speise zu seiner Zeit; du thuest deine Hand auf, und erfüllest alles, was lebet, mit Wohlgefallen. Der Herr ist gerecht in allen seinen Wegen, und heilig in allen seinen Werken. Der Herr ist nahe allen denen, die ihn anrufen,
\quellseite{397}
allen, die ihn mit Ernst anrufen. Er thut, was die Gottesfürchtigen begehren, und erhöret ihr Schreien, und hilft ihnen. Der HErr behütet alle, die ihn lieben, und wird vertilgen alle Gottlosen. Mein Mund soll dem HErrn Lob sagen, und alles Fleisch lobe seinen herrlichen Namen immer und ewiglich. Wohl dem, deß Hülfe der Gott Jacob ist; deß Hoffnung auf den HErrn, seinen Gott, stehet, der Himmel, Erd, Meer, und alles, was darinnen ist, gemacht hat; der Glauben hält ewiglich; der Recht schaffet denen, so Gewalt leiden; der die Hungerigen speiset. Der HErr löset die Gefangenen. Der HErr machet die Blinden sehend. Der HErr richtet auf, die niedergeschlagen sind. Der HErr liebet die Gerechten. Der HErr behütet die Fremdlinge und Waisen, und erhält die Wittwen, und kehret zurück den Weg der Gottlosen. Der HErr ist König ewiglich, dein Gott, Zion, für und für, Halleluja. Lobet den HErrn! Denn unsern Gott loben ist ein köstlich Ding, solch Lob ist lieblich und schön. Er heilet, die zerbrochenes Herzens sind, und verbindet ihre Schmerzen. Er zählet die Sterne, und nennet sie alle mit Namen. Unser HErr ist groß, und von großer Kraft, und ist unbegreiflich, wie er regieret. Der HErr richtet auf die Elenden, und stößet die Gottlosen zu Boden. Singet um einander dem HErrn mit Danken; lobet unsern Gott mit Harfen, der den Himmel mit Wolken bedecket, und gibet Regen auf Erden; der Gras auf Bergen wachsen lässet; der dem Viehe sein Futter gibt, den jungen Raben, die ihn anrufen. Er hat nicht Lust an der Stärke des Rosses, noch Gefallen an jemandes Beinen. Der HErr hat Gefallen an denen, die ihn fürchten, und auf seine Güte hoffen. Preise, Jerusalem, den HErrn; lobe, Zion, deinen Gott. Denn er macht fest die Riegel deiner Thore, und segnet deine Kinder darinn. Er schaffet deinen Grenzen Frieden, und sättiget dich mit dem besten Weizen. Er sendet seine Rede auf Erden; sein Wort läuft schnell. Er gibt Schnee wie Wolle; er streuet Reif wie Asche; er wirft seine Schlossen wie Bissen. Wer kann bleiben vor seinem Frost? Er spricht, so zerschmelzet es. Er lässet seinen Wind wehen, so thauet es auf. Er zeiget Jacob sein Wort, und Israel seine Rechte. Der HErr hat ein Wohlgefallen an seinem Volk. Er hilft seinen Elenden herrlich. Die Heiligen
\quellseite{398}
sollen fröhlich sein und preisen, ihr Mund soll Gott erhöhen; und sollen scharfe Schwerter in ihren Händen haben, daß sie Rache üben unter den Heiden, Strafe unter den Völkern, ihre Könige zu binden mit Ketten, und ihre Edlen mit eisernen Fesseln; daß sie ihnen thun das Recht, davon geschrieben ist. Solche Ehre werden alle seine Heiligen haben. Halleluja.

\section[450. Ein Lob der Herrlichkeit, Majestät und …]{450. Ein Lob der Herrlichkeit, Majestät und Freundlichkeit Gottes}

O du ewiger, wahrer, lebendiger Gott! der du mich zu deinem Lobe geschaffen hast, gib mir, daß ich dich würdiglich lobe, du bist ja der Herrlichste, der Löblichste, der Heiligste, der Gerechteste, der Allerfreundlichste, ja du bist der Allerwahrhaftigste, du bist gerecht in allen deinen Werken, und heilig in allen deinen Wegen, du bist der Allerweiseste, dir sind alle deine Werke von Ewigkeit her bewußt, du bist der Allerstärkste, dir mag niemand widerstehen. HErr Jehovah ist dein Name, groß von Rath und mächtig von That, deine Augen sehen auf alle Menschenkinder. Du bist allgegenwärtig, du erfüllest Himmel und Erden. Du bist unendlich, siehest, hörest, regierest alles, du hältst und trägst alles mit deinem kräftigen Wort. Du bist erschrecklich; wenn du das Urtheil lässest hören, so erschrickt das Erdreich, und wird stille. Du legst Ehre ein auf Erden, du legest Ehre ein unter den Völkern, du nimmst den Fürsten den Muth, und bist erschrecklich unter den Königen auf Erden. Du bist erschrecklich, wenn du zürnest; wer kann vor dir bestehen, wenn du zornig bist; die Heiden müssen verzagen, und die Königreiche fallen, das Erdreich muß vergehen, wenn du dich hören lässest. Du bist auch sehr barmherzig, gnädig, geduldig, von großer Güte, und gereuet dich bald der Strafe. Du zürnest nicht ewiglich und gedenkest nicht ewig der Sünden; so groß deine Allmacht ist, so groß ist auch deine Barmherzigkeit, deine Allmacht ist unendlich, und deine Barmherzigkeit hat kein Ende. O ewiges Licht! o ewiges Heil! o ewige Liebe! o ewige Süßigkeit! laß mich dich sehen, laß mich dich empfinden, laß mich dich schmecken. O ewige Lieblichkeit! o ewiger Trost! o ewige Freude! laß mich in dir ruhen, denn ich finde in dir alles, was mir in diesem Elende
\quellseite{399}
mangelt, du bist alle Fülle, und was du nicht bist, das ist eitel Armuth, Jammer und Elend, das Leben ohne dich ist der bittere Tod, deine Güte ist besser denn Leben. Ach, du überköstlicher Schatz, du ewiges Gut, du liebliches Leben, wenn werde ich vollkommenlich mit dir vereinigt werden, daß ich dich vollkommenlich schmecke, und dich vollkommenlich sehe. Heiliger Gott, unsterblicher Gott, gerechter Gott, allweiser Gott, du ewiger König, dir sei Lob, Ehre und Preis in alle Ewigkeit. Amen.

\section[451. Preis Jesu Christi von einer zu seinem …]{451. Preis Jesu Christi von einer zu seinem Lobe entzündeten Seele}

O Jesu Christe, meine Liebe, meine Freude, mein Licht, mein Heil, mein Schmuck, mein König, mein Hirte, mein Bräutigam, mein ewiger Hoherpriester, mein Leben, meine Weisheit, mein Friede und Ruhe meiner Seele, meine Gerechtigkeit, meine Heiligung, meine Erlösung, mein Opfer, mein Fürsprecher, mein Gnadenthron, mein einiger Mittler, mein einiger Nothhelfer, meine einige Hoffnung, mein höchster Schatz, mein starker Schutz und einige Zuflucht! Wie soll ich dich würdiglich loben? O du Glanz der Herrlichkeit Gottes, du allerschönstes wesentliches Ebenbild des Vaters! Heilig, gerecht und selig bin ich in dir; ohne dich ungerecht, unheilig und verdammt. Unselig bin ich, wo ich bin, wo du nicht bei mir und in mir bist. Komm, meine Freude, in mein Herz, und erfreue mich. Komm, du Trost aller Heiden, und tröste mich. Komm, meine Liebe, und erquicke mich. Komm, mein Leben, und stärke mich. Komm, mein Licht, und erleuchte mich. Komm, meine Süßigkeit, daß ich deine Freundlichkeit schmecke. Komm, du schöne, edle Gestalt, daß ich dich sehe. Komm, meine Lieblichkeit, daß ich dich höre. Komm, du edle Blume, daß ich dich rieche. Komm, du zarte Bewegung, daß ich dich empfinde. Dein edler Anblick bewege mich; dein holdseliges Anschauen erfreue mich; dein edler Geruch erquicke mich; dein heiliges Gedächtniß entzünde mich; deine edle Liebe sättige mich; deine Einwohnung belustige mich. Ach du edle Demuth, du heilige Sanftmuth, vereinige mein Herz und Gemüth mit dir. Das wird der köstliche
\quellseite{400}
Schmuck sey meiner Seele. Du bist mir theurer, denn alle Demanten; köstlicher denn alle Rubinen, lieblicher denn alle Perlen. O du immerblühende Rose, o du Paradies aller geistlichen himmlischen Wollust, o du allersüßeste Frucht aus dem Lustgärtlein Gottes! Laß mich dich ewig schmecken, fülle mich mit dem Kuß deines Mundes! O du ewiges Wort des Vaters, rede in mir, auf daß die ganze Welt in mir schweige. O du ewige Klarheit, leuchte in mir. O du ewige Wahrheit, lehre mich. O du heilsamer Arzt, heile mich. O du himmlischer Schatz, zeuch mein Herz nach dir. O du himmlischer Noah, strecke deine Hand aus, und nimm das arme Täublein meiner Seele zu dir; denn es kann nirgend Ruhe finden. Wie lieblich sind deine Wohnungen, HErr Zebaoth! Meine Seele verlanget und sehnet sich nach den Vorhöfen des HErrn. Mein Leib und Seele freuet sich in dem lebendigen Gott. Denn der Vogel hat ein Haus funden, und die Schwalbe ihr Nest, da sie Junge hecken, nämlich deine Altäre, HErr Zebaoth, mein König und mein Gott! Wohl denen, die in deinem Hause wohnen, die loben dich immerdar. Wohl den Menschen, die dich für ihre Stärke halten, und von Herzen dir nachwandeln. O HErr Jesu, wie lieblich ist deine Menschwerdung! wie wunderbar ist dein Name! wie großmächtig deine Werke! wie heilig dein Wort! wie tröstlich dein Leiden! wie sieghaft deine Auferstehung! wie herrlich deine Himmelfahrt! wie groß ist deine Ehre! wie hoch und erhaben ist dein königlicher Stuhl, da dich alle Engel Gottes anbeten! wie unaussprechlich ist deine Gewalt! wer wollte dich nicht fürchten, du König der Ehren! Machet die Thore weit, und die Thüren in der Welt hoch, daß der König der Ehren einziehe. Wer ist derselbe König der Ehren? Es ist der HErr, stark und mächtig, der HErr mächtig im Streit. Machet die Thore weit, und die Thüren in der Welt hoch, daß der König der Ehren einziehe. Wer ist derselbe König der Ehren? Es ist der HErr Zebaoth, er ist der König der Ehren, Sela. Gelobet sei, der da kommt im Namen des HErrn. Der HErr ist Gott, der uns erleuchtet. Danket dem HErrn, denn er ist freundlich, und seine Güte währet ewiglich.

\quellseite{401}

\chapter{XIII. Gebete für Kranke und Sterbende}

\section{452. Morgengebet eines Kranken}

O du getreuer Gott und Vater unsers Herrn Jesu Christi! ich sage dir Lob und Dank, daß du mich schwachen Menschen diese Nacht über so gnädiglich erhalten, geschützt und gestärket hast. Ich bitte dich von Herzen, im Namen Jesu Christi, du wollest auch diesen Tag und die übrige Zeit meines Lebens deine Barmherzigkeit zu mir wenden, mir meine Sünde vergeben, durch deine Gnade trösten, durch deine Kraft stärken, durch deinen Heiligen Geist leiten, damit ich nach deinem Willen leben und leiden möge. Ich befehle auch diesen meinen kranken Leib, und meine liebe Seele in deine heiligen Hände, dein heiliger Engel sei mit mir, daß der böse Feind keine Gewalt noch Macht an mir finde und habe, durch Jesum Christum, deinen geliebten Sohn, meinen einigen Trost und Heiland, in der Kraft des Heiligen Geistes. Amen.

\section{453. Abendgebet eines Kranken}

O Gott Vater! ich danke dir, daß du mir durch diesen Tag so gnädig geholfen, und bitte dich von Herzen, hilf mir schwachen und kranken Menschen auch durch diese Nacht mit deiner Gnade, mit deiner Kraft, mit deinem Trost, mit deinem Geist. Ich befehle dir meinen Leib und Seele in deine Hände, dein heiliger Engel sei mit mir, daß der böse Feind keine Macht an mir finde, durch Jesum Christum, deinen lieben Sohn, unsern Herrn. Amen.

\section{454. Morgensegen eines Kranken}

Ich befehle mich heute diesen Morgen und die ganze Zeit meines Lebens dir, Gott, dem ewigen Vater, der du mich erschaffen hast, dir, Gott, dem ewigen Sohn, der du mich erlöset hast, dir, Gott, dem Heiligen Geiste, der du mich in meiner Taufe geheiliget hast. O ewige Gottheit, o wahre Menschheit, o heilige Dreifaltigkeit, bereite mich heute und allezeit, wie du mich haben willst in der ewigen Seligkeit. Amen.

\quellseite{402}

Gütiger, barmherziger Gott, herzlieber Vater, durch deine Gnade habe ich abermal eine Nacht hingebracht und das erfreuliche Tageslicht wieder erlebet. Ich danke dir, daß du mein Hütter und Wächter, Tröster und Helfer gewesen, meinen schwachen Leib, samt der darin wohnenden betrübten Seele, väterlich erhalten, mich vor schnellem Tode, schweren Anfechtungen und anderem Uebel kräftiglich bewahret, und durch deine heiligen Engel wider den Teufel, den Fürsten der Finsterniß, mächtiglich beschirmet hast. Deine Gnade und Barmherzigkeit ist’s, daß ich noch lebe. Ich bitte dich von ganzem Herzen, du wollest diesen Tag auch dein göttlich Antlitz zu mir armen Creatur wenden, und mich also leiten, regieren und führen, daß ich nichts Böses gedenke, rede oder thue. Weide und leite mich, dein armes Schäflein, mit und zu deinem heiligen Worte, welches ist mein Trost in meinem Elend. Bewahre mich, o HErr, denn ich verlasse mich auf dich, du bist ja meine einige Stärke, mein Gott und höchstes Gut. Ach, richte meinen Muth und Sinn dahin, daß ich dich fürchte und zugleich auch liebe. Kehre alles ab von mir, das nicht lauter dein ist. Senke mich in dich, daß ich in meinem Kreuz und Elend meinen Willen unter dich lasse, und dir allein vertraue, weil meine und aller Creaturen Hülfe mir nicht rathen, noch mich retten kann. Du bist doch der rechte Arzt aller Kranken, Trost aller Betrübten, Helfer aller Sterbenden. Heile mich, tröste mich, hilf mir, o liebster Vater. Behüte mich sonderlich diesen Tag vor Sünden und allem Uebel, vor den Listen des Teufels, den Aergernißen der Welt und bösen Lüsten meines sündlichen Fleisches. Verleihe Geduld, Hoffnung und Beständigkeit. Und wenn meine Zeit zu leben um ist, so verleihe mir ein sanftes, seliges Ende durch JEsum Christum, deinen lieben Sohn, meinen Heiland. Amen.

\section{455. Abendsiegen eines Kranken}

Gott sei mir gnädig, gütig und barmherzig, und verleihe mir seinen göttlichen Segen. Er lasse sein Antlitz im Finstern über mir leuchten, daß ich nicht im Tod entschlafe. Die heilige Dreifaltigkeit sei mit mir, und bleibe allzeit in meinem Herzen, ich schlafe oder wache, lebe oder sterbe. Amen.

\quellseite{403}

Allmächtiger und barmherziger Gott, herzliebster Vater, der du gesagt hast: Rufe mich an in der Zeit der Noth, so will ich dich erretten; ich bitte dich bei dieser Nachtstunde, erbarme dich mein, und vergib mir alle meine Sünde, womit ich etwa diesen Tag dich beleidigt habe. Laß mich in dieser einbrechenden Nacht dir befohlen sein, daß ich in dir, wo nicht allzeit mit dem Leibe, doch mit der Seele meine Ruhe habe, und verhüte gnädiglich, daß ich nicht etwa durch einen schnellen Tod ohne Erkenntniß deines lieben Sohnes hingerissen werde. Ist es dein gnädiger Wille, so laß mich ohne Schmerzen des Leibes und Betrübniß der Seele schlafen, und den morgenden Tag wieder erleben. Sollte aber nach deinem unerforschlichen Rath dieses meine letzte Nacht in der Welt sein, und du wolltest mich darinnen abfordern, so bitte ich demüthiglich, du wollest mich zu dem ewigen Himmelstag gelangen und in deinem Reich mich den geringsten Thürhüter sein lassen. Wohlan, ich will in deinem Namen einschlafen. Ich mag zum natürlichen oder ewigen Leben wieder erwachen, so will ich dich, meinen getreuen Gott, herzlich lieben und loben. Amen. Im Namen Jesu. Amen.

\section{456. Gebet in Krankheit}

Barmherziger und gerechter Gott, der du bist ein Herr der Gesundheit und der Krankheit, des Lebens und des Todes, ich erkenne vor dir mit aufrichtigem Herzen, daß der große Mißbrauch der vielen Tage der Gesundheit und des Wohlstandes, welche du mir verliehen hast, diese gegenwärtige Heimsuchung gar billig von dir verdienet hat. O Herr, ich bin von Herzen begierig, in rechtschaffener Demuth diese Strafe meiner Sünden von deiner Hand anzunehmen, und den Zorn des Herrn zu tragen, dieweil ich wider ihn gesündiget habe. Ich bitte dich, o du gnädiger und barmherziger Vater! der du nicht das Verderben, sondern die Besserung zum Zwecke hast, wenn du züchtigest, du wollest doch auch mir durch deine Gnade diese deine Züchtigung dergestalt heiligen, damit diese Krankheit meines Leibes mir möge ein Mittel zur Gesundheit meiner Seele werden. Heile meine Seele, o Herr, die wider dich gesündiget hat, und dann, so es dein heiliger Wille ist, so heile auch meinen Leib, und hilf, daß ich inskünftige mein Leben zu deinem Lob
\quellseite{404}
anstelle, und rechtschaffene Früchte der Buße tragen möge. Solltest du aber nach deiner Weisheit es anders beschloßen haben, daß mir diese Krankheit zum Tode sein sollte, so bitte ich dich, du wollest mich dazu vorbereiten und geschickt machen. Verleihe mir eine herzliche und unverfälschte Buße, welcher du deine Gnade und Vergebung versprochen hast, ziehe mein Herz von der Welt und allen ihren vergänglichen Eitelkeiten ab, und mache mir Sehnen und Verlangen nach deinen herrlichen und beständigen Gütern, welche zu deiner Rechten sind ewiglich. Herr, erhebe du das Licht deines Antlitzes über mich, und laße in allen meinen Leibesbeschmerzen und Seelenängsten deinen Trost mein Herz erquicken, auf daß ich tüchtig sei mit Geduld zu warten, bis meine Veränderung kommet. Verleihe mir, o Herr! daß, wenn mein irdisches Haus dieser Hütten wird zerbrochen sein, ich einen Bau möge haben von Gott erbauet, ein Haus, das nicht mit Händen gemacht ist, sondern das ewig ist im Himmel. Thue solches um deßetwillen, welcher durch sein theures Blut es mir erworben hat, nämlich um JEsu Christi willen. Amen.

\section{457. Gebet eines Todkranken}

HErr Gott Vater im Himmel, erbarme dich über mich kranken und sterbenden Menschen. Weil du mich mit Leib und Seele erschaffen und nachgehends in der heiligen Taufe um JEsu Christi willen zu deinem Kinde angenommen, und gnädiglich verheißen hast: Wie sich ein Vater über die Kinder erbarmet, so wollest du dich Zeit meines Lebens über mich erbarmen; ei, so laß doch dir jetzt solch mein Leib und Seele in deine treuen Vaterhände befohlen sein, und dein erbarmendes Vaterherz über mich, dein Kind, brechen. Ach! himmlischer Vater, rechne mir doch meinen Ungehorsam, ja alle meine von Kindesbeinen an begangene Sünden nicht zur Verdammniß zu, sondern siehe an die vollgültige Bezahlung deines Sohnes JEsu Christi, um deßwillen schone, mir, armen Sünder, nicht nach Sünden lohne. Erinnere dich des theuren Eides, den du geschworen hast: So wahr ich lebe, ich will nicht den Tod des Sünders, sondern daß er sich bekehre und lebe. Gedenke nicht der Sünden meiner Jugend, noch meiner
\quellseite{405}
Uebertretung, gedenke aber mein nach deiner großen Barmherzigkeit. Siehe an meinen Jammer und Elend, und vergib mir alle meine Sünden. Laß mich dein sein und bleiben, o treuer Gott und Herr, von dir laß mich nicht treiben, halt mich bei reiner Lehr, Herr, laß mich nur nicht wanken, gib mir Beständigkeit, dafür will ich dir danken in-alle Ewigkeit. Wenn meine Augen brechen, mein Athem geht schwer aus, und kann kein Wort mehr sprechen, Herr! nimm mein Seufzen auf. Wenn mein Herz und Gedanken vergehn als wie ein Licht, das hin und her thut wanken, wenn ihm die Flamm gebricht: alsdann fein sanft und stille laß mich, Herr, schlafen ein, nach deinem Rath und Willen, wenn kömmt mein Stündelein. Herr Gott Sohn, Jesu Christe, der Welt Heiland, erbarme dich über mich kranken und sterbenden Menschen: du bist um meinetwillen vom Himmel kommen, und hast auf der Welt um meinetwillen so viel gelitten und ausgestanden, ja dein Blut so mildiglich am Kreuz für mich vergossen, und mich damit von Sünde, Tod, Teufel und Hölle erlöset. O so laß doch, lieber Herr Jesu, solche deine Mühe und Arbeit, deine Marter und Pein, dein Kreuz und Tod an mir armen Sünder nicht verloren sein, sondern bitte kraft deines Verdienstes im Himmel droben für mich, damit ich Gnade und Barmherzigkeit erlange jetzt zu der Zeit, da mir Gnade und Barmherzigkeit noth ist. Ach! laß mich nicht in dieser Noth umkommen noch verzagen, komm mir zu Hülfe, o getreuer Gott, meine Angst hilf mir auch tragen, denk’, daß ich bin am Leibe dein ein Glied und grünes Zweigelein, in Fried’ laß mich hinfahren. Schreie du mir jetzt zu, mein Trost, in mein Herz und Ohren: Fürchte dich nicht, liebes Kind, ich habe dich erlöset, ich habe dich bei deinem Namen gerufen, du bist mein. Ich bin bei dir in der Noth, ich will dich heraus reißen und zu Ehren machen, ich will dich sättigen mit langem Leben, und will dir zeigen mein Heil. Erschein mir in dem Bilde, zu Trost in meiner Noth, wie du, Herr Christ! so milde dich hast geblut’t zu Tod. Stärk’ mich mit deinem Freuden-geist, heil mich mit deinen Wunden, wasch mich mit deinem Todes-schweiß in meiner letzten Stunden, und nimm mich einst, wann dir’s gefällt, in wahrem Glauben aus der Welt zu deinen Auserwählten. Sei, werthester Heiland, mein Fürsprecher
\quellseite{406}
am Tage des strengen Gerichts Gottes, und hilf mir auch zur selben Zeit, wenn wegen deiner Wunden, daß ich im Buch der Seligkeit werd eingezeichnet funden! Daran ich denn auch zweifle nicht, denn du hast ja den Feind gericht’t, und meine Schuld bezahlet. HErr Gott Heiliger Geist, erbarme dich über mich kranken und sterbenden Menschen! Vertritt mich zuvörderst in meiner größten Schwachheit bei Gott mit unaussprechlichem Seufzen, und gib Zeugniß meinem Geist, daß ich ein Kind Gottes, und Erbe der ewigen Seligkeit sei. Stehe mir bei mit deinem kräftigen Trost wider alle Anfechtungen des bösen Geistes. Vertreibe bei mir die Furcht des Todes mit Vorstellung der fröhlichen Auferstehung und herrlichen Gnadenbelohnung. Erhalte mich in herzlicher Liebe zu Gott dem himmlischen Vater, und in wahrem Glauben an Christum JEsum, meinen Erlöser, und kürze mir meine Qual und Todesschmerzen, und hilf mir ritterlich ringen, durch Tod und Leben zu Gott dringen. O Heiliger Geist! ein Tröster heißt, an meinem End’ dein’ Hülff’ mir send’; verlaß mich nicht, wenn mich ansicht des Teufels Gewalt, des Todes Gestalt; nach deinem Wort, du treuer Hort, wollst du mir geben das ew’ge Leben, Amen. Ach du heilige und hochgelobte Dreifaltigkeit, erbarme dich über mich kranken und sterbenden Menschen, sei über mir und beschütze mich, sei neben mir, und geleite mich glücklich durch das finstere Todesthal, sei unter mir, und segne mich, sei vor mir, und bewahre mich zum ewigen Leben, Amen. O! du großer Gott! erhöre, was dein Kind gebeten hat: JEsu! den ich stets verehre, bleibe ja mein Schutz und Rath, und mein Hort, du werther Geist! der du Freund und Tröster heißt, höre doch mein sehnlich Flehen! Amen, ja es soll geschehen! HErr JEsu, dir lebe ich, HErr JEsu, dir sterbe ich, dein bin ich todt und lebendig. Amen.

\section[458. Trostsprüche, den Kranken und Sterbenden …]{458. Trostsprüche, den Kranken und Sterbenden vorzusprechen}

1 Mos. 3, 15. Des Weibes Same soll der Schlange den Kopf zertreten.

1 Mos. 32, 26. HErr, ich lasse dich nicht, du segnest mich denn.

\quellseite{407}

2 Mos. 34, 6. 7. HErr, HErr Gott, barmherzig, und gnädig, und geduldig, und von großer Gnade und Treue; der du beweisest Gnade in tausend Glied, und vergibst Missethat, Uebertretung und Sünde, und vor welchem niemand unschuldig ist.

1 Kön. 19, 4. Elias sprach: Es ist genug, so nimm nun, HErr, meine Seele; ich bin nicht besser, denn meine Väter.

Hiob 5, 17. 18. Selig ist der Mensch, den Gott strafet; darum weigere dich der Züchtigung des Allmächtigen nicht. Denn er verletzet, und verbindet; er zerschmeißet, und seine Hand heilet.

Hiob 5, 19. Aus sechs Trübsalen wird er dich erretten, und in der siebenten wird dich kein Uebel rühren.

Hiob 14, 1. 2. Der Mensch vom Weibe geboren lebt kurze Zeit, und ist voll Unruhe, gehet auf wie eine Blume, und fällt ab, fleucht wie ein Schatten, und bleibet nicht.

Hiob 19, 25—27. Ich weiß, daß mein Erlöser lebet, und er wird mich hernach aus der Erde auferwecken, und werde darnach mit dieser meiner Haut umgeben werden, und werde in meinem Fleisch Gott sehen. Denselben werde ich mir sehen, und meine Augen werden ihn schauen und kein Fremder.

Ps. 4, 9. Ich liege und schlafe ganz mit Frieden; denn allein du, HErr, hilfst mir, daß ich sicher wohne.

Ps. 10, 17. Das Verlangen der Elenden hörest du, HErr, ihr Herz ist gewiß, daß dein Ohr drauf merket.

Ps. 13, 6. Ich hoffe aber darauf, daß du so gnädig bist; mein Herz freuet sich, daß du so gerne hilfst. Ich will dem HErrn singen, daß er so wohl an mir thut.

Ps. 17, 15. Ich aber will schauen dein Antlitz in Gerechtigkeit; ich will satt werden, wenn ich erwache nach deinem Bilde.

Ps. 22, 27. Die Elenden sollen essen, daß sie satt werden; und die nach dem HErrn fragen, werden ihn preisen; euer Herz soll ewiglich leben.

Ps. 23. Der HErr ist mein Hirte; mir wird nichts mangeln. Er weidet mich auf einer grünen Aue, und führet mich zum frischen Wasser. Er erquicket meine Seele; er führet mich auf
\quellseite{408}
rechter Straße, um seines Namens willen. Und ob ich schon wanderte im finstern Thal, fürchte ich kein Unglück; denn du bist bei mir, dein Stecken und Stab trösten mich. Du bereitest vor mir einen Tisch gegen meine Feinde. Du salbest mein Haupt mit Oel, und schenkest mir voll ein. Gutes und Barmherzigkeit werden mir folgen mein Lebenlang, und ich werde bleiben im Hause des HErrn immerdar.

Ps. 25, 6. 7. Gedenke, HErr, an deine Barmherzigkeit, und an deine Güte, die von der Welt her gewesen ist. Gedenke nicht der Sünden meiner Jugend, und meiner Uebertretung; gedenke aber mein nach deiner Barmherzigkeit, um deiner Güte willen.

Ps. 27, 1. Der HErr ist mein Licht und mein Heil; vor wem sollt ich mich fürchten? Der HErr ist meines Lebens Kraft; vor wem sollte mir grauen?

Ps. 27, 9. Verbirg dein Antlitz nicht vor mir, und verstoße nicht im Zorn deinen Knecht; denn du bist meine Hülfe. Laß mich nicht, und thue nicht von mir die Hand ab, Gott, mein Heil.

Ps. 31, 6. In deine Hände befehl ich meinen Geist; du hast mich erlöset, HErr, du treuer Gott.

Ps. 34, 20. Der Gerechte muß viel leiden; aber der HErr hilft ihm aus dem allen.

Ps. 37, 4. 5. Habe deine Lust am HErrn; der wird dir geben, was dein Herz wünschet. Befiehl dem HErrn deine Wege, und hoffe auf ihn; er wirds wohl machen.

Ps. 39, 6. Siehe, meine Tage sind einer Hand breit bei dir, und mein Leben ist wie nichts vor dir. Wie gar nichts sind alle Menschen, die doch so sicher leben!

Ps. 42, 2. 3. Wie der Hirsch schreiet nach frischem Wasser, so schreiet meine Seele, Gott, zu dir. Meine Seele dürstet nach Gott, nach dem lebendigen Gott. Wann werde ich dahin kommen, daß ich Gottes Angesicht schaue?

Ps. 43, 5. Was betrübest du dich, meine Seele, und bist so unruhig in mir? Harre auf Gott; denn ich werde ihm noch danken, daß er meines Angesichts Hülfe und mein Gott ist.

\quellseite{409}

Ps. 51, 12—14. Schaffe in mir, Gott, ein rein Herz, und gib mir einen neuen gewissen Geist. Verwirf mich nicht von deinem Angesichte, und nimm deinen heiligen Geist nicht von mir. Tröste mich wieder mit deiner Hülfe, und der freudige Geist enthalte mich.

Ps. 55, 23. Wirf dein Anliegen auf den HErrn. Der wird dich versorgen, und wird den Gerechten nicht ewiglich in Unruhe lassen.

Ps. 68, 20. 21. Gelobet sei der HErr täglich. Gott legt uns eine Last auf, aber er hilft uns auch. Wir haben einen Gott, der da hilft, und den HErrn HErrn, der vom Tode errettet.

Ps. 73, 23—26. Dennoch bleibe ich stets an dir; denn du hältst mich bei meiner rechten Hand, du leitest mich nach deinem Rath, und nimmst mich endlich mit Ehren an. Wenn ich nur dich habe, so frage ich nichts nach Himmel und Erde. Wenn mir gleich Leib und Seele verschmachtet, so bist du doch, Gott, allezeit meines Herzens Trost und mein Theil.

Ps. 91, 14—16. Er begehret mein, so will ich ihm aushelfen; er kennet meinen Namen, darum will ich ihn schützen. Er ruft mich an, so will ich ihn erhören; ich bin bei ihm in der Noth, ich will ihn herausreißen, und zu Ehren machen; ich will ihn sättigen mit langem Leben, und will ihm zeigen mein Heil.

Ps. 94, 19. Ich hatte viel Bekümmernisse in meinem Herzen; aber deine Tröstungen ergötzten meine Seele.

Ps. 103, 1—3. Lobe den HErrn, meine Seele, und was in mir ist, seinen heiligen Namen; lobe den HErrn, meine Seele, und vergiß nicht, was er dir Gutes gethan hat, der dir alle deine Sünden vergibt, und heilet alle deine Gebrechen.

Ps. 103, 8—14. Barmherzig und gnädig ist der HErr, geduldig und von großer Güte. Er wird nicht immer hadern, noch ewiglich Zorn halten. Er handelt nicht mit uns nach unsern Sünden, und vergilt uns nicht nach unsrer Missethat. Denn so hoch der Himmel über der Erde ist, läßt er seine Gnade walten über die, so ihn fürchten. So ferne der Morgen ist vom Abend, lässet er unsre Uebertretung von uns sein. Wie sich ein Vater über Kinder erbarmet, so erbarmet sich der HErr über die, so ihn
\quellseite{410}
fürchten. Denn er kennet, was für ein Gemächte wir sind; er gedenket daran, daß wir Staub sind.

Ps. 143, 2. Gehe nicht ins Gericht mit deinem Knechte; denn vor dir ist kein Lebendiger gerecht.

Ps. 145, 18. 19. Der HErr ist nahe allen, die ihn anrufen, allen, die ihn mit Ernst anrufen. Er thut, was die Gottesfürchtigen begehren, und höret ihr Schreien, und hilft ihnen.

Spriche 18, 10. Der Name des HErrn ist ein festes Schloß; der Gerechte läufft dahin, und wird beschirmet.

Jes. 1, 18. Wenn eure Sünde gleich blutroth ist, soll sie doch schneeweiß werden; und wenn sie gleich ist, wie Rosinfarbe, soll sie doch wie Wolle werden.

Jes. 35, 10. Die Erlöseten des HErrn werden wiederkommen, und gen Zion kommen mit Jauchzen; ewige Freude wird über ihrem Haupte seyn; Freude und Wonne werden sie ergreifen, und Schmerz und Seufzen wird weg müssen.

Jes. 38, 17. Siehe, um Trost war mir sehr bange. Du aber hast dich meiner Seele herzlich angenommen, daß sie nicht verdürbe; denn du wirfst alle meine Sünde hinter dich zurück.

Jes. 41, 10. Fürchte dich nicht, ich bin mit dir; weiche nicht, denn ich bin dein Gott. Ich stärke dich, ich helfe dir auch, ich erhalte dich durch die rechte Hand meiner Gerechtigkeit.

Jes. 42, 3. Das zerstoßene Rohr wird er nicht zerbrechen, und das glimmende Docht wird er nicht auslöschen.

Jes. 43, 24. 25. Mir hast du Arbeit gemacht in deinen Sünden und hast mir Mühe gemacht in deinen Missethaten. Ich, ich tilge deine Uebertretung um meinetwillen, und gedenke deiner Sünden nicht.

Jes. 49, 14—16. Zion spricht: der HErr hat mich verlassen, der HErr hat mein vergessen. Kann auch ein Weib ihres Kindleins vergessen, daß sie sich nicht erbarme über den Sohn ihres Leibes? Und ob sie desselbigen vergäße, so will ich doch dein nicht vergessen. Siehe, in die Hände habe ich dich gezeichnet.

\quellseite{411}

Jes. 53, 4. 5. Fürwahr, Er trug unsre Krankheit, und lud auf sich unsre Schmerzen. Wir aber hielten ihn für den, der geplagt und von Gott geschlagen und gemartert wäre. Aber Er ist um unsrer Missethat willen verwundet, und um unsrer Sünde willen zerschlagen. Die Strafe liegt auf ihm, auf daß wir Friede hätten, und durch seine Wunden sind wir geheilet.

Jes. 54, 7. 8. Ich habe dich einen kleinen Augenblick verlassen; aber mit großer Barmherzigkeit will ich dich sammeln. Ich habe mein Angesicht im Augenblick des Zorns ein wenig von dir verborgen, aber mit ewiger Gnade will ich mich dein erbarmen, spricht der HErr, dein Erlöser.

Jes. 54, 10. Es sollen wohl Berge weichen und Hügel hinfallen; aber meine Gnade soll nicht von dir weichen und der Bund meines Friedens soll nicht hinfallen, spricht der HErr, dein Erbarmer.

Jes. 55, 8. 9. Meine Gedanken sind nicht eure Gedanken, und eure Wege sind nicht meine Wege, spricht der HErr: sondern, soviel der Himmel höher ist, denn die Erde, so sind auch meine Wege höher, denn eure Wege, und meine Gedanken, denn eure Gedanken.

Jes. 57, 15. Also spricht der Hohe und Erhabene, der ewiglich wohnet, deß Name heilig ist; der ich in der Höhe und im Heiligthum wohne, und bei denen, so zerschlagenen und demüthigen Geistes sind, auf daß ich erquicke den Geist der Gedemüthigten, und das Herz der Zerschlagenen.

Jer. 31, 20. Ist nicht Ephraim mein theurer Sohn, und mein trautes Kind? Denn ich denke noch wohl dran, was ich ihm geredet habe; darum bricht mir mein Herz gegen ihn, daß ich mich sein erbarmen muß, spricht der HErr.

Klagl. Jer. 3, 31—33. Der HErr verstößt nicht ewiglich; sondern er betrübet wohl, und erbarmet sich wieder nach seiner großen Güte; denn er nicht von Herzen die Menschen plaget und betrübet.

Ezech. 33, 11. So wahr als ich lebe, spricht der HErr HErr, ich habe keinen Gefallen am Tode des Gottlosen, sondern daß sich der Gottlose bekehre von seinem Wesen und lebe.

\quellseite{412}

Hos. 13, 14. Ich will sie erlösen aus der Hölle, und vom Tode erretten. Tod, ich will dir ein Gift sein; Hölle, ich will dir eine Pestilenz sein.

Micha 7, 18. Wo ist solch ein Gott, wie du bist? der die Sünde vergibt, und erlässet die Missethat den Uebrigen seines Erbtheils; der seinen Zorn nicht ewiglich behält; denn er ist barmherzig.

Weish. 3, 1. Der Gerechten Seelen sind in Gottes Hand, und keine Qual rühret sie an.

Matth. 5, 3. 4. 6. Selig sind, die da geistlich arm sind; denn das Himmelreich ist ihr. Selig sind, die da Leid tragen; denn sie sollen getröstet werden. Selig sind, die da hungert und dürstet nach der Gerechtigkeit; denn sie sollen satt werden.

Matth. 11, 28. Kommet her zu mir alle, die ihr mühselig und beladen seid, ich will euch erquicken.

Marc. 16, 16. Wer da glaubet und getauft wird, der wird selig werden; wer aber nicht glaubet, der wird verdammt werden.

Luc. 2, 29—32. HErr, nun lässest du deinen Diener in Frieden fahren, wie du gesagt hast; denn meine Augen haben deinen Heiland gesehen, welchen du bereitet hast vor allen Völkern. Ein Licht zu erleuchten die Heiden, und zum Preis deines Volkes Israel.

Luc. 5, 31. 32. Die Gesunden bedürfen des Arztes nicht, sondern die Kranken. Ich bin kommen, zu rufen den Sündern zur Buße, und nicht den Gerechten.

Luc. 15, 21. Vater, ich habe gesündiget in den Himmel und vor dir; ich bin fort nicht mehr werth, daß ich dein Sohn heiße.

Luc. 18, 7. 8. Sollte aber Gott nicht auch retten seine Auserwählten, die zu ihm Tag und Nacht rufen, und sollte Geduld darüber haben? Ich sage euch: Er wird sie erretten in einer Kürze.

Luc. 19, 10. Des Menschen Sohn ist kommen, zu suchen und selig zu machen, das verloren ist.

Luc. 23, 43. Und Jesus sprach zu ihm: Wahrlich, ich sage dir, heute wirst du mit mir im Paradiese sein.

Joh. 1, 29. Siehe, das ist Gottes Lamm, welches der Welt Sünde trägt.

\quellseite{413}

Joh. 3, 16. 17. Also hat Gott die Welt geliebet, daß er seinen eingebornen Sohn gab, auf daß alle, die an ihn glauben, nicht verloren werden, sondern das ewige Leben haben. Denn Gott hat seinen Sohn nicht gesandt in die Welt, daß er die Welt richte, sondern daß die Welt durch ihn selig werde.

Joh. 5, 24. Wahrlich, wahrlich, ich sage euch, wer mein Wort höret, und glaubet dem, der mich gesandt hat, der hat das ewige Leben, und kommt nicht in das Gericht, sondern er ist vom Tode zum Leben hindurch gedrungen.

Joh. 5, 25. Wahrlich, wahrlich, ich sage euch, es kommt die Stunde, und ist schon jetzt, daß die Todten werden die Stimme des Sohnes Gottes hören; und die sie hören werden, die werden leben.

Joh. 6, 37. Wer zu mir kommt, den werde ich nicht hinaus stoßen.

Joh. 6, 40. Das ist der Wille des, der mich gesandt hat, daß, wer den Sohn siehet, und glaubet an ihn, habe das ewige Leben, und ich werde ihn auferwecken am jüngsten Tage.

Joh. 6, 54. Wer mein Fleisch isset, und trinket mein Blut, der hat das ewige Leben, und ich werde ihn am jüngsten Tage auferwecken.

Joh. 6, 68. 69. Herr, wohin sollen wir gehen? Du hast Worte des ewigen Lebens; und wir haben geglaubt und erkannt, daß du bist Christus, der Sohn des lebendigen Gottes.

Joh. 8, 51. Wahrlich, wahrlich, ich sage euch, so jemand mein Wort wird halten, der wird den Tod nicht sehen ewiglich.

Joh. 10, 27. 28. Meine Schafe hören meine Stimme, und ich kenne sie, und sie folgen mir. Und ich gebe ihnen das ewige Leben; und sie werden nimmermehr umkommen, und niemand wird sie mir aus meiner Hand reißen.

Joh. 11, 25. 26. Jesus spricht zu ihr: Ich bin die Auferstehung und das Leben. Wer an mich glaubet, der wird leben, ob er gleich stürbe. Und wer da lebet und glaubet an mich, der wird nimmermehr sterben. Glaubest du das?

Joh. 12, 26. Wo ich bin, da soll mein Diener auch sein.

\quellseite{414}

Joh. 14, 6. JESus spricht: Ich bin der Weg, und die Wahrheit, und das Leben; niemand kommt zum Vater, denn durch mich.

Joh. 16, 16. Ueber ein Kleines, so werdet ihr mich nicht sehen; und aber über ein Kleines, so werdet ihr mich sehen; denn ich gehe zum Vater.

Joh. 16, 20. Wahrlich, wahrlich, ich sage euch, ihr werdet weinen und heulen; aber die Welt wird sich freuen. Ihr aber werdet traurig sein; doch eure Traurigkeit soll in Freude verkehret werden.

Joh. 16, 22. Ihr habt nun Traurigkeit; aber ich will euch wieder sehen, und euer Herz soll sich freuen, und eure Freude soll niemand von euch nehmen.

Joh. 16, 23. Wahrlich, wahrlich, ich sage euch, so ihr den Vater etwas bitten werdet in meinem Namen, so wird ers euch geben.

Joh. 16, 33. In der Welt habt ihr Angst; aber seid getrost, ich habe die Welt überwunden.

Joh. 17, 3. Das ist aber das ewige Leben, daß sie dich, daß du allein wahrer Gott bist, und den du gesandt hast, JESum Christum, erkennen.

Joh. 17, 24. Vater, ich will, daß, wo ich bin, auch die bey mir seyen, die du mir gegeben hast, daß sie meine Herrlichkeit sehen, die du mir gegeben hast; denn du hast mich geliebet, ehe denn die Welt gegründet ward.

Apostelg. 4, 12. Es ist in keinem andern Heil, ist auch kein anderer Name den Menschen gegeben, darinnen wir sollen selig werden.

Apostelg. 7, 58. Herr JESU, nimm meinen Geist auf!

Apostelg. 10, 43. Von diesem JESU zeugen alle Propheten, daß durch seinen Namen alle, die an ihn glauben, Vergebung der Sünden empfahen sollen.

Röm. 3, 23—25. Es ist hie kein Unterschied; sie sind allzumal Sünder, und mangeln des Ruhms, den sie an Gott haben sollten. Und werden ohne Verdienst gerecht aus seiner Gnade, durch die Erlösung, so durch Christum JESum geschehen ist; welchen Gott hat vorgestellet zu einem Gnadenstuhl durch den Glauben in seinem Blut.

\quellseite{415}

Röm. 3, 28. So halten wir es nun, daß der Mensch gerecht werde ohne des Gesetzes Werke, allein durch den Glauben.

Röm. 4, 5. Dem aber, der nicht mit Werken umgehet, glaubet aber an den, der die Gottlosen gerecht macht, dem wird sein Glaube gerechnet zur Gerechtigkeit.

Röm. 4, 25. Christus ist um unserer Sünden willen dahin gegeben, und um unserer Gerechtigkeit willen auferwecket.

Röm. 5, 1. 2. Nun wir denn sind gerecht worden durch den Glauben, so haben wir Friede mit Gott durch unsern Herrn JEsum Christ; durch welchen wir auch einen Zugang haben im Glauben zu dieser Gnade, darinnen wir stehen, und rühmen uns der Hoffnung der zukünftigen Herrlichkeit, die Gott geben soll.

Röm. 5, 20. Wo aber die Sünde mächtig worden ist, da ist doch die Gnade viel mächtiger worden.

Röm. 8, 1. So ist nun nichts Verdammliches an denen; die in Christo JEsu sind, die nicht nach dem Fleisch wandeln, sondern nach dem Geist.

Röm. 8, 17. 18. Sind wir denn Kinder, so sind wir auch Erben, nämlich Gottes Erben, und Miterben Christi, so wir anders mit leiden, auf daß wir auch mit zur Herrlichkeit erhoben werden. Denn ich halte es dafür, daß dieser Zeit Leiden der Herrlichkeit nicht werth sei, die an uns soll offenbar werden.

Röm. 8, 28—30. Wir wissen aber, daß denen, die Gott lieben, alle Dinge zum Besten dienen, die nach dem Vorsatz berufen sind. Denn welche er zuvor versehen hat, die hat er auch verordnet, daß sie gleich sein sollten dem Ebenbilde seines Sohns, auf daß derselbige der Erstgeborne sei unter vielen Brüdern. Welche er aber verordnet hat, die hat er auch berufen; welche er aber berufen hat, die hat er auch gerecht gemacht; welche er aber hat gerecht gemacht, die hat er auch herrlich gemacht.

Röm. 8, 31—34. Ist Gott für uns, wer mag wider uns sein? Welcher auch seines eigenen Sohns nicht hat verschonet, sondern hat ihn für uns alle dahin gegeben; wie sollt er uns mit ihm nicht
\quellseite{416}
alles schenken? Wer will die Auserwählten Gottes beschuldigen? Gott ist hie, der da gerecht macht. Wer will verdammen? Christus ist hie, der gestorben ist, ja vielmehr, der auch auferwecket ist, welcher ist zur Rechten Gottes, und vertritt uns.

Röm. 8, 37—39. In dem allen überwinden wir weit um deß willen, der uns geliebet hat. Denn ich bin gewiß, daß weder Tod noch Leben, weder Engel noch Fürstenthum, noch Gewalt, weder Gegenwärtiges noch Zukünftiges, weder Hohes noch Tiefes, noch keine andere Creatur mag uns scheiden von der Liebe Gottes, die in Christo Jesu ist, unserm HErrn.

Röm. 9, 16. So liegt es nun nicht an jemandes Wollen oder Laufen, sondern an Gottes Erbarmen.

Röm. 10, 4. Christus ist des Gesetzes Ende; wer an den glaubet, der ist gerecht.

Röm. 10, 12. 13. Es ist hie kein Unterschied unter Juden und Griechen; es ist aller zumal Ein HErr, reich über alle, die ihn anrufen. Denn wer den Namen des HErrn wird anrufen, soll selig werden.

Röm. 11, 32. Gott hat alles beschlossen unter den Unglauben, auf daß er sich aller erbarme.

Röm. 14, 8. 9. Leben wir, so leben wir dem HErrn; sterben wir, so sterben wir dem HErrn. Darum, wir leben oder sterben, so sind wir des HErrn. Denn dazu ist Christus auch gestorben, und auferstanden, und wieder lebendig worden, daß er über Todte und Lebendige HErr sei.

1 Cor. 1, 30. Christus ist uns gemacht von Gott zur Weisheit, und zur Gerechtigkeit, und zur Heiligung, und zur Erlösung.

1 Cor. 2, 9. Das kein Auge gesehen hat, und kein Ohr gehöret hat, und in keines Menschen Herz kommen ist, das Gott bereitet hat denen, die ihn lieben.

1 Cor. 10, 13. Gott ist getreu, der euch nicht lässet versuchen über euer Vermögen, sondern machet, daß die Versuchung so ein Ende gewinne, daß ihrs könnet ertragen.

\quellseite{417}

1 Cor. 15, 55—57. Der Tod ist verschlungen in den Sieg. Tod, wo ist dein Stachel? Hölle, wo ist dein Sieg? Aber der Stachel des Todes ist die Sünde; die Kraft aber der Sünde ist das Gesetz. Gott aber sei Dank, der uns den Sieg gegeben hat, durch unsern Herrn Jesum Christum.

2 Cor. 4, 17—18. Unsere Trübsal, die zeitlich und leicht ist, schaffet eine ewige und über alle Maaße wichtige Herrlichkeit, uns, die wir nicht sehen auf das Sichtbare, sondern auf das Unsichtbare. Denn was sichtbar ist, das ist zeitlich; was aber unsichtbar ist, das ist ewig.

2 Cor. 5, 1. Wir wissen aber, so unser irdisch Haus dieser Hütte zerbrochen wird, daß wir einen Bau haben von Gott erbauet, ein Haus, nicht mit Händen gemacht, das ewig ist, im Himmel.

2 Cor. 5, 21. Gott hat den, der von keiner Sünde wußte, für uns zur Sünde gemacht, auf daß wir würden in ihm die Gerechtigkeit, die vor Gott gilt.

2 Cor. 12, 9. Laß dir an meiner Gnade genügen; denn meine Kraft ist in den Schwachen mächtig.

Gal. 2, 16. Weil wir wissen, daß der Mensch durch des Gesetzes Werke nicht gerecht wird, sondern durch den Glauben an Jesum Christum: so glauben wir auch an Christum Jesum, auf daß wir gerecht werden durch den Glauben an Christum, und nicht durch des Gesetzes Werke; denn durch des Gesetzes Werke wird kein Fleisch gerecht.

Gal. 2, 20. Ich lebe aber; doch nun nicht ich, sondern Christus lebet in mir. Denn was ich jetzt lebe im Fleisch, das lebe ich in dem Glauben des Sohns Gottes, der mich geliebet hat, und sich selbst für mich dargegeben.

Gal. 3, 13. Christus hat uns erlöset von dem Fluch des Gesetzes, da er ward ein Fluch für uns.

Gal. 3, 26. 27. Ihr seyd alle Gottes Kinder durch den Glauben an Christum Jesum. Denn wie viel euer getauft sind, die haben Christum angezogen.

\quellseite{418}

Gal. 4, 4. 5. Da die Zeit erfüllet ward, sandte Gott seinen Sohn, geboren von einem Weibe, und unter das Gesetz gethan, auf daß er die, so unter dem Gesetz waren, erlöste, daß wir die Kindschaft empfingen.

Gal. 6, 14. Es sey ferne von mir rühmen, denn allein von dem Kreuz unsers HErrn JEsu Christi, durch welchen mir die Welt gekreuziget ist, und ich der Welt.

Ephes. 2, 8. 9. Aus Gnaden seid ihr selig worden durch den Glauben; und dasselbige nicht aus euch, Gottes Gabe ist es; nicht aus den Werken, auf daß sich nicht jemand rühme.

Ephes. 6, 10—17. Zuletzt, meine Brüder, seid stark in dem HErrn, und in der Macht seiner Stärke. Ziehet an den Harnisch Gottes, daß ihr bestehen könnet gegen die listigen Anläufe des Teufels. Denn wir haben nicht mit Fleisch und Blut zu kämpfen, sondern mit Fürsten und Gewaltigen, nämlich mit den Herren der Welt, die in der Finsterniß dieser Welt herrschen, mit den bösen Geistern unter dem Himmel. Um deßwillen, so ergreifet den Harnisch Gottes, auf daß ihr an dem bösen Tage Widerstand thun, und alles wohl ausrichten, und das Feld behalten möget. So stehet nun, umgürtet eure Lenden mit Wahrheit, und angezogen mit dem Krebs der Gerechtigkeit, und an Beinen geschiefelt, als fertig zu treiben das Evangelium des Friedens, damit ihr bereitet seid. Vor allen Dingen aber ergreifet den Schild des Glaubens, mit welchem ihr auslöschen könnt alle feurigen Pfeile des Bösewichtes. Und nehmet den Helm des Heils, und das Schwert des Geistes, welches ist das Wort Gottes.

Phil. 1, 6. Ich bin desselbigen in guter Zuversicht, daß der in euch angefangen hat das gute Werk, der wirds auch vollführen bis an den Tag JEsu Christi.

Phil. 1, 21. Christus ist mein Leben, und Sterben ist mein Gewinn.

Phil. 1, 23. Ich habe Lust abzuscheiden, und bei Christo zu sein.

Phil. 3, 20. 21. Unser Wandel aber ist im Himmel, von dannen wir auch warten des Heilands, JEsu Christi, des HErrn,
\quellseite{419}
welcher unsern nichtigen Leib verklären wird, daß er ähnlich werde seinem verklärten Leibe, nach der Wirkung, damit er kann auch alle Dinge ihm unterthänig machen.

Phil. 4, 7. Der Friede Gottes, welcher höher ist, denn alle Vernunft, bewahre eure Herzen und Sinne in Christo JEsu!

Col. 1, 12—14. Danksaget dem Vater, der uns tüchtig gemacht hat zu dem Erbtheil der Heiligen im Licht, welcher uns errettet hat von der Obrigkeit der Finsterniß, und hat uns versetzt in das Reich seines lieben Sohnes, an welchem wir haben die Erlösung durch sein Blut, nämlich die Vergebung der Sünden.

Col. 3, 3. 4. Euer Leben ist verborgen mit Christo in Gott. Wenn aber Christus, euer Leben, sich offenbaren wird, dann werdet ihr auch offenbar werden mit ihm in der Herrlichkeit.

1 Thess. 5, 23. 24. Der Gott des Friedens heilige euch durch und durch, und euer Geist ganz, sammt der Seele und Leib, müsse behalten werden unsträflich auf die Zukunft unsers HErrn JEsu Christi! Getreu ist er, der euch rufet, welcher wirds auch thun.

1 Tim. 1, 15. 16. Das ist je gewißlich wahr, und ein theuer werthes Wort, daß Christus JEsus kommen ist in die Welt, die Sünder selig zu machen, unter welchen ich der vornehmste bin. Aber darum ist mir Barmherzigkeit widerfahren, auf daß an mir vornehmlich JEsus Christus erzeigte alle Geduld, zum Exempel denen, die an ihn glauben sollten zum ewigen Leben.

1 Tim. 2, 4—6. Gott will, daß allen Menschen geholfen werde, und zur Erkenntniß der Wahrheit kommen. Denn es ist Ein Gott und Ein Mittler zwischen Gott und den Menschen, nämlich der Mensch Christus JEsus, der sich selbst gegeben hat für alle zur Erlösung, daß solches zu seiner Zeit geprediget würde.

1 Tim. 3, 16. Kündlich groß ist das gottselige Geheimniß: Gott ist offenbaret im Fleisch, gerechtfertiget im Geist, erschienen den Engeln, geprediget den Heiden, geglaubet von der Welt, aufgenommen in die Herrlichkeit.

2 Tim. 2, 8. Halt im Gedächtniß JEsum Christum, der auferstanden ist von den Todten.

\quellseite{420}

2 Tim. 2, 11—13. Das ist je gewißlich wahr: Sterben wir mit, so werden wir mit leben; dulden wir, so werden wir mit herrschen; verleugnen wir, so wird er uns auch verleugnen; glauben wir nicht, so bleibt er treu; er kann sich selbst nicht leugnen.

2 Tim. 4, 7. 8. Ich habe einen guten Kampf gekämpfet; ich habe den Lauf vollendet; ich habe Glauben gehalten. Hinfort ist mir beigelegt die Krone der Gerechtigkeit, welche mir der HErr an jenem Tage, der gerechte Richter, geben wird, nicht mir aber allein, sondern auch allen, die seine Erscheinung lieb haben.

2 Tim. 4, 18. Der HErr wird mich erlösen von allem Uebel, und aushelfen zu seinem himmlischen Reich; welchem sei Ehre von Ewigkeit zu Ewigkeit.

Tit. 3, 5—7. Nicht um der Werke willen der Gerechtigkeit, die wir gethan hatten, sondern nach seiner Barmherzigkeit machte er uns selig durch das Bad der Wiedergeburt und Erneuerung des Heiligen Geistes, welchen er ausgegossen hat über uns reichlich durch JEsum Christum, unsern Heiland, auf daß wir durch desselbigen Gnade gerecht, und Erben seien des ewigen Lebens nach der Hoffnung.

1 Pet. 1, 3—9. Gelobet sei Gott, und der Vater unsers HErrn JEsu Christi, der uns nach seiner großen Barmherzigkeit wiedergeboren hat zu einer lebendigen Hoffnung durch die Auferstehung JEsu Christi von den Todten, zu einem unvergänglichen, und unbefleckten, und unverwelklichen Erbe, das behalten wird im Himmel, euch, die ihr aus Gottes Macht durch den Glauben bewahret werdet zur Seligkeit, welche zubereitet ist, daß sie offenbar werde zu der letzten Zeit, in welcher ihr euch freuen werdet, die ihr jetzt eine kleine Zeit (wo es sein soll) traurig seid in mancherlei Anfechtungen, auf daß euer Glaube rechtschaffen und viel köstlicher erfunden werde, denn das vergängliche Gold, das durchs Feuer bewähret wird, zu Lobe, Preis und Ehren, wenn nun offenbaret wird JEsus Christus, welchen ihr nicht gesehen und doch lieb habt, und nun an ihn glaubet, wiewohl ihr ihn nicht sehet, so werdet ihr euch freuen mit unaussprechlicher und herrlicher Freude,
\quellseite{421}
und das Ende eures Glaubens davon bringen, nämlich der Seelen Seligkeit.

1 Pet. 1, 18. 19. Wisset, daß ihr nicht mit vergänglichem Silber oder Gold erlöset seid von eurem eiteln Wandel nach väterlicher Weise, sondern mit dem theuren Blut Christi, als eines unschuldigen und unbefleckten Lammes.

1 Pet. 1, 24. 25. Alles Fleisch ist wie Gras, und alle Herrlichkeit der Menschen wie des Grases Blume. Das Gras ist verdorret, und die Blume abgefallen; aber des Herrn Wort bleibet in Ewigkeit.

1 Pet. 2, 24. 25. Christus hat unsere Sünden selbst geopfert an seinem Leibe auf dem Holz, auf daß wir, der Sünde abgestorben, der Gerechtigkeit leben; durch welches Wunden ihr seid heil worden. Denn ihr waret wie die irrenden Schafe; aber ihr seid nun bekehret zu dem Hirten und Bischof eurer Seelen.

2 Pet. 1, 19. Wir haben ein festes prophetisches Wort, und ihr thut wohl, daß ihr drauf achtet, als auf ein Licht, das da scheinet in einem dunkeln Ort, bis der Tag anbreche, und der Morgenstern aufgehe in euren Herzen.

2 Pet. 3, 9. Gott hat Geduld mit uns, und will nicht, daß jemand verloren werde, sondern daß sich jedermann zur Buße kehre.

Pet. 3, 13. Wir warten aber eines neuen Himmels, und einer neuen Erde, nach seiner Verheißung, in welchen Gerechtigkeit wohnet.

1 Joh. 1, 7. Das Blut Jesu Christi macht uns rein von aller Sünde.

1 Joh. 1, 8. 9. So wir sagen, wir haben keine Sünde, so verführen wir uns selbst, und die Wahrheit ist nicht in uns. So wir aber unsere Sünde bekennen, so ist er treu und gerecht, daß er uns die Sünde vergibt, und reiniget uns von aller Untugend.

2 Joh. 2, 1. 2. Ob jemand sündiget, so haben wir einen Fürsprecher bei dem Vater, Jesum Christum, der gerecht ist. Und derselbige
\quellseite{422}
ist die Versöhnung für unsere Sünde; nicht allein aber für die unsere, sondern auch für der ganzen Welt.

1 Joh. 2, 25. Das ist die Verheißung, die er uns verheißen hat, das ewige Leben.

1 Joh. 3, 1. 2. Sehet, welch eine Liebe hat uns der Vater erzeiget, daß wir Gottes Kinder sollen heißen. Darum kennet euch die Welt nicht, denn sie kennet ihn nicht. Meine Lieben, wir sind nun Gottes Kinder, und ist noch nicht erschienen, was wir sein werden. Wir wissen aber, wenn es erscheinen wird, daß wir ihm gleich sein werden; denn wir werden ihn sehen, wie er ist.

1 Joh. 3, 20. So uns unser Herz verdammt, daß Gott größer ist, denn unser Herz, und erkennet alle Dinge.

1 Joh. 4, 9. Daran ist erschienen die Liebe Gottes gegen uns, daß Gott seinen eingebornen Sohn gesandt hat in die Welt, daß wir durch ihn leben sollen.

Hebr. 2, 14. 15. Nachdem die Kinder Fleisch und Blut haben, ist er gleichermaßen theilhaftig worden, auf daß er durch den Tod die Macht nähme dem, der des Todes Gewalt hatte, das ist, dem Teufel, und erlösete die, so durch Furcht des Todes im ganzen Leben Knechte sein mußten.

Hebr. 4, 9. Es ist noch eine Ruhe vorhanden dem Volk Gottes.

Hebr. 4, 15. 16. Wir haben nicht einen Hohenpriester, der nicht könnte Mitleiden haben mit unserer Schwachheit, sondern der versucht ist allenthalben gleichwie wir, doch ohne Sünde. Darum lasset uns hinzu treten mit Freudigkeit zu dem Gnadenstuhl, auf daß wir Barmherzigkeit empfahen, und Gnade finden auf die Zeit, wenn uns Hülfe noth sein wird.

Hebr. 7, 24. 25. Christus hat ein unvergänglich Priesterthum. Daher er auch selig machen kann immerdar, die durch ihn zu Gott kommen, und lebet immerdar, und bittet für sie.

Hebr. 12, 6—8. Welchen der Herr lieb hat, den züchtiget er; er stäupt aber einen jeglichen Sohn, den er aufnimmt. So ihr
\quellseite{423}
die Züchtigung erduldet, so erbeut sich euch Gott als Kindern; denn wo ist ein Sohn, den der Vater nicht züchtiget? Seid ihr aber ohne Züchtigung, welcher sie alle sind theilhaftig worden; so seid ihr Bastarde, und nicht Kinder.

Hebr. 12, 11. Alle Züchtigung aber, wenn sie da ist, dünkt sie uns nicht Freude, sondern Traurigkeit sein; aber darnach wird sie geben eine friedsame Frucht der Gerechtigkeit denen, die dadurch geübet sind.

Hebr. 12, 22—24. Ihr seid kommen zu dem Berge Zion, und zu der Stadt des lebendigen Gottes, zu dem himmlischen Jerusalem, und zu der Menge vieler tausend Engel, und zu der Gemeinde der Erstgebornen, die im Himmel angeschrieben sind, und zu Gott, dem Richter über alle, und zu den Geistern der vollkommnen Gerechten, und zu dem Mittler des neuen Testaments, Jesu, und zu dem Blut der Besprengung, das da besser redet, denn Abels:

Hebr. 13, 5. Er hat gesagt: Ich will dich nicht verlassen, noch versäumen.

Hebr. 13, 8. Jesus Christus, gestern und heute, und derselbe auch in Ewigkeit.

Jac. 1, 12. Selig ist der Mann, der die Anfechtung erduldet; denn nachdem er bewähret ist, wird er die Krone des Lebens empfahen, welche Gott verheißen hat denen, die ihn lieb haben.

Offenb. 2, 10. Sei getreu bis an den Tod, so will ich dir die Krone des Lebens geben.

Offenb. 3, 11. 12. Siehe, ich komme bald. Halte, was du hast, daß niemand deine Krone nehme. Wer überwindet, den will ich machen zum Pfeiler in dem Tempel meines Gottes, und soll nicht mehr hinaus gehen. Und will auf ihn schreiben den Namen meines Gottes, und den Namen des neuen Jerusalem, der Stadt meines Gottes, die vom Himmel hernieder kommt, von meinem Gott, und meinen Namen, den neuen.

Offenb. 3, 20. 21. Siehe, ich stehe vor der Thür, und klopfe an. So jemand meine Stimme hören wird, und die Thür aufthun, zu dem werde ich eingehen, und das Abendmahl mit ihm
\quellseite{424}
halten, und er mit mir. Wer überwindet, dem will ich geben, mit mir auf meinem Stuhl zu sitzen; wie ich überwunden habe, und bin gesessen mit meinem Vater auf seinem Stuhl.

Offenb. 7, 14. Diese sinds, die kommen sind aus großer Trübsal, und haben ihre Kleider gewaschen, und haben ihre Kleider helle gemacht im Blut des Lammes.

Offenb. 12, 10—12. Ich hörte eine große Stimme, die sprach im Himmel: Nun ist das Heil und die Kraft, und das Reich, und die Macht unsers Gottes, seines Christus worden, weil der Verkläger unserer Brüder verworfen ist, der sie verklaget Tag und Nacht vor Gott. Und sie haben ihn überwunden durch des Lammes Blut, und durch das Wort ihres Zeugnisses; und haben ihr Leben nicht geliebet bis an den Tod. Darum freuet euch, ihr Himmel, und die darinnen wohnen.

Offenb. 14, 13. Selig sind die Todten, die in dem HErrn sterben, von nun an. Ja, der Geist spricht, daß sie ruhen von ihrer Arbeit; denn ihre Werke folgen ihnen nach.

Offenb. 19, 9. Selig sind, die zum Abendmahl des Lammes berufen sind.

Offenb. 21, 7. Wer überwindet, der wirds alles ererben; und ich werde sein Gott sein, und er wird mein Sohn sein.

Offenb. 22, 20. 21. Es spricht, der solches zeuget: Ja, ich komme bald. Amen. Ja, komm, HErr JEsu. Die Gnade unsers HErrn JEsu Christi sei mit euch allen! Amen.

\section{459. Lieder und Liederverse}

\subsection{1}

\begin{verse}
Wenn ich einmal soll scheiden,\\
So scheide nicht von mir;\\
Wenn ich den Tod soll leiden,\\
So tritt du dann herfür:\\
Wenn mir am allerbängsten\\
Wird um das Herze sein,
\end{verse}

\quellseite{425}

\begin{verse}
So reiß mich aus den Aengsten,\\
Kraft deiner Angst und Pein.\\
Erscheine mir zum Schilde,\\
Zum Trost in meinem Tod,\\
Und laß mich sehn dein Bilde\\
In deiner Kreuzesnoth;\\
Da will ich nach dir blicken,\\
Da will ich glaubensvoll\\
Dich fest an mein Herz drücken;\\
Wer so stirbt, der stirbt wohl!
\end{verse}

2.

\begin{verse}
Aus tiefer Noth schrei’ ich zu dir,\\
HErr Gott! erhör’ mein Rufen!\\
Dein gnädig Ohr neig her zu mir,\\
Und meiner Bitt’ sie öffne:\\
Denn so du willst das sehen an,\\
Was Sünd’ und Unrecht ist gethan,\\
Wer kann, HErr! für dir bleiben?
\end{verse}

\begin{verse}
Bei dir gilt nichts denn Gnad’ und Gunst,\\
Die Sünde zu vergeben;\\
Es ist doch unser Thun umsonst,\\
Auch in dem besten Leben:\\
Für dir niemand sich rühmen kann,\\
Deß muß dich fürchten jedermann,\\
Und deiner Gnade leben.
\end{verse}

\begin{verse}
Darum auf Gott will hoffen ich,\\
Auf mein Verdienst nicht bauen;\\
Auf ihn mein Herz soll lassen sich,\\
Und seiner Güte trauen,\\
Die mir zusagt sein werthes Wort,\\
Das ist mein Trost und treuer Hort,\\
Deß will ich all’zeit harren.
\end{verse}

Und ob es währet bis in die Nacht, Und wieder an den Morgen,

\quellseite{426}

\begin{verse}
Doch soll mein Herz an Gottes Macht\\
Verzweifeln nicht, noch sorgen.\\
So thu’ Israel rechter Art,\\
Der aus dem Geist erzeuget ward,\\
Und seines Gott’s erharre.
\end{verse}

\begin{verse}
Ob bei uns ist der Sünden viel,\\
Bei Gott ist viel mehr Gnaden,\\
Sein’ Hand zu helfen hat kein Ziel,\\
Wie groß auch sei der Schaden.\\
Er ist allein der gute Hirt,\\
Der Israel erlösen wird\\
Aus seinen Sünden allen.
\end{verse}

\subsection{3}

\begin{verse}
Warum sollt’ ich mich denn grämen?\\
Hab’ ich doch\\
Christum noch,\\
Wer will mir den nehmen?\\
Wer will mir den Himmel rauben?\\
Den mir schon\\
Gottes Sohn\\
Beigelegt im Glauben.
\end{verse}

\begin{verse}
Nackends lag ich auf dem Boden,\\
Da ich kam,\\
Da ich nahm\\
Meinen ersten Odem;\\
Nackends werd’ ich auch hinziehn,\\
Wenn ich werd’\\
Von der Erd’\\
Als ein Schatten fliehn.
\end{verse}

\begin{verse}
Gut und Blut, Leib, Seel’ und Leben\\
Ist nicht mein,\\
Gott allein\\
Ist es, der’s gegeben:\\
Will er’s wieder zu sich kehren,
\end{verse}

\quellseite{427}

\begin{verse}
Nehm’ er’s hin,\\
Ich will ihn\\
Dennoch fröhlich ehren.
\end{verse}

\begin{verse}
Kann uns doch kein Tod nicht tödten,\\
Sondern reißt\\
Unfern Geist\\
Aus viel tausend Nöthen,\\
Schleußt das Thor der bittern Leiden,\\
Und macht Bahn,\\
Da man kann\\
Gehn zur Himmelsfreuden.
\end{verse}

\begin{verse}
Also will in süßen Schätzen\\
Ich mein Herz\\
Auf den Schmerz\\
Ewiglich ergötzen.\\
Hier ist kein recht Gut zu finden;\\
Was die Welt\\
In sich hält,\\
Muß im Hui verschwinden.
\end{verse}

\begin{verse}
Herr, mein Hirt, Brunn’ aller Freuden,\\
Du bist mein,\\
Ich bin dein,\\
Niemand kann uns scheiden.\\
Ich bin dein, weil du dein Leben\\
Und dein Blut\\
Mir zu gut\\
In den Tod gegeben.
\end{verse}

\begin{verse}
Du bist mein, weil ich dich fasse,\\
Und dich nicht\\
O mein Licht!\\
Aus dem Herzen lasse.\\
Laß mich, laß mich hingelangen,\\
Da du mich,\\
Und ich dich\\
Lieblich werd’ umfangen.
\end{verse}

\quellseite{428}

\subsection{4}

\begin{verse}
In Christi Wunden schlaf ich ein,\\
Die machen mich von Sünden rein:\\
Ja! Christi Blut und Gerechtigkeit,\\
Das ist mein Schmuck und Ehrenkleid:\\
Damit will ich vor Gott bestehn,\\
Wenn ich zum Himmel werd’ eingehen.
\end{verse}

\subsection{5}

\begin{verse}
Wer weiß, wie nahe mir mein Ende,\\
Hin geht die Zeit, her kömmt der Tod:\\
Ach! wie geschwinde und behende\\
Kann kommen meine Todesnoth!\\
Mein Gott! ich bitt’ durch Christi Blut,\\
Mach’s nur mit meinem Ende gut.
\end{verse}

\subsection{6}

\begin{verse}
Wenn mein Stündlein vorhanden ist,\\
Und ich soll fahr’n mein Straße,\\
So g’leit du mich, Herr Jesu Christ!\\
Mit Hülff’ mich nicht verlasse;\\
Mein’ Seel’ an meinem letzten End’\\
Befeh?’ ich, Herr! in deine Händ’,\\
Du wirst sie wohl bewahren.

Mein’ Sünd’ mich werden kränken sehr,\\
Mein G’wissen wird mich nagen,\\
Denn ihr’re sind viel, wie Sand am Meer,\\
Doch will ich nicht verzagen:\\
Gedenken will ich an dein’n Tod,\\
Herr Jesu! deine Wunden roth,\\
Die werden mich erhalten.

Ich bin ein Glied an deinem Leib,\\
Deß tröst’ ich mich von Herzen,\\
Von dir ich ungeschieden bleib’\\
In Todesnoth und Schmerzen.
\end{verse}

\quellseite{429}

\begin{verse}
Wenn ich gleich sterb’, so sterb’ ich dir,\\
Ein ew’ges Leben hast du mir\\
Mit deinem Tod erworben.
\end{verse}

\begin{verse}
Weil du vom Tod erstanden bist,\\
Werd’ ich im Grab nicht bleiben;\\
Mein höchster Trost dein Auffahrt ist,\\
Tod’sfurcht kann sie vertreiben;\\
Denn wo du bist, da komm’ ich hin,\\
Daß ich stets bei dir leb’ und bin,\\
Drum fahr’ ich hin mit Freuden.
\end{verse}

\begin{verse}
So fahr’ ich hin zu Jesu Christ,\\
Mein’n Arm thu’ ich ausstrecken;\\
So schlaf’ ich ein und ruhe fein,\\
Kein Mensch kann mich aufwecken;\\
Denn Jesus Christus, Gottes Sohn,\\
Der wird die Himmelsthür aufthun,\\
Mich führ’n zum ew’gen Leben.
\end{verse}

\subsection{7}

\begin{verse}
Bald will ich dir geben,\\
Du arge falsche Welt!\\
Dein sündlich böses Leben\\
Durchaus mir nicht gefällt,\\
Im Himmel ist gut wohnen,\\
Hinauf steht mein Begier;\\
Da wird Gott ewig lohnen\\
Dem, der ihm dient allhier.
\end{verse}

\begin{verse}
Rath’ mir nach deinem Herzen,\\
O Jesu, Gottes Sohn!\\
Soll ich hie dulden Schmerzen,\\
Hilf mir, Herr Christ! davon;\\
Verkürz’ mir alles Leiden,\\
Tröst’ meinen blöden Muth;\\
Laß mich selig abscheiden,\\
Setz’ mich in dein Erbgut.
\end{verse}

\quellseite{430}

\begin{verse}
In meines Herzens Grunde\\
Dein Nam’ und Kreuz allein\\
Funkelt all’ Zeit und Stunde;\\
Drauf kann ich fröhlich sein.\\
Erschein’ mir in dem Bilde,\\
Zu Trost in meiner Noth,\\
Wie du, Herr Christ! so milde\\
Dich hast geblut’t zu Tod.
\end{verse}

\begin{verse}
Verbürg’ mein Seel’ aus Gnaden\\
In deine offne Seit’,\\
Rück’ sie aus allem Schaden\\
Zu deiner Herrlichkeit!\\
Der ist wohl hie gewesen,\\
Der kömmt in’s Himmels Schloß,\\
Der ist ewig genesen,\\
Der bleibt in deinem Schooß.
\end{verse}

\begin{verse}
Schreib’ meinen Nam’n auf’s Beste\\
In’s Buch des Lebens ein,\\
Und bind’ mein’ Seel’ fein feste\\
In’s schöne Bündelein\\
Der’r, die im Himmel grünen,\\
Und vor dir leben frei,\\
So will ich ewig rühmen,\\
Daß dein Herz treue sei.
\end{verse}

\subsection{8}

\begin{verse}
Christe, du Lamm Gottes!\\
Der du trägst die Sünde der Welt,\\
Erbarm’ dich unser!
\end{verse}

\begin{verse}
Christe, du Lamm Gottes!\\
Der du trägst die Sünde der Welt,\\
Erbarm’ dich unser!
\end{verse}

\begin{verse}
Christe, du Lamm Gottes!\\
Der du trägst die Sünde der Welt,\\
Gib uns deinen Frieden! Amen.
\end{verse}

\quellseite{431}

\subsection{9}

\begin{verse}
O Lamm Gottes unschuldig!\\
Am Stamm des Kreuzes geschlachtet,\\
Allzeit erfunden geduldig,\\
Wiewohl du warest verachtet,\\
Au’ Sünd’ hast du getragen,\\
Sonst müßten wir verzagen.\\
Erbarm’ dich unser, o Jesu! o Jesu!
\end{verse}

\begin{verse}
O Lamm Gottes unschuldig!\\
Am Stamm des Kreuzes geschlachtet,\\
Allzeit erfunden geduldig,\\
Wiewohl du warest verachtet,\\
Au’ Sünd’ hast du getragen,\\
Sonst müßten wir verzagen.\\
Erbarm’ dich unser, o Jesu! o Jesu!
\end{verse}

\begin{verse}
O Lamm Gottes unschuldig!\\
Am Stamm des Kreuzes geschlachtet,\\
Allzeit erfunden geduldig,\\
Wiewohl du warest verachtet,\\
Au’ Sünd’ hast du getragen,\\
Sonst müßten wir verzagen.\\
Gib uns deinen Frieden, o Jesu! o Jesu!
\end{verse}

\subsection{10}

\begin{verse}
Stärk’ mich mit deinem Freudgeist,\\
Heil’ mich mit deinen Wunden,\\
Wasch’ mich mit deinem Todesschweiß\\
In meiner letzten Stunden,\\
Und nimm mich einst, wann dir’s gefällt,\\
In wahrem Glauben aus der Welt\\
Zu deinen Auserwählten.
\end{verse}

\subsection{11}

Hilf, Helfer, hilf in Angst und Noth, Erbarm dich mein, o treuer Gott!

\quellseite{432}

\begin{verse}
Ich bin ja doch dein liebes Kind,\\
Trotz Teufel, Welt und aller Sünd’.

Ich trau’ auf dich, mein Gott und Herr!\\
Wenn ich dich hab’, was will ich mehr?\\
Ich hab’ ja dich, Herr Jesu Christ!\\
Du mein Gott und Erlöser bist.

Deß freu’ ich mich von Herzen fein,\\
Bin gutes Muths und harre dein,\\
Verlaß mich gänzlich auf dein’n Nam’n;\\
Hilf, Helfer, hilf! drauf sprech ich: Am’n.
\end{verse}

12.

\begin{verse}
Schwing’ dich auf zu deinem Gott,\\
Du betrübte Seele!\\
Warum liegst du, Gott zum Spott,\\
In der Schwermuthshöhle?\\
Merkst du nicht des Satans List?\\
Er will durch sein Kämpfen\\
Deinen Trost, den Jesus Christ\\
Dir erworben, dämpfen.

Schüttle deinen Kopf und sprich:\\
Fleuch, du alte Schlange!\\
Was ernn’st du deinen Stich,\\
Machst mir angst und bange?\\
Ist dir doch der Kopf zerknickt,\\
Und ich bin durch’s Leiden\\
Meines Heiland’s dir entrückt\\
In den Saal der Freuden.

Wirfst du mir mein Sünd’gen für:\\
Wo hat Gott befohlen,\\
Daß mein Urtheil über mir\\
Ich bei dir soll holen:\\
Wer hat dir die Macht geschenkt,\\
Andre zu verdammen?\\
Der du selbst doch liegst versenkt\\
In der Höllen Flammen.
\end{verse}

\quellseite{433}

\begin{verse}
Hab’ ich was nicht recht gethan,\\
Ist mir’s leid von Herzen.\\
Da hingegen nehm ich an\\
Christi Tod und Schmerzen:\\
Denn das ist die Ranzion\\
Meiner Missethaten:\\
Bring’ ich dies vor Gottes Thron,\\
Ist mir wohl gerathen.
\end{verse}

\begin{verse}
Gottes Kinder säen zwar\\
Traurig und mit Thränen:\\
Aber endlich bringt das Jahr,\\
Wonach sie sich sehnen;\\
Denn es kömmt die Erntezeit,\\
Da sie Garben machen,\\
Da wird all’ ihr Gram und Leid\\
Lauter Freud’ und Lachen.
\end{verse}

\begin{verse}
Ei, so faß’, o Christenherz,\\
Alle deine Schmerzen,\\
Wird sie fröhlich hinterwärts,\\
Laß des Trostes Kerzen\\
Dich entzünden mehr und mehr,\\
Gib dem großen Namen\\
Deines Gottes Preis und Ehr’:\\
Er wird helfen, Amen.
\end{verse}

13.

\begin{verse}
Sollt’ es gleich bisweilen scheinen,\\
Als wenn Gott verließ die Seinen!\\
Ei! so glaub’ und weiß ich dies:\\
Gott hilft endlich doch gewiß.
\end{verse}

\begin{verse}
Hülfe, die er aufgeschoben,\\
Hat er drum nicht aufgehoben,\\
Hilft er nicht zu jeder Frist,\\
Hilft er doch, wenn’s nöthig ist.
\end{verse}

\quellseite{434}

\subsection{14}

\begin{verse}
HErr JEsu Christ, wahr’r Mensch und Gott!\\
Der du litt’st Marter, Angst und Spott,\\
Für mich am Kreuz auch endlich starbst,\\
Und mir dein’s Vaters Huld erwarbst.
\end{verse}

\begin{verse}
Ich bitt’ durch’s bitt’re Leiden dein,\\
Du woll’st mir Sünder gnädig sein,\\
Wenn ich nun komm’ in Sterbensnoth,\\
Und ringen werde mit dem Tod.
\end{verse}

\begin{verse}
Wenn mir vergeht all’ mein Gesicht,\\
Und meine Ohren hören nicht,\\
Wenn meine Zunge nicht mehr spricht,\\
Und mir für Angst mein Herz zerbricht.
\end{verse}

\begin{verse}
Wenn mein Verstand sich nicht besinnt,\\
Und mir all’ menschlich’ Hülff’ zerrinnt,\\
So komm, o HErr Christ! mir behend’\\
Zu Hülff’ an meinem letzten End’.
\end{verse}

\begin{verse}
Und führ’ mich aus dem Jammerthal;\\
Verkürz’ mir auch des Todes Qual;\\
Die bösen Geister von mir treib’,\\
Mit deinem Geist stets bei mir bleib’.
\end{verse}

\begin{verse}
Bis sich die Seel’ vom Leib abwend’t,\\
So nimm sie, HErr! in deine Händ’,\\
Der Leib hab’ in der Erd’ sein’ Ruh’,\\
Bis sich der jüngst’ Tag naht herzu.
\end{verse}

\begin{verse}
Ein fröhlich’ Auf’ erstehn mir verleih’,\\
Am jüngsten G’richt mein Fürsprech’r sei,\\
Und meiner Sünd’ nicht mehr gedenk’,\\
Aus Gnaden mir das Leben schenk’.
\end{verse}

\begin{verse}
Wie du hast zugesaget mir\\
In deinem Wort, das trau’ ich dir:\\
„Fürwahr, fürwahr, euch sage ich:\\
Wer mein Wort hält und glaubt an mich,
\end{verse}

\quellseite{435}

\begin{verse}
Der wird nicht kommen in’s Gericht,\\
Auch den Tod ewig schmecken nicht,\\
Und ob er gleich hie zeitlich stirbt,\\
Mit nichten er drum gar verdirbt;
\end{verse}

\begin{verse}
Sondern ich will mit starker Hand\\
Ihn reißen aus des Todes Band’,\\
Und zu mir nehmen in mein Reich,\\
Da soll er denn mit mir zugleich\\
In Freuden leben ewiglich.”\\
Dazu hilf uns ja gnädiglich!\\
Ach Herr! vergib all’ unsre Schuld,\\
Hilf! daß wir warten mit Geduld,
\end{verse}

\begin{verse}
Bis unser Stündlein kömmt herbei,\\
Auch unser Glaub’ stets wacker sei,\\
Dein’m Wort zu trauen festiglich,\\
Bis wir einschlafen seliglich.
\end{verse}

heading

\begin{verse}
Mein zu dir, HErr Jesu Christ!\\
Mein Hoffnung steht auf Erden,\\
Ich weiß, daß du mein Tröster bist,\\
Kein Trost mag mir sonst werden.\\
Von Anbeginn ist nichts erkor’n,\\
Auf Erden ward kein Mensch gebor’n,\\
Der mir aus Nöthen helfen kann,\\
Ich ruf’ dich an,\\
Zu dem ich mein Vertrauen han.
\end{verse}

\begin{verse}
Mein’ Sünd’ sind schwer und übergroß,\\
Und reuen mich von Herzen,\\
Derselben mach’ mich quitt und los\\
Durch deinen Tod und Schmerzen,\\
Und zeig’ mich deinem Vater an,\\
Daß du hast g’nug für mich gethan,\\
So werd’ ich quitt der Sündenlast.
\end{verse}

\quellseite{436}

\begin{verse}
O Herr! halt mir fest,\\
Weß du dich mir versprochen hast.\\
Gib mir nach dein’r Barmherzigkeit\\
Den wahren Christenglauben,\\
Auf daß ich deine Süßigkeit\\
Mög’ inniglich schauen:\\
Für allen Dingen lieben dich,\\
Und meinen Nächsten gleich als mich:\\
Am letzten End’ dein’ Hülf’ mir send’,\\
Thun’ mir behend’,\\
Des Teufels List sich von mir wend’.
\end{verse}

\begin{verse}
Ehr’ sei Gott in dem höchsten Thron,\\
Dem Vater aller Güte,\\
Und Jesu Christ, sein’m liebsten Sohn,\\
Der uns allzeit behüte,\\
Und Gott dem Heil’gen Geiste,\\
Der uns sein’ Hülf’ all’zeit leiste,\\
Damit wir ihm gefällig sein,\\
Hier in der Zeit\\
Und folgend in der Ewigkeit.
\end{verse}

16.

\begin{verse}
Gott der Vater wohn’ uns bei\\
Und laß uns nicht verderben; †;\\
Mach’ uns aller Sünden frei,\\
Und hilf uns selig sterben!\\
Für dem Teufel uns bewahr’,\\
Halt uns bei festem Glauben,\\
Und auf dich laß uns bauen,\\
Aus Herzensgrund vertrauen;\\
Dir uns lassen ganz und gar,\\
Mit allen rechten Christen,\\
Entfliehen Teufels Listen,\\
Mit Waffen Gott’s uns rüsten.\\
Amen, Amen, das sei wahr,\\
So singen wir: Halleluja.
\end{verse}

\quellseite{437}

\begin{verse}
O Jesu Christus wohn’ uns bei\\
Und laß uns nicht verderben ;:\\
Mach’ uns aller Sünden frei,\\
Und hilf uns selig sterben!\\
Für dem Teufel uns bewahr’,\\
Halt uns bei festem Glauben,\\
Und auf dich laß uns bauen,\\
Aus Herzensgrund vertrauen;\\
Dir uns lassen ganz und gar,\\
Mit allen rechten Christen,\\
Entfliehen Teufels Listen,\\
Mit Waffen Gott’s uns rüsten.\\
Amen, Amen, das sei wahr,\\
So singen wir: Halleluja.
\end{verse}

\begin{verse}
Heilig’ Geist der wohn’ uns bei\\
Und laß uns nicht verderben ;:\\
Mach’ uns aller Sünden frei,\\
Und hilf uns selig sterben!\\
Für dem Teufel uns bewahr’,\\
Halt uns bei festem Glauben,\\
Und auf dich laß uns bauen,\\
Aus Herzensgrund vertrauen;\\
Dir uns lassen ganz und gar,\\
Mit allen rechten Christen,\\
Entfliehen Teufels Listen,\\
Mit Waffen Gott’s uns rüsten.\\
Amen, Amen, das sei wahr,\\
So singen wir: Halleluja.
\end{verse}

17.

\begin{verse}
Wo soll ich fliehen hin,\\
Weil ich beschweret bin\\
Mit viel und großen Sünden?\\
Wo kann ich Rettung finden?\\
Wenn alle Welt verränne,\\
Mein’ Angst sie nicht wegnähme.
\end{verse}

\quellseite{438}

\begin{verse}
O Jesu! voller Gnad’,\\
Auf dein Gebet und Rath\\
Kömmt mein betrübt Gemüthe\\
Zu deiner großen Güte:\\
Laß du auf mein Gewissen\\
Ein Gnadentröpflein fließen.
\end{verse}

\subsection{18}

\begin{verse}
Ach, Jerusalem du schöne!\\
Ach! wie helle glänzest du,\\
Ach! wie lieblich Lobgetöne\\
Hört man da in sanfter Ruh’:\\
O der großen Freud’ und Wonne!\\
Jesund gehet auf die Sonne,\\
Jesund gehet an der Tag,\\
Der kein Ende nehmen mag.
\end{verse}

\subsection{19}

\begin{verse}
Auf meinen lieben Gott\\
Trau ich in Angst und Noth:\\
Er kann mich all’zeit retten\\
Aus Trübsal, Angst und Nöthen:\\
Mein Unglück kann er wenden,\\
Steht all’z. in seinen Händen.

Ob mich mein’ Sünd’ anficht,\\
Will ich verzagen nicht;\\
Auf Christum will ich bauen,\\
Und ihm allein vertrauen,\\
Ihm thu’ ich mich ergeben,\\
Im Tod und auch im Leben.

Ob mich der Tod nimmt hin,\\
Ist sterben mein Gewinn,\\
Und Christus ist mein Leben,\\
Dem thu’ ich mich ergeben,\\
Ich sterb’ heut’ oder morgen,\\
Mein’ Seel’ wird Gott versorgen.
\end{verse}

\quellseite{439}

\begin{verse}
O mein Herr Jesu Christ!\\
Der du so g'duldig bist\\
Für mich am Kreuz gestorben,\\
Hast mir das Heil erworben,\\
Auch uns allen zugleich\\
Das ewig Himmelreiche.\\
Amen! zu aller Stund'\\
Sprech'n wir aus Herzensgrund:\\
Du wollest uns thun leiten,\\
Herr Christ! zu allen Zeiten,\\
Auf daß wir deinen Namen\\
Ewiglich preisen, Amen.
\end{verse}

20.

\begin{verse}
Mach' End', o Herr! mach' Ende\\
An aller unser Noth;\\
Stärk' unsre Füß' und Hände,\\
Und laß bis in den Tod\\
Uns allzeit deiner Pflege\\
Und Treu' empfohlen sein,\\
So gehen unsre Wege\\
Gewiß zum Himmel ein.
\end{verse}

21.

\begin{verse}
Christus der ist mein Leben,\\
Sterben ist mein Gewinn,\\
Denn thu' ich mich ergeben,\\
Mit Freud' fahr' ich dahin ;:
\end{verse}

\begin{verse}
Mit Freud' fahr' ich von dannen\\
Zu Christ, dem Bruder mein,\\
Daß ich mög' zu ihm kommen\\
Und ewig bei ihm sein ;:
\end{verse}

\begin{verse}
Nun hab' ich überwunden\\
Kreuz, Leiden, Angst und Noth,\\
Durch sein' heilig' fünf Wunden\\
Bin ich versöhnt mit Gott ;:
\end{verse}

\quellseite{440}

\begin{verse}
Wenn meine Kräfte brechen,\\
Mein Athem geht schwer aus,\\
Und kann kein Wort mehr sprechen,\\
HErr! nimm mein Seufzen auf;:
\end{verse}

\begin{verse}
Wenn mein Herz und Gedanken\\
Vergehen wie ein Licht,\\
Daß hin und her thut wanken;\\
Wenn ihm die Flamm’ gebricht;:
\end{verse}

\begin{verse}
Alsdann fein sanft und stille,\\
HErr! laß mich schlafen ein\\
Nach deinem Rath und Willen,\\
Wenn kommt mein Stündelein;:
\end{verse}

\begin{verse}
Und laß mich an dir kleben,\\
Wie eine Klett’ am Kleid,\\
Und ewig bei dir leben\\
In himml’scher Wonn’ und Freud’;:
\end{verse}

\begin{verse}
Amen, das wirst du, Christe,\\
Verleihen gnädiglich,\\
Mit deinem Geist mich rüste,\\
Daß ich fahr’ seliglich;:
\end{verse}

\subsection{22}

\begin{verse}
JEsus, meine Zuversicht\\
Und mein Heiland ist mein Leben;\\
Dieses weiß ich, soll ich nicht\\
Darum mich zufrieden geben?\\
Was die lange Todesnacht\\
Mir auch für Gedanken macht.
\end{verse}

\begin{verse}
JEsus, der mein Heiland, lebt;\\
Ich werd’ auch das Leben schauen,\\
Sein, wo mein Erlöser schwebt,\\
Warum sollte mir denn grauen?\\
Lässet auch ein Haupt sein Glied,\\
Welches es nicht nach sich zieht?
\end{verse}

\quellseite{441}

\begin{verse}
Ich bin durch der Hoffnung Band\\
Zu genau mit ihm verbunden,\\
Meine starke Glaubenshand\\
Wird in ihn gelegt befunden,\\
Daß mich auch kein Todesbann\\
Ewig von ihm trennen kann.
\end{verse}

\begin{verse}
Ich bin Fleisch und muß daher\\
Auch einmal zur Asche werden.\\
Das gesteh’ ich, doch wird er\\
Mich erwecken aus der Erden,\\
Daß ich in der Herrlichkeit\\
Um ihn sein mög’ allezeit.
\end{verse}

\begin{verse}
Dann wird eben diese Haut\\
Mich umgeben, wie ich glaube,\\
Gott wird werden angeschaut\\
Dann von mir in diesem Leibe,\\
Und in diesem Fleisch werd’ ich\\
JEsum sehen ewiglich.
\end{verse}

\begin{verse}
Dieser meiner Augen Licht\\
Wird ihn, meinen Heiland, kennen,\\
Ich, ich selbst, kein Fremder nicht,\\
Werd’ in seiner Liebe brennen;\\
Nur die Schwachheit um und an\\
Wird von mir seyn abgethan.
\end{verse}

\begin{verse}
Was hier kränket, seufzt und fleht,\\
Wird dort frisch und herrlich gehen,\\
Ird’sch werd’ ich ausgesä’t,\\
Himmlisch werd’ ich auferstehen,\\
Hier geh’ ich natürlich ein,\\
Nachmals werd’ ich geistlich seyn.
\end{verse}

\begin{verse}
Seyd getrost und hoch erfreut,\\
JEsus trägt euch, meine Glieder;\\
Gebt nicht Raum der Traurigkeit,\\
Sterbt ihr, Christus ruft euch wieder,
\end{verse}

\quellseite{442}

\begin{verse}
Wenn die letzt’ Trompet’ erklingt,\\
Die auch durch die Gräber dringt.

Lacht der finstern Erdenkluft,\\
Lacht des Todes und der Höllen,\\
Denn ihr sollt euch durch die Luft\\
Eurem Heiland zugesellen,\\
Dann wird Schwachheit und Verdruß\\
Liegen unter eurem Fuß.
\end{verse}

\begin{verse}
Nur daß ihr den Geist erhebt\\
Von den Lüften dieser Erden,\\
Und euch dem schon jetzt ergebt,\\
Dem ihr beigefügt sollt werden,\\
Schickt das Herze da hinein,\\
Wo ihr ewig wünscht zu sein.
\end{verse}

\subsection{23}

\begin{verse}
Christ ist erstanden von der Marter alle,\\
Deß soll’n wir alle froh sein,\\
Christ will unser Trost sein,\\
Kyrieleis.
\end{verse}

\begin{verse}
Wär’ er nicht erstanden, so wär’ die Welt vergangen,\\
Seit daß er erstanden ist,\\
So lob’n wir den Herrn JEsum Christ.\\
Kyrieleis.
\end{verse}

\begin{verse}
Halleluja! Halleluja! Halleluja!\\
Deß soll’n wir alle froh sein,\\
Christ will unser Trost sein.\\
Kyrieleis.
\end{verse}

\subsection{24}

\begin{verse}
Nun bitten wir den Heiligen Geist\\
Um den rechten Glauben allermeist,\\
Daß er uns behüte an unserm Ende,\\
Wenn wir heimfahr’n aus diesem Elende.\\
Kyrieleis!
\end{verse}

\quellseite{443}

\begin{verse}
Du werthes Licht! gib uns deinen Schein,\\
Lehr’ uns Jesum Christ kennen allein,\\
Daß wir an ihm bleiben, dem treuen Heiland,\\
Der uns bracht hat zu dem rechten Vaterland.\\
Kyrieleis!
\end{verse}

\begin{verse}
Du höchster Tröster in aller Noth!\\
Hilf, daß wir nicht fürchten Schand’ noch Tod,\\
Daß in uns die Sinne nicht verzagen,\\
Wenn der Feind wird das Leben verklagen.\\
Kyrieleis!
\end{verse}

\subsection{25}

\begin{verse}
Aus Gnaden! — dieser Grund wird bleiben,\\
So lange Gott wahrhaftig heißt;\\
Was alle Knechte Jesu schreiben,\\
Was Gott in seinem Wort aufweist,\\
Worauf all’ unser Glaube ruht:\\
Ist Gnade durch des Lammes Blut.
\end{verse}

\begin{verse}
Aus Gnaden bleibt dem blöden Herzen\\
Das Herz des Vaters aufgethan,\\
Wenn’s unter größter Angst und Schmerzen\\
Nichts sieht und nichts mehr hoffen kann.\\
Wo nähm’ ich oftmals Stärkung her,\\
Wenn Gnade nicht mein Anker wär’!
\end{verse}

\begin{verse}
Aus Gnaden! dies hör’ Sünd’ und Teufel,\\
Ich schwinge meine Glaubensfahn’,\\
Und geh’ getrost trotz allem Zweifel\\
Durch’s rothe Meer nach Canaan.\\
Ich glaub’, was Jesu Wort verspricht,\\
Ich fühl’ es oder fühl’ es nicht.
\end{verse}

\subsection{26}

\begin{verse}
Ob sich’s anließ, als wollt’ er nicht,\\
So laß dich’s nicht erschrecken;\\
Denn wo er ist am besten mit,\\
Da will er’s nicht entdecken;
\end{verse}

\quellseite{444}

\begin{verse}
Sein Wort laß dir gewisser seyn,\\
Und ob dein Herz sprach’ lauter nein,\\
So laß doch dir nicht grauen.
\end{verse}

\subsection{27}

\begin{verse}
Ich habe nun den Grund gefunden,\\
Der meinen Anker ewig hält.\\
Wo anders als in Jesu Wunden?\\
Da lag er vor der Zeit der Welt;\\
Den Grund, der unbeweglich steht,\\
Wenn Erd’ und Himmel untergeht.
\end{verse}

\begin{verse}
Es ist das ewige Erbarmen,\\
Das alles Denken übersteigt!\\
Es sind die offnen Liebesarmen\\
Deß, der sich zu dem Sünder neigt;\\
Dem allemal das Herze bricht,\\
Wir kommen oder kommen nicht.
\end{verse}

\begin{verse}
Wir sollen nicht verloren werden,\\
Gott will, uns soll geholfen sein;\\
Deswegen kam der Sohn auf Erden\\
Und nahm hernach den Himmel ein;\\
Deswegen klopft er für und für\\
So stark an unsre Herzensthür.
\end{verse}

\begin{verse}
O Abgrund! welcher alle Sünden\\
Durch Christi Tod verschlungen hat;\\
Das heißt die Wunde recht verbinden,\\
Hier findet kein Erbarmen statt,\\
Weil Christi Blut beständig schreit:\\
Barmherzigkeit! Barmherzigkeit!
\end{verse}

\begin{verse}
Darein will ich mich gläubig senken,\\
Dem will ich mich getrost vertraun.\\
Und wenn mich meine Sünden kränken,\\
Nur bald nach Gottes Herzen schaun’;\\
Da findet sich zu aller Zeit\\
Unendliche Barmherzigkeit.
\end{verse}

\quellseite{445}

\subsection{28}

\begin{verse}
Mein Heiland nimmt die Sünder an,\\
Die unter ihrer Last der Sünden\\
Kein Mensch, kein Engel trösten kann,\\
Die nirgends Ruh’ und Rettung finden:\\
Denn’n selbst die weite Welt zu klein,\\
Die sich und Gott ein Gräuel sein,\\
Denn’n Moses schon den Stab gebrochen\\
Und sie der Höllen zugesprochen,\\
Wird diese Freistatt aufgethan:\\
Mein Heiland nimmt die Sünder an.
\end{verse}

\begin{verse}
O! solltest du sein Herze sehn,\\
Wie sich’s nach armen Sündern sehnet;\\
Sowohl, wenn sie noch irre gehn,\\
Als wenn ihr Auge vor ihm thränet.\\
Wie streckt er sich nach Zöllnern aus!\\
Wie eilt er in Zachäi Haus!\\
Wie sanft stillt er der Magdalenen\\
Den milden Fluß erpreßter Thränen!\\
Und denkt nicht, was sie sonst gethan.\\
Mein Heiland nimmt die Sünder an.
\end{verse}

\begin{verse}
Wie freundlich blickt er Petrum an!\\
Ob er gleich noch so tief gefallen.\\
Nun, dies hat er nicht nur gethan,\\
Da er auf Erden mußte wallen:\\
Nein, er ist immer einerlei;\\
Gerecht und fromm und ewig treu;\\
Und, wie er unter Schmach und Leiden,\\
So ist er auf dem Thron der Freuden\\
Den Sündern liebreich zugethan.\\
Mein Heiland nimmt die Sünder an.
\end{verse}

\begin{verse}
So komme denn, wer Sünder heißt,\\
Und wen sein Sündengraus’ betrübet,\\
Zu dem, der keinen von sich weist,\\
Der sich gebeugt zu ihm begiebet.
\end{verse}

\quellseite{446}

\begin{verse}
Wie? willst du dir im Lichte stehn?\\
Und ohne Noth verloren gehn?\\
Willst du der Sünden länger dienen?\\
Da dich zu retten er erschienen?\\
O nein! verlaß die Sündenbahn. —\\
Mein Heiland nimmt die Sünder an. ;;
\end{verse}

\begin{verse}
Sprich nicht: ich hab’s zu groß gemacht!\\
Ich hab’ die Güter deiner Gnaden\\
So lang und schändlich umgebracht;\\
Er hat mich oft umsonst geladen!\\
Wofern du’s nur jetzt redlich meinst,\\
Und deinen Fall mit Ernst beweinst;\\
So soll ihm nichts die Hände binden,\\
Und du sollst noch Gnade finden.\\
Er hilft, wenn sonst nichts helfen kann:\\
Mein Heiland nimmt die Sünder an. ;;
\end{verse}

\begin{verse}
Ja, zeuch uns selbstern recht zu dir,\\
Holdselig süßer Freund der Sünder;\\
Erfüll’ mit sehnender Begier\\
Auch uns und alle Adamskinder.\\
Zeig’ uns, bei unserm Seelenschmerz,\\
Dein aufgespaltnes Liebesherz;\\
Und wenn wir unser Elend sehn,\\
So laß uns ja nicht stille stehn,\\
Bis daß ein jeder sagen kann:\\
Gott Lob! auch mich nimmt JESus an. ;;
\end{verse}

29.

\begin{verse}
Nichts hilft nur die Gerechtigkeit,\\
Die vom Gesetz herrühret;\\
Wer sich in eig’nem Werk erfreut,\\
Wird jämmerlich verführet.\\
Des HErren JESu Werk allein,\\
Das macht’s, daß ich kann selig sein,\\
Weil ich fest an ihn glaube.
\end{verse}

\quellseite{447}

\subsection{30}

\begin{verse}
Ist Gott für mich, so trete\\
Gleich alles wider mich,\\
So oft ich ruf’ und bete,\\
Weicht alles hinter sich.\\
Hab’ ich das Haupt zum Freunde,\\
Und bin geliebt bei Gott,\\
Was kann mir thun der Feinde\\
Und Widersacher Rott’?
\end{verse}

\begin{verse}
Der Grund, da ich mich gründe,\\
Ist Christus und sein Blut,\\
Das machet, daß ich finde\\
Das ew’ge wahre Gut.\\
An mir und meinem Leben\\
Ist nichts auf dieser Erd’,\\
Was Christus mir gegeben,\\
Das ist der Liebe werth.
\end{verse}

\begin{verse}
Mein Herze geht in Springen,\\
Und kann nicht traurig sein,\\
Ist voller Freud’ und Singen,\\
Sieht lauter Sonnenschein.\\
Die Sonne, die mir lachet,\\
Ist mein Herr JESUS Christ;\\
Das, was mich singend machet,\\
Ist, was im Himmel ist.
\end{verse}

\subsection{31}

\begin{verse}
Mit Fried’ und Freud’ ich fahr’ dahin,\\
In Gottes Wille,\\
Getrost ist mir mein Herz und Sinn,\\
Sanft und stille,\\
Wie Gott mir verheißen hat,\\
Der Tod ist mein Schlaf worden.
\end{verse}

\quellseite{448}

\begin{verse}
Das macht Christus, wahr’r Gottes Sohn,\\
Der treue Heiland,\\
Den du mich, Herr, hast sehen lahn,\\
Und macht bekannt,\\
Daß er sei das Leben\\
Und Heil in Noth und Sterben.

Den hast du allen fürgestellt\\
Mit großen Gnaden,\\
Zu seinem Reich die ganze Welt\\
Heißen laden\\
Durch dein theuer heilsam Wort,\\
An allem Ort erschollen.

Er ist das Heil und selig’ Licht\\
Für die Heiden,\\
Zu erleuchten, die dich kennen nicht,\\
Und zu weiden.\\
Er ist dein’s Volks Israel\\
Der Preis, Ehr’, Freud’ und Wonne.
\end{verse}

\begin{verse}
32.\\
Gloria sei dir gesungen\\
Mit Menschen- und eng’lichen Zungen,\\
Mit Harfen und mit Cymbeln schön.\\
Von zwölf Perlen sind die Pforten\\
An deiner Stadt, wir sind Consorten\\
Der Engel hoch um deinen Thron.\\
Kein Aug’ hat je gespürt,\\
Kein Ohr hat mehr gehört\\
Solche Freude:\\
Deß sind wir froh, ja! ja! ja! ja!\\
Ewig in dulci jubilo.
\end{verse}

\begin{verse}
33.\\
Ach mein herzliebes Jesulein,\\
Mach’ dir ein rein sanft Bettelein\\
Zu ruhen in mein’s Herzens Schrein,\\
Daß ich nimmer vergesse dein.
\end{verse}

\quellseite{449}

\subsection{34}

\begin{verse}
Ich hang’, und bleib auch hangen\\
An Christo, als ein Glied;\\
Wo mein Haupt durch ist gangen,\\
Da nimmt es mich auch mit.\\
Er reißet durch den Tod,\\
Durch Welt, durch Sünd’ und Noth,\\
Er reißet durch die Höll’,\\
Ich bin stets sein Gesell.
\end{verse}

\begin{verse}
Er bringt zum Saal der Ehren,\\
Ich folg’ ihm immer nach,\\
Und darf mich gar nicht kehren\\
An einzig Ungemach;\\
Es tobe, was da kann,\\
Mein Haupt nimmt sich mein an,\\
Mein Heiland ist mein Schild,\\
Der alles Toben stillt.
\end{verse}

\begin{verse}
Er bringt uns an die Pforten,\\
Die in den Himmel führt,\\
Daran mit güldnen Worten\\
Der Reim gelesen wird:\\
Wer dort wird mit verhöhnt,\\
Wird hier auch mit gekrönt,\\
Wer dort mit sterben geht,\\
Wird hier auch mit erhöht.
\end{verse}

\subsection{35}

\begin{verse}
Ein’ feste Burg ist unser Gott,\\
Ein’ gute Wehr’ und Waffen;\\
Er hilft uns frei aus aller Noth,\\
Die uns jetzt hat betroffen.\\
Der alt’ böse Feind,\\
Mit Ernst er’s jetzt meint,\\
Groß’ Macht und viel List\\
Sein’ grausam’ Rüstung ist,\\
Auf Erd’ ist nicht sein’s gleichen.
\end{verse}

\quellseite{450}

\begin{verse}
Mit unser Macht ist nichts gethan,\\
Wir sind gar bald verloren:\\
Es streit’t für uns der rechte Mann,\\
Den Gott hat selbst erkoren.\\
Fragst du, wer der ist?\\
Er heißt JEsus Christ,\\
Der HErr Zebaoth,\\
Und ist kein ander Gott,\\
Das Feld muß er behalten.
\end{verse}

\begin{verse}
Und wenn die Welt voll Teufel wär’\\
Und wollt’ uns gar verschlingen,\\
So fürchten wir uns nicht so sehr,\\
Es soll uns doch gelingen.\\
Der Fürst dieser Welt,\\
Wie saur’ er sich stellt,\\
Thut er uns doch nicht,\\
Das macht, er ist gericht’t,\\
Ein Wörtlein kann ihn fällen.
\end{verse}

\subsection{36}

\begin{verse}
Von allem Uebel uns erlös’;\\
Es sind die Zeit und Tage böß;\\
Erlös’ uns von dem ew’gen Tod,\\
Und tröst’ uns in der letzten Noth;\\
Bescher’ uns all’n ein selig’s End’,\\
Nimm unsre Seel’ in deine Händ’.
\end{verse}

\subsection{37}

\begin{verse}
JEsu, meines Herzens Freud’,\\
Süßer JEsu!\\
Meiner Seelen Seligkeit,\\
Süßer JEsu!\\
Des Gemüthes Sicherheit,\\
Süßer JEsu,\\
JEsu, süßer JEsu!
\end{verse}

\quellseite{451}

\begin{verse}
Tausendmal gedenk’ ich dein,\\
Mein Erlöser!\\
Und begehre dich allein,\\
Mein Erlöser!\\
Sehne mich, bei dir zu sein,\\
Mein Erlöser!\\
JEsu, mein Erlöser!
\end{verse}

\begin{verse}
Weide mich und mach’ mich satt,\\
Himmelsspeise!\\
Tränke mich, mein Herz ist matt,\\
Seelenweide!\\
Sei du meine Ruhestatt,\\
Ruh’ der Seelen!\\
JEsu, Ruh’ der Seelen!
\end{verse}

\begin{verse}
Nichts ist lieblicheres als du,\\
Liebste Liebe!\\
Nichts ist freundlicheres als du,\\
Milde Liebe!\\
Auch nichts süßeres ist als du,\\
Süße Liebe,\\
JEsu, süße Liebe!
\end{verse}

\begin{verse}
Ich bin krank, komm’, stärke mich,\\
Meine Stärke!\\
Ich bin matt, erquicke mich,\\
Süßer JEsu,\\
Wenn ich sterb’, so tröste mich,\\
Du mein Tröster!\\
JEsu, du mein Tröster!
\end{verse}

38.

\begin{verse}
O Herrlichkeit der Erden!\\
Dich mag und will ich nicht,\\
Mein Geist will himmlisch werden,\\
Und ist dahin gericht’t,
\end{verse}

\quellseite{452}

Wo JEsus wird geschauet, Da sehn’ ich mich hinein, Wo JEsus Hütten bauet, Denn dort ist gut zu sein.

\subsection{39}

Nun muß ich Sünd’r von dieser Welt Hinfahr’n nach Gottes Willen, Zu meinem Gott, wenn’s ihm gefällt, Will ich ihm halten stille: Mein’ arme Seel’ Ich Gott befehl’ In meiner letzten Stunden. O frommer Gott! Sünd’, Höll’ und Tod Hast du mir überwunden.

Noch eins, HErr! will ich bitten dich, Du wirst mir’s nicht versagen: Wann mich der böse Geist anficht, Laß mich doch nicht verzagen; Hilf, steu’r’ und wehr’, Ach Gott, mein HErr! Zu Ehren deinem Namen. Wer das begehrt, Dem wird’s gewährt, Drauf sprech’ ich fröhlich: Amen.

\subsection{40}

Soll diese Nacht die letzte sein In diesem Jammerthal, So führ’ mich, HErr! in Himmel ein, Zur auserwählten Schaar. Und also leb’ und sterb’ ich dir, Du starker Zebaoth! Im Tod’ und Leben hilf du mir Aus aller Angst und Noth.

\quellseite{453}

41.

\begin{verse}
Brei’t aus die Flügel beide,\\
O Jesu! meine Freude,\\
Und nimm dein Küchlein ein.\\
Will Satan mich verschlingen,\\
So laß die Englein singen:\\
Dies Kind soll unverletzet sein.
\end{verse}

42.

\begin{verse}
Und wenn ich ende meinen Lauf,\\
So hole mich zu dir hinauf,\\
Jesu! daß ich da Fried’ und Freud’\\
Bei dir genieß’ in Ewigkeit.
\end{verse}

\begin{verse}
Jesu, erhöre meine Bitt’!\\
Jesu, verschmäh’ mein Seufzen nicht!\\
Jesu, mein’ Hoffnung steht zu dir:\\
O Jesu, Jesu, hilf du mir.
\end{verse}

43.

\begin{verse}
Es ist noch eine Ruh’ vorhanden,\\
Auf, müdes Herz! und werde Licht!\\
Du seufzest hier in deinen Banden\\
Und deine Sonne scheinet nicht.\\
Sieh’ auf das Lamm, das dich mit Freuden\\
Dort wird vor seinem Stuhle weiden,\\
Wirst hin die Last und eil’ ihm zu!\\
Bald ist der schöne Kampf geendet,\\
Bald, bald der saure Lauf vollendet,\\
So gehst du ein zu deiner Ruh’.
\end{verse}

44.

\begin{verse}
Die Allmacht Gottes, des Vaters, segne mich:\\
Die Wahrheit des Sohnes benedeie mich:\\
Die Liebe Gottes, des Heiligen Geistes, behüte mich\\
Zu dem ewigen Leben gnädiglich.\\
Amen.
\end{verse}

\quellseite{454}

\section{460. Gebet eines bußfertigen Kranken}

O allmächtiger Gott und Vater, der du um der Sünde willen die Menschen mit allerlei Plagen, und insonderheit auch mit leiblicher Krankheit heimsuchest, sie damit zu züchtigen, auf daß sie nicht sammt der Welt verdammt werden. Ich komme zu dir voller Sünden, und bekenne, daß ich dadurch nicht nur diese Krankheit, sondern auch die ewige Verdammniß wohl verdienet habe, weil ich nicht allein von Natur ein Kind des Zornes und unter die Sünde verkauft bin, sondern auch alle deine Gebote die Zeit meines Lebens unzählige Mal übertreten habe. Ich nehme aber meine Zuflucht zu deiner grundlosen Barmherzigkeit, und bitte demüthig, gehe nicht ins Gericht mit mir, denn vor dir ist kein Lebendiger gerecht. Gedenke nicht der Sünden meiner Jugend, noch aller meiner Uebertretung, sondern um deines Namens willen sei gnädig meiner Missethat, die da groß ist. Wende deine Plage von mir, und vernim mein Schreien, schweige nicht über meinen Thränen. Laß ab von mir, daß ich mich erquicke, ehe denn ich hinfahre und nicht mehr hier sei. Siehe, lieber Vater, um Trost ist mir sehr bange, nimm dich meiner Seele herzlich an, daß sie nicht verderbe, und wirf alle meine Sünde hinter dich zurück. Laß dir auch meine Leibesnoth zu Herzen gehen, und nimm entweder dieselbe weg, oder lindere sie durch deinen Trost; verleihe Geduld, hilf mir mein Kreuz tragen, oder mache mich nur selig dabei. In deinen Willen sei's gestellt, mach's lieber Gott, wie dir's gefällt, dein bin und will ich bleiben. Amen.

\section{461. Gebet um Segen der Arznei}

Herr Jesu, mein allergentreuester Nothhelfer und Arzt des Leibes und der Seele, weil du selber sprichst: Die Gesunden bedürfen des Arztes nicht, sondern die Kranken; ich aber jetzt mit Krankheit behaftet bin, so will ich die Arznei, welche du uns zum Besten aus der Erde wachsen läßt, gebrauchen. Daher bitte ich dich demüthiglich, du wollest dein kräftiges Gedeihen zu diesen Arzneimitteln, welche ich in deinem Namen jetzt gebrauchen will, geben, und durch deinen Segen sie dasjenige in mir wirken lassen, wozu du sie hast geschaffen, denn ohne dich ist ja alle Mühe und Kunst
\quellseite{455}
umsonst. Aber nach deinen Geboten und Willen müssen alle Kreaturen zu unserm Heil erscheinen, helfen und dienen, und ist keine Krankheit so groß, die du, Herr, nicht vertreiben, kein Schade so gefährlich, daß du ihn nicht heilen könntest. Erzeige demnach, Herr, deine Macht, wie du versprochen hast, mich aus aller Noth zu erretten, daß dich mein Mund ewig rühme und preise. Amen.

\section{462. Kurze Stundengebete für Kranke}

\section{1. Um ein Uhr}

O einiger, wahrer Gott, mein Schöpfer und Vater, von dem Gesundheit und Krankheit, Glück und Unglück herkommt: weil es ja dein väterlicher Wille ist, daß ich so langwieriges Kreuz und beschwerliche Leibeskrankheit tragen soll, so gib mir christliche Geduld und beständige Hoffnung. Ich weiß, daß du mein einiger lieber Vater bist, es mag mir wohl oder übel gehen. Es erforderets doch aller wahren Kirchenglieder Zustand in dieser Welt, Kreuz und Elend auszustehen. Warum wollte ich mich denn ausschließen? Hat doch mein Haupt und Heiland JEsus Christus auch meinetwegen leiden müssen. Lieber Gott und frommer Vater, erhalte mich nur bei deinem Wort, und laß dasselbe meine beste Arznei und Labsal sein. Amen.

\section{2. Um zwei Uhr}

Herzlichster Herr JEsu, mein getreuer Mittler und Heiland, ich weiß gar wohl, daß ich mit meinen Sünden diese meine Leibesbeschwerung und alles Kreuz gar wohl verdienet habe. Ich tröste mich aber deiner unerschöpflichen Gnade und theuren Genugthuung, dadurch du mich an Leib und Seel erlöst hast. Es mag dem Leib allhier gehen, wie es wolle, so bin ich doch versichert, daß Leib und Seele, von der Hölle befreit, einmal zum ewigen Leben gelangen und mit himmlischer Herrlichkeit und Gesundheit werden begabet werden. Das hilf mir, Herr JEsu. Amen.

\section{3. Um drei Uhr}

O werthster Heiliger Geist, liebster Gast meiner Seelen, und wahrhaftigster Zeuge meiner himmlischen Kindschaft und Seligkeit, ich bitte dich demüthiglich, du wollest in meiner langwierigen Leibesbeschwerung
\quellseite{456}
ferner in meinem Herzen wohnen und mit deinem göttlichen Trost mich erquicken, daß ich bei deinem reinen Wort und hochwürdigen Sacramenten beständig bleibe. Erhalte in mir Glauben, Liebe und Hoffnung, als drei Säulen meines Christenthums, und gib mir Gottes väterlichen Willen gegen mich recht zu erkennen, auf daß ich demselben stille halte, in demselben lebe und sterbe. Amen.

\section{4. Um vier Uhr}

O lieber, getreuer Gott, lehre mich doch das Ende bedenken, so werde ich nicht sündigen. Einen seligen Tod gib mir; am jüngsten Gericht absolviere mich; vor der Höllen behüte mich; in den Himmel führe mich. Im Tod liege und schlafe ich mit Frieden. Am jüngsten Gerichte hilf du mir, daß ich der Höllen entgehe und im Himmel sicher wohne. Laß mein Herz den Samen deines göttlichen Worts fassen, behalten und sich daran erquicken so lang, bis du kommen und mich zu dir nehmen wirst. Amen.

\section{5. Um fünf Uhr}

Mein allerliebster Jesu, wo soll ich in meinem Kreuz und Schmerzen anders meine Zuflucht hin nehmen, als in deine heiligen fünf Wunden? Wenn meine Sinne und Gedanken verschwinden, werde ich wohl da am besten verwahret sein. Ich bleibe in dir; ach, bleibe du auch in mir! Schmücke meine Herzenslampe mit dem wahren Glaubensöl, und mache mich durch christliche Buße bereit, bei deiner Zukunft dir mit Freuden entgegen zu gehen, und der ewigen Hochzeitsfreude sammt allen Auserwählten zu genießen. Amen.

\section{6. Um sechs Uhr}

Ich danke dir, o lieber Gott und Vater, daß du mich zu deinem Ebenbilde geschaffen, mit menschlichem Leib und vernünftiger Seele; bitte auch, du wollest mich, dein Geschöpf, in meiner Noth nicht verlassen, sondern väterlich für mich sorgen, und mich kräftiglich erretten. Viel ist meiner Trübsal, und ein jeder Tag hat seine Plage; aber ich weiß, du wirst mich aus aller Trübsal und Plage endlich erlösen. Ich hoffe auf dich; laß mich ja nicht zu Schanden werden. Amen.

\quellseite{457}

\section{7. Um sieben Uhr}

O Gott, himmlischer Vater, der du am siebenten Tage der Schöpfung geruhet hast, ach ruhe auch und halte deinen Sabbath in meinem Herzen, du, o meiner Seelen einige Ruh und Erquickung. O HErr JEsu, laß das siebente und letzte Wort, das du am Kreuz gesprochen, auch mein letztes Wort in meiner Todesnoth sein, daß ich meinen Geist Gott, deinem himmlischen Vater, befehle. O Gott, Heiliger Geist, du himmlischer Betmeister, du Geist des Gebets, entzünde in mir beständige Andacht zum Gebet, daß ich im heiligen Vaterunser, in den sieben Bußpsalmen und andern Gebeten und Seufzern Gott meine Noth klage und von ihm gnädiglich erhöret werde. Amen.

\section{8. Um acht Uhr}

Allerliebster HErr JEsu, der du am achten Tage deines zeitlichen Lebens schon angefangen, für mich Schmerzen zu leiden, und dein Blut zu vergießen: ach, gib mir einen willigen Geist, meine Leibesbeschwerden dir zu Ehren auch geduldig zu ertragen, und die Hoffnung, daß du mich und meine arme Seele in dem Wasser großer Trübsal, wie die acht Seelen im Kasten Noäh, wohl erhalten werdest. Laß mich auch unter den acht Haufen sein, die du selig preisest, daß ich arm sei im Geist, geduldig im Kreuz, sanftmüthig gegen meine Feinde, hungrig und durstig nach der Gerechtigkeit, barmherzig gegen nothleidende Christen, rein im Herzen, friedfertig gegen den Nächsten, beständig in der Verfolgung. Ach, laß mich in der Geduldsschule durch alle Klassen fortrücken, bis ich aus der letzten durch Beharrlichkeit in die hohe Freudenschule des ewigen Lebens befördert werde, da ich dich mit Vater und Heiligem Geist ewig lieben und loben werde. Amen.

\section{9. Um neun Uhr}

O mein herzlichster HErr JEsu, dein sehr schweres Leiden hat dennoch einmal ein gewünschtes Ende erlanget, da du durch einen seligen Tod denselben entrissen worden. Ach, laß mir auch ein seliges Stündlein widerfahren, und mache einmal meines langwierigen Kreuzes ein Ende. Dafür will ich dir von Herzen danken, dich loben, rühmen und preisen. Nun weiß ich, du werdest meiner
\quellseite{458}
nicht so gar vergessen, und meine Hoffnung werde nicht ewig verloren sein. Ich bin hier dein Glied in der streitenden Kirche. Ei, so laß mich auch dort in der siegenden Kirche dein Glied sein und bleiben. Amen.

\section{10. Um zehn Uhr}

O gerechter und zugleich gnädiger Gott und Vater, ich bekenne herzlich und schmerzlich, daß ich nicht allein in Sünden empfangen und geboren bin, sondern auch aus menschlicher Schwachheit deine heiligen zehn Gebote vielmals übertreten mit Gedanken, Worten und Werken, damit ich denn nicht allein diese meine Krankheit und ander zeitlich Kreuz, sondern auch die ewige Höllenstrafe gar wohl verdienet habe. Weil aber dein allerliebster Sohn, mein getreuer Heiland Jesus Christus, mit seinem theuren Verdienst und heiligen Blute für alle meine Erb- und würklichen Sünden genug gethan, so bitte ich dich, du wollest mir um desselben willen gnädig sein, es bei der zeitlichen Strafe bewenden lassen und mir die ewige aus Barmherzigkeit erlassen. Regiere mich auch mit deinem Heiligen Geiste, daß ich mich von dem großen Haufen der Gottlosen absondere, mich des wahren Glaubens und der daraus fließenden christlichen guten Werke jederzeit befleißige; sonderlich aber in meinem Kreuz geduldig sei, eine gute Ritterschaft übe, Glauben und gut Gewissen behalte, und einmal in deiner Hütte ewig wohne. Lindere mir auch meine Leibesbeschwerden nach deinem guten Willen, und machs mit meinem Kreuz, daß ichs ertragen könne, um Jesu Christ willen. Amen.

\section{11. Um eilf Uhr}

O getreuer Herr Jesu, der du deine traurigen Jünger nach deiner Auferstehung oft und kräftiglich getröstet hast, auch alle Mühseligen und Beladenen freundlich zu dir lockest, und versprichst, dieselben zu erquicken, ach, besuche mich auch in meiner Krankheit, tröste und erquicke mich, der ich meine Zuflucht allein zu dir nehme, Rath und Hülfe bei dir suche. Ich sehne mich nach dem Feierabend eines seligen Sterbestündleins, und trage Verlangen nach der fröhlichen Auferstehung meines Fleisches. Ach, laß mir beides je eher je besser widerfahren, so will ich dich mit Vater und Heiligem Geist ewig loben und preisen. Amen.

\quellseite{459}

\section{12. Um zwölf Uhr}

Liebster HErr JEsu, ich sage dir herzlichen Dank, daß du dich mir in den Schriften der heiligen Propheten und Apostel hast zu erkennen gegeben, daß ich an dich glauben und durch den Glauben das ewige Leben haben soll. Ich bitte dich, du wollest mich in wahrem Glauben beständig erhalten und mich desselben Ende, die ewige Seligkeit, erlangen lassen. Mache mich auch munter und bereit, wenn du, mein liebster Bräutigam, kommen wirst, daß ich dir mit Freuden entgegengehe, und führe mich mit dir ein in das himmlische Hochzeitshaus, da will ich dich recht lieben und loben. Amen.

\section[463. Gebet eines Kranken, der von Schwermuth …]{463. Gebet eines Kranken, der von Schwermuth angefochten wird}

O HErr JEsu, du Tröster aller traurigen und angefochtenen Herzen, der du versucht bist in allem, wie auch wir, doch ohne Sünde, auf daß du könntest mit armen, betrübten Herzen Mitleid haben, zu dir nehme ich meine Zuflucht in meinem Kummer und schweren Anfechtungen, darin ich mir selber weder rathen noch helfen kann, solches auch keinem Menschen klagen mag noch darf. Ach, rathe und hilf du mir. Laß mir Barmherzigkeit widerfahren und Gnade empfinden, jetzt, da ich Hülfe, Trost und Beistand höchst vonnöthen habe. Erbarme dich mein um deiner Versuchung und Traurigkeit willen, so du meinetwegen erduldet und überwunden, daß ich meine Traurigkeit und Schwermuth durch deine Kraft auch erdulden und überwinden möge. Schrecken hat mein Herz troffen, Angst der Hölle hat mich umfangen, Furcht des Todes ist auf mich gefallen. Ach, laß mich nicht versinken oder verzagen. Ach, laß das glimmende Töchtlein meines Glaubens nicht gar verlöschen, sondern blase es auf durch deinen heiligen Geist. Sprich zu meiner betrübten Seele: Ich bin JEsus, deine Hülfe. So will ich denn mich auch erholen, und zu ihr sagen: Was betrübst du dich, meine Seele, und bist so unruhig in mir? Harre auf Gott, denn ich werde ihm noch danken, daß er meines Angesichts Hülfe und mein Gott ist. Amen. Hilf, o JEsu, Amen.

\quellseite{460}

\section[464. Gebet in langwieriger und schmerzhafter …]{464. Gebet in langwieriger und schmerzhafter Krankheit}

Barmherziger Gott und Vater! Hier liege ich, dein armes Kind, mit grausamen Schmerzen an das ängstliche Siechbette angestellt. Wie ein Knecht sich sehnet nach dem Schatten, und ein Tagelöhner, daß seine Arbeit aus sei, also habe ich bisher einen Tag (Monat, Jahr,) nach dem andern, auf deine Hülfe und Erlösung gewartet, und elender Nächte sind mir viele geworden. Ach, wende dich doch einmal wieder zu mir, und sei mir gnädig, nachdem du dein Angesicht bisher eine (so lange) Zeit verborgen hast. Herr! die Angst meines Herzens ist groß, führe mich aus meinen Nöthen, siehe an meinen Jammer und Elend und vergib mir alle meine Sünden. Laß doch mein Schreien, Winseln und Wehklagen zu deinen Ohren dringen, und dein Herz über meine Noth brechen. Soll ich noch ferner also gequälet, und meine Plage alle Morgen da sein, ach, so laß mir auch deine Gnade über mich alle Morgen neu erscheinen, lindere meine Schmerzen, und gib christliche Geduld, bis du dieselben zu der in deinem Rath bestimmten Stunde gar hinwegnehmen wirst. Herr! sprich nur ein Wort, so wird dein armer Knecht (deine elende Magd) gesund. Nun, du wirst’s wohl machen! Weil du mein Gott und Vater bist, dein Kind wirst du verlassen nicht, du väterliches Herz! Ich bin ein armer Erdenkloß, auf Erden weiß ich keinen Trost.

Jesu! mein theurer Heiland und Erlöser! ich gedenke jetzt an dein bitteres Leiden, an deine unaussprechlichen Schmerzen. Wie groß war deine Qual und Marter, die du für mich armen Sünder ausgestanden? Wie groß aber auch deine Geduld, als du deinem Vater dich gehorsam erwiesest, bis zu dem schmerzlichen Tod am Kreuz? Du thatest, liebster Heiland, unter deinen Leiden deinen Mund nicht auf, sondern warest so geduldig, als ein Lamm, das zur Schlachtbank geführet wird. Hiermit hast du mir ein Vorbild gelassen, daß ich nachfolgen soll deinen Fußstapfen. Ach, liebster Jesu! der Geist ist hiezu willig, aber das Fleisch ist schwach in mir. Darum bitte ich um deiner großen Angst und Schmerzen willen, pflanze und erhalte durch deinen Geist in meinem Herzen gottgefällige Geduld, alles, was mir von der Hand meines himmlischen Vaters aufgelegt wird, willig zu ertragen. Zeige mir,
\quellseite{461}
mein Seligmacher, in meiner Kreuzesnoth deine heiligen Wunden und eröffnete Seite, und laß deines Leidens Heil meiner Seele Labsal sein. Gib, daß ich hier im Leiden deinem Bilde ähnlich werde, auf daß ich dort mit dir in ewiger Freude deiner Herrlichkeit genießen möge.

O Heiliger Geist, du Trost der Traurigen und Stärke der Schwachen, verlaß mich nicht und thu nicht von mir die Hand ab, habe Geduld mit meiner Schwachheit und weiche nicht von mir mit deiner himmlischen Kraft. Erquicke mit deinen süßen Tröstungen meine arme bekümmerte Seele, daß sie unter dem schweren und langwierigen Kreuze nicht verzage. Wenn ich vor Angst und Schmerzen nicht reden und beten kann, wie sich’s gebühret, so vertritt mich vor Gottes Thron mit unaussprechlichem Seufzen. Erhalte mein Herz allezeit im wahren Glauben und lebendiger Hoffnung, daß ich auch in den höchsten Nöthen und äußersten Jammer dennoch ein Kind Gottes und Erbe des Himmels sei. Amen.

\section{465. Gebet um Erlösung aus Krankheitsnoth}

O gütiger Gott und Vater! ich weiß, daß du mich liebest, weil du mich mit dieser Krankheit so väterlich heim gesucht; denn welchen du lieb hast, den pflegest du stets unter der Zuchtruthe zu halten. Ich weiß aber auch, daß du meinen Odem, meine Zeit und alle meine Wege in deiner Hand hast. Von dir kommt Gesundheit und Krankheit, du hast Gewalt über Leben und Tod; du führest in die Hölle und wieder heraus. Ist es nun dein Wille, daß ich von dieser Krankheit wiederum genese, ach Gott, so heile mich und gewähre mich meines Gebetes. Laß mir deine Barmherzigkeit widerfahren, und erlöse mich von diesen Banden. Thue mir wohl, und erhalte mich durch dein Wort, daß ich lebe, und laß mich nicht zu Schanden werden über meiner Hoffnung; stärke mich, daß ich genese, und laß meine Seele leben, daß ich dich lobe. HErr, höre und sei gnädig! HErr, sei mein Helfer! Erquicke mich auf meinem Siechbette, und hilf mir von aller meiner Krankheit. Wende deine Plage von mir, denn ich bin verschmachtet von der Strafe deiner Hand. Höre mein Gebet, HErr! und vernimm mein Schreien, und schweige nicht über meinen Thränen. Laß mein
\quellseite{462}
Fleisch wieder grünen, wie in der Jugend, und laß mich wieder jung werden, wie ein Adler. Sprich nur ein Wort, so werde ich, dein Knecht (deine Magd), gesund. Wenn es mir aber nützlicher ist, aus diesem Jammertal jetzt abzuscheiden, denn allhier lange zu wallen, so wollest du mich in festem Glauben und starker Hoffnung bis an mein seliges Ende erhalten, und alle Schmerzen des Todes überwinden helfen, daß ich zum ewigen Leben hindurchdringen möge. O Herr! ich bin in deinen Händen, was dich gut dünkt, das thue, wie du weißt, daß mir’s zu meiner zeitlichen und ewigen Wohlfahrt zum Besten dienet, um Jesu Christi willen. Amen.

\section{466. Danksagung nach Genesung}

Ich danke dir, HErr, allmächtiger Gott, daß du mich so väterlich heimgesuchet und gezüchtiget hast um meiner Sünde willen. Ja, HErr, es ist mir lieb, daß du mich gedemüthiget hast, daß ich deine Rechte lerne. O mein Gott, wie oft gedachte ich auch, wie der König Hiskia: Nun ist meine Zeit aus, und von mir aufgeräumet, wie eines Hirten Hütte, nun reißet der HErr mein Leben ab, wie ein Weber, und machet’s mit mir aus den Tag vor Abend. Ich sehe aber, mein Gott, daß meine Krankheit nicht zum Tode gewesen ist, sondern zu der Ehre Gottes, daß du, mein HErr Jesu Christe, dadurch gepriesen werdest. Denn du hast dich meiner erbarmet, und dich meiner Seele herzlich angenommen. Du hast alle meine Sünde hinter dich zurückgeworfen und mir mein Leben gefristet. Wie wohl hast du meine Krankheit gewendet, und mir so gnädiglich geholfen! Ich danke dir, mein Gott, ja ich danke dir von Grund meines Herzens, daß du mich wieder aufgerichtet und gestärket hast, daß ich deinen heiligen Tempel wieder beschauen, meinen Weg und Steg gehen, und meine Geschäfte verrichten kann. Deine Güte ist es, HErr, sonst wäre es lange aus mit mir. O, wie werde ich so oft an deine Zuchtruthe gedenken, und werde ich mich scheuen mein Lebenlang, und mich vor deinem Zorn hüten. Hilf nun, HErr, mein Gott, daß ich mit der neuen Gesundheit auch ein neues Leben anfah[e]. Gib, daß ich allezeit deinen Namen preise, und dein Lob immerdar in meinem Munde sei. Regiere
\quellseite{463}
mich durch deinen Heiligen Geist, daß ich dir zu Ehren lebe, und meine Glieder nicht mehr begebe zu Waffen der Ungerechtigkeit, der Sünde zu dienen, sondern zu Waffen der Gerechtigkeit, zu dienen dir, meinem Gott, daß ich deine Liebe singe, so lange ich lebe, und dich in deiner Gemeine lobe und preise. Amen.

\section[467. Gebet eines Kranken um Vergebung der …]{467. Gebet eines Kranken um Vergebung der Sünden vor Genießung des heil. Abendmahls}

O gnädiger Gott und Vater, ich bin an Leib und Seele krank, wie dir mein Elend unverborgen. Mein herzliches Begehren ist, an der Seele gesund zu werden; damit ich die leibliche Krankheit desto geduldiger ertragen möge. Ich erkenne demnach und bekenne dir, o Gott, von Grund meines Herzens, daß ich nicht allein in Sünden empfangen und geboren bin, sondern auch mein zeitliches Leben mit vielen wirklichen Sünden zugebracht habe und also nichts Gesundes an mir ist von der Fußsohle an bis auf den Scheitel. Daher ich auch die Krankheit des Leibes sammt anderen zeitlichen Strafen, ja die ewige Verdammniß gar wohl verdienet habe. Ach, frommer, gnädiger Gott, vergib mir solche meine Sünde aus Gnaden, um Jesu Christi, deines lieben Sohnes willen. Heile du mich, Herr, an meiner armen Seele, so werde ich heil; hilf du mir, so ist mir wohl geholfen. Du bist ja der Arzt Israels. Nun aber bedürfen nicht die Starken des Arztes, sondern die Kranken. Darum nimm meine kranke Seele auf und mache sie gesund. Ach, reiche mir, deinem gar schwachen Kinde, deine allmächtige Gnadenhand und erquicke mich durch die tröstliche Absolution, daß mir alle meine Sünde vergeben werde. Sprich zu meiner Seele: Ich bin dein Heil. Sprich nur ein Wort, so wird sie gesund. Ist sie dann wieder genesen, so mag es dem Leibe hernach gehen, wie es dein heiliger, göttlicher Rath und Wille beschlossen hat. Nun wohl, du liebster Gott und Vater, der du bist ein unerschöpflicher Brunn der Barmherzigkeit, ich zweifle nicht, du werdest auch ein Strömlein deiner Gnade über meine arme Seele ergießen und fließen lassen. Bei dir ist Gnade und viel Erlösung. O Herr, erhöre die Stimme meines Flehens, und verachte nicht das Heulen meines Herzens, um Jesu Christi willen. Amen.

\quellseite{464}

\section[468. Gebet eines Kranken nach Empfang des …]{468. Gebet eines Kranken nach Empfang des heil. Abendmahls}

Allmächtiger Gott, himmlischer Vater! ich sage dir herzlich Lob und Dank, daß du mir armen Sünder nicht allein aus Gnaden, um deines geliebten Sohnes Jesu Christi willen, alle meine Sünde verziehen, sondern auch darauf seinen Leib und Blut im heiligen Abendmahl zum Pfande gegeben hast. Wie hättest du mich mehr deiner Gnade versichern und der Verheißung des ewigen Lebens vergewissern können? Ach, HErr, verleihe mir deine Gnade, daß ich es nimmermehr aus meinem Herzen lasse, sondern alle Augenblicke bedenke, daß ich in Christo habe Gnade, die Vergebung der Sünden, das Leben und Seligkeit. Gib, daß es meinen schwachen Glauben stärke, mein blödes Gewissen tröste, mein Gemüth erquicke und zu allem Guten in mir Kräfte schaffe. Wollest dir auch meinen Leib und meine Seele allezeit befohlen sein lassen, und so es zu deinem Lob und meiner Seele Heil ist, mich nicht nur an der Seele, sondern auch an dem Leibe gesund machen. Wo aber nicht, beständigen Glauben, christliche Geduld und steife Hoffnung bis ans Ende geben, auf daß ich nach deiner Verheißung, durch Jesum Christum, dessen Leib und Blut ich darauf empfangen habe, möge einnehmen und besitzen die ewige und selige Himmelsfreude. Amen.

\section[469. Gebet eines kranken Communicanten um …]{469. Gebet eines kranken Communicanten um Verlangen nach dem ewigen Leben}

O mein liebster HErr Jesu Christe, du mein Leben! weil du mich durch die Nießung deines heiligen Abendmahls und durch das Siegel und Pfand deines Leibes und Blutes des ewigen Lebens versichert hast, so bitte ich dich demüthiglich, du wollest alle Liebe der Welt und aller eiteln Dinge in mir dämpfen und auslöschen, dagegen mein Herz lenken und wenden, daß ich dich und das ewige Leben allein suche, wünsche und begehre. Du, HErr Jesu, hast gesagt: In meines Vaters Hause sind viele Wohnungen, und ich gehe hin, euch die Stätte zu bereiten. O welch ein schönes Haus muß das Haus des himmlischen Vaters sein! Ach! wann werde ich aus dem finstern Hause dieser Welt erlöset werden, und zu diesen himmlischen Wohnungen kommen? Du hast gesagt: Ich will kommen und euch zu mir nehmen, auf daß ihr seid, wo ich bin.

\quellseite{465}

Ach, mein HErr JEsu, komme nur bald und nimm mich zu dir, denn ich habe Lust abzuscheiden, und bei dir zu sein. Einen Vorschmack habe ich empfunden bei deinem heiligen Abendmahl. Ach! laß die Freude, ach, laß das Verlangen täglich zunehmen und größer werden! Wie ein Hirsch schreiet nach frischem Wasser, so schreiet meine Seele, Gott, zu dir; meine Seele dürstet nach Gott, nach dem lebendigen Gott; wann werde ich dahin kommen, daß ich Gottes Angesicht schaue? Dich, meinen Erlöser, sammt deinem Vater und dem Heiligen Geist im ewigen Leben sehen, von Angesicht zu Angesicht, wird sein die höchste Freude, das beste Leben, die vollkommenste Seligkeit. Ach! laß mich bald dahin kommen, da ich dieser Freude, dieses Lebens und dieser Seligkeit mit allen Auserwählten werde genießen, da ich dich ewig anschauen, herzlich lieben, und ohne Verdruß, ohne Unterlaß loben und preisen werde. Das gib und verleihe mir, lieber HErr JEsu, weil du selber bist das ewige Leben. Amen. Amen.

\section[470. Tröstlicher Zuspruch nach den drei …]{470. Tröstlicher Zuspruch nach den drei Artikeln des christlichen Glaubens an einen sterbenden Menschen}

Christliches Herz! du bist nunmehr nach Gottes Willen in dem letzten Kampf mit Sünde, Tod, Teufel und Hölle. Wohlan, so kämpfe einen guten Kampf und halte Glauben, so wird dir hinfort beigelegt werden die Krone der Gerechtigkeit, welche dir der HErr, der gerechte Richter, an jenem Tage geben wird. (2 Tim. 4, 7. 9.) Sei stark in dem HErrn, und in der Macht seiner Stärke. Vor allen Dingen aber ergreife den Schild des Glaubens, mit welchem du alle feurigen Pfeile des Bösewichts wirst auslöschen können. (Eph. 6, 10. 16.) Behalte in deinem Herzen deinen christlichen Glauben, den du gelernt, und bisher beständig bekannt hast.

Du glaubest, daß der allmächtige Gott der rechte Vater sei über alles, was Kinder heißt im Himmel und auf Erden. (Eph. 3, 5.) Darum wird er auch dein Vater sein und dir das himmlische Erbtheil vorbehalten haben.

Du glaubest, daß er dich erschaffen, und dir Leib und Seel
\quellseite{466}
gegeben habe. Darum wird er auch deiner, als seines Geschöpfes, nicht vergessen. Denn er haßet nichts, was er gemacht hat. (Buch der Weisheit 11, 25.)

Du glaubest, daß der eingeborne Sohn Gottes, Jesus Christus, sei in die Welt kommen, die Werke des Teufels zu zerstören, und die Sünder selig zu machen. (1 Joh. 3, 8.) Darum hast du dich vor demselbigen nicht zu fürchten, noch an deiner Seligkeit zu zweifeln.

Du glaubest, daß Jesus Christus für alle arme Sünder gelitten, gekreuziget worden und sein unschuldiges Blut vergossen habe, darum bist auch du durch solch sein Blut von allen deinen Sünden gereiniget worden. (1 Joh. 1, 7.)

Du glaubest, daß Jesus Christus sei hinunter zur Hölle gefahren, und dieselbe allen seinen Gläubigen zerstöret habe. (Hos. 12, 14.) Darum können dich die Höllenflammen nimmermehr berühren.

Du glaubest, daß Jesus Christus an dem dritten Tage von den Todten sei auferstanden, darum hat er den Tod überwunden, daß du denselben nimmermehr schmecken sollst. Auch soll der zeitliche Tod dir nur ein Schlaf und eine Thür zu dem ewigen Leben sein.

Du glaubest, daß Jesus Christus sei aufgefahren gen Himmel, darum wirst auch du einen Eingang in denselben haben. Denn er hat zu seinem Vater gesagt: Vater, ich will, daß, wo ich bin, auch die bei mir seien, die du mir gegeben hast, daß sie meine Herrlichkeit sehen, die du mir gegeben hast. (Joh. 17, 24.)

Du glaubest, daß Jesus Christus sich zu der Rechten Gottes gesetzt habe, darum hast du allda einen Fürsprecher bei dem Vater, der dich aufs allerbeste vertritt. (Röm. 8, 34. 1 Joh. 2, 1.)

Du glaubest, daß Jesus werde wiederkommen, zu richten die Ungläubigen; dieweil du aber beständig glaubest, so wirst du nicht ins Gericht kommen. (Joh. 5, 24.)

Du glaubest an den Heiligen Geist, darum wird er deiner Schwachheit aufhelfen, und wenn du selber nicht mehr reden kannst,
\quellseite{467}
dich mit unaussprechlichem Seufzen bei Gott vertreten und rufen: Abba, lieber Vater! (Röm. 8, 15. 26.)

Du glaubest eine Vergebung der Sünden, darum wird dir geschehen, wie du glaubest. Gott hat sich deiner schon erbarmet, deine Missethat gedämpfet und alle deine Sünde in die Tiefe des Meeres geworfen. (Mich. 5, 19.)

Du glaubest eine Auferstehung des Fleisches, darum hast du dir deinen Tod nicht schwer vorzustellen, dieweil du in dem Grabe nicht bleiben, sondern nur sanft darin ausruhen sollst, bis der fröhliche Tag deiner Auferstehung wird anbrechen.

Du glaubest ein ewiges Leben, darum kommst du aus diesem vergänglichen und unseligen Leben in ein unvergängliches und seliges Leben, in welchem dir dein Herr Jesu Christus vorlängst eine Stelle bereitet hat, und dir dieselbe jetzt gnädiglich wird einräumen.

Nun, so fahre getrost hin im Frieden, in dem Namen deines himmlischen Vaters, in dem Namen deines trauten Erlösers Jesu Christi, und in dem Namen deines höchsten Trösters, des Heiligen Geistes. Halte dich fest an das Lämmlein Gottes, welches der ganzen Welt Sünde trägt, (Joh. 1, 29.) und sprich in deinem Herzen: Herr, in deine Hände befehle ich meinen Geist; du hast mich erlöset, Herr, du treuer Gott! (Ps. 31, 6.)

Der Herr, der deinen Eingang in dieses Leben behütet hat, der behüte auch deinen Ausgang, und hole dich zu sich in die ewige Freude. Amen, in Jesu Christi Namen. Amen.

\section{471. Seufzer in Todesnöthen}

Ich bin ein armer Sünder, das weißt du, mein lieber Herr; aber du hast mir lassen vorbilden durch deinen lieben Sohn Jesum Christum, daß du wollest mir gnädig sein, die Sünde vergeben, und von seinem Zorn und Verdammniß wissen, und heißest mich solches glauben und nicht zweifeln; darauf verlasse ich mich und will fröhlich darauf dahin fahren. Herr, ich weiß Niemand, weder im Himmel noch auf Erden, zu welchem ich eine tröstliche Zuflucht möchte haben, als zu dir, durch Christum; ich muß mich
\quellseite{468}
nachend ausziehen von allen fremden Werken und Verdienst. HErr, ich habe keine Zuflucht, als zu deinem göttlichen Schooß, darin der Sohn sitzet; wenn ich die Hoffnung nicht hab’, so ist’s verloren.

\section[472. Gebet um selige Auflösung durch ein …]{472. Gebet um selige Auflösung durch ein sanftes Ende}

Christe, unser lieber HErr und Heiland, sei uns gnädig, daß wir nicht in Anfechtung fallen; sondern erhalte uns rein, unsträflich, einfältig, im rechten Glauben, und erlöse uns von allem Uebel durch einen seligen Abschied von diesem Jammerthal, das ist, aus dem Reiche des leidigen Teufels und seiner Welt. Dir sei Lob und Dank, mit dem Vater und Heiligen Geist, in Ewigkeit. Amen.

\section[473. Gebet in Furcht und Schrecken des Todes …]{473. Gebet in Furcht und Schrecken des Todes und der Höllen}

Lieber HErr Christe, ob ich gleich nicht das Gesetz erfülle und ob noch wohl Sünde vorhanden ist, und ich mich vor dem Tod und der Höllen fürchte: so weiß ich doch dies aus dem Evangelio, daß du mir alle deine Werke geschenket und gegeben hast. Deß bin ich gewiß, du lügest nicht, deine Zusagen wirst du wahrhaftig halten, und deß zum Zeichen habe ich die Taufe empfangen. Weil du denn, lieber Gott, mein bist, will ich gerne sterben; denn also gefällt es dir, mein Vater, und der Tod kann mir nicht schaden, er ist verschlungen in den Sieg. Und dir, lieber HErr Gott, sei Dank, der du uns den Sieg gegeben hast, durch unsern HErrn JEsum Christum. Amen.

\section[474. Ein tröstliches Gebet in unserer letzten …]{474. Ein tröstliches Gebet in unserer letzten Stunde}

Allmächtiger, ewiger, barmherziger HErr und Gott, der du bist ein Vater unseres lieben HErrn JEsu Christi! ich weiß gewiß, daß alles, was du zugesagt hast, du auch alles halten willst und kannst, denn du kannst nicht lügen, dein Wort ist wahrhaftig. Du hast mir im Anfang deinen lieben einigen Sohn JEsum Christum zugesagt, derselbige ist kommen und hat mich vom Teufel, Tod, Hölle und Sünden erlöset, darnach zu mehrer Sicherheit aus gnädigem Willen mir die Sacramente der Taufe und des Altars geschenket, darinnen mir angeboten Vergebung der Sünden, ewiges
\quellseite{469}
Leben und alle himmlischen Güter. Auf solches sein Anbieten habe ich derselbigen gebraucht und im Glauben auf sein Wort mich feste verlassen und sie empfangen. Derhalben ich gar nicht zweifle, daß ich wohl, sicher und zufrieden bin vor dem Teufel, Tod, Hölle und Sünde. Ist dies meine Stunde und dein göttlicher Wille, so will ich friedlich, mit Freuden auf dein Wort gerne von hinnen scheiden. Amen.

\section{475. Gebet in Todesnöthen}

Ich bin ein armer Sünder, das weißt du, lieber Herr, aber du hast mir lassen vorbilden durch deinen lieben Sohn, Jesum Christum, daß du wollest mir gnädig sein, die Sünde vergeben, und von deinem Zorn und Verdammniß wissen, und heißest mich solches glauben und nicht zweifeln, darauf verlasse ich mich und will fröhlich dahin fahren. Amen.

\section{476. Gebet eines Sterbenden für die Seinen}

Barmherziger Gott, treuer himmlischer Vater, du eilest mit mir zum Ende, und willst mich in dein ewiges Himmelreich abfordern, dahin ich dir herzlich gern will folgen. Dieweil ich aber die lieben Meinigen in großer Traurigkeit hinter mir verlasse, so bitte ich dich um meines Fürbitters willen, du wollest sie in ihrer Betrübniß reichlich trösten, dein Vaterherz nicht von ihnen wenden, deine milde Hand ihnen nicht entziehen, sondern sie merklich verspüren lassen, daß du sowohl ohne mich als durch mich sie versorgen könnest, wollest sie auch zu bestimmter Zeit mir in das rechte Vaterland selig nachfolgen lassen, da wir alle mit einander dich in Freuden ohne Aufhören loben und preisen wollen. Nun, wie du mir sie gegeben und vertrauet, so gebe und vertraue ich dir sie wieder; du trautes Vaterherz wirst es besser machen, denn dieses betrübte Häuflein meinet; ich werfe alle meine und ihre Sorge auf dich, du wirst sie wohl versorgen. Amen. In Jesu Namen, Amen.

\section[477. Gebet der Umstehenden, welche die Seele …]{477. Gebet der Umstehenden, welche die Seele des Sterbenden Gott befehlen}

Allmächtiger, gütiger Gott, der du unser Leben erhältst im Tod und Sterben, wir bitten dich, du wollest die Augen deiner
\quellseite{470}
Barmherzigkeit wenden zu diesem kranken Menschen, und ihn erquicken an Leib und Seele und ihm alle Sünden aus Gnaden vergeben. Nimm das Opfer des unschuldigen Todes Jesu Christi, deines lieben Sohnes, für die Bezahlung seiner Missethat, da er auch auf dessen Namen getauft und mit dessen Blute gewaschen und gereinigt ist. So errette ihn nun von des Leibes Qual und Pein, verkürze ihm seine Schmerzen, erhalte ihn wider die Anklage des Gewissens und wider alle Anfechtung des bösen Feindes, auf daß er im Glauben ritterlich kämpfe und überwinde. Verleihe ihm eine selige Heimfahrt zum ewigen Leben, schicke deine heiligen Engel her, daß sie ihn begleiten zu der Versammlung aller Auserwählten in Christo Jesu, unserm Herrn. Amen.

\section{478. Ein Anderes}

O gütiger Gott, barmherziger Vater, du bist unsere Zuflucht, unsere Stärke und Hilfe in unseren Nöthen. Laß leuchten dein heiliges Angesicht über diesen sterbenden Menschen, der jetzt vor dir erscheinen soll. Wasche ab alle seine Sünden, o Herr Gott! mit dem theuren Blute Jesu Christi, auf daß deren im Gerichte nimmermehr gedacht werde! Mehre ihm den Glauben, errette seine Seele aus der Gewalt des brüllenden Löwen und der bösen Engel. Tröste ihn durch deinen Heiligen Geist, und gib ihm in dieser Stunde zu verstehen, daß du sein gütiger und barmherziger Vater, und daß er aus lauter Gnade dein liebes Kind sei. Laß den nicht verloren werden, den du mit deinem heiligen Blute so theuer erkauft hast. Nimm seine Seele auf in das himmlische Paradieß, wie du die Seele des Schächers am Kreuze hast aufgenommen. Laß seine Seele von den Engeln getragen werden, wie die Seele des armen Lazarus, und verleihe ihm eine fröhliche Auferstehung am jüngsten Tage. Erhöre uns, o Vater aller Gnade, an seiner Statt. Erhöre deinen Sohn, unsern einigen Erlöser und Mittler, der zu deiner Rechten sitzet, und für ihn und uns alle bittet, und sei uns gnädig um des Verdienstes seines theuren Blutes und bittern Todes willen. In solchem Vertrauen befehlen wir seine Seele in deine väterliche Hand, und rufen dich an mit dem Gebet, das uns Jesus, unser Heiland, zu beten selbst gelehret und befohlen hat. Vaterunser
\quellseite{471}
! der du bist im Himmel, geheiliget werde dein Name, dein Reich komme, dein Wille geschehe, wie im Himmel, also auch auf Erden. Unser täglich Brod gib uns heute, und vergib uns unsre Schuld, als wir vergeben unsern Schuldigern. Führe uns nicht in Versuchung, sondern erlöse uns von dem Uebel, denn dein ist das Reich und die Kraft und die Herrlichkeit in Ewigkeit. Amen.

\section[479. Eine kurze Litanei nebst Gebeten bei …]{479. Eine kurze Litanei nebst Gebeten bei einem Sterbenden}

HErr, erbarme dich, Christe, erbarme dich! HErr, erbarme dich über diesen kranken Menschen!

HErr Gott Vater im Himmel, HErr Gott Sohn, der Welt Heiland, HErr Gott Heiliger Geist, Heilige Dreifaltigkeit, einiger Gott, HErr, sei gnädig und verschone dieses sterbenden Menschen! Erbarme dich über ihn!

Für dem bösen Geist, Für des Teufels Trug und List, Für allem Uebel, Für der Höllen Pein, Behüt’ ihn, lieber HErre Gott!

Durch deine heilige Menschwerdung und Geburt, Durch deinen Todeskampf und blutigen Schweiß, Durch dein Kreuz und bittern Tod, Durch dein’ heilige Auferstehung und Himmelfahrt, Durch die Gnade des Heiligen Geistes, An seinem letzten Ende, Am jüngsten Tage, Am jüngsten Gericht, Hilf ihm, lieber HErre Gott!

Wir armen Sünder bitten, du wollest uns erhören, lieber HErre Gott! Und diesen Kranken trösten, Ihm alle seine Sünde verzeihen, Und nach diesem Elend ihm das ewige Leben geben. Erhöre uns, lieber HErre Gott!

\quellseite{472}

Christe, erhöre uns!

O du Gottes Lamm, das der Welt Sünde getragen hat, erbarme dich über diesen sterbenden Menschen, und gib ihm den ewigen Frieden!

HErr, erbarme dich! Christe, erbarme dich! HErr, erbarme dich über ihn!

Vater unser! der du bist im Himmel, geheiliget werde dein Name, dein Reich komme, dein Wille geschehe, wie im Himmel, also auch auf Erden. Unser täglich Brod gib uns heute, und vergib uns unsre Schuld, wie wir vergeben unsern Schuldigern. Führe uns nicht in Versuchung, sondern erlöse uns von dem Uebel. Denn dein ist das Reich und die Kraft und die Herrlichkeit in Ewigkeit. Amen.

\subsection{Gebet}

Allmächtiger, ewiger Gott, himmlischer, getreuer, lieber Vater! tröste und stärke diese deine arme Kreatur, und verschone ihrer durch deine Güte; hilf ihr aus aller Angst und Noth, entbinde sie in deinen Gnaden, und nimm sie zu dir in dein Reich, durch JEsum Christum, deinen lieben Sohn, unsern einigen HErrn, Heiland, Erlöser und Seligmacher. Amen.

\subsection{Desgleichen}

Allmächtiger, ewiger Gott! laß dich erbarmen das ängstliche Seufzen und Klagen dieses sterbenden Menschen, nimm ihn auf in dein Himmelreich, das du ihm und allen Gläubigen aus Gnaden von Anfang der Welt bereitet hast. Entbinde ihn gnädiglich, HErr, und nach Bezahlung menschlicher Schuld tröste ihn mit deinen Auserwählten ewiglich, durch deinen lieben Sohn, unsern HErrn JEsum Christum. Amen.

\section{480. Sterbensseufzer}

HErr JEsu, dir leb’ ich; HErr JEsu, dir sterb’ ich; HErr JEsu, dein bin ich Todt und lebendig; Mach mich selig. Amen!

\quellseite{473}

\section[481. Gebet, wenn der Kranke in dem HErrn …]{481. Gebet, wenn der Kranke in dem HErrn verschieden ist}

Jesu Christo, dem HErrn der Herrlichkeit und Fürsten des Lebens, sei jetzt und allezeit Lob und Dank gesagt, daß er diesem nunmehr seligen Menschen so gnädiglich hat abgeholfen und seine Seele in seine heiligen Hände hat aufgenommen. Er wolle die Seele an dem jüngsten Tage mit dem Leibe in himmlischer Klarheit wieder vereinigen, und uns mit einander zu der von Gott bestimmten Zeit eine selige Nachfahrt und indessen einen christlichen Wandel verleihen, auch alle betrübten Herzen mit süßem Trost erquicken und uns alle in seiner ewigen Gnade erhalten, um seiner unausführlichen Güte und Barmherzigkeit willen. Amen, o du allerheiligste und herrlichste Dreieinigkeit. Amen! Amen!

\section[482. Letzter Seufzer, einem Verurtheilten …]{482. Letzter Seufzer, einem Verurtheilten vorzusprechen}

Ach mein einiger Heiland und Seligmacher Jesu Christe, dir befehle ich mich ganz und gar, dir ergebe ich mich ganz eigen, zu dir steht alle mein Trost und Hoffnung, nach dir habe ich herzlich Verlangen, daß ich bei dir sein möge, in dir will ich fröhlich einschlafen und auferstehen zum ewigen Leben. HErr Jesu Christe, der du deinem himmlischen Vater am Kreuz deine Seele und Geist in seine Hände befohlen, ich befehle dir auch jetzt, an meinem letzten Ende, meinen Geist in deine Hände; du hast mich erlöset, HErr, du getreuer Gott! HErr, meinen Geist befehle ich dir, mein Gott, mein Gott! weiche nicht von mir, nimm mich in deine Hände, o wahrer Gott! aus aller Noth hilf mir am letzten Ende. HErr Jesu, dir leb’ ich; HErr Jesu, dir sterb’ ich; HErr Jesu, dein bin ich todt und lebendig. Amen.

\quellseite{474}

\chapter{Alphabetisches Sachregister}

Abendmahl, Gebet um Begier nach dem heil. Abendmahl 225., um würdige Empfangung des heil. Abendmahls 225. 226., vor Empfangung desselben 227. 228. 232. 233., Seufzer, wenn man den Leib und das Blut Christi empfangen will 234., Danksagung nach dem Genuß des heil. Abendmahls 234. 235. 236., für Einsetzung des heil. Abendmahls 291., tröstliche Versicherung vom heil. Abendmahl 237., Gebet nach dem Genuß des heil. Abendmahls 239., eines Bekümmerten wegen unwürdiger Empfangung des heil. Abendmahls 245., eines Kranken vor und nach dem Genuß des heil. Abendmahls 434. 435.

Abendsegen, Luthers Abendsegen 5., am Sonntag 11., Montag 16. 17., Dienstag 21. 22., Mittwoch 26. 27., Donnerstag 30. 31., Freitag 34. 35., Sonnabend 40. 41., eines Kranken 372., eines Soldaten 145. 146.

Abfall, Bewahrung vor Abfall 285.

Absolution, wenn man die Absolution empfangen will 218., Danksagung nach der Absolution 218. 219.

Alte Person, Gebet einer alten Person 110.

Amtmann, Gebet eines Amtm. 142.

Andacht, Gebet um Andacht 1., Mangel an Andacht 303.

Anfechtung, Gebet in Anfechtung zur Verzweiflung 307., satanische Anfechtung 312., nach überwindener Anfechtung 314.

Antichrist, siehe Papst.

Arbeit, Gebet vor der Arbeit 114.

Armut, 327. 329.

Arznei, um Segen der Arznei 425.

Atheisterei, 188.

Auferstehung der Todten 282.

Beichte, allgemeine 42., wenn man zum Beichtstuhl gehen will 218., Gebet eines Predigers vor der Beichthandlung 128.

Bekehrung, Gebet um Bekehrung 272., Bekehrung der Verführten 187.

Beruf, Trostgebet in Anfechtung bei seinem Beruf 76., bei Wahl eines Berufs 111., um Segen beim Beruf 114.

Besessne, Fürbitte für Besessne 173. 175.

Beständigkeit, in Bekenntniß der Wahrheit 284.

Buße, Bußgebet 201., Klage über Unbußfertigkeit 206., Gebet um wahre Buße 207., Appellation eines bußfertigen Sünders 213., Gebet eines bußfertigen Kranken 425.

Collegschüler, Gebet eines Collegeschülers 134.

Comet, bei Erscheinung eines Cometen 302.

Communicanten, Morgengebet eines Communicanten 219., Klage eines Communicanten 222.

Communionandachten, 220.

Confirmanden, Gebet für die Confirmanden 170.

Dankbarkeit, Gebet um Dankbark. 351.

Danksagung, für die Liebe Gottes 353., für die Gnadenwahl 354., für die christliche Kirche 356., für den Schutz der Engel 358., für Behütung in Gefahr 359., nach der Arbeit 362., nach Vollendung eines Hauses 362., nach vollbrachter Reise 362., nach einer Seereise 363., für Erlösung aus Trübsal 365., nach erlangtem Frieden 366., nach Genesung 433., daß uns Gott zum Gebete würdig macht 1., für das gehörte Wort Gottes 44., nach glücklicher Geburt 82., eines Predigers nach vollbrachter Amtsarbeit 131., eines Soldaten 155., nach der Absolution 218. 219., nach dem heil. Abendmahl 234. 235., für die geistliche Vermählung 237., für das Leiden Christi 243., für die Schöpfung 275., für alle erzeigte
\quellseite{475}
Wohlthaten 276., für die Erlösung 276., für die Heiligung 281., für die Auferstehung der Todten 282., für das Amt der Schlüssel 290., für Einsetzung des heil. Abendmahls 291., nach erlangtem Regen 295., nach Gewitter 301., nach überstandener Anfechtung 314.

Demuth, Gebet um Demuth 253.

Dienstboten, Gebet eines Dienstb. 80.

Dürre, Gebet zur Zeit der Dürre 294.

Ehemann, Ehefrau, 85. 86., Ehesegen 88., Gebet eines Ehemanns nach Entbindung seiner Frau von einem todten Kinde 93., eines, der eine böse Frau hat 97., einer, die einen bösen Mann hat 98., wenn ein Gemahl das andre verlassen hat 100.

Eltern, Gebet für ihre Kinder 72., wenn sie die Kinder zur Schule schicken 73., Gebet der Eltern für erwachsene Kinder, um ein frommes Ehegemahl 73., für ungerathene Kinder 74., Reisegebet eines Hausvaters 77., Gebet um die Ehre der Eltern 263., für ihre Kinder auf Reisen 163., für einen Sohn in der Fremde 164., für ein krankes Kind 165. 322.

Einwanderer, Dankgebet für Einwanderer 363.

Ende, Gebet um seliges Ende 43. 285. 439.

Erkenntniß, Gottes 247., der Sünden 207. 209.

Erlösung, vom Fluch des Gesetzes 274., Danksagung für die Erlösung 276.

Fall, Klage eines Gefallenen 209. 307.

Farmer, Gebet eines Farmers 107.

Feinde, Gebet für Feinde 194.

Feuersgefahr, Gebet in Feuersg. 331.

Friede, Gebet um Frieden 342., um Frieden in der Gemeinde 186., Danksagung nach Frieden 366.

Früchte, des Landes 192. 297.

Geburt, Gebet vor und nach der Geburt 89. 90. 92.

Geburtstag, Gebet am Geburtstage 70.

Geduld, Gebet um Geduld 342.

Geiz, wider den Geiz 267.

Gemeindeversammlung, Gebet in Gemeindeversammlungen 185. 186.

Gerechtigkeit, um Gerechtigkeit 266.

Gepolter, wider Furcht vor Gepolter 320.

Gericht, jüngstes 279.

Geschäft, wenn einer ein Geschäft anfängt 111. 113.

Geschworne, Gebet eines Geschwornen 141.

Gnade, Zuflucht zur Gnade 213., Macht der Gnade 214., um Gnade 214.

Gnadenmahl, in Christo 354.

Gottesfurcht, um Gottesfurcht 249.

Guttätigkeit, um Guttätigkeit 268.

Haushalt, Gebet für den Haushalt 115.

Hausvater, Gebet eines Hausvaters für das Gesinde 171.

Hebamme, Gebet einer Hebamme 103.

Heiligung, Danksagung für die Heiligung 281., um Heiligung des Namens Gottes 258.

Hoffnung, um Hoffnung 252.

Irrgläubige, ein Christ unter Irrgläubigen 320.

Juristen, Gebet eines Juristen 142.

Jüngling, Gebet eines Jüngl. 104. 105.

Jüngste Tag, 278. 279.

Kaufmann, Gebet eines Kaufm. 109.

Keuschheit, Gebet um Keuschheit 264.

Kinder, Gebet der Kinder 79., für erwachsene 73., ungerathene 74., für ein Kind, das auf Reisen ist 163., 164., krankes Kind 165., Fürbitte kleiner Kinder für ihren Vater auf Reisen 167. 168., für ihren kranken Vater 168., wenn Kinder krank sind 322.

Kindesnöthen. 91. 161.

Kirche, Gebet für die Kirche 182. 183. 186., Danksag. für die Kirche 356.

Kranke, Morgen- und Abendgebet eines Kranken 372. 373., Gebet in Krankheit 374., eines Todkranken 375., Törichtfrüchte für Kranke 377., eines bußfertigen Kranken 425., kurze Stundengebete für Kranke 426., in Schwermuth 430., in langwieriger
\quellseite{476}
Krankheit 431., um Erlösung aus Krankheitsnoth 432., eines Kranken um Vergebung 434., eines Kranken nach empfangenem Abendmahl 435., eines kranken Communicanten um Verlangen nach dem ewigen Leben 435., wenn der Kranke verschieden ist 444., Gebet für Kranke 172. 173.

Krieg, Gebet in Kriegsnoth 339. 340. 342.

Krüppel, Gebet eines Krüppels 333.

Landplagen, Gebet in Landplag. 342.

Leben, Verlangen nach dem ewigen Leben 435.

Lehrer, Gebet eines Lehrers 131. 132.

Leiden Christi, Betrachtung des Leidens 240. 242. 243.

Liebe, Gebet um herzliche Liebe 251., Liebe zum göttlichen Wort 259. 260.

Lieder und Liederverse, 395.

Litanei, 343., bei einem Sterbenden 442., die kleine Litanei 161.

Lob Gottes 369. 370., um Gnade Gott zu loben 258., Christi 370.

Missionsstunden, Gebet für Missionsstunden 181.

Morgengebet, Luthers 4., am Sonntage 8. 9., am Montage 13. 14., Dienstag 18. 19., Mittwoch 22. 23., Donnerstag 27. 28., Freitag 32. 33., Sonnabend 36. 37., eines Kranken 372., eines Soldaten 143. 144.

Morgenpsalme des Satans 312. 313.

Mutter, Gebet einer Mutter für ihren Säugling 162., für ungerathene Kinder 74., für ein gebrechliches Kind 84., für ein ohne Taufe gestorbenes Kind 94., einer Mutter, die ihren ersten Kirchgäng hält, 96.

Noth, Sündennoth 309., Wassersnoth 333., in großer Noth und Gefahr 334. 335., Kriegsnoth 339., Todesnoth 159.

Obrigkeit, Gebet einer obrigkeitlichen Person 138., für die Obrigkeit 191., wenn man eine neue Obrigkeit wählt 191.

Operation, Gebet vor der Operation 151.

Papst, Gebet wider den Papst 183.

Pastoralkonferenzen, Gebet auf Pastoralkonferenzen 184.

Pestzeiten, Gebet eines, der in Pestzeiten nicht fliehen kann 328., eines, der nicht gebunden ist und fliehet 339, eines Predig. in Pestzeiten 130.

Prediger, Gebet eines Predigers 122. 123. 124., Trost und Tröst eines Predigers 125., vor der Predigt 125. 126. 127., nach der Predigt 127., vor der Beichte 128., vor einem Krankenbesuch 129., in Pestzeiten 430., Danksagung eines Predigers nach vollbrachter Amtsarbeit 131., Gebet um rechtschaffene Prediger 178.

Predigt, Gebet vor der Predigt 44., um fruchtbare Anhörung der Predigt 45., nach der Predigt 46. 47., um die Ehre des Predigtamts 262.

Preß, 330.

Prüfungsfragen, 195.

Rathsherr, Gebet eines Rathsh. 142.

Regen, Gebet um Regen 295.

Reise, bei Antritt der Reise 118., auf der Reise 119., Fürbitte für Reisende 163. 164. 167. 168.

Richter, Gebet einer richterlichen Person 140.

Sacramentirer, Gebet wider die Sacramentirer 294.

See, bei Ausfahrt in die See 119.

Segen, in der Berufsarbeit 114., Ehe-Segen 87. 88., Gebet eines, der mit irdischen Gütern gesegnet ist, 117.

Seuchen und Sterbensläufe 337.

Schlacht, vor der Schlacht 148., während der Schlacht 149., nach der Schlacht 149.

Schöpfung, Danksagung für die Schöpfung 275.

Schüler, 134.

Schulc, 73.

Schwachheit des Glaubens 302.

Soldat, allgemeines Gebet eines Soldaten 142. 143., Morgengebet 143. 144., Abendgebet 145. 146., um herzhaften Muth 147., auf Wachtposten 148., vor, während, nach der
\quellseite{477}
Schlacht 148. 149., eines verwundeten 149., in großen Schmerzen 150., vor der Operation 151., in Krankheit 153., Dank- und Bußgebet eines Soldaten 155., Gebet in Gemeinschaft zu sprechen 156., für seine Angehörigen 157., der Angehörigen für den Soldat im Felde 158., in Todesfurcht 159., in Todesnoth 159.

Sonn- und Festtagsgebete, 45 bis 69.

Sonnenschein, 296.

Sorge, wider die Sorge 116.

Sterbende, Gebet eines Sterbenden 440. 441. 442., Sterbensseufzer 443. 444.

Student, Gebet eines armen Studenten der Theologie 135., eines Studenten vor dem Studiren 137., nach dem Studiren 138.

Stundengebete für Kranke 426.

Sturm, Gebet in Sturm 120.

Sünde, in großer Sündenangst 304., nach einem Sündenfall 305., in Sündennoth Zuflucht zu Christo 309., daß uns Gott vor fernem Sünden behüte 239.

Taufe, 288., Taufpathen 163. 169.

Tischgebete 6. 7.

Theurung, 325.

Tod, 438. 439. 440. 443. 375., bei Todesfurcht 159., vor bösem schnellem Tod 287.

Trostsprüche, 377., tröstlicher Zuspruch nach den drei Artikeln 436.

Unbußfertigkeit, Klage über Unbußfertigkeit 206.

Ungewitter, 298. 299. 300.

Unwürdigkeit, Klage über Unwürdigkeit 233., Trost wider Unwürdigkeit 2. 238. 245.

Vater unser, 3. 283., erste Bitte 10., zweite Bitte 14., dritte Bitte 19., vierte Bitte 24., fünfte Bitte 29., sechste Bitte 33., siebente Bitte 38.

Verbrecher, Gebet eines Verbrechers 324. 325.

Verderben, Klage über das angerichte Verderben 271. 222.

Verführung, bei Gefahr von Verführung 315.

Vergebung, Seufzer um Vergebung 214. 215. 434.

Verleumder, Gebet wider Verleumder 321.

Verlust zeitlicher Habe 331.

Versiehung in Christo 311.

Versicherung der Erhörung 3.

Versöhnung, Gebet um Versöhn. 217.

Verstockung, um Bewahrung vor Verstockung 285.

Verurtheilter, letzter Seufzer eines Verurtheilten 444.

Verwundeter, 149.

Verzweiflung, in Anfechtung zu Verzweiflung 307.

Wahl eines Berufs 111.

Wachtposten, auf dem Wachtpost. 148.

Wahrhaftigkeit, Gebet um Wahrhaftigkeit 269.

Waisenkind, 326.

Wasserbrunst, 332.

Werke, Unfruchtbarkeit in guten Werken 238.

Widerwärtigkeit, 335.

Willen Gottes zu vollbringen 238.

Wittwer, Wittwe, 100. 101. 102.

Wohlthaten, Danksagung für alle Wohlthaten 276.

Wohlthäter, Fürbitte für die Wohlthäter 193.

Wort Gottes, Danksagung für das gehörte Wort 44., um Liebe des göttlichen Worts 259. 260.

Zuhörer, Gebet eines Zuhörers 176. 177.

Zukunft Christi 194. 278. 279.

Zunge, Regierung der Zunge 269.

\quellseite{478}
\part{Hausgesangbüchlein}

\quellseite{480}

\chapter{Verzeichniß der Liederdichter}

Alber, Erasmus, Generalsuperintendent zu Neu-Brandenburg, gestorben 1553. Angelus, Johannes, genannt Silesius, eigentlich Joh. Scheffler, Arzt zu Breslau, 1624—1677. Anna Sophia, Landgräfin von Hessen, Aebtissin zu Quedlinburg, 1638—1683. Anton Ulrich, Herzog zu Braunschweig und Lüneburg, 1633—1714. Behemb, Martin, Pastor zu Lauban, 1557—1622. Birken, Sigmund von, genannt Betulius, Poet zu Nürnberg, 1626—1681. Chyträus, Nathan, Rector zu Bremen, 1543—1599. Connow, Christian Friedrich, Schuldirector zu Langermünde, 1612—1682. Debler, Wolfgang Christoph, Conrector zu Nürnberg, 1660—1722. Ebert, Jakob, Professor zu Frankfurt an der Oder, gest. 1614. Falkner, Justus, Pastor der luth. Gemeinden in Albany und New York, gest. 1723. Förtsch, Basilius, Pastor zu Gumperda in Thüringen, gest. 1619. Franck, Johann, Bürgermeister zu Guben, gest. 1677. Freylinghausen, Johann Anastasius, Pastor und Director des Waisenhauses zu Halle, 1670—1744. Geier, Martin, Oberhofprediger in Dresden, 1614—1680. Gerhardt, Paul, geb. zu Gräfenhainichen im Churkreise 1606, Diakonus in Berlin, gest. als Archidiakonus zu Lübben in der Lausitz 1676. Gesenius, Justus, Consistorialrath zu Hannover, 1601—1671. Hedinger, Johann Reinhard, Hofprediger in Stuttgart, 1664—1704. Heermann, Johann, Pastor zu Köben in Schlesien, 1585—1647. Helmbold, Ludwig, Superintendent zu Mühlhausen, 1532—1598. Hermann, Nikolaus, Cantor zu Joachimsthal, gest. 1561. Huber, Konrad, Diakonus zu Straßburg, 1542. Kunth, Sigismund, Superintendent zu Baruth, 1700—1779. Löwenstern, Matthäus Apelles von, Kaiserlicher Rath, 1594—1648. Luther, Martin, 1483—1546. Matthesius, Johann, Pfarrer zu Joachimsthal, 1504—1563. Menzer, Johann, Pfarrer zu Kemnitz in der Ober-Lausitz, 1658—1734. Moller, Martin, Pastor zu Görlitz, 1547—1613.

\quellseite{481}

Mühlmann, Johann, Professor zu Leipzig, 1573—1613.

Neumann, Kaspar, Pastor und Professor zu Breslau, 1648—1715.

Oberborn, Paul, Pastor im 16. Jahrhundert.

Olearius, Johann, General-Superintendent zu Weißenfels, 1611—1684.

Prätorius, Christoph, eigentlich Schulte, Advocat zu Stendal, in der Mitte des 17. Jahrhunderts.

Rambach, Johann Jakob, Professor und Superintendent zu Gießen, 1693—1735.

Rinckart, Martin, Archidiakonus zu Eilenburg, 1585—1649.

Rist, Johann, Pastor zu Wedel im Holsteinischen, 1607—1667.

Schade, Johann Kaspar, Diakonus zu St. Nikolai in Berlin, 1666—1698.

Schein, Johann Hermann, Cantor und Musikdirector in Leipzig, 1587—1630.

Schererz, Sigismund, zuletzt lüneburgischer Superintendent, 1584—1639.

Schmidt, Johann Eusebius, Pastor in Siebleben bei Gotha, 1670—1745.

Schmolck, Benjamin, Pastor primarius in Schweidnitz, 1672—1737.

Schwenke, Hans Christoph von, 1545—1722.

Scriba, Christian, Oberhofprediger zu Quedlinburg, 1629—1693.

Selnecker, Nikolaus, Professor und Superintendent zu Leipzig, 1532—1592.

Stegmann, Josua, Superintendent zu Rinteln, 1588—1632.

Vogt, Franz, Pastor zu Lennew, 1662—1736.

Walther, Johann, Capellmeister zu Wittenberg und Dresden, Luthers Zeitgenoss.

Weber, Jeremias, Archidiakonus und Professor zu Leipzig, 1600—1643.

Winkler, Johann Joseph, Consistorialrath und Domprediger zu Magdeburg, 1670—1722.

Woltersdorf, Ernst Gottlieb, Pastor zu Bunzlau, 1725—1761.

Ziegenspeck, Michael, Pastor zu Burg-Athanis in Osterland, um 1630.

\quellseite{482}

\chapter{I. Morgenlieder}

\subsection{1. Morgenlieder für jeden Tag in der Woche}

\subsection{Zum Anfang der Woche}

Mel. Wo Gott zum Haus nicht etc.

1. Gott Lob! die Woch’ heb’ ich jetzt an; Gott Lob! der mir viel Gut’s gethan; Gott Lob! für seine Güt’ und Treu’, die mir ist alle Morgen neu.

2. Hilf Gott, daß ich auch früh’ und spat empfinde Trost, Schutz, Hülff’ und Rath; laß all’ mein Thun gesegnet sein, erhalt’ mein Herz im Glauben rein.

3. Wend’ ab all’ Trübsal, Angst und Noth, behüt’ vor bösem, schnellem Tod’: soll diese Woch’ die letzte sein, so laß mich selig schlafen ein.

4. Ehr’ sei dem Vater und dem Sohn’, sammt Heil’gem Geist in einem Thron’; der heiligen Dreieinigkeit sei Lob und Preis in Ewigkeit!

J. Olearius.

\subsection{Am Sonntag}

Mel. Wo Gott zum Haus nicht etc.

2. Gott Lob! der Sonntag ist herbei: die Woche wird nun wieder neu, heut’ hat mein Gott das Licht gemacht, mein Heil hat mir das Leben bracht.

2. Dies ist der Tag, da Jesus Christ vom Tod’ für mich erstanden ist, und schenkt mir die Gerechtigkeit, Heil, Leben, Herzenstrost und Freud’.

3. Das ist der rechte Jubeltag, d’ran man sich nicht g’nug freuen mag, da wir mit Gott versöhnet sind, daß nun ein Christ heißt Gottes Kind.

4. Herr, führ’ mich durch dein Lebenswort hinauf zur Himmels-Ehrenpfort’, damit ich heilig leb’ allhier und dir dort danke für und für.

5. Ehr’ sei dem Vater und dem Sohn’, sammt Heil’gem Geist in einem Thron’! der heiligen Dreieinigkeit sei Lob und Preis in Ewigkeit!

J. Olearius.

\subsection{Am Montag}

Mel. Wo Gott zum Haus nicht etc.

3. Gott Lob! nun geht der Montag an, da Gott der Welt viel Gut’s gethan; heut’ ist gemacht die Himmelsfest’, Gott ist mein Gott, der mich nicht läßt.

2. Gen Himmel auf, was himmlisch ist! dem Himmel sucht ein wahrer Christ: was oft sich ändert, muß vergehn, Beständigkeit die kann besteh’n.

3. Mein Gott, erhalt’ mir diesen Sinn, so lang’ ich hier auf Erden bin,
\quellseite{483}
daß ich, im Glauben bleibe treu und bis an’s End’ beständig sey.

4. Wenn oft der Mond verändert sich, so laß mich halten stets an dich, daß ich, wenn Alles wird vergeh’n, im Himmelreich’ mög’ vor dir steh’n.

5. Ehr’ sei dem Vater und dem Sohn’, sammt Heil’gem Geist in einem Thron’; der heiligen Dreieinigkeit sei Lob und Preis in Ewigkeit!

J. Olearius.

\subsection{Am Dienstag}

Mel. Wo Gott zum Haus nicht \&c.

4. Gott Lob! der Dienstag bringt das Licht, d’rum kann ich heut’ auch schweigen nicht, mein Dienst muß heut’ sein wohl bestellt, des Herzens Dienst, der Gott gefällt.

2. Heut’ sonder’t Gott die Erd’ und Meer’, schafft Kraut und Bäum’, an Frucht nicht leer: weil Alles mir zu gut geschicht, wie soll’ ich denn Gott dienen nicht?

3. Es muß Gott dienen, was er macht; d’rum ich auch billig darnach tracht’, daß mein Thun sei zu Gott gericht’t: das ist die rechte Diensttags-pflicht.

4. Daß ich leb’ heilig und gerecht, daß ich nicht sei ein Sündenknecht: wer Gott getreulich dient allhier, der singt Gott fröhlich für und für.

5. Mein Gott, regier’ du mich allein, laß mich dein’n treuen Diener sein, bring’ endlich mich mit Fried’ und Freud’ von dieser Welt zur Herrlichkeit.

6. Ehr’ sei dem Vater und dem Sohn’, sammt Heil’gem Geist in einem Thron’; der heiligen Dreieinigkeit sei Lob und Preis in Ewigkeit!

J. Olearius.

\subsection{Am Mittwoch}

Mel. Wo Gott zum Haus nicht \&c.

5. Gott Lob! die Mittwoch kömmt herbei, sein’ Gnad’ ist mir auch heute neu; Sonn’ und Mond hat Gott heut’ gemacht, daß sie regieren Tag und Nacht.

2. Sonn’ und Mond haben ihre Zeit, die Sonne der Gerechtigkeit scheint ewig, die mein Herz erfreut, die schenkt mir alle Seligkeit.

3. Wenn ich nicht seh’ der Sonnen Licht, vergeht ihr Gnadenglanz doch nicht; Gott, laß mich fröhlich aufersteh’n, zu deinem Freudenlicht eingeh’n.

4. Da werd’ ich mitten in der Schaar der Engel seh’n dein Antlitz klar; hilf mir dorthin durch Jesum Christ, der auch mein Mittler worden ist.

5. Ehr’ sei dem Vater und dem Sohn’, sammt Heil’gem Geist in einem Thron’; der heiligen Dreieinigkeit sei Lob und Preis in Ewigkeit.

J. Olearius.

\subsection{Am Donnerstag}

Mel. Wo Gott zum Haus nicht \&c.

6. Gott Lob! heut’ ist der Donnerstag, da niemand sich g’nug freuen mag; weil Gottes Segen uns erfreut, der Alles mehret und verneut.

2. Fisch’ und Gewögel hat Gott heut’ durch sein allmächtig Wort bereit
\quellseite{484}
’s, daß annoch Alles in der Welt regieret, und uns zu gut erhält.

3. Dies ist der Tag, an dem zuletzt der Herr sein Nachtmahl eingesetzt, als er uns durch sein theures Blut erlöst hat von der Höllengluth.

4. Dies ist der Tag der Himmelfahrt, da JESus aufgenommen ward, der sitzt dem Vat’r zur rechten Hand, schenkt uns das rechte Vaterland.

5. Der noch gibt seinem Donner Kraft, daß sein Wort Buß’ und Andacht schafft, und läßt uns rufen allzumal aus dieser Welt zum Himmelssaal.

6. Dies ist der Tag der Herrlichkeit, da uns des Himmels Sitz bereit’t; hilf mir, mein Heiland, daß ich dich dort seh’ und preise ewiglich.

7. Ehr’ sei dem Vater und dem Sohn, sammt Heil’gem Geist in einem Thron’; der heiligen Dreieinigkeit sei Lob und Preis in Ewigkeit!

3. Olearius.

\subsection{Am Freitag}

Mel. Wo Gott zum Haus nicht rc.

7. Gott Lob! der Freitag ist herbei, der uns von Sünden machet frei, weil heut’ mein Heiland, JESus Christ, für meine Sünd’ gestorben ist.

2. Heut’ ist der erste Mensch gemacht, der aller Welt den Tod gebracht; heut’ stirbet für uns Gottes Sohn, erwirbet uns den Himmelsthron.

3. Heut’ ist getilgt durch deinen Tod, HErr JESu! alle meine Noth; du Lebensfürst bist selbst mein Schutz, Sünd’, Teufel, Tod und Höll’ zum Trutz.

4. Laß mich die Freiheit recht versteh’n und allem Sündendienst entgeh’n; so wird, wenn geht mein Leben aus, ein Freitag mir im Himmels-haus.

5. Ehr’ sei dem Vater und dem Sohn, sammt Heil’gem Geist in einem Thron’; der heiligen Dreieinigkeit sei Lob und Preis in Ewigkeit.

3. Olearius.

\subsection{Am Sonnabend}

Mel. Wo Gott zum Haus nicht rc.

8. Gott Lob! dies ist der Tag der Ruh’, den bring’ ich Gott zu Ehren zu; der laß mein Herz in ihm stets ruh’n, daß ich mög’ seinen Willen thun.

2. Heut’ hat geruhet JESus Christ, der für mein’ Sünd’ begraben ist, und hat die wahre Ruh’ bereit’t und alle Himmelseligkeit.

3. Führ’ mich, mein Gott, zu solcher Rast, wenn sich end’t meine Arbeitslast; gib hier Trost, Ruh’ und Einigkeit und dort die volle Herrlichkeit.

4. Ehr’ sei dem Vater und dem Sohn, sammt Heil’gem Geist in einem Thron’; der heiligen Dreieinigkeit sei Lob und Preis Ewigkeit.

Joh. Olearius.

\subsection{Zum Beschluß der Woche}

Mel. Wo Gott zum Haus nicht rc.

9. Gott Lob! die Woch’ ist auch dahin; d’rum ich dir billig dankbar bin, daß ich die Zeit hab’ überlebt, und deine Gnad’ noch ob mir schwebt.

2. Gott Lob! für seine Güt’ und Treu, die mir ist alle Morgen neu; der mir beisteht durch seine Kraft, der mir Trost, Rath und Hülfe schafft.

\quellseite{485}

3. Was ich verseh’n, vergiss, mein Gott, durch deines Sohn’s Kreuz, Hohn und Spott; gib ferner Segen, Fried’ und Freud’, und endlich deine Seligkeit.

4. Ehr’ sei dem Vater und dem Sohn, sammt Heil’gem Geist in einem Thron; der heiligen Dreieinigkeit sei Lob und Preis in Ewigkeit. J. Olearius.

\subsection{2. Allgemeine Morgenlieder}

Mel. Der lieben Sonnen Licht und Pracht.

10. Auf, Seel’! und danke deinem Herrn aus reinem Herzensgrunde! Das helle Licht vertreibt die Stern’ in dieser Morgenstunde. Die Sonn’ entdecket dir der Güter Gottes Zier, die sich, zu Gottes Kinder Schau, besiegelt in dem Morgenthau.

2. Sprich: Gott, ich bin nicht g’nug dazu, die Gnade zu erzählen, die du durch sich’re Leibesruh’ erzeigt hast meiner Seelen. Ich kann nicht zeigen an, was du mir Gut’s gethan, nebst Andern, die auch trauen dir und eines Glaubens sind mit mir.

3. Nun, meiner Seelen Hülff’ und Schutz! ich rühme deine Treue, die deinen Kindern kommt zu Nutz, die alle Morgen neue. Ich dank’ herzinniglich, daß du so gnädig mich und auch die Meinen diese Nacht durch deine Engel hast bewacht.

4. Gib, liebster Gott, daß ferner auch ich heute christlich wandle und meiner Pflicht also gebrauch’, daß, was ich thu’ und handle, gereiche mehr und mehr zu deines Namens Ehr’, zu meiner Seelen Heil und Schutz und meines Nächsten Dienst und Nutz.

5. Behüte mich für allem, was Leib und Seele kränket: hingegen gib, was angenehm, was deinen Segen schenket. Wend’ ab Gefahr und Noth, Krieg, Hunger, Feu’r und Tod; schickst du mir aber Kreuz mit zu: so gib Geduld, Trost, Hülff’ und Ruh’.

6. Ich will gern Alles nehmen an zu Lieb’ und auch zu Leide, weil deine Hand es ändern kann, und wenden Leid in Freude. Dein Herz, auf mich gericht’t, kann mich verlassen nicht: denn wer im Glauben kommt zu dir, der ist gesegnet für und für.

7. Nun, dir sei Alles heimgestellt, Seel’, Leib, Stand, Gut und Leben. Mach’ Alles, wie es dir gefällt: dir hab’ ich es ergeben. Bewahre Stadt und Land, Lehr-, Wehr- und auch Nährstand; erhalte Wort und Sacrament; gib ein vernünftig, selig End’.

8. Mein Herr und Gott, dich laß ich nicht, dir lieg’ ich hier zu Füßen, bis mir dein Mund den Segen spricht, der Alles kann durchsüßen: wenn mich dein Himmel tränkt und mir viel Gut’s einschenkt, wenn mich dein süßer Gnadengeist aus meines Jesu Wunden speist.

9. Hierauf nehm’ ich mit Freuden an, was mir ist anbefohlen. Dein Geist mich führt’ auf eb’ner Bahn, bis du mich wirst einholen in’s schöne Paradies zu deinem Lob und Preis. So ist der Tag in Gott verbracht, Anfang und Ende gut gemacht. Ehr. Scriver.

\quellseite{486}

Mel. Christus, der uns selig macht.

11. Christe! wahres Seelenlicht, deiner Christen Sonne, o du klares Angesicht, der Betrübten Wonne! Deiner Güte Lieblichkeit ist neu alle Morgen; in dir bin ich recht erfreut, darf für gar nichts sorgen.

2. Wecke mich vom Sündenschlaf, der du bist das Leben; neues Leben in mir schaff’: denn, hast du gegeben dieser Welt das Sonnenlicht, welches all’ erfreuet, willst du mich ja lassen nicht täglich uner neuet.

3. Ohne Licht steht man kein Licht; ohne Gottes Leuchten steht kein Mensch Christum nicht, der muß uns befeuchten mit dem hellen Himmels-thau seiner süßen Lehre; d’rum, Herr Christe, auf mich schau’ und dich zu mir kehre.

4. Kehre zu mir deine Güt’, freundlich mich anblicke, daß mein innerstes Gemüth in dir sich erquicke und die süße Himmelsluft mit Begierde schmecke; sonsten sei mir nichts bewußt, so da Freud’ erwecke.

5. Laß ja ferne von mir sein Hoffart, Augenweide, Fleischeslust und Heuchelschein, JEsu, mir verleide, damit sich die Welt ergößt zu ihr’s selbst Verderben; denn was sie für’s Beste schätzt, machet ewig sterben.

6. Nun, so bleibe stets mein Licht, JEsu, meine Freude, bis der frohe Tag anbricht, da, nach allem Leide, ich, in weißer Kleiderpracht, werde freudig springen und, daß Gott es wohl gemacht, ohn’ Aufhören singen. Christophs Prätorius.

Mel. Wo Gott zum Haus nicht gibt sein’ Gunst.

12. Das walt’ Gott Vater und Gott Sohn, Gott Heil’ger Geist in’s Himmels Thron. Man dankt dir, eh’ die Sonn’ aufgeht, wenn’s Licht anbricht, man vor dir steht.

2. D’rum beug’ ich diesen Morgen früh’ in rechter Andacht meine Knie, und ruf’ zu dir mit heller Stimm’, dein’ Ohren neig’, mein’ Red’ vernim[m].

3. Ich rühm’ von Herzen deine Güt’, weil du mich gnädig hast behüt’, daß ich nun hab’ die finstre Nacht in Ruh’ und Frieden zugebracht.

4. Ich schließ’ und wüst’ nicht, wie mir wär’, da schlich der Teufel um mich her. Den hast du durch dein’ Macht vertrieb’n, daß ich für ihn zur Ruh’ bin blieb’n.

5. Mein Gott, ich bitt’ durch Christi Blut, nimm mich auch diesen Tag in Hut. Laß heut’ dein’ liebste Engelein mein’ Wächter und Gefährten sein.

6. Dein Geist mein’n Leib und Seel’ regier’ und mich mit schönen Gaben zier’. Er führ’ mich heut’ auf rechter Bahn, daß ich was Gut’s verbringen kann.

7. Gib Gnad’, daß ich mein Werk und Pflicht mit Freuden diesen Tag verricht’ zu deinem Lob’ und meinem Nutz, und daß ich thu’ mein’n Nächsten Gut’s.

8. Hilf, daß ich zu regieren wiss’ mein’ Augen, Ohren, Händ’ und Füß’, mein’ Lippen, Mund und ganzen Leib, all’ bös’ Begierden von mir treib’.

\quellseite{487}

9. Bewahr’ mein Herz für Sünd’ und Schand’, daß ich, vom Uebel abgewandt, mein’ Seel’ mit Sünden nicht beschwer’ und mein Gewissen nicht versehr’.

10. Behüt’ mich heut’ und allezeit für Schaden, Schand’ und Herzeleid. Tritt zwischen mich und meine Feind’, so sichtbar und unsichtbar seind.

11. Mein’n Aus- und Eingang heut’ bewahr’, daß mir nichts Uebel’s widerfahr’. Behüt’ mich für ein’m schnellen Tod, und hilf mir, wo mir Hilf ist noth.

Mel. Lobet den HErren Alle, die ihn fürchten.

13. In dieser Morgenstund’ will ich dich loben, o Gott, mein Vater, in dem Himmel oben. Was du für Gnade mir stets thust beweisen, will ich nun preisen.

2. Du hast mich an das Tageslicht lan kommen, auch durch die Tauf’ in dein Reich aufgenommen, und lässest mich in deinem Sohn entbinden von Schuld und Sünden.

3. Du gibest mir den Heil’gen Geist daneben, daß er mich lehr’ und führ’ in meinem Leben, damit ich lern’ auf dich alleine bauen und dir vertrauen.

4. Du hast mich all’zeit väterlich ernähret, was ich bedurft, mir immerzu bescheret, und wenn mich Noth und Unglück hat berühret, bald d’raus geführet.

5. Jetzt hast du auch durch deiner Engel Schaaren mich lassen heut’ in dieser Nacht bewahren, daß mich darin kein Unfall hat erschrecket, noch aufgewecket.

6. O HErr, mit meinem Mund und Gemüthe erzähl’ und rühm’ ich solche große Güte, die ich empfangen hab’ an allen Enden aus deinen Händen.

7. Ich bitt’, du wollest über mir fort walten und mich auch heute diesen Tag erhalten, daß Satan nicht durch seine List und Pfeile mich übereile.

8. Hilf, daß ich bleib’ wie alle frommen Christen rein von der Welt und ihren Sündenlüsten, daß ich auch meinen eig’nen bösen Willen nicht mög’ erfüllen.

9. HErr, laß in deinen Schutz sein eingeschlossen Leib, Seele, Freund’, Verwandte, Hausgenossen, und was ich sonst durch deine milde Gabe auf Erden habe.

10. Die Obrigkeit erhalt’ bei gutem Stande, gib Friede, Glück und Heil im ganzen Lande, behüt’ uns allerseits, o HErr, aus Gnaden vor Schand’ und Schaden.

11. Du wollest auch an diesem Tag’ mich stärken durch deine Kraft in mein’s Berufes Werken, auf daß ich möge in denselben allen dir wohlgefallen.

12. Zuletzt verleih’, wenn ich aus diesem Leiden, es sei heut’ oder morgen, werd’ abscheiden, daß ich auf Christi Tod mög’ fröhlich sterben, den Himmel erben.

13. So bin ich denn an Leib und Seel’ genesen und wohl auf dieser Welt allhie gewesen. Dort habe ich das rechte Kleinod funden, hier überwunden.

\quellseite{488}

14. Daß hilf mir, Vater, der du mit dem Sohne und Heil'gen Geist regierst in einem Throne. Gib, daß ich dessen, was ich jetzt begehret, ja werd' gewähret.

Justus Gesenius.

In eigener Melodie.

14. Lobet den HErren Alle, die ihn fürchten, laßt uns mit Freuden seinem Namen singen und Preis und Dank zu seinem Altar bringen. Lobet den HErren!

2. Der unser Leben, das er uns gegeben, in dieser Nacht so väterlich bedeckt und aus dem Schlaf uns fröhlich auferwecket. Lobet den HErren!

3. Daß unsre Sinnen wir noch brauchen können und Hand und Füße, Zung' und Lippen regen, das haben wir zu danken seinem Segen. Lobet den HErren!

4. Daß Feuerflammen uns nicht allzusammen mit unsern Häusern unverseh'ns gefressen, daß macht's, daß wir in seinem Schooß gesessen. Lobet den HErren!

5. Daß Dieb' und Räuber unser Gut und Leiber nicht angetast't und grausamlich verletzet, darüber hat sein Engel sich gesetzet. Lobet den HErren!

6. O treuer Hüter, Brunnen aller Güter, ach, laß doch ferner über unser Leben bei Tag und Nacht dein' Hut und Güte schweben. Lobet den HErren!

7. Gib, daß wir heute, HErr, durch dein Geleite auf unsern Wegen unbehindert gehen und überall in deiner Gnade stehen. Lobet den HErren!

8. Treib' unsern Willen, dein Wort zu erfüllen, lehr' uns verrichten heilige Geschäfte, und wo wir schwach sind, da gib du uns Kräfte. Lobet den HErren!

9. Nicht' unsre Herzen, daß wir ja nicht scherzen mit deinen Strafen, sondern fromm zu werden vor deiner Zukunft uns bemüh'n auf Erden. Lobet den HErren!

10. HErr, du wirst kommen, und all' deine Frommen, die sich bekehren, gnädig dahin bringen, da alle Engel ewig, ewig singen: Lobet den HErren!

P. Gerhardt.

Mel. Wer nur den lieben Gott.

15. Mein Gott! nun ist es wieder Morgen, die Nacht vollendet ihren Lauf; nun wachen alle meine Sorgen auf einmal wieder mit mir auf. Die Ruh' ist aus, der Schlaf dahin, und ich leb' wieder, wo ich bin.

2. Ich bin noch immer auf der Erde, wo jeder Tag sein Elend hat, wo ich nur immer älter werde und häufe Sünd' und Missethat. O Gott! von dessen Brod ich zehr', wenn ich dir doch auch nütze wär'!

3. Ohn' Zweifel siehst du mich aufstehen, regier' mich auch in dieser Welt. Ich weiß nicht, wie mir's heut' wird gehen, mach' Alles so, wie dir's gefällt. Schleuß mich in deine Vorsicht ein, dein will ich todt und lebend sein.

4. Vergib mir, Vater! alle Sünden, die ich mit aus dem Bette bring', und laß mich vor dir Gnade finden, erhöre, was ich bet' und sing'; denn
\quellseite{489}
wo ich nur bei dir wohl steh’, so acht’ ich gar nicht, wie es geh’.

5. Hilf du in allen Sachen rathen, denn ich bin selber mir nicht klug. Behüte mich für Mißethaten, für böser Menschen List und Trug; laß mich den Tag wohl legen an, und Gutes schaffen, wo ich kann.

6. Behüte mir mein Leib und Le- ben, mein Weib und Kind, mein Hab und Gut, (Verwandt’ und Freunde, Hab und Gut,) und laß den Himmel Segen geben, wenn meine Hand das Ihre thut; hilf, daß ich Alles wohl verricht’; du wirst es thun, ich zweifle nicht.

Caspar Neumann.

Mel. Was willt du dich betrüben.

\subsection{Danksagung}

16. Mit Gott will ich’s anfangen, sein Name sei gepreist, nachdem die Nacht vergangen und nun der Tag sich weist. Hör’ auf, mein Herz, zu ruh’n, erkenne Gottes Güte und richte dein Gemüthe, viel Gutes heut’ zu thun.

\subsection{Bitte}

2. Herr! dir ist nicht verborgen, daß wir noch Sünder sind, und daß man alle Morgen auf Erden Kummer find’t; d’rum halt’ mein Herze rein, wend’ alle mein Verderben, und sollt’ ich heute sterben, so laß es selig sein.

\subsection{Gebet}

3. Was selig ist und nütze, gib mir von deinem Thron’, Gott Vater! mich beschütze; regier’ mich, Gott der Sohn! stär’ mich, Gott Heil’ger Geist! Herr! segne und behüte, laß leuchten deine Güte! gib Frieden allermeist!

\subsection{Fürbitte}

4. Ja, Vater, der du Allen die Sonne scheinen läßt, dein gutes Wohlgefallen steht alle Tage fest. Erbarme dich auch heut’ der From- men und der Bösen, hilf allesamt erlösen jetzt und in Ewigkeit.

Caspar Neumann.

Mel. O Ewigkeit, du Donnerwort.

17. Nun tret’ ich wieder aus der Ruh’ und geh’ dem sauren Tage zu, wie mir ist auferleget. Nicht weiß ich, was für neue Plag’ mir heute noch begegnen mag; doch weiß ich, daß mich träget mein frommer Gott in seiner Hut, daß mir die Last nicht Schaden thut.

2. Zwar meine Bürd’ ist täglich neu, doch ist mein Gott auch täglich treu, er träget meine Sinnen. Vor Abend Keiner recht vernimmmt, was ihm den Tag all ist bestimmt; es bahnt ein jeder Morgen mir einen frischen Weg zur Pein, der kann mit Gott erstiegen sein.

3. Wie wenig hab’ ich oft gedacht, daß so würd’ sein zu End’ gebracht die saure Tagesbürde; doch hab’ ich Abends wohl gespürt, daß du mich, höchster Gott! geführt, daß mich dein Schein und Würde geleitet, mir ge- leuchtet hat auf meinem Weg’ und Kreuzespfad.

4. Weil ich denn deß versichert bin, was trau’st du, verzagter Sinn! die Bürd’ auf dich zu legen? Trag’, was du kannst, Gott träget mit, der Herr der Welt der strauchelt nicht, bei ihm ist lauter Segen; mit ihme
\quellseite{490}
geh’ ich frisch daran und fürchte nicht die Leidensbahn.

5. So ist getrost mein frischer Muth, wenn ich mich in des Höchsten Hut kann eingeschlossen sehen; doch daß ich deß versichert sei, muß ich von Sünden leben frei, und Gottes Wege gehen. Mein Gott geht nimmer meinen Steg, wo ich nicht wandle seinen Weg.

6. D’rum, liebster Gott! leit’ meinen Fuß, daß ich dir folg’ in wahrer Buß’, und läut’re mich von Sünden: so kann ich als ein kühner Held bestreiten diese böse Welt, mit dir sie überwinden; so tret’ ich muthig an den Tag und scheue nicht, was kommen mag.

7. Ich leg’ auf dich, mein Gott und HErr! was mir zu tragen wird zu schwer, die Last, die mich gebogen; ich lege meine Würd’ und Stand in deine große Allmachts hand, die du mir nie entzogen; mit der hast du von Jugend auf geleitet meinen Lebenslauf.

8. Ich leg’ in deinen Schooß hinein, die meine Blutsverwandten sein, da sind sie wohl beschlossen; ich lege meine arme Seel’ in Jesu sich’re Wundenhöhl’, du wirst sie nicht verstoßen, wann sie vom Leibe scheid’t der Tod: ich leg’ mich ganz in dich, mein Gott!

Anton Ulrich, Herzog zu Braunschweig.

In eigener Melodie.

18. O Christe, Morgensterne, leucht’ uns mit hellem Schein. Schein’ uns von’s Himmels Throne an diesem dunklen Ort mit deinem reinen Wort.

2. O Jesu, Trost der Armen, mein Herz heb’ ich zu dir. Du wirst dich mein erbarmen, dein’ Gnade schenken mir, daß trau’ ich gänzlich dir.

3. Ich kann und mag nicht schlafen, ich kann nicht fröhlich sein. Mir ist verwund’t mein’ Seele und fürcht’t der Höllen Pein. O Christ, erbarm’ dich mein.

4. O Jesu, lieber HErre, du ein’ger Gottessohn, von Herzen ich begehre, du woll’st mir Hülfe thun. Du bist der Gnadenthron.

5. Du hast für mich vergossen dein rosenfarbes Blut. Das laß mich, HErr, genießen. Tröst’ mich durch deine Güt’, hilf mir, daß ist mein’ Bitt’.

6. Jesu. Ist dir verwund’t so sehre die arme Seele dein, thu’ du dich zu mir kehren, ich will dein Helfer sein, vergeben Schuld und Pein.

7. Laß du von Sünden abe und bis ein frommer Christ. Ich will dich selber laben und schenken meinen Geist, der dich zum Himmel weist.

8. Ich will dich selber speisen mit meinem Leib’ und Blut, mein’ Lieb’ an dir beweisen, und will dir theilen mit mein’n Schatz und höchstes Gut.

9. Seele. O Jesu, Lob und Ehre sing’ ich dir allezeit. Den Glauben in mir mehre, daß ich nach dieser Zeit mit dir eingeh’ zur Freud’.

10. Der uns dies Liedelein sange, so wohl gesungen hat. Gott helf, daß uns gelinge im Leben und im Tod’ durch Christi Wunden roth.

Basilius Fritsch.

\quellseite{491}

Mel. Freu’ dich sehr, o meine Seele.

19. Süßer JESU! deiner Gnaden sag’ ich Dank und deiner Güt’, daß du mich für allem Schaden diese ganze Nacht behüt’t, daß mein Seelenfeind durch List mein nicht mächtig worden ist: daß ich ganz gesund aufstehe und mit Lust die Sonne sehe.

2. Wär’ ich, HErr, ohn’ dein Erbarmen diese Nacht geschlafen ein, o wie würde denn mir Armen manches Leid geschehen sein! Satan und die böse Welt hätten mich schon hingefällt; aber weil du mich geschützet, hat ihr’ keines mich beschmitzet.

3. Nun, so laß, mein Fels und Hügel! JESU! meine Zuversicht! unter deiner Gnaden Flügel, d’rauf mein’ Hoffnung ist gericht’t, mich auch ferner diesen Tag ohn’ Gefahr und ohn’ Klag’, ganz beglückt zu Ende bringen; laß mir auch mein Thun gelingen.

4. Leib und Seel’ und alle Sinnen, die Gedanken, Werk’ und Wort, mein Thun, Lassen und Beginnen sei allein dir, treuer Hort! zu regieren heim=gestellt; mach’s mit mir, wie dir’s gefällt; denn ich mich dir ganz ergebe, ob ich sterbe oder lebe.

5. Tilge meine schweren Sünden, JESU! durch dein theures Blut; laß mich deine Gnad’ empfinden, wenn in Angst sie nöthig thut. Sei mein Reichthum, Schuh und Licht, wenn mir sonstens All’s gebricht; laß in deiner Wunden Höhlen sein die Ruhstatt’ meiner Seelen.

6. Laß auch deine Hülff empfinden, die in Angst und Nöthen sein, denen aller Trost will schwinden, geuß dich selbst in’s Herz hinein. Segne einen jeden Stand, sei auch gnädig unserm Land’, und laß deiner Engel Schaaren uns für allem Leid bewahren.

7. Ach! HErr JESU! mich regiere durch den werthen Heil’gen Geist, daß ich so mein Leben führe, wie mich dein Wort unterweist, daß ich meinen Lebenslauf richte nach dem Himmel auf; und wenn es einst kommt zum Sterben, laß mich dann den Himmel erben.

Joh. Lassenius.

\chapter{II. Abendlieder}

\subsection{1. Abendlieder für jeden Tag in der Woche}

\subsection{Am Sonntag}

Mel. Nun ruhen alle Wälder.

20. Der Sabbath ist vergangen; ich habe mein Verlangen nach Herzens Wunsch erfüllt; Gott hat mich unterweiset, mit Lebens-Brod gespeiset, und meiner Seelen Durst gestillt.

2. Gott ruht durch’s Wort im Herzen; drum leg’ ich ohne Schmerzen auch meinen Leib zur Ruh: denn allen Sünden-Schaden deckt JESUS nun in Gnaden mit seinem Purpurmantel zu.

3. O du dreieinig’s Wesen! mein
\quellseite{492}
Geist ist schon genesen, weil ich dein Tempel bin. Ich habe Licht vom Lichte, dein leuchtend Angesichte treib’ alle Finsternisse hin.

4. Du wirst schon bei mir wachen, und eine Sonne machen auch mitten in der Nacht: Bis bei den Cherubinen ein Sonntag ist erschienen, der alle Nacht zu Schanden macht.

5. Ich schlafe ganz vergnüget; denn wo mein Herze lieget, da ist der Engel Heer. Mich stört kein Weltgetümmel, es träumt mir nur vom Himmel. Ach, wer doch nur bald droben wär’!

Benj. Schmolck.

\subsection{Am Montag}

In eigener Melodie.

21. O Jesu! treu’ster Heiland mein! ich geh’ in mein Schlaf-Kämmerlein; ich will mich legen in die Ruh, schleuß du die Thür selbst nach mir zu.

2. Verzeih’ mein’ Sünd’n aus Gnaden mir, treib all’ schädliche Träum’ von hier, breit’ über mich dein’ Flügel aus, so stehn die Engel um das Haus.

3. Wenn ich nun schlaf’, wach’ du für mich, treib’ alles Unglück hinter sich; laß mich zu deinem Lob aufstehn, und fröhlich an mein’ Arbeit gehn.

4. Behüt’ für Feu’r- und Wassersnoth, für einem bösen schnellen Tod, für allem Uebel und Gefahr mich und all’ fromm’ Christen bewahr.

Zusatz.

5. Gott Vater, nimm mich in dein’ Hut! Gott Sohn, wach mich mit dein’ Blut! Gott heil’ger Geist, erleuchte mich! daß in den Himmel komme ich.

6. Und wenn du wirst nach diesem Leb’n den letzten Feierabend geb’n, die Ruh im Himmel ist bereit in alle ew’ge Ewigkeit.

Alle Abendandacht aus dem Breslauer Gesangb.

\subsection{Am Dienstag}

In eigener Melodie.

22. Wir danken dir, Herr Jesu Christ! daß du unser Erlöser bist, uns b’hütet hast für all’r Gefahr durch deiner lieben Englein Schaar.

2. Wir geb’n nun, Herr, zu unsrer Ruh, o red’n es unsrer Seel nicht zu, was wir leider! gesündigt han, vergib uns all’s, o Gottes Sohn!

3. Laß dein’ Eng’lein stets bei uns sein, die uns allzeit bewahren fein; daß der Teufel hab’ keine Macht an uns weder zu Tag noch Nacht.

4. Für Feu’r, Seelen- und Leibesnoth behüt’ uns, lieber Herre Gott! in deinem Nam’n schlafen wir ein, hilf uns, dein’n armen Kinderlein.

5. Es g’scheh, wie’s kann, wir sind ja dein, erworb’n durch dein’ fünf Wunden rein; dir leben wir, dir sterben wir, dein’ Kindlein sind wir für und für.

6. Amen! Amen! zu guter Nacht, der Engel Gottes uns bewacht: Gott Lob und Dank in Ewigkeit, wir gehn zur Ruh ohn’ alles Leid.

Nikolaus Selneccer.

\quellseite{493}

\subsection{Am Mittwoch}

Mel. Christ, der du bist der etc.

23. Ein Tag geht nach dem andern hin, da ich noch stets im Leben bin; die Zeit verläuft, und ich mit ihr, nur du, mein Gott, bist stets bei mir.

2. Du gibest, daß ich manchen Tag mit Segen hinterlegen mag; und wenn mich eine Last gedrückt, so bist du es, der mich erquickt.

3. Hingegen ist mein Wandel nicht nach deinem Willen eingericht't. Ich bin nicht werth, daß deine Gnad' so lange mich verschonet hat.

4. Doch trau' ich fest auf deine Huld, die größer ist, als meine Schuld. Ich halte dich und laß dich nicht, bis mir dein Mund den Segen spricht.

5. Ach! segne, Vater, meine Ruh, und gib mir dein Geleite zu, so schlaf ich wie ein Jacob ein, weil Mahanaim um mich sein.

6. Getreuer Jesu, steh' mir bei, und tritt des Satans Kopf entzwei, daß, wenn mich seine List betriegt, er bald zu meinem Füßen liegt.

7. Gib, guter Geist, daß mir dein Licht auch in dem Dunkeln nicht gebricht; verleih' mir einen solchen Sinn, daß ich im Schlaf auch wachend bin.

8. Du heilige Dreifaltigkeit, dein bin und bleib ich allezeit, soll diese Nacht die letzte sein, so leb' und sterb' ich dir allein.

Benj. Schmolk.

\subsection{Am Donnerstag}

Mel. Jesu, meine Freude.

24. Hirte deiner Schafe, der von keinem Schlafe etwas wissen mag, deine Wundergüte war mein Schild und Hütte den vergang'nen Tag. Sei die Nacht noch auf der Wacht, und laß mich von deinen Schaaren-um und um bewahren.

2. Decke mich von oben vor der Feinde Toben mit der Vaterhuld. Ein versöhnt Gewissen sei mein Ruhekissen, drum vergib die Schuld. Denn dein Sohn hat mich davon durch die tief geschlag'nen Wunden gnädiglich entbunden.

3. Laß auch meine Lieben keine Noth betrüben, sie sind mein und dein. Schließ uns mit Erbarmen in die Vaterarmen ohne Sorgen ein. Du bei mir und ich bei dir, also sind wir ungeschieden, und ich schlaf in Frieden.

4. Komm, verschließ die Kammer und laß allen Jammer ferne von uns sein. Sei du Schloß und Riegel, unter deine Flügel nimm dein Küchlein ein. Decke zu mit Schutz und Ruh', so wird uns mit nichten grauen vor des Satans Klauen.

5. Wie wenn ich mein Bette heut zum Grabe hätte. Wie bald roth bald todt. Drum hast du beschlossen, laß mein Ziel verflossen, kommt die Todesnoth, so will ich nicht wider dich. Leg' ich nur in Jesu Wunden, sterb' ich alle Stunden.

6. Nun wohlan, ich thue in vergnügter Ruhe Mund und Augen zu, Seele, Leib und Leben hab' ich dir ergeben, o du Hüter du. Gute Nacht, nimm mich in Acht, und erleb' ich ja den Morgen, wirst du weiter sorgen.

Benj. Schmolk.

\quellseite{494}

\subsection{Am Freitag}

Mel. O Jesu Christ, wahr’r Gottes Sohn.

25. Des Morgens, wann ich früh aufsteh’, des Abends, wann ich zu Bett geh’, sehen mein’ Augen, HErr, auf dich, HErr Jesu Christ, dir befehl ich mich.

2. In die heil’gen fünf Wunden dein, so thu’ ich ruh’n und sicher sein mit Leib und Seele, Hab’ und Gut, das ist mein Schutz dein heiligs Blut.

3. Wie’s Blut des Osterlämmleins zwar im Ländelein Gosen auch war. Ich wache oder schlafe ein, so thust du allzeit bei mir sein.

4. Drum hält dein heiligs Blut die Wacht, wed’r Teuf’l, Tod noch kein Feind ich acht’. Drum wo ich bin, bist du bei mir, all’ mein Glück und Kreuz steht bei dir.

5. Ich leb’ od’r sterb’, so bin ich dein. Das soll auch mein letzt’ Bekenntniß sein in der letzt’n Stund’ und auch im Tod. Das hilf mir, du getreuer Gott.

B. L.

\subsection{Am Sonnabend}

Eigene Melodie.

26. Der Tag hat sich geneiget, die Nacht herfür jetzt geht. Gott sei gebenedeiet, der uns beschützet hat durch seine milde Güte, erhält uns Leib und Seel’, woll’ uns fortan behüten für allem Ungfäll’.

2. Nichts ist auf dieser Erden, das da beständig bleibt. Allein die Güte des HErren die währet in Ewigkeit, steht allen Menschen offen, Gott läßt die Seinen nicht. D’rauf setz’ ich all’ mein Hoffen, mein Trost, mein Zuversicht.

3. Dem hab’ ich mich ergeben in dieser argen Welt. So ist des Menschen Leben wie Blümlein auf dem Feld. Des Morgens in dem Thauen sein sie gefärbet schön, bald werden’s abgehauen, verderben von Stund an.

4. Vergiß mir, lieber HErre, mein Sünd’ und Missethat. Ich hab’ gesündigt sehr und bitt’, HErr, um Gnad. Wenn du mir zu wollst schreiben mein’ Sünd’ und auch mein’ Schuld, wo sollt’ ich für dir bleiben? Den Tod hätt’ ich verschuldt’.

5. Ich bitt’, daß du mir gnädig durch Christus Willen seist. Mach’ mich von Sünden ledig, gib mir den Heil’gen Geist, der mich weise und lehre, ja der mich leit’ und führ’, auf daß ich nimmermehr Gottes Steg und Weg verlier’.

6. Mein’ Leib, mein’ Seel’, mein Leben, mein Haus, mein Gut und Ehr’, was du mir hast gegeben, das befehl ich dir, o HErr, in dein’ göttliche Hände. Behüt’ mich gnädiglich. Gib mir ein selig Ende und nimm mich in dein Reich.

7. Aus meines Herzens Grunde sag’ ich dir Lob und Dank in dieser Abendstunde, darzu mein Leben lang, o Gott, für deine Güte zu Ehren deinem Nam’n. Wollst uns fortan behüten, daß wir nicht zweifeln d’ran.

Paul Oderborn.

\quellseite{495}

\subsection{2. Allgemeine Abendlieder}

In eigener Melodie.

27. Der lieben Sonnen Licht und Pracht hat nun den Tag vollführet. Die Welt hat sich zur Ruh’ gemacht, thu’, Seel’, was dir gebühret. Tritt an die Himmelsthür’ und bring’ ein Lied herfür. Laß deine Augen, Herz und Sinn auf JEsum sein gericht’t hin.

2. Ihr hellen Sterne leuchtet wohl und gebet eure Strahlen. Ihr macht die Nacht des Lichtes voll, doch noch zu tausend Malen scheint heller in mein Herz die ew’ge Himmelskerz’, mein JEsus, meiner Seelen Ruhm, mein Schutz, mein Schatz und Eigen- thum.

3. Der Schlaf zwar herrschet in der Nacht bei Menschen und bei Thieren. Doch Einer ist, der oben wacht, bei dem kein Schlaf zu spüren. Es schlummert, JEsu, nicht dein Aug’ auf mich gericht’t. D’rum soll mein Herz auch wachend sein, daß JEsus wache nicht allein.

4. Verschmähe nicht dies schlechte Lied, das ich dir, JEsu, singe. In meinem Herzen ist kein Fried’, bis ich es zu dir bringe. Ich bringe, was ich kann, ach nimm es gnädig an. Es ist doch herzlich gut gemeint, o JEsu, meiner Seelen Freund.

5. Mit dir will ich zu Bette geh’n, dir will ich mich befehlen. Du wirst, mein Schutzherr, auf mich seh’n zum Besten meiner Seelen. Ich fürchte keine Noth, auch selbst den Tod. Denn wer mit JEsu schlafen geht, mit Freuden wieder aufersteht.

6. Ihr Höllengeister, packet euch, hier habt ihr nichts zu schaffen. Dies Haus gehört in JEsu Reich, laßt es nur sicher schlafen. Der Engel starke Wacht hat es in guter Acht. Ihr Heer und Lager hält ihm Schutz, d’rum sei auch allen Teufeln Trutz.

7. So will ich denn nun schlafen ein, JEsu, in deinen Armen. Dein’ Aufsicht soll mein Bette sein, mein Lager dein Erbarmen, mein Kissen deine Brust, mein Traum die süße Lust, die aus der Seiten Wunde fleußt und dein Geist in mein Herze geußt.

8. So oft die Nacht mein’ Ader schlägt, soll dich mein Geist umfangen. So viel Mal sich mein Herz bewegt, soll dies sein mein Verlangen, daß ich mit lautem Schall möcht’ rufen über- all: Ach JEsu, JEsu, du bist mein, und ich auch bin und bleibe dein.

9. Nun, matter Leib, gib dich zur Ruh’ und schlafe sanft und stille. Ihr müden Augen, schließet euch zu, denn das ist Gottes Wille. Schließet aber dies mit ein: HErr JEsu, ich bin dein. So wird der Schluß recht wohl gemacht. Nun, JEsu, JEsu, gute Nacht.

Christian Scriver.

Mel. Christe, der du bist Tag und Licht.

28. Gott Lob! der Tag ist nun dahin, d’rum ich, mein Gott, dir dankbar bin, daß ich den Abend hab’ erlebt und deine Gnad’ ob mir noch schwebt.

2. Ach, HErr, sei gnädig, hab’ Geduld, vergib mir alle Sündenschuld. Dein heil’ger Engel mich bewahr’ vor aller Angst, Noth und Gefahr.

\quellseite{496}

3. Laß mich auch mit der Morgenstund’ aufwachen fröhlich und gesund, daß ich, HErr, deine Ehr’ ausbreit’ und dich dort lob’ in Ewigkeit.

Joh. Olearius.

\subsection{Abendlied eines Kranken}

Mel. Wenn wir in höchsten.

29. Gott Lob! der Tag ist nun zum End’, gekränkter Leib, ihr schwachen Händ’, betrübter Geist, gequälter Muth, schickt euch zu Gott, dem höchsten Gut.

2. Das Herz und Lippen Opfer bringt, dem Vater euer Dankslied singt, durch den ihr wieder einen Tag habt überwunden eure Plag’.

3. Hat Satan meiner gleich begehrt, daß er in dieser Leibsbeschwerd’ mich sichten möcht’ in seinem Sieb durch Ungeduld und Sündentrieb;

4. Hat mich doch, Vater, deine Güt’ für seiner List und Macht behüt’t, daß, ob er sich schon nah’ gemacht, er mir doch keinen Schaden bracht.

5. Ward unter dieser schweren Last, die du mir aufgelegt hast, des Fleisches Schwachheit oft vermerckt, hast du doch meinen Geist gestärckt;

6. Erhalten meines Glaubens Licht, daß es nicht worden gar zu nicht; erwecket Seufzen und Gebet, daß ich getrost mit dir gered’t;

7. Und mein Herz für dir ausge-schütt’t, wenn noch so sehr der Schmerz gewüth’t; ich hatte manche Ein-d’rungsstund’, ist gleich der Leib noch nicht gesund:

8. Ja, da er noch voll Schmerzen ist, war ihm das Leben doch gefrist’t, dazu all’ Nothdurft auch geschafft, und ich nicht plötzlich hingerafft.

9. Deß sey dir Dank, ich rühme frey, daß aller Treu’ zu g’ring ich sey, zu g’ring so mancher Gnad’ und Gab’, die ich von dir empfangen hab’.

10. Bin ich heut’, als ein schwaches Kind, gefallen in so manche Sünd’, so stell’ ich mich doch wieder ein: o Vaterherz, erbarm’ dich mein!

11. Vergiß, was ich gesündigt hier, in Himmel, Vater, und für dir! Dein liebes Kind mich wieder heiß’ durch deines Sohnes Blut und Schweiß.

12. Darauf ich nun mein Leib und Seel’ dir wieder herzlich anbefehl’: wenn’s ab durch deine Hut und Wacht, was uns die Nacht endlich macht.

13. Sag’ meinen Schmerzen: bleibt in Ruh’; wo nicht, so gib Geduld dazu; dein’ Leidensnacht mir Labsal schaff’, wenn mir die Krankheit bricht den Schlaf.

14. Ist’s dein’ Will’, ei! so laß gesund erreichen mich die Morgenstund’; sprichst du: ich hab’ nicht Lust zu dir: nur wie du willt, so schick’s mit mir.

15. Soll diese Nacht die letzte sein, so laß mich fröhlich schlafen ein, und nach vollbrachtem Lebenslauf nimm meinen Geist, HErr JEsu, auf.

16. In deinen Wunden such’ ich Ruh’, dein Purpurkleid bedeck’ mich nun zu: mein Schlafküss’n sey dein’ Dornenkron’, mein Fleiß dein Blut, o Gottes Sohn!

\quellseite{497}

17. D’rauf schlaf’ ich fröhlich ein in dir, wünsch’ allen Frommen, die mit mir im Kreuz’ und Aengsten halten Wacht, durch dich auch eine gute Nacht.

Weimarisches Gesangbuch von 1744.

Mel. Geduld, die soll’n wir haben.

30. In dieser Abendstunde laßt uns mit heller Stimm’ Gott lob’n von Herzensgrunde, daß nun der Tag ist hin. Er hat uns all’ für Schaden behütet Groß und Klein. Er woll’ uns auch aus Gnaden des Tages Schuld verzeih’n.

2. Doch so wir, lieber Herre, haben was Gut’s verbracht, sei dir davon Lob, Ehre und Dank allein ge- sagt. Wir rühmen deine Gaben und bitten ferner dich, laß uns die Nacht Ruh’ haben und schlafen sicherlich.

3. Mit deiner Hand uns decke, wenn wir entschlafen seyn, daß uns nicht übel erschrecke der böse Christen- feind mit Sturm- und Feur-Geläute und anderm Ungelück. B’hüt uns auf allen Seiten für seinem Jägerstrick.

4. Beim Schlaf soll’n wir erwä- gen, was er uns bildet ab. Wenn wir uns niederlegen, das Bett ist unser Grab. Da sah’n wir an zu sterben, daß G’hör und G’sicht verschwind’t. Doch woll’n wir nicht verderben, daß Leb’n sich wieder find’t.

5. Dafür bitten wir Alle, o from- mer, treuer Gott, daß uns nicht über- falle ein böser, schneller Tod. Hilf, daß sich alle Blicke mit rechter Zu- versicht ein jeder zum Tod’ schicke, von Christo lasse nicht.

6. Was wir Liebes auf Erden all- hier nach uns verlan, deß sollt du Vater werden und dich selbst nehmen an, bis wir zusammen kommen in aller Heil’gen Schaar und preisen dei- nen Namen ewig. Das werde wahr!

Mel. Werde munter, mein Gemüthe.

31. Unsre müden Augenlieder schlie- ßen sich jetzt schläfrig zu, und des Lei- bes matte Glieder grüßen schon die Abendruh’. Denn die trüb’ und finstre Nacht hat des hellen Tages Pracht in der tiefen See verdecket und die Sterne aufgestecket.

2. Ach bekenn’, eh’ du gehst schla- fen, du, o meines Leibes Gast, ob du den, der dich erschaffen, heute nicht er- zürnet hast. Thu’, ach thu’ bei Zeiten Buß, ach, geh’ und fall’ ihm zu Fuß, und bitt’ ihn, daß er aus Gnaden dich der Strafe woll’ entladen.

3. Sprich: Herr, dir ist unver- holen, daß ich diesen Tag verbracht anders, als du mir befohlen. Ja, ich habe nicht betrachtet meines Amtes Ziel und Zweck, habe gleichfalls dei- nen Weg schändlich, o mein Gott, verlassen, bin gefolgt der Wollust Straßen.

4. Ach, Herr, laß mich Gnad’ er- langen, gib mir nicht verdienten Lohn. Laß mich deine Hut umfangen, sieh’ an deinen lieben Sohn, der für mich genug gethan. Vater, nimm den Bür- gen an. Dieser hat für mich erduldet, was mein’ Unart hat verschuldet.

5. Oeffne deiner Güte Fenster, sende deine Wacht herab, daß die schwarzen Nachtgespenster, daß des Todes finstres Grab, daß das Uebel, so bei Nacht unserm Leib zu fallen
\quellseite{498}
tragt’t, mich nicht mit dem Netz umdecke und kein böser Traum mich schrecke.

6. Laß mich, Herr, von dir nicht wanken, in dir schlaf’ ich gut und wohl. Gib mir heilige Gedanken, und bin ich gleich Schlafes voll, so laß doch den Geist in mir zu dir wachen für und für, bis die Morgenröth’ angehet und man von dem Bett’ aufstehet.

7. Vater droben in der Höhe, dessen Nam’ uns theu’r und werth, dein Reich’ komm’, dein Will’ geschehe, unser Brod werd’ uns bescheert, und vergib uns unser Schuld, schenk’ uns deine Gnad’ und Huld, laß uns nicht Versuchung tödten, hilf uns, Herr, aus allen Nöthen.

Joh. Frank.

Mel. Herzlich thut mich verlangen.

32. Wall’s Gott! mein Werk ich lasse, die Sonn’ Fey’rabend mel’t; sie hat vollendt ihr’ Straße, schleichet wieder in ihr Zelt; so mögen auch mein’ Sachen ruh’n bis zu seiner Zeit, ich will Fey’rabend machen mit schuld’ger Dankbarkeit.

2. Mein’ Augen, Herz und Hände, o Jesu! Gottes Sohn! zu dir ich nunmehr wende zum schuld’gen Tageslohn. Denn du bist selbst getreten an meine Werkstatt gut, hast mir helfen arbeiten, regieret mein’n Sinn und Muth.

3. Mein Haupt hast du gestärket, mein’n Fingern geben Kraft, hab’ dein’n Segen gemerket, der allein Frommen schafft; daher ist wohlgerathen mein’ Arbeit und mein’ Kunst; ohn’ dich geht nichts von Statten, ohn’ dich ist All’s umsonst.

4. D’rum ich von Herzensgrunde dich, Herr Gott, lob’ und preis’ in dieser Abendstunde, und bitt’ mit ganzem Fleiß, du wollest gnädig hören mein arm Betergebet, das Gut’ in mir vermehren durch dein’ Barmherzigkeit.

5. Gleichwie vor alten Zeiten du hast viel Gut’s erzeigt des Abends denen Leuten, der’r Herz sich zu dir neigt und fest auf dich gebauet, so woll’st du auch geruh’n, wie unser Herz dir trauet, uns Lieb’s und Gut’s zu thun.

6. Als Noah hat gelassen ein Taublein aus sein’m Schiff, kehrt’ es wieder sein’ Straffen und bracht’ ein’n Friedensbrief; zur Betstzeit im Munde trug es ein Oelblatt grün, daraus Noah verstunde, des Herren Zorn wär’ hin.

7. Zween heil’ge Engel kamen des Abends zu dem Lot, in ihren Schutz ihn nahmen wider die gottlos’ Rott’, erlöst’n den Propheten, bald fiel ein Schwefel und Feu’r, macht’ den gottlosen Städten ihr’ Freud’ und Frevel theu’r.

8. Gleicher Weise wir lesen, da Elias, der Prophet, im Hungerland gewesen, hört, was der Herre thät: Vögel gedienet haben zu Tisch dem Gottesmann’, Abends und Morgens Raben Brod und Fleisch brachten an.

9. So woll’st du, Herr! uns geben Abend- und Morgenbrod, und was zu diesem Leben uns allenthalben noth; dein’n Engel woll’st uns schicken, auf daß er uns bewahr’ für’s Teufels List und Stricken, so sind wir ohn’ Gefahr.

\quellseite{499}

10. Erhöre unser Bitten, ach, HErr, du treuer Gott! die Stadt woll’st du behüten für Feu’r und aller Noth; und weil die Völker toben, erregen Krieg und Streit, so sende uns von oben den Fried’ zu unsrer Zeit.

11. Ja, weil’s will finster werden um’s Wort, der Gnaden Licht, der Satan auf der Erden viel Ketzerei anricht’t; so bleib’ bei uns, HErr Christe! mit deiner Gnad’ und Schein, dein werthes Wort uns friste, alßdemm wir sicher sein.

12. Hiermit ich nun vollende mein Taggeschäfft und Sach’, und bitt’ herzlich zu Ende: HErr! den Feierabend mach’! d’rauf der Sabbath angehet, der währt viel tausend Jahr’, der ewiglich bestehet, Amen! das werde wahr!

M. Ziegenhagen.

\chapter{III. Von dem Worte Gottes}

33. Allein auf Gottes Wort will ich mein’n Grund und Glauben-bauen. Das soll mein Schatz sein ewiglich, dem ich allein will trauen. Kein’ menschlich’ Weisheit will ich nicht dem göttlich’ Wort vergleichen. Was Gottes Wort klar spricht und richt’t, soll billig Alles weichen.

2. Allein auf Gott und sein Wort rein mein Herz sich soll verlassen. Sein Wort soll mir ein’ Leuchte sein, zu geh’n auf rechter Straßen. O Gott, laß mich kein’ falsche Lehr’ von deiner Wahrheit trennen. Hilf mir um deines Namens Ehr’ die Wahrheit zu bekennen.

3. Allein auf Gott verlaß ich mich, auf seine Gnad’ und Güte. Ich hoff’, er wird mich gnädiglich für’s Teufels List behüten. Erhalt’ mich, Gott, bei deinem Wort und gönne mir solch’s reine für’s Teufels Lügen, Trug und Mord, für allem falschen Scheine.

4. Allein auf Christum, Gottes Sohn, von Ewigkeit geborn vom Vater, Gott in gleichem Thron, hab’ ich zum Trost erkoren. Den hat Gott in die Welt gesandt allen Menschen auf Erden, für alle Sünd’ solch hohes Pfand ein Opfer lassen werden.

5. Alleine Christus ist mein Trost, der für mich ist gestorben, mich durch sein Blut vom Tod erlöst, die Seligkeit erworben. Hat meine Sünd’ getragen gar, bezahlt an seinem Leibe. Das ist für Gott gewißlich wahr. Hilf, Gott, daß ich’s fest glaube.

6. Alleine Christus hat’s gethan, der mich mit Gott versühnet. Kein menschlich Werk mir helfen kann, wie hoch und schön es scheinet. Alleine Christus ist der Schatz, der mir hat bracht das Leben. Mein’ Werk’ sind nur des Mächsten Muß, die Glaubens Frucht muß geben.

7. Alleine Gott der Heilig’ Geist gibt diese Lehre reine um Christus willen allermeist der christlichen Gemeine. Vernunft, Verstand, Kunst, Witz und Kraft dies Wort nicht kann erlangen.

\quellseite{500}

Der freie Will’ hie gar nichts schafft, bleibt stets im Irrthum hangen.

8. Alleine Gott durch seinen Geist muß Gutes in uns wirken, was Wol-len und Vollbringen heißt, erleuchten und auch stärken. Der Mensch, zum Guten todt und blind, kann sich zu Gott nicht lenken, ist gar verderbet durch die Sünd’, daß er nichts Gut’s kann denken.

9. Alleine Gottes Geist und Stärk’ das Herz zum Guten rühret. Die Neugeburt ist Gottes Werk, die zu der Wahrheit führet. O Gott, mein Herr, erleuchte mich, mein Herz auch zu dir wende. Bei deinem Wort mich seliglich erhalt’ bis an mein Ende.

10. Gott Vater, Sohn und Heil’ger Geist, hilf, daß mein Glaub’ dich preise, mein Fleisch dem Geist’ Gehorsam leist’, des Glaubens Frucht beweise. Hilf, Herre Christ, aus aller Noth, wenn ich von hinnen scheide, und führe mich auch aus dem Tod’ zur Seligkeit und Freude.

Joh. Walther.

Mel. Gott des Himmels und der rc.

34. Herr, ich preise dein Erbarmen, deine Treu und Gütigkeit, daß du mich Unwürd’gen, Armen in dem Volk der Christenheit, wo dein Wort uns selig macht, hast auf diese Welt gebracht.

2. Ich bin auf der fetten Weide! denn dein Evangelium zeigt mir Leben, Fried’ und Freude in dem wahren Christenthum. Und dies helle, süße Licht schauen viele Tausend nicht.

3. Wenn die Heiden ohne Wahrheit zu den stummen Götzen gehn; wenn die Juden ohne Klarheit tief in der Verblendung stehn; wenn auch viele Christen blind und voll Aberglauben sind:

4. So hab’ ich dein Wort in Händen, daß ich selber lesen kann. So darf mich kein Irrthum blenden. Dein Wort zeigt die Wahrheit an. Deine Schrift, das helle Licht, gibt den Blinden das Gesicht.

5. Wer dein Wort nicht liest und hört, lebt auf Erden wie ein Thier. Wen dein Heil’ger Geist nicht lehrt, der kommt nimmermehr zu dir. Denn er öffnet den Verstand, und macht uns dein Heil bekannt.

6. O so laß mich nicht im Dunkeln, wie der blinde Haufe irrt. Laß dein Wort im Herzen funkeln, dessen Kraft die Seele führt! Mach’ es mir zur Lebenskost, süß, wie lauter Milch und Most.

7. So kann ich im Lichte wandeln, als ein sel’ges Christenkind; fröhlich glauben, heilig handeln. So werd’ ich wie du gesinnt. O Herr Jesu, durch dein Blut hab’ ich’s hier und ewig gut.

E. G. Woltersdorf.

Mel. Lobt Gott ihr Christen allzugleich.

35. Lobt Gott, ihr Christen allzugleich, in seinem höchsten Thron, sein Wort zeigt uns sein Himmelreich, durch Christum, seinen Sohn.

2. Was er zuvor zu vieler Zeit, auf so viel Weis’ und Art, durch die Propheten ausgebreit’t, ist uns nun offenbart.

\quellseite{501}

3. Sein Sohn hat uns nun selbst gelehrt den Weg zur Seligkeit, daß wir durch ihn zu Gott bekehrt, einkeh’rn zur Himmelsfreud’.

4. Das ist die Wahrheit und das Licht, das ist der Morgenstern, denn was Gott zusagt, das geschicht, er hilft von Herzen gern.

5. Und wann die ganze Welt vergeht, so bleibt mir dieses fest, daß Gottes Wort allein besteht, daß mich Gott nicht verläßt.

6. Ehr’ sei dem Vater und dem Sohn, Ehr’ sei dem Heil’gen Geist, gelobt sei Gott in’s Himmels Thron, Mein Gott, sei hochgepreist.

Aus J. Crüger’s Praxis pietatis v. J. 1693.

Mel. Erhalt’ uns, HErr, bei deinem Wort.

36. O Gott, du höchster Gnadenhort, verleih’, daß uns dein göttlich Wort von Ohren so zu Herzen dring’, daß es sein’ Kraft und Schein verbring’.

2. Der einig’ Glaub’ ist diese Kraft, der steif an Jesu Christ behast. Die Werk’ der Lieb’ sind dieser Schein, dadurch wir Christi Jünger sein.

3. Verschaff’ bei uns auch, lieber HErr, daß wir durch deinen Geist je mehr in dein’r Erkenntniß nehmen zu und endlich bei dir finden Ruh’.

Konrad Huber.

Mel. Herzlich thut mich verlangen.

37. Wohl dem, der Jesum liebet, und dessen Himmelswort! der wird niemals betrübet von’s Teufels Höllenwuth. Wo JEsus sich befindet, da stehet alles wohl; wer sich auf JEsum gründet, der lebet lebensvoll.

2. Bist du vielleicht verirret vom rechten Lebensport, hat dich die Welt verwirret, komm, hier ist Gottes Wort! dies wird dir klärlich weisen die rechte Lebensbahn, darauf du müßest reisen, wenn du willst himmelan.

3. Bist du vielleicht betrübet, ja wirfst du fort und fort in Kreuz und Noth geübet, komm, hier ist Gottes Wort! Dies wird dich schon erquicken, daß, wenn gleich Höll’ und Welt dich wollten unterdrücken, du doch behältst das Feld.

4. Hast du dich lassen blenden, so daß bald hier bald dort du tappest an den Wänden, komm, hier ist Gottes Wort! dies machet, daß die Blinden sich zu dem rechten Steg hinwieder können finden von ihrem Irrweg.

5. Wirst du gleich auch geführet durch den stockfinstern Ort, da sonst der Tod regieret, komm, hier ist Gottes Wort! Dies ist der Stab und Stecken, mit diesem kannst du dich vor’s Teufels List und Schrecken beschützen mächtiglich.

6. Hilf, JEsu, daß ich liebe dein seligmachend Wort, daß ich mich stets drin übe; hilf, liebster Seelenhort, daß ich’s in meinem Herzen bewahr durch deine Huld, damit in Kreuzes Schmerzen es Frucht trag’ in Geduld!

Anna Sophia, Landgr. von Hessen.

\quellseite{502}

\chapter{IV. Von der heiligen Taufe}

Mel. Wer nur den lieben Gott läßt walten.

38. Ich bin getauft auf deinen Namen, Gott Vater, Sohn und Heil’ger Geist. Ich bin gezählt zu deinem Samen, zum Volk, das dir geheiligt heißt. Ich bin in Christum eingesenkt, ich bin mit seinem Geist beschenkt.

2. Du hast zu deinem Kind und Erben, mein lieber Vater, mich erklärt. Du hast die Frucht von deinem Sterben, mein treuer Heiland, mir gewährt. Du willst in aller Noth und Pein, o guter Geist, mein Tröster sein.

3. Doch hab’ ich dir auch Furcht und Liebe, Treu’ und Gehorsam zugesagt. Ich habe mich aus reinem Triebe dein Eigenthum zu sein gewagt. Hingegen sagt’ ich bis in’s Grab des Satans schnöden Werken ab.

4. Mein treuer Gott, auf deiner Seite bleibt dieser Bund wohl feste steh’n. Wann aber ich ihn überschreite, so laß mich nicht verloren geh’n. Nimm mich, dein Kind, zu Gnaden an, wenn ich hab’ einen Fall gethan.

5. Ich gebe dir, mein Gott, auf’s neue Leib, Seel’ und Herz zum Opfer hin. Erwecke mich zu neuer Treue und nimm Besitz von meinem Sinn. Es sei in mir kein Tropfen Blut, der nicht, HErr, deinem Willen thut.

6. Weich’, weich’, du Fürst der Finsternissen, ich bleibe mit dir unvermenget. Hier ist zwar ein befleckt Gewissen, jedoch mit Jesu Blut besprengt. Weich’, eit’le Welt, du Sünde weich’, Gott hört es, ich entsage euch.

7. Laß diesen Vorsatz nimmer wanken, Gott Vater, Sohn und Heil’ger Geist. Halt’ mich in deines Bundes Schranken, bis mich dein Wille sterben heißt. So leb’ ich dir, so sterb’ ich dir, so lob’ ich dich dort für und für.

J. J. Rambach.

Mel. Es woll’ uns Gott genädig sein.

39. O welch ein unvergleichlich’s Gut gibst du, HErr, deinen Kindern! Das Wasser und zugleich das Blut verehrest du den Sündern. Drei Dinge sind, welch’ allermeist auf Erden Zeugniß geben: das Blut, das Wasser und der Geist, die können uns erheben zu deinem Freudenleben.

2. Dies Sacrament ist selbst durch dich geheiligt und beschlossen, daß, wie du, HErr! bist sichtbarlich mit Wasser ganz begossen im Jordan durch Johannis Hand; so soll auch uns rein machen dein heilig Blut, das theure Pfand, das lauter Himmelsachen kann wirken in uns Schwachen.

3. Du hast uns durch dies Sacrament der Kirche einverleibet, also, daß man uns Christen nennt, und in dein Buch ist schreibet: Dies Wassersbad hat uns im Wort auch rein gemacht von Sünden; dein guter Geist, der will hinfort die Herzen recht entzünden, und Lieb’ in ihnen gründen.

4. Wir sind, HErr! in dein Gnadenreich durch diesen Bund gesetzet, der uns an Leib und Seel’ zugleich recht inniglich ergötzet. Du hast uns
\quellseite{503}
durch dies reine Bad so trefflich schön bekleidet, daß auch hinfort von deiner Gnad’ uns selbst der Tod nicht scheidet, noch Alles, was uns neidet.

5. Aus Höllenkindern sind wir schon der Gnaden Kinder worden; dies ist der Christen schönste Kron’ und Schmuck in ihrem Orden; ja Christus selber und sein Blut, sein Tod und Sieg darneben, ist nunmehr unser eig’nes Gut, das er uns hat gegeben, mit ihm dadurch zu leben.

6. Er hat uns auch das Kinderrecht der Seligkeit geschenket; durch solches ist die Sünde schlecht in’s tiefe Meer versenket. Was können Teufel, Hölle, Tod, welch’ uns stets widerstunden, weil JEsus Christus alle Noth sammt ihnen überwunden? Nun ist das Heil gefunden!

7. HErr! laß uns doch, den Neben gleich, auch gute Früchte bringen, und aus der Welt nach deinem Reich im Glauben eifrig ringen; laß uns durch wahre Reu’ und Buß’ auch täglich mit dir sterben, demnach der alte Adam muß bis auf den Grund verderben, soll man dein Reich ererben.

8. Hilf, daß wir diesen Gnadenbund der Taufe nie vergessen, und sich kein freches Herz noch Mund zu schmähen ihn vermessen. Die Taufe muß in Angst und Pein, ja, wenn wir geh’n von hinnen, HErr! unser Trost und Freude sein, das heißt der Welt entrinnen, den Himmel zu gewinnen.

Joh. Rist.

\chapter{V. Vom heiligen Abendmahl}

Mel. Ich hab’ mein’ Sach’ Gott.

40. Ich weiß ein Blümlein hübsch und fein, das thut mir wohl gefallen. Es liebt mir in dem Herzen mein das Blümelein, für andern Blümlein allen.

2. Das Blümlein ist das göttlich’ Wort, das uns Gott hat gegeben. Es leucht’t uns durch die enge Pfort’, ja hier und dort, wohl in das ew’ge Leben.

3. Er ist der Weg, das Licht, die Pfort’, die Wahrheit und das Leben. Wer Reu’ für seine Sünde hat, und bitt’t um Gnad’, dem sind’s im Glaub’n vergeben.

4. Er spricht: Kommt Alle her zu mir, All’, die ihr seid beladen; ich will nach eures Herzens Begier, das glaubet mir, heilen all’ euren Schaden.

5. Nehmt hin und eßt, das ist mein Leib, den ich euch jetzt thu’ schenken; ich verschreib’ euch all’ mein Gut dabei, das glaubet frei, daß ihr mein sollt gedenken.

6. Nehmt hin, trinket, das ist mein Blut, das ich für euch vergossen, welch’s g’nug für eure Sünde thut, so oft ihr’s thut, wie ich’s euch hab’ gelassen.

7. Wir bitten dich, HErr JEsu Christ! wohl durch dein bitter Leiden, weil du für uns gestorben bist, HErr JEsu Christ! du woll’st nicht von uns scheiden.

8. Nimm uns für deine Kinder an,
\quellseite{504}
daß wir dich all’zeit loben, dein Wort bekenn’n vor jedermann, auf rechter Bahn, durch Jesum Christum, Amen.

Basilius Fritsch.

\subsection{Wenn man communiciren will}

Mel. Herr Jesu Christ, mein’s Lebens Licht.

41. In Gottes Namen geh’ ich hin zum Herren, dem ich täglich dien’, daß er sein’ Gut’ an mir beweis’ und mit sein’m Fleisch und Blut mich speis’.

2. Viel tausend Engel umher steh’n, die all’ gelüstet, ihn zu seh’n. Soll ich denn nicht voll Andacht sein, daß ich würdig für ihn erschein’?

3. Herr Christ, mein leiblich Aug’ nicht sieht, es kann’s Vernunft auch fassen nicht. Jedoch mein Glaube jetzt erblickt dein Fleisch und Blut, welch’s mich erquickt.

4. Ich bin verderbt durch’s Sündengift, d’rum hast du Hilf’ und Heil gestift’t, daß meiner Seel’ werd’ Rath geschafft, dein Fleisch und Blut wirkt solche Kraft.

5. Ich bin ein armer Patient, mein Herz im Leib’ ist mir entbrennt, an Leib und Seel’ ich bin verwund’t, dein Fleisch und Blut mach’ mich gesund.

6. Ich komm’ zu dir, weil ich bin schwach, und fühl’ an mir groß Ungemach, nach dir ich Durst und Hunger hab’, dein Fleisch und Blut mich stärkt und lab’.

7. Gar arm und dürftig bin ich hier, nach dem Gut steht all’ mein’ Begier, des Himmels Schatz ich mir erwähl’, dein Fleisch und Blut mach’ reich mein’ Seel’.

8. Ohn’ dich bin ich ganz geistlich blind, daß ich das Gut’ nicht seh’ noch find’; ich thu’ das Bös’, welch’s mich gut deucht, dein Fleisch und Blut mein Herz erleucht’.

9. Ich hab’ nicht Lust zu Gottes Zucht, d’rum bin ich durch die Sünd’ verrucht; kein Wohlthat ist auf meinen Weg’n, dein Fleisch und Blut bringt mir den Seg’n.

10. Ich war verstoßen hin im Zorn, daß ich sollt’ ewig sein verlor’n, weil ich viel Sünden hab’ vollbracht, dein Fleisch und Blut mich selig macht.

11. Nimm mich ins Himmelreich zu dir, mein Leib und Seel’ mit Klarheit zier’; dahin steht all’ mein Herz und Sinn, dein Fleisch und Blut helf’ mir dahin.

Torgauisches Handbuch von 1676.

Mel. Ach Gott und Herr!

42. Mein’ Seel’! dich freu’, und lustig sei, mit Glauben wohl gezieret; zur Mahlzeit schön wirst du heut’ geh’n, zu der dich Christus führet.

2. Merk’ auf mit Fleiß, die werth’ Speis’, sein Leib, für dich gegeben; der Trank ist gut, sein theures Blut stärkt dich zum ew’gen Leben.

3. Wenn deine Sünd’ dich woll’n geschwind’ zur Höllen niederdrücken, dies’r edle Tisch machet dich frisch, thut dich lieblicherquicken.

4. Ihr Sünder schwer! kommt doch hieher, die Last legt von dem Herzen; der Arzt so reich, Christus, wird euch heilen all eure Schmerzen.

5. Zum Gastmahl hier all’ pflegen wir mit Kleidern uns zu schmücken; schmück’ dich auch fein im Herzen dein, thu’ dich recht darzu schicken.

\quellseite{505}

6. Mit Himmelbrod versieht dich Gott bei diesem Tisch von oben; der Fels dich tränkt, das Leb’n dir schenkt, d’rum thu’ ihn herzlich loben.

7. Von Herzen ich will freuen mich, daß mich der HErr geladen; er ist ja mein und ich bin sein, mein’r Seel’ ist wohl gerathen.

8. Sehr ich beklag’, daß Mancher mag nicht diese Freud’ bedenken, und läßt sich nicht mit Zuversicht aus die- sem Kelch so tränken.

9. All’ ich komm’ heut’ mit großer Freud’, HErr Christ, zu deinem Tische, und stell’ mich ein zur Mahlzeit rein, mein’n Leib und Seel’ erfrische.

10. Wasch’ mich auch rein von Sünden mein, du höchster Gott! aus Gnaden, so kann mir nicht der Bösewicht an Leib und Seele schaden.

11. Ein ängstlich Herz, zerschlag’n mit Schmerz von wegen seiner Sün- den, bitt’ sehnlich ich, veracht’ ja nicht, laß mich Verzeihung finden.

12. Kein Opfer sonst ich deiner Gunst auf dieser Welt kann bringen. O HErr! laß mein erschrock’n Gebein für Freuden nieder springen.

13. Die Mahlzeit ist, HErr Christ! und mich vereint’t himmlischer Weise; bleib’ doch in mir, und ich in dir, daß ich dich ewig preise.

14. Gewiß ich weiß, daß da mit Fleiß viel tausend Engel stehen, wenn wir so sein im Glauben rein zum Tisch des HErren gehen.

15. Wenn in der Näh’ ich recht anseh’ den Kelch in dieser Stunden, denk’ ich, dein Blut fließ’, als ein’ Fluth, aus deinen heil’gen Wunden.

(Die Vernunft kann dies hohe Sacrament nicht ergründen.)

16. Wie es zugeh’, ich nicht ver- steh’, und will nicht disputiren. Wort und Element ein Sacrament heilig constituiren.

17. In dieser Sach’ bin ich zu schwach, gar wunderbarer Weise der Glaub’ allein (Glaube rein) führt’s Herze mein zu der herrlichen Speise.

18. Aus Priesters Hand dies edle Pfand mir Armen wird gegeben, da- durch ich mich stark’ sicherlich, daß ich werd’ ewig leben.

19. Ach Gott! wie stark durch Bein und Mark die Freud’ mir ist thut dringen! Wie sehn’ ich mich, HErr Christ! durch dich nach himm- lischen Dingen!

20. Mich dünkt, als sei der Him- mel frei, wie Stephano, mir offen; sein’ Liebe hat in höchster Gnad’ mein’ Seel’ so süß getroffen.

21. Zur Freude mein Violen rein beweglich laßt erklingen; von Herzen thut Motetten gut in Harmonie bald singen.

22. Die Orgeln auch nach altem Brauch im Tempel Gottes schöne, Posaunenklang zier’ den Gesang mit lieblichen Getöne.

23. Auf daß die Gäst’ zu diesem Fest hierüber jubiliren, und sich mög’ heur’ ihr’ Traurigkeit aus ihrem G’müth verlieren.

24. Ei! daß nur bald die Engel all’ (die Engl’ mit Schall) mein’ Seel’ führt’n aus dem Leben in Himmelsthron, da mir ein’ Kron’ die Hand des HErrn wird geben.

\quellseite{506}

25. Herr! laß die Freud’ zu keiner Zeit aus meinem Herzen weichen; dein’s Geists gewiß wahrlich ist dies, der in mir wohnt, ein Zeichen.

26. O! daß ich sollt’, wie gern ich wollt’, dein Antlitz nur bald schauen! doch ich das will in Gottes Still’ erwarten mit Vertrauen.

27. Indessen ich ergeb’ dir mich, und laß im Kreuz dich walten. Ich werd’ einmal im Himmelsaal die ew’ge Tafel halten.

28. Dasselb’ weiß ich ganz sicherlich, Christus hat mir’s erworben, als er für mich ganz bitterlich ist an dem Kreuz’ gestorben.

\chapter{VI. Von der Buße}

Mel. Zion klagt mit Angst und Schmerzen.

43. Ach, wie will es endlich werden, ach, wie will es laufen ab, daß ich stets hang’ an der Erden und so wenig Andacht hab’, ja daß ich im harten Sinn ohne Buß’ geh’ immer hin. JEsu, mir mein Herze rühre, mich zur wahren Buße führe.

2. Ach, ich fühle keine Reue, und mein Sündenherz ist kalt. Ich erkenne ohne Scheue meine Fehler mannigfalt, und doch soner Traurigkeit lebe ich in Sicherheit. JEsu, mich doch so regiere, daß ich wahre Buße spüre.

3. Ach, ich kann an Gott nicht denken, wie ich sollt’ und gerne wollt’. Denn mein Geist, der läßt sich lenken dahin, wo er fliehen soll’, und klebt dem zu feste an, was zum Abgrund ziehen kann. JEsu, hilf mein Herz aus Gnaden dieser schweren Last entladen.

4. Ach, weil denn mein Geist erstorben und mein Fleisch hält Ueberhand, seh’ ich kaum, was ich erworben, da ich mich von Gott gewandt. Mein Verderben merk’ ich kaum, weil ich gab dem Fleische Raum.

5. Ach, so wache doch von Sünden auf, o du mein sich’rer Geist! Suche Gott, weil er zu finden, weil er noch dein Vater heißt. Komme, weil die Gnadenthür’ nun noch stehet offen dir. JEsu, zu so gutem Werke mich mit Herzensandacht stärke.

6. Ach, mein Herz, bedenkt’ die Schuldte, schau’, wie weit du bist von Gott, wie du seine Gnad’ und Hulde hast verworffen und verspott’t, und wie du aus seiner Ruh’ selbst der Höllen liefst zu. JEsu, mir mein Herze lenke, daß es diese Noth bedenke.

7. Ach, weil denn noch Zeit, zu kehren, so verlaß ich diesen Stand. Du, mein JEsu, woll’st mich hören und mir bieten deine Hand. Hilf mir Schwachen, hilf mir auf, steu’re meinem Sündenlauf. Herr, um deiner Mutter willen willst du meinen Wunsch erfüllen.

8. Ach, ich kann schon tröstlich finden, JEsu, hier in deinem Blut die Abwaschung meiner Sünden.

\quellseite{507}

Dieses stärket meinen Muth, dies erfrischet meinen Sinn, ob ich gleich ein Sünder bin. Jesu, mich forthin regiere, daß ich ein fromm Leben führe. Joh. Angelus.

Mel. Nun sich der Tag geendet.

\subsection{44}

Mein Gott! das Herze bring’ ich dir zur Gabe und Geschenk; du forderst dieses ja von mir, deß bin ich eingedenk.

2. Gib mir, mein Sohn! dein Herz, sprichst du, das ist mir lieb und werth; du findest anders auch nicht Ruh’ im Himmel und auf Erd’.

3. Nun du, mein Vater! nimm es an, mein Herz, veracht’ es nicht; ich geb’s, so gut ich’s geben kann, kehr’ zu mir dein Gesicht.

4. Zwar ist es voller Sündenwust und voller Eitelkeit; des Guten aber unbewußt, der wahren Frömmigkeit.

5. Doch aber steht es nun in Reu’, erkennt sein’n Uebelstand, und träget \unsicher{übend} vor dem Scheu, d’ran es zuvor Lust fand.

6. Hier fällt und lieget es zu Fuß, und schreit: Nur schlage zu! zerknirsch’, o Vater, daß ich Buß’ rechtschaffen vor dir thu’.

7. Zermalm’ mir meine Härtigkeit, mach’ mürbe meinen Sinn, daß ich in Seufzer, Reu’ und Leid und Thränen ganz zerrinn’.

8. Sodann nimm mich, mein Jesu Christ, tauch’ mich tief in dein Blut; ich glaub’, daß du getreuzigt bist mir und der Welt zu gut.

9. Stärk’ mein’ sonst schwache Glaubenshand, zu fassen auf dein Blut, als der Vergebung Unterpfand, das Alles machet gut.

10. Schenk’ mir nach deiner Jesus-Huld Gerechtigkeit und Heil, und nimm auf dich mein’ Sündenschuld und mein verdientes Theil.

11. In dich woll’st du mich kleiden ein, dein’ Unschuld ziehen an, daß ich, von allen Sünden rein, vor Gott bestehen kann.

12. Gott Heil’ger Geist, nimm du auch mich in die Gemeinschaft ein; ergieß um Jesu willen dich tief in mein Herz hinein.

13. Dein göttlich Licht schütt’ in mich aus, und Brunst der reinen Lieb’; lösch’ Finsterniß, Haß, Falschheit aus, schenk’ mir stets deinen Trieb.

14. Hilf, daß ich sei von Herzen treu im Glauben meinem Gott; daß mich im Guten nicht mach’ scheu der Welt List, Macht und Spott.

15. Hilf, daß ich sei von Herzen fest im Hoffen und Geduld; daß, wenn du nur mich nicht verläßt, mich tröste deine Huld.

16. Hilf, daß ich sei von Herzen rein im Lieben, und erweis’, daß mein Thun nicht sei Augenschein, durch’s Werk zu deinem Preis.

17. Hilf, daß ich sei von Herzen schlecht, aufrichtig, ohn’ Betrug; daß meine Wort’ und Werke recht, und niemand schelt’ ohn’ Fug.

18. Hilf, daß ich sei von Herzen klein, Demuth und Sanftmuth halt’, daß ich, von aller Weltlieb’ rein, vom Fall’ aufstehe bald.

19. Hilf, daß ich sei von Herzen fromm, ohn’ alle Heuchelei; damit
\quellseite{508}
mein ganzes Christenthum dir wohlgefällig sey.

20. Nimm gar, o Gott, zum Tempel ein mein Herz hier in der Zeit; ja, laß es auch dein Wohnhaus sein in jener Ewigkeit.

21. Dir geb’ ich’s ganz zu eigen hin, brauch’s, wozu dir’s gefällt; ich weiß, daß ich der Deine bin, der Deine, nicht der Welt.

22. D’rum soll sie nun und nimmermehr dies richten aus bei mir; sie lock’ und droh’ auch noch so sehr, daß ich soll dienen ihr.

23. In Ewigkeit geschieht das nicht, du falsche Teufelsbraut! gar wenig mich, Gott Lob! anschar dein’ glänzend’ Schlangenhaut.

24. Weg, Welt! weg, Sünd’! dir geb’ ich nicht mein Herz: nur JESU! dir ist dies Geschenke zugericht’t, behalt’ es für und für.

J. C. Schade.

\chapter{VII. Von Glauben und Rechtfertigung}

Mel. Ich ruf’ zu dir, HErr JESU Christ.

45. Also hat Gott die Welt geliebt, daß er Christum hat geben, sein’n lieben Sohn, in Tod betrübt, daß wir nun ewig leben. O Gott, wie groß ist deine Lieb’, die kein Mensch kann von Herzen ganz ermessen! O HErr Gott, gib Gnad’, daß wir des nimm’r vergessen.

2. Christus hat uns zum Himmel bracht ohn’ aller Menschen Stärke. Allein der Glaube selig macht ohn’ alle unsre Werke. Darum, wer glaubt an JESum Christ, der hat das ewig’ Leben, merkt gar eben, dann wo nicht Glaube ist, mag kein Mensch selig werden.

3. Wer nicht glaubt, der ist schon verdammt, als Christus selbst thut sagen: nimmermehr er sich des erbarmt, der trotzt auf Werk’ ohn’ Glauben. An Leib und Seel’ wird er geplagt, kann auch kein’n Trost erlangen, liegt gefangen allhier ohn’ alle Gnad’ ewig in Todesbanden.

4. Wer glaubig bleibt bis an sein End’, der wird g’wiß selig werden, und Christum frei ohn’ Furcht bekennt für aller Welt auf Erden, denn wird er auch im höchsten Thron vor Gott mit Chr’n bekennen, und ihn nennen seinen Miterben schon, deß Freud’ wird sein kein Ende.

Altes Lied aus Crügers Praxis pietatis.

Mel. Gott Lob! es geht nunmehr.

46. Wenn so viel tausend Teufel wären, als an den Ufern Körnchen Sand, als Tropfen sind in allen Meeren, als Sonnenstäubchen durch das Land; so fürcht’ ich doch dieselben nicht, denn JESus ist mein Heil und Licht.

2. Und wären zehnmal so viel Höllen, als zwischen Bergen Klüfte sind, als Nitz’ in Felsen, die zerfielen, als man in Wäldern Höhlen find’t; so fürcht’ ich doch dieselben nicht, denn JESus ist mein Heil und Licht.

3. Wenn so viel Feinde zu mir
\quellseite{509}
bringen, als das Gewässer Fische hegt, als Vögel unterm Himmel singen, als Luft und Erde Wärmer trägt; so fürcht’ ich doch dieselben nicht, denn JEsus ist mein Heil und Licht.

4. Ja gäb’ es auch mehr Lästerzungen, als man auf Ackern Aehren zählt, als jemals Gräschen sind entsprungen, als man in Hecken Stacheln sieht; so fürcht’ ich doch dieselben nicht, denn JEsus ist mein Heil und Licht.

5. Brüllt, ihr vermaledeyten Drachen! Ihr Höllengluthen, macht euch groß! Ihr Feinde, laßt die Stücke krachen! Ihr Lügner, schmäht nur auf mich los! Ich fürchte mich für euch doch nicht, denn JEsus ist mein Heil und Licht.

6. Was werft, ihr Teufel! mir die Sünden mit so ergrimmtem Schreien für? Wo könnt ihr was von Schulden finden? Ich habe ja die Quittung hier; denn JEsus ist mein Heil und Licht, d’rum fürcht’ ich euer Schreien nicht.

7. Ihr sprüht umsonst, ihr Höllenflammen! ich werde durch euch nicht erschreckt, ihr sollt mich nimmermehr verdammen, weil mich des Heilands Unschuld deckt. Ist JEsus nur mein Heil und Licht, so fürcht’ ich euren Abgrund nicht.

8. Was wollt ihr Feinde mit mir machen? Verfolgt, verjagt und ängstet mich, verbrennt und raubt mein bißchen Sachen, erwürgt, zerreißt mich grausamlich. Ich fürchte dieses Alles nicht, denn JEsus ist mein Heil und Licht.

9. Und laßt doch seh’n, ihr losen Mäuler! was habt ihr Junges ausgedreckt? Hier steht mein fester Ehrenpfeiler, der hält mich vor euch unbefleckt. Mein JEsus ist mein Heil und Licht, d’rum acht’ ich euer Lästeren nicht.

10. So biet’ ich nun den bösen Geistern, der Höll’ und allen Feinden Trutz! Trutz sey auch allen Lügenmeistern! denn JEsus ist mein sich’rer Schutz. Ist der nur stets mein Heil und Licht, so fürcht’ ich das geringste nicht.

11. Ach JEsu! laß doch mein Vertrauen auf dich ja nicht vergeblich seyn, so wird mir nie vor etwas grauen. Erbarm’, erbarm’, erbarm’ dich mein! Bleib’, JEsu! du mein Heil und Licht, so fürcht’ ich mich für’m Tode nicht.

J. Mencer.

\chapter{VIII. JEsus-Lieder}

Mel. Wer nur den lieben Gott läßt walten.

47. Allein, und doch nicht ganz alleine bin ich in meiner Einsamkeit; denn wenn ich ganz verlassen scheine, vertreibt mir JEsus selbst die Zeit. Ich bin bei ihm und er bei mir, so kommt mir gar nichts einsam für.

2. Komm ich zur Welt, man red’t von Sachen, die nur auf Eitelkeit gericht’t; da muß sich lassen der verlachen, der etwas von dem Himmel spricht; d’rum wünsch’ ich lieber ganz allein, als bei der Welt ohn’ Gott zu sein.

\quellseite{510}

3. Verkehrte können leicht verkehren; wer greift doch ohne Kleben an? Wie sollt’ ich denn dahin begehren, wo man Gott bald vergessen kann? Gesellschaft, die verdächtig scheint, wird oftmals nach dem Fall beweint.

4. Zudem kann sich ein Mensch verstellen; wer will in aller Herzen seh’n? Man sieht oft heimliche Gesellen, die sich nur nach dem Winde dreh’n, daß der, so vor voll Zucker wär, bald eine Schlange d’rauf gebar.

5. D’rum kann mich niemand hier verdenken, wenn ich in meiner Einsamkeit mich also suche zu beschränken, daß Gott allein mein Herz erfreut. Die Welt ist voller Trug und List; wohl dem, der Gott verbunden ist!

6. Ein Erbfeind mag Gesellschaft suchen; ich liebe Gott in stiller Ruh’. Und sollten mir die Neider fluchen, so schließ’ ich meine Kammer zu und nehme Gott mit mir hinein; so wird der Neid betrogen sein.

7. Mit Gott red’ ich in seinem Worte, und durch sein Wort red’t er mit mir: Bet’ ich an einem stillen Orte, so findet er sich bald zu mir; brauch’ ich in meinem Kreuze Rath, er ist, der mich getröstet hat.

8. Mach’ ich im Stillen meine Sachen, so hat er seine Hand bei mir; ich mag nun schlafen oder wachen, so stellt er mir sein Bildniß für. Auch in dem Traume spielet er, als wenn ich wachend bei ihm wär’.

9. Wer wollte denn nun nicht erkennen, daß ich stets in Gesellschaft bin? Und will die Welt mich einsam nennen, so thue sie es immerhin. G’nug, daß bei mir, wenn ich allein, Gott und viel tausend Engel sein.

Benz. Schmold.

In eigener Melodie.

48. Ich laß dich nicht! du mußt mein Jesus bleiben. Will raube Noth, Welt, Höll’ und Tod mich aus dem Feld’ erbitter Treue treiben? Nur her! ich halte mich, mein starker Held, an dich. Hör’, was mein Herze spricht: Du mußt mein Jesus bleiben, ich laß dich nicht! ich laß dich nicht!

2. Ich laß dich nicht, du allerhöchste Liebe! Wann Zweifel sich setzt wider mich, ich weiß, wie dich die keusche Flamme triebe. Du trugest Schuld und Pein, sollt’ ich verwurtheilt sein an jenem Weltgericht? Du allerhöchste Liebe, ich laß dich nicht! ich laß dich nicht!

3. Ich laß dich nicht, du süße Seelenstärke, die mich stets labt, mit Kraft begabt, wenn ich in mir des Glaubens Schwachheit merke. Macht Krankheit gleich den Leib durch Schmerzenächte schwach, so spricht die Seele doch: Du süße Seelenstärke, ich laß dich nicht! ich laß dich nicht!

4. Ich laß dich nicht, du Hülff’ in allen Nöthen! Leg’ dich auf Joch, ich hoffe doch, auch wenn es scheint, als wolltest du mich tödten. Mach’s, wie du willst, mit mir, ich weiche nicht von dir; verfehle dein Gesicht, du Hülff’ in allen Nöthen, ich laß dich nicht! ich laß dich nicht!

5. Ich laß dich nicht! Sollt’ ich den Segen lassen? Nein, Jesu, nein,
\quellseite{511}
du bleibest mein. Dich halt’ ich noch, wenn ich nichts mehr kann fassen. Nach kurzer Nächte Lauf geht mir der Segen auf von dir, dem Segenslicht. Sollt’ ich den Segen lassen? Ich laß dich nicht! ich laß dich nicht!

6. Ich laß dich nicht! Führ’ mich nach deinem Willen, ich folge nach durch Wohl und Ach, dein weiser Schluß kann allen Kummer stillen. Dir, Jesu, hang’ ich an, und achte keine Bahn, wo mich die Distel sticht; führ’ mich nach deinem Willen, ich laß dich nicht! ich laß dich nicht!

7. Ich laß dich nicht auch in dem Schooß der Freude; denn wenn ich mich seh’ ohne dich, so ist die Lust mir eine Wermuthweide. Mir graut für ihrer Kost, wenn nicht von deinem Trost mein Herz durchfesset spricht, auch in dem Schooß der Freude: Ich laß dich nicht! ich laß dich nicht!

8. Ich laß dich nicht! was will die Hölle haben? Herr, ich bin dein, wie kann ich sein als deine Taub’ auch bei verdammten Raben? Mich reiniget dein Blut: was drohet dann mit Gluth ihr rauchendes Gesicht? Was will die Hölle haben? Ich laß dich nicht! ich laß dich nicht!

9. Ich laß dich nicht, mein Gott! mein Herr! mein Leben! Mich reißt das Grab von dir nicht ab, der du dich hast für mich in Tod gegeben; du starbst aus Liebe mir, ich sag’ in Liebe dir, auch wenn das Herz zerbricht, mein Gott! mein Herr! mein Leben! Ich laß dich nicht! ich laß dich nicht!

W. E. Tesker.

Mel. Jesu, der du meine Seele.

49. Jesu, frommer Menschenheerden guter und getreuer Hirt! Laß mich auch dein Schäflein werden, daß dein Stab und Stimme führt. Ach! du hast aus Lieb’ dein Leben für die Schafe hingegeben, und du gabst es auch für mich; laß mich wider lieben dich.

2. Heerden ihren Hirten lieben, und ein Hirt liebt seine Heerd’. Laß uns auch so Liebe üben, du im Himmel, ich auf Erd’. Schallet deine Lieb’ hernieder, soll dir meine schallen wieder; wenn du rufst: Ich liebe dich! ruft mein Herz: Dich liebe ich!

3. Schafe ihren Hirten kennen, denn sie auch sind wohl bekannt. Laß mich auch nach dir nur rennen, wie du samst zu mir gerannt. Als des Höllenvolkes Rachen eine Beut’ aus mir wollt’ machen, riefest du: Ich kenne dich. Ich ruf’ auch: Dich kenne ich.

4. Heerden ihren Hirten hören, folgen seiner Stimm’ allein. Hirten auch zur Heerd’ sich kehren, wenn sie rufen Groß und Klein: Laß mich hören, wenn du schreiest; laß mich laufen, wenn du bräuest; laß mich horchen stets auf dich; Jesu, höre du auch mich.

5. Höre, Jesu! und erhöre meine, deines Schäfleins, Stimm’; mich auch zu dir schreien lehre, wenn sich naht des Wolfes Grimm. Laß mein Schreien dir gefallen, deinen Trost herniederschallen; wenn ich bete, höre mich; Jesu sprach: Ich höre dich!

6. Höre, Jesu, und erhöre, wenn ich ruf’, anklopf’ und schrei’; Jesu!

\quellseite{512}

dich nicht von mir kehre, steh’ mir bald in Gnaden bei. Ja, du hörst, in deinem Namen ist ja Alles Ja und Amen. Nun, ich glaub’ und fühle schon deinen Trost, o Gottes Sohn!

Sigmund v. Birken.

Mel. Den die Hirten lobten sehre.

50. Kommt und laßt uns Christum ehren, Herz und Sinnen zu ihm kehren. Singet fröhlich, laßt euch hören, werthes Volk der Christenheit.

2. Sünd’ und Hölle mag sich grämen, Tod und Teufel mag sich schämen. Wir, die unser Heil annehmen, werfen allen Kummer hin.

3. Sehet, was hat Gott gegeben, seinen Sohn zum ew’gen Leben. Dieser kann und will uns heben aus dem Leid in’s Himmels Freud’.

4. Seine Seel’ ist uns gewogen, Lieb’ und Gunst hat ihn gezogen, uns, die Satanas betrogen, zu besuchen aus der Höh’.

5. Jacobs Stern ist aufgegangen, stillt das sehnliche Verlangen, bricht den Kopf der alten Schlangen und zerstört der Höllen Reich.

6. Unser Kerker, da wir saßen und mit Sorgen ohne Maßen uns das Herze selbst abfraßen, ist entzwei und wir sind frei.

7. O du hochgesegn’te Stunde, da wir das von Herzensgrunde glauben und mit unserm Munde danken dir, o Jesulein!

8. Schönstes Kindlein in dem Stalle, sei uns freundlich, bring’ uns Alle dahin, da mit süßem Schalle dich der Engel Heer erhöht.

P. Gerhardt.

Mel. Schmücke dich, o liebe Seele.

51. König, dem kein König gleichet, dessen Ruhm kein Mund erreichet, dem als Gott das Reich gebühret, der als Mensch das Scepter führet, dem das Recht gehört zum Throne als des Vaters ew’gem Sohne, den so viel Vollkommenheiten krönen, zieren und begleiten.

2. Himmel, Wasser, Luft und Erde nebst der ungezählten Heerde der Geschöpfe in den Feldern, in den Seen, in den Wäldern, sind, HErr über Tod und Leben, dir zum Eigenthum gegeben. Thiere, Menschen, Geister scheuen, Menschensohn, dein mächtig Dräuen.

3. In des Gnadenreiches Grenzen stehst man dich am schönsten glänzen, wo viel tausend treue Seelen dich zu ihrem Haupt erwählen, die durchs Scepter deines Mundes nach dem Recht des Gnadenbundes sich von dir regieren lassen und wie du das Unrecht hassen.

4. In dem Reiche deiner Ehren kann man stets dich loben hören von dem himmlischen Geschlechte, von der Menge deiner Knechte, die dort ohne Furcht und Grauen dein verklärtes Antlitz schauen, die dich unermüdet preisen und dir Ehr’ und Dienst erweisen.

5. O Monarch in dreien Reichen, dir ist niemand zu vergleichen an dem Ueberfluß der Schätze, an der Ordnung der Gesetze, an Vortrefflichkeit der Gaben, welche deine Bürger haben. Du beschüßest deine Freunde, du bezwingest deine Feinde.

\quellseite{513}

6. Herrsche auch in meinem Herzen über Zorn, Furcht, Lust und Schmerzen. Laß mich deinen Schutz genießen, laß mich dich im Glauben küssen, ehren, fürchten, loben, lieben und mich im Gehorsam üben, hier mit Leiden, kämpfen, schwitzen, dort mit auf dem Throne sitzen.

J. J. Rambach.

In eigener Melodie.

52. Wer Jesum bei sich hat, kann feste stehen, wird’ auf dem Unglücksmeer nicht untergehen. Wer Jesum bei sich hat, was kann dem schaden? sein Herz ist überall mit Trost beladen.

2. Wer Jesum bei sich hat, der hat den Himmel, wünscht zu verlassen nur das Weltgetümmel. Wer Jesum bei sich hat, der lebt vergnüget mit dem, was Gott und Glück ihm zugefüget.

3. Wer Jesum bei sich hat, der mag nicht haben die Eitelkeit der Welt und ihre Gaben. Wer Jesum bei sich hat, hat g’nug auf Erden und mag in Ewigkeit nicht reicher werden.

4. Wer Jesum bei sich hat, kann sicher reisen, er wird ihm schon den Weg zum Himmel weisen. Wer Jesum bei sich hat in höchsten Nöthen, den kann sein Teufel nicht, noch Mörder tödten.

5. Wer Jesum bei sich hat, ist wohl beschützet, wenn heftig donnert es und schrecklich blitzet. Wer Jesum bei sich hat, darf nicht erschrecken, wenn seine Sünd’ ihm Furcht und Angst erwecken.

6. Wer Jesum bei sich hat, darf nicht verzagen, und kann den Teufel auch leicht von sich jagen. Wer Jesum bei sich hat, wird nicht verderben, wer Jesum bei sich hat, kann fröhlich sterben.

7. Wer Jesum bei sich hat, kann nicht verderben; wer Jesum bei sich hat, kann selig sterben; wer Jesum bei sich hat, hat alle Fülle; wer Jesum bei sich hat, ist ruhig, stille.

E. F. Connnow.

Melodie: Seelenbräutigam.

53. Wer ist wohl wie du, Jesu, süße Ruh’, unter Vielen auserkoren, Leben derer, die verloren, und ihr Licht dazu, Jesu, süße Ruh’.

2. Leben, das den Tod, mich aus aller Noth zu erlösen, hat geschmecket, meine Schulden zugedecket und mich aus der Noth hat geführet zu Gott.

3. Glanz der Herrlichkeit, du bist vor der Zeit zum Erlöser uns geschenket, und in unser Fleisch versenket in der Füll’ der Zeit, Glanz der Herrlichkeit.

4. Größer Siegesheld, Tod, Sünd’, Höll’ und Welt, alle Kraft des großen Drachen hast du woll’n zu Schanden machen durch das Lösegeld deines Blut’s, o Held!

5. Höchste Majestät, König und Prophet, deinen Scepter will ich küssen, ich will sitzen dir zu Füßen, wie Maria thät, höchste Majestät.

6. Laß mich deinen Ruhm als dein Eigenthum durch des Geistes Licht erkennen, stets in deiner Liebe brennen, als dein Eigenthum, allerschönster Ruhm.

7. Beug mich ganz in dich, daß vor Liebe ich ganz zerinne und zerschmelze und auf dich mein Elend wälze, daß stets drücket mich. Beug mich ganz in dich.

\quellseite{514}

8. Deiner Sanftmuth Schild, deiner Demuth Bild mir anlege, in mich präge, daß kein Zorn noch Stolz sich rege. Vor dir sonst nichts gilt, als dein eigen Bild.

9. Steure meinem Sinn, der zur Welt will hin, daß ich nicht mög’ von dir wanken, sondern bleiben in den Schranken. Sei du mein Gewinn, gib mir deinen Sinn.

10. Wecke mich recht auf, daß ich meinen Lauf unverrückt zu dir fortsetze, und mich nicht in seinem Netze Satan halte auf. Fördre meinen Lauf.

11. Deines Geistes Trieb in die Seele gib, daß ich wachen mög’ und beten, freudig vor dein Antlitz treten. Ungefärbte Lieb’ in die Seele gib.

12. Wenn der Wellen Macht in der trüben Nacht will des Herzens Schifflein decken, woll’st du deine Hand ausstrecken. Habe auf mich Acht, Hüter in der Nacht!

13. Einen Heldenmuth, der da Gut und Blut gern um deinetwillen lasse und des Fleisches Lüste hasse, gib mir, höchstes Gut, durch dein theures Blut.

14. Soll’s zum Sterben geh’n, woll’st du bei mir steh’n, mich durch’s Todesthal begleiten und zur Herrlichkeit bereiten, daß ich einst mag seh’n mich zur Rechten steh’n.

J. A. Freylinghausen.

In eigener Melodie.

54. Wo ist JEsus, mein Verlangen, mein Geliebter und mein Freund? Wo ist er denn hingegangen? Wo mag er zu finden sein? Meine Seel’ ist sehr betrübt, mit viel Sünden abgematt’t; wo ist JEsus, den sie liebt? den begehrt sie Tag und Nacht.

2. Ach! ich ruf’ vor Angst und Schmerzen: Wo ist denn mein JEsus hin? Kein Ruh’ ist in meinem Herzen, so lang’, bis ich bei ihm bin. Ach! wer gibt mir Taubenflügel, daß ich kann zu jeder Frist fliegen über Berg’ und Hügel, suchen, wo mein JEsus ist?

3. Er vertreibet Angst und Schmerzen, er vertreibet Sünd’ und Tod, wenn sie quälen in dem Herzen; er hilft Jedem aus der Noth! Darum will ich nicht nachlassen, will bald laufen hin und her; auf den Feldern, auf den Straßen will ich suchen mehr und mehr.

4. Liebster JEsu! laß dich finden, meine Seele schreit in mir, thu’ ihr mit den Augen winken, laß sie eilends sein bei dir! Ach, laß mich die Gnad’ erlangen, allerliebster JEsu mein, und nimm meine Seel’ gefangen, laß sie ewig bei dir sein.

5. Ach! ich sterb’ vor lauter Freuden, ich find’ JEsum, meinen Schatz; alle Weltlust will ich meiden, bei ihm will ich finden Platz. Nunmehr soll mich nichts betrüben, was mich vor betrübet hat; nichts als JEsum will ich lieben, den mein’ Seel’ gefunden hat.

6. Liebster JEsu! meine Freude, meiner Seelen schönste Zier! Du bist meines Herzens Weide, mich verlangt ja stets nach dir. Mein Herz, Sinnen und Gedanken sollen dir ergeben sein; laß mich nimmer von dir wanken, du sollt ewig bleiben mein.

\quellseite{515}

\chapter{IX. Vom christlichen Leben}

Mel. Mach’s mit mir, Gott, nach deiner Güt’.

55. Auf, Christenmenschi auf, auf zum Streit! auf, auf zum Ueberwinden! In dieser Welt, in dieser Zeit ist keine Ruh’ zu finden. Wer nicht will streiten, trägt die Kron’ des ew’gen Lebens nicht davon.

2. Der Teufel kommt mit seiner List, die Welt mit Pracht und Prangen, das Fleisch mit Wollust, wo du bist, zu fällen dich und fangen. Streit’st du nicht wie ein tapfrer Held, so bist du hin und schon gefällt.

3. Gedenke, daß du zu der Fahn’ dein’s Feldherrn hast geschworen; denk’ ferner, daß du als ein Mann zum Streit bist auserkoren; ja denke, daß ohn’ Streit und Sieg nie keiner zum Triumph aufstieg.

4. Wie schmählich ist’s, wenn ein Soldat dem Feind den Rücken kehret; wie schändlich, wenn er seine Statt verläßt, und sich nicht wehret; wie spöttlich, wenn er noch mit Fleiß aus Zagheit wird dem Feind’ zum Preis!

5. Bind’ an! der Teufel ist bald hin, die Welt wird leicht verjaget; das Fleisch muß endlich aus dem Sinn, wie sehr dich’s immer plage. O ew’ge Schande, wenn ein Held vor diesen dreien Bub’n fällt!

6. Wer überwind’t, und kriegt den Raum der Feinde, die vermessen, der wird im Paradies vom Baum des ew’gen Lebens essen. Wer überwind’t, den soll kein Leid noch Tod berühr’n in Ewigkeit.

7. Wer überwind’t, und seinen Lauf mit Ehren kann vollenden, dem wird der Herr alsbald darauf verborg’nes Manna senden, ihm geben einen weißen Stein und einen neuen Namen d’rein.

8. Wer überwind’t, bekommt Gewalt, wie Christus, zu regieren, mit Macht die Völker mannigfalt in einer Schnur zu führen. Wer überwind’t, bekommt vom Herrn zum Feldpanier den Morgenstern.

9. Wer überwindet, der soll dort in weißen Kleidern gehen, sein guter Name soll sofort im Buch des Lebens stehen, ja Christus wird denselben gar bekennen vor der Engel Schaar.

10. Wer überwind’t, soll ewig nicht aus Gottes Tempel gehen, vielmehr d’rin, wie ein englisch Licht und güld’ne Säule, stehen; der Name Gottes, unsers Herrn, soll leuchten von ihm weit und fern.

11. Wer überwind’t, soll auf dem Thron mit Christo Jesu sitzen, soll glänzen wie ein Gottessohn, und wie die Sonne blitzen, ja ewig herrschen und regier’n, und immerdar den Himmel zier’n.

12. So streit’ denn wohl, streit’ fest und kühn, daß du mög’st überwinden! Streng’ an die Kräfte, Muth und Sinn, daß du dies Gut mög’st finden! Wer nicht will streiten um die Kron’, bleibt ewiglich in Spott und Hohn.

Joh. Angelus.

\quellseite{516}

Mel. Meine Hoffnung steht feste.

56. Auf, ihr Christen, Christi Glieder, die ihr noch hangt an dem Haupt! Auf! wacht auf! ermannt euch wieder, eh’ ihr werdet hingeraubt! Satan beut an den Streit Christo und der Christenheit.

2. Auf! folgt Christo, eurem Helde, trauet seinem starken Arm! Liegt der Satan gleich zu Felde mit dem ganzen Höllenschwarm, sind doch der’r noch viel mehr, die da stets sind um uns her.

3. Nur auf Christi Blut gemaget mit Gebet und Wachsamkeit! Dieses machet unverzaget und recht tapfre Kriegesleut’. Christi Blut gibt uns Muth wider alle Teufelsbrut.

4. Christi Heeres Kreuzesfahne, so da weiß und roth besprengt, ist schon auf dem Siegesplane uns zum Troste aufgehängt. Wer hier kriegt, nie erliegt, sondern unter’m Kreuze siegt.

5. Diesen Sieg hat auch empfunden vieler Heil’gen starker Muth, da sie haben überwunden fröhlich durch des Lammes Blut: sollten wir denn allhier nicht auch streiten mit Begier?

6. Wer die Sklaverei nur liebet in der Zeit und Ewigkeit, und den Sünden sich ergibet, der hat wenig Lust zum Streit; denn die Nacht, Satans Macht, hat ihn in den Schlaf gebracht.

7. Aber wen die Weisheit lehret, was die Freiheit für ein Theil, dessen Herz zu Gott sich kehret, seinem allerhöchsten Heil, sucht allein ohne Schein Christi freier Knecht zu sein.

8. Denn vergnügt auch wohl das Leben, so der Freiheit mangeln muß? Wer sich Gott nicht ganz ergeben, hat nur Müh’, Angst und Verdruß. Der, der kriegt recht vergnügt, wer sein Leben selbst besiegt.

9. D’rum auf! laßt uns überwinden in dem Blute Jesu Christ’, und an unsre Stirne binden sein Wort, so ein Zeugniß ist, das uns deckt und erweckt und nach Gottes Liebe schmeckt.

10. Unser Leben sei verborgen mit Christo in Gott allein, auf daß wir an jenem Morgen mit ihm offenbar auch sei’n, da das Leid dieser Zeit werden wird zu lauter Freud’,

11. Da Gott seinen treuen Knechten geben wird den Gnadenlohn, und die Hütten der Gerechten stimmen an den Siegesthon; da fürwahr Gottes Schaar ihn wird loben immerdar.

Justus Falckner.

In eigener Melodie.

57. Fahre fort;:; Zion, fahre fort im Licht! Mache deinen Leuchter helle, laß die erste Liebe nicht, suche stets die Lebensquelle! Zion, dringe durch die enge Pfort’, fahre fort ;:!

2. Leide dich ;:; Zion, leide ohne Scheu Trübsal, Angst mit Spott und Hohne, sei bis in den Tod getreu, siehe auf die Lebenskrone. Zion, fühlest du der Schlangen Stich, leide dich ;:!

3. Folge nicht ;:; Zion, folge nicht der Welt, die dich suchet groß zu machen, achte nichts ihr Gut und Geld, nimm nicht an den Stuhl des Drachen! Zion, wenn sie dir viel Lust verspricht, folge nicht ;:!

4. Prüfe recht ;:; Zion, prüfe recht den Geist, der dir ruft zu beidern Seiten; thue nicht, was er dich heißt, ...

\quellseite{517}

Laß nur deinen Stern dich leiten! Zion, beide, das, was fromm und schlecht, \footnote{Schlicht, gerade.} prüfe recht :;

5. Dringe ein :; Zion, dringe ein in Gott! Stärke dich mit Geist und Leben, sei nicht, wie die Andern, todt, sei du gleich den grünen Reben! Zion, in die Kraft für Heuchelschein dringe ein :;!

6. Brich herfür :; Zion, brich herfür in Kraft, weil die Bruderliebe brennet; zeige, was der in dir schafft, der als seine Braut dich kennet! Zion, durch die dir gegeb'ne Thür brich herfür :;!

7. Halte aus :; Zion, halte deine Treu’, laß dich ja nicht laulich finden. Auf, daß Kleinod rück herbei! Auf, verlasse, was dahinten! Zion, in dem letzten Kampf und Strauß halte aus :;!

Joh. Euseb. Schmidt.

Mel. Zion klagt mit Angst und Schmerzen.

58. Halte, was du hast empfangen, mein so theu’r erkaufter Christ, da viel Geister ausgegangen, die durch ihre schnöde List schändlich alle wollen gern dir den hellen Morgenstern nehmen oder dunkel machen. Ach, es ist hie Zeit zu wachen.

2. Ach, es ist hie Zeit zu wachen, weil es Seelenheil betrifft, wo die Schlang’ aus ihrem Rachen speiet solchen Otterngift. D’rum so sei doch nicht so blind, schlag’ es ja nicht in den Wind, als wär’ nicht viel d’ran gelegen, ob du gehst auf Irrwegen.

3. Ob du gehst auf Irrwegen oder seist auf rechter Bahn, pflegst du wohl zu überlegen. Du stehst bei dir selbst oft an, wenn dein Gang nur in der Welt ist aufs Zeitlich’ angestellt. Warum wollt’st du hier erhalten, recht den Himmelsweg zu halten?

4. Recht den Himmelsweg zu halten, fordert ja Vorsichtigkeit, da die Lehre so zerspalten mitten in der Christenheit. Der weiß hie, der dort hinaus nach des lieben Vaters Haus. Ach, da Alles so zerrissen, wer mag hie den Weg recht wissen?

5. Wer mag hie den Weg recht wissen, den man soll zum Himmel geh’n? setzet oft ein zart Gewissen. Sieh’, be’m Worte ist gut steh’n. Du bist ja so aber nicht; Gold erkennst du am Gewicht. Falsche Waar’ von guter Seiden weißt du ja zu unterscheiden.

6. Weißt du ja zu unterscheiden reinen Weizen von der Spreu; was das Beste unter beiden, was sei alt und was sei neu; oder wo dem schönen Wein Hefen untermengt sein: alles dies kannst du bald wissen, weil du dich darauf befliessen.

7. Weil du dich darauf befliessen, daß doch eitle Dinge sein, so laß hier auch dein Gewissen gehen in den Grund hinein. Brauche gleiche Sorg’ und Fleiß auch auf deiner Himmelsreis’. Sieh’ dich vor vor bösen Haufen, die dir Spreu für Korn verkaufen.

8. Die dir Spreu für Korn verkaufen, sind gestreuet durch die Welt. Gott läßt sie auch an dich laufen, wenn er auf die Probe stellt deinen Glauben, Lieb’ und Treu’, ob du auch ohn’ alle Scheu werdest bei dem Worte
\quellseite{518}
bleiben, wenn dich solche Schwätzer treiben.

9. Wenn dich solche Schwätzer treiben, treibt dich Gott zur Schrift hinein. Was dem Eisen nützt das Reiben auf dem harten Wetzestein, was das Sieb beim Weizen thut, das ihn nur macht rein und gut: so kann Gott bei falschen Lehren Licht, Verstand und Glauben mehren.

10. Licht, Verstand und Glauben mehren müssen auch durch ihren Wind alle, die die Schrift verehren, oft bei einem Gotteskind. Wie der Wölfe Rauben macht, daß ein Hirte treulich wacht, und der Feind mit seinem Schmeißen freudig heißt das Schwert ergreifen.

11. Freudig heißt das Schwert ergreifen gleichfalls auch der Ketzer Schwarm, die mit ihren Lügenpfeifen blasen allzu oft Alarm. Aber ein geübter Sinn wirft als eitel Spreu dahin und schlägt aus der Schrift darnieder, was läuft Gottes Wort zuwider.

12. Was läuft Gottes Wort zuwider, was im Tod’ hält seinen Stich, woraus wahre Jesusglieder auch nicht können trösten sich, was verlieret seinen Schein, wenn es auf dem Prüfestein nach der Schrift wird untersuchet, solche Lehre bleibt verfluchet.

13. Solche Lehre bleibt verfluchet, käm’ sie auch vom Engel her. Wer sie nun zu pflanzen suchet, Jesu, dem steur’ selbst und wehr’. Bring’ zurecht, du treuer Hirt, alle Seelen, die verirrt. Pflanze, stärke, gründ’ und mehre deine Kirch’ und reine Lehre.

Franz Vogt.

Mel. Freu’ dich sehr, o meine Seele.

59. Herr, auf dich will ich fest hoffen, Gut’s zu thun, verleihe mir, deine Wohnung sei mir offen, redlich da zu dienen dir. An dir laß mich haben Lust, Weltlust sei mir unbewußt. Was mein Herz ann noch verlanget, gib mir, weil es an dir hanget.

2. Dir befehl’ ich meine Wege, auf dich hoff’ ich, Herr, mein Gott, zeige mir die rechten Stege, laß mich nirgends in der Noth. Alles wirst du machen gut, stärke nur den schwachen Muth, daß ich dir sei ganz ergeben, und in dir recht möge leben.

3. Herr, nach deinem Wohlgefallen lehre mich thun immerdar; denn du bist mein Gott in allen, dein Geist meiner nehme wahr, führe mich auf eb’ner Bahn, und bring’ mich gen Himmel an. Ach, Herr, laß mich nimmer irren, noch des Satans List verwirren.

4. Lehre mich, Herr, wohl bedenken, daß es mit mir haben muß bald ein Ende, da sich’s lenken muß zum Ziel, und da mein Fuß muß dann aus dieser Zeit hin zur lautern Ewigkeit. Ach! laß mich, Herr, um dich schweben, ewig wohl bei dir zu leben.

Martin Geier.

In eigener Melodie.

60. Höchster Priester, der du dich selbst geopfert hast für mich, laß doch, bitt’ ich, noch auf Erden auch mein Herz dein Opfer werden.

2. Denn die Liebe nimmt nichts an, was du, Liebe, nicht gethan. Was durch deine Hand nicht gehet, wird zu Gott auch nicht erhöhet.

\quellseite{519}

3. D’rum so tödt’ und schlachte hin meinen Willen, meinen Sinn. Reiß mein Herz aus meinem Herzen, sollt’s auch sein mit tausend Schmerzen.

4. Trage Holz’ auf dein’n Altar und verbrenn’ mich ganz und gar. O du tausendliebste Liebe, wenn doch nichts von mir mehr bliebe!

5. Also wird es wohl gescheh’n, daß der Herr es an wird sehn. Also werd’ ich noch auf Erden Gott ein liebes Opfer werden.

Joh. Angelus.

Mel. O daß ich tausend Zungen hätte.

61. Ich will dich lieben, meine Stärke, ich will dich lieben, meine Zier. Ich will dich lieben mit dem Werke und immerwährender Begier. Ich will dich lieben, schönstes Licht, bis mir das Herze bricht.

2. Ich will dich lieben, o mein Leben, als meinen allerbesten Freund. Ich will dich lieben und erheben, so lange mich dein Glanz bescheint. Ich will dich lieben, Gotteslamm, als meinen Bräutigam.

3. Ach, daß ich dich so spät erkennet, du hochgelobte Schönheit du, und dich nicht eher mein genennet, du höchstes Gut und wahre Ruh’. Es ist mir leid und bin betrübt, daß ich so spät geliebt.

4. Ich lief verirrt und war verblendet, ich suchte dich und fand dich nicht. Ich hatte mich von dir gewendet und liebte das geschaff’ne Licht. Nun aber ist’s durch dich gescheh’n, daß ich dich hab’ ersehn.

5. Ich danke dir, du wahre Sonne, daß mir dein Glanz hat Licht gebracht.

Ich danke dir, du Himmelswonne, daß du mich froh und frei gemacht. Ich danke dir, du güld’ner Mund, daß du mich machst gesund.

6. Erhalte mich auf deinen Stegen und laß mich nicht mehr irre geh’n. Laß meinen Fuß in deinen Wegen nicht straucheln oder stille steh’n. Erleucht’ mir Leib und Seele ganz, du starker Himmelsglanz.

7. Gib meinen Augen süße Thränen, gib meinem Herzen keusche Brunst. Laß meine Seele sich gewöhnen, zu üben in der Liebe Kunst. Laß meinen Sinn, Geist und Verstand stets sein zu dir gewandt.

8. Ich will dich lieben, meine Krone, ich will dich lieben, meinen Gott. Ich will dich lieben, ohne Lohne auch in der allergrößten Noth. Ich will dich lieben, schönstes Licht, bis mir das Herze bricht.

Joh. Angelus.

In eigener Melodie.

62. Liebe, die du mich zum Bilde deiner Gottheit hast gemacht; Liebe, die du mich so milde nach dem Fall hast wiederbracht; Liebe, dir ergeb’ ich mich, dein zu bleiben ewiglich.

2. Liebe, die du mich erfreun, eh’ als ich erschaffen war; Liebe, die du Mensch geboren und mir gleich wardst ganz und gar; Liebe, dir ergeb’ ich mich, dein zu bleiben ewiglich.

3. Liebe, die für mich gelitten und gestorben in der Zeit; Liebe, die mir hat erstritten ew’ge Lust und Seligkeit; Liebe, dir ergeb’ ich mich, dein zu bleiben ewiglich.

\quellseite{520}

4. Liebe, die du Kraft und Leben, Licht und Wahrheit, Geist und Wort; Liebe, die sich bloß ergeben mir zum Heil und Seelenhort; Liebe, dir ergeb’ ich mich, dein zu bleiben ewiglich.

5. Liebe, die mich hat gebunden an ihr Joch mit Leib und Sinn; Liebe, die mich überwunden und mein Herze hat dahin; Liebe, dir ergeb’ ich mich, dein zu bleiben ewiglich.

6. Liebe, die mich ewig liebet, die für meine Seele bitt’t; Liebe, die das Lösegeld gibt und mich kräftiglich vertritt; Liebe, dir ergeb’ ich mich, dein zu bleiben ewiglich.

7. Liebe, die mich wird erwecken aus dem Grab’ der Sterblichkeit; Liebe, die mich wird umstecken mit dem Laub’ der Herrlichkeit; Liebe, dir er- geb’ ich mich, dein zu bleiben ewiglich.

Joh. Angelus.

Mel. O Gott, du frommer.

\subsection{63. Welch’ eine Sorg’ und Furcht}

soll nicht bei Christen wachen und sie behutsamlich und wohlbedächtig machen! Mit Furcht und Zittern, heißt’s, schafft eurer Seelen Heil, wenn kaum der Fromme bleibt, wie denn der sünd’ge Theil?

2. Der Satan geht umher und sucht zu verschlingen, legt tausend Netz’ und Strick’ in unvermerkten Dingen. Die Welt ist voll verführt und bind’t ihr selbst die Ruth’; ja gar der ärgste Feind ist unser Fleisch und Blut.

3. Man kann so manche Sünd’ un- wissentlich begehen; vor Gott kommt der Gedank’ gleich einem Werk zu stehen; ein einzig räudig Schaf ver- derbt den ganzen Stall. Wer steht, der sehe zu, daß er nicht plötzlich fall’.

4. Ihr sollet, saget Gott, wie ich bin, heilig leben, mir eure Seele ganz, nicht halb getheilet, geben; vom Adel und Gewalt hab’ ich nicht viel’ er- wählt, viel’ Herr-Herr-Sager sind der Höllen zugezählt.

5. Und wird ein Frommer schlimm, so soll ihm das nichts dienen, daß er vor solcher Zeit rechtschaffen gut er- schienen; der Knecht, der es nicht thut, den Willen aber weiß, macht sich viel größ’re Schlag’ durch seiner Bosheit Fleiß.

6. Ein ein’ger Apfelbiß konnt’ uns so heftig schaden, und Christus mußt’ sich selbst in Gottes Zorne baden; ein Bruch an deiner Tauf’ spricht dir den Meineid zu; vielleicht kömmt Tod und End’ in diesem Blick und Nu.

7. Man hat genug zu thun, die ein’ge Seel’ zu retten; wer noch viel anders hat, wie kann er das vertreten? Je größer Amt und Gut und Pfund und Gaben sein, je größer’ Rechen- schaft bild’ man sich fädelich ein.

8. Die ganze erste Welt mußt jäm- merlich vergehen, acht Seelen sind allein vor Gott gerecht ersehen; nicht zehn Gerechte sind in Sodoms Nach- barschaft; des Samens vierter Theil geht nur in Frucht und Kraft.

9. Es sind nur Etliche in Canaan gegangen, auch aus den Zwölfen selbst ist Judas aufgehangen, der Zehnte danket nur, daß er vom Aus- satz rein; ach! möchten Fünf doch klug, und Fünf nur thöricht sein!

\quellseite{521}

10. Des Richters Zukunft wird gleich einem Blitz geschehen; was unrein und gemein, darf nicht in Himmel gehen; es fühlt die letzte Plag’ Egyptens erstes Kind; es wird kein Haus fast sein, da nicht Verdammte sind.

11. Herr Gott! so oft ich dies nach aller Läng’ bedenke, so schneidet Angst und Furcht durch Adern und Gelenke, so schauert mir die Haut, so gellt und klingt das Ohr, so bebet Mund und Herz und hebet sich empor.

12. Ich lebe zwar getrost durch Glauben, Lieb’ und Hoffen, und weiß, daß deine Gnad’ noch allen Menschen offen, und deine Lieb’ und Treu’ mich tröstet und erhält; doch steh’ ich desto mehr für mich und alle Welt.

13. Senk’ deine wahre Furcht in aller Menschen Herzen, laß niemand mit der Buß’ und wahrem Glauben scherzen, thu’ allen Kitzel weg, Trägheit, Vermessenheit, Verstockung, Heuchelei, Bosheit, Unheiligkeit.

14. Hilf siegen über das, was du schon überwunden, zerstör’ des Teufels Reich, laß uns ihn seh’n gebunden, laß wir mit Fleiß und Sorg’ verleugnen, widersteh’n, anhalten mit Gebet, entfliehen und entgeh’n.

15. Verleih’ Geduld und Trost im Kämpfen und im Ringen, sorgfält’ge Wachsamkeit laß mit Verstand durchdringen, laß uns in heil’ger Furcht und in Bereitschaft steh’n, daß wir mit Freudigkeit vor deine Augen geh’n.

3. A. Hedinger.

\chapter{X. Gebets-, Lob- und Dank-Lieder}

Mel. Es war einmal ein reicher Mann.

64. Für allen Freuden auf Erden kann niemand fein feiner werden, denn die ich geb’ mit mein’m Singen und mit manchem süßen Klingen.

2. Sie kann nicht sein ein böser Muth, wo da singen Gesellen gut. Sie bleibt kein Zorn, Zank, Haß noch Neid, weichen muß alles Herzeleid.

3. Geiz, Sorg’ und was sonst hart anleit, fährt hin mit aller Traurigkeit. Auch ist ein jeder deß wohl frei, daß solche Freud’ kein’ Sünde sei.

4. Sondern auch Gott viel daß gefällt, denn alle Freud’ der ganzen Welt. Dem Teufel sie sein Werk zerstört und verhindert viel böser Mörd.

5. Das zeugt David des Kön’ges That, der dem Saul oft gewehret hat mit gutem süßem Harfenspiel, daß er in großem Mord nicht fiel.

6. Zum göttlichen Wort und Wahrheit macht sie das Herz still und bereit. Solch’s hat Elisäus bekannt, da er den Geist durch’s Harfen fand.

7. Die beste Zeit im Jahr ist mein, da singen alle Vögelein. Himmel und Erden ist derr’ voll, viel gut Gesang da lautet wohl.

8. Voran die liebe Nachtigall macht Alles fröhlich überall mit ihrem lieblichen Gesang. Deß muß sie haben immer Dank.

9. Viel mehr der liebe Herre
\quellseite{522}
Gott, der sie also geschaffen hat, zu sein die rechte Sängerin, der Musiken ein’ Meisterin.

10. Dem singt und springt sie Tag und Nacht, sein’s Lebens sie nichts müde macht. Den ehrt und lobt auch mein Gesang und sagt ihm ein’n ewigen Dank.

Dank für Genesung von Krankheit.

In eigener Melodie.

65. Herr, du Kraft und Arzt der Schwachen, der du in die Hölle stößt und auch wieder d’raus erlößt, der du den gesund kannst machen, der auf deine Gnade traut und auf deine Güte baut:

2. Ach, es fehlet mir an Worten, daß ich deinen Ruhm und Preis nicht g’nug auszusprechen weiß, deinen Ruhm, der aller Orten durch die weite Welt erklingt und des Todes Reich bezwingt.

3. Ach, o hätt’ ich Engelszungen, daß jetzt deiner Hülfe Macht durch derselben Zier und Pracht würde nach Gebühr besungen, weil du deine Vaterhand mir so gnädig zugewandt.

4. Das Gebet kann dich bewegen, wenn es nur von Herzen geht und auf festem Glauben steht, daß du deinen Zorn mußt legen, daß, anstatt der Strafen, du uns den Segen schickest zu.

5. Ist ein Vater gleich erbittert; wenn das Kind zu Füßen fällt und vor ihm sich kläglich stellt, wenn es weinet, seufzet, zittert: wird das Herz ihm so erweicht, daß er Hülff’ und Hand darreicht.

6. Nun, es soll mein schlechtes Singen, weil du aus der Krankheit Nacht mich mir selbst hast wiederbracht, Vater, dir zu Ehren klingen, und was du, o Schöpfer willst, werd’ hinfort an mir erfüllt.

Dank für erlangten Frieden.

Mel. Nun danket Alle Gott.

66. Herr Gott, dich loben wir, regier’, Herr, unsre Stimmen, laß deines Geistes Gluth in unsern Herzen glimmen. Komm, komm, o edle Flamm’, ach komm zu uns allhier, so singen wir mit Lust: Herr Gott, dich loben wir.

2. Herr Gott, dich loben wir, wir preisen deine Güte, wir rühmen deine Macht mit herzlichem Gemüthe. Es steiget unser Lied bis an die Himmelsthür’ und tönt mit großem Schall: Herr Gott, dich loben wir.

3. Herr Gott, dich loben wir für deine große Gnaden, daß du das Vaterland von Kriegsblast entladen, daß du uns blicken läßt des güldnen Friedens Zier, d’rum jauchzet alles Volk: Herr Gott, dich loben wir.

4. Herr Gott, dich loben wir, die wir in langen Jahren der Waffen schweres Joch und frechen Grimm erfahren. Jetzt rühmet unser Mund mit herzlicher Begier: Gott Lob, wir sind in Ruh’! Herr Gott, wir danken dir.

5. Herr Gott, dich loben wir, daß du die Pfeil’ und Bogen, Schilde, Bogen, Spieß und Schwert zerbrochen und zerschlagen. Der Strick
\quellseite{523}
ist nun entzwei, darum so singen wir mit Herzen, Zung’ und Mund: Herr Gott, wir danken dir.

6. Herr Gott, dich loben wir, daß du uns zwar gestrafet, jedoch in dei- nem Zorn nicht gar hast weggeraffet. Es hat die Vaterhand uns deine Gnadenthür jetzt wieder aufgethan. Herr Gott, wir danken dir.

7. Herr Gott, wir danken dir, daß du Land, Kirch’ und Häuser, den edlen Rautenstamm und dessen grüne Reiser bisher erhalten hast. Gib ferner Gnad’ allhier, daß auch die Nachwelt sing’: Herr Gott, wir danken dir.

8. Herr Gott, wir danken dir, und bitten, du woll’st geben, daß wir auch künftig stets in stolzer Ruhe leben. Krön’ uns mit deinem Gut, erfüll’ jetzt nach Begier, o Vater, unsern Wunsch. Herr Gott, wir danken dir.

9. Herr Gott, wir danken dir mit Orgeln und Trompeten; mit Harfen und Pandor, Posaunen, Geigen, Flö- ten, und was nur Odem hat, ertön’ jetzt für und für: Herr Gott, dich loben wir, Herr Gott, wir danken dir!

Joh. Franck.

Mel. Freu’ dich sehr, o meine Seele.

67. Hilf uns, Herr! in allen Dingen, daß wir unser Amt und Werk wohl anfangen und vollbringen; gib uns Weisheit, Kraft und Stärke! Ohne deine Hülff und Gunst ist all’ unser Werk umsonst! Hilf uns, Herr! in allen Dingen, und laß Alles wohl gelingen.

2. Hilf uns, Herr, an allen Or- ten, wo wir dein bedürftig sein, brich der Höllen Macht und Pforten und gib deinem Häufelein, deiner armen Christenheit, Liebe, Fried’ und Einig- keit! Hilf uns, Herr! \&c.

3. Hilf uns, Herr, in allen Nö- then, aller Trübsal und Gefahr; laß uns keine Drangsal tödten, nimm doch unser also wahr, daß Kreuz, Elend und Verdruß uns zum Besten dienen muß! Hilf uns, Herr! \&c.

4. Hilf uns, Herr! aus allem Lei- den in der letzten Todesnoth; laß uns fahren hin mit Freuden, und durch deinen bittern Tod kommen in das Paradeis, uns zur Freude, dir zum Preis! Hilf uns, Herr! \&c.

5. Hilf uns, Herr! in letzten Zü- gen, hilf nach unsrer Zuversicht, laß uns ritterlich obsiegen, und zu Schan- den werden nicht. O Herr Jesu! deine Hand leist’ uns Hülfe und Bei- stand, daß wir, nach vollbrachtem Ringen, dir Lob und Dank ewig singen.

Martin Rindart.

In eigener Melodie.

68. Nun preiset Alle Gottes Barm- herzigkeit, lob’ ihn mit Schalle, wertheste Christenheit. Er läßt dich freundlich zu sich laden, freue dich, Israel, seiner Gnaden.

2. Der Herr regieret über die ganze Welt, was sich nur rühret, ihme zu Füßen fällt. Viel tausend Engel um ihn schweben, Psalter und Harfe ihm Ehre geben.

3. Wohlauf, ihr Heiden, lasset das Trauren sein, zur grünen Weiden stellet euch willig ein. Da läßt er uns sein Wort verkünden, machet uns ledig von allen Sünden.

\quellseite{524}

4. Er gibet Speise reichlich und überall, nach Vaters Weise sättigt er All’ zumal. Er schaffet früh und späten Regen, füllet uns Alle mit seinem Segen.

5. D’rum preist’ und ehre seine Barmherzigkeit, sein Lob vermehre, wertheste Christenheit. Uns soll hinfort kein Unfall schaden, freue dich, Israel, seiner Gnaden.

M. Apelles von Löwenstern.

\subsection{Bitte um gute Witterung}

Mel. Wo Gott zum Haus.

69. Gott Vater, der du deine Sonn’ läßt scheinen über Böh’ und Fromm’ und der ganzen Welt damit leucht’st, mit Regen und Thau die Erde befeucht’st.

2. Die Berg’ machst du von oben naß, und läßt d’rauf wachsen Laub und Gras. In Gäng’ und Fels gut Erz du legst, Fried’, Schutz und Recht du selber hegst.

3. Du gibst auch reichlich Brod und Wein, daß Menschenherz kann fröhlich sein; du deckst auch unsre Sünden zu, dein Wort bringt uns Trost, Fried’ und Ruh’.

4. So bitten wir nun dein’ Gnad’ und Güt’, im Wort und Fried’ uns stets behüt’; die Frücht’ der Erden uns bewahr’, und gib uns heur’ ein reiches Jahr.

5. Ein fruchtbar Wetter uns bescher’, dem Hag’l und Ungewitter wehr’; Schnee, Regen, Wind und Sonnenschein all’zeit dein’m Wort gehorsam sein.

6. Heuschreck’n und Raupen sind dein’ Ruth’, und All’s, was Schad’n an Früchten thut, solch Ungeziefer, Herr, vertreib’, daß dein’ Gab’ unbeschädigt bleib’.

7. Denk’, daß wir arme Würmelein dein’ G’schöpf’, Erbaut und Kinder sein, und warten auf dein’ milde Hand, uns aus dein’m Wort und Werk bekannt.

\subsection{Um Sonnenschein}

8. Die liebe Sonn’ uns scheinen laß, heiß’ wachsen Erz, Brod, Wein und Gras, daß Leut’ und Vieh ihr Nahrung hab’n und dich erkenn’n aus deinen Gab’n.

\subsection{Um Regen}

9. Durch Christ, dein’n Sohn, hör’ unser’ Bitt’, theil’ uns ein’n gnäd’gen Regen mit, und frön’ das Jahr mit deiner Hand, mit dein’ Fußstapfen düng’ das Land.

10. Den Herrn von Zion man dich nennt, in aller Welt dein’ Güt’ man kennt, hörst unser’ Bitt’ und hilfst allein; gib Gnad’, daß wir dir dankbar sei’n.

M. Hermann.

\subsection{Dank für gute Ernte}

Mel. Nun lob’ mein’ Seel’ den x.

70. Nun lasset uns Gott danken für seine Gnade und Wohlthat; nun lasset uns bedenken, was er uns ist erzeiget hat. Die Früchte hat er geben gar reichlich auf dem Land’ und auch durch seinen Segen geliefert in die Hand, daß Alles wohl gerathen, was man gesäet aus, und durch seine Gutthaten auch wohl gebracht zu Haus.

2. Wer soll’ das nicht erkennen für Gottes Macht und weisen Rath?

\quellseite{525}

Wer sollt’ es nicht bekennen für eine göttliche Wohlthat? Es ist nichts, der da säet, nichts, der da schneidet ein, nichts, der damit umgehet, nichts, der da führet heim, Gott muß mit seinem Segen auch selber dabei sein; an dem ist es gelegen, wenn’s wohl soll kommen ein.

3. Und wenn er denn uns geben die Frucht in’s Feld und in die Stadt durch seinen Vatersegen, daß man Alles wohl einbracht hat, als soll man ihn ja preisen und rühmen die Wohlthat, ihm Ehr’ allezeit erweisen für solche große Gnad’, von Herzensgrund ihm danken und loben mit dem Mund’, in Werken auch nicht wanken von ihm zu keiner Stund’.

4. Daß solches nun geschehe, so hilf du uns, getreuer Gott! Laß’ uns auf deinem Wege, seg’n’ uns und unser täglich Brod; laß uns es auch genießen christlich in unserm Haus, daß unsre Bächlein fließen zu den Armen heraus, und wir nicht nur hier zeitlich an Früchten werden reich, sondern auch dorten selig erben das Himmelreich.

H. Erem. Weber.

\subsection{Schließlich}

Mel. Nun laßt uns Gott \&c.

71. Amen! Gott Vat’r und Sohne sei Lob in’s Himmels Throne; sein Geist stärk’ uns im Glauben und mach’ uns selig, Amen.

2. Amen! es wird geschehen, wir werden Christum sehen in den Wolken herkommen, uns mitzunehmen, Amen.

3. Amen! und ewig währe die Freude, Gott die Ehre, bringt alle Sprach’n zusammen in einem Glauben, Amen.

4. Amen! kein Tod soll schrecken, Christus will uns erwecken, der selbst zuvor begraben, und lebet ewig, Amen.

5. Amen! Gott sei gepriesen, der Geist auf Christum weiset, der helf’ uns All’n zusammen in’s ew’ge Leben, Amen.

Z. Helmhold.

\chapter{XI. Von der Kirche}

In eigener Melodie.

72. Auf, ihr Heiden, lobet Gott, der euch aus des Todes Schatten, so euch in der Finsterniß vormals ganz umgeben hatten, zu dem wahren Licht gebracht, und euch Arme nicht veracht’t! Lobet solche große Güte stets mit dankbarem Gemüthe.

2. Alle Völker, preiset Gott, der aus unverdienter Gnade hat durch seinen Sohn gemacht, daß der alte Fall nicht schade, der euch hilft durch Christi Tod aus der Sünd’ und aller Noth! Preiset solche große Güte stets mit dankbarem Gemüthe.

3. Gnad’ und Wahrheit, die da ist durch den Herren Christum kommen, so wir auch von seiner Füll’ durch den Glauben angenommen, walte mächtig allezeit über uns in Ewigkeit! Rühmet solche große Güte stets mit dankbarem Gemüthe.

Aus Rauer’s Sammlung.

\quellseite{526}

Mel. O Ewigkeit, du Donnerwort.

\subsection{73}

Die Macht der Wahrheit bricht herfür und klopft an vieler Herzen Thür’, daß man sie woll’ einlassen; des Reichs Erkenntniß gehet auf und führet ihren schnellen Lauf durch aller Völker Straßen; ja das Erkenntniß sammt dem Reich eröffnet sich nun allzugleich.

2. Seht, was der HErr für Wunder thut! Er giebt seinen Boten Muth, und Kraft den Wahrheitszeugen, die er aussendet spät und früh; was sie gehört, posaunen sie, und können’s nicht verschweigen; sie breiten aus zu seinem Ruhm das ew’ge Evangelium.

3. Ich lasse, spricht das A und O, vom ew’gen Evangelio die Botschaft nun ertönen; ich laß ausrufen Fried’ und Freud’, wie alles Volk soll sein bereit, mit mir sich zu versöhnen. Mein Engel, der die Botschaft bringt, sich mitten durch den Himmel schwingt.

4. Er wird gesehen überall, man höret seinen Friedenshall; den Völ- lern, Sprachen, Zungen wird meine Gnade weit und breit, der Reichthum meiner Gütigkeit, verkündigt und ge- sungen; man rufet aus, bald hier, bald dar, das große Hall- und Jubel- jahr.

5. Es hören die Geheimnisse die Würd’ge und Unwürdige, die hier auf Erden leben; sie hören diese Freuden- post, dies große Zeugniß nun getrost, und können sich ergeben dem großen König aller Welt, der Liebe, Bund und Gnade hält.

6. Ich habe die Unwissenheit bis- her in der vergang’nen Zeit vielfältig übersehen; nunmehro, da aufgeht das Licht, werd’ ich’s so übersehen nicht, wie bis daher geschehen. Die Wahr- heit will, wie sich’s gebühret, mit Dank sein zu Gemüth geführet.

7. Ich, ich will machen Alles neu; die Kreatur soll werden frei vom Dienst der Eitelkeiten. Wohl dem, der dieses fassen kann, und nimmt in meinem Sohn mich an zu diesen fro- hen Zeiten! Wie selig ist, der diesen Tag des Heils anjezt erkennen mag!

Aus Raumer’s Sammlung.

Mel. Freu’ dich sehr, o meine Seele.

\subsection{74}

Freuet euch, ihr Gotteskinder, preiset mit mir Gottes Macht; freuet euch, ihr armen Sünder, laßt der Welt den Stolz und Pracht: freu’t euch, daß euch Gott regiert und zu seiner Kirche führt, daß wir armen Heiden können uns nun Gottes Kin- der nennen.

2. Seine Herrschaft sich erstrecket über Alles, was da lebt; seine Hand schützt und bedecket, was sich regt, lebt und schwebt. Wenn er spricht, so geht sein Wort schnell wie Blitz und Donner fort; das durchdringet Geist und Leben, das kann Trost und Leben geben.

3. Nun, ihr Sünder hier auf Er- den, wollt ihr’s endlich haben gut, wollt ihr Gottes Kinder werden, ei, so ändert Sinn und Muth; preiset Gottes Güt’ und Macht er wird euch auch geben Kraft, daß ihr zu ihm könnet kommen und euch freuen mit den Frommen.

\quellseite{527}

4. Aller Heiden Trost ist kommen, JEsus, unser Freudenlicht, der Beschirmer aller Frommen, unser Schutz und Zuversicht; preis’t seine Güt’ und Treu’, die uns doch wird immer neu; so wird er euch Allen geben, hier und dort mit ihm zu leben.

Joh. Olearius.

paragraph

Mel. Erhalt’ uns, HErr, bei deinem Wort.

heading

75. Gleichwie sein Haus der Vogel baut, die Schwall’ auch auf ihr Nestlein schaut, also trag’ ich nun ganz und gar, HErr JEsu Christ, nach dein’m Altar.

2. Hilf mir, o HErr, durch’s Jammerthal! Gib deinen Dienern überall, daß sie dein’ Brünnlein leiten recht, daß wir behalt’n den Glauben schlecht. *)

3. Viel besser ist, in deiner Gemein’ der Allergeringst’ und -wenigst’ sein, denn sonst der Herrlichst’ in der Welt, voll Reichthum, Glück, Ehr’, Gut und Geld.

4. Du bist der HErr, Licht, Sonn’ und Schild, gibst Gnad’ und Ehr’, bist gut und mild. Hilf, HErr, den Glauben in uns mehr’ und unser Herz zu dir bekehr’!

Kü. Selnecker.

Mel. Von Gott will ich nicht lassen.

76. Groß ist an allen Enden, HErr Gott, dein Gnadenreich; wohin wir uns nur wenden, da prangt die Herrschaft gleich. Dies Reich weiß gar kein Ziel, d’rum kann es auch auf Erden nicht ausgerottet werden, es hat der Völker viel.

2. Kein Reich ist je gewesen, daß ewig könnte steh’n; ein jedes, wie wir lesen, mußt’ endlich doch vergeh’n, ja gänzlich hören auf; nur Christus Reich muß bleiben, denn das kann nicht vertreiben der Zeiten Macht noch Lauf.

3. Der HErr hat den erwecket, den er versprochen hat, der gnädig zugedecket Israels Missethat. Dies ist sein Nam’ allein, daß man ihn HErr muß nennen, den die Gerechten kennen, und der soll König sein.

4. Aus Sion mußt’ er führen sein herrlich’s Regiment und königlich regieren ein Volk, das ihn nur kennt. Sein Scepter ist das Wort, das uns zum Himmel bringet, durch alle Welt sich schwinget und siegt an manchem Ort.

5. Aus Sion ist gesendet der Scepter in die Welt, der hat zu sich gewendet ein Volk, das ihm gefällt, ein Volk, das heilig lebt, ein Volk, das Christum liebet, das sich im Glauben übet und an der Wahrheit klebt.

6. O König hochgeboren, o Gott von Ewigkeit zum Regiment erkoren der Kirchen in der Zeit! O Kraft, Licht, Stärk’ und Strahl! Du gibst ja das Gedeihen, wenn deine Diener schreien: Bekehrt euch allzumal!

7. Wohl uns, die Christus nennet sein Volk so gnädiglich, sein Volk, das ihn nur kennet und selbst verleugnet sich. Wir sind ja nicht hinfort, als eh’mals, blinde Heiden, itzt kann gar nichts uns scheiden von Gott und seinem Wort.

8. O was für Gnad’ und Segen hat uns sein Reich gebracht! Wir
\quellseite{528}
geb’n auf sichern Wegen, nicht bei der finstern Nacht; wir wandeln stets im Licht, demnach wir sind gelehret und so zu Gott bekehret, daß wir leicht irren nicht.

9. Hier gilt noch Knecht, noch Freier, hier gilt noch Weib, noch Mann, denn Christus, der Getreuer, nimmt Alles gnädig an; er macht den, der ihn liebt, zu seinem Reichsgenossen, kein Mensch wird ausgeschlossen, der ihm sein Herz nur gibt.

10. Auf! auf! ihr Müßeseelen, kommt, daß wir allzugleich uns vor den König stellen in seinem Gnadenreich, und thun, was er uns lehrt; kommt, laßt uns ihm vertrauen; kommt, laßt uns freudig schauen, wie sich sein Reich gemehrt.

11. Der Herr will Bürger haben, die fromm und heilig sind, die stets im Glauben traben, der Gott sein Herz gewinnt. Er liebt nur das Geschlecht, das nach dem Himmel strebet, das hier auf Erden lebet keusch, heilig und gerecht.

12. So schwingt nun die Gemüther, ihr Christen allzumal, zu nehmen an die Güter, welch’ uns, fast ohne Zahl, der große König schenkt, ja, welch’ uns ganz entbinden vom Teufel, Tod’ und Sünden. Wohl dem, der dies bedenkt!

Dresdner Kirchen- und Hausbuch vom J. 1694.

Mel. Es woll’ uns Gott genädig.

77. Herr Jesu, gib uns Gnad’ und Stärk’, daß wir der blinden Heiden unzählig Wesen, Thun und Werk verleugnen stets und meiden; erleucht’ auch, die im Irrweg sein, bring’ sie zu deinen Heerden, führ’ sie in deinen Schaffstall ein, daß Aller Zung’ auf Erden dich, Jesum Christ, bekenne.

2. Wie du vor diesem hast dein Wort durch heil’ge fromme Leute gebreitet aus an manchem Ort, so thu’ es, Herr, auch heute; gib treue Hirten, die zu dir fortan auch Andre bringen, mit reiner Lehre für und für durch’s Herze Allen dringen, und Alle zu dir führen.

3. Hilf, daß sie fleißig allzeit ermahnen, trösten, strafen; im Glauben, Lieb’, Gottseligkeit vorgehen ihren Schafen. Laß uns ja alle Thäter sein des Wortes, nicht nur Hörer, daß überall und allgemein wir und all’ unsre Lehrer des Glaubens End’ erhalten.

Aus Raumer’s Sammlung.

Mel. O Herrre Gott, dein göttlich Wort.

78. Hilf Gott, mein Herr, wo kommt’s doch her, daß niemand hie kann dulden uns arme Leut’, die ungescheut dein’m Wort mit Glauben hulden, geduldig still in B’schwerung viel uns hüten gern für Sünden, da doch ohn’ Leid Paß, Jud’ und Heid’ ihr’n Raum und Platz stets finden.

2. Das macht, o Christ, des Teufels List und Grimm zu diesen Zeiten, der Lügen treibt und kräftig bleibt mit Mord auf allen Seiten, und hilft mit Schein den Dienern sein, schafft durch sie sein Verlangen. Allein Gott’s Wort kann seinen Ort in dieser Welt erlangen.

3. Das ist die Art dies Schifflein zart, d’rin Noe ist erhalten, welch’s
\quellseite{529}
schwebet fort, jetzt hie, jetzt dort, im Sturmwind unzerspalten. Lust, Wasser, Feu’r ganz ungeheur’ sind ihm feindlich entgegen. Doch muß Satan es schweben lan, ohn’ Gott gilt kein Vermögen.

4. Denn Jesus Christ mit drinnen ist, und ob er wohl scheint schlafen, so wacht er doch und sorget noch, und wird uns nicht verlassen. Zu seiner Zeit wird ganz bereit sein’ Hülf’ sich lassen finden, dagegen halb des Teufels G’walt mit seinem Reich verschwinden.

5. Lob’ noch so sehr, du tolles Meer, bedeck’ dies Schifflein kleine. Du alter Drach’, halt’ kein Gemach, spei’ Feu’r in Christi G’meine. Du Antichrist, brauch’ Macht und List, Trotz, Schwert und Feu’r daneben, schaff’ Gleisnerei, Haß, Krieg, Untreu’; Gott kennt die Seinen eben.

6. Dies Schifflein muß dennoch sein und bleiben unversehret, und soll’s auch dir ein Marter sein, wirst du’s doch Gott nicht wehren. Hie ist der Herr, der Wind und Meer und All’s nach seinem Willen regiert und hält, wie’s ihm gefällt; wer will sein’n Eifer stillen?

7. Auch ist nicht fern der Morgenstern, der diese Nacht wird enden."Maß ist das Land" und wohl bekannt, da wird dies Schiff hin wenden. Hilf, treuer Gott, aus aller Noth, daß wir den Port erlangen, nach welcher Zeit dein’ Christenheit so sehnlich thut verlangen.

Nathan Chyträus.

Mel. Ihr Gestirn, ihr hohen Lüfte.

79. Jauchze Tyrus, Saba lache, freue dich, du Heidenschaft! Du be- trübtes Volk erwache und empfinde neue Kraft! Denn des Herrn Glanz und Prangen ist in dir jetzt auf- gegangen.

2. Denen, die den Götzen dienen, und dem Volk der Finsterniß, ist ein großes Licht erschienen. Nun sind wir, Gott Lob! gewiß, daß uns Niemand kann verdammen; Sem und Japhet wohnt beisammen.

3. Heute geht uns armen Heiden ein Gestirn in Jacob auf; d’rum so tragen wir mit Freuden unser Opfer jetzt zu Hauf, sind ganz fertig, sammt den Weisen unser Jesulein zu preisen.

4. Nimm doch, Jesu, nimm die Gabe, nimm die schlechte Gabe hin, die ich dir bereitet habe und die ich, wie arm ich bin, hier in Japhet’s kleinen Hütten dir in deinen Schooß will schütten.

6. Ob zwar Weihrauch, Gold und Myrrhen ich dir jetzt nicht geben kann, ei, so laß es dich nicht irren, siehe mei- nen Glauben an: dieser glänzet un- verloschen. Ach, nimm an den güld’nen Groschen!

6. Weihrauch will ich dir gewähren, wenn ich bete, ruf’ und schrei’; und die herben Myrrhenzähren der Geduld und Buß’ und Reu’. Dieses ist nun das Geschenke, das ich dir zu opfern denke.

7. Jesu, laß es dir behagen, nimm doch, nimm für lieb und gut, was ich dir jetzt vorgetragen; gib mir neuen Geist und Muth, daß ich mag in allen Dingen dir ein reines Opfer bringen.

Jch. Franck.

\quellseite{530}

Mel. Ein' feste Burg ist unser Gott.

80. O Jesu, einzig wahres Haupt der heiligen Gemeine, die an dich, ihren Heiland, glaubt und nur auf dir alleine als ihrem Felsen steht, der nie untergeht, wenn gleich die ganze Welt zertrümmert und zerfällt: erhör', erhör' uns, Jesu!

2. Laß uns, dein kleines Häufelein, das sich zu dir bekennet, dir ferner anbefohlen sein, erhalt' uns ungetrennet; Wort, Tauf' und Abendmahl laß in seiner Zahl und ersten Reinigkeit bis an den Schluß der Zeit zu unserm Troste bleiben.

3. Hilf, daß wir dir zu aller Zeit mit reinem Herzen dienen; laß uns das Licht der Seligkeit, daß uns bisher geschienen, zur Buß' kräftig sein und zum hellen Schein, der unserm Glauben mehrt, der Sünden Macht zerstört und fromme Christen machet.

4. Laß uns beim Evangelio Gut, Blut und Leben wagen, mach' uns dadurch getrost und froh, das schwerste Kreuz zu tragen. Gib Beständigkeit, daß uns Lust und Leid von dir nicht scheiden mag, bis wir den Jubeltag bei dir im Himmel halten.

5. Erbarm' dich deiner Christenheit, vermehre deine Heerde; für uns, dein armes Häuflein, streit', daß es erhalten werde; den Ärgernissen wehr'; was dich haßt, bekehr'; was sich nicht beugt, zerbrich; mach' endlich seliglich an aller Noth ein Ende.

6. Ach, Jesu! ach, wir bitten dich in deinem Jesus-Namen, erhör', erhör' uns gnädiglich, sprich, Jesu, Ja und Amen. Willt du uns Jesu sein, sind wir, Jesu, dein; so halt' dein Jesuswort und laß uns hier und dort darüber jubiliren.

Joh. Menger.

In eigener Melodie.

81. Die ist mir lieb, die werthe Magd, und kann ihr'r nicht vergessen. Lob, Ehr' und Zucht von ihr man sagt, sie hat mein Herz besessen. Ich bin ihr hold, und wenn ich sollt' groß Unglück han, da liegt nicht an, sie will mich deß ergötzen mit ihrer Lieb' und Treu' an mir, die sie zu mir will setzen und thun all' mein' Begier.

2. Sie trägt von Gold so rein ein' Kron', da leuchten in zwölf Sterne. Ihr Kleid ist wie die Sonne schön, das glänzet hell und ferne, und auf dem Mon' ihr Füße ston. Sie ist die Braut dem Herrn vertraut. Ihr ist weh' und muß gebären ein schönes Kind, den edlen Sohn und aller Welt ein Herrn, dem sie ist unterthon.

3. Das thut dem alten Drachen Zorn und will das Kind verschlingen. Sein Toben ist doch ganz verlorn, es kann ihm nicht gelingen. Das Kind ist doch gen Himmel hoch genommen hin und lässet ihn auf Erden fast sehr wüthen. Die Mutter muß gar sein allein, doch will sie Gott behüten und der recht' Vater sein.

Luther.

\quellseite{531}

\chapter{XII. Berufs- und Zeilieder}

Mel. Erschienen ist der herrlich’ Tag.

82. Das walte Gott, der helfen kann, mit Gott fang’ ich die Arbeit an, mit Gott es gebe glücklich fort, d’rum ist auch dies mein erstes Wort: Das walte Gott!

2. All’ mein Beginnen, Thun und Werk erfordert von Gott Kraft und Stärk’; mein Herz zu Gott ist stets gericht’t, d’rum auch mein Mund mit Freuden spricht: Das walte Gott!

3. So Gott nicht hilft, so kann ich nichts; wo Gott nicht gibt, da gebricht’s; Gott gibt und thut mir alles Gut’s, d’rum sprech’ ich auch nun gutes Muth’s: Das walte Gott!

4. Will Gott mir etwas geben hier, so will ich dankbar sein dafür; auf sein Wort werf’ ich aus mein Netz und sag’ in meiner Arbeit stets: Das walte Gott!

5. Anfang und Mittel sammt dem End’ ich stell’ allein in Gottes Händ’; er gebe, was mir nützlich ist, d’rum sprech’ ich auch zu jeder Frist: Das walte Gott!

6. Legt Gott mir seinen Segen bei, nach seiner großen Güt’ und Treu’, so g’nüget mir zu jeder Stund’, d’rum sprech’ ich auch von Herzensgrund: Das walte Gott!

7. Trifft mich ein Unglück, unverzagt! ist doch mein Werk mit Gott gemacht; er wird mir gnädig stehen bei, d’rum dies auch meine Losung sei: Das walte Gott!

8. Er kann mich segnen früh und spat, bis all’ mein Thun ein Ende hat; er gibt und nimmt, macht, wie er’s will, d’rum sprech’ ich auch fein in der Still’: Das walte Gott!

9. Gott steht mir bei in aller Noth und gibt mir auch mein bißchen Brod; nach seinem alten Vaterbrauch thut er mir Gut’s, d’rum sprech’ ich auch: Das walte Gott!

10. Ihn’ ihn ist all’ mein Thun umsonst, nichts hilft Verstand, Wiz oder Kunst; mit Gott geht’s fort, geräth auch wohl, laß ich dann sagen glaubensvoll: Das walte Gott!

11. Theilt Gott was mit aus Gütigkeit, so acht’ ich keiner Feinde Neid; laß haffen, wer’s nicht lassen kann, so stimm’ ich doch mit Freuden an: Das walte Gott!

12. Thun’ ich denn was mit Gottes Rath, der mir beistehet früh’ und spat, dann Alles wohl gerathen muß, d’rum sprech’ ich auch nun zum Beschluß: Das walte Gott!

Treuengebräner Gesangbuch von 1722.

\subsection{Wiegenlied}

Mel. Lobet Gott, ihr Christen.

83. Du liebes Kind, wie klein du bist und alles Witzes los, weißt nicht, was gut und sündlich ist, liegst nackend, arm und bloß.

2. So bist du durch die Taufe doch Gott wieder einverleibt; der schicket’s, daß sein Engel nach dein Schutz und Hüter bleibt.

3. Es haben deine Eltern zwar in Sünden dich gezeugt; Gott aber hat zu dir sich gar mit Gnaden hingeneigt.

\quellseite{532}

4. Der ist dein Vater, du sein Kind, er liebt dich herzlich sehr; er wächet dir ab deine Sünd’ und wirft sie ganz in’s Meer.

5. Des Herren Christi rothes Blut hat dich schneeweiß gemacht; in der hochheil’gen Taufe Fluth, ist dir das Leben bracht.

6. Denn Gottes Sohn ward dir zu gut ein schwaches Kindelein, auf daß du von der Höllengluth befreiet möchtest sein.

7. D’rum hat der Satan keine Macht noch Ansprach’ mehr an dir; du bist zum Leben wieder bracht und selig für und für.

8. Deß kannst du nun in Sicherheit einschlafen und in Ruh’, denn Gott schenkt dir die Seligkeit und alles Gut’s dazu.

9. Die Kirche schleuset dich mit ein in ihr gemein Gebet; das muß bei Gott erhöret sein, weil es aus Andacht geht.

10. Du Satan, weiche von uns weit! dies Kind ist schon getauft, und schon vor langer, lieber Zeit durch Christi Blut erkauft.

11. So schlaf’, du liebes Kindelein! Der heil’gen Engel Schaar benimmt dich aller Angst und Pein und wendet die Gefahr.

12. Du frommer Gott, verlaß uns nicht, wir trauen deiner Kraft; Herr, sei du unser Heil und Licht auf unsrer Wanderschaft.

13. Laß dieses Kind zu deinem Preis erwachsen, steh’ ihm bei, daß es hiernächst auf dein Geheiß fromm und gott’fürchtig sei.

14. Bis endlich seine Seele du durch einen sanften Tod versetzest in die Himmelsruh’ und endest seine Noth.

15. Dank sei dir, süßer Jesu Christ, der du uns auch zu gut ein schwaches Kindlein worden bist. Nimm uns in deine Hut.\footnote{A. S. B., aus Crüger’s Praxis pietatis.}

\subsection{Sommerlied}

Mel. Heut singt die liebe Christenheit.

84. Geh’ aus, mein Herz, und suche Freud’ in dieser lieben Sommerzeit an deines Gottes Gaben. Schau’ an der schönen Gärten Zier und siehe, wie sie mir und dir sich ausgeschmücket haben.

2. Die Bäume stehen voller Laub, das Erdreich decket seinen Staub mit einem grünen Kleide. Narcissus und die Tulipan die ziehen sich viel schöner an als Salomonis Seide.

3. Die Lerche schwingt sich in die Luft, das Täublein fleugt aus seiner Kluft und macht sich in die Wälder. Die hochbegabte Nachtigall ergötzt und füllt mit ihrem Schall Berg’, Hügel, Thal und Felder.

4. Die Glucke führt ihr Völklein aus, der Storch baut und bewohnt sein Haus, das Schwalblein speist ihr’ Jungen. Der schnelle Hirsch, das leichte Reh ist froh und kommt aus seiner Höh’ in’s tiefe Gras gesprungen.

5. Die Bächlein rauschen in dem Sand und malen sich und ihren Rand mit schattenreichen Myrthen. Die Wiesen liegen hart dabei und klingen ganz von Lustgeschrei der Schaf und ihrer Hirten.

\quellseite{533}

6. Die unverdroß'ne Bienenschar zeugt hin und her, sucht hier und dar ihr' edle Honigspeise. Des süßen Weinstocks starker Saft kriegt täglich neue Stärk' und Kraft in seinem schwachen Reise.

7. Der Weizen wächst mit Gewalt, darüber jauchzet Jung und Alt und rühmt die große Güte deß, der so über- flüssig läßt und mit so manchem Gut begabt das menschliche Gemüthe.

8. Ich selbst kann und mag nicht ruh'n, des großen Gottes großes Thun erweckt mir alle Sinnen. Ich singe mit, wenn Alles singt und lasse, was dem Höchsten klingt, aus meinem Her- zen rinnen.

9. Ach, denk' ich, bist du hie so schön und läßt du's uns so lieblich geh'n auf dieser armen Erden, was will doch wohl nach dieser Welt dort in dem reichen Himmelszelt und güld'nen Schlosse werden!

10. Welch' hohe Luft, welch heller Schein wird wohl in Christi Garten sein! Wie muß es da wohl klingen, da so viel tausend Seraphim mit ein- gestimmtem Mund und Stimm' ihr Halleluja singen.

11. O wär' ich da, o stünd' ich schon, ach süßer Gott, vor deinem Thron und trüge meine Palmen, so wollt' ich nach der Engel Weis' erhöhen deines Namens Preis mit tausend schönen Psalmen.

12. Doch will ich gleichwohl, weil ich noch hier trage dieses Leibs Joch, auch nicht gar stille schweigen. Mein Herze soll sich fort und fort an diesem und an allem Ort zu deinem Lobe neigen.

13. Hilf nur und segne meinen Geist mit Segen, der vom Himmel fleußt, daß ich dir stetig blühe. Gib, daß der Sommer deiner Gnad' in meiner Seele früh und spat viel Glaubensfrucht' erziehe.

14. Mach' in mir deinem Geiste Raum, daß ich dir werd' ein guter Baum, und laß mich wohl bekleiden. Verleihe, daß zu deinem Ruhm ich deines Gartens schöne Blum' und Pflanze möge bleiben.

15. Erwähle mich zum Paradeis und laß mich bis zur letzten Reis' an Leib und Seele grünen. So will ich dir und deiner Ehr' allein und sonsten Keinem mehr hier und dort ewig dienen. P. Gerhardt.

\subsection{Reiselied}

Mel. Dies sind die heil'gen zehn Gebot'.

85. In Gottes Namen fahren wir, sein heil'ger Engel geh' uns für, wie dem Volk in Egyptenland, das ent- ging Pharaonis Hand. Kyrieleis.

2. Herr, du woll'st unser Geleits- mann sein und mit uns gehen aus und ein und zeigen alle Steig' und Steg'. Wehre dem Unfall auf dem Weg'. Kyrieleis.

3. So wird kein Berg noch tiefes Thal, kein Wasser uns irr'n überall. Fröhlich komm'n wir an unserm Ort, wenn du uns gnädig hilfest fort. Kyrieleis.

4. Herr Christ, du bist der rechte Weg zum Himmel und der einig' Steg. Hilf uns Pilgrim in's Vaterland, weil du dein Blut hast d'ran gewandt. Kyrieleis. Nicol. Hermann.

\quellseite{534}

\subsection{Wiegenlied}

In eigener Melodie.

86. Nun schlaf’, mein liebes Kindelein, und thu’ dein’ Aeuglein zu; denn Gott, der will dein Vater sein, d’rum schlaf’ mit guter Ruh’.

2. Dein Vater ist der liebe Gott und will’s auch ewig sein, der Leib und Seel’ dir geben hat wohl durch die Eltern dein.

3. Und da du warst in Sünd’ gebor’n, wie Menschenkinder all’, und lagst dazu in Gottes Zorn und Adams Sünd’ und Fall:

4. Da schenkt dir Gott sein’n lieben Sohn, den senkt er in den Tod. Der kam auf Erd’ vom Himmels Thron, half dir aus aller Noth.

5. Ein Kindlein klein, ward er gebor’n, am Kreuz’ sein Blut vergoß. Damit stillt er sein’s Vaters Zorn, macht dich von Sünden los.

6. Hör’, was dir Christ erworben hat mit seiner Marter groß, die heil’ge Tauf’, das selig’ Bad, aus seiner Seiten floß.

7. Darin bist du nun neu gebor’n durch Christi Wunden roth, verschlung’n ist Gottes grimmig Zorn, dein’ Schuld ist quitt und todt.

8. Mit seinem Geist’ er dich auch krönt aus lauter Lieb’ und Treu’, der in dein zartes Herzlein ströhmt und macht dich gar spannen.

9. Er send’t dir auch sein’ Engelein zu Hütern Tag und Nacht, daß sie bei deiner Wiegen sein’ und halten gute Wacht.

10. Damit der böse Geist kein’n Theil an deinem Seelchen find’. Das bringt dir alles Christus Heil, d’rum bist ein selig Kind.

11. Dem Vater und der Mutter rein befiehlt er dich mit Fleiß, daß sie dein’ treue Pfleger sein, zieh’n dich zu Gottes Preis.

12. Dazu das liebe Jesulein, das g’stellt sich zu dir fein, will dein Emmanuelchen sein und liebes Brüderlein.

13. D’rum schlaf, mein liebes Kindelein, preiß Gott, den Vater dein, wie Zacharias Hänselein, so wirst du selig sein.

14. Der heilig’ Christ, der segne dich, bewahr’ dich allezeit. Sein heil’ger Nam’ behüte dich, schütz’ dich für allem Leid.

15. Amen, amen, ja das ist wahr, das sagt der Heilige Geist. Geb’ Gott, daß du von heut’ zu Jahr ein gotts-sel’g Kindlein sei’st. Joh. Rattkefuß.

\subsection{87. Voller Wunder, voller Kunst,}

Mel. Nicht so traurig, nicht so sehr.

voller Weisheit, voller Kraft, voller Hulde, Gnad’ und Gunst, voller Labsal, Trost und Saft, voller Wunder, sag’ ich noch, ist der keuschen Liebe Joch.

2. Die sich nach dem Angesicht niemals hier bevor gekannt, auch sonst im Geringsten nicht mit Gedanken zugewandt, derer Herzen, derer Hand knüpft Gott in ein Liebesband.

3. Dieser Vater zeugt sein Kind, jener sein’s dagegen auf; beide treibt ihr sonder Wind, ihre sond’re Bahn und Lauf; aber wenn die Zeit nun war, wird’s ein wohlgerath’nes Paar.

4. Hier wächst ein geschickter Sohn, dort ein’ edle Tochter zu; Eines ist des
\quellseite{535}
Andern Kron’, Eines ist des Andern Ruh’, Eines ist des Andern Licht, wissen’s aber Beide nicht.

5. Bis so lang’ es dem beliebt, der die Welt im Schooße hält und zur rechten Stunde gibt Jedem, was ihm wohlgefällt: da erscheint im Werk und That der so tief verborg’ne Rath.

6. Da wählt Ahasverus Blick sich die stille Esther aus, den Tobias führt das Glück in der frommen Sara Haus, David’s bald gewandter Will’ holt die flug’ Abigail.

7. Jacob fleucht vor Esau’s Schwert und trifft seine Rahel an, Joseph dient auf fremder Erd’ und wird Asnath Herr und Mann, Mose spricht bei Jethro ein, da wird die Zipora sein.

8. Jeder findet, jeder nimmt, was der Höchst’ ihm auserseh’n: was im Himmel ist bestimmt, pflegt auf Erden zu gescheh’n, und was denn nun so geschicht, das ist sehr wohl ausgericht’t.

9. Oest’res denkt man: dies und dies hätte können besser sein; aber wie die Finsterniß nicht erreicht den Sonnenschein: also geht auch Menschen-sinn hinter Gottes Weisheit hin.

10. Laß zusammen, was Gott fügt; der weiß, wie’s am Besten sei; unser Denken fehlt und trügt, sein Gedank’ ist mangelfrei; Gottes Werk hat festen Fuß, wenn sonst Alles fallen muß.

11. Siehe frommen Kindern zu, die im heil’gen Stande steh’n, wie so wohl Gott ihnen thu’, wie so schön er lasse geh’n alle Thaten ihrer Händ’ auf ein gutes, sel’ges End’.

12. Ihrer Tugend werther Ruhm steht in steter voller Blüt’; wenn sonst aller Liebe Blum’ als ein Schatten sich verzieht, und wenn aufhört alle Treu’, ist doch ihre Treue neu.

13. Ihre Lieb’ ist immer frisch und verjüngt sich fort und fort; Liebe zieret ihren Tisch und verzuckert alle Wort; Liebe gibt dem Herzen Rast in der Müh’ und Sorgenlast.

14. Geht’s nicht all’zeit, wie es soll ist doch diese Liebe still, hält sich in dem Kreuze wohl, denkt, es sei des Herrn Will’, und versichert sich mit Freud’ einer künftig bessern Zeit.

15. Unterdessen geht und fleußt Gottes reicher Segenbach, speist die Leiber, tränkt den Geist, stärkt des Hauses Grund und Dach, und was klein, gering’ und bloß, macht er mächtig, viel und groß.

16. Endlich, wenn nun ganz vollbracht, was Gott hier in dieser Welt frommen Kindern zugedacht, nimmt er sie in’s Himmelszelt und drückt sie mit großer Lust selbst an seinen Mund und Brust.

17. Nun, so bleibt ja voller Gunst, voller Labsal, Trost und Saft, voller Wunder, voller Kunst, voller Weisheit, voller Kraft, voller Wunder, sag’ ich noch, bleibt der keuschen Liebe Joch.

Paul Gerhardt.

Mel. Wie schön leuchtet der Morgenstern.

88. Wie schön ist’s doch, Herr Jesu Christ, im Stande, da. dein Segen ist, im Stande heil’ger Ehe! Wie steig’t und neigt sich deine Gab’ und alles Gut’ so mild herab aus deiner heil’gen Höhe, wenn sich an dich
\quellseite{536}
fleißig halten Jung’ und Alten, die im Orden eines Lebens einig worden.

2. Wenn Mann und Weib sich wohl begehn und unverrückt beisammen steh’n im Bunde reiner Treue, da geht das Glück in vollem Lauf, da sieht man, wie der Engel Hauf’ im Himmel selbst sich freue. Kein Sturm, kein Wurm kann zerschlagen, kann zernagen, was Gott gibet dem Paar, das sich ewig liebet.

3. Vor Allen gibt er seine Gnad’, in deren Schooß er früh und spat sein’ Hochgeliebten heget. Da spannt sein Arm sich täglich aus, da faßt er uns und unser Haus, gleich als ein Vater pfleget. Da muß ein Fuß nach dem andern geh’n und wandern, bis sie kommen in das Zelt und Sitz der Frommen.

4. Der Mann wird einem Baume gleich, an Aesten schön, an Zweigen reich; das Weib gleich einem Reben, der seine Träublein trägt und nährt und sich je mehr und mehr vermehret mit Früchten, die da leben. Wohl dir, o Zier, Mannes Sonne, Hauses Wonne, Ehrenkrone, Gott denkt dein bei seinem Throne.

5. Dich, dich hat er ihm auserkorn, daß aus dir werd’ herausgeborn das Volk, das sein Reich bauet. Sein Wunderwerk geht immerfort, und seines Mundes starkes Wort macht, daß dein Auge schauet schöne Söhne und die Töchter, die den Rocken fein abspinnen und mit Kunst die Zeit gewinnen.

6. Sei gutes Muth’s; wir sind es nicht, die diesen Orden aufgericht’t, es ist ein höh’rer Vater. Der hat uns je und je geliebt und bleibt, wenn unsre Sorg’ uns trübt, der beste Freund und Rathgeber. Anfang, Ausgang aller Sachen, die zu machen wir gedenken, wird er wohl und weislich lenken.

7. Zwar bleibt’s nicht aus, es kommt ja wohl ein Stündlein, da man leidensvoll die Thränen lässet schießen. Jedennnoch wer sich in Geduld ergibt, deß Leid wird Gottes Huld in großen Freuden schließen. Sitze, schwitze nur ein wenig. Unser König wird behende machen, daß die Angst sich wende.

8. Wohl, her, mein König, nah’ herzu, gib Rath im Kreuz, in Nöthen Ruh’, in Aengsten Trost und Freude. Deß sollst du haben Ruhm und Preis; wir wollen singen bester Weis’ und danken alle beide, bis wir bei dir, deinen Willen zu erfüllen, deinen Namen ewig loben werden. Amen.

P. Gerhardt.

\chapter{XIII. Kreuz- und Trostlieder}

\subsection{Gebet um Frieden}

In eigener Melodie.

89. Du Friedefürst, HErr JESU Christ, wahr’r Mensch und wahrer Gott, ein starker Nothhelfer du bist im Leben und im Tod’. D’rum wir allein im Namen dein zu deinem Vater schreien.

2. Nicht große Noth uns stößet an von Krieg und Ungemach. Daraus
\quellseite{537}
und niemand helfen kann, denn du, d’rum führ’ die Sach’. Dein’n Vater bitt’, daß er ja nit im Zorn mit uns wollt’ fahren.

3. Gedenk’, HErr, jezund an dein Amt, daß du ein Friedfürst bist, und hilf uns gnädig allesammt jezund zu dieser Frist. Laß uns hinfort dein göttlich Wort im Fried noch länger schallen.

4. Verdienet haben wir Alles wohl und leiden’s mit Geduld. Doch deine Gnad’ größer sein soll, denn unser’ Sünd’ und Schuld. Darum vergib nach deiner Lieb’, die du fest zu uns trägst.

5. Es ist groß Elend und Gefahr, wo Pestilenz regieret. Aber viel größer ist fürwahr, wo Krieg geführet wird. Da wird veracht’t und nicht betracht’t, was recht und löblich wäre.

6. Da fragt man nicht nach Ehrbarkeit, nach Zucht und nach Gericht. Dein Wort leít\footnote{liegt.} auch zu solcher Zeit und geht im Schwange nicht. D’rum hilf uns, HErr, treib’ von uns fern Krieg und all schädlich’s Wesen.

7. Erleucht’ auch unsern Sinn und Herz durch den Geist deiner Gnad’, daß wir nicht treiben d’raus ein Scherz, der unsrer Seelen schad’. O JEsu Christ, allein du bist, der solch’s wohl kann ausrichten.

Zusatz.

8. Seid nun getrost und wohlgemuth. Wie thut ihr so verzagen? Ihr seid erkauft durch Christi Blut, sollt’s fühnlich drauf wagen. Der Engel Schutz sei euer Trutz im Leben und im Sterben.

9. Du wirst auch thun nach deinem Wort und uns verlassen nicht, weil wir auf dich, du höchster Hort, setz’n unser’ Zuversicht. D’rum hoffen wir allein zu dir. Hilf uns, o Helfer. Amen.

90. Frisch auf, mein’ Seel’, in Noth, vertrau’ allein auf Gott und laß bestellen walten; er wird dich wohl erhalten und dich aus allen Nöthen wohl wissen zu erretten.

2. Steckst du in Armuth sehr oder in Krankheit schwer, und mußt Verfolgung leiden, bedrängt auf allen Seiten, so thu’ nur fleißig beten, Gott wird dich wohl erretten.

3. Hast du viel Herzeleid in dieser bösen Zeit, und mußt dich lassen plagen, sollt du d’rum nicht verzagen; wenn du wirst fleißig beten, wird Gott dich wohl erretten.

4. Bist du der Welt ein Spott, und steckst in großer Noth, von jedermann verlassen, thu’ ein frisch Herze fassen; wenn du wirst fleißig beten, wird Gott dich wohl erretten.

5. Denn wer ein Christ will sein, der schid’ sich nur darein und sei im Kreuz geduldig, geb’ sich vor Gott nur schuldig, der ihn aus allen Nöthen gar wohl weiß zu erretten.

6. Denn es zu dieser Frist also beschaffen ist: Wenn ein Kreuz nur aufbäret, bald ein ander’s einfähret; d’rum mußt du fleißig beten zu Gott in deinen Nöthen.

\quellseite{538}

7. Je größer ja die Noth, je näher ist uns Gott, thut sich oft zu uns nähern, eh’ wir uns sein versehen, und hilft aus allen Nöthen, wenn wir nur fleißig beten.

8. Weil außerhalb der Noth unser Gebet ist todt, so muß Gott Trübsal senden, daß wir uns zu ihm wenden und all’zeit fleißig beten, daß er uns tröst’ in Nöthen.

9. Ein Christ macht’s, wie er will, so muß er leiden viel, wird unschuldig ang’geben, man thut ihm widerstreben; d’rum muß er fleißig beten, daß ihn Gott tröst’ in Nöthen.

10. Denn wer hier in Geduld befiehlt Gott sein’ Unschuld, der sie an Tag kann bringen, dem wird’s endlich gelingen, wenn er in seinen Nöthen nur stets thut fleißig beten.

11. Und weil Trübsal und Noth der Christen täglich Brod, Herzeleid in allen Ständen, wo man sich hin thut wenden, so muß man fleißig beten, daß Gott uns tröst’ in Nöthen.

12. Hat’s jemals Noth gethan, daß man Gott rufe an, so ist jetzt Noth vorhanden, Elend in allen Landen; d’rum laßt uns fleißig beten, daß uns Gott helf’ aus Nöthen.

13. Denn der Christen Gebet nimmermehr leer abgeht; der lieb’ Gott es erhöret, sie ihrer Bitt’ gewähret; wenn sie nur fleißig beten, so hilft er ihn’n aus Nöthen.

14. Ob’s gleich nicht bald geschicht, muß man ablassen nicht, sondern getrost anhalten und dem Gott lassen walten, der uns aus allen Nöthen gar wohl weiß zu erretten.

15. Denn wir durch viel Trübsal in diesem Jammerthal uns hindurch müssen dringen, und nach dem Himmel ringen, da Gott aus allen Nöthen uns endlich wird erretten.

16. So ist dieser Zeit Leid nicht werth der Herrlichkeit, die offenbar soll werden an und nach den Beschwerden; wenn wir nur hier in Nöthen zum HErren fleißig beten.

17. Darum, o frommer Christ, der Gott ergeben ist, laß dir für’m Kreuz nicht grauen, thu’ auf den HErren bauen, der dich aus allen Nöthen wird wissen zu erretten.

18. Wandle auf Gottes Weg, zum Guten sei nicht träg’, fahr’ fort, leid’ dich geduldig, ob du gleich bist unschuldig, und bet’ herzlich in Nöthen, Gott wird dich wohl erretten.

19. Und ob’s hier nicht geschicht, so thu’ verzagen nicht; es wird nach dein’m Begehren im Himmel besser werden, da uns Gott wird erretten aus allen unsern Nöthen.

20. Da denn all’ unser Leid sich kehren soll in Freud’, und wir werden zusammen preisen des HErren Namen, befreit von allen Nöthen, vor Gottes Antlitz treten.

Josua Stegmann.

In eigener Melodie.

91. Gib dich zufrieden und sei stille in dem Gott deines Lebens. In ihm ruht aller Freuden Fülle, ohn’ ihn mühst du dich vergebens. Er ist dein Quell und deine Sonne, scheint täglich hell zu deiner Wonne. Gib dich zufrieden.

\quellseite{539}

2. Er ist voll Lichtes, Trost und Gnaden, ungefärbten treuen Herzens. Wo er steht, thut dir kein'n Schaden auch die Pein des größten Schmerzens. Kreuz, Angst und Noth kann er bald wenden, ja auch den Tod hat er in Händen. Gib dich zufrieden.

3. Wie dir's und Andern oft er- gebe, ist ihm wahrlich nicht verborgen. Er sieht und kennet aus der Höhe der betrübten Herzen Sorgen. Er zählt den Lauf der heißen Thränen und faßt zu Hauf' all' unser Sehnen. Gib dich zufrieden.

4. Wenn gar kein Ein'ger mehr auf Erden, dessen Treue du darfst trauen, alsdann will er dein Treu'ster werden und zu deinem Besten schauen. Er weiß dein Leid und heimlich Grä- men, auch weiß er Zeit, dir's zu be- nehmen. Gib dich zufrieden.

5. Er hört die Seufzer deiner See- len und des Herzens stilles Klagen; und was du Keinem darfst erzählen, magst du Gott gar kühnlich sagen. Er ist nicht fern, steht in der Mitten, hört bald und gern der Armen Bitten. Gib dich zufrieden.

6. Laß dich dein Elend nicht be- zwingen, halt an Gott, so wirst du siegen. Ob alle Fluthen einher- gingen, dennoch mußt du oben liegen. Denn wenn du wirst zu hoch beschwe- ret, hat Gott, dein Fürst, dich schon erhöret. Gib dich zufrieden.

7. Was sorgst du für dein armes Leben, wie du's halten woll'st und nähren? Der dir das Leben hat ge- geben, wird auch Unterhalt bescheren. Er hat ein' Hand voll aller Gaben, da See und Land sich muß von laben. Gib dich zufrieden.

8. Der allen Vöglein in den Wäl- dern ihr bescheid'nes Körnlein weiset, der Schaf' und Rinder in den Fel- dern alle Tage tränkt und speiset, der wird ja auch dich Ein'gen füllen und deinen Bauch zur Nothdurft stillen. Gib dich zufrieden.

9. Sprich nicht: ich sehe keine Mittel, wo ich such', ist nichts zum Besten. Denn das ist Gottes Ehren- titel, helfen, wenn die Noth am größ- ten. Wenn ich und du ihn nicht mehr spüren, da schickt er zu, uns wohl zu führen. Gib dich zufrieden.

10. Bleibt gleich die Hilf' in etwas lange, wird sie dennoch endlich kom- men. Macht dir das Harren angst und bange, glaube mir, es ist dein Frommen. Was langsam schleicht, faßt man gewisser, und was verzeucht, ist desto süßer. Gib dich zufrieden.

11. Nimm nicht zu Herzen, was die Rotten deiner Feinde von dir dich- ten. Laß sie nur immer neidlich spot- ten, Gott wird's hören und recht rich- ten. Ist Gott dein Freund und dei- ner Sachen, was kann dein Feind, der Mensch, groß machen? Gib dich zufrieden.

12. Hat er doch selbst auch wohl das Seine, wenn er's sehen könnt' und wollte. Wo ist ein Glück so klar und reine, dem nicht etwas fehlen sollte? Wo ist ein Haus, das könnte sagen: ich weiß durchaus von keinen Plagen? Gib dich zufrieden.

13. Es kann und mag nicht anders werden, alle Menschen müssen leiden.

\quellseite{540}

Was webt und lebet auf der Erden, kann das Unglück nicht vermeiden. Des Kreuzes Stab schlägt unsre Lenden bis in das Grab, da wird sich’s enden. Gib dich zufrieden.

14. Es ist ein Ruhetag vorhanden, da uns unser Gott wird lösen. Er wird uns reißen aus den Banden dieses Leib’s und allem Bösen. Es wird einmal der Tod herspringen und aus der Qual uns sämtlich bringen. Gib dich zufrieden.

15. Er wird uns bringen zu den Schaaren der Erwählten und Getreuen, die hier mit Frieden abgefahren, sich auch nun in Friede freuen, da sie den Grund, der nicht kann brechen, den ew’gen Mund selbst hören sprechen: Gib dich zufrieden.

P. Gerhardt.

\subsection{Abendlied in Sterbensläuften}

Mel. Mein’ Wallfahrt ich \&c.

92. Gott, der uns diesen Tag bewacht, dem sei nun Lob gesungen ;;: daß er durch seine starke Macht uns heute beigesprungen, damit uns nicht ein schneller Tod hat plötzlich überfallen, und in so schwerer großer Noth erhöret unser Lallen.

2. Dir, großer Gott, bekennen wir die Menge unsrer Sünden ;;: wir bitten aber: laß für dir uns doch Gnade finden, und dich der großen Niederlag’, HErr Jesu, doch erbarmen, und treib’ der Seuchen schwere Plag’ von uns verlaß’nen Armen.

3. Wend’, HErr, von uns die Pestilenz, die um uns her thut schleichen ;;: laß sie doch unser Haus und Grenz’, o Jesu, nicht erreichen. Steh’ du uns bei mit deiner Gnad’, die kräftig in den Schwachen, und hilf, daß uns die Pest nicht schad’, reiß’ uns aus ihrem Rachen.

4. Du starker Gott, dem folgen muß Wind, Meer und auch die Erden ;;: wir fallen dir durch Buß’ zu Fuß und wollen frömmer werden. Sprich nur ein Wort durch deinen Mund, so heilet, was geschlagen, und bleibet frisch, was noch gesund von diesen schweren Plagen.

5. Weil uns die schwarze Nacht umgibt, so laß uns sicher schlafen ;;: wer Gott vertraut und selber liebt, dem können solche Strafen auch krümmen nicht ein einzig’s Haar, es sei denn Gottes Wille. D’rum bet’ und schick’ dich zur Gefahr, und halt’ dem HErren stille.

6. Verzage nicht in deinem Leid, Gott wird dich nicht verderben ;;: du bist des HErren allezeit im Leben und im Sterben. Sollt’ du dann sterben, sterbe hin, thu’ Gott dein’ Seel’ ergeben; ist doch der Tod uns ein Gewinn und Christus unser Leben.

7. Hat doch der Mensch kein Fried noch Ruh’, dieweil er lebt auf Erden ;;: bis daß er thut die Augen zu, da möcht’ es besser werden. Dann ist er von der harten Schlacht des bösen Feind’s entsprungen und frei von allem Kreuz gemacht, mit dem er oft gerungen.

8. Darum geb’ ich mich willig d’rein und hoff’ auf Gott, den HErren ;;: es muß einmal gestorben sein; was soll ich mich viel sperren? Wer heut’ entschläft, ist morgen schon aus diesem Weltgetümmel und prangt
\quellseite{541}
mit unverwelkter Kron’ bei seinem Gott im Himmel.

9. Dir, Jesu, sei es heimgestellt, mein Leben und mein Sterben; mach du’s mit mir, wie’s dir gefällt, nur daß ich mög’ ererben die hohe Gnad’, o Jesulein, daß ich in jenem Leben nur möge dein Thürhüter sein. Das wollest du mir geben.

10. Und sterb’ ich gleich, so sterb’ ich dir, du bist ja auch gestorben; und hast ein ewig’s Leben mir durch deinen Tod erworben. Dies ist mein Trost, darauf ich mich beherzt nun lege schlafen. Wer Buße thut und hofft auf dich, ruht unter deinen Waffen.

Marburger Gesangbuch.

Mel. Wer nur den lieben Gott läßt walten.

93. Je größer Kreuz, je näher Himmel, wer ohne Kreuz, ist ohne Gott; bei Sündenlust und Weltgetümmel vergißt man Hölle, Fluch und Tod. O selig ist der Mann geschätzt, den Gott in Kreuz und Trübsal setzt!

2. Je größer Kreuz, je besser Christe; Gott streicht uns an den Probestein; wie mancher Garten lieget wüste, wo keine Thränenregen sein! Das Gold wird auf dem Feuerheerd, ein Christ in mancher Noth bewährt.

3. Je größer Kreuz, je stär’rer Glaube; die Palme wächst bei der Last, die Süßigkeit fleußt aus der Traube, wenn du sie wohl gekeltert hast. Am Kreuze wächst uns der Muth, wie Perlen in gesalz’ner Fluth.

4. Je größer Kreuz, je größ’re Liebe; der Wind bläst nur die Flammen auf, und scheinet gleich der Himmel trübe, so lachet doch die Sonne d’rauf.

Das Kreuz vermehrt der Liebe Gluth, gleichwie das Oel im Feuer thut.

5. Je größer Kreuz, je mehr Gebete; gerieb’ne Kräuter riechen wohl; wenn um das Schiff kein Sturmwind wehte, so fragte man nicht nach dem Pol. Wo kämen Davids Psalmen her, wenn er nicht auch versucht wär’?

6. Je größer Kreuz, je mehr Verlangen; im Thale steiget man bergan; wer durch die Wüsten oft gegangen, der sehnet sich nach Kanaan. Das Täublein, find’t es hier nicht Ruh’, so fleucht es nach der Arche zu.

7. Je größer Kreuz, je lieber Sterben; man freut sich recht auf seinen Tod, denn man entgehet dem Verderben, es stirbt auf einmal alle Noth. Das Kreuze, das die Gräber ziert, bezeugt, man habe triumphirt.

8. Je größer Kreuz, je schön’re Krone, die Gottes Schatz uns beigelegt, und die einmal vor seinem Throne der Ueberwinder Scheitel trägt. Ach, dieses theu’re Kleinod macht, daß man das größte Kreuz nicht acht’t.

9. Gekreuzigter, laß mir dein Kreuze je länger und je lieber sein! Daß mich die Ungeduld nicht reize, so pflanz’ ein solches Herz mir ein, das Glaube, Liebe, Hoffnung hegt, bis dort mein Kreuz die Krone trägt.

Beyj. Schmold.

\subsection{Morgenlied in Sterbensläuften}

Mel. Ach HErr, mich armen.

94. In diesen schweren Zeiten, wer weiß, wer heute steht; der jetzt ohn’ alles Leiden aus seinem Hause geht, dem kann’s gar leicht geschehen, daß er
\quellseite{542}
kommt heim und stirbt und, wo er nicht versehen, mit Leib und Seel' verdirbt.

2. Wer selig hofft zu sterben und nicht in Ewigkeit ;; will Leib und Seel' verderben, der mache sich bereit all' Stund' und Augenblicke durch wahre Buß' zum Tod', daß, wie es Gott ihm schicke, die Seele sei aus Noth.

3. Denn, Sünder, deine Sünden, die dich und deinen Gott ;; von Freundschaft thun entbinden, ursachen dir den Tod und solche schwere Plagen, die unser Land und Stadt jetzund so hart muß tragen um ihre Missethat.

4. Ich schid' mich, abzuscheiden gar gern von dieser Welt ;; entbind' mich von dem Leiden, wann dir's, mein Gott, gefällt. Soll ich an dieser Plage, die ich sehr wohl verschuld't, beschließen meine Tage, so gib mir nur Geduld.

5. Hie schneide, brenn' und säge und mich mit Qual und Pein ;; so hart du willst, belege: laß nur das Leiden dein an mir nicht sein vergebens und nimm mich nach der Pein, du Friedefürst des Lebens, in deinen Himmel ein.

Marburger Gesangbuch.

Mel. Nun danket Alle Gott.

95. Stell' dich, Herr, wie du willst, ich fahre fort mit Schreien; du wirst in Angst und Noth mit Hülfe mich erfreuen. Du hast mir's zugesagt, d'rum muß es auch gescheh'n; ich will noch meine Lust an deiner Hülfe seh'n.

2. Laß dich, Herr Jesu Christ, durch mein Gebet bewegen, komm in mein Herz und Haus und bringe mir den Segen. All' Arbeit, Müh' und Kunst ohn' dich nichts richtet aus; wo du mit Gnaden bist, kömmt Segen in das Haus.

3. Fünf tausend Mann hast du mit wenig Brod gespeiset, die dir aus Liebe sind, mein Jesu, nachgereiset; was zag' ich denn nun viel, wenn Hülfe mir gebricht? Wer sich auf dich verläßt, o! der verdirbet nicht.

4. Ach! liebster Heiland, ach! sieh', wie mein Herz sich kränket, wenn es an meine Sünd' und deinen Tod gedenket; doch sprichst du ja zu mir: Du bist mein liebes Kind, aus Gnad' erlaß ich dir die Strafe sammt der Sünd'.

5. Sehr langsam kömmt wohl oft, Herr Jesu, deine Stunde; es scheint, als ging' nunmehr mein Hoffen gar zu Grunde; doch ist's gewiß, sie kömmt, und kömmt zu rechter Zeit, ersetzt auch den Bezug mit Fried' und Seligkeit.

Chemnitzer Gesangbuch.

\subsection{Gebet im Kriege}

Mel. Da Christus geboren war

96. Treuer Wächter Israel, deß sich freuet meine Seel', ach, du weißest alles Leid deiner armen Christenheit. O du Wächter, der du nicht schläffst noch schlummerst, zu uns richt' dein hilfreiches Angesicht.

2. Schau', wie große Noth und Qual trifft dein Häuflein überall. Täglich wird der Trübsal mehr. Hilf, ach hilf, schütz' deine Lehr'! Wir verderben, wir vergehn, nichts wir sonst vor Augen sehn, wo du nicht bei uns wirst steh'n.

3. Hoherpriester Jesu Christ, der
\quellseite{543}
du eingegangen bist in den heil’gen Ort zu Gott, durch dein Kreuz und bittern Tod uns versöhnt, mit deinem Blut ausgelöscht der Höllen Gluth, wiederbracht das höchste Gut.

4. Sitz’t auch heut’ in’s Vaters Reich, ihm an Macht und Ehren gleich, unser Mittler und Patron, seine höchste Freud’ und Kron’, den er in dem Herzen trägt, wie sich selbst zu lieben pflegt, dem er seine Bitt’ abschlägt.

5. Kläglich schreien wir zu dir, klopfen an die Gnadenthür’, wir, die du mit höchstem Ruhm dir erkauft zum Eigenthum. Deines Vaters Zorn abwend’, der wie lauter Feu’r jetzt brennt und schier alle Welt durchbrennt.

6. Zeig’ ihm seine Wunden roth, red’ von deinem Kreuz’ und Tod’, und was du mehr hast gethan, zeig’ ihm unserwegen an. Sage, daß du unser’ Schuld hast bezahlt in Geduld, uns erlanget Gnad’ und Huld.

7. Jesu, der du Jesus heißt, als ein Jesus Hilfe leist’. Hilf mit deiner starken Hand, Menschenhülf’ hat sich gewandt. Eine Mauer um uns bau’, daß dem Feinde davor grau’ und mit Zittern sie anschau’.

8. Liebster Schatz Immanuel, du Beschützer meiner Seel’, Gott mit uns in aller Noth, neben uns und in uns Gott, Gott für uns in aller Zeit, Trutz dem, der uns thut ein Leid, Gottes Straf’ ist ihm bereit.

9. Deines Vaters starker Arm, komm und unser dich erbarm’. Laß jetzt sehen deine Macht, d’rauf wir hoffen Tag und Nacht. Aller Feinde Koppel trenn’, daß dich alle Welt erkenn’, aller Herren Herren nenn’.

10. Andre trau’n auf ihre Kraft, auf ihr Glück und Ritterschaft. Deine Christen trau’n auf dich, auf dich trau’n sie festiglich. Laß sie werden nicht zu Schand’, bleib’ ihr Helfer und Beistand, sind sie dir doch all’ bekannt.

11. Gürte dein Schwert an die Seit’ als ein Held und für sie streit’, und zerschmetter’e deine Feind’, so viel ihr’r auf Erden seind. Auf die Hölle tritt du ihn’n, leg’ sie dir zum Schemel hin und brich ihren stolzen Sinn.

12. Du bist ja der Held und Mann, der den Kriegen steuern kann, der da Spieß und Schwert zerbricht, der den Bogen macht zunicht, der die Wagen gar verbrennt und der Menschen Herzen wend’t, daß der Krieg gewinnt ein End’.

13. Jesu, wahrer Friedefürst, der der Schlangen hat zerknirscht ihren Kopf durch seinen Tod, wiederbracht den Fried’ bei Gott, gib uns Frieden gnädiglich, so wird dein Volk freuen sich, dafür ewig preisen dich.

Joh. Heermann.

Mel. Wenn wir in höchsten Nöthen.

97. Wenn Menschenhülf’ scheint aus zu sein, so stellt sich Gottes Hülfe ein; wenn niemand hilft, so hilft er und macht mein Leiden nicht zu schwer.

2. Was trägt’ ich lang’ nach Menschengunst, die doch vergehet wie ein Dunst? Es ist in dieser Welt kein Freund, der stets es gut mit einem meint.

\quellseite{544}

3. Nimm deine Zuflucht nur zu Gott, der dir kann nehmen deine Noth; such’ den zum Freund’, der dir allein mit seiner Hülf’ kann nützlich sein.

4. Wenn Gott ist Freund, alsdann dein Feind ist ohne Macht und ganz verkleint; und wären noch viel’ tausend hier, so könnt’ doch keiner schaden dir.

5. Es muß doch geh’n, wie’s Gott gefällt, wenn sich gleich Alles gegenstellt. Laß Gott nur machen, wie er will, und halte seinem Willen still.

Anton Ulrich, Herzog zu Braunschweig und Lüneburg.

6. Drum strecken auch all’ Element hier wider uns aus ihre Händ’; Angst kömmt uns aus der Tief’ und See, Angst kömmt uns aus der Luft und Höh’.

7. Es ist ein’ hochbetrübte Zeit, man plagt und jagt die armen Leut’, eh’, als es Zeit, zur Grube zu, und gönnet ihnen keine Ruh’.

8. Drum trauert auch der Freudenquell, die Sonn’, und scheint uns nicht so hell; die Wolken gießen allzumal die Thränen, ohne Maaß und Zahl.

9. Ach! mein’ auch du, o Menschenkind! und trauere über deine Sünd’; halt doch mit deinen Lastern ein, und mache dich durch Buße rein.

10. Fall’ auf die Knie, fall’ in die Arm’ des HErrn, daß sich sein Herz erbarm’, und der so wohl verdienten Rach’ in Gnaden bald ein Ende mach’.

11. Er ist ja fromm und bleibet fromm, begehrt nichts mehr, als daß man komm’, und, mit geneigter Furcht und Scheu, ihn bitt’ um Gnad’ und Vatertreu.

12. Ach Vater! Vater! höre doch, und lös’ uns aus dem Sünden-Joch, und zeuch uns aus der Welt herfür, und kehr’ uns selbsten ganz zu dir.

13. Erweiche unsern harten Muth, und mach’ uns Böse fromm und gut; wen du bekehrst, der wird bekehrt, und wer dich hört, der wird erhört. Jer. 31,18.

14. Laß deine Augen freundlich sein, und nimm mit gnäd’gen Ohren ein das Angstgeschrei, das von der Erd’ aus unsern Herzen zu dir fährt.

\subsection{Gebet bei langwierigen Regen}

Mel. Wenn wir in höchsten Nöthen.

98. 1. O Herrscher in dem Himmelszelt! was ist es doch, das unser Feld, und was es uns hervorgebracht, so ungefallt und traurig macht?

2. Nichts anders, traun, als daß die Schaar der Menschen sich so ganz und gar bis in den tiefsten Grund verkehrt und täglich ihre Schuld vermehrt.

3. Die, so, als Gottes Eigenthum, stets preisen sollten Gottes Ruhm und lieben seines Wortes Kraft, sind gleich der blinden Heidenschaft.

4. D’rum wird uns auch der Himmel blind, des Firmamentes Glanz verschwind’t; wir warten, wenn der Tag anbricht, auf’s Tageslicht, und kommt doch nicht.

5. Man zankt noch immer fort und fort, es bleibet Krieg an allem Ort, in allen Winkeln Haß und Neid, in allen Ständen Streitigkeit.

\quellseite{545}

15. Reiß weg das schwarze Zorngewand, erquicke uns und unser Land, und der so schönen Früchte Kranz mit süßem warmen Sonnenglanz.

16. Verleih’ uns, bis in unsern Tod, alltäglich unser liebes Brod, und dermaleinst, nach dieser Zeit, das süße Brod der Ewigkeit. P. Gerhardt.

\chapter{XIV. Sterbe- und Ewigkeit-Lieder}

\subsection{Sterbelied für Kinder}

Mel. Herr Jesu Christ, mein’s Lebens Licht.

99. Ach, mein herzliebstes JESulein, wie gern wollt’ ich bald bei dir sein! Ist doch in diesem Jammerthal Angst, Noth und Trübsal überall.

2. W’scher mir ein selig’s Stündelein, daß ich fein sanft mög’ schlafen ein und mit den lieben Engelein anschau’n das lieblich Ang’sicht dein.

3. Wenn mir gleich Leib und Seel’ verschmacht’t, nichts mehr versteh’ oder betracht’, so bist du doch mein’s Herzens Trost, HErr JESu, der du mich erlöst.

4. Ich werd’ nicht sterben, sondern leb’n, dem HErren Christ Lob und Preis geb’n. HErr JESu Christ, in deine Händ’ b’fehl ich mein’ Seel’ am letzten End’.

5. Du bist mein’s Lebens Kraft und Sinn, Sterben in Gott ist mein Gewinn. Ich werd’ wohl fein sanft schlafen ein, beim HErren Christo ewig sein.

6. Ade, du arge, falsche Welt, das Himmelreich ist mir bestellt. Fahr’ hin, du elend Jammerthal, dein acht’ ich nichts mehr überall. Joh. Herm. Schein.

Mel. Christus, der ist mein Leben.

100. Ach, wär’ ich doch schon droben! Mein Heiland, wär’ ich da, wo dich die Schaaren loben, und säng’: Halleluja!

2. Wo wir dein Antlitz schauen, da sehn’ ich mich hinein; da will ich Hütten bauen, denn dort ist gut zu sein.

3. Da werd’ ich Alles sehen: den großen Schöpfungsrath; was durch dein Blut geschehen; und deines Geistes That.

4. Da feiern die Gerechten, die ungezählte Schaar, mit allen deinen Knechten das große Jubeljahr.

5. Mit göttlich süßen Weisen wird mein verklärter Mund dich ohne Sünde preisen, du, meines Lebens Grund.

6. Da werden meine Thränen ein Meer voll Freude sein. Ach, stille bald mein Sehnen, und hole mich hinein. E. G. Woltersdorf.

Mel. Wie wohl ist mir, o Freund der Seelen.

101. Es ist noch eine Ruh’ vorhanden; auf! müdes Herz, und werde licht. Du seufzest hier in deinen Banden, und deine Sonne scheinet nicht. Sieh’ auf das Lamm, das dich mit Freuden dort wird vor seinem Stuhle weiden. Wirf hin die Last und eil’
\quellseite{546}
herzu. Bald ist der schwere Kampf geendet, bald, bald der saure Lauf vollendet, so gehst du ein zu deiner Ruh’.

2. Die Ruhe hat Gott auserkoren, die Ruhe, die kein Ende nimmt. Es hat, da noch kein Mensch geboren, die Liebe sie uns so bestimmt. Das Lämmlein wollte darum sterben, uns diese Ruhe zu erwerben. Es ruft, es lockt weit und breit: Ihr müden Seelen und ihr Frommen, versäumet nicht, heut’ einzukommen zu meiner Ruhe Lieblichkeit.

3. So kommet denn, ihr matten Seelen, die manche Last und Bürde drückt. Eilt, eilt aus euren Kummerhöhlen, geht nicht mehr krumm und sehr gebückt. Ihr habt des Tages Last getragen, dafür läßt euch das Lämmlein sagen: Ich selbst will eure Ruh’statt sein. Ihr seid mein Volk, gezeugt von oben. Ob Sünde, Welt und Teufel toben, seid nur getrost und gehet ein.

4. Was mag wohl einen Kranken laben und einen müden Wandersmann? Wo jener nur ein Bettlein haben und sanfte darauf ruhen kann, wenn dieser sich darf niedersetzen, an einem frischen Trunk ergötzen, wie sind die beide so vergnügt! Doch dies sind kurze Ruhestunden, es ist noch eine Ruh’ erfunden, da man auf ewig stille liegt.

5. Da wird man Freudengarden bringen, denn unsre Thränen saat ist aus. O welch ein Jubel wird erklingen und süßer Ton in’s Vaters Haus. Schmerz, Seufzen, Leid, Tod und dergleichen wird müssen flieh’n und von uns weichen, wir werden auch das Lämmlein seh’n. Es wird beim Brünnlein uns erfrischen, die Thränen von den Augen wischen. Wer weiß, was sonst noch soll gescheh’n?

6. Kein Durst noch Hunger wird uns schwächen, denn die Erquickungszeit ist da. Die Sonne wird uns nicht mehr stechen, das Lamm ist seinem Volke nah’. Es will selbst über ihnen wohnen und ihre Treue wohl belohnen mit Licht und Trost, mit Ehr und Preis. Es werden die Gebeine grünen, der große Sabbath ist erschienen, da man von keiner Arbeit weiß.

7. Da ruhen wir und sind im Frieden und leben ewig sorgenlos. Ach, fasset dieses Wort, ihr Müden, legt euch dem Lamm in seinen Schooß. Ach, Flügel her! wir müssen eilen und uns nicht länger hier verweilen, dort wartet schon die frohe Schaar. Fort, fort, mein Geist, zum Jubiliren, begürtet euch zum Triumphiren! Auf, auf, es kommt das Ruhejahr!

Joh. Sigismund Kunth.

Mel. Freut’ euch, ihr Gotteskinder all’.

102. Ihr lieben Christen, freut’ euch nun, bald wird erscheinen Gottes Sohn, der unser Bruder worden ist, das ist der lieb’-Herr Jesus Christ.

2. Der jüngste Tag ist nun nicht fern, komm Jesu Christe, lieber Herr. Kein Tag vergeht, wir warten dein, und wollten gern bald bei dir sein.

3. Verrathen ist der Widerchrist, sein Heuchelei und arge List sind offenbar und gar am Tag’. Deß führt er täglich große Klag’.

\quellseite{547}

4. Du treuer Heiland, JEsu Christ, dieweil die Zeit erfüllet ist, die uns verkündet Daniel, so komm, lieber Immanuel.

5. Sanct Simeon wart’t auch auf dich und deiner Zukunft freuet sich. Er ward auch seiner Bitt’ gewährt, da er sah, was sein Herz begehrt.

6. Er sprach: Nun will ich sterben gern, weil ich geseh’n hab’ meinen Herrn. Doch soll es nicht gestorben sein, sondern im Frieden fahr’ ich fein.

7. So warten wir nun auch der Stund’ und bitten dich von Herzens=Grund, du wollest nicht ausbleiben lang’, und straf’ einmal die alte Schlang’.

8. Der alle Welt ermordet hat und kann nicht Lügens werden satt, den nimm sammt seiner Lästerschul’ und wirf sie in den feur’gen Pfuhl.

9. Dein’ lieben Kinder warten all’, wann doch einmal die Welt zerfall’, und wann des Teufels Reich vergeh’ und er in ew’gen Schanden steh’.

10. Er ist’s, der deinen Namen schänd’t und der die armen Leut’ verblend’t. Der böse Geist sucht seinen Ruhm, und hindert, daß dein Reich nicht komm’.

11. Was du befiehlst, das lästert er und tobt darüber greulich sehr. Was uns beschwert dein’ milde Hand, daß nähm’ uns gern der Höllenbrand.

12. Der Satan läßt nicht ab, zu wehr’n, daß sich so wenig Leut’ bekehr’n. Er wend’t die Leut’ von deinem Wort und richtet an Haß, Neid und Mord.

13. Der Teufel brächt’ uns gern zu Fall’ und wollt’ uns gern verschlingen all’. Er tracht’t nach Leib, Seel’, Gut und Ehr. Herr Christ, dem rothen Drachen wehr’!

14. Die Welt kann nun nicht länger steh’n, ist schwach und alt, sie muß vergeh’n. Sie tracht an allen Orten sehr und kann die Last nicht tragen mehr.

15. Die Creatur nicht länger kann der Eitelkeit sein unterthan, und wollt’ gern wieder werden frei von’s Türken Mord und Heuchelei.

16. Der Papst hat sie so hart beschwert und all’ gut’ Ordnung verkehrt. D’rum wär’ sie gern sammt uns erlöst, wir hoffen all’ auf deinen Trost.

17. Die alten Väter warten all’, wenn du erscheinst mit großem Schall, mit aller lieben Engel Schaar, d’rauf warten sie manch hundert Jahr.

18. Eil’, lieber Herr, eil’ zum Gericht, laß seh’n dein herrlich Angesicht, das Wesen der Dreifaltigkeit. Deß helf’ uns Gott in Ewigkeit!

Erasmus Alber.

Mel. O JEsu Christ, wahr’r Gottes Sohn.

103. Ich armer Sünder gar nichts bin, allein Gott’s Sohn ist mein Gewinn. Daß er Mensch worden, ist mein Trost, er hat mich durch sein Blut erlöst.

2. O Gott Vater, regier’ du mich mit deinem Geiste stetiglich. Laß deinen Sohn, mein’n Trost und Leb’n, all’zeit in meinem Herzen schweb’n.

3. Und wenn die Stund’ vorhanden ist, nimm mich zu dir, Herr JEsu Christ. Denn ich bin dein und
\quellseite{548}
du bist mein, wie gern wollt’ ich doch bei dir sein!

4. Herr Jesu Christe, hilf du mir, daß ich ein Zweiglein bleib’ in dir, und nachmals mit dir aufersteh’, zu deiner Herrlichkeit eingeh’.

5. Mit deinen Engeln in deinem Reich dich lob’ und preise ewiglich. Amen, amen, Herr Jesu Christ, allein du mein Nothhelfer bist.

Joh. Herm. Schein.

Zu eigener Melodie.

104. Mach’s mit mir, Gott, nach deiner Güte, hilf mir in meinem Leiden. Ruf’ ich dich an, versag’ mir’s nit. Wenn sich mein’ Seel’ will scheiden, so nimm sie, Herr, in deine Händ’. Ist Alles gut, wenn gut das End’.

2. Gern will ich folgen, liebster Herr, du wirst mir’s nicht verderben. Ach, du bist doch von mir nicht fern, wenn ich gleich hier muß sterben, verlassen meine liebsten Freund’, die’s mit mir herzlich gut gemeint.

3. Ruht doch der Leib sanft in der Erd’, die Seel’ zu dir sich schwinget. In deiner Hand sie unversehrt durch’n Tod in’s Leben bringet. Hier ist doch nur ein Thränenthal, Angst, Noth, Müh’, Arbeit überall.

4. Tod, Teufel, Höll’, die Welt, die Sünd’ mir können nichts mehr schaden. An dir, o Herr, ich Rettung find’, ich tröst’ mich deiner Gnaden. Dein einz’ger Sohn aus Lieb’ und Huld für mich bezahlt hat alle Schuld.

5. Was wollt’ ich denn lang’ traurig sein, weil ich so wohl bestehe, bekleid’t mit Christi Unschuld rein, wie eine Braut hergehe. Gehab’ dich wohl, du schnöde Welt, bei Gott zu leben, mir gefällt.

Joh. Herm. Schein.

6. Zu guter Nacht, ihr lieben Freund’, Gott woll’ euch wohl bewahren, die ihr mir lieb im Herzen seind, kein Böses soll euch widerfahren. Zu guter Nacht vergebt eur’ Leid, ich komm zu Gott in die ewig’ Freud’.

Mel. O Jesu Christ, wahr’r Gottes Sohn.

105. O Jesu, Gottes Lämmlein, ich leb’ od’r sterb’, so bin ich dein. Ich bitt’, laß mich mit dir zugleich ein Erbe sein in deinem Reich.

2. Denn was wär’ sonst dein’ Sterbensnoth, so viel Striemen und Wunden roth, wenn ich nicht auch der Seligkeit genießen sollt’ in Ewigkeit?

3. Warum hätt’st du dein Leib und Leb’n im Grab’ verschlossen und aufgeg’n, wenn nicht mein Tod durch deinen Tod verjagt sollt’ werd’n, du treuer Gott?

4. Darum, o Jesu, steh’ mir bei, gewissen Trost und Hilf’ verleih’. Verlaß den nicht, Herr Jesu Christ, der mit deinem Blut gewaschen ist.

5. Im Friede laß mich schlafen ein und in dir haben Ruhe sein. Dein Antlitz laß mich sehen, Herr, ein selig’s Ende mir bescher’.

6. Das bitt’ ich durch das Leiden dein, laß das mein letztes Wünschen sein. So will ich loben allezeit dich, o Herr Gott, in Ewigkeit.

Mart. Moller.

Zu eigener Melodie.

106. Wird das nicht Freude sein, nach gläubigem Vertrauen, dort selbst den Heiland schauen in unserm Fleisch
\quellseite{549}
und Bein, mit seinen holden Blicken und Worten sich erquicken? Wird das nicht Freude sein?

2. Wird das nicht Freude sein, wenn, was der Tod entnommen, uns wird entgegen kommen, und jauchzend holen ein; wenn man wird froh umfassen, was thränend man verlassen? Wird das \&c.

3. Wird das nicht Freude sein, seh’n untern Füßen liegen, womit man hier muß kriegen, Gott dienen engelrein, von Schmerzen, Leid, Verdrüssen nicht das geringste wissen? Wird das \&c.

4. Wird das nicht Freude sein, was unaussprechlich, hören, des Höchsten Lob vermehren, die Engel stimmen ein, wenn sie mit süßem Klingen ihr dreimal Heilig! singen? Wird das \&c.

5. O das wird Freude sein! Weg, Güter dieser Erden, ihr Ehren voll Beschwerden, ihr Freuden auf den Schein! Gehabt euch wohl, ihr Lieben, muß ich euch jetzt betrüben, denkt, dort wird Freude sein!

S. E. v. Schweinitz.

\chapter{Alphabetisches Liederverzeichniß}
\begin{originalinhalt}

\quellseite{550}

\inhaltsgruppe{A.}

\inhaltseintrag{Ach, mein herzliebstes}{68}

\inhaltseintrag{Ach, wär’ ich doch schon droben}{68}

\inhaltseintrag{Ach, wie will es endlich}{29}

\inhaltseintrag{Allein auf Gottes Wort}{22}

\inhaltseintrag{Allein, und doch nicht}{32}

\inhaltseintrag{Also hat Gott}{31}

\inhaltseintrag{Amen! Gott Vat’r und}{48}

\inhaltseintrag{Auf, Christenmensch}{38}

\inhaltseintrag{Auf, ihr Christen}{39}

\inhaltseintrag{Auf, ihr Heiden, lobet}{48}

\inhaltseintrag{Auf, Seel’! und danke}{8}

\inhaltsgruppe{C.}

\inhaltseintrag{Christe, wahres Seelenlicht}{9}

\inhaltsgruppe{D.}

\inhaltseintrag{Das walte Gott, der helfen}{54}

\inhaltseintrag{Das walt’ Gott Vater}{9}

\inhaltseintrag{Der lieben Sonnen Licht}{18}

\inhaltseintrag{Der Sabbath ist vergangen}{14}

\inhaltseintrag{Der Tag hat sich geneiget}{17}

\inhaltseintrag{Des Morgens, wann ich früh}{17}

\inhaltseintrag{Die Macht der Wahrheit}{49}

\inhaltseintrag{Du Friedefürst}{59}

\inhaltseintrag{Du liebes Kind}{54}

\inhaltsgruppe{E.}

\inhaltseintrag{Ein Tag geht nach dem andern}{16}

\inhaltseintrag{Es ist noch eine Ruh}{68}

\inhaltsgruppe{F.}

\inhaltseintrag{Fahre fort, Zion}{39}

\inhaltseintrag{Freuet euch, ihr Gotteskinder}{49}

\inhaltseintrag{Frisch auf, mein Seel’}{60}

\inhaltseintrag{Für allen Freuden auf Erden}{44}

\inhaltsgruppe{C.}

\inhaltseintrag{Geh’ aus, mein Herz}{55}

\inhaltseintrag{Gib dich zufrieden}{61}

\inhaltseintrag{Gleichwie sein Haus}{50}

\inhaltseintrag{Gott, der uns diesen Tag}{63}

\inhaltseintrag{Gott Lob! die Woch’ heb’ ich jetzt an}{5}

\inhaltseintrag{Gott Lob! der Sonntag ist herbei}{5}

\inhaltseintrag{Gott Lob! nun geht der Montag an}{5}

\inhaltseintrag{Gott Lob! der Dienstag bringt das Licht}{6}

\inhaltseintrag{Gott Lob! die Mittwoch kömmt herb}{6}

\inhaltseintrag{Gott Lob! heut’ ist der Donnerstag}{6}

\inhaltseintrag{Gott Lob! der Freitag ist herbei}{7}

\inhaltseintrag{Gott Lob! dies ist der Tag der Ruh’}{7}

\inhaltseintrag{Gott Lob! die Woch’ ist auch dahin}{7}

\inhaltseintrag{Gott Lob! der Tag ist nun dahin}{18}

\inhaltseintrag{Gott Lob! der Tag ist nun zum End’}{19}

\inhaltseintrag{Gott Vater, der du deine}{47}

\inhaltseintrag{Groß ist an allen Enden}{50}

\inhaltsgruppe{H.}

\inhaltseintrag{Halte, was du hast}{40}

\inhaltseintrag{Herr, auf dich will ich sehn}{41}

\inhaltseintrag{Herr, du Kraft und Arzt}{45}

\inhaltseintrag{Herr Gott, dich loben wir}{45}

\inhaltseintrag{Herr, ich preise dein Erbarmen}{23}

\inhaltseintrag{Herr Jesu, gib uns Gnad’}{51}

\inhaltseintrag{Hilf Gott, mein Herr}{51}

\inhaltseintrag{Hilf uns, Herr! in allen}{46}

\inhaltseintrag{Hirte deiner Schafe}{16}

\inhaltseintrag{Höchster Priester}{41}

\inhaltsgruppe{J.}

\inhaltseintrag{Jauchze, Eyrus}{52}

\inhaltseintrag{Ich armer Sünder}{70}

\inhaltseintrag{Ich bin getauft auf deinen}{25}

\quellseite{551}

\inhaltseintrag{Ich laß dich nicht}{33}

\inhaltseintrag{Ich weiß ein Blümlein}{26}

\inhaltseintrag{Ich will dich lieben}{42}

\inhaltseintrag{Je größer Kreuz}{64}

\inhaltseintrag{Jesu, frommer Menschenherden}{34}

\inhaltseintrag{Ihr lieben Christen, freu’t}{69}

\inhaltseintrag{In diesen schweren Zeiten}{64}

\inhaltseintrag{In dieser Abendstunde}{20}

\inhaltseintrag{In dieser Morgenstund’}{10}

\inhaltseintrag{In Gottes Namen fahren wir}{56}

\inhaltseintrag{In Gottes Namen geh’ ich hin}{27}

\inhaltsgruppe{K.}

\inhaltseintrag{Kommt und laßt uns Christum}{35}

\inhaltseintrag{König, dem kein König}{35}

\inhaltsgruppe{L.}

\inhaltseintrag{Liebe, die du mich zum}{42}

\inhaltseintrag{Lobet den Herren alle}{11}

\inhaltseintrag{Lobt Gott, ihr Christen}{23}

\inhaltsgruppe{M.}

\inhaltseintrag{Mach’s mit mir, Gott}{71}

\inhaltseintrag{Mein Gott, das Herze}{30}

\inhaltseintrag{Mein Gott! nun ist es wieder}{11}

\inhaltseintrag{Mein’ Seel! dich freu’}{27}

\inhaltseintrag{Mit Gott will ich’s anfangen}{12}

\inhaltsgruppe{N.}

\inhaltseintrag{Nun lasset uns Gott}{47}

\inhaltseintrag{Nun preiset Alle Gottes}{46}

\inhaltseintrag{Nun schlaf, mein liebes}{57}

\inhaltseintrag{Nun tret’ ich wieder aus der Ruh’}{12}

\inhaltsgruppe{D.}

\inhaltseintrag{O Christe, Morgenstern}{13}

\inhaltseintrag{O Gott, du höchster Gnadenhort}{24}

\inhaltseintrag{O Herrscher in dem Himmelszelt}{67}

\inhaltseintrag{O Jesu, einig wahres Haupt}{53}

\inhaltseintrag{O Jesu, Gottes Lämmlein}{71}

\inhaltseintrag{O Jesu, treu’ster Heiland}{15}

\inhaltseintrag{O welch ein unvergleichlich’s}{25}

\inhaltsgruppe{S.}

\inhaltseintrag{Sie ist mir lieb}{53}

\inhaltseintrag{Stell’ dich, Herr, wie du}{65}

\inhaltseintrag{Süßer Jesu! deiner Gnaden}{14}

\inhaltsgruppe{T.}

\inhaltseintrag{Treuer Wächter Israel}{65}

\inhaltsgruppe{U.}

\inhaltseintrag{Unsre müden Augenlider}{20}

\inhaltsgruppe{V.}

\inhaltseintrag{Voller Wunder}{57}

\inhaltsgruppe{W.}

\inhaltseintrag{Walt’s Gott! mein Werk}{21}

\inhaltseintrag{Welch’ eine Sorg’ und}{43}

\inhaltseintrag{Wenn Menschenhülf’}{66}

\inhaltseintrag{Wenn so viel tausend Teufel}{31}

\inhaltseintrag{Wer Jesum bei sich hat}{36}

\inhaltseintrag{Wer ist wohl wie du}{36}

\inhaltseintrag{Wie schön ist’s doch}{58}

\inhaltseintrag{Wir danken dir, Herr Jesu}{15}

\inhaltseintrag{Wird das nicht Freude sein}{71}

\inhaltseintrag{Wohl dem, der Jesum}{24}

\inhaltseintrag{Wo ist Jesus, mein Verlangen}{37}

\end{originalinhalt}

\chapter{Zugabe einiger Melodien}

\quellseite{552}
\notensatz{facsimiles/552}{Noten und Liedtexte im Original: 1. Auf, ihr Heiden, lobet Gott; 2. Christe, wahres Seelenlicht. Mit der Anmerkung zur Melodienauswahl.}

\quellseite{553}
\notensatz{facsimiles/553}{Noten und Liedtexte im Original: 3. Gib dich zufrieden; 4. Höchster Priester, der du dich; 5. Ich laß dich nicht.}

\quellseite{554}
\notensatz{facsimiles/554}{Noten und Liedtexte im Original: 6. O Christ, Morgensterne; 7. Sie ist mir lieb, die werthe Magd; 8. Wer Jesum bei sich hat.}

\quellseite{555}
\notensatz{facsimiles/555}{Noten und Liedtexte im Original: 9. Treuer Wächter Israel; 10. Wird das nicht Freude sein; 11. Wo ist JEsus, mein Verlangen.}
